\documentclass[12pt,a4paper]{article}
\usepackage[utf8]{inputenc}
\usepackage[portuguese]{babel}
\usepackage{hyperref}

\title{Eurovisao Todos}
\author{Tomás Pinto - Scripting no Processamento de Linguagem Natural}
\date{\today}

\begin{document}

\maketitle

\begin{abstract}
Com base no modelo de n-grams desenvolvido e no método de scoring aplicado, as três frases selecionadas do texto são:
\begin{itemize}
    \item Também todos os países de leste e a Tunísia transmitiram o festival.
    \item O país vencedor do festival foi a Suécia, com a canção "Euphoria" interpretada por Loreen.
    \item Esta foi a primeira vez que o concurso foi realizado na Irlanda .
\end{itemize}
\end{abstract}

\section{Texto Original com Entidades (NER)}
Festival Eurovisão da Canção 1956



O \textbf{Festival Eurovisão da Canção 1956} (MISC) (em inglês : \textbf{Eurovision Song Contest 1956} (ORG) e em francês : \textbf{Concours Eurovision} (MISC) de la chanson 1956 ) foi o primeiro \textbf{Festival Eurovisão da Canção} (MISC) e teve lugar no dia 24 de \textbf{Maio de 1956} (MISC) na cidade suíça de \textbf{Lugano} (LOC) . \textbf{Lohengrin Filipello} (PER) foi o apresentador do festival, e \textbf{Fernando Paggi} (PER) o maestro principal de todo o concurso, que foi ganho por \textbf{Lys Assia} (PER) em representação da \textbf{Suíça} (LOC) com a canção " \textbf{Refrain} (PER) " (em português : \textbf{Refrão} (LOC) ). No primeiro de muitos eventos eurovisivos, participaram sete países, entre estes, encontravam-se alguns dos países mais afectados pela recente guerra . No entanto, o número de participantes poderia ter sido maior, chegando aos dez (estes três países que não conseguiram participar neste primeira edição, participaram na segunda edição, em 1957). Nesta primeira edição, a emissão \textbf{Festival Eurovisão da Canção} (MISC) foi sobretudo realizada como se um programa de rádio se tratasse, apesar de existirem algumas câmaras no estúdio para os poucos europeus que tinham televisão naquela altura. A década de 1950 foi marcada por quase toda a \textbf{Europa} (LOC) pela reconstrução, após a \textbf{Segunda Guerra Mundial} (MISC) e ter televisão era um luxo que a grande maioria não podia ter, sendo assim, o festival acabou por ter que ser essencialmente transmitido via rádio, e não televisão, como era o objectivo primário e praticamente único dos inventores e organizadores do concurso. Os cenários pouco tiveram de luxuosos. A sua decoração baseou-se praticamente e quase exclusivamente em arranjos florais. Logo na primeira edição da \textbf{Eurovisão} (LOC) , foi fortemente recomendado que em cada país participante organizasse um festival nacional para eleger o seu representante e respectiva canção. A \textbf{Áustria} (LOC) e a \textbf{Dinamarca} (LOC) foram desqualificados da fase final por se terem inscrito fora do prazo previsto para a entrega das participações. Também se chegou a pensar que o \textbf{Reino Unido} (LOC) também tinha sido desclassificado nas mesmas condições que a \textbf{Áustria} (LOC) e a \textbf{Dinamarca} (LOC), mas em janeiro de 2017 , a \textbf{União Europeia de Radiodifusão} (ORG) (\textbf{UER} (ORG)), revelou que tudo isso não passava de um "mito" criado pelos fãs. A \textbf{UER} (ORG) continuou a explicar que o "\textbf{Festival of British Popular Song} (MISC)" ("\textbf{Festival da Canção Popular Britânica} (MISC)"), um concurso criado pela \textbf{BBC} (ORG) para o \textbf{Reino Unido} (LOC), foi a inspiração que trouxe mudanças no formato do concurso a partir de 1957 . Todos os países participantes enviaram dois membros do júri a \textbf{Lugano} (LOC), excepto o \textbf{Luxemburgo} (PER) que delegou directamente os seus votos à \textbf{Suíça} (LOC). De salientar que neste festival os membros do júri podiam votar no seu próprio país - a primeira e última vez que tal se sucedeu. Este factor, também é tido por muitos, como uma ajuda para explicar o sucesso suíço. Este também foi o único festival em que cada país pode enviar duas canções. A partir de 1957, devido ao acréscimo do número de países participantes, cada país passou a ter direito a apenas a enviar uma única canção. O \textbf{Festival Eurovisão da Canção 1956} (MISC), foi o mais simples e económico de todos os festivais eurovisão, que foram realizados durante todos os anos posteriores a este primeira edição.

Esta primeira edição do certame europeu incluiu práticas que nunca mais se repetiram: duas canções concorrentes por país, voto secreto, dupla votação de um país em nome de outro, voto no próprio país, apenas o título "\textbf{Grand Prix} (MISC)" de recepção para a vencedora, e um único apresentador do sexo masculino.



\textbf{História} (MISC)O festival foi aberto em italiano pelo apresentador \textbf{Lohengrin Filipello} (PER) . O programa era inicialmente conhecido como "\textbf{Eurovision Grand Prix} (MISC)". Este nome, "\textbf{Grand Prix} (MISC)", foi adoptado pelos países francófonos, onde o festival ficou conhecido como "\textbf{Le Grand-Prix Eurovision de la Chanson Européenne} (MISC)". O processo para organizar a primeira edição do Festival Eurovisão da Canção começou dois anos antes da realização do mesmo, em 1954 , quando a \textbf{União Europeia de Rádiodifusão} (ORG) (\textbf{EBU} (MISC)) designou um comité para organizar um programa de televisão emitido ao mesmo tempo (em directo), para todos os países membros da organização. Ao princípio, pensava-se que o festival iria ser utilizado como um espetáculo de variedades, e os italianos chegaram a propor organizar um evento internacional similar ao \textbf{Festival de Sanremo} (PER) , que já organizavam havia alguns anos, no entanto com canção de todos os países.

O comité, dirigido por \textbf{Marcel Bezençon} (PER) , director da televisão suíça, começou em \textbf{Janeiro} (LOC) de 1955 no \textbf{Mónaco} (LOC) a estudar sériamente a idea, e a 19 de \textbf{Outubro} (MISC) do mesmo ano, essa mesma ideia foi proposta em \textbf{Roma} (LOC) . Um periodista britânico foi quem pôs (involuntáriamente), o actual nome ao festival, \textbf{Eurovision Song Contest} (ORG) (traduzido para português \textbf{Festival Eurovisão da Canção} (MISC) ) ao concurso, já que naquela altura não existia nenhum termo \textbf{Eurovision} (MISC) , que actualmente usa-se para se referir à \textbf{EBU} (MISC) e ao festival. A \textbf{EBU} (MISC) traduziu o nome para francês com o nome de \textbf{Grand Prix Eurovision de la Chanson} (MISC) , que foi o utilizado oficialmente no festival. O regulamento começou então a ser escrito já nos finais de 1955 .

Poucos meses depois, deu-se a primeira edição do Festival Eurovisão da Canção, na \textbf{Suíça} (LOC), com sete estreantes (que poderiam ser dez, caso não tivessem estregue as candidaturas tarde de mais).



\textbf{Votação} (MISC)

O sistema de votação do \textbf{Festival Eurovisão da Canção 1956} (MISC) permitiu que \textbf{júris} (PER) pudessem votar em qualquer canção a concurso, incluindo aquelas que estivessem a representar o seu próprio país. Cada país/delegação tinha apenas dois \textbf{júris} (LOC), que atribuíram uma pontuação de 1 a 10 pontos às suas dez canções preferidas. Este sistema nunca mais voltou a ser utilizado em nenhuma edição posterior do festival.

Adicionalmente, o júri suíço estava permitido a votar como parte da delegação belga, outro facto, que nunca mais se repetiu na história do concurso. Este foi outro factor que levou aque outras delegações se queixassem do sistema de votação, que determinou como vencedor o país anfitrião, a \textbf{Suíça} (LOC) (com a cantora \textbf{Lys Assia} (PER) ).

Os resultados deste festival, nunca foram tornados públicos, o que também lançou numerosos rumores sobre outros posicionamentos na tabela de resultados. A segunda canção interpretada por \textbf{Lys Assia} (PER) nunca foi comercializada, o que causou sugestões, sobre a música ter ficado num lugar muito baixo da tabela. Outras sugestões, incluem a possibilidade de a \textbf{Alemanha} (LOC) , com \textbf{Walter Andreas Schwarz} (PER) , ter ficado em segundo lugar (ou \textbf{Freddy Quinn} (ORG) , nalgumas versões do rumor), estas justificadas com o facto de a \textbf{Alemanha} (LOC) ter sido a anfitriã da segunda edição do festival, e a possibilidade de um terceiro lugar de uma das canções provenientes da \textbf{Bélgica} (LOC) .



\textbf{Local} (MISC)O \textbf{Festival Eurovisão da Canção 1956} (MISC) ocorreu em \textbf{Lugano} (LOC) , na \textbf{Suíça} (LOC) . \textbf{Lugano} (LOC) é uma comuna da \textbf{Suíça} (LOC) , no \textbf{Cantão Tessino} (LOC) , do qual é a maior cidade, embora a capital do cantão seja \textbf{Bellinzona} (LOC) . A língua oficial nesta comuna é o italiano . A cidade é banhada pelo \textbf{Lago de Lugano} (LOC) . O nome \textbf{Lugano} (LOC), provém, provavelmente da palavra latina \textbf{Lucus} (PER), que significa madeira ou madeira sagrada. \textbf{Lugano} (LOC) é a cidade mais povoada do \textbf{Cantão Tessino} (LOC) e possui uma área de cerca de 48 km, com uma população actual, superior a cinquenta mil habitantes.

O festival em si, ocorreu no \textbf{Teatro Kursaal} (LOC) , um dos mais importantes e prestigiados locais de espetáculo na altura da realização deste.



\textbf{Visual} (MISC)A primeira edição do festival, não teve uma grande produção visual, em parte por também não ser necessário, visto o festival ter sido transmitido essencialmente via rádio, e também pelo tempo em que decorreu, ter sido pós-guerra, e o festival visava a harmonia entre os países. No entanto, a sala do espetáculo contou com uma pequena audiência, que encheu o local até aos balcões. O palco estava essencialmente decorado com flores e plantas, e o fundo coberto por cortinas que desciam do tecto até ao chão. De um lado, estava reservado o espaço para os artistas cantarem, e do outro estava toda a orquestra, composta por vinte e quatro elementos, dispostos (à excepção do pianista), numa espécie de escadaria, com três degraus. As luzes foram simples, percebendo-se apenas a existência de um holofote para além das luzes normais da sala. A cobertura televisiva, foi efectuada pela televisão do país anfitrião, e através de apenas duas câmaras, uma virada para o público, e outra para o palco, que era virada para um lado e para o outro, para mostrar a orquestra e o artista em palco. No fim, o prémio (um ramo de flores) foi atribuído a \textbf{Lys Assia} (PER) , por um casal de crianças. O logotipo do festival, era apenas a palavra \textbf{Eurovision} (MISC) , escrita em forma de circunferência, com um luz brilhante a sair do topo das letras.



\textbf{Formato} (MISC)A abertura do \textbf{Festival Eurovisão da Canção 1956} (MISC), foi das mais pequenas de sempre (e mais simples também). Tendo durado apenas um minuto e dez segundos, onde a orquestra presente no local do evento tocou um tema com as características que as músicas tinham nos anos 50.

O intervalo esteve a cargo dos \textbf{Les Joyeux Rossignols} (ORG) e dos \textbf{Les Trois Ménestrels} (LOC) , tendo actuado primeiro estes últimos. A segunda actuação, a dos \textbf{Les Joyeux Rossignols} (LOC) teve de ser prolongada devido a atrasos no processo de votação, que mesmo assim apenas acabou por ver revelado o vencedor, e nada mais.



\textbf{Participações individuais} (MISC)Logo na primeira edição, cada país escolheu o seu representante de forma diferente e especifica. Nos artigos da "caixa" em baixo, pode se ler mais sobre o estilo de selecção de cada país, e os concorrentes que concorreram nas selecções nacionais. \textbf{Predefinição:Participações Individuais} (MISC) ESC1956



\textbf{Participações Dos} (MISC) sete países participantes na primeira edição do \textbf{Festival Eurovisão da Canção} (MISC) , apenas \textbf{Suíça} (LOC) e \textbf{Luxemburgo} (LOC) enviaram o mesmo artista, para as suas duas canções representativas. Uma das representantes holandesas, \textbf{Corry Brokken} (PER) , que interpretou o tema " \textbf{Voorgoed Voorbij} (MISC) ", voltou a representar os \textbf{Países Baixos} (LOC) , e venceu o \textbf{Festival Eurovisão da Canção 1957} (MISC) , a segunda edição da \textbf{Eurovisão da Canção} (MISC). Por outro lado, o representante belga , \textbf{Fud Leclerc} (PER) voltou a representar o seu país mais três vezes no \textbf{Festival Eurovisão da Canção 1958} (MISC) , 1960 e 1962 , e por sua vez, \textbf{Lys Assia} (PER) (que cantou com cinco cantores em segundo plano, backsingers ), a primeira representante da \textbf{Suíça} (LOC) , no primeiro evento do festival, realizado na \textbf{Suíça} (LOC), e vencido pela cantora, voltou a representar o seu país, nos dois anos seguintes (em 1957 e 1958 . Sendo assim, foram três os cantores que voltaram a representar o seu país noutra edição do Festival Eurovisão da Canção. Na tabela abaixo, é possível observar as canções (e respectivas traduções) e os artistas representantes escolhidos para cada país, assim como a forma da escolha dos mesmos e a data do respectivo evento para tal selecção.



\textbf{Festival} (PER)No primeiro Festival Eurovisão da Canção , as canções não superaram os três minnutos e meio (não havendo nesta altura qualquer limite para a duração das músicas), e os artistas, foram acompanhados musicalmente por uma orquestra, composta por 24 músicos, dirigidos pelos maestros \textbf{Fernando Paggi} (PER) , \textbf{Jacques Lassry} (PER) , \textbf{Franck Pourcel} (PER) , \textbf{Leo Souris} (PER) e \textbf{Gian Stellari} (PER) . Na primeira participação, não houve nenhum grupo musical, já que este género de arte musical, esteve proibido até aos anos 70. A ordem de actuação foi determinada por sorteio, e os países podiam escolher a música que eles queriam ver cantar primeiro. Cada país pôde nomear seu próprio maestro ou confiar nos serviços do maestro do país anfitrião.

Na tabela abaixo, é possível observar a ordem de entrada em palco das quatorze canções no concurso, representantes de sete países europeus (2 músicas por cada um).

Segundo o regulamento da edição , a orquestra foi composta por 24 músicos, distribuídos por:

Em baixo encontra-se a lista de maestros que conduziram a orquestra, na respectiva actuação de cada país concorrente. Os maestros estão organizados pela ordem de apresentação das músicas de cada país: 1 para a primeira canção interpretada por determinado participante e 2 para a segunda participação desse mesmo participante.



\textbf{Transmissão do Festival} (MISC)O primeiro \textbf{Festival Eurovisão da Canção} (MISC), foi transmitido já via televisão , no entanto, a grande maioria das pessoas acompanhou o evento através de rádios . Isto deveu-se ao facto de naquela altura, a maior parte da população não possuir televisão, devido ao custo destas (de lembrar que o mundo também tinha acabado de sair da sua segunda grande guerra mundial). Hoje em dia, apenas existe o vídeo da actuação vencedora (de \textbf{Lys Assia} (PER) , representante da \textbf{Suíça} (LOC) ), e um registo em áudio, com oitenta minutos de duração (no entanto o festival durou mais vinte minutos, porém essa parte da gravação - a dos votos - perdeu-se ou nunca foi tornada pública).



Festival Eurovisão da Canção 1957



O \textbf{Festival Eurovisão da Canção 1957} (MISC) (em inglês : \textbf{Eurovision Song Contest} (ORG) 1957 , em francês : \textbf{Concours Eurovision} (MISC) de la chanson 1957 e em alemão : \textbf{Großer Preis der Eurovision} (PER) 1957 ) foi o segundo \textbf{Festival Eurovisão da Canção} (MISC) e realizou-se a 3 de \textbf{Março} (PER) de 1957 , em \textbf{Frankfurt am Main} (LOC) . \textbf{Anaïd Iplikjan} (PER) foi a apresentadora do evento, que foi ganho por \textbf{Corry Brokken} (PER) representando os \textbf{Países Baixos} (LOC) com a canção " \textbf{Net Als Toen} (MISC) ". \textbf{Corry} (PER) recebeu trinta e um pontos com a sua música, obtendo assim uma distância de catorze pontos do segundo classificado, e realizando assim a primeira vitória dos \textbf{Países Baixos} (LOC) na \textbf{Eurovisão} (LOC). O segundo festival foi ainda principalmente um programa de rádio, mas houve um notável aumento de pessoas na \textbf{Europa} (LOC) com televisão , o que levou a que o festival fosse mais cuidadoso a nível visual. Se bem que a \textbf{Suíça} (LOC) tivesse ganho o \textbf{Festival Eurovisão da Canção 1956} (MISC) , em 1957 o festival não se realizou na \textbf{Suíça} (LOC) mas na \textbf{Alemanha} (LOC), pois ainda não tinha sido adoptada a regra que o país vencedor organiza o evento no ano seguinte, como acontece actualmente.

Durante algum tempo, existiu o rumor de que a canção alemã "\textbf{Im Wartesaal} (MISC) zum großen Glück" interpretada por \textbf{Walter Andreas Schwarz} (PER) , no \textbf{Festival Eurovisão da Canção 1956} (MISC) , teria ficado em segundo lugar, e era por isso que a \textbf{Alemanha} (LOC) organizou o evento em 1957 (de lembrar que em 1956, apenas o vencedor foi revelado, e nem sequer a pontuação deste foi). De facto, a regra de que o país vencedor deveria organizar o evento eurovisivo do ano seguinte ainda não estava concebida, e estava planeado naquela altura que cada país participante, seria sede do festival através de um sistema "giratório", isto é, um ano num país e noutro, num outro país. No entanto, a medida que mais países participavam no evento, e outros mostravam interesse em participar, este regra tornou-se impraticável.

Neste festival, devido ao facto de a canção italiana ter 5:09 minutos, enquanto que a do \textbf{Reino Unido} (LOC) tinha apenas 1:52 minutos, levou a que fosse implementada uma regra, já no festival do ano seguinte, regra essa que dura até hoje. Cada música apenas passou a poder ter três minutos de duração máxima, com possibilidade de desqualificação se a regra não fosse cumprida. Os cantores dinamarqueses \textbf{Birthe Wilke} (PER) e \textbf{Gustav Winckler} (PER) , depois de acabarem de interpretar a sua canção trocaram o maior beijo da história do festival até hoje - se bem que apenas fosse visto pelas poucas pessoas que tinham televisão. Isto deveu-se ao facto de que o membro do pessoal da produção esqueceu-se de mostrar o pré-aviso, num cartaz (dirigido aos artistas), de que o beijo deveria acabar. No entanto, esta ocorrência não ficou mal vista pelos países \textbf{Europeus} (LOC) a concurso. No festival deste ano, outra novidade também foi o facto de os júris terem sido contactados por telefone. No ano anterior, todas as delegações enviaram os seus jurados a \textbf{Lugano} (LOC) , facto que não se teve que repetir nunca mais na história do festival. Com isto, a edição deste ano, também se tornou na primeira em que o sistema de votação foi público, o que hoje é reconhecido e característico do \textbf{Festival Eurovisão da Canção} (MISC).



\textbf{Local} (PER)

O local destinado para o festival, foram os estúdios da televisão anfitriã alemã, \textbf{Großer Sendesaal des Hessischen Rundfunks} (MISC) , em \textbf{Frankfurt} (LOC) .

O \textbf{Großer Sendesaal des Hessischen Rundfunks} (LOC) é um salão de música e ex-estúdio de televisão, que servia de sede para a \textbf{Hessischer Rundfunk} (ORG) .

Quanto à cidade, é uma cidade independente da \textbf{Alemanha} (LOC) , no estado de \textbf{Hessen} (LOC) . Localiza-se às margens do \textbf{rio Meno} (LOC) ( Main em alemão). \textbf{Frankfurt am Main} (LOC) é actualmente a quinta maior cidade da \textbf{Alemanha} (LOC) e a maior cidade no estado de \textbf{Hessen} (LOC) . \textbf{Frankfurt} (LOC) foi fundada pelos romanos no século I . Desde 1875 é uma metrópole , com mais de 100 mil habitantes.



Visual

O visual/aspecto do \textbf{Festival Eurovisão da Canção 1957} (MISC), foi bastante simples, comparado com os elaborados festivais organizados hoje em dia, e principalmente durante a década de 2000, no entanto, para o ano em questão, 1957, este foi um dos maiores eventos do ano, para os países participantes, especialmente para a \textbf{Alemanha} (LOC), anfitriã do evento. O evento foi realizado num estúdio de televisão da estação alemã, a qual tentou aproveitar ao máximo o espaço disponível, colocando um palco em forma de escada, onde os maestros e os artistas desciam antes da sua actuação. O "pano de fundo" do local de actuação, foi uma grande harpa convenientemente decorada para a ocasião, e o centro tinha um grafismo diferente para cada música. A orquestra acompanhou todos os artistas, do lado direito do local de actuação destes.

Para a votação, foi "reservado" um espaço no estúdio, onde foi colocada uma secretária, com um placar por trás, com o mapa da \textbf{Europa} (LOC), e o nome dos países participantes, onde quando se fazia a ligação telefónica, acendia-se uma luz no respectivo país. Ao lado desse placar encontrava-se outro onde as pontuações eram inseridas. Este quadro possuía o nome das músicas a concurso, e as pontuações eram alteradas manualmente por membros da produção enquanto a votação decorria. O sistema de votação desta edição, foi o mais simples e curto alguma vez realizado na história da \textbf{Eurovisão} (LOC), sem contar com a primeira edição do evento, onde o sistema de votação nunca foi tornado público.



\textbf{Formato} (MISC)A abertura do \textbf{Festival Eurovisão da Canção 1957} (MISC), esteve a cargo da orquestra presente na sala do evento. A \textbf{Tanzorchester des Hessischen Rundfunks} (ORG) , abriu o evento com um tema típico da década de 50.

A apresentadora do festival, \textbf{Anaïd Iplicjian} (MISC) , que utilizou predominantemente alemão para apresentar o espetáculo, foi mais tarde, como que forçada a utilizar inglês e francês na parte da votação. Tal como a maioria dos artistas, \textbf{Anaid} (LOC) vestiu um normal e formal vestido de noite, típico daquela época, e não à semelhança dos dias de hoje, a apresentadora utilizou a mesma indumentária durante todo o festival. A apresentação das canções foi também ela feita com alguma rapidez, com um pequeno espaço temporal entre cada actuação. Por vezes, a orquestra começava a tocar a introdução da música seguinte em palco, ainda antes dos próprios concorrentes se encontrarem em palco. Esta apresentação muito eficiente, é um claro contraste, se compararmos com as edições da \textbf{Eurovisão} (LOC) de hoje em dia, onde cada participação tem trinta segundos para entrar em palco, com um cartão postal a introduzir a sua actuação.



Sistema de votação

Sem a famosa recapitulação das músicas a concurso e um intervalo, os júris internacionais tiveram muito pouco tempo para compilar os seus votos antes de receberem a chamada de \textbf{Anaïd Iplicjian} (MISC) , e da sua assistente telefónica. Cada país tinha um júri constituido por dez membros. Cada júri tinha direito a um voto, que atribuiam à sua actuação preferida, que nunca poderia ser a do seu próprio país. Posteriormente, cada país atribui os seus dez votos aos outros países a concurso, e sem tradução da maioria do sistema de votação, a sequência de votos de 1957 foi a mais curta alguma vez feita na história da \textbf{Eurovisão} (LOC). Os quadros de pontuação não continham os nomes dos países, mas em vez disso, continham os títulos das canções apresentadas naquela noite. O primeiro júri a votar foi o da \textbf{Suíça} (LOC), o qual atribuiu 7 dos seus 10 pontos aos \textbf{Países Baixos} (LOC), e depois da votação de apenas três países, já era óbvio que naquela noite os \textbf{Países Baixos} (LOC) venceriam o certame. Sendo assim, as atenções viraram-se essencialmente para a disputa do segundo lugar.



\textbf{Participações individuais} (MISC)Cada país escolheu o seu representante de forma diferente e especifica. Nos artigos da "caixa" em baixo, pode se ler mais sobre o estilo de selecção de cada país, e os concorrentes que concorreram nas selecções nacionais. \textbf{Predefinição:Participações Individuais} (MISC) ESC1957



\textbf{Participações} (MISC)Em 1956, a \textbf{Áustria} (LOC) , a \textbf{Dinamarca} (LOC) e o \textbf{Reino Unido} (LOC) estrearam-se no Festival Eurovisão da Canção. Estes três países já haviam tentado participar na edição de 1956, no entanto entregaram as candidaturas fora do prazo pré-estipulado pela \textbf{União Rádiodifusora Europeia} (ORG) (\textbf{EBU} (MISC)) , e tiveram de ser desqualificados, tendo que esperar pela edição deste ano. Os três países tornaram-se ao longo dos anos alguns dos países com mais número de participações no festival. Em 1957, a \textbf{Suíça} (LOC) voltou a enviar a mesma representante que em 1956 (assim como faria na próxima edição do festival, em 1958). \textbf{Lys Assia} (PER) , que havia vencido a última edição do festival, e ao mesmo tempo a primeira, voltava a representar pela segunda vez o seu país nos palcos da \textbf{Eurovisão} (LOC), porém não foi tão bem sucedida como na sua estreia. Outro país a repetir o seu representante, foram os \textbf{Países Baixos} (LOC) , que enviaram \textbf{Corry Brokken} (PER) à \textbf{Alemanha} (LOC) para representar o país. A aposta nesta cantora holandesa valeu a vitória aos \textbf{Países Baixos} (LOC) no Festival Eurovisão da Canção 1957. A edição deste ano, foi juntamente com a edição de 1958, a menor em número de músicas apresentadas. Apenas dez concorrentes/países foram à \textbf{Alemanha} (LOC) lutar pelo título europeu. No entanto, se se olhar para o número de países, esta edição (juntamente com a posterior) fica em segundo lugar no que toca a número de países participantes, pois em 1956, foram apenas sete contra os dez países desta edição. Porém, esta não deixou de ser a edição mais curta do festival alguma vez realizada, com apenas 69 minutos.



\textbf{Festival} (PER)O festival foi aberto pela segunda vez consecutiva por uma música em língua holandesa . Durante este festival, foi apresentada a maior e a mais pequena canção alguma vez apresentada e levada à \textbf{Eurovisão} (LOC). A \textbf{Itália} (LOC) apresentou uma música com mais de cinco minutos, e por sua vez, o \textbf{Reino Unido} (LOC) apresentou uma música com pouco mais de um minuto. A música vencedora também exedeu os três minutos (que na altura não eram regra), e consistia numa forte balada, com um solo de violino. A entrada alemã deste ano, foi a primeira a utilizar adereços em palco, tendo levado um telefone para a sua actuação (visto a sua música intitular-se \textbf{Telefone} (MISC), \textbf{Telefone em português} (MISC)). O único duo em 1957, foram os dinamarqueses \textbf{Birthe Wilke} (PER) \& \textbf{Gustav Winckler} (PER), que também exederam o tempo de actuação, porém não por culpa da música que levaram a concurso, mas sim por causa do beijo que trocaram no final da actuação. A fechar o concurso, esteve a \textbf{Suíça} (LOC), com a mesma cantora e o mesmo compositor do ano anterior, esperando assim conseguir mais uma vitória no evento. A sua música estava muito próxima da do ano anterior, de modo a que o mesmo do ano passado se repetisse. Quando os dez países acabaram de actuar, passou-se logo para a votação, que durou pouco mais de doze minutos. Posto isto, e desvendado o vencedor, os \textbf{Países Baixos} (LOC), representados pela cantora \textbf{Corry Brokken} (PER) , voltaram a actuar a música vencedora, e a encerrar a segunda edição do Festival Eurovisão da Canção.

A ordem de votação, foi ao contrário da ordem de actuação dos países no festival. Sendo assim, a ordem de votação no Festival Eurovisão da Canção 1957, foi a seguinte:

Em baixo encontra-se a lista de maestros que conduziram a orquestra, na respectiva actuação de cada país concorrente.



\textbf{Transmissão do Festival} (MISC)Na segunda edição do evento, já havia mais pessoas na \textbf{Europa} (LOC), com televisão próprio, o que levou a um substancial aumento de telespectadores, no entanto, a maioria das pessoas continnuavam sem televisão, pelo que o festival ainda foi em grande parte seguido pela rádio (aparelho muito comum nos lares europeus). Os canais de televisão responsáveis pela difussão do concurso quer via televisão, quer via rádio foram as seguintes cadeias televisivas:



Artistas repetentes

Logo na segunda edição do festival, alguns artistas eram repetentes na sua experiência \textbf{Eurovisiva} (PER). Em 1957, os repetentes foram:



Festival Eurovisão da Canção 1958



\textbf{Países Participantes Países} (LOC) que participaram no passado, mas não em 1958

O \textbf{Festival Eurovisão da Canção 1958} (MISC) (em inglês : \textbf{Eurovision Song Contest 1958} (ORG) , em francês : \textbf{Concours Eurovision} (MISC) de la chanson 1958 ; e em neerlandês : \textbf{Europæiske Melodigrandprix 1958} (MISC) ) foi o 3. \textbf{Festival Eurovisão da Canção} (MISC) e realizou-se em 12 de \textbf{Março de 1958} (MISC) na cidade de \textbf{Hilversum} (LOC) nos \textbf{Países Baixos} (LOC) . A apresentadora do festival foi \textbf{Hannie Lips} (PER).

A edição teve a vitória de \textbf{André Claveau} (PER) , que representou a \textbf{França} (LOC) com a canção " \textbf{Dors} (MISC), mon amour " ( Dorme meu amor ), esta foi a primeira vez que terra francesas venceram o festival. Esta foi a primeira vez na história da \textbf{Eurovisão} (LOC) que o país vencedor num ano organizou o festival no ano seguinte, esta regra tem-se mantido até à actualidade, com algumas excepções. A canção italiana não foi bem ouvida em alguns países, isto significou que depois de todas as outras canções terem sido apresentadas, \textbf{Domenico Modugno} (PER) teve de cantar parcialmente a canção porque os júris não estavam em estúdio como nos anos anteriores. Pela segunda vez, os \textbf{júris} (LOC) estavam nos seus países a assistir ao espectáculo e enviaram os seus resultados através do telefone. A canção italiana foi um grande sucesso nos \textbf{Estados Unidos da América} (LOC) e teve dois \textbf{Prémios Grammy} (MISC) em 1958. O grupo \textbf{Gipsy Kings} (ORG) em 1989 no álbum \textbf{Mosaique} (MISC) adaptou esse grande sucesso. Foi uma das poucas canções \textbf{Eurovisivas} (LOC), que ficou em primeiro lugar nos tops de música na \textbf{América} (LOC). No espetáculo organizado para celebrar os 50 anos da \textbf{Eurovisão} (LOC), em 2005 intitulado de " \textbf{Congratulations} (MISC): 50 \textbf{Anos do Festival Eurovisão da Canção} (LOC) ", a música italiana de 1958 foi escolhida como a 2. melhor música da \textbf{Eurovisão da sempre} (MISC) (atrás da canção dos \textbf{ABBA} (ORG) " \textbf{Waterloo} (MISC) "). Em 1958, foi o último ano e que o \textbf{Reino Unido} (LOC) não participou no festival até aos dias de hoje (conta já com 54 participações no festival). Ao contrário da primeira edição do festival, em 1958 continuou a regra imposta em 1957, onde um país só poderia enviar uma única canção para o concurso. Em 1958 foi adoptada a regra do limite das datas para as canções. Juntamente com o \textbf{Festival Eurovisão da Canção 1956} (MISC), são os únicos festivais na história que não tiveram a concurso uma única canção interpretada em inglês . A única estreia foi a da \textbf{Suécia} (LOC) , um país que ao longo dos anos tornar-se-ia num dos mais bem sucedidos países no concurso. Os júris não se encontravam no estúdio, tal como aconteceu em 1956. À semelhança da edição passada, em 1957, neste ano os júris ficaram e deram as votações a partir do seu próprio país, ouvindo também o festival de lá. Uma vez a apresentação das músicas concluída, os júris enviaram os resultados via telefone . O intervalo esteve a cargo da orquestra presente na sala do evento, liderada pelo maestro

\textbf{Dolf van der Linden} (PER) . Pela primeira vez na histroia do festival, ocorreram dois intervalos, um a meio das actuações e outro no fim de todas as apresentações estarem concluídas. Pelo terceiro ano consecutivo, foi escolhida uma noite diferente para a realização do concurso, e seguida de um domingo em 1956 e uma quinta-feira em 1957, a edição de 1958 teve lugar a uma quarta-feira. Todos os sete países originais do festival, regressaram ao enredo outravez, e a estes juntaram-se a \textbf{Áustria} (LOC) e a \textbf{Dinamarca} (LOC), que obtiveram boas classificações na edição anterior. O único país que não regressou foi o \textbf{Reino Unido} (LOC). Seguido do desapontamento com o resultado obtido no ano anterior, a \textbf{BBC} (ORG) decidiu que não entraria no festival em 1958. No entanto, a \textbf{Suécia} (LOC) estreou-se na \textbf{Eurovisão} (LOC) nesse mesmo ano, e mais uma vez competiram dez países no concurso (o mais baixo número de músicas num festival da \textbf{Eurovisão} (LOC) de sempre).



História

Em 1958, com apenas duas edições até então realizadas, o evento estabeleceu-se como um evento anual de grande sucesso e prestigio, e começou a atrair as atenções de outros países por todo o continente \textbf{Europeu} (LOC) . A única coisa que o festival ainda não tinha conseguido, era produzir um grande sucesso internacional, pois este era conhecido nesta altura como um evento para escolher a melhor canção popular do ano, o que foi um grande desapontamento para os organizadores. Contudo a terceira edição do Festival Eurovisão da Canção estava prestes a mudar isso e da mais forma espectacular possível. O \textbf{Reino Unido} (LOC) originalmente foi o primeiro escolhido para organizar a edição de 1958. No entanto, após o falhanço das negociações com sindicatos artísticos, a \textbf{BBC} (ORG) desistiu no verão de 1957 e não participou pela segunda e última vez na \textbf{História} (MISC). Por isso, para a edição de 1958, um novo formato foi posto em cima da mesa e executado. O país vencedor de um ano, passaria a ter a oportunidade, honra e privilégio de poder organizar o concurso nas suas terras nacionais. Esta nova regra/medida foi aceite com grande honra pelos países participantes, visto que até a própria estação televisiva anfitriã, tinha que pagar altas cotas para poder hostear o concurso. Esta medida do país vencedor organizar o concurso no ano seguinte, continua até aos dias de hoje, com poucas excepções ao longo da história. Seguida da vitória de " \textbf{Net Als Toen} (MISC) ", em 1957, a televisão holandesa \textbf{NTS} (MISC) assumiu a responsabilidade de organizar o evento, no entanto a sua escolha para local do festival não foi vista como a mais óbvia por muita gente que vivia fora dos \textbf{Países Baixos} (LOC) . A cidade rural de \textbf{Hilversum} (PER) , a vinte milhas de \textbf{Amesterdão} (LOC) , é cercada por pântanos, lagos, florestas e outras vilas mais pequenas. Até aos nossos dias, esta foi a cidade mais pequena que alguma vez hosteou o festival. A razão pela qual \textbf{Hilversum} (PER) foi escolhida como sede do evento, foi por ser conhecida como a " \textbf{Cidade dos} (LOC) \textbf{Media} (LOC) " e é o centro tradicional da rádio e televisão holandesa. Hoje em dia, apenas 80 mil pessoas vivem na cidade, mas voltando a 1958, mais de 100 mil viviam no local, graças não só a localização das televisões, mas também por ser um importante centro da agricultura nos \textbf{Países Baixos} (LOC) .



\textbf{Local} (MISC)O \textbf{Festival Eurovisão da Canção 1958} (MISC) ocorreu em \textbf{Hilversum} (PER) , nos \textbf{Países Baixos} (LOC) . \textbf{Hilversum} (PER) é uma municipalidade e cidade dos \textbf{Países Baixos} (LOC) , na província da \textbf{Holanda do Norte} (LOC) . Localizada na região denominada "'t \textbf{Gooi} (LOC)", é a maior cidade daquela área. É cercada por charnecas , florestas , campinas , lagos e vilarejos . \textbf{Hilversum} (PER) faz parte do \textbf{Randstad} (LOC) , uma das maiores conurbações da \textbf{Europa} (LOC) . \textbf{Hilversum} (PER) fica a 30 km a sudeste de \textbf{Amesterdão} (LOC) e 25 km ao norte de \textbf{Utrecht} (LOC) . Na altura do festival, a cidade possuía uma população superior a cem mil residentes. O festival em si, ocorreu nos \textbf{AVRO Studios} (ORG) , os estúdios televisivos situados na \textbf{Cidade de Hilversum} (LOC) . Esta arena foi escolhida para o \textbf{Festival Eurovisão da Canção 1958} (MISC), por ser os estúdios principais da televisão holandesa, tipicamente moderna,, no entanto um edificio um pouco anónimo, que se encontra em uso ainda nos dias correntes. Pelas características do local, um estúdio de televisão, era muito pequeno, e a imagem do palco era bastante similar aos dos anos anterior, utilizando uma escadaria para introduzir os artistas.



\textbf{Formato} (PER) e Visual

O palco do Festival Eurovisão da Canção 1958 foi muito similar aos das edições passadas. Aparte de algumas decorações típicas dos \textbf{Países Baixos} (LOC), o palco estava coberto de flores, o que deu um impacto limitado, ao fazer parecer que o espetáculo era transmitido a preto e branco. O tema floral foi repetido no vestido da apresentadora, \textbf{Hannie Lips} (PER). \textbf{Hannie} (PER) era bem conhecida a nível local e era também uma apresentadora de televisão bem sucedida, era também cantora e acordeonista , no entanto em 1958 \textbf{Lips} (PER) não teve a oportunidade de mostrar os seus dotes artísticos à \textbf{Europa} (LOC). Mais uma vez, o vestido de noite (de gala) foi o adoptado pela grande maioria dos concorrentes, à excepção da entrada \textbf{Alemã} (LOC) e \textbf{Sueca} (ORG) que apresentaram a sua música com roupas um pouco mais curtas que o habitual para os anos 1950 . Tal como em 1957, seguida de uma apresentação muito simples e breve, cada um dos dez concorrentes desceu a escadaria presente no palco e começaram imediatamente a interpretar a sua canção. Os artistas, tiveram pela primeira vez um intervalo um pouco maior entre as actuações em relação aos anos anteriores. Tal acontecimento ajudou também a uma melhor percepção das canções a concurso por parte dos júris internacionais. Outra razão para que a edição 1958 da \textbf{Eurovisão} (LOC) tenha sido maior do que a do ano anterior, foi o facto de ter havido dois intervalos, um a meio das actuações e outro no fim das actuações e antes dos votos. A \textbf{Metropole Orchestra} (ORG) providenciou um pouco de música característica da época, o que ajudar a acalmar a tensão que se fazia sentir em estúdio, e ajudou a estabilizar o intervalo na \textbf{Eurovisão} (LOC), que se tornou um símbolo e umas das partes importantes do festival desde então.

A abertura do \textbf{Festival Eurovisão da Canção 1958} (MISC) voltou a ser uma abertura simples e curta, executada pela orquestra presente na sala.

Durante o intervalo houve música da \textbf{Metropole Orkest} (PER) , com direcção de orquestra por \textbf{Dolf van der Linden} (PER) . Em 1958 houve dois intervalos: um no meio e outro depois de as canções terem sido interpretadas.



\textbf{Votação Seguida} (MISC) da segunda actuação da música italiana, devido a problemas técnicos e de um curto intervalo, iniciou-se a sequência de voto. Mais uma vez os países votaram pela ordem contrária à da sua actuação, algo que foi feito aparentemente para acelarar o espetáculo, pois o último país a actuar não podia (nem pode) atribuir pontos a si mesmo, sendo assim poderiam acabar a sua votação a tempo de serem os primeiros a serem chamados. O sistema de votação foi o mesmo utilizado no \textbf{Festival Eurovisão da Canção 1957} (MISC), em que cada país tinha um júri composto por dez membros, e cada um desses dez jurados atribuiam um ponto à sua canção preferida. A votação de 1958 pode ser descrita como uma votação onde ocorreram vários "altos e baixos" para vários países. A estreante no festival \textbf{Suécia} (LOC) liderou o início da votação, no entanto foi ultrapassada depois de a \textbf{Áustria} (LOC) atribuir sete dos seus dez pontos à \textbf{França} (LOC), o que fez com que a \textbf{França} (LOC) passa-se para a frente. Depois de quatro júris terem votado, a \textbf{França} (LOC) e a \textbf{Itália} (LOC) estavam empatadas em primeiro lugar, no entanto, o júri que se seguiu, da \textbf{Dinamarca} (LOC) atribui-o nove pontos à \textbf{França} (LOC), e esta passou novamente a liderar, na altura em que já se sabia que a canção " \textbf{Dors} (MISC), mon amour " era a grande candidata a tentar vencer pelos outros concorrentes. A entrada \textbf{Italiana} (LOC) que até aqui se encontrava em segundo, desceu algumas posições durante a segunda parte da votação. Com apenas um júri em falta, o resultado era altamente delicado. A \textbf{França} (LOC) tinha então vinte e um pontos, apenas mais um do que a \textbf{Suíça} (LOC). O último país a votar foi a \textbf{Itália} (LOC), e visto a \textbf{Itália} (LOC) ser vizinha da \textbf{Suíça} (LOC) e com melhores relacionamentos com esta do que com a \textbf{França} (LOC), muitos pensaram que a \textbf{Suíça} (LOC) já tivesse vencido o evento. No entanto, e ignorando todos os outros países, a \textbf{Itália} (LOC) deu seis pontos à \textbf{França} (LOC) e apenas quatro à \textbf{Suíça} (LOC), dando assim à \textbf{França} (LOC) a sua primeira vitória no \textbf{Festival Eurovisão da Canção} (MISC). O grande desapontamento da noite (outra novidade da edição de 1958), foi a participante dos \textbf{Países Baixos} (LOC), \textbf{Corry Brokken} (PER), que apenas conseguiu um ponto, partilhando o último lugar com o \textbf{Luxemburgo} (LOC).



\textbf{Participações individuais} (MISC)Cada país a concurso, executou a escolha do seu representante de maneira diferente. Para características específicas sobre a entrada de cada participante no festival, seguir a ligação com o nome dos países na caixa em baixo. \textbf{Predefinição:Participações Individuais} (MISC) ESC1958



\textbf{Participantes} (MISC)Em 1958, participaram apenas dez países no \textbf{Festival Eurovisão da Canção} (MISC), naquela que foi a terceira edição do evento. Este seria, juntamente com a segunda edição do evento , o espetáculo com menos canções a concurso (não de países pois em 1956 concorreram apenas sete países, no entanto cada país a concurso levou nesse ano duas músicas a \textbf{Lugano} (LOC)). No ano anterior o \textbf{Reino Unido} (LOC) estreou-se na \textbf{Eurovisão} (LOC) (após ter tentado fazê-lo na edição antes, porém já entregou a sua participação demasiado tarde, a ponto de esta ser recusada pela \textbf{União Europeia de Rádiodifusão} (ORG) ) no entanto não obteve os resultados que espera e cria, o que levou a \textbf{BBC} (ORG) a desistir e retirar-se do \textbf{Festival} (LOC) em 1958, voltando apenas em 1959 (esta foi a única vez na história do evento, que o \textbf{Reino Unido} (LOC) "perdeu" uma edição do concurso). Com a saída do \textbf{Reino Unido} (LOC) ficaram apenas nove países, no entanto outro país do norte europeu quis entrar na competição. Tratava-se da \textbf{Suécia} (LOC) , que viria ao longo da história a ter uma das mais bem sucedidas participações na \textbf{Eurovisão} (LOC). Sendo assim, a terceira edição do Festival Eurovisão da Canção voltava a ter dez nações participantes. Sete dos dez países a concurso organizaram um festival nacional para escolherem o seu representante e mais uma vez, \textbf{Itália} (LOC) enviou o vencedor do \textbf{Festival de San Remo} (PER) . A edição de 1958, foi também a primeira (e até aos dias de hoje a única), onde dois vencedores de festivais anteriores concorreram. Pela \textbf{Suíça} (LOC), \textbf{Lys Assia} (PER), vencedora da primeira edição em 1956, e pelos \textbf{Países Baixos Corry Brokken} (LOC), vencedora da edição 1957. Em 1958 foi a primeira vez que se começou a ouvir avaliações e julgamentos sobre uma música antes do festival (o que hoje em dia é muito comum, até meses antes do festival). A canção italiana " \textbf{Nel} (ORG) blu dipinto di blu ", co-escrita por \textbf{Domenico Modugno} (PER) , que já havia atingido os tops italianos e já era um pouco conhecida pelos países vizinhos da \textbf{Itália} (LOC), era mencionada como aquela a não perder no festival, e talvez mesmo fosse a vencedora.



\textbf{Festival} (PER)Embora mais uma vez o concurso de 1958 possa não ter refletido a música que dominou os charts da época, certamente tinha mais variedade que o concurso do ano anterior, e várias das canções populares apresentadas continuam ainda a ser conhecidas por muitos. A música de abertura e a favorita foi a canção italiana " \textbf{Nel} (ORG) blu dipinto di blu " (\textbf{Azul} (PER) pintado em azul), realizada e co-escrita por \textbf{Domenico Modugno} (PER) . Apesar dos quatro minutos de duração, o desempenho da música foi sempre dramático, ajudado por movimentos expressivos da mão de \textbf{Domenico} (PER). A canção que ganhou o \textbf{Festival de San Remo} (MISC) um mês antes, havia-se tornado já um enorme sucesso na \textbf{Itália} (LOC). "Nel blu dipinto di blu" tornar-se-ia mais conhecida como "\textbf{Volare} (PER)" (\textbf{Voo} (PER)) e no ano seguinte tornaria-se um dos maiores sucessos na história do \textbf{Festival Eurovisão da Canção} (MISC), ficando em primeiro nos tops musicais de toda a \textbf{Europa} (LOC) e chegar a número um no \textbf{Top} (ORG) americano 100. Em 2005, os telespectadores de toda a \textbf{Europa} (LOC) viriam a selecioná-la como a segunda melhor canção alguma vez apresentada no \textbf{Festival Eurovisão da Canção} (MISC) (derrotada pelos \textbf{ABBA} (ORG) , com " \textbf{Waterloo} (MISC) "), no especial \textbf{Congratulations} (MISC): 50 \textbf{Anos do Festival} (LOC) Eurovisão da Canção . A segunda canção a ser apresentada em 1958 foi a entrada do anfitrião, e os holandeses decidiram repetir a mesma fórmula que lhes havia dado sucesso no ano anterior. Pelo terceiro ano consecutivo, e doze meses após dar aos \textbf{Países Baixos} (LOC) a sua primeira vitória no \textbf{Festival} (LOC), \textbf{Corry Brokken} (PER) voltaria a representar o país. A sua canção " \textbf{Heel de wereld} (MISC) " (O mundo inteiro) foi, porém, uma pálida imitação da canção de 1957, e apesar de usar uma roupa muito semelhante, e dando uma outra boa actuação, ficou claro que esta não seria como a outra entrada que havia apresentado no ano anterior. De seguida foi a vez de se ouvir o francês em palco, com a música " \textbf{Dors Mon Amour} (MISC) " (Dorme meu amor), cantada por \textbf{André Claveau} (PER) . A música era uma balada romântica tipicamente francesa e incluía uma introdução alargada, combinando com estilo descontraído de \textbf{Claveau} (PER) e como muitos descreveram, uma voz quente , que o tornou num concorrente sério. A entrada do \textbf{Luxemburgo} (LOC), que se seguiu à da \textbf{França} (LOC), foi em um estilo similar, mas sem uma forte melodia. A canção foi avaliada por muitos como já desgastada , devido ao facto de ser do mesmo estilo que tanto se ouviu durante vários anos na \textbf{Europa} (LOC), no entanto ajudou a tornar a apresentação da \textbf{França} (LOC) melhor. A estreia da \textbf{Suécia} (LOC) na \textbf{Eurovisão} (LOC) foi memorável em dois aspectos, a cantora \textbf{Alice Babs} (PER) preferiu ignorar os típicos vestidos de noite que as outras concorrentes também usavam, e resolveu utilizar um típico traje sueco, que a ajudou a ultrapassar o problema de cantar numa nova língua no festival, a abertura da música tornar-se-ia então memorável na história do concurso, utilizando o universal "la-la-la" para os primeiros 30 segundos da actuação. A canção " \textbf{Lilla} (PER) stjärna " foi, porém, muito doce e encantadoramente executada, como relatam os registos da altura. Os vizinhos escandinavos foram os próximos a actuar, mas, apesar de um belo truque de adoptar o título de "dagbog \textbf{Jeg} (PER) rev et blad ud af min" (Eu arranquei uma página do meu diário), pouco mais sobre a música foi especialmente memorável.

A belga \textbf{Fud Leclerc} (PER) , que tinha competido no primeiro concurso em 1956 retornou em 1958 com um estilo muito diferente de música, dois anos depois. O jazz chamado de " \textbf{Ma} (LOC) Petite Chatte " (Meu pequeno doce) foi uma das melhores performances da noite. Não pela última vez, a entrada alemã baseou-se fortemente em chamarizes para grande parte do seu recurso. A letra da canção fora completamente alterada em relação à versão de estúdio e à apresentada no país, para permitir à cantora \textbf{Margot Hielscher} (PER) vestir-se como "\textbf{Miss Alemanha} (MISC)", com tiara e faixa de miss, para executar/interpretar " \textbf{Für zwei Groschen Musik} (ORG) " (Música cachola). No entanto isso não foi o único artifício inventado pela \textbf{Alemanha} (LOC); \textbf{Margot} (LOC) levou três 7" singles para palco com ela, para não deixar morrer as letras das canções que mencionavam três estilos musicais (jazz, folk e ópera). O desempenho foi certamente memorável, e a canção bastante aplaudida, no entanto não arrecadou uma boa posição na final, e pensou-se na altura que tivesse sido um pouco diferente para os júris internacionais. De seguida a vizinha \textbf{Áustria} (LOC) tentou algo um pouco mais tradicional, com um desempenho mais convencional, executado por \textbf{Liane Auguestine} (PER) , com " \textbf{Die ganze} (MISC) \textbf{Welt} (PER) braucht \textbf{Liebe} (PER) " (O mundo inteiro precisa de amor), uma canção que devia mais do que um pouco aos dois primeiros vencedores da \textbf{Eurovisão} (LOC), tanto em termos de estilo e desempenho. Pelo segundo ano consecutivo a \textbf{Suíça} (LOC) teve a canção (lugar de actuação) final na competição, pelo terceiro ano consecutivo também, \textbf{Lys Assia} (PER) representou a \textbf{Suíça} (LOC). A canção " \textbf{Giorgio} (MISC) " foi um estilo muito diferente dos seus esforços anteriores. Cantada, principalmente, em italiano, \textbf{Giorgio} (LOC) é um conto de charme desenvolto num fim de semana pelo \textbf{lago Maggiore} (LOC), onde \textbf{Lys} (PER) e seu parceiro, jantaram e onde ela se recusou a beber café expresso em troca de vinho. \textbf{Assia} (PER) teve um bom desempenho ao interpretar a canção em questão, e até hoje muitas pessoas consideram que esta foi a melhor das canções que levou à \textbf{Eurovisão} (LOC). Em baixo fica a tabela com as principais e mais gerais características de cada nação a concurso.

A ordem de votação, foi ao contrário da ordem de actuação dos países no festival. Sendo assim, a ordem de votação no Festival Eurovisão da Canção 1958, foi a seguinte:

Em baixo encontra-se a lista de maestros que conduziram a orquestra, na respectiva actuação de cada país concorrente.



\textbf{Artistas Repetentes} (PER)

Com apenas três edições, o \textbf{Festival Eurovisão da Canção 1958} (MISC) viu vários dos seus artistas passados regressarem ao concurso. Em baixo fica uma tabela com esses mesmos artistas repetentes, incluindo o seu desempenho e principais características noutros festivais.



Transmissão

Os canais de televisão responsáveis pela difussão do concurso quer via televisão, quer via rádio foram as seguintes cadeias televisivas:



\textbf{Problemas} (PER)

Enquanto o espetáculo pareceu bastante atraente nos ecrãs de televisão pela \textbf{Europa} (LOC) fora, a única coisa que o concurso de 1958, não consegui-o foi qualquer senso de atmosfera. O pequeno público presente no estúdio é uma das mais baixas audiências na história da competição e é impossível avaliar quais as músicas que foram bem recebidas e que não o tenham sido, embora que, as boas-vindas para a entrada da representante do país anfitrião tenha sido bem mais audível. Mesmo a sequência de votação, que foi muito mais emocionante do que no ano anterior, parece que não consegui-o gerar uma reação maior, talvez por causa da falha da entrada da \textbf{Suíça} (LOC), país afitrião. O outro problema com o espetáculo, é um obstáculo técnico que significou/resultou que o desempenho da música de abertura da \textbf{Itália} (LOC) não fosse visto em vários países e teve de ser repetida no final da ordem de execução. Globalmente, o \textbf{Festival Eurovisão da Canção 1958} (MISC) é uma organização eficiente, mas tem uma apresentação um pouco maçante, assim descrito pelos da época. No entanto, quando se pensa nos concursos anteriores, é no padrão geral da música, uma das mais elevadas edições nos primeiros anos da competição.



Festival Eurovisão da Canção 1959



\textbf{Países Participantes Países} (LOC) que participaram no passado, mas não em 1959

O \textbf{Festival Eurovisão da Canção 1959} (MISC) (em inglês : \textbf{Eurovision Song Contest} (ORG) 1959 e em francês : \textbf{Concours Eurovision} (MISC) de la chanson 1959 ) foi o 4 \textbf{Festival Eurovisão da Canção} (MISC) e realizou-se no dia 11 de março de 1959 na cidade da \textbf{Riviera Francesa} (LOC) de \textbf{Cannes} (LOC) . \textbf{Jacqueline Joubert} (PER) foi a apresentadora do festival que foi ganho pela holandesa \textbf{Teddy Scholten} (PER) que representou os \textbf{Países Baixos} (LOC) com a canção " \textbf{Een beetje} (MISC) " (Um pouco), dando assim, a segunda vitória aos \textbf{Países Baixos} (LOC), o primeiro país a ganhar duas vezes o concurso e também o único a ganhar mais do que uma vez na década de 1950 . Em segundo lugar posicionaram-se os britânicos \textbf{Pearl Carr} (ORG) \& \textbf{Teddy Johnson} (PER) , com a canção " \textbf{Sing Little Birdie} (MISC) ", que obteve 16 pontos, dando assim, o primeiro dos 15 2s lugares que o \textbf{Reino Unido} (LOC) viria a ter e também a melhor classificação até 1967 . Em terceiro lugar, com 15 pontos, foi a canção francesa " \textbf{Oui oui oui oui} (MISC) ", cantada por \textbf{Jean Philippe} (PER) . Estas três canções foram interpretadas no final do festival, sendo a única vez que tal aconteceu, que durou 73 minutos. A letra da canção vencedora foi escrita por \textbf{Willy van Hemert} (PER) , o qual também escreveu a letra da canção holandesa de 1957. Uma nova regra foi criada na \textbf{Eurovisão} (LOC) obrigou que apenas editores ou compositores profissionais podiam fazer parte dos júris nacionais. Este festival foi a primeira retirada de \textbf{Luxemburgo} (PER), depois de terminar com um ponto no ano anterior, a estreia do \textbf{Mónaco} (LOC) e o retorno do \textbf{Reino Unido} (LOC) após sua ausência em 1958, assim que a \textbf{UER} (ORG) alterou as regras de concurso para admitir 11 participantes. O representante monegasco, \textbf{Jacques Pills} (PER) , terminou em último, com um ponto, mas sua filha \textbf{Jacqueline Boyer} (PER) venceu o festival no ano seguinte, pela \textbf{França} (LOC). Na imprensa, foi sugerido que \textbf{Itália} (LOC) e a \textbf{França} (LOC) deram mais pontos á canção holandesa porque nenhum destes queria que o outro obtivesse a vitória final.



\textbf{Local} (MISC)O \textbf{Festival Eurovisão da Canção 1959} (MISC) ocorreu em \textbf{Cannes} (LOC) , na \textbf{França} (LOC) . A estação emissora francesa \textbf{RTF} (MISC) decidiu organizar a 4 edição do concurso em \textbf{Cannes} (LOC), em vez de \textbf{Paris} (LOC), o que não foi como uma grande surpresa por parte dos franceses. \textbf{Cannes} (LOC) já era bem utilizada para encenar grandes eventos internacionais, incluindo o mais famoso, festival de cinema anual de \textbf{Cannes} (LOC) , que tinha sido executado desde 1946. A televisão francesa tinha um estúdio local, com base nas proximidades e a combinação de um local prontamente disponível e do clima mediterrânico fez a centro de vela popular de \textbf{Cannes} (LOC) o natural favorito para sediar a competição, apesar de sua população relativamente pequena de pouco mais de 50.000. A escolha do local também foi nenhuma surpresa. O famoso \textbf{Palais des Festivals} (MISC) (à esquerda), no coração da cidade, era na época o maior local para organizar o concurso. O original \textbf{Palais des Festivals} (MISC) foi construído em 1949 para sediar o festival de cinema. O edifício era localizado na famosa avenida de \textbf{La Croisette} (LOC), no local atual do \textbf{Hotel Noga Hilton} (LOC). Apesar de também encenar o \textbf{Festival Eurovisão da Canção 1961} (MISC) , hoje o antigo \textbf{Palais des Festivals} (MISC) já não existe. Em resposta ao crescente sucesso do festival de cinema e o advento de várias convenções de grandes empresas, a cidade de \textbf{Cannes} (LOC) decidiu construir um novo e maior espaço em 1979. O festival de cinema mudou-se para o novo \textbf{Palais des Festivals et des Congrès} (MISC) em 1983 e logo após o antigo \textbf{Palais des Festivals} (MISC) foi destruído. Apesar de sua maior capacidade e conceito moderno, o design do atual \textbf{Palais des Festivals} (MISC) foi originalmente condenado por muitos como semelhante a um bunker gigante, em comparação com o slendour do original \textbf{Palais des Festivals} (MISC).



\textbf{Formato} (PER) e Visual

A única coisa que atingiu a todos assistir o \textbf{Festival Eurovisão da Canção 1959} (MISC), é a encenação muito elaborada e atractiva do espectáculo, ao contrário do que se passou nos dois concursos anteriores. A apresentadora do espectáculo \textbf{Jacqueline Joubert} (PER) (abaixo) foi uma das apresentadoras de televisão mais experientes da \textbf{Europa} (LOC), tendo sido a primeira locutora de continuidade quando a TV francesa foi criada em 1949. \textbf{Joubert} (PER) utilizou o francês durante todo o show, só falando em inglês por parte do votação. A antiga atriz de teatro e cinema, \textbf{Joubert} (PER) (mãe do apresentadpr "\textbf{Eurotrash" Antoine de Caunes} (MISC)) começou o festival com uma introdução individual de cada concorrente das 11 nações participantes, enquanto usava um ponteiro gigante para indicar a posição de cada país no placar das votações. Cada interprete apareceu numa das três tendas giratórias, cada um com um cenário combinado para o país que representavam. Embora esta foi uma introdução ao invés prolixo, foi certamente visualmente atraente e bastante teatral. A apresentação de todos os cantores como a introdução para o concurso, foi feito em várias ocasiões, tendo a última sido na edição de 2018 . Após a introdução inicial, \textbf{Joubert} (PER) também apresentou cada cantor antes da sua apresentação. Estes desempenhos foram, mais uma vez feitos a partir das tendas rotativas. A apresentação das músicas era rápida e a recepção do público foi certamente mais entusiasmante do que nas edições anteriores. \textbf{Frank Pourcell} (PER) tinha conduzido o vencedor da \textbf{Eurovisão 1958} (MISC) e foi sua orquestra que forneceu o acompanhamento musical para esta edição. Curiosamente seis dos onze canções foram conduzidas por \textbf{Pourcell} (MISC), um recorde para um concurso. Apesar da apresentação de abertura longa e das apresentações individuais, o \textbf{Festival Eurovisão da Canção 1959} (MISC) durou apenas 73 minutos, ajudado por não ter tido nenhum intervalo, e uma sequência de votação muito suave. Mais uma vez, o padrão geral das canções foi bastante elevado com várias novidades recebendo um passeio, com resultados muito diferentes. O cenário era composto por três tendas giratórias pequenas, de onde vieram os cantores, que cantaram com fotografias típicas do país de fundo.



\textbf{Votação} (MISC)Cada país tinha um júri constituído por 10 membro e cada um deles dava 1 ponto à sua canção preferida. Por alguma razão, uma nova regra foi criada para este \textbf{Festival Eurovisão da Canção} (MISC), garantindo que não há editores profissionais ou compositores fossem autorizados a estar nos júris nacionais. Essa regra não durou por muito tempo. No entanto, é muito improvável que esta mudança de regra tinha nada a ver com o facto de que esta foi uma das sequências de voto mais renhidas e mais emocionantes na primeira década do concurso. A apresentadora \textbf{Jacqueline Joubert} (PER) adotou uma abordagem bastante formal para o processo de votação, usando um ponteiro grande para mostrar onde o placar estava a vir. Mais uma vez, os países considerados em ordem inversa á sua interpretação e cada país tinha dez jurados, cada um dos quais atribuído um ponto para sua música favorita. \textbf{Bélgica} (LOC) deu uma vantagem esmagadora aos seus vizinhos holandeses, mas depois de quatro países terem votado, os \textbf{Países Baixos} (LOC) foi apenas um dos três países que estavam empatados; com a \textbf{Suíça} (LOC) eo \textbf{Reino Unido} (LOC) que tinha uma ligeira vantagem, tendo já votado. Quando o júri holandês entregou ao \textbf{Reino Unido} (LOC) cinco dos seus dez pontos, o que pareceu tornar-se decisivo. No entanto, foi o júri italiano o centro das atenções com os votos críticos, quando concedeu sete pontos para os \textbf{Países Baixos} (LOC). No entanto, após não ter sido pontuado pelo júri dinamarquês (o penúltimo a votar), os \textbf{Países Baixos} (LOC) era apenas um dos quatro países que ainda poderia ganhar, dependendo dos votos do júri da \textbf{França} (LOC). Quando o juri francês concedeu quatro pontos para a canção holandesa, ficou claro que os \textbf{Países Baixos} (LOC) estava para vencer pela segunda vez o \textbf{Festival Eurovisão da Canção} (MISC) em três anos. Pelo segundo ano consecutivo a participação italiana foi a grande decepção, terminando em 6 lugar. Excepcionalmente, os países constituintes do top3, voltaram a interpretar as suas canções, dando ao público uma segunda oportunidade para ouvir a participação de casa e o vice-campeão do \textbf{Reino Unido} (LOC), bem como a vencedora holandesa \textbf{Teddy Scholten} (PER) .



\textbf{Participações individuais} (MISC)Cada país a concurso, executou a escolha do seu representante de maneira diferente, tendo sido marcado pela final nacional sueca (o \textbf{Melodifestivalen} (ORG) ), que foi considerado um pouco unusual, por ter sido vencido um interprete e para o festival ter sido escolhido outro. Para características específicas sobre a entrada de cada participante no festival, seguir a ligação com o nome dos países na caixa em baixo. \textbf{Predefinição:Participações Individuais} (MISC) ESC1959



\textbf{Participantes} (MISC)Bem como ter o maior local e audiência até à data, o \textbf{Festival Eurovisão da Canção 1959} (MISC) também teve o maior número de países participantes até à data. Depois de ter falhado a competição de 1958, o \textbf{Reino Unido} (LOC) voltou, e aproveitando a competição que estava sendo encenado em sua porta, o pequeno principado de \textbf{Mônaco} (LOC) fez sua estreia na \textbf{Eurovisão} (LOC). A única surpresa foi a decisão do \textbf{Luxemburgo} (LOC) de se retirar, na sequência do resultado decepcionante no ano anterior. \textbf{Luxemburgo} (PER), assim, tornou-se a primeira das sete nações originais da \textbf{Eurovisão} (LOC), a perder uma edição. Pela primeira vez, a \textbf{Suíça} (LOC) não foi representada pela vencedora de 1956 \textbf{Lys Assia} (PER) e os \textbf{Países Baixos} (LOC) não foram representados pela vencedora de 1957 \textbf{Corry Brokken} (PER) , que havia perdido um recorde de quarta particição na \textbf{Eurovisão} (LOC) no final nacional holandesa. Dois veteranos da \textbf{Eurovisão} (LOC) fizeram no entanto um retorno, de 1957 a cantora dinamarquesa \textbf{Birthe Wilke} (PER) , desta vez como solista e o italiano \textbf{Domenico Modugno} (PER) , pondo a deceção acontecida no ano anterior para trás, para tentar por uma segunda vez. No ano de intervir, \textbf{Modugno} (LOC) tinha-se tornado uma das maiores estrelas da música internacional, graças ao sucesso de "\textbf{Nel} (ORG) blu dipinto di blu (\textbf{Volare} (PER))" e mais uma vez a \textbf{Itália} (LOC) começou como o favorito na pré-competição, com nação anfitriã \textbf{França} (LOC) também altamente classificada.

A música alemã foi originalmente intitulada "\textbf{Heut 'Möcht' Ich Bummeln} (PER)" ( Hoje eu quero dar um passeio' ). As irmãs \textbf{Kessler} (PER), as representantes alemãs, pediram para modificar algumas palavras. O título então se tornou "\textbf{Heute} (LOC) abend woll'n wir tanzen geh'n" ( Hoje vamos dançar ).



\textbf{Festival} (PER)Pela primeira vez, o país anfitrião foi o responsável por abrir o desfile das canções. Tendo ganho no ano anterior com uma balada, a \textbf{França} (LOC) apresentou-se em 1959 com uma canção latin-tempo " \textbf{Oui} (MISC), oui, \textbf{oui} (MISC), oui " (\textbf{Sim} (ORG), sim, sim, sim). A canção interpretada de forma enérgica por \textbf{Jacques Philippe} (PER) , incorporou duas ideias que estariam na base da chamada "canção eurovisiva": um refrão orelhado e referência a lugares (neste caso com \textbf{Capri} (LOC) , \textbf{Bengala} (LOC) e \textbf{Havai} (LOC) ). Seguiu-se \textbf{Birthe Wilke} (PER) , que, após representar a \textbf{Dinamarca} (LOC) no \textbf{Festival Eurovisão da Canção 1957} (MISC) com uma música de amor em dueto, regressou com a música up-tempo " \textbf{Uh} (MISC), \textbf{jeg ville ønske jeg var dig} (PER) " (\textbf{Oh, Quem} (MISC) me dera ser tu) interpretada com charme e uma forte orquestração. Pelo segundo ano consecutivo, \textbf{Domenico Modugno} (PER) voltou a representar a \textbf{Itália} (LOC), desta vez com " \textbf{Piove} (ORG) (\textbf{Ciao} (LOC), ciao bambina) " (Está a chover (\textbf{Adeus} (PER), adeus menina)), música dramática e não usual que se tornaria outro sucesso de \textbf{Modugno} (LOC). Porém, não conseguiu atingir o sucesso de " \textbf{Nel} (ORG) blu dipinto di blu ". Após três músicas fortes, chegou a vez da estreia do \textbf{Mónaco} (LOC), representada pela canção " \textbf{Mon ami Pierrot} (MISC) " (Meu amigo \textbf{Pierrot} (MISC)), interpretada por \textbf{Jacques Pils} (PER) , cuja filha, \textbf{Jacqueline Boyer} (PER) , tornar-se-ia um ano mais tarde a vencedora do \textbf{Festival Eurovisão da Canção} (MISC) , pela \textbf{França} (LOC). A quinta canção a subir ao palco foi a canção representante dos \textbf{Países Baixos} (LOC), interpretada por \textbf{Teddy Scholten} (PER) . " Een beetje " (Um pouco) uma canção com toques de up-tempo e com um refrão cativante, tornar-se-ia a canção vencedora da edição. A participação alemã traria à \textbf{Eurovisão} (LOC) a primeira "família", as irmãs gémeas \textbf{Alice \& Ellen Kessler} (ORG) , que interpretaram a canção swing-tempo " \textbf{Heute Abend} (PER) wollen wir tanzen geh'n " (Hoje à noite queremos ir dançar) com o refrão cativante "hallo boys".

Com mais de metade das canções já interpretadas, a \textbf{Suécia} (LOC) trouxe outra música up-temp ao certame. Brita \textbf{Borg} (PER) usou várias expressões faciais para contar a história de " \textbf{Augustin} (MISC) ", um rapaz mulherengo. Em contraste, \textbf{Christa Williams} (PER) apresentou a balada operática " \textbf{Irgendwoher} (MISC) " (De algures). A canção austríaca " Der K und K \textbf{Kalypso aus Wien} (MISC) " (O \textbf{K} (MISC) e \textbf{K Calipso de Viena} (ORG)) cantado por \textbf{Ferry Graf} (PER) , protagonizou uma das atuações mais bizarras da história do \textbf{Eurovisão} (LOC), combinando um calipso caribenho com iodelei , truques orquestrais e nenhum refrão óbvio. O \textbf{Reino Unido} (LOC) regressou ao certame com o casal \textbf{Pearl Carr} (ORG) \& \textbf{Teddy Johnson} (PER) , cuja atuação da canção " \textbf{Sing} (PER), \textbf{Little Birdie} (PER) " (Canta, \textbf{Passarinho} (PER)), incluiu um pássaro de brincar. Apesar da atuação pouco usual, tornou-se uma das fortes candidatas à vitória. Em contrapartida, a canção que encerrou o desfile, " \textbf{Hou} (MISC) toch van mij " (Por favor ama-me), representando a \textbf{Bélgica} (LOC), e interpretado por \textbf{Bob Benny} (PER) , foi mais convencional, aproximando-se à música vencedora do ano anterior.

A ordem de votação, foi ao contrário da ordem de actuação dos países no festival. Sendo assim, a ordem de votação no Festival Eurovisão da Canção 1959, foi a seguinte:



\textbf{Artistas Repetentes} (PER)

Já na quarta edição do festival, alguns artistas eram repetentes na sua experiência \textbf{Eurovisiva} (PER). Em 1959, os repetentes foram:



Transmissão

Os canais de televisão responsáveis pela difusão do concurso quer via televisão, quer via rádio foram as seguintes cadeias televisivas:



Festival Eurovisão da Canção 1960



O \textbf{Festival Eurovisão da Canção 1960} (MISC) (em inglês : \textbf{Eurovision Song Contest} (ORG) 1960 e em francês : \textbf{Concours Eurovision} (MISC) de la chanson 1960 ) foi o \textbf{5 Festival} (MISC) Eurovisão da Canção e teve lugar a 29 de março de 1960 em \textbf{Londres} (LOC) . \textbf{Katie Boyle} (PER) foi a apresentadora do evento que foi ganho por \textbf{Jacqueline Boyer} (PER) que representou a \textbf{França} (LOC) com a canção "\textbf{Tom Pillibi} (MISC)".

Apesar de terem ganho o festival em 1959, os \textbf{Países Baixos} (LOC) declinaram realizar o festival em 1960 pois já tinham realizado o Festival Eurovisão da Canção 1958 e entregaram essa honra ao \textbf{Reino Unido} (LOC) , segundo classificado em 1959 . A início, a emissora britânica estava relutante em abraçar o \textbf{Festival Eurovisão da Canção} (MISC) , mas com o país considerando a sua primeira tentativa de se juntar à recém-criada \textbf{Comunidade Económica Europeia} (ORG) (o precursor da atual \textbf{União Europeia} (ORG) ) surgiu a oportunidade para o \textbf{Reino Unido} (LOC) mostrar as suas credenciais europeias a uma audiência de televisão continental.

Foi nesse ano que se instaurou a tradição do vencedor do ano anterior entregar o troféu ao vencedor do ano, com \textbf{Teddy Scholten} (PER) a deslocar-se a \textbf{Londres} (LOC) para entregar o prémio a \textbf{Jacqueline Boyer} (PER) , vencedora de 1960 .



\textbf{Local} (MISC)O \textbf{Festival Eurovisão da Canção} (MISC) 1960 ocorreu em \textbf{Londres} (LOC) , na \textbf{Reino Unido} (LOC) . Dado que a \textbf{BBC} (ORG) está sediada em \textbf{Londres} (LOC) , uma cidade cheia de locais adequados para sediar o \textbf{Festival Eurovisão da Canção} (MISC) , a decisão de hostear o certame na capital britânica não foi surpreendente. No entanto, decidir realizar o concurso no \textbf{Royal Festival Hall} (LOC) (à esquerda) foi bastante ambicioso, já que este foi de longe o maior local para sediar o evento até aquela data.

O \textbf{Royal Festival Hall} (ORG) , visualmente um tanto rígido e modernista, foi construído na margem sul do \textbf{rio Tâmisa} (LOC) e inaugurado em 1951 para o "\textbf{Festival da Grã-Bretanha} (MISC)". Tem sido um local popular para concertos e conferências desde então. Para a \textbf{Eurovisão} (LOC), a audiência ao vivo excede as 2000 pessoas. Hoje, o local faz parte do \textbf{Southbank Centre} (LOC) e desde a realização do concurso, o \textbf{Royal Festival Hall} (ORG) passou por diversas transformações.

Como se para coincidir com o grande local, este foi o maior concurso até à data, tanto em termos de países participantes (13, com o regresso do \textbf{Luxemburgo} (LOC) e a estreia da \textbf{Noruega} (LOC) ) e também em termos de audiência de televisão ao vivo. No início dos anos 1960 , os aparelhos de televisão estavam-se a tornar mais comuns nos lares europeus e acredita-se que este foi o \textbf{Festival Eurovisão da Canção} (MISC) onde a audiência de televisão excedeu a audiência da rádio (acredita-se que cerca de 30 milhões de pessoas acompanharam o concurso, transmitido em 14 países, os concorrentes, mais a \textbf{Finlândia} (LOC) ). . Por outro lado, foi o último \textbf{Festival Eurovisão da Canção} (MISC) a ter lugar a meio da semana. A partir de 1961 , introduziu-se a tradição de ser realizado ao sábado.



\textbf{Formato} (PER) e Visual

Embora a encenação do Festival Eurovisão da Canção de 1960 possa não ter sido tão atraente quanto a competição do ano anterior em \textbf{Cannes} (LOC), não há dúvida de que a produção da \textbf{BBC} (ORG) foi uma das melhores da história da \textbf{Eurovisão} (LOC). Mais uma vez o espetáculo começou com uma introdução multilingue dos cantores. A apresentadora encarregue foi \textbf{Katie Boyle} (PER) , uma mulher que ainda hoje é uma lenda nos círculos da \textbf{Eurovisão} (LOC) . Este foi o primeiro de quatro edições apresentados por ela. Nascida na \textbf{Itália} (LOC) , \textbf{Boyle} (PER) voltaria a apresentar o certame em 1963 , 1968 e 1974 , estabelecendo um recorde que ainda não foi batido.

Após as apresentações iniciais, cada uma das músicas foi cantada em rápida sucessão, com uma peça incomum de teatro vendo a cortina do palco cair no final de cada apresentação. Apesar do facto de que não houve nenhum intervalo, a edição durou 82 minutos, atrasado ligeiramente por uma longa sequência de votação. Enquanto esta edição pode não ter tido o glamour do concurso do ano anterior, as músicas da mais alta qualidade, mas com uma sequência de votação muito empolgante para compensá-lo, foram certamente apresentadas com eficiência e bem encenadas e como o primeiro concurso apresentado pela \textbf{BBC} (ORG) para os outros nos próximos anos.

O sorteio da ordem de atuação foi feita a 28 de março , com os ensaios a decorrer entre esse dia e 29 de março .



\textbf{Votação} (MISC)Cada país tinha um júri constituído por 10 membros e cada um destes escolhia a melhor canção atribuindo 1 ponto.

Mais uma vez não houve intervalo, então após uma breve explicação do processo de votação, chegou a hora de os júris serem chamados. Infelizmente, o quadro operado manualmente teve alguns problemas, e a votação levou muito mais tempo do que nos anos anteriores. Como em anos anteriores, os países votaram na ordem inversa à da sua atuação, e o júri francês foi o primeiro a votar. Talvez influenciados pelo título francês, eles colocaram a \textbf{Alemanha} (LOC) na dianteira, mas ignoraram completamente a música italiana que seria a principal rival da \textbf{França} (LOC) . O júri italiano adotou uma tática semelhante, ignorando a canção francesa e iniciando especulações de votação tática deliberada que persiste até hoje. O primeiro movimento dramático no quadro veio quando a \textbf{Alemanha} (LOC) outorgou sete de seus dez pontos para o \textbf{Mónaco} (LOC), colocando-os numa liderança muito forte. No entanto, com altos votos dos \textbf{Países Baixo} (LOC) e \textbf{Suíça} (LOC) , o \textbf{Reino Unido} (LOC) logo assumiu a liderança, chegando a liderar por 8 pontos.

Os dois favoritos tiveram destinos diferentes ao longo da votação. A canção italiana nunca teve qualquer impacto significativo na votação, mas a francesa começou a fazer um forte progresso a meio da votação. Com quatro países por votar, os dinamarqueses concederam quatro pontos à \textbf{França} (LOC) e ignoraram o \textbf{Reino Unido} (LOC) , mudando as posições dos líderes, mas o \textbf{Luxemburgo} (LOC) rapidamente reverteu isso com cinco votos para o \textbf{Reino Unido} (LOC) e apenas um para a \textbf{França} (LOC) . A audiência no \textbf{Royal Festival Hall} (ORG) ficou cada vez mais animada, ao ponto da apresentadora pedir calma ao público, pois estava a tornar-se cada vez mais difícil ouvir os votos. Com dois países por votar, um único ponto separou o \textbf{Reino Unido da França} (LOC) , sem mais ninguém numa disputa séria. Quando a \textbf{Suécia} (LOC) deu quatro pontos à \textbf{França} (LOC) , a empolgação chegou ao fim, pois o país seguinte a votar, e +ultimo, o \textbf{Reino Unido} (LOC) tinha sido relegado para segundo lugar e não podia votar em si própria.

Após pensar-se que o país anfitrião iria conseguir uma vitória em casa pela primeira vez desde 1956 , a \textbf{França} (LOC) acabou por ganhar, seguida do \textbf{Reino Unido} (LOC) e do \textbf{Mónaco} (LOC). Mais uma vez, uma canção italiana altamente popular não conseguiu empolgar os júris internacionais, com " \textbf{Romantica} (MISC) " classificada em 8 lugar. No entanto, isso não impediria que a música se tornasse um sucesso internacional considerável nos meses seguintes. Parecendo-se tratar do primeiro "nul points" da história, a \textbf{Dinamarca} (LOC) foi pela primeira vez pontuada a seis países do fim da votação, terminando em 10 lugar com 4 pontos. Já o \textbf{Luxemburgo} (LOC) , que se fez representar pela primeira vez com uma canção em luxemburguês , terminou me último lugar com 1 ponto, mostrando a desvalorização das línguas minoritárias no certame.



\textbf{Participações individuais} (MISC)Cada país a concurso, executou a escolha do seu representante de maneira diferente, tendo sido marcado mais uma vez pela final nacional sueca (o \textbf{Melodifestivalen} (ORG) ), que foi considerado um pouco usual, por ter sido vencido um interprete e para o festival ter sido escolhido outro. Para características específicas sobre a entrada de cada participante no festival, seguir a ligação com o nome dos países na caixa em baixo. \textbf{Predefinição:Participações Individuais} (MISC) ESC1960



\textbf{Participantes Todos} (MISC) os onze países que haviam competido em \textbf{Cannes} (LOC) no ano anterior voltaram em 1960. Além disso, o \textbf{Luxemburgo} (LOC) , que não participou em 1959, voltou e a \textbf{Noruega} (LOC) estreou-se na \textbf{Eurovisão} (LOC), o que significa que pela primeira vez três países escandinavos estiveram no concurso.

Embora vários dos compositores estivessem fazendo um regresso ao evento, apenas um concorrente, o belga \textbf{Fud Leclerc} (PER) , repetiria a sua experiência eurovisiva. Curiosamente, \textbf{Leclerc} (LOC) também competiu na final nacional suíça e poderia ter representado os dois países diferentes em \textbf{Londres} (LOC) . Notavelmente, tal conflito de interesse só foi anulado nos últimos anos, depois que o grupo polaco \textbf{Ich Troje} (MISC) ter em 2003 estado perto de representar a \textbf{Alemanha} (LOC) e a \textbf{Polónia} (LOC) em simultâneo.



\textbf{Festival} (PER)Pelo segundo ano consecutivo, o país anfitrião foi o primeiro a atuar. " \textbf{Looking High} (MISC), \textbf{High High} (MISC) " (\textbf{Olhando Para o Alto} (MISC), \textbf{Alto} (LOC), \textbf{Alto} (LOC)) uma música alegre no estilo do tema da série de televisão "\textbf{Robin Hood} (MISC)". Talvez inspirado por seu irmão \textbf{Teddy} (PER) no ano anterior, o cantor \textbf{Bryan Johnson} (PER) escolheu assobiar em parte da canção. \textbf{Siw Malmkvist} (PER) , representando a \textbf{Suécia} (LOC) , continuou o clima de otimismo com " \textbf{Alla Andra Får Varann} (PER) " (Todos os Outros Apanham-se um ao \textbf{Outro} (MISC)) pode realmente ter tido uma melodia forte, mas liricamente fraco, tendo uma parte instrumental no meio que durou quase um minuto.

Excepcionalmente, o \textbf{Luxemburgo} (LOC) optou por usar a sua língua nacional, o luxemburguês para " \textbf{So laang we's du do bast} (MISC) " (\textbf{Tão Longe} (MISC) que tu estás) cantada por \textbf{Camillo Felgen} (PER) , estrela de sucesso no seu país natal e na \textbf{Alemanha} (LOC) , porém, nunca chegou a gravar esta música, devido à sua fraca classificação. Foi também o primeiro luxemburguês a representar o \textbf{Luxemburgo} (LOC) . A dinamarquesa \textbf{Katy Bødtger} (PER) , cantando sobre um passado sentimental, interpretou " \textbf{Det var en yndig tid} (MISC) " (Foi \textbf{Um Tempo Amoroso} (MISC)). O belga \textbf{Fud Leclerc} (PER) voltou à \textbf{Eurovisão} (LOC) pela terceira vez, com " \textbf{Mon} (PER) amour pour toi " (\textbf{O Meu Amor Para Ti} (MISC)), que pode ser melhor descrita como uma típica balada francesa.

A estreia da \textbf{Noruega} (LOC) , " \textbf{Voi Voi} (MISC) ", viu a cantora \textbf{Nora Brockstedt} (PER) numa versão bastante bizarra do traje tradicional da \textbf{Lapónia} (LOC) , numa música que se destacou após o que tinha acontecido antes. \textbf{Harry Winter} (PER) interpretou a canção austríaca " Du hast mich so fasziniert " (\textbf{Tu Fascinas-me Tanto} (MISC)), que foi mais uma tentativa de ganhar os votos com a clássica balada. A canção monegasca " Ce soir-là " (\textbf{Aquela Noite} (MISC)) foi composta por \textbf{Hubert Giraud} (PER) , que também compôs " \textbf{Dors} (MISC), mon amour ", vencedora em 1958 e, à semelhança desta, foi uma música forte, especialmente no desempenho do cantor francês \textbf{François Deguelt} (PER) .

A música " \textbf{Cielo} (LOC) e terra " (\textbf{Céu} (LOC) e \textbf{Terra} (LOC)) foi a primeira vez que a \textbf{Suíça} (LOC) foi representada por uma canção totalmente em italiano e, embora entusiasticamente interpretada por \textbf{Anita Traversi} (PER) , estava seriamente ausente de qualquer refrão cativante. Os \textbf{Países Baixos} (LOC) foram representados por \textbf{Rudi Carrell} (PER) com a canção " \textbf{Wat een geluk} (MISC) " (Que \textbf{Sorte} (MISC)), que juntou os compositores dos dois anteriores vencedores holandeses. Um desempenho bastante estranho por \textbf{Rudi} (PER) não disfarçou que esta era uma das canções mais curtas e insubstanciais da \textbf{História do certame} (MISC). A canção alemã interpretada por \textbf{Wyn Hoop} (PER) , com um título francês " \textbf{Bonne} (LOC) Nuit Ma Chérie " (\textbf{Boa Noite} (LOC), \textbf{Minha Querida} (MISC)), que combinava estilo latino com balada. A combinação entre músicas em alemão com título em francês iria resultar na vitória da \textbf{Áustria} (LOC) no Festival Eurovisão da Canção 1966 .

Os grandes favoritos à vitória foram os últimos a atuar. A \textbf{Itália} (LOC) , com a canção " \textbf{Romantica} (MISC) " (\textbf{Romântica} (LOC)) foi cantada pelo letrista \textbf{Renato Rascel} (PER) . Esta música, tal como as anteriores representantes italianas, foi um sucesso internacional e combinava arranjos orquestrais mistos e uma duração superior a 4 minutos. A última canção, e uma das mais aplaudidas, veio da \textbf{França} (LOC) . " \textbf{Tom Pillibi} (PER) ", interpretado pela jovem de 18 anos \textbf{Jacqueline Boyer} (PER) , que obteve algum sucesso além fronteiras, tocando nas rádios de países francófonos e não só, durante o ano de 1960.

A ordem de votação, foi ao contrário da ordem de actuação dos países no festival. Sendo assim, a ordem de votação no Festival Eurovisão da Canção 1960, foi a seguinte:



Transmissão

Os canais de televisão responsáveis pela difussão do concurso quer via televisão, quer via rádio foram as seguintes cadeias televisivas:



Festival Eurovisão da Canção 1961



O \textbf{Festival Eurovisão da Canção 1961} (MISC) (em inglês : \textbf{Eurovision Song Contest} (ORG) 1961 e em francês : \textbf{Concours Eurovision} (MISC) de la chanson 1961 ) foi o sexto \textbf{Festival Eurovisão da Canção} (MISC) e realizou-se na noite do dia 18 de março de 1961 em \textbf{Cannes} (LOC) , em \textbf{França} (LOC) . \textbf{Jacqueline Joubert} (PER) foi a apresentadora do festival que foi ganho pelo cantor francês \textbf{Jean-Claude Pascal} (PER) que representou o \textbf{Luxemburgo} (LOC) com a canção " \textbf{Nous les amoureux} (MISC) " (Nós os amantes), que falava, implicitamente, sobre um amor homossexual, algo que ainda era criminalizado em alguns países da \textbf{Europa} (LOC) à época. \textbf{Jean-Claude Pascal} (PER) foi também o primeiro vencedor a representar um país que não era o seu, visto que o cantor era francês e representou o \textbf{Luxemburgo} (LOC) .

Como em 1959 , o \textbf{Festival} (LOC) teve lugar no \textbf{Palais des Festivals et des Congrès} (MISC) , em \textbf{Cannes} (LOC) , novamente apresentada por \textbf{Jacqueline Joubert} (PER).

O número de participantes subiu para 16 e pela primeira vez, o certame realizou-se a um sábado, tradição que se mantêm até hoje, apenas interrompida em 1962 , que se realizou a um domingo.



\textbf{Local} (MISC)O \textbf{Festival Eurovisão da Canção} (MISC) em 1961 ocorreu, mais uma vez, em \textbf{Cannes} (LOC) , na \textbf{França} (LOC) . O famoso \textbf{Palais des Festivals} (MISC) (à esquerda), no coração da cidade, era na época o maior local para organizar o concurso. O original \textbf{Palais des Festivals} (MISC) foi construído em 1949 para sediar o festival de cinema. O edifício era localizado na famosa avenida de \textbf{La Croisette} (LOC), no local atual do \textbf{Hotel Noga Hilton} (LOC). Apesar de também encenar o \textbf{Festival Eurovisão da Canção 1959} (MISC) , hoje o antigo \textbf{Palais des Festivals} (MISC) já não existe. Em resposta ao crescente sucesso do festival de cinema e o advento de várias convenções de grandes empresas, a cidade de \textbf{Cannes} (LOC) decidiu construir um novo e maior espaço em 1979. O festival de cinema mudou-se para o novo \textbf{Palais des Festivals et des Congrès} (MISC) em 1983 e logo após o antigo \textbf{Palais des Festivals} (MISC) foi destruído. Apesar de sua maior capacidade e conceito moderno, o design do atual \textbf{Palais des Festivals} (MISC) foi originalmente condenado por muitos como semelhante a um bunker gigante, em comparação com o slendour do original \textbf{Palais des Festivals} (MISC).



\textbf{Formato} (PER) e Visual

O vídeo introdutório foi inteiramente semelhante ao de 1959 : primeiro uma visão do mar e do \textbf{Palais des Festivals et des Congrès} (MISC) ; depois, uma visão do local, cabines do comentadores, o público, a orquestra e, finalmente, uma cortina, que por detrás desvendava um cenário, maior do que nos anos anteriores, representando um jardim de um grande castelo mediterrâneo , com escadas de pedra, flores e tapeçarias e com o quadro de votações posicionado à esquerda do palco. Abriu então, desvelando os artistas, todos reunidos no palco, pela primeira vez história do concurso. Surgiu então a apresentadora, \textbf{Jacqueline Joubert} (PER), que se tornou a primeira pessoa a apresentar duas vezes o concurso. Começou com o clássico "Boa noite, \textbf{Europa} (LOC)!" E conclui a sua introdução com estas palavras: "Para aproximar as pessoas, um pouco de refrão às vezes é muito melhor do que um longo discurso." \textbf{Joubert} (PER) deu então as boas vindas aos participantes, que, um por vez, avançaram para o microfone e soletraram seu nome. Depois, a própria \textbf{Joubert} (PER) introduziu cada música.

Como em 1956, houve um intervalo de alguns minutos entre a última música e a votação. A \textbf{RTF} (MISC) contratou dois dançarinos da \textbf{Ópera de Paris} (LOC) para dar uma breve representação de seus talentos, enquanto os jurados contavam seus votos.

A emissão durou 1 hora e 35 minutos, com o evento a ser seguido por 35 milhões de pessoas.

O sorteio da ordem de atuação foi feita a 16 de março .

Como resultado do tempo limite do concurso, a música vencedora não foi transmitida no \textbf{Reino Unido} (LOC) .



\textbf{Votação} (MISC)Cada país tinha um júri constituído por 10 membros e cada um destes escolhia a melhor canção atribuindo 1 ponto. Foi notado que, durante a votação, o \textbf{Luxemburgo} (LOC) outorgou 8 pontos ao \textbf{Reino Unido} (LOC) e a \textbf{Noruega} (LOC) 8 à \textbf{Dinamarca} (LOC) . Este foi o maior número de pontos dados a um só país desde 1958 , quando a \textbf{França} (LOC) recebeu 9 pontos da \textbf{Dinamarca} (LOC) .

O \textbf{Reino Unido} (LOC) assumiu a liderança desde o início da votação, em especial graças ao júri do \textbf{Luxemburgo} (LOC) , que lhe atribuiu 8 votos. Ocorreu um erro durante a votação do júri belga, que concedeu 1 voto ao \textbf{Reino Unido} (LOC) , apesar de terem aparecido 4 no quadro. O \textbf{Luxemburgo} (LOC) retomou o seu atraso na segunda parte da votação e passou à frente quando o erro foi corrigido. De fato, após o fim da votação do júri austríaco, os 25 pontos do \textbf{Reino Unido} (LOC) foram reduzidos para 21. \textbf{Jacqueline Joubert} (PER) perguntou à audiência se houve um erro e o público respondeu afirmativamente. Ela respondeu: "\textbf{Vocês} (LOC) contam melhor do que eu!" Imediatamente adicionando: "É verdade que estou muito comovida para chamar todos esses países. Se vocês estivessem no meu lugar... eu nunca falei com tantas pessoas ao mesmo tempo."



\textbf{Participantes} (MISC)De entre os concorrentes, encontrava-se a alemã \textbf{Lale Andersen} (PER) , conhecida pela sua canção "\textbf{Lili Marleen} (MISC)", que fez sucesso durante a \textbf{Segunda Guerra Mundial} (MISC) . Para espanto de todos, \textbf{Lale Andersen} (PER) também interpretou parte da sua música em francês. Era a mais velha participante na edição, com 55 anos no dia do certame.

Pela primeira vez, um cantor de origem grega participava na \textbf{Eurovisão} (LOC), tratando-se de \textbf{Jimmy Makulis} (PER) , que representou a \textbf{Áustria} (LOC) .

Um total de cinco músicas tiveram como tema principal as estações do ano , com a \textbf{França} (LOC) , o \textbf{Mónaco} (LOC) e a \textbf{Suécia} (LOC) a cantarem sobre a primavera e a \textbf{Bélgica} (LOC) e a \textbf{Noruega} (LOC) a cantarem sobre o \textbf{Verão} (MISC) .



\textbf{Festival} (PER)A ordem de votação, foi ao contrário da ordem de actuação dos países no festival. Sendo assim, a ordem de votação no Festival Eurovisão da Canção 1961, foi a seguinte:



Transmissão

Os canais de televisão responsáveis pela difussão do concurso quer via televisão, quer via rádio foram as seguintes cadeias televisivas:



Festival Eurovisão da Canção 1962



O \textbf{Festival Eurovisão da Canção 1962} (MISC) (em inglês : \textbf{Eurovision Song Contest} (ORG) 1962 e em francês : \textbf{Concours Eurovision} (MISC) de la chanson 1962 ) foi o 7 \textbf{Festival Eurovisão da Canção} (MISC) e foi realizado a 18 de março de 1962 no \textbf{Luxemburgo} (LOC) . \textbf{Mireille Delannoy} (PER) foi a apresentadora do festival que foi ganho pela cantora francesa \textbf{Isabelle Aubret} (PER) que representou a \textbf{França} (LOC) , com a canção " \textbf{Un premier amour} (MISC) ", fazendo da \textbf{França} (LOC) o primeiro país a ganhar o certame três vezes. Esta edição foi a única que se realizou a um \textbf{Domingo} (MISC). Pela primeira vez, o top 3 foi constituído por canções interpretadas na mesma língua, o francês .

Devido ao fraco sistema de votação, onde apenas 3 países (em 16) eram pontuados por cada júri nacional, a \textbf{Áustria} (LOC) , a \textbf{Bélgica} (LOC) , a \textbf{Espanha} (LOC) e os \textbf{Países Baixos} (LOC) partilharam o último lugar, com 0 pontos. Pela primeira vez na história, havia "nul points".



\textbf{Local} (MISC)O \textbf{Festival Eurovisão da Canção} (MISC) 1962 ocorreu na capital do \textbf{Luxemburgo} (LOC) . A cidade do \textbf{Luxemburgo} (LOC) é uma comuna de \textbf{Luxemburgo} (LOC) com status de cidade , pertencente ao distrito de \textbf{Luxemburgo} (LOC) e ao cantão de \textbf{Luxemburgo} (LOC) . \textbf{Luxemburgo} (LOC) é uma cidade muito desenvolvida em seu comércio e nas indústrias. \textbf{Luxemburgo} (LOC) é uma das cidades mais ricas da \textbf{Europa} (LOC) e tornou-se um importante centro financeiro e administrativo. A cidade do \textbf{Luxemburgo} (LOC) contém várias instituições da \textbf{União Europeia} (ORG) , incluindo o \textbf{Tribunal de Justiça Europeu} (ORG) , o \textbf{Tribunal de Contas} (ORG) e o \textbf{Banco Europeu de Investimento} (ORG) .

O festival em si realizou-se no \textbf{Villa Louvigny} (LOC) , no sul do \textbf{Luxemburgo} (LOC) , que serviu de sede da \textbf{Compagnie Luxembourgeoise de Télédiffusion} (MISC) , o precursor da \textbf{RTL Group} (MISC) . Está localizado no \textbf{Parque Municipal} (LOC) , no bairro \textbf{Ville Haute} (LOC) do centro da cidade.



\textbf{Formato} (LOC) e Visual

O vídeo introdutório começou com uma vista aérea da \textbf{Villa Louvigny} (LOC) . Enquanto a orquestra interpretou o instrumental de " \textbf{Nous les amoureux} (MISC) ", a canção vencedora do ano anterior, a câmara ofereceu uma visão da sala, do público e, finalmente, da orquestra.

A apresentadora foi \textbf{Mireille Delannoy} (PER) , que abriu a noite com o inevitável "Boa noite, \textbf{Europa} (LOC)!" , cumprimentando em seguida todos os países participantes nas suas línguas oficiais. Concluiu com: "Boa noite, milhões de amigos!"

A decoração lembrava as mansões do século XVI com enfeites e flores. O salão brilhante era composto por personalidades, principalmente francesas, convidadas pela \textbf{Télé Luxembourg} (LOC) . A orquestra, os artistas e o quadro de votação foram colocados no mesmo palco, a orquestra à esquerda, os artistas no centro e o quadro à direita. O pano de fundo por trás dos artistas incluía uma entrada, uma balaustrada de ferro forjado, colunas, bacias e janelas cruzadas. Das janelas podia-se ver um céu pontilhado de estrelas cintilantes.

Uma nova regra foi introduzida, limitando a duração das músicas a três minutos. Esta regra ainda está em vigor hoje.

As músicas participantes foram interpretadas sem apresentação individual. Um pequeno incidente ocorreu logo após a passagem da canção francesa, um corte de poder geral mergulhou o auditório no escuro. Os espectadores europeus tiveram que esperar alguns segundos na frente de uma tela preta. O mesmo tinha acontecido antes, durante a interpretação da canção holandesa .

A \textbf{Alemanha} (LOC) , com \textbf{Conny Froboess} (PER) e sua canção " \textbf{Zwei kleine Italiener} (MISC) " (Dois pequenos italianos) era uma das grandes favoritas à vitória. Contudo, viria a classificar-se em 6 lugar, com 9 pontos.

Depois da última música, pensou-se que iria aparecer em palco um 17 concorrente. Na verdade, tratava-se do famoso palhaço \textbf{Achille Zavatta} (PER) , responsável pelo intervalo. Ele representou "\textbf{Zavattaland} (LOC)" e tocou diferentes instrumentos musicais que ele não hesitou em roubar alguns músicos da orquestra, antes de ser liberado da cena por dois seguranças.



\textbf{Votação} (MISC)

Um novo sistema de votação foi utilizado. Assim, cada país tinha 10 júris que atribuíram 3, 2 e 1 ponto às suas três canções favoritas, por ordem de preferência; somaram os totais de cada país e deram à canção mais votada 3 pontos, à segunda 2 pontos e à terceira 1 ponto. Os resultados dos votos foram anunciados por ordem decrescente de votos: 3, 2 e 1 voto.

A votação só conheceu um engano. Depois de agradecer ao porta-voz do júri sueco, \textbf{Mireille Delannoy} (PER) , contactando logo o próximo país a votar, a \textbf{Dinamarca} (LOC) , disse: "\textbf{Olá} (MISC), \textbf{Dinamarca} (LOC). \textbf{Olá Estocolmo} (LOC). Você nos ouve?". O porta-voz do júri dinamarquês disse: "Não, \textbf{Copenhague} (LOC)!". \textbf{Delannoy} (PER) desculpou-se. Então o porta-voz leu os votos dinamarqueses na ordem de passagem dos participantes. \textbf{Delannoy} (PER) perdeu-se na votação, acreditando que a \textbf{Dinamarca} (LOC) havia concedido duas vezes os seus 2 votos. O porta-voz teve que repetir os resultados dinamarqueses pela segunda vez.



\textbf{Participantes} (MISC)O representante da \textbf{Suíça} (LOC) , \textbf{Jean Philippe} (PER) , tornou-se o primeiro artista a representar dois países, tendo representado a \textbf{França} (LOC) em 1959 , com " \textbf{Oui} (MISC), oui, \textbf{oui} (MISC), \textbf{oui} (MISC) ", que terminou em 3 lugar.



\textbf{Festival} (MISC)A ordem de votação, foi ao contrário da ordem de actuação dos países no festival. Sendo assim, a ordem de votação no Festival Eurovisão da Canção 1962, foi a seguinte:

Os países que receberam 5 pontos foram os seguintes:



Transmissão

Os canais de televisão responsáveis pela difussão do concurso quer via televisão, quer via rádio foram as seguintes cadeias televisivas:



Festival Eurovisão da Canção 1963



O \textbf{Festival Eurovisão da Canção 1963} (MISC) (em inglês : \textbf{Eurovision Song Contest} (ORG) 1963 e em francês : \textbf{Concours Eurovision} (MISC) de la chanson 1963 ) foi o 8 \textbf{Festival Eurovisão da Canção} (MISC) , que teve lugar no \textbf{Reino Unido} (LOC) . A \textbf{Dinamarca} (LOC) venceu pela primeira vez com a canção " Dansevise " interpretado pelo duo \textbf{Grethe \& Jørgen Ingmann} (ORG) . Por seu turno, a \textbf{Finlândia} (LOC) , a \textbf{Noruega} (LOC) a \textbf{Suécia} (LOC) e os \textbf{Países Baixos} (LOC) não obtiveram qualquer ponto, com este último a tornar-se o primeiro país a não obter qualquer pontuação dois anos consecutivos.

Depois da \textbf{França} (LOC) (vencedora do ano anterior), ter recusado organizar pela terceira vez o certame, passou-se a vez ao segundo classificado, o \textbf{Mónaco} (LOC) que também recusou, assim como o \textbf{Luxemburgo} (LOC) (terceiro classificado). Em outubro de 1962 , foi finalmente anunciado que seria o \textbf{Reino Unido} (LOC) a organizar a edição de 1963.

Esta edição apresenta uma alteração na regra de solistas e duetos. Estão autorizados a ser acompanhado por artistas principais e em segundo plano um coro de até três pessoas.

Esta edição ficou ainda marcada pela controversa votação da \textbf{Noruega} (LOC) (que votou primeiro de uma maneira, e, quando foi novamente chamada a votar, modificou os seus resultados) e, pela primeira vez, uma mulher, \textbf{Yvonne Littlewood} (PER) , responsável pela produção do espetáculo.

Esta foi a primeira edição do Festival Eurovisão da Canção que foi transmitida em direto para \textbf{Portugal} (LOC) , através da \textbf{RTP} (ORG) . Em edição foi vista por 50 milhões de pessoas em todo o mundo.

\textbf{Nana Mouskouri} (PER) , \textbf{Françoise Hardy} (PER) , \textbf{Esther Ofarim} (PER) e \textbf{Alain Barrière} (PER) foram alguns dos artistas famosos que participaram nesta edição, com o último a ter o seu tema (" \textbf{Elle} (MISC) était si jolie ") a transformar-se num êxito internacional.



\textbf{Local} (MISC)O \textbf{Festival Eurovisão da Canção} (MISC) 1963 ocorreu em \textbf{Londres} (LOC) , no \textbf{Reino Unido} (LOC) . \textbf{Londres} (LOC) é uma das maiores cidades da \textbf{Europa} (LOC) e já liderou a lista das maiores cidades do mundo . De uma cidade romana chamada \textbf{Londínio} (LOC) , a capital da região administrativa romana \textbf{Britannia} (LOC), a cidade acabou por se tornar o centro do império britânico, contribuindo hoje com 17\% do PNB de um país que é a quarta maior economia do mundo. \textbf{Londres} (LOC) tem sido um dos mais importantes centros da política e do comércio mundiais por quase dois milénios (apesar de a capital de \textbf{Inglaterra} (LOC) ter sido \textbf{Winchester} (LOC) durante grande parte da \textbf{Idade Média} (MISC) ).

O festival em si realizou-se na \textbf{BBC Television Centre} (LOC) . Inaugurado em 1960 , a \textbf{BBC Television Centre} (ORG) é uma das instalações mais facilmente reconhecíveis de seu tipo, tendo surgido como pano de fundo para muitos programas da \textbf{BBC} (ORG) . Permaneceu como uma das maiores instalações do mundo até o fechamento em março de 2013.



\textbf{Visual Como o estúdio} (MISC) \textbf{TC1} (MISC), à época o maior estúdio da \textbf{Europa} (LOC), não estaria pronto até abril de 1964 , foram usados os estúdios \textbf{TC3} (MISC) e \textbf{TC4} (MISC). Um terceiro estúdio, o \textbf{TC5} (MISC), foi usado pelo júri britânico. Dois estúdios (\textbf{TC3} (MISC) e \textbf{TC4} (MISC)) foram usados: um para a plateia, para o quadro de votações e para a apresentadora, \textbf{Katie Boyle} (PER) ; e outro, com melhor acústica, para a orquestra e para os artistas. A decoração do primeiro estúdio consistia numa escadaria ladeada por pilastras douradas e grandes espelhos. O quadro de votação foi colocado à direita do palco. A decoração do segunda estúdio consistia num piso feito de vários pódios hexagonais de diferentes tamanhos. Para isso, foram adicionadas uma escada, um pórtico e outros acessórios, conforme as atuações dos vários artistas.

A principal novidade técnica da edição de 1963 (para além do uso de duas salas separadas) foi o uso de microfones suspensos. Estes, invisíveis na tela, permitiram maior liberdade de movimento aos artistas e uma encenação mais elaborada. Cada artista tinha, portanto, o direito a uma apresentação personalizada. Esta técnica deu origem a persistentes, mas imprecisos rumores de que os artistas não interpretaram as músicas ao vivo. Pela primeira vez, efeitos visuais foram usados para as atuações: close-ups, imagens de rastreamento, imagens sobrepostas e filtros animados. Além disso, pela primeira vez, uma máquina de vento foi usada. O objetivo desta edição foi inovar para criar um novo visual.

O sorteio da ordem de atuação foi feita a 26 de novembro de 1962 , muito mais cedo do que em edições anteriores, de modo a que os artistas posteriormente na ordem de atuação não precisassem de estar presentes durante os três dias de ensaios.

Os ensaios começaram a 21 de março , com os primeiros 8 países a terem um ensaio de 55 minutos com a orquestra, composta por 45 elementos.



\textbf{Formato} (PER)O vídeo introdutório começou com a apresentação dos nomes dos países participantes no logótipo da \textbf{Eurovisão} (LOC), nas suas línguas oficiais. Depois, começou com uma vista aérea da \textbf{BBC Television Center} (ORG) e depois dos dois estúdios. Após as saudações habituais, \textbf{Katie Boyle} (PER) chamou os artistas de cada vez, que se apresentaram ao público

Cada música foi introduzida por uma visão em um mapa da \textbf{Europa} (LOC). Uma lâmpada luminosa, colocada no local da capital, indicava o país participante. O nome deste, da música e dos artistas apareceu então sobreposto no mapa.

Inicialmente, planeava-se inserir entre cada música trinta segundos de sketches com os bonecos \textbf{Pinky \& Perky} (MISC) . Essas sequências foram, portanto, roteirizadas e gravadas nnm estúdio em \textbf{Manchester} (LOC) , e até pré-publicitados, mas a produção mudou no último momento, não sendo transmitidas.

Para o intervalo, os equilibristas \textbf{Ola \& Barbro} (MISC) interpretaram uma atuação de circo. Barbro fez algumas figuras de equilíbrio, no guidão de uma bicicleta dirigida por \textbf{Ola} (PER). Esta última então realizou outras acrobacias, sozinha, na mesma bicicleta. Ambos concluíram com números de dueto.



\textbf{Votação} (MISC)

Um novo sistema de votação foi utilizado. Assim, a votação baseou-se em 20 júris por país que atribuíram 5, 4, 3, 2 e 1 ponto às suas cinco canções favoritas, por ordem de preferência; somaram os totais de cada país e deram à canção mais votada 5 pontos, à segunda 4 pontos, à terceira 3 pontos, à quarta 2 pontos e à quinta 1 ponto. Uma outra novidade foi o facto da apresentadora também estar em ligação telefónica com um escrutinador da \textbf{EBU} (MISC) , um árbitro cuja decisão caberia em caso de dificuldades.

Aconteceu algo de controverso durante a votação. Durante a votação norueguesa, o porta-voz do júri enganou-se ao revelar as votações esquecendo-se de ordená-las (pois a votação era dada pela ordem de actuação e não de forma decrescente). \textbf{Katie Boyle} (PER) , a pedido do porta-voz, adiou a divulgação dos resultados noruegueses para o final da votação dos outros júris nacionais. Quando voltou para dar os votos, o júri norueguês alterou-os e deu 5 pontos ao \textbf{Reino Unido} (LOC) ; 4 pontos à \textbf{Dinamarca} (LOC) (anteriormente dados à \textbf{Itália} (LOC) ), 3 pontos à \textbf{Itália} (LOC) (anteriormente dados à \textbf{Suíça} (LOC) ), 2 pontos à \textbf{Alemanha} (LOC) (anteriormente dados à \textbf{Dinamarca} (LOC) ) e 1 ponto à \textbf{Suíça} (LOC) (anteriormente dado à \textbf{Alemanha} (LOC) ), roubando o primeiro lugar da \textbf{Suíça} (LOC) para a "vizinha" \textbf{Dinamarca} (LOC). Ainda hoje não se sabe, ao certo, se foi uma confusão ou se a alteração foi intencional para beneficiar outro país escandinavo . No entanto, sabe-se que os resultados do júri norueguês não estavam prontos quando foi chamado pela primeira vez por \textbf{Katie Boyle} (PER) . O seu presidente ainda estava ocupado a somar os votos dos jurados. Surpreendido, \textbf{Roald Øyen} (PER) teria lido os resultados provisórios.

Após a segunda votação da \textbf{Noruega} (LOC) , o \textbf{Mónaco} (LOC) também foi chamado a divulgar os seus votos pela segunda vez, pois tinha atribuído duas vezes um ponto (para o \textbf{Reino Unido} (LOC) e para o \textbf{Luxemburgo} (LOC) ). Após a releitura dos votos, foi retirado um ponto a o \textbf{Luxemburgo} (LOC), que não alterou em nada a classificação dos países

Também se especulou se os júris estavam de facto no outro lado da linha ou no próprio estúdio, dada a clareza com que suas vozes podiam ser ouvidas, em vez de soar como se estivessem sendo redirecionadas por uma linha telefónica.

Esta foi também uma das votações mais renhidas da \textbf{História} (MISC), visto que ora liderava a \textbf{Dinamarca} (LOC) , ora liderava a \textbf{Suíça} (LOC) , com a primeira acabando por vencer por apenas 2 pontos.



\textbf{Participantes} (MISC)De entre os participantes encontravam-se vários nomes de renome internacional, como \textbf{Nana Mouskouri} (PER) (representando o \textbf{Luxemburgo} (LOC) ), \textbf{Françoise Hardy} (PER) (representando o \textbf{Mónaco} (LOC) ), \textbf{Esther Ofarim} (PER) (representando a \textbf{Suíça} (LOC) ) e \textbf{Alain Barrière} (PER) (representando a \textbf{França} (LOC) ).



\textbf{Festival} (PER)Pela primeira vez, a ordem de votação, respeitou a ordem de atuação. A \textbf{Noruega} (LOC) , apesar de ter sido a quinta a votar, aparece a tabela em último, por ter sido novamente chamada a votar. Sendo assim, a ordem de votação no Festival Eurovisão da Canção 1963, foi a seguinte:

Os países que receberam 5 pontos foram os seguintes:



Transmissão

Os canais de televisão responsáveis pela difussão do concurso quer via televisão, quer via rádio foram as seguintes cadeias televisivas:



Festival Eurovisão da Canção 1964



\textbf{Países Participantes Países} (LOC) que participaram no passado, mas não em 1964

O \textbf{Festival Eurovisão da Canção 1964} (MISC) (em inglês : \textbf{Eurovision Song Contest 1964} (ORG) , em francês : \textbf{Concours Eurovision} (MISC) de la chanson 1964 e em dinamarquês : \textbf{Eurovision Melodi Grand-prix 1964} (MISC) ) foi o 9 \textbf{Festival Eurovisão da Canção} (MISC) e teve lugar no dia 21 de março de 1964 , em \textbf{Copenhaga} (LOC) , na \textbf{Dinamarca} (LOC) . \textbf{Lotte Wæver} (PER) foi a apresentadora do festival que foi ganho por \textbf{Gigliola Cinquetti} (PER) que representou a \textbf{Itália} (LOC) com a canção " \textbf{Non ho l'età} (MISC) " ( Não tenho idade ), e que foi assistido por 100 milhões de pessoas.

\textbf{Portugal} (LOC) participou pela primeira vez com \textbf{António Calvário} (PER) cantando \textbf{Oração} (LOC) , contudo a estreia não foi auspiciosa: a canção não teve qualquer voto (tornou-se o primeiro país a estrear-se sem qualquer ponto; nesta edição, a \textbf{Alemanha} (LOC) , a \textbf{Jugoslávia} (ORG) e a \textbf{Suíça} (LOC) também não teriam qualquer ponto) e a política ditatorial de \textbf{António Oliveira Salazar} (PER) e a \textbf{Guerra Colonial} (MISC) , contribuiu para uma recepção fria por parte do público presente.

Este festival foi marcado por um protesto político depois da participação suíça : um homem entrou em palco e exibiu um cartaz a protestar contra \textbf{Francisco Franco} (PER) e \textbf{António de Oliveira Salazar} (PER) , cartaz esse que dizia "\textbf{Boicotem Franco} (PER) e \textbf{Salazar} (PER)". Enquanto isto se passava, os espectadores viram o quadro de votações, depois do homem ser retirado, a competição continuou sem mais sobressaltos. Curiosamente, não existe qualquer cópia das imagens deste festival, ao que parece, na década de 1970 , um fogo terá destruído os estúdios da \textbf{DR} (ORG) (a televisão estatal dinamarquesa ), onde estava a fita deste festival. Não existe qualquer gravação do festival na íntegra por parte das outras televisões (ainda houve rumores de que a \textbf{BBC} (ORG) tinha), pois estas só tinham o hábito de gravar o final, quando o cantor \textbf{vencedor} (PER) voltava a cantar a canção vencedora, o único registo vídeo que existe deste festival (a parte sonora áudio existe na totalidade) é o princípio (onde se podem ver os soldados reais dinamarqueses tocando o hino da \textbf{Eurovisão} (LOC)) até à apresentadora \textbf{Lotte Wæver} (PER) dizer umas palavras em dinamarquês e depois dizer "\textbf{Good Evening Europe} (MISC)" e a repetição da canção vencedora " \textbf{Non ho l'età} (MISC) " de \textbf{Gigliola Cinquetti} (PER) , tudo o resto desapareceu. No entanto, em dezembro de 2021, o utilizador do \textbf{Reddit DYLCWS} (ORG) carregou um vídeo de 3 minutos do festival que comprou à emissora finlandesa \textbf{YLE} (MISC) , onde aparece \textbf{Gigliola Cinquetti} (PER) a ser declarada a vencedora, \textbf{Lotte Wæver} (PER) a apresentar \textbf{Svend Pedersen} (PER) , que entregou o prémio a \textbf{Gigliola} (PER), um clipe do prémio e do público, e um trecho da canção vencedora.



\textbf{Local} (MISC)O \textbf{Festival Eurovisão da Canção 1964} (MISC) ocorreu em \textbf{Copenhaga} (LOC) , na \textbf{Dinamarca} (LOC) . \textbf{Copenhaga} (LOC) é a capital e a maior cidade da \textbf{Dinamarca} (LOC) e dispõe de privilégios de condado . Situa-se na costa do mar Báltico em frente da \textbf{Suécia} (LOC) . A sua principal atracção turística é a escultura \textbf{Sereiazinha} (LOC) ou \textbf{Pequena Sereia} (PER) situada no porto da cidade.

O festival em si realizou-se no \textbf{Tivolis Koncertsal} (MISC) . Inaugurada em 1956 e remodelado em 2005 , o \textbf{Tivolis Koncertsal} (MISC) é uma sala de espetáculos situado nos \textbf{Jardins de Tivoli} (LOC) . O edifício foi desenhado pelos arquitectos \textbf{Frits Schlegel} (MISC) e \textbf{Hans Hansen} (PER) . Vários artistas pop e rock já atuaram no \textbf{Tivolis} (LOC), como \textbf{The Grateful Dead} (ORG) , \textbf{Elton John} (PER) , \textbf{Procol Harum} (ORG) , \textbf{Uriah Heep} (ORG) , \textbf{Cream} (ORG) , \textbf{Jerry Lee Lewis} (PER) e \textbf{Jethro Tull} (PER) . Atualmente, recebe concertos de música clássica . Em 2005 , o edifício foi renovado e foi estendido pela empresa 3XN, onde o estilo clássico da década de 1950 , o principal auditório incluindo uma cor característica de vermelho, azul, amarelo e verde, foi restaurada, enquanto as instalações aos visitantes melhoraram e expandiram-se.



\textbf{Formato} (LOC) e visual

A apresentadora desta edição foi \textbf{Lotte Wæver} (PER) , que apresentou o festival maioritariamente em dinamarquês , mas também em inglês e francês . Antes de cada atuação, introduzia o diretor de orquestra, como o intérprete, mas também explicava cada tema.

A orquestra escolhida para esta edição foi a \textbf{Grand Prix Orchestra} (MISC), de 42 membros, dirigida por \textbf{Kai Mortensen} (PER) .

Pela primeira vez na história do festival, um intérprete foi autorizado a retornar ao palco após sua apresentação: após a interpretação de \textbf{Gigliola Cinquetti} (PER) , o público respondeu com aplausos prolongados e entusiasmados e \textbf{Cinquetti} (PER) pôde voltar para saudar novamente os espectadores.

O festival abriu com a atuação da \textbf{Guarda Real Dinamarquesa} (ORG), seguindo-se a apresentação de \textbf{Lotte Wæver} (PER) .

O intervalo foi preenchido com um balé intitulado "\textbf{Harlequinade} (LOC)", no qual os amores de \textbf{Arlequim} (LOC) e \textbf{Colombina} (LOC) foram encenados. A coreografia foi criada por \textbf{Niels Bjørn Larsen} (PER) , então diretor do \textbf{Ballet Real da Dinamarca} (MISC). Os principais papéis foram desempenhados por \textbf{Solveig Oestergaard} (PER) e \textbf{Niels Kehlet} (PER), que foram acompanhados pelos bailarinos do \textbf{Ballet Real da Dinamarca} (MISC).

O palco era simplesmente decorado com buquês de flores , painéis e medalhões exibindo vistas típicas da \textbf{Dinamarca} (LOC). Os artistas fizeram a sua entrada por uma escada em espiral, localizada na parte de trás do palco. A orquestra foi colocada num buraco, à frente e à direita do pódio. O quadro de votação estava localizado à esquerda do palco.

O festival de 1964 juntamente com o de 1956 são os únicos cujo registo em vídeo não existe e apenas são conhecidos os registos em áudio de cada um destes festivais. Julga-se que o desaparecimento do ficheiro foi devido a um incêndio nos arquivos da televisão dinamarquesa \textbf{DR} (ORG). Também se afirma que esta edição nunca foi gravada, aparentemente não havia tecnologia disponível para gravar este programa.



\textbf{Votação} (MISC)

Um novo sistema de votação foi utilizado. Assim, cada país tinha 10 júris e cada um tinha 9 votos para atribuir. A canção que obtivesse mais votos dentro do júri obtinha 5 pontos, o segundo 3 e a canção em 3 lugar 1 ponto. Se apenas uma canção tivesse todos os votos entre os júris obtinha 9 pontos e se apenas duas canções fossem escolhidas, a canção mais votada obtinha 6 pontos e a outra canção 3 pontos. A vitória da \textbf{Itália} (LOC) nunca foi posta em causa, principalmente após receber a pontuação máxima, 5 pontos, dos primeiros dois países a votar.

Os resultados dos votos foram anunciados oralmente, de acordo com a ordem crescente dos votos: 1, 3 e depois 5 votos.

Pela primeira vez, a \textbf{UER} (ORG) delegou um supervisor no local: \textbf{Miroslav Vilček} (LOC) , encarregado de monitorizar o processo de votação e controlar os resultados. Ele sentou-se com uma assistente junto ao quadro de votação. Ele deu o sinal de partida para a votação. Ele interveio duas vezes: primeiro, para acrescentar à mesa, uma votação da \textbf{Itália} (LOC) para a \textbf{Espanha} (LOC) , que havia sido esquecida; segundo, para que o porta-voz espanhol repetir a atribuição da pontuação máxima, pois o aplauso do público impediu-o de ouvir o nome do país beneficiário.

Pela primeira vez na história do concurso, os votos não foram escritos no tabuleiro, mas como flechas graduadas.



\textbf{Participantes Para} (MISC) além da já mencionada \textbf{Gigliola Cinquetti} (PER) , convém ainda realçar os \textbf{Países Baixos} (LOC) , que se tornaram o primeiro país a enviar uma representante de origem não europeia , visto que \textbf{Anneke Grönloh} (PER) era de origem indonésia . \textbf{Espanha} (LOC) também foi representada por um grupo não europeu, \textbf{Los TNT} (ORG) , que se tornaram a primeira banda com três ou mais elementos a atuar na \textbf{Eurovisão} (LOC). \textbf{Nelly} (ORG), a vocalista da banda, foi também a única a recorrer a um acessório durante a sua atuação, uma concha , que ilustrava o tema da canção, " \textbf{Caracola} (LOC) ". Representando a \textbf{Áustria} (LOC) estava \textbf{Udo Jürgens} (PER) , que dois anos mais tarde viria a conquistar a primeira vitória para a \textbf{Áustria} (LOC) , que se tornou o primeiro artista que se fez acompanhar por um piano em palco.

A resposta imediata do público de \textbf{Koncertsal} (PER) à canção italiana foi marcadamente entusiástica e prolongada e, mais incomum para uma atuação no concurso, depois de deixar o palco, \textbf{Gigliola Cinquetti} (PER) foi autorizada a retornar. A sua atuação foi repetida no dia seguinte na \textbf{BBC} (ORG) . No evento, ela conquistou a vitória mais esmagadora na história do concurso até então, com uma pontuação quase três vezes maior que a de seu rival mais próximo, um feito extremamente improvável de ser derrotado pelo sistema de pontuação pós-1974. " \textbf{Non ho l'età} (MISC) " figurou no \textbf{top20} (MISC) em \textbf{Itália} (LOC) , \textbf{França} (LOC) , \textbf{Reino Unido} (LOC) , \textbf{Alemanha} (LOC) , \textbf{Países Baixos} (LOC) , \textbf{Bélgica} (LOC) e \textbf{Noruega} (LOC) , e com algum sucesso \textbf{em Portugal} (LOC) . A canção também foi gravada em inglês , francês , castelhano , alemão e japonês .

Destaque também para o título da música alemã, ainda é o mais longo título gravado na \textbf{Eurovisão} (LOC): 34 letras no total.



\textbf{Festival} (PER)A ordem de votação no \textbf{Festival Eurovisão da Canção 1964} (MISC), foi a seguinte:

Os países que receberam 5 pontos foram os seguintes:



Transmissão

Os canais de televisão responsáveis pela difussão do concurso quer via televisão, quer via rádio foram as seguintes cadeias televisivas:



Festival Eurovisão da Canção 1965



O \textbf{Festival Eurovisão da Canção 1965} (MISC) (em inglês : \textbf{Eurovision Song Contest} (ORG) 1965 , em francês : \textbf{Concours Eurovision} (MISC) de la chanson 1965 e em italiana : \textbf{Gran Premio Eurovisione della Canzone 1965} (MISC) ) foi o 10 Festival Eurovisão da Canção e realizou-se em 20 de março de 1965 em \textbf{Nápoles} (LOC) , na \textbf{Itália} (LOC) . \textbf{Renata Mauro} (PER) foi apresentadora do evento que foi ganho pela cantora francesa \textbf{France Gall} (PER) que representou o \textbf{Luxemburgo} (LOC), com a canção " \textbf{Poupée} (MISC) de cire, poupée de son ", tornando-se a primeira canção pop a vencer o concurso e a primeira canção vencedora desde 1959 a não ser uma balada. Participaram 18 países, o maior número até então.

Nesta edição, a \textbf{Suécia} (LOC) tornou-se o primeiro país a apresentar uma música cantada na totalidade num idioma não oficial no seu país (o inglês ), começando aqui a polémica sobre a língua a usar no festival.

Pela primeira vez, o \textbf{Festival Eurovisão da Canção} (MISC) transmitido não só pela \textbf{União Europeia de Radiodifusão} (ORG) , como também pela \textbf{Intervisão} (LOC) , seu homólogo da \textbf{Europa} (LOC) de Leste , em países como a \textbf{Checoslováquia} (LOC) , \textbf{Alemanha de Leste} (LOC) , \textbf{Polónia} (LOC) , \textbf{Hungria} (LOC) , \textbf{Roménia} (LOC) e \textbf{União Soviética} (ORG) .



\textbf{Local} (MISC)O \textbf{Festival Eurovisão da Canção 1965} (MISC) ocorreu em \textbf{Nápoles} (LOC) , \textbf{Itália} (LOC) . \textbf{Nápoles} (LOC) é uma comuna do sul de \textbf{Itália} (LOC) , da região da \textbf{Campânia} (LOC) , província de \textbf{Nápoles} (LOC) , com cerca de 1 000 000 habitantes (\textbf{cens} (LOC). 2001) e com cerca de 4 400 000 habitantes na região metropolitana (que compreende áreas na \textbf{província de Caserta} (LOC) , \textbf{Avellino} (LOC) e \textbf{Salerno} (LOC) ). \textbf{Nápoles} (LOC) é a terceira cidade mais populosa da \textbf{Itália} (LOC) após \textbf{Roma} (LOC) e \textbf{Milão} (LOC) e tem a segunda ou terceira maior (dependendo dos dados) região metropolitana do país. Estende-se por uma área de 117 km , de densidade populacional 849 hab./km. É conhecida mundialmente pela sua história, sua música, seus encantos naturais e por ser a terra natal da pizza . O centro histórico de \textbf{Nápoles} (LOC) é \textbf{Património Mundial} (LOC) da \textbf{UNESCO} (ORG) .

O festival em si realizou-se na \textbf{Sala di Concerto della RAI} (LOC) , inaugurado poucos anos antes desta edição. Localizado em \textbf{Viale Marconi} (LOC), no distrito de \textbf{Fuorigrotta} (LOC), a estrutura tem três estúdios de televisão, ocupando um total de 1227 m e com uma capacidade para 370 pessoas. O arquivo napolitano de músicas encontra-se aqui.



\textbf{Formato} (PER) e visual

A apresentadora desta edição foi \textbf{Renata Mauro} (PER) , que apresentou o festival maioritariamente em italiano , mas também em inglês e francês . Ela mesma apresentou os participantes, explicando o significado de suas músicas.

O palco foi muito simples, com a orquestra foi colocada no centro do palco, os artistas à esquerda e o quadro de votação à direita. Atrás dos artistas, foi colocado em segundo plano uma ampliação da sigla \textbf{Eurovisão} (LOC) .

A competição abriu em um close-up do órgão do local. O plano, ampliando, revelou a orquestra, dirigida por \textbf{Gianni Ferrio} (PER) , que então fez sua entrada sob o aplauso do público. A orquestra tocou uma breve introdução musical.

No intervalo, antes das votações, cantou o tenor italiano \textbf{Mario del Monaco} (PER) .

O festival foi visto por 150 milhões de pessoas, estabelecendo um novo recorde.

Pela primeira vez na história da competição, os ensaios foram interrompidos por um incidente ocorrido entre a orquestra e a delegação luxemburguesa. Os músicos não apreciaram a atitude em relação a eles do compositor da música luxemburguesa \textbf{Serge Gainsbourg} (PER) . Alguns compararam sua pontuação ao som do galope de um cavalo e outros vaiaram. \textbf{Gainsbourg} (LOC), furioso, bateu a porta dos ensaios e ameaçou retirar sua música da competição. Um compromisso foi finalmente encontrado, mas persistiu uma certa tensão, que se refletiu na atitude e desempenho da \textbf{France Gall} (ORG) , desestabilizada pelo incidente. Anos mais tarde, \textbf{France Gall} (PER) viria a dizer que a sua vitória foi uma "triste experiência", revelando que o seu namorado de então, quando ganhou, terminou o seu relacionamento com ela, acrescentando que na sua atuação final não chourou lágrimas de alegria, mas lágrimas de tristeza. Também revelou que foi agredida pela representante britânica, \textbf{Kathy Kirby} (PER) , que sentiu que a competição foi falsificada.



\textbf{Votação} (MISC)

O sistema de votação manteve-se. Assim, cada país tinha 10 júris e cada um tinha 9 votos para atribuir. A canção que obtivesse mais votos dentro do júri obtinha 5 pontos, o segundo 3 e a canção em 3 lugar 1 ponto. Se apenas uma canção tivesse todos os votos entre os júris obtinha 9 pontos e se apenas duas canções fossem escolhidas, a canção mais votada obtinha 6 pontos e a outra canção 3 pontos. A vitória da \textbf{Itália} (LOC) nunca foi posta em causa, principalmente após receber a pontuação máxima, 5 pontos, dos primeiros dois países a votar. Os resultados dos votos foram anunciados oralmente, de acordo com a ordem crescente dos votos: 1, 3 e depois 5 votos.

A tabela de votação utilizada foi idêntica à do ano anterior: os resultados foram indicados na forma de setas graduadas e não na forma de números.

Mais uma vez, o supervisor delegado pela \textbf{UER} (ORG) foi \textbf{Miroslav Vilček} (PER) , auxiliado por cinco assistentes. Ele interveio apenas uma vez, pedindo ao porta-voz do júri espanhol para repetir seus votos.

O \textbf{Luxemburgo} (LOC) liderou a votação do começo ao fim, empatando com o \textbf{Mónaco} (LOC) após o segundo país votar.



\textbf{Participantes Para} (MISC) além da já mencionada \textbf{France Gall} (PER) (representando o \textbf{Luxemburgo} (LOC) ), \textbf{Guy Mardel} (PER) (representando a \textbf{França} (LOC) ) \textbf{Udo Jürgens} (PER) (representando a \textbf{Áustria} (LOC) ) e \textbf{Bobby Solo} (PER) (representando a \textbf{Itália} (LOC) ) foram os artistas mais conhecidos desta edição, com este último a ser o único que se fez acompanhar por coristas. \textbf{Simone de Oliveira} (PER) , artista bastante conhecida no mundo lusófono, também esteve presente representando \textbf{Portugal} (LOC) com a canção " \textbf{Sol de Inverno} (MISC) ", conquistando o primeiro ponto para \textbf{Portugal} (LOC) , vindo do \textbf{Mónaco} (LOC) .



\textbf{Festival} (PER)A ordem de votação no \textbf{Festival Eurovisão da Canção 1965} (MISC), foi a seguinte:

Os países que receberam 5 pontos foram os seguintes:



Transmissão

Os canais de televisão responsáveis pela difussão do concurso quer via televisão, quer via rádio foram as seguintes cadeias televisivas:



Festival Eurovisão da Canção 1966



O \textbf{Festival Eurovisão da Canção 1966} (MISC) (em inglês : \textbf{Eurovision Song Contest} (ORG) 1966 e em francês : \textbf{Concours Eurovision} (MISC) de la chanson 1966 ) foi o 11 Festival Eurovisão da Canção e realizou-se a 5 de março de 1966 no \textbf{Luxemburgo} (LOC) . \textbf{Josiane Shen} (PER) foi a apresentadora do evento que foi ganho pelo representante da televisão austríaca ( \textbf{ORF} (ORG) ) \textbf{Udo Jürgens} (PER) , com a canção " \textbf{Merci} (LOC), \textbf{Chérie} (PER) ", que foi escrita pelo cantor em parceria com \textbf{Thomas Hörbiger} (PER) .

A regra que obrigava que cada país cantasse na sua língua oficial foi criada neste ano. Esta regra terá sido criada porque no ano anterior o intérprete sueco cantou o seu tema em inglês e não em sueco , dando origem a contestações de alguns países que se sentiam prejudicados.

A representante dos \textbf{Países Baixos} (LOC) , \textbf{Milly Scott} (PER) , foi a primeira pessoa negra a participar no certame, como também a primeira participante a usar microfone de mão.



\textbf{Local} (MISC)O \textbf{Festival Eurovisão da Canção} (MISC) 1966 ocorreu na capital do \textbf{Luxemburgo} (LOC) . A cidade do \textbf{Luxemburgo} (LOC) é uma comuna de \textbf{Luxemburgo} (LOC) com status de cidade , pertencente ao distrito de \textbf{Luxemburgo} (LOC) e ao cantão de \textbf{Luxemburgo} (LOC) . \textbf{Luxemburgo} (LOC) é uma cidade muito desenvolvida em seu comércio e nas indústrias. \textbf{Luxemburgo} (LOC) é uma das cidades mais ricas da \textbf{Europa} (LOC) e tornou-se um importante centro financeiro e administrativo. A cidade do \textbf{Luxemburgo} (LOC) contém várias instituições da \textbf{União Europeia} (ORG) , incluindo o \textbf{Tribunal de Justiça Europeu} (ORG) , o \textbf{Tribunal de Contas} (ORG) e o \textbf{Banco Europeu de Investimento} (ORG) .

O festival em si realizou-se no \textbf{Villa Louvigny} (LOC) , no sul do \textbf{Luxemburgo} (LOC) , que serviu de sede da \textbf{Compagnie Luxembourgeoise de Télédiffusion} (MISC) , o precursor da \textbf{RTL Group} (MISC) . Está localizado no \textbf{Parque Municipal} (LOC) , no bairro \textbf{Ville Haute} (LOC) do centro da cidade. Já tinha sido a sede da \textbf{Eurovisão} (LOC) em 1962 .



\textbf{Formato Para} (PER) além da obrigatoriedade em que cada país tinha de apresentar a sua canção na sua língua oficial, foi também revogada a regra de que os júris não poderiam incluir profissionais da música.

Como já mencionado, 1966 tornou-se a primeira edição a contar com a presença de uma mulher negra, a representante dos \textbf{Países Baixos} (LOC) , \textbf{Milly Scott} (PER) , que também se tornou a primeira a usar um microfone portátil. Esta edição marcou a última presença da \textbf{Dinamarca} (LOC) no festival, sendo representada por uma canção que na sua atuação contou com a presença de dois bailarinos, pela primeira vez na história. O país voltaria ao certame em 1978 .

Esta foi também a primeira edição em que uma canção não foi acompanhada por uma orquestra. A canção italiana, " \textbf{Dio} (ORG), come ti amo ", interpretada por \textbf{Domenico Modugno} (PER) , foi modificada desde a sua final nacional e quebrou oficialmente a regra da \textbf{UER} (ORG) que afirmava que os arranjos deveriam estar finalizados com bastante antecedência. Durante o ensaio da tarde de sábado, \textbf{Modugno} (PER) realizou o novo arranjo com três de seus próprios músicos, em oposição à orquestra, que ultrapassou o limite de três minutos. Após seu ensaio, \textbf{Modugno} (PER) foi confrontado pelos produtores do programa sobre exceder o limite de tempo e foi solicitado a usar o arranjo original com a orquestra. \textbf{Modugno} (PER) estava tão insatisfeito com a orquestra, que ameaçou desistir do concurso. Tanto os produtores quanto o supervisor-executivo da \textbf{UER} (ORG), \textbf{Clifford Brown} (PER) , sentiram que o tempo era muito curto para pedir a \textbf{Gigliola Cinquetti} (PER) (intérprete original da canção) para se deslocar ao \textbf{Luxemburgo} (LOC) para representar a \textbf{Itália} (LOC) . Então, a \textbf{UER} (ORG) cedeu e permitiu que \textbf{Modugno} (PER) usasse seu próprio conjunto em vez da orquestra. Apesar dos sites e do programa oficial listarem \textbf{Angelo Giacomazzi} (PER) como o maestro, \textbf{Giacomazzi} (PER) na verdade tocou o piano.



\textbf{Visual} (MISC)O vídeo introdutório mostrava diferentes visões do \textbf{Luxemburgo} (LOC) à noite, terminando com uma vista exterior da \textbf{Villa Louvigny} (LOC) . A orquestra tocou a partitura da música vencedora do ano anterior, " \textbf{Poupée} (MISC) de cire, poupée de son ". A câmera então mostrou uma visão da orquestra e concluiu com um close-up em \textbf{France Gall} (PER) , sentada junto do público.

A apresentadora desta edição foi \textbf{Josiane Shen} (PER) , que apresentou o festival maioritariamente em francês , mas também usando, esporadicamente, o inglês . Conclui sua primeira intervenção com as palavras: "\textbf{Tenho} (LOC) a honra de declarar aberta a \textbf{Grand Prix Eurovision de la Chanson} (MISC) 1966".

A orquestra foi colocada à esquerda do palco e o quadro de votação à direita. No centro foi reservado espaço para artistas. Ela consistia de uma escada, parcialmente oculta, e, no fundo, objetos móveis, inspirado na arte pop . Foi a primeira vez na história do concurso que o cenário permaneceu em movimento durante as apresentações dos artistas.

No intervalo, atuou o grupo de jazz \textbf{Les Haricots Rouges} (PER) .



\textbf{Votação} (MISC)

O sistema de votação manteve-se. Assim, cada país tinha 10 júris e cada um tinha 9 votos para atribuir. A canção que obtivesse mais votos dentro do júri obtinha 5 pontos, o segundo 3 e a canção em 3 lugar 1 ponto. Se apenas uma canção tivesse todos os votos entre os júris obtinha 9 pontos e se apenas duas canções fossem escolhidas, a canção mais votada obtinha 6 pontos e a outra canção 3 pontos. A vitória da \textbf{Itália} (LOC) nunca foi posta em causa, principalmente após receber a pontuação máxima, 5 pontos, dos primeiros dois países a votar. Os resultados dos votos foram anunciados oralmente, de acordo com a ordem crescente dos votos: 1, 3 e depois 5 votos.

O supervisor delegado pela \textbf{UER} (ORG) foi, pela primeira vez, \textbf{Clifford Brown} (PER) .

A tabela de votação na forma de setas graduadas foi abandonada, retornando a com números.

\textbf{Josiane Chen} (PER) e o porta-voz britânico protagonizaram um dos momentos mais divertidos da edição. A apresentadora saudou o \textbf{Reino Unido} (LOC) com as palavras "\textbf{Good Night} (MISC), \textbf{London} (PER)", corrigindo imediatamente para "\textbf{Good Evening} (MISC), \textbf{London} (ORG)". Mas o porta-voz respondeu com "\textbf{Good Morning} (MISC), \textbf{Luxembourg} (LOC)", despertando risos da plateia e da própria apresentadora.

Pela primeira vez, foi bem evidente os votos em bloco que viriam a caracterizar todas as edições até hoje, com a troca de pontos entre os países escandinavos e os países ibéricos . \textbf{Lill Lindfors} (PER) , representante da \textbf{Suécia} (LOC) , afirmou anos depois que acreditava que a troca de votos entre os países escandinavos foi uma forma de combater os países de língua francesa .

No inicio da repetição da canção vencedora, \textbf{Udo Jürgens} (PER) disse "\textbf{Merci} (LOC), \textbf{Jury} (LOC)" (Obrigado, júri).



\textbf{Participações} (MISC)O regresso dos cantores de sucesso internacional, como \textbf{Udo Jürgens} (PER) , representando a \textbf{Áustria} (LOC) e \textbf{Domenico Modugno} (PER) , representante da \textbf{Itália} (LOC) . A canção italiana era originalmente cantada por \textbf{Gigliola Cinquetti} (PER) , mas \textbf{Domenico Modugno} (PER) (seu autor e compositor) não permitiu que a canção fosse ao festival caso não fosse ele a interpreta-la. Apesar de se ter posicionado em último lugar, esta música foi gravada em inúmeros países, tornando-se um grande sucesso. Também a representante da \textbf{Noruega} (LOC) , \textbf{Åse Kleveland} (PER) , marcou este festival por ter sido a primeira mulher a atuar na \textbf{Eurovisão} (LOC) sem vestido. Ela viria a apresentar a \textbf{Eurovisão 20} (MISC) anos mais tarde. Outra futura apresentadora concorrente nesta edição foi \textbf{Lill Lindfors} (PER) , que representou a \textbf{Suécia} (LOC) , juntamente com \textbf{Svante Thuresson} (PER) , que viria a apresentar a edição de 1985 .

No \textbf{Mónaco} (LOC) , durante a seleção interba, a já famosa artista francesa \textbf{Isabelle Aubret} (PER) e a artista croata \textbf{Tereza Kesovija} (PER) ficaram empatadas em primeiro lugar. Foi a \textbf{princesa Grace} (PER) , ela própria membro do júri, que teve o voto decisivo. Ela explicou que as duas músicas eram de igual qualidade e ela preferiu dar uma chance a um artista iniciante.

\textbf{Dominique Walter} (PER) é filho de \textbf{Michèle Arnaud} (PER) , que representou o \textbf{Luxemburgo} (LOC) em 1956 e curiosamente fez parte do júri que selecionou a canção francesa.

O representante do \textbf{Reino Unido} (LOC) , \textbf{Kenneth McKellar} (PER) , foi o primeiro artista a apresentar-se em palco com um kilt .



\textbf{Festival} (PER)A ordem de votação no \textbf{Festival Eurovisão da Canção 1966} (MISC), foi a seguinte:

Os países que receberam 5 pontos foram os seguintes:



Transmissão

Os canais de televisão responsáveis pela difussão do concurso quer via televisão, quer via rádio foram as seguintes cadeias televisivas:



Festival Eurovisão da Canção 1967



\textbf{Países Participantes Países} (LOC) que participaram no passado, mas não em 1967

O \textbf{Festival Eurovisão da Canção} (MISC) 1967 (em inglês : \textbf{Eurovision Song Contest} (ORG) 1967 , em francês : \textbf{Concours Eurovision} (MISC) de la chanson 1967 e em alemão : \textbf{Großer Preis der Eurovision} (PER) 1967 ) foi o 12 Festival Eurovisão da Canção e realizou-se a \textbf{8 de Abril de 1967} (ORG) , em \textbf{Viena} (LOC) , \textbf{Áustria} (LOC) . \textbf{Erica Vaal} (PER) foi a apresentadora do evento que foi ganho por \textbf{Sandie Shaw} (PER) que representou o \textbf{Reino Unido} (LOC) com a canção " \textbf{Puppet on a String} (MISC) ". A canção de \textbf{Luxemburgo} (PER), " \textbf{L'amour est bleu} (MISC) ", cantada por \textbf{Vicky Leandros} (PER) , ficou em quarto lugar, embora ainda se tornou o maior sucesso internacional da \textbf{Eurovisão} (LOC) de 1967, e um ano depois seria um grande sucesso instrumental para o músico francês, \textbf{Paul Mauriat} (PER), sob a versão \textbf{em Inglês} (MISC), "\textbf{Love is Blue} (MISC)".

\textbf{Portugal} (LOC) enviou a canção " O vento mudou " que foi interpretada por \textbf{Eduardo Nascimento} (PER) , o primeiro cantor negro na história do \textbf{Festival Eurovisão da Canção} (MISC) , a primeira cantora negra fora \textbf{Milly Scott} (PER) , curiosamente no ano anterior, 1966 . Surgiram imediatamente rumores de que a escolha de \textbf{Eduardo Nascimento} (PER) para representar \textbf{Portugal} (LOC) era uma forma de \textbf{Salazar} (PER) mostrar à \textbf{Europa} (LOC) que não era racista, mostrando também que \textbf{Portugal} (LOC) era um país multirracial.

A \textbf{Dinamarca} (LOC) optou por não participar e deixou a competição, para retornar em 1978. A razão foi que o novo diretor para o departamento de entretenimento de TV da \textbf{DR} (ORG) pensava que o dinheiro poderia ser gasto de uma maneira melhor.



\textbf{Local} (MISC)O \textbf{Festival Eurovisão da Canção} (MISC) 1967 ocorreu em \textbf{Viena} (LOC) , na \textbf{Áustria} (LOC) . \textbf{Viena} (LOC) é a capital da \textbf{Áustria} (LOC) , centro cultural e político do país. É também um dos nove estados ( \textbf{Bundesland Wien} (ORG) ); com 1.681.469 habitantes, era em 2008 o mais populoso deles, ainda que seus 414 km façam dele o menor, sendo também a maior cidade sobre o \textbf{rio Danúbio} (LOC) . \textbf{Viena} (LOC) é cercada pelo \textbf{Estado da Baixa Áustria} (LOC) . A sua aglomeração urbana tem 2,3 milhões de habitantes. Segundo a pesquisa "\textbf{Qualidade de Vida} (MISC) no \textbf{Mundo 2007} (MISC)", realizada pela consultoria de recursos humanos \textbf{Mercer} (LOC) , \textbf{Viena} (LOC) é a melhor cidade do mundo para viver-se, seguida por: \textbf{Zurique} (LOC) , \textbf{Genebra} (LOC) , \textbf{Vancouver} (LOC) e \textbf{Auckland} (LOC) .

O festival em si realizou-se no \textbf{Großer Festsaal} (LOC) der \textbf{Wiener Hofburg} (PER) . O \textbf{Palácio de Hofburg} (LOC) foi construído entre o século XIII e o século XIX , e foi a residência vienense dos \textbf{Habsburgos} (PER) . Atualmente, serve como residência oficial do \textbf{Presidente} (MISC) da \textbf{Áustria} (LOC) .



\textbf{Formato} (MISC)A configuração do palco foi um pouco incomum este ano. A orquestra foi colocada na parte esquerda, os artistas na parte direita. O pano de fundo atrás dos artistas consistia em três espelhos giratórios que começavam a rodar no início de cada música e paravam no final de cada música. O quadro de votação foi colocado à direita do palco.

250 milhões de espetadores assistiram a esta edição, com a então província de \textbf{Angola} (LOC) a acompanhar atentamente o festival. Na falta de televisores, as estações de rádio locais passaram todas as músicas da \textbf{Eurovisão} (LOC), em disco, pela ordem em que atuaram em \textbf{Viena} (LOC). Depois da votação, as rádios depressa deram a notícia da classificação..



\textbf{Visual} (MISC)O vídeo introdutório começou com um tiro fixo num dos grandes lustres da prefeitura, antes de desvendar o palco e a orquestra, que começou a tocar várias músicas conhecidas. Depois, o maestro \textbf{Johannes Fehring} (PER) , passou para \textbf{Udo Jürgens} (PER) , que dirigiu a orquestra com o instrumental de " \textbf{Merci} (LOC), \textbf{Chérie} (PER) ", organizado em valsa.

\textbf{Erica Vaal} (MISC) abriu o festival apresentando em 7 línguas: alemão , francês , inglês , italiano , castelhano , russo e polaco . Conclui a sua mensagem com as palavras: "Que são as melhores e mais belas canções que ganham, este é o nosso desejo!", arrependendo-se também de não ter tido tempo de aprender a língua de outros países, mas certamente a encontraria quando a competição fosse novamente organizada em \textbf{Viena} (LOC) .

Pela primeira vez, legendas foram usadas: os títulos das músicas foram legendados em alemão , inglês e francês .

Para o intervalo, a \textbf{ÖRF} (ORG) decidiu oferecer o que \textbf{Viena} (LOC) poderia fazer de melhor: interpretar \textbf{Johann Strauss} (PER) , o rei da valsa que, na voz \textbf{Coral dos Meninos de Viena} (LOC) , interpretam a mais famosa das suas valsas, o \textbf{Danúbio Azul} (ORG) , mas também \textbf{Land der Berge} (MISC), \textbf{Land am Strome} (MISC) , o hino nacional da \textbf{Áustria} (LOC) .

A apresentadora, \textbf{Erika Vaal} (PER) terminou o programa por felicitar a vencedora, e despedindo-se em várias línguas, dizendo que "A música não conhece fronteiras". Esta foi a última edição a ser transmitida em preto e branco.



\textbf{Votação} (MISC)

Mais uma vez os países votaram pela mesma ordem à da sua actuação, algo que foi feito aparentemente para acelarar o espetáculo, pois os países a actuar não podiam (nem podem) atribuir pontos a si mesmo, sendo assim poderiam acabar a sua votação a tempo de serem os primeiros a serem chamados. O sistema de votação foi o mesmo utilizado em edições anteriores, em que cada país tinha um júri composto por dez membros, e cada um desses dez jurados atribuiam um ponto à sua canção preferida.

O supervisor delegado pela \textbf{UER} (ORG) foi, mais uma vez, \textbf{Clifford Brown} (PER) , que teve que intervir quatro vezes para que os votos dos \textbf{júris monegasco} (PER), jugoslavo, italiano e irlandês fossem corrigidos.

Pela primeira vez, imagens da sala que depois se tornaria conhecida como \textbf{green room} (ORG) foram exibidas, com a aparição de vários artistas, como \textbf{Sandie Shaw} (PER) , \textbf{Therese Steinmetz} (ORG) , \textbf{Peter Horten} (PER) a assinar um autógrafo para um jovem corista e \textbf{Kirsti Sparboe} (PER), com um dos \textbf{Meninos de Viena} (LOC) ao colo.

A apresentadora confundiu-se durante a votação para declarar vitória britânica antes de ouvir os votos do último júri. \textbf{Clifford Brown} (PER) teve que intervir pela quinta e última vez: "Ainda esperando pelo voto irlandês". \textbf{Erica Vaal} (PER) pediu desculpas e ligou para \textbf{Dublin} (LOC) . O porta-voz do júri irlandês respondeu às suas saudações com estas palavras: "Eu pensei que íamos ficar de fora". A plateia então aplaudiu calorosamente. Mas no final o resultado não se alterou após a última votação e o \textbf{Reino Unido} (LOC) acabou por se sagrar campeão, com mais do dobro dos pontos que a \textbf{Irlanda} (LOC) , o segundo classificado. Em adição, o quadro de votações avariou-se durante a votação do voto de \textbf{Mónaco} (LOC) , e as pontuações começaram a ser exibidas incorretamente, depois corrigidas.



\textbf{Participantes} (MISC)A música luxemburguesa, interpretada por \textbf{Vicky Leandros} (PER) nos seus primórdios, seria a música com mais sucesso. O mais famoso é a versão instrumental da orquestra de \textbf{Paul Mauriat} (PER) , que chegou a número um nos \textbf{Estados Unidos} (LOC) .

O representante português, \textbf{Eduardo Nascimento} (PER) , como já referido, foi o primeiro artista de origem angolana e o primeiro artista negro a participar no concurso.

A representante britânica, \textbf{Sandie Shaw} (PER) , foi a primeiro artista na história da competição a cantar descalça no palco. \textbf{Shaw} (PER) viria a revelar que não gostava da música que fez dela a vencedora da edição.

A participação monegasca marcaria o regresso de \textbf{Serge Gainsbourg} (PER) à \textbf{Eurovisão} (LOC), depois de ganhar pelo \textbf{Luxemburgo} (LOC) em 1965 com " \textbf{Poupée} (MISC) de cire, poupée de son ". O maestro monegasco, \textbf{Aimé Barelli} (PER) , não era mais do que o pai do intérprete, \textbf{Minouche Barelli} (PER) .



\textbf{Festival} (PER)A ordem de votação no \textbf{Festival Eurovisão da Canção} (MISC) 1967, foi a seguinte:

Em baixo encontra-se a lista de maestros que conduziram a orquestra, na respectiva actuação de cada país concorrente.



Transmissão

Os canais de televisão responsáveis pela difussão do concurso quer via televisão, quer via rádio foram as seguintes cadeias televisivas:



Festival Eurovisão da Canção 1968



\textbf{Países Participantes Países} (LOC) que participaram no passado, mas não em 1968

O \textbf{Festival Eurovisão da Canção 1968} (MISC) (em inglês : \textbf{Eurovision Song Contest} (ORG) 1968 e em francês : \textbf{Concours Eurovision} (MISC) de la chanson 1968 ) foi o 13 \textbf{Festival Eurovisão da Canção} (MISC) e realizou-se a 6 de abril de 1968 , em \textbf{Londres} (LOC) , no \textbf{Reino Unido} (LOC) . \textbf{Katie Boyle} (PER) foi a apresentadora do evento, pela terceira vez na história (já tinha apresentado as edições de 1960 e de 1963 ), que foi ganho pela canção espanhola " La, la, la " por \textbf{Massiel} (PER) . Originalmente, a \textbf{Espanha} (LOC) seria representada por \textbf{Joan Manuel Serrat} (PER) , mas o seu desejo em cantar em catalão (a língua proibida de \textbf{Francisco Franco} (PER) ) foi uma afronta ao regime ditatorial franquista que jamais permitiria semelhante ousadia. \textbf{Serrat} (LOC) acabaria por ser substituído por \textbf{Massiel} (PER) que cantou a mesma canção em castelhano . \textbf{Massiel} (PER), que estava em turnê no \textbf{México} (LOC) , só teve 10 dias para se preparar antes do concurso.



\textbf{Local} (MISC)O \textbf{Festival Eurovisão da Canção} (MISC) 1968 ocorreu em \textbf{Londres} (LOC) , no \textbf{Reino Unido} (LOC) , que já tinha sido sede das edições de 1960 e de 1963 . \textbf{Londres} (LOC) é uma das maiores cidades da \textbf{Europa} (LOC) e já liderou a lista das maiores cidades do mundo . De uma cidade romana chamada \textbf{Londínio} (LOC) , a capital da região administrativa romana \textbf{Britannia} (LOC), a cidade acabou por se tornar o centro do império britânico, contribuindo hoje com 17\% do PNB de um país que é a quarta maior economia do mundo. \textbf{Londres} (LOC) tem sido um dos mais importantes centros da política e do comércio mundiais por quase dois milénios (apesar de a capital de \textbf{Inglaterra} (LOC) ter sido \textbf{Winchester} (LOC) durante grande parte da \textbf{Idade Média} (MISC) ).

O festival em si realizou-se no \textbf{Royal Albert Hall} (LOC) , uma sala de concertos situado na \textbf{Cidade de Westminster} (LOC) em \textbf{Londres} (LOC) , a capital e a maior cidade da \textbf{Inglaterra} (LOC) e do \textbf{Reino Unido} (LOC) . O \textbf{Royal Albert Hall} (LOC) é conhecido por ali terem lugar várias exposições de vários artistas de vários géneros, desportos, cerimónias de entrega de prémios, concertos diversos e outros eventos desde a sua abertura em 1871 e tornou-se um dos mais importantes edifícios de \textbf{Londres} (LOC) .



\textbf{Formato} (PER)

1968 foi a primeira vez em que o \textbf{Festival Eurovisão da Canção} (MISC) foi transmitido a cores. Os países que o emitiram a cores foram a \textbf{França} (LOC) (através da \textbf{Antenne} (LOC) 2 ), a então \textbf{Alemanha Ocidental} (LOC) , \textbf{Países Baixos} (LOC) , a \textbf{Finlândia} (LOC) , a \textbf{Noruega} (LOC) , a \textbf{Suíça} (LOC) , a \textbf{Suécia} (LOC) e o \textbf{Reino Unido} (LOC) . A título de curiosidade, em \textbf{Portugal} (LOC) só no dia 7 de março de 1980 estrearam as emissões regulares a cores, com a transmissão do \textbf{Festival RTP da Canção} (MISC) desse ano. Também todos os países de leste e a \textbf{Tunísia} (LOC) transmitiram o festival. \textbf{Katie Boyle} (PER) apresentou o festival pela terceira vez.

No primeiro dia de ensaios, um trio de albaneses invadiu o gabinete do chefe de produção da \textbf{BBC} (ORG) e disseram: "somos albaneses e queremos apresentar a nossa canção ao \textbf{Festival} (PER)", proposta que foi de imediato recusada. De repente, começaram a tocar " \textbf{Congratulations} (MISC) "" em versão jazz. Era apenas uma brincadeira, ou uma jogada de marketing, de alguns responsáveis da \textbf{BBC} (ORG) .

Pela \textbf{Jugoslávia} (ORG) , o grupo \textbf{Dubrovački Trubaduri} (LOC) , apareceu em trajes de jograis medievais . Foi um grupo originalmente formado por seis pessoas. No entanto, o regulamento só permitia que solistas e duplas competissem, então oficialmente cinco foram para \textbf{Londres} (LOC), dois deles intérpretes oficiais, \textbf{Luči Kapurso} (MISC) \& \textbf{Hamo Hajdarhodžić} (PER) , e os outros três como coristas.

Como acontecia de ano para ano, o número de espetadores aumentou: cerca de 300 milhões de pessoas viram o espetáculo, não só nos países concorrentes como também no leste europeu e na \textbf{América} (LOC) , onde o certame passou em diferido. \textbf{Em Portugal} (LOC), os cafés tornaram a encher para assistir à final.

Em maio de 2008, um documentário do realizador espanhol \textbf{Montse Fernández Villa} (PER) 1968. Yo viví el mayo español ("Eu vivi o maio espanhol)", centrado nos efeitos do maio de 1968 na \textbf{Espanha} (LOC) franquista alegou que o \textbf{Festival Eurovisão da Canção 1968} (MISC) foi fraudado pelo ditador espanhol , o \textbf{General Francisco Franco} (PER) , que teria enviado funcionários por toda a \textbf{Europa} (LOC) oferecendo dinheiro e prometendo comprar séries de televisão e contratar artistas desconhecidos. A alegação foi baseada no testemunho do jornalista \textbf{José María Íñigo} (PER) , um funcionário da \textbf{TVE} (ORG) na época que reclamou que a fraude era de conhecimento geral e sugeriu que a gravadora espanhola da canção de \textbf{Massiel} (PER) oferecia-se para lançar álbuns de artistas da então \textbf{Checoslováquia} (LOC) e da \textbf{Bulgária} (LOC) (países que na época não faziam parte da \textbf{União Europeia de Radiodifusão} (ORG) ).

O documentário considera que o festival deveria ter sido ganho pela canção do britânica - " \textbf{Congratulations} (MISC) " interpretada por \textbf{Cliff Richard} (PER) - que terminou em segundo lugar a apenas um ponto de distância. \textbf{Massiel} (PER), a intérprete da canção vencedora, sentiu-se ofendida pelas acusações e disse que "outros cantores que eram mais afiados pelo regime franquista teriam sido beneficiados". \textbf{José María Iñigo} (PER), autor da declaração do documentário, defendeu pessoalmente \textbf{Massiel} (PER) e disse que ele tinha espalhado o rumor. Tanto \textbf{Massiel} (PER) e \textbf{Iñigo} (LOC) acusaram o canal \textbf{La Sexta} (MISC) , emissor do documentário, de fabricar o escândalo.



\textbf{Visual} (MISC)O vídeo introdutório começou com uma visão do palco, que se ampliou, revelando a orquestra e o público. A orquestra tocou um cover da música vencedora do ano anterior, " \textbf{Puppet on a String} (MISC) ".

A apresentadora da edição foi novamente \textbf{Katie Boyle} (PER) , que apresentava o certame pela terceira vez, depois de 1960 e de 1963 . Mais tarde, ela contou que teve de enfrentar a edição toda, pois o seu casamento estava a passar por uma fase difícil e que tinha chorado tanto que no dia do certame ela teve que deitar-se várias horas antes com batatas cruas sobre os olhos.

A orquestra foi instalada num buraco ao pé do palco, que apresentava três pódios separados, um no centro para os artistas, um à esquerda para os coristas e outro à direita para a apresentadora. Os artistas fizeram a sua apresentação pela parte de trás do palco, passando por uma ampliação metálica e dourada do logótipo da \textbf{Eurovisão} (LOC), que serviu de fundo. Atrás do logótipo, uma tela mostrava um retrato dos artistas, antes de tomar tons azulados. A mesa do supervisor executivo foi posto à direita do palco. O palco estava cercado por cortinas azuis claras. \textbf{Katie} (PER) apresentou as canções sem aparecer em frente às câmaras.

O intervalo foi preenchido por um vídeo intitulado "\textbf{Impressions from London} (MISC)". Consistia em vários pontos turísticos da capital britânica, dia e noite, como o \textbf{Palácio de Westminster} (LOC) , o \textbf{Big Ben} (LOC) , a \textbf{Tower Bridge} (LOC) , a \textbf{Torre de Londres} (LOC) , o \textbf{Palácio de Buckingham} (LOC) , a \textbf{Hyde Park} (LOC) , a \textbf{Catedral de São Paulo} (LOC) , a \textbf{Abadia de Westminster} (LOC) , a \textbf{Trafalgar Square} (LOC) e a \textbf{Piccadilly Circus} (MISC) . O acompanhamento musical foi tocado ao vivo pela orquestra.



\textbf{Votação} (MISC)

Mais uma vez os países votaram pela mesma ordem à da sua actuação, algo que foi feito aparentemente para acelarar o espetáculo, pois os países a actuar não podiam (nem podem) atribuir pontos a si mesmo, sendo assim poderiam acabar a sua votação a tempo de serem os primeiros a serem chamados. O sistema de votação foi o mesmo utilizado em edições anteriores, em que cada país tinha um júri composto por dez membros, e cada um desses dez jurados atribuiam um ponto à sua canção preferida.

O supervisor delegado pela \textbf{UER} (ORG) foi, mais uma vez, \textbf{Clifford Brown} (PER) , que teve que intervir, no final da votação.

A votação começou com a \textbf{França} (LOC) na liderança, antes de ser ultrapassada pelo \textbf{Reino Unido} (LOC) , sob os aplausos do público. A \textbf{Espanha} (LOC) gradualmente diminuiu a diferença e ganhou a vitória, depois de um final surpreendente. O penúltimo júri, o júri alemão, concedeu dois votos ao \textbf{Reino Unido} (LOC) e seis na \textbf{Espanha} (LOC). A \textbf{Espanha} (LOC) então totalizava 29 votos e o \textbf{Reino Unido} (LOC) 28. Mas o último júri, o júri jugoslavo, não atribuiu nenhum ponto a esses dois países. Houve gritos de surpresa no público. Naquele momento, \textbf{Clifford Brown} (PER) pediu a \textbf{Katie Boyle} (PER) que se lembrasse do júri jugoslavo que havia concedido uma votação supranumerária. Depois de uma comunicação confusa com \textbf{Skopje} (LOC) , os resultados foram validados pelo supervisor e \textbf{Espanha} (LOC) foi proclamada vencedora.

Pela primeira vez desde 1961, nenhum país terminou a votação com "\textbf{null point} (LOC)".



\textbf{Participantes} (MISC)O processo de seleção de participantes foi realizado nos meses anteriores ao festival. Dos 17 países participantes, 10 organizaram uma final nacional e 7 elegeram seu representante internamente.

Entre os representantes mais conhecidos estavam \textbf{Wencke Myhre} (PER) , cantora norueguesa que tentou representar seu país várias vezes e foi finalmente escolhida para representar a \textbf{Alemanha Ocidental} (LOC) ; \textbf{Isabelle Aubret} (PER) , vencedora do Festival Eurovisão da Canção 1962 ; \textbf{Sergio Endrigo} (PER) , vencedor do \textbf{Festival de Sanremo} (PER) naquele ano; \textbf{Karel Gott} (PER) , cantor checoslovaco que representou a \textbf{Áustria} (LOC) ; e \textbf{Cliff Richard} (PER) , o grande favorito para ganhar este festival de acordo com uma grande parte da imprensa europeia e que alcançaria uma importante projeção internacional. Uma ainda não muito conhecida \textbf{Massiel} (PER) , embora já tivesse dois sucessos no \textbf{México} (LOC) e participações em outros festivais europeus, se consolidou na \textbf{Espanha} (LOC) e na \textbf{América Latina} (LOC) após sua vitória na \textbf{Eurovisão} (LOC).

A canção representante do \textbf{Mónaco} (LOC) , " À chacun sa chanson ", participara na final nacional francesa.



\textbf{Festival} (PER)A ordem de votação no \textbf{Festival Eurovisão da Canção 1968} (MISC), foi a seguinte:

Em baixo encontra-se a lista de maestros que conduziram a orquestra, na respectiva actuação de cada país concorrente.



Festival Eurovisão da Canção 1969



\textbf{Países Participantes Países} (LOC) que participaram no passado, mas não em 1969

O \textbf{Festival Eurovisão da Canção} (MISC) 1969 (em inglês : \textbf{Eurovision Song Contest} (ORG) 1969 , em francês : \textbf{Concours Eurovision} (MISC) de la chanson 1969 e em castelhano : \textbf{Festival de la Canción de Eurovisión} (MISC) 1969 ) foi o 14 \textbf{Festival Eurovisão da Canção} (MISC) e realizou-se em 29 de março de 1969 em \textbf{Madrid} (LOC) . \textbf{Laurita Valenzuela} (PER) foi a apresentadora do evento. O \textbf{Festival} (LOC) terminou com um empate entre quatro intérpretes: \textbf{Salomé} (PER) que representou a \textbf{Espanha} (LOC) com a canção " \textbf{Vivo Cantando} (MISC) ", \textbf{Lulu} (PER) que representou o \textbf{Reino Unido} (LOC) , com a canção " \textbf{Boom Bang-a-Bang} (MISC) ", \textbf{Lenny Kuhr} (PER) que representou os \textbf{Países Baixos} (LOC) com a canção " De troubadour " e \textbf{Frida Boccara} (PER) que representou a \textbf{França} (LOC) com a canção " \textbf{Un jour} (MISC), un enfant ".

\textbf{Portugal} (LOC) , representado por \textbf{Simone de Oliveira} (PER) , com a música " \textbf{Desfolhada Portuguesa} (MISC) ", poema de \textbf{Ary dos Santos} (PER) , foi considerada tantos pelos jornalistas como pelos seus pares, como uma candidata à vitória. Contudo, para além do penúltimo lugar conquistado, no dia anterior à edição, \textbf{Simone} (PER) e uma das coristas foram hospitalizadas por intoxicação alimentar.

Este foi o primeiro \textbf{Festival Eurovisão da Canção} (MISC) em que ocorreu um empate no 1 lugar, com 4 países obtendo 18 pontos. O problema foi que na época não havia uma regra em caso de empate, e então os quatro países foram declarados vencedores, provocando uma onda de contestação, com vários países a desistirem na edição seguinte.

A \textbf{Áustria} (LOC) não participou neste \textbf{Festival} (PER), porque não desejava enviar um cantor a um país que era governado por um ditador . \textbf{Francisco Franco} (PER) , que governava a \textbf{Espanha} (LOC) na época. \textbf{Salvador Dalí} (PER) foi responsável pela publicidade deste festival.

O \textbf{Liechtenstein} (LOC) desejou participar neste festival e escolheu mesmo uma canção " \textbf{Un beau matin} (MISC) " (Uma bela manhã), mas \textbf{Liechtenstein} (LOC) não tinha uma televisão pública e ainda por cima não era membro da \textbf{UER} (ORG) . Desta forma foi impedido de participar no evento. No entanto, tudo isso não passou de um rumor criado pelo animador e comediante francês, \textbf{Jacques Martin} (PER) . O \textbf{País de Gales} (LOC) tinha intenção de participar, mas foi-lhe recusado, devido ao facto de fazer parte do \textbf{Reino Unido} (LOC) .

Durante a edição, algumas delegações pediram ao governo espanhol a libertação de certos presos políticos.



\textbf{Local} (MISC)O \textbf{Festival Eurovisão da Canção} (MISC) 1969 ocorreu em \textbf{Madrid} (LOC) , em \textbf{Espanha} (LOC) . \textbf{Madrid} (LOC) é a capital e a maior cidade em \textbf{Espanha} (LOC) , tal como no município de \textbf{Madrid} (LOC) e na \textbf{Comunidade autónoma de Madrid} (ORG) . A cidade foi edificada nas margens do \textbf{rio Manzanares} (LOC) , no centro do país. Devido à sua localização geográfica e histórica, é juntamente com \textbf{Lisboa} (LOC) o centro financeiro e político da \textbf{Península Ibérica} (LOC) . No seguimento da restauração da democracia, em 1975 , e a adesão à \textbf{CEE} (ORG) , em 1985 , a cidade de \textbf{Madrid} (LOC) tem vindo a desempenhar um papel importante na economia europeia, tornando-se num dos principais focos financeiros \textbf{do Sul} (LOC) da \textbf{Europa} (LOC).

O festival em si realizou-se no \textbf{Teatro Real} (LOC) , uma casa de ópera localizado em \textbf{Madrid} (LOC) . O teatro foi reaberto em 1966 como um teatro de concertos e da principal sala de concertos da \textbf{Orquestra Nacional} (ORG) espanhol e da \textbf{Orquestra Sinfónica RTVE} (ORG) .



\textbf{Formato} (PER)O artista surrealista \textbf{Salvador Dalí} (PER) foi o responsável pela acção publicitária que foi feita ao Festival Eurovisão de 1969, como também o desenho da escultura de metal usado no palco.

\textbf{Áustria} (LOC) optou por não participar pois não queria enviar nenhuma representação sua a um país que vivia sobre o domínio de um ditador ( \textbf{Franco} (PER) ). O \textbf{Liechtenstein} (LOC) tentou fazer a sua estreia, mas, como não tinha uma televisão de serviço público, nem era membro da \textbf{UER} (ORG) , não pôde participar. A música escolhida seria " \textbf{Un beau matin} (MISC) ", interpretada por \textbf{Vetty} (PER) . Além disso, a música foi escrita em francês , enquanto a língua nacional do \textbf{Liechtenstein} (LOC) é o alemão . Mais tarde, foi revelado de que tudo não passava de uma piada, montada pelo animador e comediante francês, \textbf{Jacques Martin} (PER) . O \textbf{País de Gales} (LOC) também teve intenção de se estrear, através da \textbf{BBC Cymru} (MISC) , fazendo inclusive uma final nacional, \textbf{Cân i Gymru} (MISC) , mas foi impedida de participar por não ser um país soberano . Apenas a \textbf{BBC} (ORG) tinha o direito exclusivo de representar o \textbf{Reino Unido} (LOC) .

A produção e organização do concurso foi um desafio para a \textbf{TVE} (ORG), que transmitiu o festival a cores, ainda em fase de testes, para vários países. Como a \textbf{TVE} (ORG) não tinha câmaras a cores, teve que alugá-los na televisão pública alemã. Finalmente, os espectadores espanhóis viram o concurso em preto e branco, pois a televisão espanhola ainda não se convertera à cor. O evento foi seguido ao vivo nos países membros da \textbf{UER} (ORG) e \textbf{Intervisão} (LOC), bem como por satélite no \textbf{Chile} (LOC) , \textbf{Porto Rico} (LOC) e, pela primeira vez, no \textbf{Brasil} (LOC) .

Este foi o primeiro \textbf{Festival Eurovisão da Canção} (MISC) em que ocorreu um empate no 1 lugar, com 4 países obtendo 18 pontos. O problema foi que na época não havia uma regra em caso de empate, e então os quatro países foram declarados vencedores. Isto causou um problema com as medalhas entregues aos vencedores não serem suficientes para todos. Apenas as intérpretes receberam medalhas, com os compositores a receber após a edição. Muitos países ficaram irritados com o sistema de votação, levando à sua desistência no ano seguinte, acusando-o de favorecer os grandes países. No caso de \textbf{Espanha} (LOC) , foi a primeira vez que um país ganhou dois anos consecutivos.

\textbf{Madrid} (LOC) animou-se com a \textbf{Eurovisão} (LOC): foram várias as festas, jantares, cocktails, e até visitas a museus preparados para os artistas.



\textbf{Visual} (MISC)O vídeo introdutório começou com um close-up da escultura de metal no palco definitivo criado pelo artista surrealista espanhol, \textbf{Salvador Dalí} (PER) , feita especialmente para a ocasião. O plano se expande para revelar o órgão que tocou o " \textbf{Te Deum} (MISC) " de \textbf{Marc-Antoine Charpentier} (PER) . Depois, a orquestra começou a partitura de " La, la, la ", a música vencedora do ano anterior.

A apresentadora da edição foi \textbf{Laurita Valenzuela} (PER) , que apresentou em castelhano , inglês e francês . Ela cumprimentou todos os países participantes nos seus idiomas e, em seguida, introduziu cada música. A competição foi produzida em cores para a transmissão internacional, enquanto que na televisão espanhola não foi sequer passada para a cor. Os espectadores espanhóis viram este concurso em preto e branco.

A orquestra, instalada num buraco ao pé do palco, foi dirigida por \textbf{Augusto Algueró} (PER) . O palco foi decorado com quatro camadas de flores rosa. No fundo, foram instalados um órgão e a escultura de \textbf{Dali} (PER). Na parede esquerda, havia uma reprodução colorida da sigla \textbf{Eurovisão} (LOC) e na parede à direita, o quadro de votação.

Depois das músicas, \textbf{Laurita} (PER) convidou o público e os espectadores a assistir a um filme filmado para a ocasião, feito por \textbf{Salvador Dalí} (PER) , intitulado "La España diferente". Foi dedicado aos quatro elementos que compõem a \textbf{Espanha} (LOC) : terra , ar , água e fogo . \textbf{Inspirado} (LOC) pelo surrealismo, o filme explorou os temas dos quatro elementos, através de vistas turísticas de paisagens e monumentos espanhóis. \textbf{Alhambra} (LOC) , em \textbf{Granada} (LOC) e a \textbf{Sagrada Família} (LOC) apareceram na tela. O acompanhamento musical foi composto por \textbf{Luis} (PER) de \textbf{Pablo} (PER) .



\textbf{Votação} (MISC)

Mais uma vez os países votaram pela mesma ordem à da sua actuação, algo que foi feito aparentemente para acelerar o espetáculo, pois os países a actuar não podiam (nem podem) atribuir pontos a si mesmos, sendo assim poderiam acabar a sua votação a tempo de serem os primeiros a serem chamados. O sistema de votação foi o mesmo utilizado em edições anteriores, em que cada país tinha um júri composto por dez membros, e cada um desses dez jurados atribuía um ponto à sua canção preferida.

O supervisor delegado pela \textbf{UER} (ORG) foi, mais uma vez, \textbf{Clifford Brown} (PER) , acompanhado por dois secretários, que interveio para pedir aos porta-vozes espanhol e monegasco para repetir os votos. O voto foi efetuado sem mais incidentes, mas no final ocorreu um evento totalmente imprevisto. Nada menos do que quatro países empatados em primeiro lugar. A plateia prendeu a respiração e \textbf{Laurita} (MISC) não sabia quem seria a vencedora. Seguiu-se um diálogo em francês entre a apresentadora e o supervisor, que acabou com \textbf{Brown} (PER) a confirmar a vitória dos quatro países. Não havia disposições no regulamento para decidir sobre qualquer empate.

Se a regra de desempate implementada em 1991 (onde ocorreu uma situação semelhante) tivesse sido aplicada em 1969, os \textbf{Países Baixos} (LOC) seriam os vencedores, pois receberam 6 pontos da \textbf{França} (LOC) , em 2 ficaria o \textbf{Reino Unido} (LOC) , por ter recebido 5 pontos da \textbf{Suécia} (LOC) , em 3 ficaria a \textbf{França} (LOC) , por ter recebido 4 pontos da \textbf{Irlanda} (LOC) e do \textbf{Reino Unido} (LOC) e em 4 \textbf{Espanha} (LOC) . Por outro lado, com a atual regra de desempate em vigor (ou seja, a música que recebeu votos de mais países, então a música que recebeu pontuações mais elevadas em caso de outro empate), a \textbf{França} (LOC) teria sido a vencedora, com \textbf{Espanha} (LOC) em 2 lugar. Ambos os países receberam votos de 9 países, mas a \textbf{França} (LOC) recebeu 4 pontos de 2 países, enquanto a \textbf{Espanha} (LOC) recebeu 3 pontos como seu maior voto. Em 3 ficaria o \textbf{Reino Unido} (LOC) , que foi pontuado por 8 países e em 4 os \textbf{Países Baixos} (LOC) , que foram pontuados por 7 países.

A tabela seguinte apresenta, de forma mais resumida o resultado de acordo com as regras de desempate aplicadas ao longo dos anos.



\textbf{Participantes} (MISC)Entre os representantes mais conhecidos estavam \textbf{Simone de Oliveira} (PER) , que voltaria a representar \textbf{Portugal} (LOC) , com a sua canção de maior sucesso " \textbf{Desfolhada Portuguesa} (MISC) "; \textbf{Salomé} (PER) , representando a \textbf{Espanha} (LOC) , que usava um vestido de 14 quilos feito de porcelana. Além disso a cantora dançou o que quebrava uma das regras da altura, causando uma onda de protestos em inúmeros países; \textbf{Jean Jacques Bortolai} (PER) , que representou o \textbf{Mónaco} (LOC) , tinha apenas 12 anos, tornando-se um dos artistas mais novos a atuar no palco da \textbf{Eurovisão} (LOC), numa altura que não existia nenhuma regra de limite mínimo de idade. \textbf{Muriel} (PER) Day , representando a \textbf{Irlanda} (LOC) , foi a primeira artista feminina e a primeira artista de origem da \textbf{Irlanda do Norte} (LOC) a representar o seu país; \textbf{Lenny Kuhr} (PER) , representado os \textbf{Países Baixos} (LOC) , escreveu sua própria música. A ideia chegara-lhe a meio da noite. Ela então acordou os seus pais para lhes mostras a sua nova composição; \textbf{Kirsti Sparboe} (PER) , representando a \textbf{Noruega} (LOC) pela terceira vez, causou polémica pelo seu vestido feito de pele de baleia; \textbf{Frida Boccara} (PER) , representado a \textbf{França} (LOC) , não pôde comparecer aos ensaios, pois contraíra uma infecção na garganta que foi curada por injecção, horas antes do início do concurso; \textbf{Romuald} (PER) , representando o \textbf{Luxemburgo} (LOC) , já havia representado o \textbf{Mónaco} (LOC) no Festival Eurovisão da Canção 1964 ; finalmente, \textbf{Siw Malmkvist} (ORG) , representando a \textbf{Alemanha Ocidental} (LOC) , já havia representado a \textbf{Suécia} (LOC) no Festival Eurovisão da Canção 1960

A canção representante de \textbf{Portugal} (LOC) , " \textbf{Desfolhada Portuguesa} (MISC) ", interpretada por \textbf{Simone de Oliveira} (PER) , foi uma das músicas que marcou esta edição, sendo mesmo considerada uma das fortes candidatas à vitória, inclusive entre os seus pares. \textbf{Em Portugal} (LOC) , os teatros chegaram mesmo a mudar a hora dos seus espetáculos para que nenhum acontecesse ao mesmo tempo do certame, tal era a expectativa em volta da canção portuguesa. Contudo, \textbf{Portugal} (LOC) acabou a votação em 15 e penúltimo lugar, recebendo apenas 4 pontos. \textbf{Mal} (PER) acabou a votação, a equipa do \textbf{Diário de Notícias} (ORG) foi ao camarim de \textbf{Simone} (PER), onde assistiu a dois momentos impressionantes: o primeiro foi quando o ministro da Informação e Turismo espanhol foi pedir desculpa a \textbf{Simone} (PER); depois foi um jornalista espanhol dizer-lhe que era a vencedora moral. "Não é costume pedir desculpa quando não se cometem erros…", conclui o jornal. A cantora estava calma: "que hei de fazer? \textbf{Sigo} (MISC) cantando", disse ela, que tinha a certeza de ter dado o melhor de si. Com o avançar da conversa, partilhou com o jornalista a desconfiança de ter havido uma irregularidade na votação, porque foi o primeiro empate e logo entre quatro canções. Para além disso, a participação portuguesa esteve em risco, após \textbf{Simone de Oliveira} (PER) e uma das coristas terem ido parar ao hospital no dia anterior por intoxicação alimentar. Chegou-se mesmo a falar em sabotagem. Estas situações provocaram uma grande manifestação de apoio à nossa intérprete na estação de \textbf{Santa Apolónia} (LOC) . Quando \textbf{Simone de Oliveira} (PER) chegou, foi recebida em festa pelos milhares de populares, na maior recepção a um artista português jamais feita \textbf{em Portugal} (LOC), só repetida pela receção a \textbf{Salvador} (LOC) e \textbf{Luísa} (LOC) Sobral 48 anos mais tarde, quando \textbf{Portugal} (LOC) venceu o \textbf{Festival Eurovisão da Canção 2017} (MISC) . Esta edição foi especialmente recordada na série da \textbf{RTP1} (ORG) " \textbf{Conta-me} (ORG) como Foi ", nalguns episódios da 2 temporada.



\textbf{Festival} (PER)A ordem de votação no \textbf{Festival Eurovisão da Canção} (MISC) 1969, foi a seguinte:



Festival Eurovisão da Canção 1970



\textbf{Países Participantes Países} (LOC) que participaram no passado, mas não em 1970

O \textbf{Festival Eurovisão da Canção 1970} (MISC) (em inglês : \textbf{Eurovision Song Contest} (ORG) 1970 , em \textbf{Francês} (MISC) : Concours Eurovision de la chanson 1970 e em holandês : \textbf{Eurovisiesongfestival 1970} (MISC) ) foi o 15 Festival Eurovisão da Canção e teve lugar em 21 de março de 1970 em \textbf{Amesterdão} (LOC) com a canção " \textbf{All Kinds of Everything} (MISC) ".

Devido ao empate de 4 países que ocorreu em 1969 houve grande controvérsia acerca de qual dos países vencedores devia realizar o festival em 1970. Acabou por se realizar nos \textbf{Países Baixos} (LOC) . A \textbf{Finlândia} (LOC) , \textbf{Noruega} (LOC) , \textbf{Portugal} (LOC) e \textbf{Suécia} (LOC) boicotaram este festival, porque não ficaram satisfeitos nem com os seus resultados nem com o sistema de votação do ano anterior, com \textbf{Portugal} (LOC) a anunciar a sua desistência em setembro de 1969 . Apesar disso, em maio , dois meses após o certeme europeu, \textbf{Portugal} (LOC) voltaria a organizar o \textbf{Grande Prémio TV da Canção 1970} (MISC) , que seria ganho por \textbf{Sérgio Borges} (PER) , com a canção " Onde vais rio que eu canto ".

Para evitar um incidente semelhante ao de 1969 foi criada uma nova regra. Se duas ou mais canções obtivessem o mesmo número de pontos, cada canção deveria ser interpretada novamente. Depois disso, todos os júris, excepto os países visados reuniam-se para escolher o vencedor. No final, se o empate persistisse seriam ambas declaradas como vencedoras.

\textbf{Julio Iglesias} (PER) , representante da \textbf{Espanha} (LOC) cantou vestido com um fato azul que não tinha bolsos para evitar que o cantor interpretasse o tema com as mãos nos bolsos, tal como era seu hábito.

Alguns dias antes do festival toda a imprensa já tinha dado a vitória à representante da \textbf{Irlanda} (LOC) , \textbf{Dana} (PER) , o que acabou por se confirmar na grande noite do festival.



\textbf{Local} (MISC)O \textbf{Festival Eurovisão da Canção} (MISC) 1970 ocorreu em \textbf{Amesterdão} (LOC) , nos \textbf{Países Baixos} (LOC) . \textbf{Amesterdão} (LOC) é a capital e a cidade mais populosa do \textbf{Reino dos Países Baixos} (LOC) . O seu estatuto de capital holandesa é garantido pela \textbf{Constituição dos Países Baixos} (ORG) , embora não seja a sede do governo holandês, que fica em \textbf{Haia} (LOC) . Originária de uma pequena vila de pescadores que surgiu no final do século XII, \textbf{Amesterdão} (LOC) se tornou um dos portos mais importantes do mundo durante a \textbf{Século de Ouro} (MISC) dos \textbf{Países Baixos} (LOC) (século XVII), como resultado de seus desenvolvimentos inovadores no comércio . Durante essa época, a cidade era o principal centro financeiro e de diamantes do mundo. Nos séculos XIX e XX a cidade se expandiu e muitos novos bairros e subúrbios foram planejados e construídos. Os canais de \textbf{Amesterdão} (LOC) e a \textbf{Linha de Defesa de Amesterdão} (LOC) são considerados \textbf{Patrimônios Mundiais} (MISC) pela \textbf{UNESCO} (ORG) . Como a capital comercial dos \textbf{Países Baixos} (LOC) e um dos principais centros financeiros da \textbf{Europa} (LOC) , \textbf{Amesterdão} (LOC) é considerada uma cidade global alfa . A cidade é também a capital cultural do país. Muitas grandes instituições holandesas mantêm suas sedes na cidade e sete das 500 maiores empresas do mundo, incluindo \textbf{Philips} (ORG) e \textbf{ING} (ORG) , baseiam-se na capital holandesa. Em 2012, \textbf{Amsterdã} (LOC) foi classificada como a segunda melhor cidade para se viver pela \textbf{Economist Intelligence Unit} (ORG) (\textbf{EIU} (ORG)).

O festival em si realizou-se no \textbf{RAI Congrescentrum} (LOC) , um complexo de salas de conferências e exposições, inaugurado em 1961 .



\textbf{Formato Depois} (PER) de terem sido eleitas quatro vencedoras no ano anterior, a \textbf{UER} (ORG) sorteou em qual delas se realizaria o \textbf{Festival} (LOC). A sorte coube aos \textbf{Países Baixos} (LOC) .

Os produtores holandeses foram obrigados a organizar o evento, apesar de apenas 12 nações fazerem a viagem a \textbf{Amesterdão} (LOC). O resultado foi um formato que perdurou quase até os dias atuais. Uma sequência de abertura estendida (filmada em \textbf{Amesterdão} (LOC)) abriu a edição, enquanto cada país foi introduzido por um pequeno "cartão postal" de vídeo apresentando cada um dos artistas participantes, ostensivamente na sua própria nação. No entanto, os "cartões postais" para a \textbf{Suíça} (LOC) , \textbf{Luxemburgo} (LOC) e \textbf{Mónaco} (LOC) foram todos filmados em \textbf{Paris} (LOC) (assim como o cartão postal francês). O longo filme de introdução (com mais de quatro minutos de duração) foi seguido pelo que provavelmente é uma das introduções mais curtas feitas por uma apresentadora. \textbf{Willy Dobbe} (PER) apenas apresentou os telespectadores em inglês , francês e holandês , terminando sua introdução depois de apenas 24 segundos. Legendas em tela introduziam cada participação, com os títulos das músicas listados em letras minúsculas e os nomes dos artistas e compositores/autores, todos em maiúsculas.

Para evitar um incidente como em 1969, foi criada uma regra de empate, que consistia em que se duas ou mais músicas obtivessem o mesmo número de votos e estivessem empatadas no primeiro lugar, cada música teria que ser interpretada novamente. Depois disso, cada júri nacional (que não fosse o júri dos países envolvidos) teria que votar na que julgavam ser a melhor. Se os países empatassem novamente, eles dividiriam o primeiro lugar.



\textbf{Visual} (MISC)O design do cenário foi desenvolvido por \textbf{Roland de Groot} (PER) ; um projeto simples, de inspiração modernista, era composto de cinco barras curvas horizontais e sete esferas de prata suspensas que podiam ser movidas de várias maneiras diferentes, bem como uma entrada de metal por meio do qual os artistas interpretaram as canções. Além disso, os pontos autorizados a dar cores específicas em segundo plano e obter o jogo de luz nas barras e esferas. Esse foi o primeiro uso verdadeiramente moderno das possibilidades oferecidas pela televisão a cores. À direita do palco estavam o painel de votação e a mesa do supervisor, e ao pé, no fosso, a orquestra.

O vídeo introdutório mostrava vistas aéreas de \textbf{Amesterdão} (LOC) , intercaladas com vistas turísticas das ruas e monumentos da capital holandesa. As vistas de \textbf{Congrescentrum} (LOC) se seguiram dia e noite. A câmara varreu o hall de entrada, duas recepcionistas abriram uma porta para ela e entrou na sala.

A apresentadora foi \textbf{Willy Dobbe} (PER) . Após o vídeo introdutório, ela cumprimentou brevemente os espectadores, em inglês , francês e holandês , e depois retornou apenas para apresentar a votação.

A orquestra foi conduzida por \textbf{Dolf} (PER) \textbf{van der Linden} (PER) .

No intervalo, a \textbf{NOS} (ORG) pediu aos bailarinos de \textbf{Don Lurio} (PER) para realizarem uma coreografia. Consistia de um número de dança moderna e rítmica, apresentado a música "\textbf{Rich Man's Frug} (MISC)", do musical \textbf{Sweet Charity} (ORG) .

Cada cartão postal mostrava artistas a passear pela capital de seu país. Os representantes monegasca, suíço, e \textbf{luxemburguês} (PER) filmaram em \textbf{Paris} (LOC) .



\textbf{Votação} (MISC)

Mais uma vez os países votaram pela mesma ordem à da sua actuação, algo que foi feito aparentemente para acelarar o espetáculo, pois os países a actuar não podiam (nem podem) atribuir pontos a si mesmo, sendo assim poderiam acabar a sua votação a tempo de serem os primeiros a serem chamados. O sistema de votação foi o mesmo utilizado em edições anteriores, em que cada país tinha um júri composto por dez membros, e cada um desses dez jurados atribuiam um ponto à sua canção preferida.

Para evitar um incidente como em 1969, foi criada uma regra de empate, que consistia em que se duas ou mais músicas obtivessem o mesmo número de votos e estivessem empatadas no primeiro lugar, cada música teria que ser interpretada novamente. Depois disso, cada júri nacional (que não fosse o júri dos países envolvidos) teria que votar na que julgavam ser a melhor. Se os países empatassem novamente, eles dividiriam o primeiro lugar.

O supervisor delegado pela \textbf{UER} (ORG) foi, mais uma vez, \textbf{Clifford Brown} (PER) , acompanhado por dois secretários e pela apresentadora.

A \textbf{Irlanda} (LOC) tomou a votação do começo ao fim. O porta-voz do júri irlandês, que falou por último, cumprimentou \textbf{Willy Dobbe} (PER) com as palavras: "\textbf{Olá} (MISC), \textbf{Amsterdão} (LOC)! Esta é uma chamada alegre de \textbf{Dublin} (LOC)!"

Durante a votação, a câmera fez duas filmagens dos artistas nos bastidores. Ela mostrava \textbf{Dana} (LOC) ( \textbf{Irlanda} (LOC) ), acompanhada por por \textbf{Henri Dès} (PER) ( \textbf{Suíça} (LOC) ) e \textbf{Mary Hopkin} (PER) ( \textbf{Reino Unido} (LOC) ).

\textbf{David Alexandre Winter} (PER) ( \textbf{Luxemburgo} (PER) ) exigiu completo silêncio durante a votação, o que desagradou os demais concorrentes. No final do anúncio dos resultados, ele se viu ridículo, terminando em último com 0 pontos.

Esta foi a única vez neste sistema de votação que um júri concedeu nove votos a uma única música. Este é o júri belga, nove membros votaram a favor da música irlandesa.

Pela primeira vez, o \textbf{Luxemburgo} (LOC) terminou em último lugar com 0 pontos.



\textbf{Participantes Dos artistas participantes} (MISC), destacaram-se a galesa \textbf{Mary Hopkin} (PER) , que era uma das favoritas, que já tinha uma carreira de considerável sucesso, com o single " \textbf{Those Were the Days} (MISC) ", de 1968 produzido por \textbf{Paul McCartney} (PER) ; \textbf{David Alexandre Winter} (PER) , nascido em \textbf{Amsterdão} (LOC) e representando o \textbf{Luxemburgo} (LOC) , deixara os \textbf{Países Baixos} (LOC) para ir para \textbf{Paris} (LOC) em 1969, onde gravara o seu maior sucesso "\textbf{Oh Lady Mary} (MISC)" que se tornou rapidamente disco de ouro; \textbf{Dominique Dussault} (PER) , representando o \textbf{Mónaco} (LOC) , cantou uma homenagem a \textbf{Marlene Dietrich} (PER) , que desejava ser como ela, tendo no terceiro verso, citado várias frases da estrela em inglês e alemão .

\textbf{Pelos} (LOC) \textbf{Países Baixos} (LOC) , um trio, \textbf{Hearts of Soul} (MISC) , foi apresentado. No entanto, as regras só permitiam a participação de solistas e duos, acompanhados por coros. Assim, a vocalista do trio, \textbf{Patricia Maessen} (PER) , cantou sozinha, acompanhada pelas restantes elementos do grupo, que foram coristas. Este foi o último ano em que este regulamento esteve em vigor.

\textbf{Julio Iglesias} (PER) , então desconhecido, representou \textbf{Espanha} (LOC) .



\textbf{Festival} (PER)A ordem de votação no \textbf{Festival Eurovisão da Canção 1970} (MISC), foi a seguinte:

Em baixo encontra-se a lista de maestros que conduziram a orquestra, na respectiva actuação de cada país concorrente.



Artistas repetentes

Em 1970, pela primeira vez na história, não houve repetentes.



Festival Eurovisão da Canção 1971



\textbf{Países Participantes Países} (LOC) que participaram no passado, mas não em 1971

O \textbf{Festival Eurovisão da Canção 1971} (MISC) (em inglês : \textbf{Eurovision Song Contest} (ORG) 1971 ; em francês : \textbf{Concours Eurovision} (MISC) de la chanson 1971 ; em irlandês : \textbf{Comórtas Amhránaíochta} (LOC) na \textbf{hEoraifíse} (MISC) 1971 ) foi o 16 edição anual do evento e teve lugar a 3 de abril de 1971 em \textbf{Dublin} (LOC) . A apresentadora do evento foi Bernadette Ní Ghallchóir . \textbf{Séverine} (PER) foi a vencedora do festival com a canção \textbf{Un banc} (MISC), un arbre, une rue . Os países classificados em 2 ( \textbf{Espanha} (LOC) ) e 3 lugares ( \textbf{Alemanha} (LOC) ) também receberam prémios no palco. A pedido de vários países, a \textbf{EBU} (MISC) alterou o sistema de votação, assegurando assim que nenhuma canção obteria 0 pontos, consistindo em 2 júris por país que avaliaram individualmente todas as canções concorrentes e atribuíram-lhes 5, 4, 3, 2 ou 1 pontos consoante o seu agrado. No entanto o sistema causou muita polémica pois o total de pontos atribuído por cada país não era o mesmo. A título de exemplo, o \textbf{Luxemburgo} (LOC) concedeu um total de 43 pontos, enquanto que a \textbf{França} (LOC) concedeu um total de 107 pontos. Apesar da polémica, o sistema vira a ser novamente utilizado em 1972 e 1973 . Pela primeira vez, grupos até 6 elementos foram permitidos participar, com a \textbf{Suécia} (LOC) (\textbf{Family Four} (MISC)) e a \textbf{Suíça} (LOC) ( \textbf{Peter} (LOC), \textbf{Sue \& Marc} (LOC) ) a estrearem-se nesse registo.

A vitória do \textbf{Mónaco} (LOC) foi sua primeira e única vitória. A música foi interpretada por uma cantora francesa, morando na \textbf{França} (LOC), cantado em francês, conduzido por um francês e escrito por uma equipa francês. \textbf{Séverine} (PER) afirmou mais tarde que nunca visitou o \textbf{Mónaco} (LOC) antes ou depois da vitória - uma afirmação facilmente refutada pelo vídeo de pré-visualização enviado pela \textbf{Télé Monte Carlo} (LOC) com a cantora num local no \textbf{Principado} (LOC).

\textbf{Malta} (LOC) fez sua estreia neste concurso, enquanto a \textbf{Áustria} (LOC) , a \textbf{Finlândia} (LOC) , a \textbf{Noruega} (LOC) , \textbf{Portugal} (LOC) e a \textbf{Suécia} (LOC) voltaram depois de uma breve ausência. Isso trouxe o número total de países para 18.



Localização

O local escolhido foi o \textbf{Gaiety Theatre} (LOC) , na capital do país, \textbf{Dublin} (LOC) . Esta foi a primeira vez que o concurso foi realizado na \textbf{Irlanda} (LOC) .

O \textbf{Gaiety Theatre} (PER) , construída em 1871 , é um espaço que acolhe vários eventos musicais, tais como o pantomima de \textbf{Natal} (MISC) todos os anos desde 1874 .

Quanto à cidade, localiza-se na costa oriental da ilha, na província de \textbf{Leinster} (LOC) . Desde 1898 possui nível administrativo de condado ( county-boroughs ). Seus limites são os condados de \textbf{Fingal} (LOC) a norte, \textbf{Dublin} (LOC) meridional a sudoeste e \textbf{Dun Laoghaire-Rathdown} (LOC) a sudeste. Encontra-se na foz do \textbf{rio Liffey} (LOC) , na \textbf{baía de Dublin} (LOC). Tem uma população de 505.739 habitantes na cidade e sua área metropolitana tem 1.661.185 habitantes. Foi fundada pelos \textbf{Viquingues} (LOC) que a dominaram até 1170 .



\textbf{Formato} (PER)Pela primeira vez, cada emissora participante era obrigada a transmitir todas as músicas em "pré-visualizações" antes da final ao vivo, para ser transmitido em todas as emissoras. Era proibido a publicação das letras e músicas das canções até essa data, caso contrário eram excluídas. A pré-visualização da \textbf{Bélgica} (LOC) apresentou a \textbf{Nicole \& Hugo} (ORG) interpretando a música " \textbf{Goeiemorgen} (LOC), \textbf{Morgen} (LOC) ", mas \textbf{Nicole} (LOC) foi atingida com uma doença súbita dias antes da final do concurso. Ambos foram substituídos por \textbf{Jacques Raymond} (PER) e \textbf{Lily Castel} (PER) . Os relatórios sugerem que \textbf{Castel} (LOC) não teve tempo suficiente para comprar um vestido adequado para o show.

A \textbf{BBC} (ORG) estava preocupada com a possível reação da audiência à canção do \textbf{Reino Unido} (LOC) devido às hostilidades que grassam na \textbf{Irlanda do Norte} (LOC) . Eles selecionaram especificamente uma cantora da \textbf{Irlanda do Norte} (LOC) , \textbf{Clodagh Rodgers} (ORG) , que era popular tanto no \textbf{Reino Unido} (LOC) como na \textbf{República da Irlanda} (LOC) , para aliviar qualquer mal-estar do público de \textbf{Dublin} (LOC) . No entanto, \textbf{Rodgers} (PER) ainda recebeu ameaças de morte do \textbf{IRA} (ORG) por representar o \textbf{Reino Unido} (LOC) .

Grupos de até seis pessoas foram autorizados a realizar pela primeira vez, com a regra em concursos anteriores de realizar solos ou duetoa abolido.

Esta era apenas a segunda transmissão externa da \textbf{RTÉ} (MISC) e a primeira a cores. O concurso foi transmitido também na \textbf{Islândia} (LOC) , \textbf{Estados Unidos} (LOC) e \textbf{Hong Kong} (LOC) dias depois.



\textbf{Visual} (MISC)A decoração era de inspiração celta , com a orquestra localizado num buraco ao pé do palco. O palco foi decorado por uma borda metálica oval. Um espaço circular foi organizado no centro para os artistas. No fundo, foi suspensa uma escultura composta por três círculos de padrões concêntricos. Holofotes de cor foram usados, permitindo dar uma atmosfera diferente a cada atuação.

O espetáculo começou com um filme romântico, uma carruagem com um casal a bordo para uma ótima noite, atravessando \textbf{Dublin} (LOC) , as suas ruas e parques para estacionar em frente ao \textbf{Gaiety Theatre} (PER) .

Foi também uma estreia para a \textbf{RTÉ} (LOC) , que com este espetáculo produziu o seu primeiro evento a cores. Para as diferentes fases da apresentação, Bernadette Ní Ghallchóir foi instalada numa bancada na sala e dirigiu-se aos telespetadores em irlandês , francês e inglês . As músicas foram precedidas por um cartão postal turístico de seus países. Os artistas não apareceram nesses filmes.

Para o intervalo, foi produzido um filme mostrando danças, músicas e tradições irlandesas, no \textbf{Castelo de Bunratty} (LOC) , juntamente com a recriação de um banquete medieval.

\textbf{Bernadette} (PER) chamou o júri para indicar seus votos, que haviam sido instalados no palco e, à chamada de cada país, todos mostravam a nota que haviam atribuído, em grupos de três países.



Sistema de votação

Um novo sistema de votação foi introduzido no concurso deste ano: cada país enviou dois membros do júri, um com mais de 25 anos e o outro com menos de 25 anos (com pelo menos dez anos de diferença entre as idades), com ambos a pontuar cada país (exceto o seu próprio) com uma pontuação entre 1 e 5 pontos.

Isso significava que nenhum país poderia ter menos de 34 pontos. No entanto, isso deu origem a um grande problema: alguns membros do júri tendiam a atribuir apenas um ou dois pontos, para aumentar as chances de conquista de seus respectivos países. Apesar da polémica, o sistema vira a ser novamente utilizado em 1972 e 1973 .



\textbf{Festival} (MISC)Os países votaram em blocos de 3 (\textbf{Áustria} (LOC), \textbf{Malta} (LOC) e \textbf{Mónaco} (LOC); \textbf{Suíça} (LOC), \textbf{Alemanha} (LOC) e \textbf{Espanha} (LOC)...)

Os países que receberam a pontuação perfeita de 10 pontos foram os seguintes:

Em baixo encontra-se a lista de maestros que conduziram a orquestra, na respectiva actuação de cada país concorrente.



Artistas repetentes

Alguns artistas repetiram a sua experiência \textbf{Eurovisiva} (PER). Em 1971, os repetentes foram:



Festival Eurovisão da Canção 1972



\textbf{Países Participantes Países} (LOC) que participaram no passado, mas não em 1972

O \textbf{Festival Eurovisão da Canção 1972} (MISC) (em inglês : \textbf{Eurovision Song Contest} (ORG) 1972 e em francês : \textbf{Concours Eurovision} (MISC) de la chanson 1972 ) foi o 17 \textbf{Festival Eurovisão da Canção} (MISC) realizou-se em 25 de março de 1972 em \textbf{Edimburgo} (LOC) , \textbf{Escócia} (LOC) , no \textbf{Reino Unido} (LOC) . Esta foi a única vez que o certame se realizou na \textbf{Escócia} (LOC) . \textbf{Moira Shearer} (PER) foi a apresentadora do concurso \textbf{O Festival} (MISC) foi ganho por \textbf{Vicky Leandros} (PER) que representou o \textbf{Luxemburgo} (LOC) com a canção " \textbf{Après Toi} (MISC) " (Depois de \textbf{Ti} (PER)).

O \textbf{Mónaco} (LOC) vencedor do Festival Eurovisão da Canção 1971 foi incapaz de realizar o evento, cabendo ao \textbf{Reino Unido} (LOC) a tarefa de produzir mais um festival. O \textbf{Reino Unido} (LOC) pela primeira vez, escolheu um lugar fora da \textbf{Inglaterra} (LOC) .

A \textbf{Irlanda} (LOC) pela primeira e única vez enviou uma canção interpretada em língua irlandesa . Das outras vezes foi sempre em língua inglesa .

\textbf{Yves Desca} (PER) foi o escritor das canções vencedoras de 1971 e 1972, tornando-se a primeira pessoa a vencer por dois países diferentes e também a primeira a vencer duas vezes seguidas.



\textbf{Local} (MISC)O \textbf{Festival Eurovisão da Canção} (MISC) 1971 ocorreu em \textbf{Edimburgo} (LOC) , na \textbf{Escócia} (LOC) , no \textbf{Reino Unido} (LOC) . \textbf{Edimburgo} (LOC) , é uma cidade , capital da \textbf{Escócia} (LOC) , no \textbf{Reino Unido} (LOC) , com poderes próprios, situada na margem sul do estuário do \textbf{rio Forth} (LOC) ( \textbf{Firth of Forth} (LOC) ). É a capital escocesa desde 1492 , sendo sede do parlamento escocês desde 1999 . \textbf{Edimburgo} (LOC) é também uma das 32 \textbf{Áreas de Conselho da Escócia} (LOC) , subdivisão administrativa similar aos estados brasileiros, desde 1996 quando este tipo de subdivisão foi atribuída para substituir os distritos e regiões que existiam desde 1975 . A cidade é dominada pelo \textbf{Castelo de Edimburgo} (LOC) construído sobre uma rocha de origem vulcânica. Após a unificação do parlamento da \textbf{Escócia} (LOC) com o da \textbf{Inglaterra} (LOC), \textbf{Edimburgo} (LOC) perdeu sua importância política mas permaneceu um importante centro económico e cultural. A cidade é mundialmente conhecida pelo \textbf{Festival de Edimburgo} (PER) que acontece durante três semanas no mês de agosto. A cidade ainda possui uma das mais prestigiadas universidades da \textbf{Europa} (LOC) e do mundo, a \textbf{Universidade de Edimburgo} (LOC) , pioneira na informática e gerenciamentos.

O festival em si realizou-se no \textbf{Usher Hall} (LOC) , é uma sala de concertos, inaugurado em 1914 . É o centro da vida musical da cidade. É la sede oficial do \textbf{Festival de Edimburgo} (PER).



\textbf{Formato} (PER)Originalmente, os vencedores do ano anterior, o \textbf{Mónaco} (LOC) , queriam sediar o evento de 1972 e propuseram um evento ao ar livre a ser realizado em junho de 1972 , em vez de ser na primavera, como era costume. Havia também um salão adequado a ser construído, que a \textbf{Radio Monte Carlo} (LOC) tentou acelerar. No entanto, em julho de 1971 , a \textbf{Radio Monte Carlo} (LOC) finalmente declarou que não havia condições para organizar o evento.

Sabia-se que a \textbf{TVE} (ORG) ( \textbf{Espanha} (LOC) ) e a \textbf{ARD} (ORG) ( \textbf{Alemanha} (LOC) ), como segundo e terceiro colocados de 1971, recusariam a oportunidade de realizar o concurso de 1972, e que a \textbf{EBU} (MISC) estava "chamando alto" os voluntários. e mais uma vez a \textbf{BBC} (ORG) decidiu aceitar o desafio de organizar o concurso. Tendo reduzido a escolha a dois locais possíveis, eles optaram por encenar o concurso pela primeira vez em \textbf{Edimburgo} (LOC) , no \textbf{Usher Hall} (LOC) . O \textbf{Mónaco} (LOC) tornou-se assim o único vencedor que nunca recebeu a competição no seu solo.

O sorteio da ordem de atuação foi feito em \textbf{Londres} (LOC) na quarta-feira, 1 de dezembro de 1971 .

Os ensaios começaram com os artistas na quarta-feira, 22 de março , com cada país tendo um ensaio inicial de 50 minutos com a orquestra de 44 elementos. O concurso foi transmitido em 28 países.

Pela primeira vez, o programa chegou a cores ao emissor da \textbf{RTP} (ORG) , em \textbf{Monsanto} (LOC) . Era evidente que só podia ver a cores quem tivesse um televisor para o efeito: só havia 30 \textbf{em Portugal} (LOC).



\textbf{Visual} (MISC)O design do palco incluiu uma tela para apresentar e acompanhar as atuações concorrentes no palco, e para mostrar o intervalo e a sequência de votação que foram feitas no \textbf{Castelo de Edimburgo} (LOC) . Antes da atuação de cada país, uma imagem de cada artista, juntamente com seus nomes e o título da música, foram projetadas na tela, e durante cada apresentação, formas espirais animadas foram projetadas como efeito visual adicional. O conjunto foi desenhado por \textbf{Brian Trediggen} (PER) . Os jurados estavam localizados na segurança do castelo e assistiram às apresentações concorrentes em \textbf{Usher Hall na TV} (MISC).

A orquestra, liderada por \textbf{Malcolm Lockyer} (PER) e com 44 músicos, localizou-se num buraco ao pé do palco. Este último era circular e emoldurado por duas paredes em ângulos retos. As paredes eram compostas de painéis brancos alternados, aparelhos luminosos e colunas de luz. Três entradas, esquerda, centro e direita, foram tapadas por cortinas, agitadas por fãs. Durante as apresentações, a cena foi animada por jogos de sombras e luzes. Na parede à direita, figuras azuladas foram projetadas. Quanto à parede esquerda, incluía uma tela grande, que permitia apresentar os artistas e sua música.

Esta foi a primeira vez na história da competição que um ecrã foi utilizado, localizando-se à esquerda do palco.

O vídeo introdutório mostrou vários pontos turísticos de \textbf{Edimburgo} (LOC) . A câmera então mostrou o teto da sala, depois o público. A orquestra tocou então a partitura de " \textbf{Un banc} (MISC), un arbre, une rue ", a música vencedora do ano anterior.

Pela primeira vez desde 1970, não havia cartões postais. A câmara fez um plano na tela da parede, mostrando os rostos dos artistas e o título da música deles.

O intervalo mostrou um vídeo intitulado \textbf{Tatto} (PER) no \textbf{Castelo de Edimburgo} (LOC) . Foi um desfile militar, ao som da gaita de \textbf{foles} (MISC) , executado por oito regimentos escoceses, 5 que tiveram lugar na esplanada à entrada do \textbf{Castelo de Edimburgo} (LOC) . O evento não foi gravado especificamente para o concurso: os produtores simplesmente reutilizaram um arquivo de vídeo de 1968.



\textbf{Votação} (MISC)Cada país enviou dois membros do júri, um com mais de 25 anos e o outro com menos de 25 anos (com pelo menos dez anos de diferença entre as idades), com ambos a pontuar cada país (exceto o seu próprio) com uma pontuação entre 1 e 5 pontos. O júri ficou instalado no \textbf{Castelo de Edimburgo} (LOC) , alegadamente, para evitar fraudes no voto.

O supervisor no local delegado pela \textbf{EBU} (MISC) foi \textbf{Clifford Brown} (PER) , que repetia cada vez em inglês o número de votos e o total de votos para cada país.

Durante a votação, a câmera fez vários close-ups dos artistas. Em particular, \textbf{Vicky Leandros} (PER) , \textbf{Sandra Reemer} (PER) , \textbf{The New Seekers} (ORG) , \textbf{Carlos Mendes} (PER) , \textbf{Véronique Müller} (PER) e \textbf{Beatrix Neundlinger} (PER) apareceram.

1972 foi também o primeiro ano em que não ocorreu empates na votação.



\textbf{Festival} (MISC)Os países votaram em blocos de 3 (\textbf{Alemanha} (LOC), \textbf{França} (LOC) e \textbf{Irlanda} (LOC); \textbf{Espanha} (LOC), \textbf{Reino Unido} (LOC) e \textbf{Noruega} (LOC)...). Assim, ordem de votação no Festival Eurovisão da Canção 1972, foi a seguinte:

Os países que receberam a pontuação perfeita de 10 pontos foram os seguintes:

Em baixo encontra-se a lista de maestros que conduziram a orquestra, na respectiva actuação de cada país concorrente.



Festival Eurovisão da Canção 1973



\textbf{Países Participantes Países} (LOC) que participaram no passado, mas não em 1973

O \textbf{Festival Eurovisão da Canção 1973} (MISC) (em inglês : \textbf{Eurovision Song Contest} (ORG) 1973 e em francês : \textbf{Concours Eurovision} (MISC) de la chanson 1973 ) foi o 18 Festival Eurovisão da Canção , realizado em 7 de abril de 1973 no \textbf{Luxemburgo} (LOC) . \textbf{Helga Guitton} (PER) foi a apresentadora do concurso dessa noite. \textbf{Anne-Marie David} (PER) venceu o \textbf{Festival} (PER) desse ano representando o \textbf{Luxemburgo} (LOC) com a canção " \textbf{Tu te reconnaîtras} (MISC) ", com a maior pontuação já alcançada na \textbf{Eurovisão} (LOC) sob qualquer formato de votação, registando 129 pontos de 160 possíveis; pontuando pouco abaixo de 81\% do máximo possível, mas em parte devido a um sistema de pontuação que garantiu a todos os países pelo menos dois pontos a cada outro país.

Depois do sucedido no ano anterior nos \textbf{Jogos Olímpicos de Munique} (MISC) , o festival de 1973 foi marcado pela tónica da segurança, devido a ameaças terroristas por parte de grupos palestinos. O motivo para essa segurança relacionou-se com \textbf{Israel} (LOC) , o primeiro país não europeu a participar, participar pela primeira vez no \textbf{Festival Eurovisão da Canção} (MISC) .

Pela primeira vez na história do concurso, membros de uma família real governante estavam presentes no salão. A câmera mostrou vários membros da família real, como a princesa \textbf{Maria Astrid} (PER) , o então princesa herdeiro \textbf{Henrique} (PER) e a princesa \textbf{Margarida} (PER) .



\textbf{Local} (MISC)O \textbf{Festival Eurovisão da Canção} (MISC) 1973 ocorreu na capital do \textbf{Luxemburgo} (LOC) . A cidade do \textbf{Luxemburgo} (LOC) é uma comuna de \textbf{Luxemburgo} (LOC) com status de cidade , pertencente ao distrito de \textbf{Luxemburgo} (LOC) e ao cantão de \textbf{Luxemburgo} (LOC) . \textbf{Luxemburgo} (LOC) é uma cidade muito desenvolvida em seu comércio e nas indústrias. \textbf{Luxemburgo} (LOC) é uma das cidades mais ricas da \textbf{Europa} (LOC) e tornou-se um importante centro financeiro e administrativo. A cidade do \textbf{Luxemburgo} (LOC) contém várias instituições da \textbf{União Europeia} (ORG) , incluindo o \textbf{Tribunal de Justiça Europeu} (ORG) , o \textbf{Tribunal de Contas} (ORG) e o \textbf{Banco Europeu de Investimento} (ORG) .

O festival em si realizou-se no \textbf{Grand Théâtre de Luxembourg} (LOC) , inaugurado em 1964 como \textbf{Théâtre Municipal de la Ville} (LOC) de \textbf{Luxembourg} (LOC) , é uma das maiores salas de espetáculos do país para teatro , ópera e ballet . Passou por obras de renovação em 2002-2003, resultando em melhorias substanciais nas instalações de tecnologia de palco, acústica e iluminação.



\textbf{Formato} (PER)A regra que obrigava os países a participar com canções interpretadas na sua língua oficial foi abandonada, daí que alguns países como a \textbf{Suécia} (LOC) tenham escolhido cantar em inglês e não na sua língua materna .

Pela primeira vez, foi permitido que fosse utilizado playback instrumental, desde que os instrumentos também estivessem a ser tocados em palco. Isso foi levantado por \textbf{Cliff Richard} (PER) , representando o \textbf{Reino Unido} (LOC) , que queria que playback instrumental fosse utilizado para que os seus músicos não precisassem de atuar, o que causou alguma controvérsia.

\textbf{Monique Dominique} (PER) tornou-se a primeira mulher a dirigir a orquestra, pela \textbf{Suécia} (LOC) . \textbf{Nurit Hirsch} (PER) , de \textbf{Israel} (LOC) , tornar-se-ia alguns mnuitos depois a segunda mulher a dirigir a orquestra.



Visual

O vídeo introdutório mostrava os pontos turísticos do \textbf{Luxemburgo} (LOC) , seguidos de um resumo dos preparativos para a competição. A orquestra tocou a partitura de " \textbf{Après} (LOC) \textbf{Toi} (MISC) ", a música vencedora do ano anterior. A câmera terminou com uma foto de \textbf{Vicky Leandros} (PER) , vencedora da edição de 1972, sentada na sala ao lado de seu pai, \textbf{Leo Leandros} (PER) .

A apresentadora foi \textbf{Helga Guitton} (PER) , que falou aos espectadores em francês , inglês e alemão , acrescentando duas saudações em luxemburguês .

O palco tinha uma cor cinza uniforme. A orquestra estava na arquibancada, colocada na parte de trás do palco, atrás dos artistas. Estes últimos faziam a sua entrada por um pórtico, decorado com um medalhão representado pelo leão luxemburguês. O conjunto consistia em dois pódios circulares, sendo o da esquerda reservado para os coristas, o da direita para os maestros. O lado esquerdo do palco foi decorado com seis colunas de vidro contendo buquês de flores. Durante as apresentações, o pórtico foi iluminado por peças de luz.

Como no ano anterior, não havia cartões postais. A câmera mostrou apenas uma foto dos artistas tiradas durante os ensaios.

O intervalo foi ocupado pelo palhaço \textbf{Charlie Rivel} (PER) . Foi anunciado por \textbf{Helga Guitton} (PER) como "a grande diva \textbf{Carlotta Rivelo} (PER)". Disfarçado de cantora de ópera, vestido com um vestido de paetês vermelho e um gigantesco chapéu, \textbf{Charlie Rivel} (PER) parodiou um recital de diva caprichosa. Seu pianista começou a dar à diva uma rosa que ela não podia sentir o cheiro. Ela então recebeu a tampa do piano em sua mão direita e começou a chorar. Depois veio uma discussão com o pianista e um concerto de trompa, emitido pelo opulento peito da diva. Parecia então que ela, não tendo colocado óculos, não conseguia decifrar sua pontuação. O recital finalmente começou, brevemente interrompido por uma queda do pianista, assustado por um grito da diva, depois pela perda da partitura, jogado no chão num momento de emoção e depois lido para trás. A diva acaba se perdendo em suas anotações e um de seus seios, estourando, concluindo o intervalo.



\textbf{Controvérsias} (PER) e incidentes

O evento foi assinalado por um escândalo , com a canção da \textbf{Espanha} (LOC) " Eres tú " ser acusada de plágio , sugeriu-se que essa canção não era mais que uma cópia da canção que representara a \textbf{Jugoslávia} (ORG) em 1966 (" \textbf{Brez} (LOC) besed " interpretado por \textbf{Berta Ambrož} (PER) . Apesar das acusações, o certo é que a canção não foi desqualificada. Tem-se referido que a sua não desqualificação deveu-se a razões políticas, por a \textbf{Espanha de Franco} (LOC) ter mais poder do que a \textbf{Jugoslávia} (ORG) de \textbf{Josip Broz Tito} (PER) . A canção terminou em segundo lugar e tornou-se um enorme sucesso internacional. Curiosamente a censura evitou que na \textbf{Espanha} (LOC) se tivesse conhecimento destes factos.

A letra da canção de \textbf{Portugal} (LOC) " \textbf{Tourada} (LOC) " era uma crítica implícita à decadente ditadura que então vigorava no país. Ainda hoje não se sabe como é que os censores, tão zelosos quase sempre, não compreenderam a mensagem da canção.

Depois do sucedido no ano anterior nos \textbf{Jogos Olímpicos de Munique} (MISC) , o festival de 1973 foi marcado pela tónica da segurança, devido a ameaças terroristas por parte de grupos palestinos. O motivo para essa segurança relacionou-se com \textbf{Israel} (LOC) , o primeiro país não europeu a participar, participar pela primeira vez no \textbf{Festival Eurovisão da Canção} (MISC) .

A \textbf{Irlanda} (LOC) também esteve em destaque nesta edição. A cantora, \textbf{Maxi} (PER) , teve um grande desentendimento com o maestro, pois a orquestra estava a tocar muito rapidamente. Por essa razão, chegou a temer-se que a cantora desistisse. Por isso, a \textbf{RTÉ} (MISC) chamou ao \textbf{Luxemburgo} (LOC) a cantora \textbf{Tina Reynolds} (PER) , para a substituir. Contudo, chegou-se a acordo e \textbf{Maxi} (PER) representou a \textbf{Irlanda} (LOC) , não sendo preciso a presença de \textbf{Tina Reynolds} (PER) .



\textbf{Votação} (MISC)Cada país enviou dois membros do júri, um com mais de 25 anos e o outro com menos de 25 anos (com pelo menos dez anos de diferença entre as idades), com ambos a pontuar cada país (exceto o seu próprio) com uma pontuação entre 1 e 5 pontos. O júri ficou instalado nos estúdios da \textbf{Villa Louvigny} (LOC) .

O supervisor no local delegado pela \textbf{EBU} (MISC) foi \textbf{Clifford Brown} (PER) , que repetia cada vez em francês o número de votos e o total de votos para cada país.

Durante a votação, a câmera fez vários close-ups dos artistas. Em particular, \textbf{Anne-Marie David} (PER) e \textbf{Mocedades} (PER) apareceram.

A votação ficou marcada pelo jurado sénior da \textbf{Suíça} (LOC) , com os seus gestos a fazerem rir os seus colegas.



\textbf{Participações Malta} (MISC) esteve para participar, tendo inclusive sido colocada como o 6 país a atuar. Contudo, desistiu devido aos fracos resultados dos anos anteriores.



\textbf{Festival} (PER)Os países votaram em blocos de 3 (\textbf{Finlândia} (LOC), \textbf{Bélgica} (LOC) e \textbf{Portugal} (LOC); \textbf{Alemanha} (LOC), \textbf{Noruega} (LOC) e \textbf{Mónaco} (LOC)...). Assim, ordem de votação no Festival Eurovisão da Canção 1973, foi a seguinte:

Os países que receberam a pontuação perfeita de 10 pontos foram os seguintes:

Em baixo encontra-se a lista de maestros que conduziram a orquestra, na respectiva actuação de cada país concorrente. \textbf{Monica Dominique} (PER) e \textbf{Nurit Hirsh} (PER) foram as primeiras mulheres a dirigirem a orquestra na história da \textbf{Eurovisão} (LOC).



Festival Eurovisão da Canção 1974



\textbf{Países Participantes Países} (LOC) que participaram no passado, mas não em 1974

O \textbf{Festival Eurovisão da Canção 1974} (MISC) (em inglês : \textbf{Eurovision Song Contest} (ORG) 1974 e em francês : \textbf{Concours Eurovision} (MISC) de la chanson 1974 ) foi o 19. \textbf{Festival Eurovisão da Canção} (MISC) e realizou-se em 6 de abril de 1974 , em \textbf{Brighton} (LOC) , \textbf{Reino Unido} (LOC) . \textbf{Katie Boyle} (PER) foi pela quarta e última vez a apresentadora do festival que foi ganho pela banda sueca \textbf{ABBA} (ORG) , com a canção " \textbf{Waterloo} (MISC) ", que, em 2005 , seria considerada no especial \textbf{Congratulations} (MISC): 50 \textbf{Anos do Festival} (LOC) Eurovisão da Canção , a melhor canção de todos os tempos. Esta canção também foi o ponto de partida para os \textbf{ABBA} (ORG) se lançarem numa carreira internacional de sucesso. A roupa do diretor de orquestra sueco, a lembrar \textbf{Napoleão Bonaparte} (PER) , também ficou para a história.

O \textbf{Luxemburgo} (LOC) , vencedor do ano anterior não queria realizar o festival pelo segundo ano consecutivo, então o \textbf{Reino Unido} (LOC) concordou em realizar o evento. Quem representou o \textbf{Reino Unido} (LOC) nesse festival foi a cantora pop britânica \textbf{Olivia Newton-John} (PER) , que ficou em 4 lugar com a canção " \textbf{Long Live Love} (MISC) ". Ela pretendia cantar uma canção diferente, mas a canção " \textbf{Long Live Love} (MISC) " foi escolhida pelo público.

A \textbf{França} (LOC) tinha escolhido a sua canção para a \textbf{Eurovisão "} (LOC) La vie a vingt cinq ans " interpretada por \textbf{Dani} (PER) , que seria a 14 canção a subir ao palco, mas desistiu porque o presidente francês \textbf{Georges Pompidou} (PER) morreu na semana do festival e o funeral seria a dia 6 de abril , dia de festival. A competição foi gravada e transmitida na semana seguinte, seguida da transmissão da música que representaria a \textbf{França} (LOC). \textbf{Anne-Marie David} (PER) , que entregaria o troféu, também estava ausente.

Em 25 de abril de 1974 , a canção " E depois do adeus ", representante de \textbf{Portugal} (LOC) e interpretada por \textbf{Paulo de Carvalho} (PER) , foi usada como senha para o começo das operações dos tanques e o início da \textbf{Revolução dos Cravos} (ORG) , movimento dos capitães do \textbf{Exército português} (ORG) que derrubou a ditadura salazarista que vigorava no país desde 1926. Esta é considerada a única canção da história da \textbf{Eurovisão} (LOC) a começar literalmente uma revolução.

A \textbf{Malta} (ORG) tinha escolhido \textbf{Enzo Guzman} (PER) e a canção " \textbf{Paċi} (PER) Fid Dinja " (\textbf{Paz no Mundo} (MISC)) para representar o país, mas desistiram por razões desconhecidas.

Já a canção \textbf{Sì} (MISC) , representante da \textbf{Itália} (LOC) e interpretada por \textbf{Gigliola Cinquetti} (PER) , ficou em segundo lugar, com 18 pontos, mas foi motivo de polêmica em seu país: ela foi gravada na mesma época em que o governo italiano fez um referendo para tentar revogar parcialmente a \textbf{Lei do Divórcio} (MISC). Apesar de ser uma canção romântica, seu refrão foi associado ao "\textbf{Sim} (PER)", lado apoiado pela \textbf{Igreja Católica} (ORG) e por partidos de direita. O festival coincidiu com o dia da votação do referendo, o que levou a \textbf{RAI} (ORG) , emissora pública italiana, a suspender a transmissão direta do \textbf{Eurovisão} (LOC) para o país, a fim de evitar influência no resultado. O programa foi gravado e exibido no dia 12 de maio . Já o referendo foi vencido pelo "Não", apoiado pelos partidos de esquerda, e a \textbf{Lei do Divórcio} (MISC) foi mantida, o que representou uma forte derrota sofrida pelo governo e pelo \textbf{Vaticano} (LOC) .



\textbf{Local} (MISC)O \textbf{Festival Eurovisão da Canção} (MISC) 1974 ocorreu em \textbf{Brighton} (LOC) , no \textbf{Reino Unido} (LOC) . \textbf{Brighton} (LOC) é a parte mais importante da cidade de \textbf{Brighton} (LOC) e \textbf{Hove} (LOC) (formados a partir das cidades da antiga \textbf{Brighton} (LOC), \textbf{Hove} (LOC) , \textbf{Portslade} (LOC) e várias outras aldeias), em \textbf{East Sussex} (LOC) , no \textbf{Reino Unido} (LOC) na costa sul da \textbf{Inglaterra} (LOC) . O antigo povoado de \textbf{Brighthelmston} (LOC) data de antes do \textbf{Livro de Domesday} (MISC) (1086), mas emergiu como estância de saúde durante o século XVIII e tornou-se um destino turístico após a chegada da ferrovia em 1841. \textbf{Brighton} (PER) experimentou um rápido crescimento populacional atingindo um pico de mais de 160 000 em 1961. A moderna \textbf{Brighton} (LOC) faz parte de uma aglomeração se estende ao longo da costa, com uma população de cerca de 480 mil.

O festival em si realizou-se no \textbf{Brighton Dome} (MISC) , um espaço de artes que abriga o \textbf{Concert Hall} (LOC), o \textbf{Corn Exchange} (ORG) e o \textbf{Pavilion Theatre} (LOC). Todos os três locais estão ligados ao resto do \textbf{Royal Pavilion Estate} (ORG) por um túnel subterrâneo para o \textbf{Royal Pavilion} (ORG) em \textbf{Pavilion Gardens} (MISC) e através de corredores compartilhados para \textbf{Brighton Museum} (LOC), como todo o complexo foi construído para o príncipe regente (mais tarde rei \textbf{George IV} (PER) ) e concluída em 1805.



\textbf{Formato} (PER)Pela primeira vez, os vídeos promocionais das músicas selecionadas foram apresentados numa transmissão especial, transmitida simultaneamente na maioria dos países participantes. O programa, intitulado \textbf{Auftakt für Brighton} (MISC) (\textbf{Prelúdio} (ORG) para \textbf{Brighton} (PER)), foi coordenado pela \textbf{ARD} (ORG) , a televisão pública alemã, e apresentado por \textbf{Karin Tietze-Ludwig} (PER) . Foi o primeiro programa do tipo "preview" a ser transmitido em muitos países europeus simultaneamente (tradicionalmente cada emissora nacional organizava o seu próprio programa de visualização). O programa também foi notável por ser a estreia na televisão europeia dos vencedores, \textbf{ABBA} (ORG) , que foram particularmente creditados em prévias como "\textbf{The Abba} (MISC)".



\textbf{Votação Inicialmente} (MISC), uma extensão do método de votação utilizado desde 1971 seria aplicado. Esse método consistia em que cada país tivesse 10 jurados a pontuar cada canção de 1 a 5. Como, nos ensaio, viu-se que este método de votação faria com que a votação demorasse mais de uma hora, a \textbf{BBC} (ORG) , a 48 horas da emissão, fez voltar o método de votação utilizado até 1970 , fazendo com que várias delegações, incluindo a portuguesa e a britânica, protestassem contra a \textbf{BBC} (ORG) . Assim, cada país tinha um júri composto por dez membros, e cada um desses dez jurados atribuíam um ponto à sua canção preferida. Os diferentes júris foram contatados por telefone, de acordo com uma ordem determinada por sorteio, a primeira na história do concurso.

Antes do início da votação, \textbf{Katie Boyle} (PER) explicou duas outras regras. Primeiro, os jurados tiveram que manter suas deliberações longe de qualquer aparelho de rádio ou televisão. Eles não podiam ver o desenrolar da votação nem ouvir os resultados dos outros júris. Em segundo lugar, se a comunicação com um dos júris fosse perdida, seria lembrada no final da votação. Se fosse impossível restaurar a comunicação naquele momento, os votos do júri em questão seriam considerados nulos e sem efeito.

O supervisor delegado pela \textbf{UER} (ORG) foi, mais uma vez, \textbf{Clifford Brown} (PER) .

Durante a votação, a câmera fez vários close-ups dos artistas. Em particular, os \textbf{ABBA} (ORG) , \textbf{Gigliola Cinquetti} (PER) e \textbf{Maggie MacNeal} (PER) apareceram.



\textbf{Festival} (PER)O vídeo introdutório traçou brevemente a história da cidade de \textbf{Brighton} (LOC) , do século XIX até 1974 , através de postais, fotos e vídeos. Seguiram-se alguns pontos turísticos da cidade litorânea. A câmera então mostrou a sala, as cabines dos comentadores, o quadro de votação, o palco e finalmente a orquestra.

A apresentadora da edição foi novamente \textbf{Katie Boyle} (PER) , que apresentava o certame pela quarta e última vez, depois de 1960 , de 1963 e de 1968 , que se dirigiu aos espetadores em inglês e francês .

A orquestra, dirigida por \textbf{Ronnie Hazlehurst} (PER) , estava num fosso ao pé do palco; o quadro de votação e a mesa do supervisor, à direita do palco. O palco em si era composto por uma decoração de painéis com formas orgânicas e curvas, de cor neutra. O pano de fundo consistia numa cortina. Vários degraus, um fosso e um pódio deram continuidade à decoração. Os painéis estavam alinhados com lâmpadas. Durante as apresentações, manchas coloridas permitiram dar àao palco diferentes atmosferas, essencialmente azuladas, rosa e laranja.

Os cartões postais começaram com pontos turísticos dos países participantes. Eles então mostraram os artistas, dando entrevistas, ensaiando no palco do \textbf{Dome} (MISC) e andando por \textbf{Brighton} (PER) .

O intervalo foi ocupado pelo famoso grupo britânico \textbf{The Wombles} (PER) , personagens no programa de televisão infantil da \textbf{BBC} (ORG) . Os \textbf{Wombles} (MISC) cantaram a sua música, "\textbf{Remember You} (MISC) A Womble", enquanto passeavam por \textbf{Brighton} (PER) para conhecer as crianças da cidade. Uma sequência mostrou-os no porto da cidade numa lancha. No final do vídeo, um \textbf{Womble} (PER) entrou no palco, com um cartaz a dizer "\textbf{Vote for the Wombles} (MISC)" (\textbf{Votem} (PER) nos \textbf{Wowbles} (LOC)), oferecendo depois uma rosa à apresentadora, \textbf{Katie Boyle} (PER) , que exclamou: "Bem, claro, nunca esquecerei que sou uma \textbf{Womble} (PER) agora".

Inicialmente, uma extensão do método de votação utilizado desde 1971 seria aplicado. Esse método consistia em que cada país tivesse 10 jurados a pontuar cada canção de 1 a 5. Como, nos ensaio, viu-se que este método de votação faria com que a votação demorasse mais de uma hora, a \textbf{BBC} (ORG) , a 48 horas da emissão, fez voltar o método de votação utilizado até 1970 , fazendo com que várias delegações, incluindo a portuguesa e a britânica, protestassem contra a \textbf{BBC} (ORG) . Assim, cada país tinha um júri composto por dez membros, e cada um desses dez jurados atribuíam um ponto à sua canção preferida. Os diferentes júris foram contatados por telefone, de acordo com uma ordem determinada por sorteio, a primeira na história do concurso.

Antes do início da votação, \textbf{Katie Boyle} (PER) explicou duas outras regras. Primeiro, os jurados tiveram que manter suas deliberações longe de qualquer aparelho de rádio ou televisão. Eles não podiam ver o desenrolar da votação nem ouvir os resultados dos outros júris. Em segundo lugar, se a comunicação com um dos júris fosse perdida, seria lembrada no final da votação. Se fosse impossível restaurar a comunicação naquele momento, os votos do júri em questão seriam considerados nulos e sem efeito.

O supervisor delegado pela \textbf{UER} (ORG) foi, mais uma vez, \textbf{Clifford Brown} (PER) .

Durante a votação, a câmera fez vários close-ups dos artistas. Em particular, os \textbf{ABBA} (ORG) , \textbf{Gigliola Cinquetti} (PER) e \textbf{Maggie MacNeal} (PER) apareceram.Pela primeira vez, um sorteio foi feito para determinar a ordem de votação. Assim, ordem de votação no Festival Eurovisão da Canção 1974, foi a seguinte:

Em baixo encontra-se a lista de maestros que conduziram a orquestra, na respectiva actuação de cada país concorrente.



Festival Eurovisão da Canção 1975



\textbf{Países Participantes Países} (LOC) que participaram no passado, mas não em 1975

O \textbf{Festival Eurovisão da Canção} (MISC) 1975 (em inglês : \textbf{Eurovision Song Contest} (ORG) 1975 , em francês : \textbf{Concours Eurovision} (MISC) de la chanson 1975 e em sueco : \textbf{Eurovisionens Melodi Festival} (MISC) 1975 ) foi o 20 Festival Eurovisão da Canção e realizou-se em 22 de março de 1975 em \textbf{Estocolmo} (LOC) . \textbf{Karin Falck} (PER) foi a apresentadora do festival que foi ganho pela banda \textbf{Teach-In} (ORG) que representou os \textbf{Países Baixos} (LOC) , com a canção " \textbf{Ding-a-dong} (PER) ". Esta edição do festival é marcada pela consequência de diversos conflitos políticos que ocorriam a época. Um dos casos mais notórios é o de \textbf{Portugal} (LOC) , que na altura vivia em pleno \textbf{PREC} (MISC) (\textbf{Processo Revolucionário} (MISC) em \textbf{Curso} (LOC),que é o período de transição entre a \textbf{Revolução dos Cravos} (ORG) e a \textbf{Constituinte de 1976} (LOC) ), foi representado pelo \textbf{Capitão Duarte Mendes} (PER) , um dos "\textbf{Capitães de Abril} (MISC)". \textbf{Mendes} (PER),queria se apresentar de uniforme militar,algo que quebrava as regras do concurso,já que a \textbf{EBU} (MISC) prefere distanciar-se dos acontecimentos políticos que acontecem no mundo durante a época do \textbf{Festival} (PER).Para não infrigir o regulamento,\textbf{Mendes} (PER) se apresentou com um cravo na lapela.

Foi neste ano que o sistema de votação ainda hoje em vigor foi criado, com os países a pontuarem as canções com 12, 10, 8, 7, 6, 5, 4, 3, 2 e 1 ponto, por ordem de preferência.



\textbf{Local} (MISC)O \textbf{Festival Eurovisão da Canção} (MISC) 1975 ocorreu em \textbf{Estocolmo} (LOC) , na \textbf{Suécia} (LOC) . \textbf{Estocolmo} (LOC) é a capital e maior cidade da \textbf{Suécia} (LOC) . É a sede do governo sueco, representado na figura do \textbf{Riksdagen} (LOC) , o parlamento nacional do país, além de ser a residência oficial dos membros da monarquia sueca. \textbf{Estocolmo} (LOC) é o maior e mais importante centro urbano , cultural , político , financeiro , comercial e administrativo da \textbf{Suécia} (LOC) desde o século XIII . Sua localização estratégica sobre 14 ilhas no centro-sul da costa leste da \textbf{Suécia} (LOC), ao longo do \textbf{Lago Malar} (LOC) , é historicamente importante. Uma vez que a capital sueca está situada sobre ilhas conhecidas por sua beleza, a cidade é destino de turistas de todo o mundo, tendo sido apelidada nos últimos anos de " \textbf{Veneza do Norte} (LOC)" . \textbf{Estocolmo} (LOC) é conhecida pelos seus edifícios e monumentos extremamente bem preservados, por seus arborizados parques , por sua riquíssima vida cultural e gastronômica, e pela gigantesca qualidade de vida que oferece a seus moradores. Há décadas, \textbf{Estocolmo} (LOC) figura como uma das cidades mais visitadas dos países nórdicos , com mais de um milhão de turistas internacionais anualmente. Nos últimos anos, tem sido citada entre as cidades mais "habitáveis" do mundo, sendo uma das mais limpas , organizadas e seguras do mundo.

O local escolhido foi o \textbf{Stockholmsmässan} (LOC) , que é o maior pavilhão de feiras e exposições do país. O hall central foi construído em 1941 e encontra-se em \textbf{Älvsjö} (LOC) , um subúrbio no sul de \textbf{Estocolmo} (LOC) .A época a capacidade deste hall era de 4 mil pessoas.



\textbf{Formato} (PER)Com o regresso de \textbf{Malta} (LOC) e \textbf{França} (LOC) , a estreia da \textbf{Turquia} (LOC) (o que forçou a retirada da \textbf{Grécia} (LOC) ,motivada pela \textbf{Invasão turca de Chipre} (MISC) no ano anterior).

Esta edição é marcada pela aumento das medidas de segurança,devido a avisos feitos pelo exército do \textbf{Reino Unido} (LOC) de que eventualmente o festival seria alvo de atentados terroristas. No entanto, foi a embaixada da \textbf{República Federal da Alemanha} (LOC) em \textbf{Estocolmo} (LOC) , que foi vítima, um mês depois, do ataque e de uma sangrenta tomada de reféns.

Antes mesmo do inicio do festival,as autoridades suecas estavam insatisfeitas com o fato de que estariam arcando sozinhas com os custos do evento e queriam dividir os custos de forma justa entre todos os países participantes.Mas,as negociações não avançaram e a emissora local,foi forçada a arcar com os custos. Este fato causou diversas criticas por parte da opinião pública do país,o que o posteriormente causou uma onda de protestos e motivou a realização de um festival alternativo em \textbf{Estocolmo} (LOC) . O mote deste festival era de que a \textbf{Eurovisão} (LOC) havia se tornado algo "comercial" demais. Diversos artistas que estavam envolvidos com a competição foram convidados,mas devido aos compromissos contratuais não puderam participar.Alguns meses mais tarde,a \textbf{EBU} (MISC) foi informada de que devido aos altos custos relacionados a organização deste concurso a \textbf{SR} (PER) não iria participar do festival do ano seguinte.Isso forçou uma grande mudança de regras para o 1977 em que os custos de organização seriam repartidos entre as emissoras participantes.Mas,mesmo com essas mudanças nas regras a \textbf{SR} (PER) não transmitiu o festival de 1976

O sorteio da ordem de execução ocorreu em 24 de janeiro de 1975 na sede da \textbf{UER} (LOC) em \textbf{Genebra} (LOC) . Os ensaios começaram na quarta-feira, 19 de março de 1975 , três dias antes da final.



\textbf{Visual} (MISC)O filme introdutório apresentava a história da \textbf{Suécia} (LOC), desde dos vikings ao pin-up contemporâneo, terminando com uma vista do porto de \textbf{Estocolmo} (LOC) .

A orquestra, dirigida por \textbf{Mats Olsson} (PER) , estava num buraco ao pé do palco. O palco, desenhado por \textbf{Bo Ruben Hedwall} (PER) , tinha um fundo azul escuro, na frente do qual foram colocados cinco painéis móveis trapezoidais. Esses painéis estavam numa gradação de cores, variando do azul escuro ao branco, passando também por bandas de céu azul, azul escuro e prata.

A apresentadora foi \textbf{Karin Falck} (PER) . Ela falou aos espectadores em sueco , inglês e francês , apresentando algumas dificuldades neste dois últimos idiomas.

Os cartões postais mostravam os artistas que participavam numa oficina de pintura. Numa tela branca, eles foram levados a pintar um auto retrato com as cores da bandeira de seus países

O intervalo foi ocupado por um vídeo intitulado "\textbf{The World of John Bauer} (MISC) " (\textbf{O Mundo de John Brauer} (MISC)). Foi uma sucessão de imagens, folheando várias pinturas do pintor sueco \textbf{John Bauer} (PER) , todas expostas no museu de \textbf{Jönköping} (LOC) , a cidade berço do artista. O acompanhamento musical foi feito ao vivo pela orquestra.

O festival durou duas horas e 14 minutos, sendo a primeira vez em que o mesmo ultrapassou duas horas. Ao mesmo tempo, foi transmitido para 34 países, mais notavelmente para os países da então \textbf{Cortina de Ferro} (PER) , \textbf{Hong Kong} (LOC) , \textbf{Islândia} (LOC) , \textbf{Japão} (LOC) , \textbf{Jordânia} (LOC) , entre outros. Os técnicos da televisão sueca recusaram-se a transmitir o festival via satélite para o \textbf{Chile} (LOC) , onde o \textbf{TVN} (MISC) (membro associado da \textbf{União Europeia de Radiodifusão} (ORG) ) estava se preparando para a transmissão. A recusa por parte dos técnicos suecos foi uma forma de protesto à \textbf{Ditadura} (ORG) militar de \textbf{Augusto Pinochet} (PER) , que regia o país desde o 1973



\textbf{Votação} (MISC)

Depois de diversas controvérsias e polémicas, um novo sistema de votação foi criado. Assim, cada país tinha um júri composto por 11 elementos (seis tinham de ter entre 16 e 26 anos e cinco entre 26 e 55), em que cada jurado teve que alocar entre um e cinco pontos para cada música (excepto obviamente a do seu país),sendo vedadas abstenções. No fim, a canção melhor classificada receberia 12 pontos, a segunda 10, a terceira 8, a quarta 7, e assim sucessivamente até à décima mais votada, que receberia 1 ponto. Os primeiros "doze ponto" da história da competição foram concedidos pelos \textbf{Países Baixos} (LOC) ao \textbf{Luxemburgo} (LOC) . Ao passar do tempo,este procedimento provou ser mais longo, mas mais justo e emocionante para o seu desenvolvimento.

Ao contrário do que acontece atualmente,a pontuação não era divulgada em ordem crescente (de 1 a 12),mas sim na ordem na qual as canções se apresentavam.O atual processo de divulgação das notas começaria três anos mais tarde em \textbf{1980,em Haia} (LOC) . Esta forma seria mantida até 1996 , mesmo havendo mudanças no número de membros do júri (entre 1975 e 1987 eram 11 e entre 1987 e 1996 eram 16) e em alguns anos as notas dadas por cada jurado variavam entre 1 a 10. Em 1997,os jurados de cinco países foram substituídos pelo televoto e no período de 1998 a 2008 todos os países foram incentivados,quando possível a realizarem suas votaçõesp or meio de televoto.

Na final da edição de 2009, os júris foram reintroduzidos e tinham o peso de 50\% das decisões finais de cada país participantes,os outros 50\% vinham do televoto.Mas mesmo assim,o sistema "12 pontos" foi mantido até 2015.Em 2016,um novo sistema de votação foi introduzido,em que o jurado técnico de cada país e o tele-voto operam de forma separada,dando dois conjuntos de 12 por país,entretanto,ele foi modelado após este processo.

O supervisor delegado pela \textbf{UER} (ORG) foi, mais uma vez, \textbf{Clifford Brown} (PER) .

Devido a problemas técnicos,foram registrados dois grandes problemas durante a votação.O primeiro foi no momento em que o júri israelense seria contactado e o sinal telefônico caiu diversas vezes. Em segundo lugar, uma falha eletrônica no placar,quando o júri monegasco anunciou que dava 3 pontos a \textbf{Turquia} (LOC),mas que foram marcados para \textbf{Israel} (LOC).

Por diversos momentos da votação, a apresentadora mostrou algumas dificuldade com o novo método,changando ao ponto de não usar o francês por diversos momentos chegando a perguntar para o supervisor "Como é que se diz sete em francês?".

Uma grande novidade neste ano,foi a transmissão dos artistas na "\textbf{green} (MISC) room",durante vários momentos,o foco era nos artistas que lideravam a votação como foi o caso de \textbf{Géraldine} (ORG) , \textbf{The Shadows} (ORG) e \textbf{Teach-In} (MISC) apareceram.

Ao contrário do ano anterior, a ordem de votação correspondeu à ordem de atuação.



\textbf{Festival} (MISC)A ordem de votação no \textbf{Festival Eurovisão da Canção} (MISC) 1975, foi a seguinte:

Os países que receberam 12 pontos foram os seguintes:

Em baixo encontra-se a lista de maestros que conduziram a orquestra, na respectiva actuação de cada país concorrente.



Festival Eurovisão da Canção 1976



\textbf{Países Participantes Países} (LOC) que participaram no passado, mas não em 1976

O \textbf{Festival Eurovisão da Canção 1976} (MISC) (em inglês : \textbf{Eurovision Song Contest} (ORG) 1976 , em francês : \textbf{Concours Eurovision} (MISC) de la chanson 1976 e em holandês : \textbf{Eurovisiesongfestival} (MISC) 1976 ) foi o 21 \textbf{Festival Eurovisão da Canção} (MISC) e realizou-se a 3 de abril de 1976 em \textbf{Haia} (LOC) , nos \textbf{Países Baixos} (LOC) . \textbf{Corry Brokken} (PER) foi a apresentadora, tornando-se assim na primeira ex-participante a viver a experiência de apresentar o \textbf{Festival Eurovisão da Canção} (MISC). O concurso foi ganho pelo grupo \textbf{Brotherhood of Man} (MISC) que representou o \textbf{Reino Unido} (LOC) , com a canção " \textbf{Save Your Kisses for Me} (MISC) ", que se converteu num grande êxito, vendendo seis milhões de exemplares, sendo o tema vencedor de \textbf{Eurovisão} (LOC) mais vendido até hoje.



\textbf{Local} (MISC)O \textbf{Festival Eurovisão da Canção} (MISC) 1976 ocorreu em \textbf{Haia} (LOC) , nos \textbf{Países Baixos} (LOC) . \textbf{Haia} (LOC) é a terceira mais populosa cidade dos \textbf{Países Baixos} (LOC) (depois de \textbf{Amesterdão} (LOC) e \textbf{Roterdão} (LOC) ), com uma população de 489 375 (2010) (população da área metropolitana: 600 000) e com uma área aproximada de 100 quilômetros quadrados. Está localizada no oeste do país, na província da \textbf{Holanda do Sul} (LOC) , da qual também é capital. A cidade da \textbf{Haia} (LOC), assim como \textbf{Amesterdão} (LOC), \textbf{Roterdão} (LOC) e \textbf{Utrecht} (LOC) , é parte do conglomerado urbano de \textbf{Randstad} (LOC) , com uma população cerca de 7,6 milhões habitantes. \textbf{A Haia} (LOC) é a sede de facto do governo do país: todavia, oficialmente, não é a capital dos \textbf{Países Baixos} (LOC), pois, de acordo com a constituição, a capital é \textbf{Amesterdão} (LOC). \textbf{A Haia} (LOC) é a sede do \textbf{Eerste Kamer} (LOC) (primeira câmara) e da \textbf{Tweede Kamer} (PER) (segunda câmara), respetivamente as câmaras alta e baixa, que formam o \textbf{Staten Generaal} (LOC) (literalmente, os "\textbf{Estados Gerais} (LOC)"). O rei \textbf{Guilherme Alexandre} (PER) dos Países Baixos vive e trabalha na \textbf{Haia} (LOC). Todas as embaixadas e ministérios estão localizados na cidade, assim como a \textbf{Hoge Raad der Nederlanden} (PER) (\textbf{A Suprema Corte} (MISC)), o \textbf{Raad van State} (PER) (\textbf{Conselho do Estado} (LOC)) e muitas organizações lobistas .

O festival em si realizou-se no \textbf{Congresgebouw} (PER) , este era um local de concertos e centro de convenções em \textbf{Haia} (LOC) , \textbf{Países Baixos} (LOC) , perto dos edifícios do \textbf{Tribunal Penal Internacional} (ORG) para a \textbf{ex-Jugoslávia} (LOC) , da \textbf{Organização para a Proibição de Armas Químicas} (ORG) e um dos escritórios administrativos do \textbf{Bacharelado Internacional} (LOC) . Foi inaugurado em 1969 e foi projetado no estilo funcional holandês pelo arquiteto \textbf{Jacobus Johannes Pieter Oud} (PER). O seu filho, \textbf{Hans Oud} (PER) , completou a construção após a morte de seu pai em 1963 . Em 2006 , uma parte do centro de convenções, incluindo o \textbf{Statenhal} (LOC), foi demolida para dar lugar ao edifício da Europol . Muitos concertos e festivais foram realizados lá antes, como o \textbf{Festival de Jazz do Mar do Norte} (MISC), e os \textbf{Festivais Eurovisão da Canção de 1976} (MISC) e 1980.



\textbf{Formato Suécia} (PER),\textbf{Malta} (LOC) e \textbf{Turquia} (LOC) se retiraram da competição,reduzindo o número de países participantes para 16.Este número subiu para 18,quando em um estágio tardio \textbf{Áustria} (LOC) e \textbf{Grécia} (LOC) ,anunciaram seus retornos.

\textbf{Malta} (LOC) iria participar com a canção " \textbf{Sing Your Song} (MISC), \textbf{Country Boy} (MISC) ", interpretada por \textbf{Enzo Guzman} (PER) , mas desistiu por razões desconhecidas, O país só voltaria ao certame em 1991 .

No ano anterior ,um debate acalorado surgiu na \textbf{Suécia} (LOC) sobre a participação do país na \textbf{Eurovisão} (LOC) devido aos custos organizacionais e operacionais daquela edição e com isso a emissora local não teria condições de bancar a participação sueca do ano seguinte.Além disso,a \textbf{SR} (PER) ,estava sob fogo cruzado da opinião pública local,sob as alegações de que a \textbf{Eurovisão} (LOC) se tornou "comercial demais". Somados estes fatores,não restou a emissora desistir de sua participação. Como resultado, a \textbf{EBU} (MISC) decidiu modificar as regras do concurso: a partir de agora, cada emissora teria que pagar uma parte. Isso permitiu que a emissora anfitriã não assumisse mais os custos organizacionais completos.

A \textbf{Turquia} (LOC) seria o participante de número 17,mas foi forçada a desistir de sua participação em protesto a participação grega,a música " \textbf{Panagia Mou} (MISC), \textbf{Panagia Mou} (MISC) " (\textbf{Minha} (PER) senhora,minha senhora) que falava abertamente da \textbf{Invasão Turca} (MISC) no \textbf{Chipre} (LOC) ,ocorrida em 1974 e criticava explicitamente a política externa do país.A emissora turca TRT transmitiu a final em delay no dia 3 de abril daquele ano.Entretanto,a participação grega foi censurada e durante a sua apresentação foi apresentada a canção "\textbf{Memleketim} (MISC)" (literalmente "\textbf{Minha terra",que} (MISC) é a versão turca da tradicional música yiddish "\textbf{Rabbi Elimelekh} (PER)") que se tornou um dos hinos da invasão turca na ilha.Devido a isto,a \textbf{Turquia} (LOC) ficou ausente do festival por dois anos .

Uma outra regra importante do concurso foi alterada: pela primeira vez, a reprodução instrumental foi permitida, mas desde que certas passagens da composição não pudessem ser reproduzidas pela orquestra.

Notavelmente,esta edição teve um dos maiores esquemas de segurança da história do \textbf{Festival} (LOC) até então.Foram montadas diversas barreiras de segurança no perímetro do \textbf{Centro de Congressos} (LOC) e mais de 200 policiais e militares estavam a paisana pela cidade.Para evitar desgastes,as delegações estavam hospedadas em um hotel dentro do complexo e iriam para o \textbf{Congresgebouw} (ORG) por um túnel conectado aos bastidores.O controle de entrada era tão restrito que as autoridades sabiam a identidade de cada pessoa que estava dentro do auditório.



\textbf{Visual} (MISC)O vídeo introdutório mostrou paisagens turísticas dos \textbf{Países Baixos} (LOC) , intercaladas com placas e sinalizações de trânsito que indicam a direção da cidade, aviões, comboios e autocarros e, finalmente, os episódios da vida social e cultural. Também foi mostrada a chegada da rainha \textbf{Juliana dos Países Baixos} (PER) à abertura anual dos \textbf{Estados Gerais} (LOC) que acontecem anualmente na cidade. O vídeo termina com várias vistas exteriores e interiores do \textbf{Congresgebouw} (ORG).

A orquestra, dirigida por \textbf{Jan Steulen} (PER) , ficava ao pé do palco e cercada por um recinto luminoso. O quadro de votação e a mesa do supervisor foram instalados à direita do palco. Combinou três elementos distintos: Primeiro, no chão, um círculo enorme e de cor preta. A metade da frente do círculo foi fixada e decorada com um círculo vermelho menor. A metade traseira tinha três partes móveis, retrateis no chão e cada uma alinhada com uma faixa de luz. Essas três partes podem surgir de forma independente, para criar um pódio de um, dois ou três pessoas. Em segundo lugar, pendurado, um quatro com barras verticais de forma retangular, sete barras horizontais de forma trapezoidal, um semicírculo e um disco. As quatro barras verticais formaram a armação do palco. As sete barras horizontais e o semicírculo se juntaram para criar uma forma oval. Todos esses elementos eram articulados, móveis e iluminados por dentro. Para cada atuação, formaram figuras diferentes e assumiram novas cores. Terceiro, um fundo neutro que também tomou cores diferentes para cada atuação. O palco foi da autoria de \textbf{Roland de Groot} (PER) , que já tinha projetado o palco do \textbf{Festival Eurovisão da Canção 1970} (MISC) .

A apresentadora foi \textbf{Corry Brokken} (PER) , vencedora do \textbf{Festival Eurovisão da Canção 1957} (MISC) , que falou aos espectadores em holandês , inglês e francês .

Todos os participantes receberam flores ao final de suas participações

Os cartões postais mostravam os artistas passeando em locais turísticos de seu país, geralmente na capital. Todos os vídeos também incluíram um plano de perfil do participante.

O intervalo foi ocupado pelo \textbf{The Dutch Swing College Band} (ORG) , um grupo cujos membros eram de \textbf{Aruba} (LOC) ,dos \textbf{Países Baixos Caribenhos} (LOC) e do \textbf{Suriname} (LOC) .

Pela primeira vez na história do concurso, os artistas eram entrevistados nos bastidores. Durante o intervalo, o apresentador \textbf{Hans von Willigenburg} (PER) fez duas perguntas, querendo saber se estavam satisfeitos com a sua atuação e qual música preferiam. Todos os artistas se declararam felizes pela sua atuação. \textbf{Catherine Ferry} (PER) disse que a sua música favorita era a do \textbf{Reino Unido} (LOC) ; \textbf{Ruthie Holzman} (PER) e \textbf{Braulio} (PER) elogiaram os seus concorrentes; \textbf{Waterloo \& Robinson} (MISC) elogiaram a \textbf{Jugoslávia} (ORG) ; e \textbf{Pierre Rapsat} (PER) hesitou entre \textbf{Mónaco} (LOC) , a \textbf{França} (LOC) , \textbf{Portugal} (LOC) e \textbf{Grécia} (LOC) .



\textbf{Votação} (MISC)Cada país tinha um júri composto por 11 elementos, que atribuiu 12, 10, 8, 7, 6, 5, 4, 3, 2, 1 pontos às dez canções mais votadas.

O supervisor delegado pela \textbf{UER} (ORG) foi, mais uma vez, \textbf{Clifford Brown} (PER) .

Durante a votação do júri francês, a escrutinadora da \textbf{EBU} (MISC) acidentalmente não contabilizou os 4 pontos dados à \textbf{Jugoslávia} (ORG), assim, durante muitos anos,só eram contabilizados os 6 pontos dados por outros jurados.

Durante a votação,as cameras focalizavam aqueles artistas que lideravam a votação naquele momento,o que aconteceu com \textbf{Catherine Ferry} (PER) , \textbf{Brotherhood of Man} (MISC) , \textbf{Mariza Koch} (PER) , \textbf{Mary Christy} (PER) , \textbf{Al Bano} (PER) e \textbf{Romina Power} (PER) .



\textbf{Festival} (PER)A ordem de votação no \textbf{Festival Eurovisão da Canção 1976} (MISC), foi a seguinte:

Os países que receberam 12 pontos foram os seguintes:

Em baixo encontra-se a lista de maestros que conduziram a orquestra, na respectiva actuação de cada país concorrente.



Festival Eurovisão da Canção 1977



\textbf{Países Participantes Países} (LOC) que participaram no passado, mas não em 1977

O \textbf{Festival Eurovisão da Canção 1977} (MISC) (em inglês : \textbf{Eurovision Song Contest} (ORG) 1977 e em francês : \textbf{Concours Eurovision} (MISC) de la chanson 1977 ) foi o 22 \textbf{Festival Eurovisão da Canção} (MISC) e realizou-se em 7 de maio de 1977 em \textbf{Londres} (LOC) . A apresentadora foi a jornalista \textbf{Angela Rippon} (PER) . O \textbf{Festival} (LOC) foi ganho pela portuguesa \textbf{Marie Myriam} (PER) que representava a \textbf{França} (LOC) com a canção " \textbf{L'oiseau et l'enfant} (MISC) " (O pássaro e a criança). Originalmente,a data prevista para o festival era o dia 2 de abril de 1977,mas devido a uma greve de funcionários da \textbf{BBC} (ORG),o festival foi adiado para o dia 7 de maio.

Neste ano, a \textbf{Tunísia} (LOC) iria participar pela primeira vez no certame e se apresentaria na posição de número 4. Mas, por razões desconhecidas, o país acabou desistindo de sua participação.



\textbf{Local} (MISC)O \textbf{Festival Eurovisão da Canção} (MISC) 1977 ocorreu \textbf{Londres} (LOC) no \textbf{Reino Unido} (LOC) , que já tinha sido sede das edições de 1960 , de 1963 e de 1968 . \textbf{Londres} (LOC) é a segunda maior \textbf{Europa} (LOC) e durante muitos anos foi a maior cidade do mundo . De uma cidade romana chamada \textbf{Londínio} (LOC) , a capital da região administrativa romana \textbf{Britannia} (LOC), a cidade acabou por se tornar o centro do império britânico, contribuindo hoje com 17\% do PIB de um país que é a quarta maior economia do mundo. \textbf{Londres} (LOC) tem sido um dos mais importantes centros da política e do comércio mundiais por mais de 2 mil anos(apesar de a capital do \textbf{Reino Unido} (LOC) ter sido \textbf{Winchester} (LOC) durante grande parte da \textbf{Idade Média} (MISC) ).

O local escolhido para o certame foi o recém-inaugurado \textbf{Centro de Convenções de Wembley} (LOC) . Este foi o primeiro centro multiúso do \textbf{Reino Unido} (LOC) e foi demolido em 2006.



\textbf{Formato} (PER)A regra de os cantores interpretarem os temas na língua oficial de cada país voltou depois de em 1973 a 1976 ter sido permitido cantarem em qualquer língua (ou seja os países escandinavos e do norte da \textbf{Europa} (LOC) preferiam cantar em inglês). Contudo, a \textbf{Alemanha} (LOC) e a \textbf{Bélgica} (LOC) tiveram a possibilidade de cantarem em inglês , porque já tinham escolhido as canções antes de a regra entrar em vigor.

A data inicial para a final foi o sábado 2 de abril de 1977 . Mas,devido a uma greve dos funcionários da \textbf{BBC} (ORG),nem mesmo a final britânica pode ser realizada.Ao entender que esta edição poderia ser a primeira da história do certame a correr risco de cancelamento,outras emissoras se propuseram a organizar o certame,o que foi o caso dos \textbf{Países Baixos} (LOC) que estava propondo realizar a segunda edição consecutiva no país,propondo \textbf{Amsterdã} (LOC) ,a \textbf{Espanha} (LOC) ,que estava novamente propondo \textbf{Málaga} (LOC) e \textbf{Portugal} (LOC) com \textbf{Lisboa} (LOC) . Faltando 3 dias para a primeira data marcada,a \textbf{BBC} (ORG) e a \textbf{EBU} (MISC) anunciaram oficialmente que o \textbf{Festival} (LOC) seria realizado no dia 7 de maio,cinco semanas depois da data original. No entanto, os sindicatos e técnicos holandeses declararam sua solidariedade aos colegas britânicos e ameaçaram entrar em greve por sua vez. Finalmente, um acordo, mediado pela \textbf{UER} (ORG) e pela \textbf{BBC} (ORG) foi efetivado.



\textbf{Contexto Geopolítico} (MISC)Novamente,uma música portuguesa esteve envolvida em várias polêmicas,a maior parte delas causadas pelo chefe de produção da \textbf{BBC} (ORG), \textbf{Bill Cotton} (PER) ,que se demonstrou constrangido pelo fato de que a música portuguesa estar se referindo publicamente "às novas liberdades", declarando que a política entrou de vez no contexto do festival,isto aconteceu pelo fato de que o letrista da música ser \textbf{Ary dos Santos} (PER) , um comunista assumido. No entanto, este mesmo havia estado presente em diversas edições do \textbf{Festival} (LOC),o que foi questionado pela imprensa portuguesa.Quando questionado por parte da imprensa internacional defendeu-se, dizendo: "a nossa canção é sobre morrer pelo seu país e renascer novamente em liberdade".

Uma outra controvérsia aconteceu com a delegação austríaca que era representada pelo grupo \textbf{Schwetterlinge} (LOC) ,que assumidamente tinha tendências esquerdistas. No entanto,eles já haviam anunciado que não iriam fazer nenhuma manifestação política durante a final,pois, a música não tinha esse contexto.Cabe ressaltar,que este até então era o ano mais controverso da história do Festival. Além da desconfiança de que a ideologia de esquerda estava presente nas as músicas de \textbf{Portugal} (LOC) e \textbf{Áustria} (LOC),houve a greve na \textbf{BBC} (ORG) e pairavam suspeitas de que para abafar a greve,a \textbf{BBC} (ORG) corrompeu os jurados de diversos países para lhe dar a vitória,algo que quase foi comprovado durante a votação,quando durante uma boa parte da votação,o \textbf{Reino Unido} (LOC) esteve na liderança do placar.



\textbf{Visual} (MISC)Ao contrário das edições anteriores,a produção foi realizada de uma forma mais simples,o vídeo introdutório apresentou cada um países constituintes do \textbf{Reino Unido} (LOC) apresentados por tomadas aéreas. Cada país foi introduzida pela sua bandeira e brasão. O vídeo termina com uma vista do \textbf{Centro de Conferências de Wembley} (LOC) .

A orquestra, dirigida por \textbf{Ronnie Hazlehurst} (PER) , estava localizada num buraco na parte de trás do palco. Esta cova foi cercada por uma trilha circular e encimada por um arco móvel, formado por três arcos. Dois arcos menores, localizados em ambos os lados do buraco, foram integrados ao fundo. Na frente, três círculos sobrepostos acolheram os artistas. Todos os elementos da decoração eram de cor neutra e forrados com faixas de luz rosa. Holofotes de cores tornaram possível dar todos os tons de azul celeste, azul escuro, rosa ou laranja. A mesa do supervisor e o quadro de votação ficaram à direita do palco, no alto de uma escada. Durante o intervalo, a faixa circular ao redor da orquestra foi decorada com as bandeiras dos dezoito países participantes.

A apresentadora foi a jornalista \textbf{Angela Rippon} (PER) , que falou aos espectadores em inglês e francês .

Devido as dificuldades durante a produção,a filmagem dos cartões postais foi cancelada e pela primeira vez desde 1973,eles não exisitiram,durante o tempo de transição de cada uma das performances,foram mostradas imagens do público no local.

O intervalo foi ocupado por um filme feito durante a semana de ensaios nos bares de \textbf{Londres} (LOC), onde é possível ouvir música jazz.



\textbf{Votação} (MISC)Cada país tinha um júri composto por 11 elementos, que atribuiu 12, 10, 8, 7, 6, 5, 4, 3, 2, 1 pontos às dez canções mais votadas.

O supervisor delegado pela \textbf{UER} (ORG) foi, mais uma vez, \textbf{Clifford Brown} (PER) , que teve de intervir três vezes durante a votação,todas relacionadas a falhas técnicas,que atrapalharam a condução de \textbf{Angela Rippon} (LOC),quando a mesma teve que fazer as conexões relacionadas a língua francesa. Em primeiro lugar,durante a divulgação dos resultados do \textbf{juri monegasco} (PER),que deu 3 pontos aos \textbf{Países Baixos} (LOC),a mesa atribuiu dois pontos em vez de três, sem que o escrutinador interviesse.Isso provocou uma reação inesperada do porta-voz neerlandês,que pediu para que os pontos acumulados por seu país possem retificados imediatamente.Após essa situação,houve uma discussão entre \textbf{Rippon} (LOC) e \textbf{Brown} (LOC).\textbf{Rippon} (PER) contrariada propôs que a votação continuasse normalmente. O erro foi finalmente corrigido apenas depois dos pontos do \textbf{Reino Unido} (LOC) terem sido revelados. Então, o porta-voz britânico teve que ser interrompido pelo supervisor.Assim, a mesa teve que interromper imediatamente a votação. Para piorar,a situação, o porta-voz do júri israelense acabou se esquecendo de divulgar que a \textbf{Suíça} (LOC) tinha ganhado 4 pontos de seu país. O supervisor então enfatizou seu esquecimento e a pontuação foi corrigida.

Durante a votação, a câmera fez vários close-ups dos artistas. Em particular,aqueles que durante algum momento lideraram a votação como: \textbf{Lynsey de Paul} (PER) , \textbf{Mike Moran} (PER) e \textbf{Marie Myriam} (PER) apareceram.

Pela primeira vez, a host-broadcaster colocou as bandeiras dos países participantes no placar.



\textbf{Festival} (MISC)A ordem de votação no \textbf{Festival Eurovisão da Canção 1977} (MISC), foi a seguinte:

Os países que receberam 12 pontos foram os seguintes:

Em baixo encontra-se a lista de maestros que conduziram a orquestra, na respectiva actuação de cada país concorrente.



Festival Eurovisão da Canção 1978



\textbf{Países Participantes Países} (LOC) que participaram no passado, mas não em 1978

O \textbf{Festival Eurovisão da Canção 1978} (MISC) (em inglês : \textbf{Eurovision Song Contest} (ORG) 1978 e em francês : \textbf{Concours Eurovision} (MISC) de la chanson 1978 ) foi o 23 \textbf{Festival Eurovisão da Canção} (MISC) e realizou-se em 22 de abril de 1978 em \textbf{Paris} (LOC) . Os apresentadores do festival foram \textbf{Denise Fabre} (PER) e \textbf{Leon Zitrone} (PER) , a primeira vez que duas pessoas apresentaram o certame e a primeira vez desde 1956 que um homem apresentou o festival. O festival foi ganho por \textbf{Izhar Cohen \& The Alphabeta} (ORG) que representaram \textbf{Israel} (LOC), com a canção " A-Ba-Ni-Bi ".

Durante a votação, ao se configurar uma vitória israelita , diversos países do \textbf{Norte da África} (LOC) e do \textbf{Oriente Médio} (ORG) cortaram de forma abrupta a transmissão. Um dos casos mais conhecidos foi o da televisão da \textbf{Jordânia} (LOC) , que ao cortar o sinal da transmissão colocou no ar a foto de um ramo de narcisos e em seguida anunciou a vitória da \textbf{Bélgica} (LOC) (que na verdade foi a segunda colocada). Nos dias seguintes, a imprensa do país se negava a reconhecer a vitória israelense. Esta foi a primeira vez em que se falou do conflito entre palestinos e israelenses no festival.

A \textbf{Noruega} (LOC) acabou em último lugar pela quinta vez, sendo o primeiro país a receber o enfadonho "\textbf{null points} (MISC)" no então sistema de votação, mas o intérprete, \textbf{Jahn Teigen} (PER) , tornou-se popular no seu país natal por cauda desse feito.



\textbf{Local} (MISC)O \textbf{Festival Eurovisão da Canção 1978} (MISC) ocorreu em \textbf{Paris} (LOC) , em \textbf{França} (LOC) . \textbf{Paris} (LOC) é a capital e a mais populosa cidade da \textbf{França} (LOC) , bem como a capital da região administrativa de \textbf{Ilha de França} (LOC) . A cidade se situa em um dos meandros do \textbf{Sena} (LOC) , no centro da bacia parisiense, entre os confluentes do \textbf{Marne} (LOC) e do \textbf{Sena} (LOC) rio acima, e do \textbf{Oise} (ORG) e do \textbf{Sena} (LOC) rio abaixo. Como a antiga capital de um império estendido pelos cinco continentes , é, hoje, a capital do mundo francófono . A posição de \textbf{Paris} (LOC) numa encruzilhada entre os itinerários comerciais terrestres e fluviais no coração de uma rica região agrícola a tornou uma das principais cidades da \textbf{França} (LOC) ao longo do século X , beneficiada com palácios reais, ricas abadias e uma catedral . Ao longo do século \textbf{XII} (MISC) , \textbf{Paris} (LOC) se tornou um dos primeiros focos europeus do ensino e da arte. Ao fixarem-se os \textbf{Reis de França} (MISC) e, pois, também a corte (o que incluía grande parte da alta nobreza francesa ), na cidade, sua importância económica e política não cessou de crescer. Assim, no início do s século XIV , \textbf{Paris} (LOC) era a mais importante cidade de todo o mundo ocidental . No século XVII , era a capital da maior potência política europeia; no século XVIII , era o centro cultural da \textbf{Europa} (LOC) e, no século XIX , era a capital da arte e do lazer, a \textbf{Meca da Belle Époque} (LOC) . Sua arquitetura, seus parques, suas avenidas e seus museus fazem-na, pelo ano de 2004, a cidade mais visitada do mundo francófono, com cerca de 25 milhões de turistas , aproximadamente 500 000 a mais do que em 2003, segundo a \textbf{Secretaria de Turismo} (MISC) e de \textbf{Congressos de Paris} (LOC). As margens parisienses do \textbf{Sena} (LOC) foram inscritas, em 1991, na lista do \textbf{Património Mundial da Organização das Nações Unidas para a} (MISC) \textbf{Educação} (LOC), a \textbf{Ciência} (LOC) e a \textbf{Cultura} (LOC) . \textbf{Paris} (LOC) é a capital económica e comercial da \textbf{França} (LOC), onde os negócios da \textbf{Bolsa} (ORG) e das finanças se concentram. A densidade da sua rede ferroviária, rodoviária e da sua estrutura aeroportuária -- um hub da rede aérea francesa e europeia -- fazem-na um ponto de convergência para os transportes internacionais. Essa situação resultou duma longa evolução, em particular das concepções centralizadoras das monarquias e das repúblicas, que dão um papel considerável à capital do país e, nela, tendem a concentrar, ao extremo, todas as instituições. Desde os anos 1960, os governos sucessivos têm desenvolvido políticas de desconcentração e de descentralização a fim de reequilibrar o país. Abrigando numerosos monumentos, por seu considerável papel político e económico, \textbf{Paris} (LOC) é também uma cidade importante na história do mundo . Símbolo da cultura francesa , a cidade atrai quase 30 milhões de visitantes por ano, ocupando, também, um lugar preponderante no mundo da moda e do luxo. Em 2007, a população intramuros (dentro do limite dos antigos muros ) de \textbf{Paris} (LOC) era de 2 193 031 habitantes pelo recenseamento do \textbf{Instituto Nacional de Estatísticas e Estudos Econômicos} (ORG) . Porém, ao longo do século XX , a área metropolitana de \textbf{Paris} (LOC), se desenvolveu largamente fora dos limites da comuna original. A \textbf{Grande Paris} (LOC) é, com seus 11 836 970 habitantes, uma das maiores aglomerações urbanas da \textbf{Europa} (LOC) e da \textbf{União Europeia} (ORG). Com um PIB de 813 364 milhões de dólares a \textbf{Região Parisiense} (LOC) é um ator econômico europeu de primeira grandeza, sendo a primeira região econômica europeia.

O festival em si realizou-se no \textbf{Palais des Congrès} (MISC) , um centro de eventos, espetáculos e convenções localizado na cidade de \textbf{Paris} (LOC), \textbf{França} (LOC), construído pelo arquiteto \textbf{Guillaume Gillet} (PER) e inaugurado em 1974 .



\textbf{Formato Nesta} (PER) edição participaram 20 países,não havendo nenhuma desistência entre a edição do ano anterior e esta.Dois países retornaram ao \textbf{Festival:a Dinamarca} (MISC) e a \textbf{Turquia} (LOC) .Esta edição acabou entrando para a história da \textbf{Eurovisão} (LOC) como o festival com maior número de países participantes até então. Pela primeira vez este festival foi apresentado por um duo, cabendo assim a \textbf{Denise Fabré} (PER) e a \textbf{Léon Zitrone} (PER) a apresentação deste certame.

O sueco \textbf{Björn Skifs} (PER) estava insatisfeito com a regra de que todos os países teriam que interpretar as suas canções na sua língua oficial. Ele planejava cantar em inglês de qualquer maneira, mas mudou de ideia no último momento, fazendo com que ele se esquecesse da letra. Por isso, ele murmurou as primeiras linhas antes de encontrar as palavras novamente.

Pela primeira vez, \textbf{Grécia} (LOC) e \textbf{Turquia} (LOC) participaram na mesma edição, desde a \textbf{Invasão turca de Chipre} (MISC) de 1974 . Em 1975 , com a estreia da \textbf{Turquia} (LOC), a \textbf{Grécia} (LOC) desistira em protesto, acontecendo o mesmo com a \textbf{Turquia} (LOC) em 1976 e 1977.

Juntamente com os 20 países participantes, o certame também foi transmitido ao vivo na então \textbf{Jugoslávia} (ORG) , \textbf{Tunísia} (LOC) , \textbf{Argélia} (LOC) , \textbf{Marrocos} (LOC) , \textbf{Jordânia} (LOC) , \textbf{Alemanha Oriental} (ORG) , \textbf{Polónia} (LOC) , \textbf{Hungria} (LOC) , \textbf{Checoslováquia} (LOC) , \textbf{Dubai} (LOC) , \textbf{Hong Kong} (LOC) , \textbf{União Soviética} (ORG) e \textbf{Japão} (LOC) .



Visual

O vídeo introdutório mostrou pontos turísticos de \textbf{Paris} (LOC) . Uma imagem final do \textbf{Arco do Triunfo} (LOC) incluía uma visão da orquestra e do palco.

A orquestra, dirigida por \textbf{François Rauber} (PER) , estava localizada na parte de trás do palco, numa estrutura móvel de branco, em forma de barco, simbolizando o brasão de armas de \textbf{Paris} (LOC) . Esta foi a única vez na história da competição em que a orquestra se movimentou durante a transmissão, a estrutura girando sobre si mesma no início de cada apresentação. O palco apresentava um pódio principal composto de três níveis de \textbf{lajes} (LOC) brancas emolduradas por duas bordas negras. As lajes eram iluminadas por dentro. Um segundo pódio circular, redondo e de cor preta, ficava à direita do primeiro. Foi destinado para os coristas. A decoração apresentava em primeiro plano dois arcos assimétricos, encontrando-se acima do pódio principal e sustentando uma gigantesca bola de espelhos. A parte superior das asas desses arcos foi decorada com asas brancas. O fundo do palco, de cor neutra, foi enquadrado por três pilares verticais, quatro barras horizontais e duas formas côncavas. Todo o palco foi iluminado com tons de rosa, vermelho e azul. Finalmente, à esquerda do palco estava a área dos apresentadores e à direita o quadro de votação e a mesa do supervisor.

Os apresentadores foram \textbf{Denise Fabre} (PER) e \textbf{Léon Zitrone} (PER) , que falaram aos espectadores em francês e inglês . Pela primeira vez, duas pessoas apresentaram o certame e a primeira vez desde 1956 um homem apresentou o festival.

Pela primeira vez, cartões postais foram filmados ao vivo. Eles mostraram a entrada dos artistas. Eles usaram um corredor, em seguida, um elevador que os levou ao palco. Lá, eles saudaram os participantes anteriores. A câmera também filmou várias cenas do público, incluindo \textbf{Jane Birkin} (PER) e \textbf{Serge Gainsbourg} (PER) .

O intervalo foi ocupado por um video de tributo para o violinista \textbf{Stéphane Grappelli} (PER) , por ocasião do seu aniversário. Na primeira parte, \textbf{Grappelli} (PER) tocou no violino uma peça de jazz instrumental intitulada "\textbf{My Heart Stood Still} (MISC)". Ele foi acompanhado ao piano por \textbf{Oscar Peterson} (PER) , na bateria de \textbf{Kenny Clarke} (PER) e no baixo por \textbf{Niels-Henning Ørsted Pedersen} (ORG) . Na segunda parte, tocou no violino uma peça clássica chamada "\textbf{Pick Yourself} (MISC) Up"". Ele foi acompanhado no violino por \textbf{Yehudi Menuhin} (PER) .



\textbf{Votação} (MISC)Cada país tinha um júri composto por 11 elementos, que atribuiu 12, 10, 8, 7, 6, 5, 4, 3, 2, 1 pontos às dez canções mais votadas.

O supervisor executivo da \textbf{EBU} (MISC) foi, pela primeira vez, \textbf{Frank Naef} (PER) .

Durante a votação, a câmera fez vários close-ups dos artistas. Em particular, \textbf{Carole Vinci} (PER) , \textbf{Joël Prévost} (PER) , \textbf{Jean Vallée} (PER) e \textbf{Izhar Cohen \& the Alphabeta} (ORG) apareceram.

Pela primeira vez, um país (neste caso, \textbf{Israel} (LOC) ) recebeu cinco pontuações máximas consecutivas, um recorde para a época, apenas igualado em 1997 ( \textbf{Reino Unido} (LOC) ) e 2012 ( \textbf{Suécia} (LOC) ).



\textbf{Festival} (PER)A ordem de votação no \textbf{Festival Eurovisão da Canção 1978} (MISC), foi a seguinte:

Os países que receberam 12 pontos foram os seguintes:

Em baixo encontra-se a lista de maestros que conduziram a orquestra, na respectiva actuação de cada país concorrente.



Festival Eurovisão da Canção 1979



\textbf{Países Participantes Países} (LOC) que participaram no passado, mas não em 1979

O \textbf{Festival Eurovisão da Canção 1979} (MISC) (em inglês : \textbf{Eurovision Song Contest} (ORG) 1979 , em francês : \textbf{Concours Eurovision} (MISC) de la chanson 1979 e em hebraico :  1979 ) foi o 24 Festival Eurovisão da Canção , que se realizou em 31 de março de 1979 em \textbf{Jerusalém} (LOC) , a última vez que o certame se realizou em março . Os apresentadores foram \textbf{Daniel Pe'er} (PER) e \textbf{Yardena Arazi} (PER) . Representando \textbf{Israel} (LOC) , \textbf{Gali Atari} (ORG) e \textbf{Milk and Honey} (ORG) foram os vencedores desse ano ao cantarem o tema " \textbf{Hallelujah} (LOC) ".

A apresentação do festival deste ano foi original, pois os países foram apresentados com um sketch mostrando as características de cada país. \textbf{Portugal} (LOC) por exemplo surgiu com barcos com pescadores, e claro o célebre vinho do \textbf{Porto} (LOC) , não podia faltar.

A \textbf{Turquia} (LOC) , por razões políticas desistiu de participar, outros países muçulmanos pressionaram a \textbf{Turquia} (LOC) para não participar, pois o festival realizava-se em \textbf{Israel} (LOC) . O tema que seria o 11 a subir ao palco, intitulava-se " \textbf{Seviyorum} (LOC) " e era interpretada por \textbf{Maria Rita Epik} (PER) \& 21.Peron .



\textbf{Local} (MISC)O \textbf{Festival Eurovisão da Canção} (MISC) 1979 ocorreu em \textbf{Jerusalém} (LOC) , em \textbf{Israel} (LOC) . \textbf{Jerusalém} (LOC) (em hebraico :  , \textbf{Yerushaláyim} (PER) ; em árabe :  , \textbf{al-Quds} (LOC) ; em grego : \textbf{} (LOC) , \textbf{Ierossólyma} (LOC) ), localizada em um planalto nas montanhas da \textbf{Judeia} (LOC) entre o \textbf{Mediterrâneo} (LOC) e o \textbf{mar Morto} (LOC) , é uma das cidades mais antigas do mundo . É considerada sagrada pelas três principais religiões abraâmicas -- judaísmo , cristianismo e islamismo . Israelenses e palestinos reivindicam a cidade como sua capital , mas \textbf{Israel} (LOC) mantém suas principais instituições governamentais em \textbf{Jerusalém} (LOC), enquanto o \textbf{Estado da Palestina} (LOC) , em última instância, apenas a prevê como a sua futura sede política; nenhuma das reivindicações, no entanto, é amplamente reconhecida pela comunidade internacional.

O festival em si realizou-se no \textbf{Centro Internacional de Convenções} (ORG) , uma sala de concerto e centro de convenções , que é o maior do \textbf{Médio Oriente} (LOC) .



\textbf{Formato Israel} (PER) , tendo vencido a edição de 1978 , assumiu a organização da edição de 1979. Esta foi a primeira vez que a competição foi realizada geograficamente fora do continente europeu e também foi a última vez que a final aconteceu em março. Outra grande novidade é que esta foi a primeira produção de televisão a cores a ser produzida e exibida no país.

Grupos religiosos ortodoxos repetidamente protestaram contra a organização do concurso, pois o contrato assinado entre a \textbf{IBA} (ORG) e a \textbf{EBU} (MISC) determinava a realização do certame em um sábado, que é o dia dia do \textbf{Shabat} (LOC) , dia de descanso consagrado na religião judaica . Fora isto, o contexto internacional não ajudava e eram recorrente as ameaças da \textbf{Organização para a Libertação da Palestina} (ORG) , que se seguiram à assinatura do \textbf{Tratado de paz israelo-egípcio} (MISC) , resultaram num fortalecimento drástico das medidas de segurança.

Por ser um país muçulmano, a \textbf{Turquia} (LOC) recebeu muita pressão dos países integrantes da \textbf{Liga Árabe} (ORG) para não enviar sua delegação para um evento que estava sendo em \textbf{Israel} (LOC) . A pressão deu resultado e faltando algumas semanas para o começo dos ensaios a participação turca foi cancelada. A música que se chama " \textbf{Seviyorum} (LOC) " e é interpretada por \textbf{Maria Rita Epik} (PER) \& 21.Peron ,havia sido sorteada para se apresentar na posição de número 11.A desistência do país consequentemente reduziu o número de países participantes de 20 para 19.

As novas tecnologias presentes a época, se fizeram presentes em dois momentos, até então inéditos na história do concurso. O primeiro foi quando tanto \textbf{IBA} (MISC) e \textbf{EBU} (MISC) decidiram fazer a sua transmissão por via satélite, algo extremamente avançado a época (ressaltando que este foi o primeiro concurso a usar essa tecnologia) e o segundo foi quando o grupo italiano \textbf{Matia Bazar} (PER) foi o primeiro ato que abriu mão da orquestra e optou por usar uma base musical pré-gravada para se apresentar (algo que se tornaria recorrente a partir dali). e



\textbf{Visual} (MISC)O vídeo introdutório mostrou pontos turísticos de \textbf{Jerusalém} (LOC) , seguidos por cenas da vida espiritual e religiosa das comunidades judaica, muçulmana e cristã da cidade.Ele se concluiu com imagens de diversos rostos dos habitantes da cidade.

A orquestra, dirigida pelo experiente maestro \textbf{Izhak Graziani} (PER) , estava localizada num fosso localizado ao pé do palco que teve decorações desenhadas por \textbf{Dov Ben David} (PER) . Este último tinha a forma de uma meia esfera aberta. O espaço dentro da meia esfera foi ocupado por uma gigantesca escultura móvel colocada num pódio de quatro degraus. Esta escultura era composta por três círculos concêntricos e duas bases circulares, feitas de metal cromado e que reproduziam o logotipo da emissora e desta edição. Finalmente, dois grupos de três contrafortes foram localizados em ambos os lados da parede do hemisfério. O fundo tomou cinco tons diferentes, dependendo da atuação: azul, amarelo, laranja, rosa e vermelho. O espaço dos apresentadores, decorado com plantas, ficava à esquerda do palco; enquanto que à direita estavam localizadas a mesa do escrutinador geral, do supervisor do evento e o placar que neste ano estava em uma escala menor do que o dos anos anteriores.

Os apresentadores foram \textbf{Daniel Pe'er} (PER) e cantora \textbf{Yardena Arazi} (PER) . , que falaram aos espectadores em hebraico , inglês e francês .

Os cartões postais mostraram os países apresentados com um sketch sobre os clichês de cada país. \textbf{Portugal} (LOC), por exemplo surgiu com barcos com pescadores, e o célebre vinho do \textbf{Porto} (LOC) .

O intervalo foi ocupado por um número de dança ,que foi produzido e executado por um grupo cultural especialmente formato para o evento que ganhou o nome de \textbf{Shalom} (PER) '79, e que teve a responsabilidade de apresentar um medley de 5 músicas tradicionais da cultura judaica,acompanhados por um grupo de dança.Em mais uma quebra de tradições vigentes,esta foi a primeira vez na história do festival em que o país organizador decidiu apresentar a sua cultura no ato de intervalo da grande final. Propositalmente,foi escolhida para terminar o ato, a tradicional \textbf{Hevenu Shalom Alekhem} (PER) (    ) e acompanhados pelo espírito acolhedor desta música, pela primeira vez desde a edição de 1961 todos os envolvidos no evento retornaram ao palco.



\textbf{Votação} (MISC)Cada país tinha um júri composto por 11 elementos, que atribuiu 12, 10, 8, 7, 6, 5, 4, 3, 2, 1 pontos às dez canções mais votadas.

O supervisor executivo da \textbf{EBU} (MISC) foi \textbf{Frank Naef} (PER) . Ao contrário dos anos anteriores,a votação deste ano foi marcada pela tensão e pela emoção,já que ao contrário dos anos anteriores,o vencedor não era claro até o último jurado. Nas três primeiras rodadas,foram registradas a liderança de três países diferentes (\textbf{França} (LOC),\textbf{Israel} (LOC) e \textbf{Reino Unido} (LOC)) e parecia que a vitória fosse se encaminhar novamente para o \textbf{Reino Unido} (LOC).Mas, \textbf{Israel} (LOC) conseguiu disparar em um certo momento e parecia ter a sua segunda vitória consecutiva consolidada. No entanto, esta vitória começou a ficar ameaçada após a divulgação dos resultados do sétimo júri nacional (\textbf{Grécia} (LOC)) em que o país começou a receber uma sequência de notas extremamente baixas, enquanto que a \textbf{Espanha} (LOC) que discretamente se posicionava em nono lugar e começava a receber uma sequência de notas altas vindas majoritariamente da \textbf{Grécia} (LOC),do \textbf{Benelux} (ORG),da \textbf{Europa Central} (LOC) e de alguns países da \textbf{Escandinávia} (LOC) e progressivamente foi subindo de forma vertiginosa para os primeiros lugares, assumindo a liderança no décimo primeiro júri (França),chegando em alguns momentos a liderar com folga,dando sinais de que a terceira vitória espanhola no certame seria evidente. No décimo quarto júri (\textbf{Países Baixos),Israel} (PER) perdeu o segundo lugar para a \textbf{França} (LOC) e ambos países ficaram nestas posições nos dois juris seguintes (\textbf{Suécia} (LOC) e Noruega).Em uma reviravolta do destino, ao chegar ao décimo quinto jurado (\textbf{Suécia),Israel} (PER) começou uma incrível sequência de 3 notas 12 (além de \textbf{Suécia} (LOC) e \textbf{Noruega} (LOC),houve o \textbf{Reino Unido),enquanto} (LOC) seus rivais mais próximos (\textbf{França} (LOC) e \textbf{Espanha} (LOC)) ganharavam notas muito baixas. E ao chegar ao final desta sequência, \textbf{Israel} (LOC) havia retomado a liderança, mas novamente com um ponto de diferença em relação a \textbf{Espanha} (LOC). No penúltimo júri (Áustria),a situação ainda se acirraria mais, pois a situação se inverteria novamente com a \textbf{Espanha} (LOC) novamente liderando por apenas um ponto. Com um clima extremamente tenso e para azar dos espanhóis, eles seriam o último júri a votar e caso \textbf{Israel} (LOC) recebesse uma pontuação superior a um ponto, o país ganharia o certame pela segunda vez e de forma seguida,conseguindo terceiro "back-to-back" (quando alguém ganha algum evento por dois anos consecutivos) da história (o primeiro havia sido a própria \textbf{Espanha} (LOC) em 1968 e em 1969 e o segundo havia sido o \textbf{Luxemburgo} (LOC) em 1972 e em 1973).O ar de mistério e de incerteza tomou conta do \textbf{Centro Internacional de Convenções de Jerusalém} (ORG) pelos próximos minutos. A situação estava tão imprevisível que durante a divulgação dos resultados deste júri ,a apresentadora \textbf{Yardena Arazi} (PER) se confundiu de tal forma que anunciou que \textbf{Portugal} (LOC) havia ganhado os 10 pontos do júri espanhol e erroneamente eles foram adicionados no placar, motivando a revisão e a correção dos resultados em seguida (\textbf{Ao invés} (MISC) dos 10 pontos, o júri espanhol tinha-lhe dado os 6 pontos, reduzido a sua pontuação para 64,mas deixando a nona posição intocada). Assim, o silêncio permaneceu no local até que o nome \textbf{Israel} (LOC) foi lido pelo porta-voz espanhol e em seguida "10 points". Naquele exato momento,\textbf{Israel} (ORG) havia ganho o \textbf{Festival Eurovisão da Canção} (MISC) pela segunda vez e desta vez em seu próprio solo, o que dava a condição da \textbf{Eurovisão} (LOC) permanecer por mais um ano no país, o que acabou não acontecendo.Alguns dias mais tarde, surgiram rumores junto a imprensa espanhola de que a votação do júri espanhol foi realizada de forma estratégica para que \textbf{Betty Missiego} (PER) não ganhasse, já que a \textbf{TVE} (ORG) havia divulgado que se em caso de uma eventual vitória espanhola ,a emissora não teria condições financeiras para organizar a próxima edição.

Os pontos foram repetidos em inglês , francês e hebraico .

Durante a votação, a câmera fez vários close-ups dos artistas. Em particular, \textbf{Gali Atari} (PER) e \textbf{Milk and Honey} (ORG) , \textbf{Betty Missiego} (PER) e \textbf{Anne-Marie David} (PER) apareceram.



\textbf{Festival} (MISC)A ordem de votação no \textbf{Festival Eurovisão da Canção 1979} (MISC), foi a seguinte:

Os países que receberam 12 pontos foram os seguintes:

Em baixo encontra-se a lista de maestros que conduziram a orquestra, na respectiva actuação de cada país concorrente.



Festival Eurovisão da Canção 1980



\textbf{Países Participantes Países} (LOC) que participaram no passado, mas não em 1980

O \textbf{Festival Eurovisão da Canção 1980} (MISC) (em inglês : \textbf{Eurovision Song Contest} (ORG) 1980 , em francês : \textbf{Concours Eurovision} (MISC) de la chanson 1980 e em holandês : Eurovisiesongfestival 1980 ) foi o 25 \textbf{Festival Eurovisão da Canção} (MISC) e realizou-se em 19 de abril de 1980 em \textbf{Haia} (LOC) . A apresentadora do festival foi \textbf{Marlous Fluitsma} (PER) , que foi ganho por \textbf{Johnny Logan} (PER) com a canção, \textbf{What's Another Year} (MISC) , uma balada romântica, que recebeu o troféu do grande prémio de \textbf{Marcel Bezençon} (PER) , fundador do concurso.

Neste festival, as canções foram apresentadas por apresentadores de cada dum dos países participantes na respetiva língua nacional.



\textbf{Local} (MISC)O \textbf{Festival Eurovisão da Canção} (MISC) 1980 ocorreu em \textbf{Haia} (LOC) , nos \textbf{Países Baixos} (LOC) . \textbf{Haia} (LOC) é a terceira mais populosa cidade dos \textbf{Países Baixos} (LOC) (depois de \textbf{Amesterdão} (LOC) e \textbf{Roterdão} (LOC) ), com uma população de 489 375 (2010) (população da área metropolitana: 600 000) e com uma área aproximada de 100 quilômetros quadrados. Está localizada no oeste do país, na província da \textbf{Holanda do Sul} (LOC) , da qual também é capital. A cidade da \textbf{Haia} (LOC), assim como \textbf{Amesterdão} (LOC), \textbf{Roterdão} (LOC) e \textbf{Utrecht} (LOC) , é parte do conglomerado urbano de \textbf{Randstad} (LOC) , com uma população cerca de 7,6 milhões habitantes. \textbf{A Haia} (LOC) é a sede de facto do governo do país: todavia, oficialmente, não é a capital dos \textbf{Países Baixos} (LOC), pois, de acordo com a constituição, a capital é \textbf{Amesterdão} (LOC). \textbf{A Haia} (LOC) é a sede do \textbf{Eerste Kamer} (LOC) (primeira câmara) e da \textbf{Tweede Kamer} (PER) (segunda câmara), respetivamente as câmaras alta e baixa, que formam o \textbf{Staten Generaal} (LOC) (literalmente, os "\textbf{Estados Gerais} (LOC)"). O rei \textbf{Guilherme Alexandre} (PER) dos Países Baixos vive e trabalha na \textbf{Haia} (LOC). Todas as embaixadas e ministérios estão localizados na cidade, assim como a \textbf{Hoge Raad der Nederlanden} (PER) (\textbf{A Suprema Corte} (MISC)), o \textbf{Raad van State} (PER) (\textbf{Conselho do Estado} (LOC)) e muitas organizações lobistas .

O festival em si realizou-se no \textbf{Congresgebouw} (ORG) , é um local de concertos e centro de convenções em \textbf{Haia} (LOC) , \textbf{Países Baixos} (LOC) , perto dos edifícios do \textbf{Tribunal Penal Internacional} (ORG) para a \textbf{ex-Jugoslávia} (LOC) , da \textbf{Organização para a Proibição de Armas Químicas} (ORG) e um dos escritórios administrativos do \textbf{Bacharelado Internacional} (LOC) . Foi inaugurado em 1969 e foi projetado no estilo funcionalismo holandês pelo arquiteto \textbf{Jacobus Johannes Pieter Oud} (PER). O seu filho, \textbf{Hans Oud} (PER) , completou a construção após a morte de seu pai em 1963 . Em 2006 , uma parte do centro de convenções, incluindo o \textbf{Statenhal} (LOC), foi demolida para dar lugar ao edifício da Europol . Muitos concertos e festivais foram realizados lá antes, como o \textbf{Festival de Jazz do Mar do Norte} (MISC), e os \textbf{Festivais Eurovisão da Canção de 1976} (MISC) e 1980.



\textbf{Formato Israel} (PER) , que venceu o festival em 1979 , renunciou ao direito de sediar o evento, no dia 21 de agosto do mesmo ano. A \textbf{IBA} (ORG) alegou que não poderia financiar outra produção internacional num curto espaço de tempo. Além disso, naquele ano, o governo israelita tinha cortado o orçamento do canal. Para piorar, a data marcada pela \textbf{União Europeia de Radiodifusão} (ORG) coincidiu com o \textbf{Yom Hazikaron} (MISC) que é o dia em memória as vitimas de guerra no país. Assim, \textbf{Israel} (LOC) foi forçado a desistir do festival nesse ano. Assim, a organização foi oferecida à \textbf{BBC} (ORG) , como foi feito em algumas edições anteriores. Sem embargo, a \textbf{BBC} (ORG) também renunciou ao evento, da mesma forma que a \textbf{TVE} (ORG) , também renunciou ---- apesar de que a \textbf{Secretária de Turismo da Costa del Sol} (LOC) já estava trabalhando com a possibilidade do evento ser realizado no recém inaugurado \textbf{Palácio de Congresos de Torremolinos} (LOC) em \textbf{Málaga} (LOC) . Após isso, o evento foi oferecido para todas as emissoras que participaram do ano anterior até que a televisão nacional dos \textbf{Países Baixos} (LOC) , a \textbf{NOS} (ORG), decidiu assumir a organização do evento com um prazo apertado e em escala reduzida. De acordo com \textbf{Yair Lapid} (PER) , filho de \textbf{Tommy Lapid} (PER) que era o diretor geral da \textbf{IBA} (MISC) a época, \textbf{Lapid} (PER) conseguiu convencer a \textbf{Televisão dos Países Baixos} (ORG), a assumir a ingrata missão de organizar evento de forma inesperada, quando ele percebeu que a realização consecutiva do evento poderia quebrar a emissora. .

O local que sediou o 1976 , o \textbf{Centro de Congressos dos Países Baixos} (LOC), foi escolhido. Além disso, a mesma equipe de produção chefiada pelo design \textbf{Roland de Groot} (PER) foi recontratada. Devido ao baixo orçamento disponível,a organização holandesa optou por reciclar diversas partes do vídeo de abertura e do cenário de 1976. Como nos festivais de 1977 e 1978, não houve vídeos pré-gravados para a apresentação dos participantes, sendo que a \textbf{NOS} (ORG) optou por usar apresentadores vindos de cada país participante. Apesar disso, a apresentadora \textbf{Marlous Fluitsma} (PER) apresentou o certame na sua maioria em holandês . A \textbf{NOS} (ORG) conseguiu gastar apenas \textbf{US\$} (LOC) 725 mil, um valor muito abaixo dos valores gastos pelas emissoras que organizaram o festival nos anos anteriores.

Devido aos fatores supracitados,a organização local optou por não produzir os já conhecidos postcards e por ser a edição de número 25 do certame,a \textbf{EBU} (MISC) e a \textbf{NOS} (ORG) decidiram usar um introdutor para cada uma das 19 canções participantes. Os introdutores dirigiram-se aos telespectadores em uma das suas línguas nacionais A introdução da \textbf{Irlanda} (LOC) foi feita em gaélico e foi a única a usar um idioma diferente da música do seu país, que era em inglês. \textbf{Dinamarca} (LOC) , \textbf{Finlândia} (LOC) , \textbf{Portugal} (LOC) , \textbf{Suécia} (LOC) e \textbf{Turquia} (LOC) recorreram aos seus próprios comentaristas,enquanto outros 13 optaram por apresentadores de suas emissoras. A introdutora da participação local foi a própria apresentadora \textbf{Marlous Fluitsma} (PER) .

Enquanto as músicas iam sendo introduzidas, as fotos dos seus autores, compositores e artistas participantes passavam no canto inferior direito da tela. Essas fotos foram tiradas durante as viagens promocionais dos participantes no período prévio ao concurso em diversos locais dos \textbf{Países Baixos} (LOC) .



Visual

O vídeo introdutório começou com uma cena numa praia holandesa. Um violinista apareceu a tocar \textbf{Te Deum} (PER) por \textbf{Marc-Antoine Charpentier} (PER) . Das ondas, em seguida, emerge a sigla "25", lembrando o jubileu de prata do concurso que se comemorava naquele ano. Em seguida, o vídeo introdutório mostrado em 1976 foi mostrado.

A orquestra, dirigida por \textbf{Rogier van Otterloo} (PER) , estava novamente ao pé do palco e cercada por um recinto luminoso. O espaço dos introdutores, decorado com duas tulipas estilizadas, estava localizado à esquerda do palco; o quadro de votação e a mesa do supervisor, à direita. O palco, da autoria de \textbf{Roland de Groot} (PER) , foi emoldurado por um frontão à maneira dos templos gregos. O espaço era ocupado por um pódio de dois níveis, uma escada recuada de cinco degraus e três faixas horizontais curvas. Acima do palco havia móveis suspensos: seis meias elipses (duas das quais decoradas com um quarto de círculo vermelho) e dois trapézios. Todos esses elementos eram de cor neutra e revestidos de faixas claras. Eles mudavam, de acordo com as atuações, para cores creme, cinza, azul, rosa ou lilás.

A apresentadora foi \textbf{Marlous Fluitsma} (PER) , que falou aos espectadores em holandês , tendo usado o inglês e o francês apenas na votação.

Neste ano não houve os famosos postcards. Mas para marcar os vinte e cinco anos da competição, os organizadores pediram a um membro de cada delegação para apresentar a música do seu país. Os introdutores dirigiram-se aos telespectadores na sua própria língua nacional. A introdução da \textbf{Irlanda} (LOC) foi feita em gaélico e foi o único a usar um idioma diferente da música do seu país, que era em inglês. \textbf{Dinamarca} (LOC) , \textbf{Finlândia} (LOC) , \textbf{Portugal} (LOC) , \textbf{Suécia} (LOC) e \textbf{Turquia} (LOC) recorreram aos seus próprios comentadores, tendo os outros países recorrido a outros.

Enquanto as músicas iam sendo introduzidas, as fotos dos seus autores, compositores e artistas participantes passavam no canto inferior direito da tela. Essas fotos foram tiradas durante várias viagens nos \textbf{Países Baixos} (LOC) .

O intervalo foi ocupado pelo \textbf{The Dutch Swing College Band} (ORG) , um grupo cujos membros eram de \textbf{Aruba} (LOC) , \textbf{Países Baixos Caribenhos} (LOC) e do \textbf{Suriname} (LOC) . Eles interpretaram seu título de "\textbf{San Fernando} (LOC)", acompanhado pelos \textbf{Dançarinos} (LOC) de \textbf{Lee Jackson} (PER), enquanto atrás deles, os motivos do cenário permaneceram em movimento.

Pela segunda vez na história do concurso, os artistas foram entrevistados nos bastidores. Durante o intervalo, o apresentador \textbf{Hans von Willigenburg} (PER) falou com \textbf{Katja Ebstein} (PER) , \textbf{Sophie \& Magaly} (ORG) , o grupo \textbf{Prima Donna} (PER) , \textbf{Johnny Logan} (PER) , \textbf{Sverre Kjelsberg} (PER) e \textbf{Maggie MacNeal} (PER) , perguntando se já sabiam o nome do vencedor da noite. Todos responderam que não.



\textbf{Votação} (MISC)Cada país tinha um júri composto por 11 elementos, que atribuiu 12, 10, 8, 7, 6, 5, 4, 3, 2, 1 pontos às dez canções mais votadas.Esta foi a primeira vez que se anunciaram os resultados por ordem crescente (1,2,3...)

O supervisor executivo da \textbf{EBU} (MISC) foi \textbf{Frank Naef} (PER) , que teve que intervir três vezes para corrigir erros cometidos por \textbf{Marlous Fluitsma} (PER) .

Como um aceno às primeiras edições do concurso, a produção comunicou \textbf{Marlous Fluitsma} (PER) e os porta-vozes por meio de telefones, um diferente para cada país.

Durante a votação, a câmera fez vários close-ups dos artistas. Em particular, \textbf{Maggie MacNeal} (LOC) , \textbf{Katja Ebstein} (PER) e \textbf{Johnny Logan} (PER) apareceram.



\textbf{Festival} (MISC)A ordem de votação no \textbf{Festival Eurovisão da Canção 1980} (MISC), foi a seguinte:

Os países que receberam 12 pontos foram os seguintes:

Em baixo encontra-se a lista de maestros que conduziram a orquestra, na respectiva actuação de cada país concorrente.



Festival Eurovisão da Canção 1981



\textbf{Países Participantes Países} (LOC) que participaram no passado, mas não em 1981

O \textbf{Festival Eurovisão da Canção} (MISC) 1981 (em inglês : \textbf{Eurovision Song Contest} (ORG) 1981 , em francês : \textbf{Concours Eurovision} (MISC) de la chanson 1981 e em irlandês : \textbf{Comórtas Amhránaíochta} (LOC) na \textbf{hEoraifíse} (MISC) 1981 ) foi o 26 \textbf{Festival Eurovisão da Canção} (MISC) e realizou-se em 4 de abril de 1981 em \textbf{Dublin} (LOC) . A apresentadora foi Doireann Ní Bhriain . Os \textbf{Bucks Fizz} (ORG) foram os vencedores desse ano com a canção " \textbf{Making Your Mind} (MISC) Up ".



\textbf{Local} (MISC)O \textbf{Festival Eurovisão da Canção} (MISC) 1981 ocorreu em \textbf{Dublin} (LOC) , na \textbf{Irlanda} (LOC) . \textbf{Dublin} (LOC) é a capital e maior cidade da \textbf{Irlanda} (LOC) . O nome em inglês deriva da palavra irlandesa " \textbf{Dubhlinn} (LOC) " (ocasionalmente também grafada \textbf{Duibhlinn} (PER) ou \textbf{Dubh Linn} (PER) ), que significa "\textbf{Lago Negro} (LOC)". Localiza-se na província de \textbf{Leinster} (LOC) próxima ao ponto mediano da costa leste da \textbf{Irlanda} (LOC), sendo cortada pelo \textbf{Rio Liffey} (LOC) e o centro da região de \textbf{Dublin} (LOC). Desde 1898 possui nível administrativo de condado ( county-boroughs ). Seus limites são os condados de \textbf{Fingal} (LOC) a norte, \textbf{Dublin} (LOC) meridional a sudoeste e \textbf{Dun Laoghaire-Rathdown} (LOC) a sudeste. Tem uma população de 527 612 habitantes na cidade, e sua área metropolitana tem 1 804 156 habitantes. Fundada como um assentamento viquingue , foi o centro do \textbf{Reino de Dublin} (LOC) e se tornou a principal cidade da \textbf{Ilha} (LOC) após a invasão dos \textbf{Normandos} (LOC) . A cidade cresceu de maneira rápida durante o século XVII ; se tornou na época a segunda maior cidade do \textbf{Império Britânico} (ORG) e a quinta maior da \textbf{Europa} (LOC) . \textbf{Dublin} (LOC) entrou em um período de estagnação após o \textbf{Ato de União de 1800} (MISC) , mas continuou o centro econômico da \textbf{Ilha} (LOC). Após a \textbf{Partição da Irlanda} (MISC) em 1922, virou a capital do \textbf{Estado Livre Irlandês} (MISC) , e mais tarde, da \textbf{República da Irlanda} (LOC) . \textbf{Dublin} (LOC) é reconhecida como uma cidade global , com um ranking "\textbf{Alpha-} (MISC)", colocando a cidade entre as 30 mais globalizadas do mundo. Atualmente é o principal centro histórico, cultural, econômico, industrial e educacional da \textbf{Irlanda} (LOC).

O festival em si realizou-se no \textbf{RDS Simmonscourt Pavilion} (MISC) , uma sala polivalente, pertencente ao complexo de recreação da \textbf{Royal Dublin Society} (ORG) e destina-se a sediar festivais equestres e agrícolas. O pavilhão pode acomodar até 7 000 pessoas.



\textbf{Formato} (LOC) e visual

O vídeo introdutório foi sobre o encontro entre o passado e o presente da \textbf{Irlanda} (LOC) . Misturou imagens de vestígios arqueológicos, paisagens irlandesas, castelos, falésias, episódios de vida social e cultural, percorrendo assim a história do país, da Antiguidade ao período contemporâneo.

A orquestra de 46 membros, dirigida por \textbf{Noel Kelehan} (PER) , causou certos problemas para algumas canções por não ter nenhum instrumento de sopro, pois \textbf{Kelehan} (PER) não acreditava que o instrumento seria orquestral, o que causou grande preocupação com a música do \textbf{Reino Unido} (LOC) , pois um saxofone aparecia na parte instrumental da música. \textbf{Andy Hill} (PER) - o produtor do single - disse que, se soubesse, eles teriam retirado um dos dois coristas para ser substituídos por um saxofonista, pois na gravação no estúdio foram usados dois

Do ponto de vista da produção, esta edição é considerada um ponto sem volta na história, pois a \textbf{RTÉ} (LOC) aumentou de forma gigantesca o evento, e foi construído o maior palco até então. A decoração de \textbf{Michael Grogan} (PER) foi dividida em três partes distintas. Pela primeira vez a orquestra foi colocada em uma arquibanca. Esta parte foi delimitada por uma escadaria curva cujas bordas dos degraus eram luminosas. Os maestros tiveram que descer para chegar a um círculo. Acima da orquestra pendia uma faixa curva. No centro estava o palco em que os artistas se apresentavam. Consistia em três círculos interligados, sustentando cinco pódios circulares de diferentes alturas, interligados por degraus. \textbf{Círculos} (PER), degraus e pódios tinham bordas brilhantes. A decoração por trás dos artistas era uma representação estilizada de um motivo tradicional celta. Era uma faixa curva, dividida em quartos e sustentando quatro círculos simétricos. Este elemento foi atravessado por costelas, às quais foram anexadas lâmpadas. Todo o conjunto assumia cores diferentes, dependendo das atuações. À direita, finalmente, estava o comitê e o placar, localizado no topo de um arco curvo. Dois pódios foram desenhados, um para a apresentadora e outro para a mesa do supervisor. No arco, logo abaixo da mesa, havia uma porta de correr que deveria abrir no final da votação para deixar o vencedor entrar.

Outro ponto sem volta, foi a criação da "Green-room" (\textbf{Sala Verde} (LOC)), local onde os artistas participantes aguardavam a sua passagem no palco e os resultados. A cor verde está associada ao país e a esperança da vitória. Foi também criada a "\textbf{Eurovision Village} (MISC)", um local onde os artistas poderiam se divertir, interagir com o público e se apresentar com músicas de seu repertório. A \textbf{RTÉ} (MISC) queria mostrar o frenesi que a realização do evento no país e não somente a realização de um mero programa de televisão.

A apresentadora foi Doireann Ní Bhriain , que falou aos espectadores em irlandês , inglês e francês .

Os cartões postais começaram com uma visão de um globo estilizado. O país participante cresceu então, até ocupar a tela inteira. Vídeos então mostraram os artistas e compositores passeando por \textbf{Dublin} (LOC) , entre os ensaios.

O intervalo foi ocupado pelo grupo folk \textbf{Plantxy} (MISC) e sua composição "\textbf{Timedance} (PER)". Paralelamente ao vídeo introdutório, o intervalo encenou a cultura irlandesa através dos tempos e explorou em três pinturas,todos os aspectos do evento por meio de manifestações. A música foi especialmente projetada para a ocasião por \textbf{Bill Whelan} (PER) . A coreografia foi feita por \textbf{Iain Montague} (PER) e executada pelos bailarinos da \textbf{Dublin City Ballet} (LOC). A primeira pintura foi uma peça neolítica; a segunda, uma peça inédita composta no estilo dos harpistas tradicionais; a terceira, uma peça que destaca a influência da música contemporânea na música folclórica. Após a competição, o \textbf{Timedance} (MISC) alcançou o terceiro lugar em vendas recordes na \textbf{Irlanda} (LOC).



\textbf{Votação} (MISC)Cada país tinha um júri composto por 11 elementos, que atribuiu 12, 10, 8, 7, 6, 5, 4, 3, 2, 1 pontos às dez canções mais votadas.

A votação foi marcada por três situações. Em primeiro lugar, o júri austríaco votou com a ordem dos países alterados, e o supervisor (\textbf{Frank Naef} (PER)) pediu-lhe para dar os votos em ordem crescente. Em segundo lugar, quando a apresentadora fez o contato com o júri da \textbf{Jugoslávia} (ORG) e pediu os votos, a porta-voz ( \textbf{Helga Vlahović} (PER) , apresentadora da edição de 1990 ) respondeu "não os tenho" e a apresentadora agradeceu. Finalmente, num dos momentos em que a \textbf{Irlanda} (LOC), o país anfitrião, acrescentou pontos, por engano eles levantaram mais de 300 pontos para o placar, o que provocou o riso do público.

Durante a votação, a câmera fez vários close-ups dos artistas. Em particular, \textbf{Jean Gabilou} (PER) , \textbf{Sheeba} (ORG) , \textbf{Bucks Fizz} (ORG) e \textbf{Lena Valaitis} (PER) apareceram.



\textbf{Festival} (MISC)A ordem de votação no \textbf{Festival Eurovisão da Canção} (MISC) 1981, foi a seguinte:

Os países que receberam 12 pontos foram os seguintes:

Em baixo encontra-se a lista de maestros que conduziram a orquestra, na respectiva actuação de cada país concorrente.



Festival Eurovisão da Canção 1982



\textbf{Países Participantes Países} (LOC) que participaram no passado, mas não em 1982

O \textbf{Festival Eurovisão da Canção 1982} (MISC) (em inglês : \textbf{Eurovision Song Contest} (ORG) 1982 e em francês : \textbf{Concours Eurovision} (MISC) de la chanson 1982 ) foi o 27 Festival Eurovisão da Canção e realizou-se em 24 de abril de 1982 em \textbf{Harrogate} (LOC) , \textbf{Inglaterra} (LOC) . A apresentadora foi \textbf{Jan Leeming} (PER) . \textbf{Nicole} (LOC) foi a vencedora desse ano com a canção " \textbf{Ein bißchen Frieden} (MISC) " (Um pouco de paz).



\textbf{Local} (MISC)O \textbf{Festival Eurovisão da Canção 1982} (MISC) ocorreu no auditório principal do \textbf{Centro de Convenções Internacional de Harrogate} (LOC) , no \textbf{Reino Unido} (LOC) , o que, na altura do seu anúncio, causou uma certa surpresa, pois até então,o \textbf{Festival} (PER) nunca aconteceu em uma cidade de pequeno porte. A capacidade do auditório era de 2 mil pessoas. A ideia da \textbf{BBC} (ORG) era a de aproveitar a infraestrutura do seu centro de convenções recém inaugurado e que a época era um dos mais modernos do mundo. \textbf{Harrogate} (LOC) é uma cidade da \textbf{Inglaterra} (LOC) ( \textbf{Reino Unido} (LOC) ), com cerca 85 000 habitantes. Fica situada no condado de \textbf{North Yorkshire} (LOC) e a este da cidade de \textbf{York} (LOC) . Possui indústrias alimentares e é centro turístico e estância balnear (famosa pelas suas cerca 100 nascentes de águas sulfurosas e ferruginosas). A população está dividida em \textbf{High} (MISC) e \textbf{Low Harrogate} (ORG). A cidade fica perto de \textbf{York} (LOC) , \textbf{Ripon} (LOC) , \textbf{Knaresborough} (LOC) , \textbf{West Yorkshire} (LOC) , \textbf{Leeds} (LOC) e \textbf{Wetherby} (LOC) . Votada como a melhor povoação para viver em \textbf{Yorkshire} (LOC) , \textbf{Harrogate} (LOC) é uma considerada uma joia do norte inglês. O seu refinamento, apesar da sua proximidade a cidades industriais, reflete-se no bairro de \textbf{Montpellier} (LOC), que acolhe cerca de 80 lojas e converteu-se num centro de referência artística e de venda de antiguidades. É casa do famoso salão de chá \textbf{Bettys} (LOC), os seus banhos turcos são uma alternativa para esquecer o stress .



\textbf{Formato} (MISC)Após a mudança de pessoal na emissora pública francesa, como consequência das eleições do ano anterior na \textbf{França} (LOC) , os novos líderes da televisão pública francesa queriam introduzir a cultura em todos os nível e, por ser considerada algo "não cultural" pela emissora \textbf{TF1} (ORG) , a \textbf{França} (LOC) acabou a retirar-se do concurso. No entanto, a decisão da retirada foi anunciada em cima da hora e devido a falta de tempo útil, nenhuma outra emissora do sistema público de televisão pôde assumir a participação, nem transmitir o evento. Quando a decisão foi anunciada, esta caiu como uma bomba no país, levando a diversas ondas de manifestações públicas no país. \textbf{Jacques Lang} (PER) , então \textbf{Ministro da Cultura} (MISC), chegou a descrever o concurso como um monumento à idiotice , com o chefe de entretenimento \textbf{Pierre Bouteiller} (PER) a acompanhá-lo, dizendo: "A falta de talento e a mediocridade das músicas são incomodas." No ano seguinte, outro canal público tornou-se pela responsável pela participação, a \textbf{Antenne 2} (MISC) .

\textbf{A Grécia} (LOC) também se retirou, apesar de fazer parte da lista de participantes e ser 2 a subir ao palco. O tema que seria apresentado em \textbf{Harrogate} (LOC) intitulava-se " Sarantapente Kopelies " e era interpretada por \textbf{Themis Adamantidis} (PER) . Tal aconteceu porque a então ministra da cultura grega, a atriz \textbf{Melina Mercouri} (PER) , achou a música de gosto duvidoso. Por medo de que arranhasse a imagem do país no exterior, \textbf{Mercouri} (LOC) decidiu monocraticamente pela retirada grega.

A \textbf{Alemanha} (LOC) sagrou-se pela primeira vez como vencedora. A vitória de \textbf{Nicole} (LOC) foi muito previsível pois esta obteve 161 pontos contra os 100 pontos alcançados por \textbf{Avi Toledano} (PER) . No final, ao repetir a canção, cantou em inglês , francês , holandês e alemão , recebendo um estrondoso aplauso de toda a audiência.

A tradição dos vencedores do ano anterior, entregando o prêmio aos atuais vencedores, não foi seguida por \textbf{Bucks Fizz} (PER), vencedores em 1981.

A banda irlandesa \textbf{Chips} (MISC) perdeu na final nacional, que, se tivessem sido bem sucedidas, teriam levado à situação única de duas bandas na mesma edição com o mesmo nome (sendo a outra a \textbf{Suécia} (LOC)).

O participante da \textbf{Finlândia} (LOC), \textbf{Kojo} (PER) , veio com uma música de protesto contra as bombas nucleares e um apelo à paz. No entanto, não recebeu um único voto.

Apesar de a cidade não ter deixado a sua vida habitual e de o país estar a seguir atentamente a crescente tensão nas \textbf{Malvinas} (MISC) , que parecia estar a caminho da guerra, não deixou de haver uma recepção oficial. Esta aconteceu no castelo de \textbf{Howard} (LOC) e implicou uma viagem no tempo: as delegações foram transportadas em liteiras, havia empregados vestidos de acordo com a moda do século XVII e músicos a tocar canções dessa época.



Visual

O vídeo introdutório começou com uma visão de um mapa da \textbf{Europa} (LOC) desenhado num pergaminho. Os diferentes países participantes iluminados sucessivamente com laranja e em cada um deles, saía a mesma pergunta, formulada na língua nacional do país: "Onde está \textbf{Harrogate} (LOC)?". A câmera então fez uma visão ampliada do \textbf{Reino Unido} (LOC) , depois de \textbf{Yorkshire} (LOC) . \textbf{Harrogate} (LOC) apareceu no centro, marcado por um círculo laranja. O vídeo continuou com gravuras e pontos turísticos da cidade. Em seguida, foram mostrados os preparativos para o evento, as vitrines decoradas das lojas e uma exposição de flores. O vídeo concluiu com vistas noturnas do \textbf{Centro de Conferências} (LOC) e a chegada dos carros com as delegações ao local de competição

A orquestra, dirigida por \textbf{Ronnie Hazlehurst} (PER) , localizava-se à esquerda em duas arquibancadas. À direita estava o comitê de votação. A mesa do supervisor foi montada pela primeira vez na sala, de frente para o quadro de votações, e não mais no palco, debaixo do quadro de votações. No centro estava o pódio para os artistas. Consistia em quatro níveis, os dois primeiros mais estreitos, os dois últimos maiores. Esses níveis alternavam superfícies reflexivas e superfícies foscas. Os dois níveis superiores foram decorados com um quadrado brilhante. Todos foram alinhados com bandas de luz. O primeiro nível deu as boas-vindas à extremidade direita, o microfone do apresentador. O terceiro nível tinha dois pódios retangulares móveis, de cor branca e interligados entre si. O quarto nível apresentava duas decorações rotativas, consistindo de espelhos e prateleiras retangulares. O fundo foi decorado com elementos retangulares e pilastras luminosas. Todo o conjunto assumiu cores vermelho, azul e verde, dependendo das atuações.

A apresentadora foi \textbf{Jan Leeming} (PER) , que falou aos espectadores em inglês e francês e apresentou os países participantes, maestros e artistas.

Os cartões postais começaram com planos na bolha de cada emissora nacional, antes de entrar em pontos turísticos do país participante, com algumas notas do hino nacional de fundo. Os artistas foram mostrados em seu próprio país e depois em \textbf{Harrogate} (LOC) durante a semana do evento.

O intervalo foi ocupado com imagens de \textbf{Yorkshire} (PER) e \textbf{Castle Howard} (PER) , com os melhores momentos da recepção oficial a serem mostrados, com fogos de artifícios no final.



\textbf{Votação} (MISC)Cada país tinha um júri composto por 11 elementos, que atribuiu 12, 10, 8, 7, 6, 5, 4, 3, 2, 1 pontos às dez canções mais votadas.

Durante a votação, a câmera fez vários close-ups dos artistas. Em particular, \textbf{Nicole} (LOC) , \textbf{Ralph Siegel} (PER) , \textbf{Bill van Dijk} (PER) , \textbf{Aska} (PER) , \textbf{Bardo} (PER) , \textbf{Arlette Zola} (PER) , \textbf{Kojo} (PER) e \textbf{Anna Vissi} (PER) apareceram.

O público manifestou a sua surpresa duas vezes: quando a \textbf{Suécia} (LOC) atribuiu os seus "doze pontos" à \textbf{Jugoslávia} (ORG) e quando a \textbf{Áustria} (LOC) atribuiu apenas um ponto à \textbf{Alemanha} (LOC) .



\textbf{Festival} (PER)A ordem de votação no \textbf{Festival Eurovisão da Canção 1982} (MISC), foi a seguinte:

Os países que receberam 12 pontos foram os seguintes:

Em baixo encontra-se a lista de maestros que conduziram a orquestra, na respectiva actuação de cada país concorrente.



Festival Eurovisão da Canção 1983



\textbf{Países Participantes Países} (LOC) que participaram no passado, mas não em 1983

O \textbf{Festival Eurovisão da Canção 1983} (MISC) (em inglês : \textbf{Eurovision Song Contest} (ORG) 1983 , em francês : \textbf{Concours Eurovision} (MISC) de la chanson 1983 e em alemão : \textbf{Liederwettbewerb der Eurovision 1983} (MISC) ) foi o 28 Festival Eurovisão da Canção e realizou-se a 23 de abril de 1983 , em \textbf{Munique} (LOC) . \textbf{Marlene Charell} (PER) foi a apresentadora do evento, que foi ganho por \textbf{Corinne Hermès} (PER) representando o \textbf{Luxemburgo} (LOC) com a canção " \textbf{Si La Vie Est Cadeau} (MISC) ", uma balada romântica.

O concurso deste ano foi o primeiro a ser transmitido na \textbf{Austrália} (LOC), via canal 0/28 (agora SBS ) em \textbf{Sydney} (LOC) e \textbf{Melbourne} (LOC) , embora este país não tivesse até 2015 permissão para participar. No entanto, vencedores anteriores como os \textbf{ABBA} (ORG) e \textbf{Bucks Fizz} (ORG) tinham sido bem sucedidos na \textbf{Austrália} (LOC).

Devido à campanha eleitoral para as eleições legislativas portuguesas de 25 de abril , o certame foi transmitido em diferido \textbf{em Portugal} (LOC) . O tempo de antena dos partidos apenas acabou às 20h40, e só a essa hora arrancou a transmissão da \textbf{Eurovisão} (LOC), com comentários de \textbf{Eládio Clímaco} (LOC) . \textbf{Portugal} (LOC) foi o único país a não receber o evento em direto.



\textbf{Local} (MISC)O \textbf{Festival Eurovisão da Canção} (MISC) 1983 ocorreu em \textbf{Munique} (LOC) , na \textbf{Alemanha} (LOC) . \textbf{Munique} (LOC) é uma cidade da \textbf{Alemanha} (LOC) , capital do estado alemão da \textbf{Baviera} (LOC) , no sudeste do país. Conta atualmente cerca de 1,3 milhão de habitantes (2012 ), enquanto a sua região metropolitana , que engloba diversas cidades vizinhas ou próximas a \textbf{Munique} (LOC), abriga mais de 2,6 milhões de pessoas. É assim a cidade mais populosa da \textbf{Baviera} (LOC) e do sul da \textbf{Alemanha} (LOC) , e a terceira cidade mais populosa do país , depois da capital, \textbf{Berlim} (LOC) , e de \textbf{Hamburgo} (LOC) . \textbf{Munique} (LOC) é uma cidade independente ( \textbf{kreisfreie Stadt} (ORG) ) ou distrito urbano ( \textbf{Stadtkreis} (LOC) ), ou seja, possui estatuto de distrito ( \textbf{Kreis} (LOC) ). Adicionalmente, \textbf{Munique} (LOC) é também sede do governo do distrito administrativo da "\textbf{Alta Baviera} (LOC)" ( \textbf{Oberbayern} (LOC) em alemão ) bem como do distrito territorial ( \textbf{Landkreis} (LOC) ) de \textbf{Munique} (LOC). Cidades grandes próximas são \textbf{Zurique} (LOC) ( \textbf{Suíça} (LOC) ), a 315 km a oeste, \textbf{Praga} (LOC) (\textbf{República Checa} (LOC)), a cerca de 380 km a nordeste, \textbf{Viena} (LOC) ( \textbf{Áustria} (LOC) ) a cerca de 440 km a leste, \textbf{Milão} (LOC) (\textbf{Itália} (LOC)) a 490 km a sul e \textbf{Berlim} (LOC) , a cerca de 590 km a norte. Foi fundada em 1158. O número de habitantes da cidade de \textbf{Munique} (LOC) ultrapassou por volta de 1854 os cem mil, tendo nessa altura obtido o estatuto de cidade grande ( \textbf{Grosstadt} (LOC) ). A cidade foi destruída pela metade durante a \textbf{Segunda Guerra Mundial} (MISC) , porém reconstruída nas décadas posteriores ao fim do conflito. Desde os anos 1960 , alcançou a marca de um milhão de habitantes, estabelecendo-se desde então como a terceira mais populosa cidade alemã (entre os anos 60 e 80, a segunda ou a terceira mais populosa cidade da \textbf{Alemanha Ocidental} (LOC) ). \textbf{Munique} (LOC) é atravessada pelo \textbf{rio Isar} (LOC) . É em \textbf{Munique} (LOC) que é realizada anualmente a \textbf{Oktoberfest} (LOC) , uma tradicional festa alemã, que é a maior do mundo, sendo o evento um dos principais alicerces turísticos da \textbf{Alemanha} (LOC). A \textbf{Munique} (LOC) moderna é um importante e desenvolvido centro financeiro , urbano , logístico , cultural e político da \textbf{Alemanha} (LOC) e da \textbf{Europa} (LOC) continental . É sede de diversas empresas de renome mundial, incluindo a montadora \textbf{BMW} (ORG) . Entre 2011 e 2012, \textbf{Munique} (LOC) foi posicionada na 4 posição entre as "\textbf{Cidade Mais Habitáveis do Mundo} (LOC)" , segundo estudos da consultoria internacional \textbf{Mercer} (LOC) . A partir de 2006, o lema da cidade passou a ser "\textbf{München} (LOC) mag dich" ( \textbf{Munique} (LOC) gosta de ti (\textbf{Pt} (ORG)) ou \textbf{Munique} (LOC) ama você(Br) , em alemão ). Até 2005, o lema era "\textbf{Weltstadt mit Herz} (MISC)" ( \textbf{Cidade} (LOC) cosmopolita com coração ).

O festival em si realizou-se no \textbf{Rudi-Sedlmayer-Halle} (LOC) , uma arena multiusos, que, nos \textbf{Jogos Olímpicos} (MISC) de 1972 sediara principalmente os eventos de levantamento de peso.



\textbf{Formato 21} (MISC) países estiveram inicialmente inscritos para o concurso, com o regresso da \textbf{França} (LOC) , \textbf{Itália} (LOC) e \textbf{Grécia} (LOC) , o que faria desta edição a maior até então. No entanto, a emissora irlandesa \textbf{RTÉ} (MISC) desistiu do concurso devido a restrições financeiras, deixando 20 países para a final em \textbf{Munique} (LOC).

Os ensaios começaram na segunda-feira 18 de abril , com cada país tendo 40 minutos de ensaio. No final da semana, cada um teria um segundo ensaio de 20 minutos de duração.

O concurso era para durar 165 minutos, mas durou pouco mais de três horas, principalmente devido à apresentadora \textbf{Marlene Charell} (PER) ter que repetir todas as suas apresentações, e os anúncios de votação em três idiomas ( alemão , inglês e francês ).

O concurso deste ano foi o primeiro a ser transmitido na \textbf{Austrália} (LOC), via canal 0/28 (agora SBS ) em \textbf{Sydney} (LOC) e \textbf{Melbourne} (LOC) , embora este país não tivesse até 2015 permissão para participar. No entanto, vencedores anteriores como os \textbf{ABBA} (ORG) e \textbf{Bucks Fizz} (ORG) tinham sido bem sucedidos na \textbf{Austrália} (LOC).

Nesta edição do festival foi utilizado pela primeira vez microfones sem fio, e foi a primeira vez que a audiência ao vivo da \textbf{Eurovisão} (LOC) excedeu a capacidade de 10 000 pessoas.

O núcleo da equipa técnica foi formado por \textbf{Rainer Bertram} (PER) (direção), \textbf{Christian Hayer} (PER) e \textbf{Günther Lebram} (PER) (produção executiva), \textbf{Hans Gailling} (PER) e \textbf{Marlies Frese} (PER) (cenografia) e \textbf{Dieter Reith} (ORG) (direção musical).



\textbf{Visual} (MISC)A introdução durou quase quinze minutos, um dos mais longos de sempre , que consistia em duas partes. Primeiro, um vídeo foi projetado, começando com uma visão de mapa da \textbf{Europa} (LOC), expandindo para a \textbf{Alemanha Ocidental} (LOC) e \textbf{Munique} (LOC) , onde foram mostrados muitos pontos turísticos do país. O vídeo terminou com uma passagem aérea sobre a \textbf{Vila Olímpica} (ORG) em \textbf{Munique} (LOC) e um mapa de \textbf{Rudi-Sedlmayer-Halle} (ORG) . Em segundo lugar, todas as delegações subiram ao palco, uma a uma, até preencherem completamente o pódio.

A orquestra, dirigida por \textbf{Dieter Reith} (PER) , localizava-se num fosso, numa curva interna do arco do palco. O palco, da autoria de \textbf{Hans Gailling} (PER) , consistia num vasto pódio em forma de arco circular, com nove degraus nas suas extremidades. A decoração atrás do pódio consistia de três colunas de nove, onze e treze quadrados, respectivamente. Estes quadrados foram feitos de uma armação de metal polida que suporta fios de luz. Estes fios foram dispostos horizontalmente ou verticalmente, alternadamente quadrados. A borda do pódio e as extremidades dos quadrados tinham lâmpadas, que piscavam de forma a acompanhar a melodia de cada canção. Vislumbravam-se assim os primeiros avanços tecnológicos que incutiam ao certame e às interpretações maior brilho e vivacidade. O quadro de votação estava localizado à direita do palco. A mesa do supervisor não foi mostrada.Uma outra inovação técnica foi introduzida: o uso de microfones sem fio.

A apresentadora foi \textbf{Marlene Charell} (PER) , que falou aos espectadores em alemão , inglês e francês . Essa apresentação trilingue, que se manteria ao longo de toda a emissão, fez com que se ocorresse vários erros de linguagem, como trocar nomes de artistas e maestros (como \textbf{Charell} (PER) apresentou a cantora finlandesa \textbf{Ami Aspelund} (PER) como "\textbf{Ami Aspesund} (MISC)", apresentou o maestro norueguês \textbf{Sigurd Jansen} (PER) como "\textbf{... Johannes ...} (MISC) Skorgan ...", tendo sido forçada a inventar um nome no momento depois de esquecer o nome do maestro) e trocar nomes e países no momento da votação (conceder pontos a "\textbf{Schweden} (LOC)" ( \textbf{Suécia} (LOC) ) que foram originalmente dados a "\textbf{Schweiz} (LOC)" ( \textbf{Suíça} (LOC) ).

Ao contrário dos anos anteriores, não houve os tradicionais cartões postais. Em vez disso, e confrontada com a impossibilidade de filmar como tinha sido uma tradição na \textbf{Eurovisão} (LOC), a televisão alemã limitou-se a projetar o nome do país participante nas três línguas reguladoras numa imagem de fundo do palco. Os comentadores das diferentes televisões aproveitaram os quarenta e cinco segundos que esta sequência estática durou para fazer suas respectivas apresentações. Em seguida, \textbf{Marlene Charell} (PER) introduziu nos três idiomas, o título da música e os nomes dos autores, compositores, intérpretes e maestros. À sua direita estavam dispostos buquês de flores compostas pela própria, a lembrar as bandeiras dos países participantes.

O intervalo foi ocupado por número de dança definido para um medley de canções alemãs que se tornaram internacionalmente famosas, incluindo \textbf{Strangers in the Night} (MISC). A apresentadora, \textbf{Marlene Charell} (PER) , foi a bailarina principal.



\textbf{Votação} (MISC)Cada país tinha um júri composto por 11 elementos, que atribuiu 12, 10, 8, 7, 6, 5, 4, 3, 2, 1 pontos às dez canções mais votadas.

O supervisor executivo da \textbf{EBU} (MISC) foi \textbf{Frank Naef} (PER) , que apareceu no palco para parabenizar a realização do certame. Ele foi auxiliado em sua tarefa por \textbf{Marie-Claire Vionnet} (PER) .

Durante a votação, a câmera fez vários close-ups dos artistas. Em particular, \textbf{Ofra Haza} (MISC) , \textbf{Riccardo Fogli} (PER) , \textbf{Hoffmann \& Hoffmann} (ORG) , \textbf{Sweet Dreams} (ORG) , \textbf{Carola} (PER) , \textbf{Guy Bonnet} (PER) , \textbf{Danijel} (PER) e \textbf{Corinne Hermès} (PER) apareceram.

O procedimento foi interrompido por novos erros de linguagem de \textbf{Marlène Charell} (PER) .



\textbf{Festival} (PER)A ordem de votação no \textbf{Festival Eurovisão da Canção 1983} (MISC), foi a seguinte:

Os países que receberam 12 pontos foram os seguintes:

Em baixo encontra-se a lista de maestros que conduziram a orquestra, na respectiva actuação de cada país concorrente.



Festival Eurovisão da Canção 1984



\textbf{Países Participantes Países} (LOC) que participaram no passado, mas não em 1984

O \textbf{Festival Eurovisão da Canção 1984} (MISC) (em inglês : \textbf{Eurovision Song Contest} (ORG) 1984 e em francês : \textbf{Concours Eurovision} (MISC) de la chanson 1984 ) foi o 29 Festival Eurovisão da Canção e realizou-se em 5 de \textbf{Maio de 1984} (MISC) em \textbf{Luxemburgo} (LOC) . A apresentadora foi \textbf{Désirée Nosbusch} (PER) . Nosbusch, com apenas 19 anos de idade na época, abriu o certame de forma descontraída, o que era bastante incomum na época. Manifestou a sua fluência em 5 línguas, quando ela mudou entre inglês , francês , alemão , italiano e luxemburguês no decorrer de seu discurso, muitas vezes na mesma frase. O trio sueco \textbf{Herreys} (PER) foi o vencedor deste ano, interpretando o tema " \textbf{Diggi-Loo Diggi-Ley} (MISC) ". Foi a primeira edição em que foi transmitida em estéreo em alguns países.



\textbf{Local} (MISC)O \textbf{Festival Eurovisão da Canção} (MISC) 1984 ocorreu na capital do \textbf{Luxemburgo} (LOC) . A cidade do \textbf{Luxemburgo} (LOC) é uma comuna de \textbf{Luxemburgo} (LOC) com status de cidade , pertencente ao distrito de \textbf{Luxemburgo} (LOC) e ao cantão de \textbf{Luxemburgo} (LOC) . \textbf{Luxemburgo} (LOC) é uma cidade muito desenvolvida em seu comércio e nas indústrias. \textbf{Luxemburgo} (LOC) é uma das cidades mais ricas da \textbf{Europa} (LOC) e tornou-se um importante centro financeiro e administrativo. A cidade do \textbf{Luxemburgo} (LOC) contém várias instituições da \textbf{União Europeia} (ORG) , incluindo o \textbf{Tribunal de Justiça Europeu} (ORG) , o \textbf{Tribunal de Contas} (ORG) e o \textbf{Banco Europeu de Investimento} (ORG) .

O festival em si realizou-se no \textbf{Grand Théâtre de Luxembourg} (LOC) , inaugurado em 1964 como \textbf{Théâtre Municipal de la Ville} (LOC) de \textbf{Luxembourg} (LOC) , é uma das maiores salas de espetáculos do país para teatro , ópera e ballet . Passou por obras de renovação em 2002-2003, resultando em melhorias substanciais nas instalações de tecnologia de palco, acústica e iluminação.



\textbf{Formato Inicialmente} (PER), existiram rumores de que o \textbf{Luxemburgo} (LOC) não poderia receber a edição de 1984 devido ao facto de não ter um local adequado para o realizar. No entanto, optou-se pelo \textbf{Grand Théâtre de Luxembourg} (LOC) , que já tinha recebido a edição de 1973 , mas o acesso era limitado: apenas foram permitidos entrar as delegações, a imprensa e alguns convidados VIP . De entre as personalidades que assistiram ao festival ao vivo no recinto estavam o futuro grão-duque \textbf{Henrique de Luxemburgo} (PER) e a sua esposa \textbf{Maria Teresa de Luxemburgo} (PER) .

\textbf{Grécia} (LOC) e \textbf{Israel} (LOC) desistiram do concurso, este último devido ao facto do certame ter tido lugar no mesmo dia do que o \textbf{Yom HaZikaron} (MISC) , o \textbf{Dia da Lembrança dos Soldados Mortos de Israel} (MISC).

O certame foi transmitido em mais de 30 países, entre os quais a \textbf{União Soviética} (ORG) , com alguns deles a transmitir pela primeira vez em estéreo. Esta edição foi a menor em cerca de 10 anos, com apenas 2 horas e 13 minutos de emissão.

O núcleo da equipa técnica foi formado por \textbf{René Steichen} (PER) (direção), \textbf{Ray van Kant} (PER) , \textbf{Mariette Dumont} (PER) , \textbf{Patrick Seyler} (PER) e \textbf{Arlette Goebel} (PER) (produção executiva), \textbf{Roland de Groot} (PER) (cenografia) e \textbf{Pierre Cao} (PER) (direção musical).



\textbf{Visual} (MISC)A introdução consistia na transmissão de imagens aéreas do país anfitrião, terminando com o maestro \textbf{Pierre Cao} (PER) a dirigir a orquestra com o instrumental das canções vencedoras pelo \textbf{Luxemburgo} (PER).

A orquestra, dirigida por \textbf{Pierre Cao} (PER) , localizava-se num fosso, em frente ao palco. O palco, da autoria de \textbf{Roland de Groot} (PER) (que já tinha sido o autor dos palcos das edições de 1970, 1976 e 1980), consistia em painéis móveis cujas diferentes montagens deram profundidade ao pequeno palco, com três pódios com escadas, bem como uma série de figuras geométricas planas, que a princípio formavam um "4", referindo-se ao ano em curso: essas figuras mudavam de posição de acordo com as atuações; todos esses elementos foram iluminados por dentro. O quadro de votação estava localizado à esquerda do palco. A mesa do supervisor estava localizada à direita. Pela primeira vez, a produção utilizou animações geradas por computador, especialmente para abertura e para os cartões postais.

A apresentadora foi \textbf{Désirée Nosbusch} (PER) , que conduziu o certame de forma descontraída, falando aos espectadores em inglês , francês , alemão , italiano e luxemburguês e, devido à sua grande fluência nestes idiomas, falava estas 5 línguas, muitas vezes na mesma frase. Curiosamente, o diretor musical, \textbf{Pierre Cao} (PER) , foi seu professor de música na escola.

Os cartões postais consistiram em imagens geradas por computador em que estereótipos foram intercalados com carimbos do país correspondente, sempre em tom humorístico. Não em vão, estes sketches foram realizados por um grupo de comediantes parisienses batizados para a ocasião com o nome de "\textbf{Turistas} (LOC)". Da mesma forma, e como aconteceu em \textbf{Harrogate} (LOC) em 1982, antes do início de cada apresentação, as câmeras focaram no comentarista da respectiva televisão, quando muitas delas aproveitaram a oportunidade para saudar os espectadores.

O intervalo foi ocupado pelos artistas da \textbf{Prague Theatre of Illuminated Drawings} (PER) , com um número de mímica que girava em torno do diálogo silencioso entre um palhaço e um cavalo, ambos manipulados por cordas.

Pela primeira vez na história do concurso, uma música e seus artistas foram vaiados pelo público. Tratpu-se da canção inglesa e o grupo \textbf{Belle \& The Devotions} (ORG) , composto por \textbf{Kit Rolfe} (PER) , \textbf{Laura James} (PER) e \textbf{Linda Sofield} (PER) . A causa deste incidente ainda é difícil de determinar. Duas hipóteses são geralmente avançadas: em primeiro lugar, uma reação dos espectadores de \textbf{Luxemburgo} (PER) aos eventos que ocorreram a 16 de novembro de 1983 , pois nesse dia, depois de uma partida de futebol entre as equipes inglesa e luxemburguesa, hooligans britânicos saquearam a cidade de \textbf{Luxemburgo} (LOC), causando danos generalizados e indignação do público luxemburguês; em segundo lugar, uma reação dos espectadores a um truque presumido da atuação, pois parecia que apenas \textbf{Rolfe} (LOC) estava a cantar.



\textbf{Votação} (MISC)Cada país tinha um júri composto por 11 elementos, que atribuiu 12, 10, 8, 7, 6, 5, 4, 3, 2, 1 pontos às dez canções mais votadas.

Durante a votação, a câmera fez vários close-ups dos artistas. Em particular, a banda \textbf{Bravo} (MISC) , \textbf{Herreys} (PER) e \textbf{Linda Martin} (PER) apareceram.

Na primeira parte do procedimento, o placar apresentava na tela um fundo preto. Mas depois do voto belga, o fundo ficou azul.



\textbf{Festival} (MISC)A ordem de votação no \textbf{Festival Eurovisão da Canção 1984} (MISC), foi a seguinte:

Os países que receberam 12 pontos foram os seguintes:

Em baixo encontra-se a lista de maestros que conduziram a orquestra, na respectiva actuação de cada país concorrente.



Festival Eurovisão da Canção 1985



\textbf{Países Participantes Países} (LOC) que participaram no passado, mas não em 1985

O \textbf{Festival Eurovisão da Canção 1985} (MISC) (em inglês : \textbf{Eurovision Song Contest} (ORG) 1985 , em francês : \textbf{Concours Eurovision} (MISC) de la chanson 1985 e em sueco : \textbf{Eurovisionens Melodi Festival 1985} (MISC) ) foi o 30 \textbf{Festival Eurovisão da Canção} (MISC) e realizou-se em 4 de maio de 1985 em \textbf{Gotemburgo} (LOC) . A apresentadora foi \textbf{Lill Lindfors} (PER) . O duo \textbf{Bobbysocks} (LOC) ( \textbf{Hanne Krogh} (PER) e \textbf{Elisabeth Andreassen} (PER) ) convenceu o júri europeu, arrecadando 123 pontos. Assim, após vários anos entre os últimos lugares, a \textbf{Noruega} (LOC) consegue finalmente atingir a vitória.



\textbf{Local} (MISC)O \textbf{Festival Eurovisão da Canção} (MISC) 1985 ocorreu em \textbf{Gotemburgo} (LOC) , na \textbf{Suécia} (LOC) . \textbf{Gotemburgo} (LOC) é uma cidade da \textbf{Suécia} (LOC) do \textbf{condado de Västra Götaland} (LOC) . Com cerca de 513 000 habitantes na cidade e 879 000 na área metropolitana, é a segunda maior cidade sueca depois de \textbf{Estocolmo} (LOC) . A cidade situa-se na foz do \textbf{rio Gota} (LOC) , que desagua no \textbf{estreito de Categate} (LOC) . O rio divide a cidade em duas partes, sendo a margem norte a ilha \textbf{Hisingen} (LOC) . O local é apropriado para a presença de um porto : a cidade expandiu o seu porto marítimo, tornando-se no maior dos países nórdicos .

O festival em si realizou-se no \textbf{Scandinavium} (LOC) , então o maior recinto da \textbf{Suécia} (LOC) . Com a capacidade para receber 12 000-14 000 pessoas, foi inaugurado em 1971, quando a cidade comemorou os 350 anos da sua existência. Está situado numa parte central de \textbf{Gotemburgo} (LOC), junto às \textbf{Torres Gothia Towers} (MISC) e ao \textbf{Centro de Exposições e} (LOC) \textbf{Congressos da Suécia} (MISC) , e na proximidade do estádio de futebol \textbf{Ullevi} (ORG) , do parque de diversões de \textbf{Liseberg} (LOC) , do pavilhão científico \textbf{Universeum} (LOC) e do \textbf{Museu da Cultura Mundial} (LOC) .



\textbf{Formato} (MISC)A Grécia e \textbf{Israel} (LOC) regressaram ao certame, enquanto que os \textbf{Países Baixos} (LOC) retiraram-se porque a data do concurso coincidiu com o \textbf{Dia da Memória Nacional} (MISC) em homenagem às vítimas da \textbf{Segunda Guerra Mundial} (MISC) e a \textbf{Jugoslávia} (ORG) desistiu também pelo dia coincidir com o dia do aniversário da morte do \textbf{marechal Tito} (PER) , apesar de ter escolhido a música "\textbf{Pokora} (MISC)" ( \textbf{Penitência} (LOC) ), composto por \textbf{Ivo Pupačić} (PER) e \textbf{Zvonimir Pupačić} (PER) e cantado por \textbf{Zorica Kondža} (PER) e \textbf{Josip Genda} (PER) .

A esta edição, regressaram, nada mais nada menos do que 13 artistas, entre eles as vencedoras ( \textbf{Hanne Krogh} (PER) e \textbf{Elisabeth Andreassen} (PER) ).

O \textbf{Festival} (LOC) foi transmitido pela primeira vez via satélite e, pela primeira vez, os bilhetes para assistir os ensaios foram colocados à venda.

O núcleo da equipa técnica foi formado por \textbf{Steen Priwin} (PER) (direção e produção executiva), \textbf{Ingemar Wiberg} (PER) , \textbf{Ebbe Börjesson} (PER) , \textbf{Jan Johansson} (PER) , \textbf{Dick Johansson} (PER) , \textbf{Göran Hallberg} (PER) , \textbf{Richard Farbotko} (PER) e \textbf{Sven Orsman} (MISC) (cenário) e \textbf{Curt-Eric Holmquist} (PER) (endereço musical).



\textbf{Visual} (MISC)A abertura da competição começou com um vídeo mostrando os pontos turísticos de \textbf{Gotemburgo} (LOC) e concluindo com um amplo plano sobre o \textbf{Scandinavium} (MISC) . A câmera então revelou o palco e a orquestra. Depois, apareceu a apresentadora, \textbf{Lill Lindfors} (PER) a cantar a música e jazz "\textbf{My Joy is Building Bricks of Music} (MISC)". A orquestra e a plateia aplaudiam ritmicamente durante a sua atuação. Esta foi a primeira vez na história da competição em que uma apresentadora cantou. .

\textbf{Lill Lindfors} (PER) então introduziu as apresentações habituais, dizendo: "A música não conhece fronteiras, e isso é especialmente verdadeiro hoje à noite", lembrando também que o concurso comemorava a sua trigésima edição, que tinha na sua plateia a primeira vencedora do certame, \textbf{Lys Assia} (PER) , tendo a orquestra tocado parte do refrão da sua música " \textbf{Refrain} (PER) ".

A orquestra, dirigida por \textbf{Curt-Eric Holmquist} (PER) , ficava à esquerda do palco, sob uma parede decorada com formas trapezoidais e uma reprodução gigantesca do logótipo. O quadro de votação e a mesa do supervisor foram colocados à direita do palco num pódio separado. A cena em si, de cor branca, consistia em várias bandejas, formas geométricas e diferentes alturas, entrelaçadas entre si. Eles estavam cobertos com faixas pretas e bordas brilhantes. Uma escadaria larga, com degraus luminosos, permitia o acesso aos bastidores. A decoração consistia em duas estruturas principais. À esquerda, uma escultura feita de três triângulos dourados e uma cerca cúbica. À direita, uma estrutura tubular que suporta padrões na forma de uma vela de barco e palmeiras.

A apresentadora foi \textbf{Lill Lindfors} (PER) , que falou aos espectadores quase exclusivamente em inglês , mas também em francês .

Os cartões postais foram focados nos autores e compositores. Um pequeno vídeo mostrou-os, passeando por \textbf{Gotemburgo} (LOC) , durante a semana de ensaio. \textbf{Lill Lindfors} (PER) então introduziu as canções, os artistas e os maestros, muitas vezes na língua nacional do país em questão. A cada vez, ela segurava na mão uma reprodução em miniatura da bandeira do país.

O intervalo foi ocupado pelos artistas do "\textbf{Guitars Unlimited} (MISC)", um duo de guitarra composto por \textbf{Peter Almqvist} (PER) e \textbf{Ulf Wakenius} (PER) , que tocaram algumas canções de jazz na guitarra.

É após o intervalo que aconteceu aquele que é considerado um dos momentos mais marcantes da \textbf{História do certame} (MISC). Após o intervalo, a apresentadora \textbf{Lill Lindsfors} (PER) chocou toda a \textbf{Europa} (LOC) quando a sua saia ficou presa num elemento do cenário, deixando assim a sua roupa interior à vista. Depois de alguns segundos de riso, entende-se que todo o acidente estava previamente ensaiado e a apresentadora rapidamente recompôs o seu vestido.



\textbf{Votação} (MISC)Cada país tinha um júri composto por 11 elementos, que atribuiu 12, 10, 8, 7, 6, 5, 4, 3, 2, 1 pontos às dez canções mais votadas.

Durante a votação, a câmera fez vários close-ups dos artistas. Em particular, a banda \textbf{Bobbysocks} (ORG) e \textbf{Wind} (LOC) apareceram.

Uma pequena novidade no quadro de voto foi o facto de à esquerda do nome dos diferentes países estava a classificação provisória. Então, após cada votação, era mais fácil ler a situação geral.



\textbf{Festival} (PER)A ordem de votação no \textbf{Festival Eurovisão da Canção 1985} (MISC), foi a seguinte:

Os países que receberam 12 pontos foram os seguintes:

Em baixo encontra-se a lista de maestros que conduziram a orquestra, na respectiva actuação de cada país concorrente.



Festival Eurovisão da Canção 1986



\textbf{Países Participantes Países} (LOC) que participaram no passado, mas não em 1986

O \textbf{Festival Eurovisão da Canção 1986} (MISC) (em inglês : \textbf{Eurovision Song Contest} (ORG) 1986 , em francês : \textbf{Concours Eurovision} (MISC) de la chanson 1986 e em norueguês : \textbf{Europeiske Melodi Grand Prix} (MISC) 1986 ) foi o 31 \textbf{Festival Eurovisão da Canção} (MISC) e realizou-se em 3 de maio de 1986 na cidade de \textbf{Bergen} (LOC) , na \textbf{Noruega} (LOC) . A apresentadora foi \textbf{Åse Kleveland} (PER) . A jovem belga \textbf{Sandra Kim} (PER) foi a vencedora desse ano, interpretando a canção " \textbf{J'aime la vie} (MISC) ". Com apenas 13 anos, \textbf{Sandra} (LOC) foi até ao momento a mais jovem vencedora da competição. As regras atuais exigem que os concorrentes têm de ter pelo menos 15 anos, portanto a não ser que as regras sejam alteradas, o seu recorde nunca será batido. A irreverente representante portuguesa, \textbf{Dora} (PER) , apresentou-se em palco com uma espalhafatosa saia verde-alface aos folhos e umas botas de tropa, o que lhe valeu enumeras fotos nas páginas de vários jornais por toda \textbf{Europa} (LOC).



\textbf{Local} (MISC)O \textbf{Festival Eurovisão da Canção 1986} (MISC) ocorreu em \textbf{Bergen} (LOC) , na \textbf{Noruega} (LOC) . \textbf{Bergen} (LOC) é a segunda maior cidade da \textbf{Noruega} (LOC) , com uma população de aproximadamente 250 mil habitantes. A cidade está cercada por sete montanhas, o que lhe confere uma bela paisagem, mas também o título de cidade mais chuvosa da \textbf{Europa} (LOC) . \textbf{Bergen} (LOC) é um centro de cultura, comércio e universitário na costa oeste da \textbf{Noruega} (LOC). O compositor \textbf{Edvard Grieg} (PER) nasceu e viveu na cidade e lá compôs várias de suas mundialmente famosas peças. É uma cidade muito popular quer para turistas norugueses, quer para turistas estrangeiros, sendo um do principais pontos de paragens dos cruzeiros dos mares do norte da \textbf{Europa} (LOC) . É o ponto de partida do barco \textbf{Hurtigruten} (MISC) , que viaja pelas costas norueguesas até ao \textbf{Kirkenes} (ORG) , sendo estação final da linha ferroviária de \textbf{Bergen} (LOC) ( \textbf{Bergensbanen} (LOC) ), que a liga a \textbf{Oslo} (LOC) , por paisagens de grande beleza natural. \textbf{Bergen} (LOC) teve a sua própria universidade ( \textbf{Universidade de Bergen} (LOC) ) em 1946 . Esta tem hoje cerca de 17 000 estudantes. A montanha amarela é maior das sete montanhas que rodeiam a cidade, com 987 metros de altura. Foi fundada em 1070 , por \textbf{Olav Kyrre} (PER) , tendo sido a capital da \textbf{Noruega} (LOC) até 1299 , quando \textbf{Oslo} (LOC) passou a ser a capital. Ainda assim, continuou a ser a maior cidade da \textbf{Noruega} (LOC) até à década de 1830 . Entre 1350 e 1750 pertenceu à \textbf{Liga Hanseática} (ORG) . A parte antiga da cidade ligada à baía, chamada \textbf{Bryggen} (LOC) , circundada por casas datando do tempo da \textbf{Liga Hanseática} (ORG) , passou a integrar a lista do património da humanidade da \textbf{Unesco} (ORG) , em 1979 .

O festival em si realizou-se no \textbf{Grieghallen} (LOC) , sala de concertos com capacidade de 1.500 lugares. O salão tem o nome do compositor norueguês \textbf{Edvard Grieg} (PER) nascido em \textbf{Bergen} (LOC), que foi diretor musical da \textbf{Orquestra Filarmónica de Bergen} (ORG) entre 1880 e 1882. O \textbf{Grieghallen} (LOC) também é famosa como o estúdio de gravação e da comunidade black metal, como muitas bandas importantes como \textbf{Burzum} (ORG), \textbf{Mayhem} (ORG) ou \textbf{Immortal} (ORG) gravando lá seus primeiros registros com discografia \textbf{Deathlike Silence Productions} (PER) e \textbf{Pytten} (PER) como técnico de som.



\textbf{Formato} (MISC)A \textbf{Islândia} (LOC) competiu pela primeira vez, com a emissora nacional \textbf{RÚV} (MISC) , enquanto que \textbf{a Grécia} (LOC) retirou-se, além de ser o 18 país a apresentar a sua canção. A razão para a retirada foi que a data desta edição da \textbf{Eurovisão} (LOC) coincidiu com o \textbf{Sábado Santo} (LOC). A sua canção seria "Wagon-lit" ( ) e seria interpretada por \textbf{Polina} (PER). A emissora italiana, \textbf{RAI} (ORG) , simplesmente decidiu não enviar uma delegação para \textbf{Bergen} (LOC) .

A organização do concurso foi uma fonte de orgulho nacional em toda a \textbf{Noruega} (LOC) . Portanto, várias cidades norueguesas apresentaram a sua candidatura para sediar o certame. No outono de 1985 , a \textbf{NRK} (ORG) decidiu escolheu \textbf{Bergen} (LOC) para organizar a edição, derrotando \textbf{Oslo} (LOC) , \textbf{Stavanger} (LOC) , \textbf{Sandnes} (LOC) e \textbf{Trondheim} (LOC) . A televisão pública norueguesa divulgou um orçamento muito alto, a fim de apresentar a \textbf{Noruega} (LOC) à melhor luz para o resto do mundo.

O núcleo da equipa técnica foi formado por \textbf{John Andreassen} (PER) (direção), \textbf{Per Selstrøm} (PER) e \textbf{Harald Tusberg} (PER) (produção executiva), \textbf{Per Fjeld} (PER) (cenário), \textbf{Asbjørn W. Hagen} (PER) , \textbf{Tom Sundli} (PER) e \textbf{Willy Bettvik} (PER) (iluminação) e \textbf{Egil-Monn Iversen} (PER) e \textbf{Terje Fjærn} (LOC) (direção musical).



\textbf{Visual} (MISC)A abertura da competição começou com um vídeo mostrando os fiordes noruegueses, com pinturas das mesmas paisagens, com a orquestra a tocar em direto " Na \textbf{Gruta do Rei da Montanha} (LOC) ", da peça \textbf{Peer Gynt} (MISC) , de \textbf{Grieg} (PER) . O vídeo terminou com as vistas de \textbf{Bergen} (LOC) e do \textbf{Grieghallen} (LOC) .

\textbf{Åse} (PER) \textbf{Kleveland} (LOC) entrou e cumprimentou os espectadores cantando "\textbf{Welcome to Music} (MISC)", em inglês , francês e alemão . A orquestra tocou um trecho de \textbf{Te Deum} (PER) de \textbf{Marc-Antoine Charpentier} (PER) . Depois, fez as apresentações habituais, começando por sublinhar que os \textbf{Jogos Olímpicos} (MISC) e o \textbf{Festival Eurovisão da Canção} (MISC) têm pelo menos uma coisa em comum: o importante não é ganhar, mas participar.

A orquestra, dirigida por \textbf{Egil Monn-Iversen} (PER) , ficava à esquerda do palco, num pódio separado. A mesa do supervisor foi instalada na plateia, ao pé de uma tela gigante, do mesmo lado da orquestra. O quadro de votação estava à direita do palco. Este último consistia numa grande plataforma elevada de forma irregular, delimitada por néons luminosos, criando um gradiente de azul a rosa. O acesso ao palco era feito por uma escada de nove degraus, forrada de luzes azuis e cor-de-rosa. Esta escadaria foi encimada por um arco móvel. O arco e outros elementos decorativos reproduziam cristais de rocha gigantes, translúcidos e iluminados por dentro. Toda a cena assumiu cores azul, rosa e laranja, dependendo das atuações. O cenário do palco foi inspirado num icebergue .

A apresentadora foi \textbf{Lill Lindfors} (PER) , que falou aos espectadores em inglês e francês . A primeira roupa que ela usava era um vestido incrustado de diamantes, feito especialmente para a ocasião e pesando nada menos que 6,8 kg.

Os cartões postais consistiam em vídeos turísticos da \textbf{Noruega} (LOC), terminando com a bandeira do país participante. Em seguida, apareceu um verdadeiro cartão postal, em que foi escrito na língua nacional do país: "Saudações para \textbf{Noruega} (LOC) de [nome do país]". Os nomes dos autores e compositores, assim como o título da música, foram escritos no mapa. Finalmente, os intérpretes apareceram na tela, com, no fundo, uma imagem turística. \textbf{Åse Kleveland} (PER) então apresentou os intérpretes e maestros, dando algumas explicações adicionais sobre eles e segurando na sua mão uma reprodução em miniatura da bandeira dos países participantes.

O intervalo foi ocupado pela apresentação ao público internacional da cantora de \textbf{Bergen Sissel Kyrkjebø} (PER) , que interpretou, junto com o músico \textbf{Steinar Ofsdal} (PER) e um coro infantil, um vídeo filmado na cidade durante toda a mesma semana e em o que interpretou uma seleção de canções entre a população local. O intervalo terminou com a jovem no palco tocando, rodeada pela orquestra, um fragmento de "Udsigter fra \textbf{Ulriken} (LOC)", o hino regional de \textbf{Bergen} (LOC) popularmente conhecido com nomes como "\textbf{Nystemten} (MISC)" e "\textbf{Bergensiana} (LOC)".

A família real norueguesa esteve presente no recinto, representador pelo príncipe \textbf{Haroldo} (PER) , princesa \textbf{Sónia} (LOC) , o príncipe \textbf{Haquino Magno} (PER) e a princesa \textbf{Marta Luísa} (PER) .

A edição foi marcada por três eventos: pela homenagem feita a \textbf{Ossi Runne} (ORG) , maestro finlandês, que recebeu uma ovação especial porque estava a comemorar a sua vigésima participação no concurso; pelo facto da atuação norueguesa ter um grupo de drag queens , vestidos com roupas estilizadas com reminiscências do século XVIII; e também pela vencedora. \textbf{Sandra Kim} (PER) que afirmava ter 15 anos para poder participar, tinha apenas 13 quando se sagrou vencedora, tornando-se assim na mais jovem vencedora do \textbf{Festival Eurovisão da Canção} (MISC). Após se ter descoberto esta mentira, a \textbf{Suíça} (LOC) (2 lugar) exigiu de imediato e, sem êxito, a desclassificação da \textbf{Bélgica} (LOC).



\textbf{Votação} (MISC)Cada país tinha um júri composto por 11 elementos, que atribuiu 12, 10, 8, 7, 6, 5, 4, 3, 2, 1 pontos às dez canções mais votadas.

Durante a votação, a câmera fez vários close-ups dos artistas. Em particular, \textbf{Sandra Kim} (PER) , \textbf{Daniela Simmons} (PER) e \textbf{Ingrid Peters} (PER) apareceram.

O porta-voz do júri holandês foi o único a fazer um comentário durante o processo. Depois de outorgar um ponto para a \textbf{Finlândia} (LOC) , dois para a \textbf{Noruega} (LOC) , três para a \textbf{Suécia} (LOC) , quatro para a \textbf{Dinamarca} (LOC) e cinco para a \textbf{Islândia} (LOC) , acrescentou: "Isso deve ser suficiente para a \textbf{Escandinávia} (ORG), eu acho".



\textbf{Festival} (MISC)A ordem de votação no \textbf{Festival Eurovisão da Canção 1986} (MISC), foi a seguinte:

Os países que receberam 12 pontos foram os seguintes:

Em baixo encontra-se a lista de maestros que conduziram a orquestra, na respectiva actuação de cada país concorrente.



Festival Eurovisão da Canção 1987



\textbf{Países Participantes Países} (LOC) que participaram no passado, mas não em 1987

O \textbf{Festival Eurovisão da Canção} (MISC) 1987 (em inglês : \textbf{Eurovision Song Contest} (ORG) 1987 , em francês : \textbf{Concours Eurovision} (MISC) de la chanson 1987 e em holandês : Eurovisiesongfestival 1987 ) foi o 32 \textbf{Festival Eurovisão da Canção} (MISC) e realizou-se em 9 de maio de 1987 em \textbf{Bruxelas} (LOC) . Pela premira vez, o \textbf{Festival} (LOC) coincidiu com o \textbf{Dia da Europa} (MISC) , justamente no ano em que se comemoravam os 30 anos da assinatura dos \textbf{Tratados de Roma} (ORG) , embrião da \textbf{Comunidade Económica Europeia} (ORG) . A apresentadora foi \textbf{Viktor Lazlo} (PER) . Foi com uma típica balada irlandesa que \textbf{Johnny Logan} (PER) se sagrou vencedor do \textbf{Festival Eurovisão da Canção} (MISC) de 1987. \textbf{Johnny Logan} (PER) já tinha participado e vencido em 1980, tornando-se assim no primeiro concorrente a ganhar duas vezes este certame.



\textbf{Local} (MISC)O \textbf{Festival Eurovisão da Canção} (MISC) 1987 ocorreu em \textbf{Bruxelas} (LOC) , na \textbf{Bélgica} (LOC) . \textbf{Bruxelas} (LOC) é a capital de facto da \textbf{União Europeia} (ORG) (\textbf{UE} (ORG)) e a maior área urbana na \textbf{Bélgica} (LOC) . É composta por 19 comunas , incluindo a \textbf{Cidade de Bruxelas} (LOC) , propriamente dita, que é a capital da \textbf{Bélgica} (LOC) , \textbf{Flandres} (LOC) e da \textbf{Comunidade francesa} (ORG) da \textbf{Bélgica} (LOC) . \textbf{Bruxelas} (LOC) cresceu de uma fortaleza no século X , fundada por um descendente de \textbf{Carlos Magno} (PER) , para uma metrópole de mais de um milhão de habitantes. A área metropolitana da cidade tem uma população de mais de 1,8 milhões de habitantes, tornando-a maior da \textbf{Bélgica} (LOC) . Desde o final da \textbf{Segunda Guerra Mundial} (MISC) ( 1939 - 1945 ), \textbf{Bruxelas} (LOC) foi um importante centro de política internacional . A presença das principais instituições da \textbf{União Europeia} (ORG) , bem como a sede da \textbf{Organização do Tratado do Atlântico Norte} (ORG) (\textbf{Otan} (ORG)) fez da cidade uma casa poliglota de muitas organizações internacionais, políticos, diplomatas e de funcionários públicos. Embora historicamente seja uma região de falantes do neerlandês , \textbf{Bruxelas} (LOC) tornou-se uma cidade com cada vez mais falantes da língua francesa ao longo dos séculos XIX e XX . Hoje a maioria dos habitantes são nativos falantes do francês , embora ambas as línguas tenham estatuto oficial. Tensões linguísticas continuam e as leis de linguagem dos municípios em torno de \textbf{Bruxelas} (LOC) são um tema de muita controvérsia na \textbf{Bélgica} (LOC).

O festival em si realizou-se no \textbf{Palais du Centenaire} (MISC) , área no norte de \textbf{Bruxelas} (LOC) , conhecida por ter sido o local da \textbf{Exposição Universal de 1935} (MISC) e da \textbf{Expo} (MISC) 58 .



\textbf{Formato} (PER)Esta foi a maior edição até então: 22 países participantes, contando com o regresso da \textbf{Itália} (LOC) e da \textbf{Grécia} (LOC). De todos os países participantes até então no certame, apenas \textbf{Malta} (LOC), \textbf{Marrocos} (LOC) e \textbf{Mónaco} (LOC) não participaram na edição de 1987.

Inicialmente, a \textbf{RTBF} (ORG) e a \textbf{VRT} (ORG) iriam organizar a edição em conjunto. No entanto, devido a desentendimentos entre as duas emissoras, a \textbf{VRT} (ORG) decidiu não co-organizar, deixando a \textbf{RTBF} (ORG) sozinha na organização. O orçamento necessário era tão importante que uma nova lei teve que ser adotada, permitindo o uso de publicidade para financiar os canais públicos belgas. Foi a primeira vez que patrocinadores ajudaram a fazer o concurso e apareceram na tela. Apesar de a \textbf{RTBF} (ORG) ser o canal televisivo organizador do festival, não teve a responsabilidade de escolher o seu representante, cabendo à \textbf{VRT} (ORG) essa escolha. Este facto ocorreu devido à partilha da participação no Festival Eurovisão da Canção, por parte dos dois canais belgas que intercalam anualmente as suas participações. Assim, a \textbf{RTBF} (ORG) foi o único canal a organizar o festival sem ter o seu representante nessa mesma edição.

A canção israelita, " \textbf{Shir} (LOC) habatlanim ", causou polémica devido ao seu conteúdo humorístico, tendo a atuação sido criticada pelo \textbf{Ministro da Cultura} (MISC).

O título original da canção sueca era "\textbf{Fyra Bugg Och en Coca-Cola} (MISC)" (Quatro pastilhas e um coca-cola), no entanto a \textbf{EBU} (MISC) impôs uma mudança no tema devido à publicidade à famosa marca de refrigerantes " \textbf{Coca-Cola} (PER) ". Assim o título foi alterado para "\textbf{Boogaloo} (LOC)" e onde se ouvia "\textbf{Fyrabugg} (MISC) och en coca-cola" passou a ouvir-se "\textbf{Boogaloo} (LOC), dansa rock'n rolla".

O núcleo da equipa técnica foi formado por \textbf{Jacques Bourton} (PER) (endereço), \textbf{Michel Gehu} (PER) (produtor executivo), \textbf{Daniel Georges} (PER) (cenário) e \textbf{Jo Carlier} (PER) (direção musical).



\textbf{Visual} (MISC)A abertura da competição começou com uma animação por computador. O logótipo do concurso apareceu no espaço estrelado, logo juntou-se o foguete lunar projetado por \textbf{Hergé} (PER) em " Objetivo \textbf{Lua} (LOC) " e o avião de " O \textbf{Caranguejo das Tenazes de Ouro} (ORG) ". O avião então mergulhou na atmosfera e nuvens, lançando um vídeo turístico da \textbf{Bélgica} (LOC) . Isso terminou com uma vista do \textbf{Palais du Centenaire} (MISC) .

A câmera então revelou a orquestra e o palco. \textbf{Viktor Lazlo} (PER) fez sua entrada e desceu as escadas enquanto cantava "\textbf{Breathless} (LOC)". Ela então fez as apresentações, lembrando que, aquele dia, se tratava do aniversário da \textbf{Declaração Schuman} (MISC) , e do 30 aniversário da \textbf{União Europeia} (ORG) .

A orquestra, dirigida por \textbf{Jo Carlier} (PER) , ficava à esquerda do palco, num espaço reservado. O palco construído pela \textbf{RTBF} (ORG) , inspirado em formas geométricas tridimensionais, respondeu a uma estética vanguardista e rompeu com a cenografia dos anos anteriores. A tecnologia de iluminação incorporou como principal novidade o uso de lasers, que, apesar de terem sido utilizados em tempo útil nalguma edição anterior, adquiriram especial destaque nesta ocasião.

A apresentadora foi \textbf{Viktor Lazlo} (PER) , que falou aos espectadores em inglês e francês .

Os cartões postais consistiam em vídeos turísticos da \textbf{Bélgica} (LOC) , terminando com a bandeira do país participante. Em seguida, a bandeira foi substituída por um quadradinho de banda desenhada , na qual estava escrito na língua nacional do país: "Para [nome do país], com amizade". Finalmente, os intérpretes se moveram para dentro do quadradinho , antes de saudar os telespectadores. \textbf{Viktor Lazlo} (PER) então apresentou os intérpretes e maestros, dando algumas explicações adicionais sobre eles e segurando na sua mão uma reprodução em miniatura da bandeira dos países participantes.

O intervalo foi por um vídeo dedicado às conquistas da \textbf{União Europeia} (ORG) . As estrelas da bandeira europeia podiam ser vistas viajando através do continente, através de suas paisagens e locais culturais, enquanto que, em paralelo, bailarinos executavam balé contemporâneo. O vídeo começou e terminou com uma cena do músico \textbf{Marc Grauwels} (PER) , que tocou variações no estilo da flauta sobre o tema musical da \textbf{Nona Sinfonia de Beethoven} (MISC) .

A família real belga esteve presente no recinto, neste caso os príncipes herdeiros príncipe \textbf{Alberto} (PER) e a \textbf{princesa Paula} (PER) .

Os cartões postais desta edição consistiu numa forma de arte que é património cultural da \textbf{Bélgica} (LOC) : a banda desenhada .



\textbf{Votação} (MISC)Cada país tinha um júri composto por 11 elementos, que atribuiu 12, 10, 8, 7, 6, 5, 4, 3, 2, 1 pontos às dez canções mais votadas.

Durante a votação, a câmera fez vários close-ups dos artistas. Em particular, \textbf{Johnny Logan} (PER) e a banda \textbf{Wind} (MISC) apareceram.

Esta foi a última edição em que apareceu um quadro de votação físico.



\textbf{Festival} (PER)A ordem de votação no \textbf{Festival Eurovisão da Canção} (MISC) 1987, foi a seguinte:

Os países que receberam 12 pontos foram os seguintes:

Em baixo encontra-se a lista de maestros que conduziram a orquestra, na respectiva actuação de cada país concorrente.



Festival Eurovisão da Canção 1988



\textbf{Países Participantes Países} (LOC) que participaram no passado, mas não em 1988

O \textbf{Festival Eurovisão da Canção 1988} (MISC) (em inglês : \textbf{Eurovision Song Contest} (ORG) 1988 , em francês : \textbf{Concours Eurovision} (MISC) de la chanson 1988 e em irlandês : \textbf{Comórtas Amhránaíochta} (LOC) na \textbf{hEoraifíse} (MISC) 1988 ) foi o 33. \textbf{Festival Eurovisão da Canção} (MISC) e realizou-se em 30 de abril de 1988 , em \textbf{Dublin} (LOC) , na \textbf{Irlanda} (LOC) . Os apresentadores foram \textbf{Pat Kenny} (PER) e \textbf{Michelle Rocca} (PER) .

A futura super-estrela canadiana \textbf{Céline Dion} (PER) , apenas famosa no mundo expressão francesa, foi a vencedora do \textbf{Festival} (PER) desse ano, com a canção " \textbf{Ne} (PER) partez pas sans moi ", representando a televisão suíça.



\textbf{Local} (MISC)O \textbf{Festival Eurovisão da Canção} (MISC) 1988 ocorreu em \textbf{Dublin} (LOC) , na \textbf{Irlanda} (LOC) . \textbf{Dublin} (LOC) é a capital e maior cidade da \textbf{Irlanda} (LOC) . O nome em inglês deriva da palavra irlandesa " \textbf{Dubhlinn} (LOC) " (ocasionalmente também grafada \textbf{Duibhlinn} (PER) ou \textbf{Dubh Linn} (PER) ), que significa "\textbf{Lago Negro} (LOC)". Localiza-se na província de \textbf{Leinster} (LOC) próxima ao ponto mediano da costa leste da \textbf{Irlanda} (LOC), sendo cortada pelo \textbf{Rio Liffey} (LOC) e o centro da região de \textbf{Dublin} (LOC). Desde 1898 possui nível administrativo de condado ( county-boroughs ). Seus limites são os condados de \textbf{Fingal} (LOC) a norte, \textbf{Dublin} (LOC) meridional a sudoeste e \textbf{Dun Laoghaire-Rathdown} (LOC) a sudeste. Tem uma população de 527.612 habitantes na cidade, e sua área metropolitana tem 1.804.156 habitantes. Fundada como um assentamento viquingue , foi o centro do \textbf{Reino de Dublin} (LOC) e se tornou a principal cidade da \textbf{Ilha} (LOC) após a invasão dos \textbf{Normandos} (LOC) . A cidade cresceu de maneira rápida durante o século XVII ; se tornou na época a segunda maior cidade do \textbf{Império Britânico} (ORG) e a quinta maior da \textbf{Europa} (LOC) . \textbf{Dublin} (LOC) entrou em um período de estagnação após o \textbf{Ato de União de 1800} (MISC) , mas continuou o centro económico da \textbf{Ilha} (LOC). Após a \textbf{Partição da Irlanda} (MISC) em 1922, virou a capital do \textbf{Estado Livre Irlandês} (MISC) , e mais tarde, da \textbf{República da Irlanda} (LOC) . \textbf{Dublin} (LOC) é reconhecida como uma cidade global , com um ranking "\textbf{Alpha-} (MISC)", colocando a cidade entre as 30 mais globalizadas do mundo. Atualmente é o principal centro histórico, cultural, económico, industrial e educacional da \textbf{Irlanda} (LOC).

O festival em si realizou-se no \textbf{RDS Simmonscourt Pavilion} (MISC) , uma sala polivalente, pertencente ao complexo de recreação da \textbf{Royal Dublin Society} (ORG) e destina-se a sediar festivais equestres e agrícolas. O pavilhão pode acomodar até 7000 pessoas.



\textbf{Formato} (MISC)A representante da \textbf{Suíça} (LOC), \textbf{Céline Dion} (PER) , tinha nacionalidade canadiana e era, até à data, uma simples cantora desconhecida para a grande parte dos países europeus. Após a sua vitória e devido ao pouco sucesso que a sua canção obteve na \textbf{Europa} (LOC), ninguém previa que \textbf{Céline Dion} (PER) se iria tornar numa das grandes divas internacionais da música, tal como é considerada actualmente. Também nesta edição, mas pelo \textbf{Luxemburgo} (LOC) , \textbf{Lara Fabian} (PER) concorreu, tendo ficado em 4 lugar.

O \textbf{Chipre} (LOC) , que era para ser o segundo país a subir ao palco, já tinha escolhido os seus representantes, seriam \textbf{Yiannis Dimitrou} (PER) e \textbf{Scott Adams} (PER) , com a canção de " \textbf{Thimame} (LOC) " (\textbf{Recordar} (LOC)), mas não puderam representar o seu país pois já tinham lançado o seu tema anos antes e isso é uma violação das regras do concurso.

Esta edição foi a primeira desde 1979 que teve dois apresentadores.

O logótipo do show, foi desenhado por \textbf{Maria Quigley} (PER). A \textbf{Eurovisão V} (MISC) representa o símbolo da harpa celta da \textbf{Irlanda} (LOC). O conjunto de \textbf{Michael Grogan} (PER) , que já havia sido criado em 1981, consistia em uma cena na forma de uma perspectiva brilhante. No fundo um céu estrelado.

O núcleo da equipa técnica foi formado por \textbf{Declan Lowney} (PER) (direção), \textbf{Liam Miller} (PER) (produtor executivo), \textbf{Paula Farrell} (PER) e \textbf{Michael Grogan} (PER) (cenário) e \textbf{Noel Kelehan} (PER) e \textbf{Bill Whelan} (PER) (direção musical).



\textbf{Visual} (MISC)A abertura da competição começou com um vídeo, resumindo a história da \textbf{Irlanda} (LOC) , desde a pré-história até os dias atuais, passando por pontos turísticos e evidências arqueológicas. O vídeo termina com um close-up de uma lua cheia, que é então exibida na parede da tela. A câmera então revelou o palco, com \textbf{Johnny Logan} (PER) , o vencedor da edição de 1987, a cantar " Hold Me \textbf{Now} (MISC) ", a música que lhe valeu a vitória em \textbf{Bruxelas} (LOC) no ano anterior. Depois, os apresentadores \textbf{Michelle Rocca} (PER) e \textbf{Pat Kenny} (PER) fizeram as apresentações habituais, observando que \textbf{Dublin} (LOC) comemorou o seu milénio, e que também estava a ser comemorado o \textbf{Ano Europeu do Cinema e Televisão} (MISC). \textbf{Michelle Rocca} (PER) cumprimentou os participantes e a abertura termina com um segundo vídeo mostrando a recepção dos participantes em \textbf{Dublin} (LOC) na semana de ensaio. A música que acompanhava era a música "\textbf{Heartland} (PER)", interpretada por \textbf{Johnny Logan} (PER) .

A orquestra, dirigida por \textbf{Noel Kelehan} (PER) , estava à esquerda do palco. A \textbf{RTÉ} (MISC) , seguindo o exemplo do seu antecessor, investiu nos avanços tecnológicos da época para oferecer um grande programa de televisão. Para isso, a emissora depositou a gestão do programa nas mãos de \textbf{Declan Lowney} (PER) , um jovem produtor de apenas 28 anos de idade, com experiência em programas para jovens e concertos de grande escala. O palco, que no princípio dava a impressão de estar suspenso, já foi chamado de "espetacular". Era uma plataforma dividida em grades, cujas linhas eram tubos de néon que mudavam de cor dependendo do tema em competição e que, graças em parte a uma conclusão cuidadosa, conferida à televisão, tanto ao local quanto ao próprio palco, mais profundidade e dimensões maiores do que realmente tinham. Ele também foi flanqueado em ambos os lados por duas grandes telas de vídeo que projetavam dinamicamente imagens do \textbf{logótipo} (LOC) e do artista no palco, pela primeira vez na história.

Os apresentadores foram \textbf{Michelle Rocca} (PER) e \textbf{Pat Kenny} (PER) , que falou aos espectadores em inglês e francês .

Os cartões postais consistiam nos participantes a fazer coisas na \textbf{Irlanda} (LOC) , de cultura a tradição, desporto ou turismo.

O intervalo foi preenchido pela banda de rock irlandesa \textbf{Hothouse Flowers} (PER) com a música "\textbf{Do Not Go} (MISC)", sobre imagens filmadas em diferentes cidades da \textbf{Europa} (LOC) e foi o videoclipe mais caro já produzido na \textbf{Irlanda} (LOC) na época. A produção deste clipe foi subsidiada pela \textbf{Comissão Europeia} (ORG) .



\textbf{Votação} (MISC)Cada país tinha um júri composto por 16 elementos, que atribuiu 12, 10, 8, 7, 6, 5, 4, 3, 2, 1 pontos às dez canções mais votadas.

Durante a votação, a câmera fez vários close-ups dos artistas. Em particular, \textbf{Lara Fabian} (PER) , \textbf{Scott Fitzgerald} (PER) e \textbf{Celine Dion} (PER) apareceram.

Pela primeira vez, um quadro de votação digital foi usado. Também pela primeira vez, um ranking dos cinco principais países foi mostrado na tela.

Durante o processo de votação, o \textbf{Reino Unido} (LOC) manteve-se no 1 lugar com alguma vantagem sobre a \textbf{Suíça} (LOC). No entanto, o júri da \textbf{Jugoslávia} (ORG) (o último a dar os votos) entregou 6 pontos à \textbf{Suíça} (LOC) e para surpresa de todos não pontuou o \textbf{Reino Unido} (LOC), dando assim a vitória à \textbf{Suíça} (LOC) por apenas 1 ponto de diferença.



\textbf{Festival} (PER)A ordem de votação no \textbf{Festival Eurovisão da Canção 1988} (MISC), foi a seguinte:

Os países que receberam 12 pontos foram os seguintes:

Em baixo encontra-se a lista de maestros que conduziram a orquestra, na respectiva actuação de cada país concorrente.



Festival Eurovisão da Canção 1989



Erro: é necessário especificar uma imagem na primeira linha

\textbf{Países Participantes Países} (LOC) que participaram no passado, mas não em 1989

O \textbf{Festival Eurovisão da Canção 1989} (MISC) (em inglês : \textbf{Eurovision Song Contest 1989} (ORG) e em francês : \textbf{Concours Eurovision} (MISC) de la chanson 1989 ) foi o 34. \textbf{Festival Eurovisão da Canção} (MISC) e realizou-se em 6 de maio de 1989 , em \textbf{Lausana} (LOC) , na \textbf{Suíça} (LOC) . Os apresentadores foram \textbf{Lolita Morena} (PER) e \textbf{Jacques Deschenaux} (PER) .

A banda jugoslava \textbf{Riva} (PER) foi a vencedora desse ano, com a canção " Rock Me ", canção tipo pop-rock e que nada de novo trouxe à competição. Esta vitória foi vista por muitos como sendo política, visto que a \textbf{Jugoslávia} (ORG) era vista nessa época (o \textbf{muro de Berlim} (MISC) só caiu em Novembro desse ano) como o país mais "democrático" dos países de leste.



\textbf{Local} (MISC)O \textbf{Festival Eurovisão da Canção 1989} (MISC) ocorreu em \textbf{Lausana} (LOC) , na \textbf{Suíça} (LOC) . \textbf{Lausana} (LOC) é uma cidade na \textbf{Suíça} (LOC) romanda , a parte francófona da \textbf{Suíça} (LOC) , e é a capital do cantão de \textbf{Vaud} (LOC) . Sede do distrito de \textbf{Lausana} (LOC) , a cidade está situada às margens do \textbf{Lago Léman} (LOC) (em francês : \textbf{Lac Léman} (PER) ). Limitada pela cidade francesa de Évian-les-Bains ao sul do lago, com as montanhas \textbf{Jura} (LOC) a noroeste. \textbf{Lausana} (LOC) está localizada 62 km a nordeste de \textbf{Genebra} (LOC) . \textbf{Lausana} (LOC) tinha, em dezembro de 2011, uma população de 129 383, tornando-se a quarta maior cidade do país, com a área de aglomeração com 336 400 habitantes. A \textbf{Região Metropolitana de Lausana-Genebra} (LOC) possui cerca de 1,2 milhão de habitantes. A sede do \textbf{Comitê Olímpico Internacional} (ORG) (\textbf{COI} (ORG)) e o \textbf{Museu Olímpico} (LOC) estão localizados na cidade,além de inúmeras federações e organismos esportivos estão sediados na cidade - o \textbf{COI} (ORG) reconhece oficialmente a cidade como a \textbf{Capital Olímpica} (ORG) - assim como a sede do \textbf{Tribunal Arbitral do Esporte} (LOC) . Encontra-se no meio de uma região do vinho . A cidade tem um sistema de metrô com 28 estações, tornando-se a menor cidade do mundo a ter um sistema de transporte rápido .

O festival em si realizou-se no \textbf{Palais} (MISC) de \textbf{Beaulieu} (LOC) , centro de convenções onde, em outubro de 1986 , foi decidido a cidade-sede dos \textbf{Jogos Olímpicos de Verão de 1992} (MISC) .



\textbf{Formato Nesta} (PER) edição competiram 22 países, igualando o recorde de 1987 , com o regresso do \textbf{Chipre} (LOC).

Esta edição marcou a despedida de \textbf{Ossi Runne} (PER) , o famoso maestro finlandês que, após 23 anos, se despedia do palco eurovisivo e pela primeira vez, um país usou dois maestros para conduzirem a orquestra durante a sua atuação, pois \textbf{Henrik Krogsgård} (PER) dirigiu a orquestra apenas metade da música, antes de se juntar a \textbf{Birthe Kjær} (PER) no palco e \textbf{Benoît Kaufman} (PER) , que era o maestro do país anfitrião, substitui \textbf{Krogsgård} (LOC) e dirigiu a orquestra até ao final da música. Foi a única vez na história em que uma canção teve dois maestros.

O logo da competição representava o \textbf{Matterhorn} (ORG) . A torre da \textbf{Catedral de Lausana} (LOC) apareceu na linha sublinhando o título.

\textbf{Ricardo Landum} (PER) , membro dos \textbf{Da Vinci} (PER) , viu-se impossibilitado de representar \textbf{Portugal} (LOC), devido a um acidente sofrido pouco antes do \textbf{Festival Eurovisão da Canção} (MISC). Sendo assim, substituído por \textbf{Tó} (LOC).

De destacar que dois intérpretes deste festival, \textbf{Nathalie Pâque} (PER) e \textbf{Gili Natanael} (PER) tinham respectivamente apenas 11 e 12 anos. Devido aos protestos que surgiram, a \textbf{União Europeia de Radiodifusão} (ORG) introduziu uma regra segundo a qual apenas podem participar cantores com mais de 16 anos.

Como curiosidade, vale sublinhar que o autor das músicas da \textbf{Alemanha} (LOC) e da \textbf{Áustria} (LOC) tratava-se de \textbf{Dieter Bohlen} (PER) , membro dos \textbf{Modern Talking} (ORG) .

Foi adotado um novo método de desempate, consistindo em proclamar vencedor o país que tiver recebido maior número de vezes 12 ponto e assim sucessivamente.

O núcleo da equipa técnica foi formado por \textbf{Alain Bloch} (PER) (realização e produção executiva), \textbf{Charles-André Grivet} (PER) (realização), \textbf{Raymond Zumsteg} (PER) (produção executiva), \textbf{Georges Mani} (PER) e \textbf{Paul Waelti} (PER) (cenografia) e \textbf{Benoît Kaufman} (PER) (direção musical).



\textbf{Visual} (MISC)A abertura da competição começou com um vídeo, com uma jovem rapariga. Na sua companhia, a câmera viajou por paisagens e cidades suíças e desvendou certos aspectos da vida social, cultural e desportiva do país. O vídeo termina com uma a rapariga a descer as escadas do \textbf{Palais} (MISC) de \textbf{Beaulieu} (LOC) . Uma limusine preta parou, cuja jovem abriu a porta dos. Então saiu \textbf{Céline Dion} (PER) , a vencedora do ano anterior. Juntas, subiram as escadas e entraram no recinto. \textbf{Céline Dion} (PER) , segurando a rapariga pela mão, entrou no palco e cantou uma parte de " Ne partez pas sans moi ", música vencedora do ano anterior. Ela se juntou aos apresentadores que agradeceram e congratularam-na pela sua vitória. Depois, \textbf{Céline} (PER) cantou seu primeiro single solo em inglês, " \textbf{Where Does My Heart Beat Now} (MISC) ". Essa atuação chamou a atenção dos líderes da \textbf{Disney Studios} (ORG) e permitiu a \textbf{Dion} (PER) prosseguir a sua carreira em países de língua inglesa. "\textbf{Where Does My Heart Beat Now} (MISC)" foi lançado no ano seguinte nos \textbf{Estados Unidos} (LOC) e no início de 1991 no resto do mundo. O sucesso seguirá com a banda sonora de \textbf{A Bela e o Monstro} (MISC) , que em 1992 receberia o \textbf{Oscar} (MISC) de melhor banda sonora original e entraria no top 10 no ranking de música do \textbf{Canadá} (LOC) , \textbf{Estados Unidos} (LOC) , \textbf{Nova Zelândia} (LOC) e \textbf{Reino Unido} (LOC) . Depois, os apresentadores enviaram mais duas saudações especiais a \textbf{Jean-Pascal Delamuraz} (PER) , \textbf{Presidente da Confederação Helvética} (MISC) , que estava presente na audiência, e a \textbf{Benoît Kaufman} (ORG) , o maestro. A abertura termina com uma peça executada ao vivo pela orquestra.

A orquestra, dirigida por \textbf{Benoît Kaufman} (PER) , estava à esquerda do palco. O palco consistia em três pódios separados, de formato triangular e vermelho escuro. O primeiro, à esquerda, foi reservado para os apresentadores. O segundo, no centro, acolheu os artistas. Foi fornecido em ambos os lados com duas entradas na forma de pirâmides truncadas, de aparência metálica e iluminadas por neons azuis. Na parte de trás, um elemento de plexiglass, na forma de um custódio, também era iluminado, iluminado com luzes de néon azuis. O fundo consistia de duas formas triangulares, uma das quais estava coberta com faixas claras, assim como tubos de néon azul, rosa e verde, que se entrelaçavam. Um elemento móvel, na forma de asas, também poderia ser implantado atrás do quarto traseiro. Finalmente, o terceiro pódio, à direita, foi para a orquestra. Duas telas completaram o palco.

Os apresentadores foram \textbf{Jacques Deschenaux} (PER) e \textbf{Lolita Morena} (PER) , que falou aos espectadores em francês , inglês , alemão , italiano , como também algumas palavras em castelhano e \textbf{Romanche} (PER) .

Os cartões postais começaram com uma nota dos apresentadores, explicando o título e o tema da música, e então apresentando os artistas e o maestro, tendo, nas canções da \textbf{Itália} (LOC) , \textbf{Alemanha} (LOC) , \textbf{Áustria} (LOC) , \textbf{Espanha} (LOC) e parte da de \textbf{Israel} (LOC) , para além das dos países anglófonos e francófonos, apresentado nas línguas oficiais. Em seguida, seguiu um vídeo, mostrando os participantes na descoberta de pontos turísticos da \textbf{Suíça} (LOC) .

O intervalo foi preenchido pelo artista de circo \textbf{Guy Tell} (PER) , que fez um número baseado na lenda de \textbf{William Tell} (PER) . Usando uma besta , \textbf{Guy Tell} (PER) alcançou um alvo mascarado, cortou uma rosa, depois um jornal e cortou os fios de quatro balões. Então ele estourou cinco outros balões de uma só vez. O final foi particularmente especiais, pois \textbf{Guy Tell} (PER) fez disparar 16 bestas sucessivamente. O último ladrilho deveria ter acertado numa maçã. Mas, ao executar o número ao vivo, o bloco perdeu a maçã. A produção, no entanto, antecipou esse evento e registou uma tentativa bem-sucedida. Este último foi transmitido mais tarde, logo após a tentativa fracassada.



\textbf{Votação} (MISC)Cada país tinha um júri composto por 11 elementos, que atribuiu 12, 10, 8, 7, 6, 5, 4, 3, 2, 1 pontos às dez canções mais votadas.

O supervisor executivo da \textbf{EBU} (MISC) foi \textbf{Frank Naef} (PER) , que estava a comemorar os seus 30 anos na \textbf{União Europeia de Radiodifusão} (ORG) .

Durante a votação, a câmera fez vários close-ups dos artistas. Em particular, \textbf{Thomas Forstner} (PER) , a banda \textbf{Riva} (ORG) e \textbf{Tommy Nilsson} (PER) apareceram.



\textbf{Festival} (MISC)A ordem de votação no \textbf{Festival Eurovisão da Canção 1989} (MISC), foi a seguinte:

Os países que receberam 12 pontos foram os seguintes:

Em baixo encontra-se a lista de maestros que conduziram a orquestra, na respectiva actuação de cada país concorrente.



Artistas repetentes

Alguns artistas repetiram a sua experiência \textbf{Eurovisiva} (PER). Em 1989, os repetentes foram:



\textbf{Festival Eurovisão da Canção 1990} (MISC)



\textbf{Erro} (PER): é necessário especificar uma imagem na primeira linha

\textbf{Países Participantes Países} (LOC) que participaram no passado, mas não em 1990

O \textbf{Festival Eurovisão da Canção 1990} (MISC) (em inglês : \textbf{Eurovision Song Contest 1990} (ORG) , em francês : \textbf{Concours Eurovision} (MISC) de la chanson 1990 e em \textbf{Servo-croata} (LOC) : Pjesma Eurovizije 1990 ) foi o 35. \textbf{Festival Eurovisão da Canção} (MISC) e realizou-se em 5 de maio de 1990 , em \textbf{Zagreb} (LOC) , na \textbf{Jugoslávia} (ORG) , atualmente na \textbf{Croácia} (LOC) . Os apresentadores do festival foram \textbf{Helga Vlahović} (PER) e \textbf{Oliver Mlakar} (PER) . \textbf{Toto Cutugno} (PER) foi o vencedor do festival deste ano com a canção " \textbf{Insieme} (ORG): 1992 " que preconizava a unidade europeia . \textbf{Cutugno} (LOC) tinha 46 anos e 302 dias no momento de sua vitória, tornando-se o mais velho vencedor do concurso até então, recorde esse ultrapassado 10 anos depois.

As letras de várias canções celebravam os ventos de mudança e democratização dos países do leste da \textbf{Europa} (LOC), focando em especial o momento da queda do \textbf{Muro de Berlim} (MISC) ocorrida em 9 de novembro de 1989. Todavia, a letra da canção vencedora era uma antecipação do mercado único europeu na \textbf{União Europeia} (ORG) que teria lugar no final de 1992.



\textbf{Local} (MISC)O \textbf{Festival Eurovisão da Canção 1990} (MISC) ocorreu em \textbf{Zagreb} (LOC) , na \textbf{Jugoslávia} (ORG) , atualmente na \textbf{Croácia} (LOC) . \textbf{Zagreb} (LOC) é a capital e maior cidade da \textbf{Croácia} (LOC) . Localizada entre a margem do \textbf{rio Sava} (LOC) e a encosta do \textbf{monte Medvednica} (LOC), situa-se 120 metros acima do nível do mar , nas coordenadas 4548'N 15\textbf{58'E.} (LOC) Uma encruzilhada entre a \textbf{Europa Central} (LOC) e o \textbf{Mar Adriático} (LOC) , \textbf{Zagreb} (LOC) concentra indústria, instituições científicas, órgãos administrativos nacionais e ministérios.

O festival em si realizou-se no \textbf{Koncertna} (MISC) dvorana \textbf{Vatroslav Lisinski} (PER) , é uma grande sala de concertos e um centro de convenções, que tinha acabado de ter a sua primeira remodelação em 1989. Em 1992, a cobertura do telhado de cobre do salão foi completamente substituída. O seu nome é devido a \textbf{Vatroslav Lisinski} (PER) , um compositor croata do século XIX . O edifício tem duas salas: uma grande com 1841 lugares e uma pequena com 305 lugares. Um grande salão também serve como uma grande área de exposição.



\textbf{Formato Nesta} (PER) edição do Festival Eurovisão da Canção foram notórias as diversas mensagens políticas e os apelos à paz presentes nas letras dos temas apresentados. Este facto deveu-se aos acontecimentos ocorridos por toda a \textbf{Europa} (LOC) , com especial destaque para a queda do \textbf{muro de Berlim} (MISC) em 9 de novembro de 1989.

Pela primeira vez na história do \textbf{Festival Eurovisão da Canção} (MISC) foi introduzida uma mascote: o \textbf{Eurocat} (MISC), criado por \textbf{Joško Marušić} (PER) . Este desenho animado apareceu diversas vezes ao longo do certame e a sua principal função foi animar os separadores entre as actuações dos países concorrentes, tendo em vista o \textbf{Ano Europeu do Turismo} (MISC).

Durante os ensaios, os apresentadores sentiram-se ofendidos por alguns comentários dos jornalistas acerca da sua idade avançada e decidiram renegar o cargo de apresentadores. Foram entretanto substituídos temporariamente por \textbf{Rene Medvešek} (PER) e \textbf{Dubravka Marković} (PER) . No entanto, desfizeram-se todos os desentendidos e assim \textbf{Helga Vlahović} (PER) e \textbf{Oliver Mlakar} (PER) voltaram a tomar as suas posições na apresentação do certame.

A actuação espanhola, foi marcada por uma enorme falha no playback instrumental que acidentalmente arrancou com alguns segundos de avanço, causando alguma perplexidade nas intérpretes, que por ordem da sua delegação, abandonaram o palco. Alguns segundos depois e após uma salva de palmas, o playback instrumental recomeçou, desta vez, de início e toda a actuação decorreu sem mais incidentes.

\textbf{Malta} (LOC) desejou voltar ao certame pela primeira vez em 15 anos, chegando a escolher \textbf{Maryrose Mallia} (PER) , com a canção " \textbf{Our Little World of Yesterday} (MISC) ", com o objetivo de concorrer. No entanto, devido à regra que obrigava que o número máximo de países concorrentes seriam 22, \textbf{Malta} (LOC) não pode concorrer.

Após a controvérsia da edição de 1989, a \textbf{EBU} (MISC) introduziu uma nova regra. A partir de agora, todos os candidatos teriam que ter 16 anos no dia da sua participação no concurso.



\textbf{Visual} (MISC)A abertura da competição começou com um vídeo sobre música em \textbf{Zagreb} (LOC) e na \textbf{Croácia} (LOC) , apresentando a riqueza musical cultural da capital croata, do clássico ao pop , passando pelo tradicional e folclórico , concluindo com o \textbf{Te Deum} (MISC), de \textbf{Marc-Antoine Charpentier} (PER) , por vários cantores e músicos diferentes. Após o vídeo, \textbf{Eurocat} (PER) entrou no \textbf{Koncertna} (LOC) dvorana \textbf{Vatroslav Lisinski} (PER) e depois os apresentadores saudaram os espetadores. A orquestra executou brevemente a partitura de " Rock Me ", a música vencedora do ano anterior.

A orquestra, dirigida por \textbf{Igor Kuljerić} (PER) , estava à esquerda do palco. O cenário no palco era caracterizado por um luminoso piso quadriculado que mudava de cor, flanqueado por duas grandes telas de vídeo e, ao fundo, uma estrutura de luzes de néon também na forma de um tabuleiro de xadrez.

Os apresentadores foram \textbf{Helga Vlahović} (PER) e \textbf{Oliver Mlakar} (PER) , que falaram aos espectadores em inglês e francês .

Os cartões postais apresentavam o \textbf{Eurocat} (MISC) a viajar pelos países concorrentes, em comemoração do \textbf{Ano Europeu do Turismo} (MISC). Começou com um breve esboço do \textbf{Eurocat} (MISC), sobre um fundo azul, com um lembrete do nome do país participante. No final de cada cartão postal, o \textbf{Eurocat} (MISC) tirou uma foto em preto e branco aos artistas.

O intervalo foi preenchido por um vídeo intitulado "\textbf{Yugoslav Changes} (PER)", dirigido por \textbf{Ivo Laurenčić} (PER) , um filme apresentando a riqueza turística e a diversidade cultural da \textbf{Jugoslávia} (ORG) , aparecendo \textbf{Split} (LOC) , \textbf{Piran} (LOC) , \textbf{Dubrovnik} (LOC) , \textbf{Sarajevo} (LOC) , \textbf{Belgrado} (LOC) e \textbf{Zagreb} (LOC) .



\textbf{Votação} (MISC)Cada país tinha um júri composto por 11 elementos, que atribuiu 12, 10, 8, 7, 6, 5, 4, 3, 2, 1 pontos às dez canções mais votadas.

Durante a votação, a câmera fez vários close-ups dos artistas. Em particular, \textbf{Toto Cutugno} (PER) e \textbf{Liam Reilly} (PER) apareceram.



\textbf{Festival} (MISC)A ordem de votação no \textbf{Festival Eurovisão da Canção 1990} (MISC), foi a seguinte:

Os países que receberam 12 pontos foram os seguintes:

Em baixo encontra-se a lista de maestros que conduziram a orquestra, na respectiva actuação de cada país concorrente.



Artistas repetentes

Alguns artistas repetiram a sua experiência \textbf{Eurovisiva} (PER). Em 1990, os repetentes foram:



Festival Eurovisão da Canção 1991



\textbf{Países Participantes Países} (LOC) que participaram no passado, mas não em 1991

O \textbf{Festival Eurovisão da Canção 1991} (MISC) (em inglês : \textbf{Eurovision Song Contest} (ORG) 1991 , em francês : \textbf{Concours Eurovision} (MISC) de la chanson 1991 e em italiana : \textbf{Concorso Eurovisione della Canzone} (PER) 1991 ) foi o 36. Festival Eurovisão da Canção e realizou-se a 4 de maio de 1991 em \textbf{Roma} (LOC) . Os apresentadores do festival foram os dois antigos vencedores italianos do \textbf{Festival Eurovisão da Canção} (MISC) , \textbf{Gigliola Cinquetti} (PER) (1964) e \textbf{Toto Cutugno} (LOC) (1990). \textbf{Carola} (PER) a cantora representante da \textbf{Suécia} (LOC) foi a vencedora do de 1991, com a canção \textbf{Fångad} (PER) av en stormvind ( \textbf{Capturada} (LOC) por uma tempestade de vento ).



\textbf{Local} (MISC)O \textbf{Festival Eurovisão da Canção} (MISC) 1991 ocorreu em \textbf{Roma} (LOC) , em \textbf{Itália} (LOC) . \textbf{Roma} (LOC) é uma cidade e uma comuna especial (chamada " \textbf{Roma Capitale} (LOC) ") da \textbf{Itália} (LOC) . É a capital do país, da província homônima e também da região do \textbf{Lácio} (LOC) . Com 2,7 milhões de habitantes em 1.285,3 quilômetros quadrados de área, também é a maior cidade italiana e a quarta cidade mais populosa da \textbf{União Europeia} (ORG) . A área urbana de \textbf{Roma} (LOC) se estende além dos limites administrativos da cidade com uma população de cerca de 3,8 milhões. Entre 3,2 e 4,2 milhões de pessoas vivem na área metropolitana da capital italiana. A cidade está localizada na porção centro-ocidental da \textbf{península itálica} (LOC) , cortada pelo \textbf{rio Tibre} (LOC) , dentro do \textbf{Lácio} (LOC). \textbf{Roma} (LOC) é a única cidade no mundo que tem em seu interior um país inteiro, o enclave do \textbf{Vaticano} (LOC) . A história de \textbf{Roma} (LOC) abrange mais de dois mil e quinhentos anos , desde a sua fundação lendária em \textbf{753 a.} (MISC)C. Roma é uma das mais antigas cidades continuamente ocupadas na \textbf{Europa} (LOC) e é conhecida como "A \textbf{Cidade Eterna} (LOC)", uma ideia expressa por poetas escritores da \textbf{Roma Antiga} (LOC) . No mundo antigo , foi sucessivamente a capital do \textbf{Reino de Roma} (LOC) , da \textbf{República Romana} (LOC) e do \textbf{Império Romano} (LOC) e é considerada um dos berços da civilização ocidental . Desde o século I, a cidade é a sede do papado e no século VIII a cidade tornou-se a capital dos \textbf{Estados Pontifícios} (LOC) , que duraram até 1870. Em 1871, \textbf{Roma} (LOC) se tornou a capital do \textbf{Reino da Itália} (LOC) e em 1946 da \textbf{República Italiana} (LOC) . Após a \textbf{Idade Média} (MISC) , \textbf{Roma} (LOC) foi governada pelos papas \textbf{Alexandre VI} (PER) e \textbf{Leão X} (PER) , que transformaram a cidade em um dos principais centros do \textbf{Renascimento} (LOC) italiano , juntamente com \textbf{Florença} (LOC) . A versão atual da \textbf{Basílica de São Pedro} (LOC) foi construída e a \textbf{Capela Sistina} (LOC) foi pintada por \textbf{Michelangelo} (PER) . Artistas famosos e arquitetos, como \textbf{Bramante} (PER) , \textbf{Bernini} (PER) e \textbf{Rafael} (PER) , residiu por algum tempo em \textbf{Roma} (LOC), contribuindo para a sua arquitetura renascentista e barroca. \textbf{Roma} (LOC) é considerada uma cidade global . Em 2007, \textbf{Roma} (LOC) foi a 11. cidade mais visitada do mundo, a terceira mais visitada da \textbf{União Europeia} (ORG) e a atração turística mais popular na \textbf{Itália} (LOC). A cidade tem uma das melhores "marcas" da \textbf{Europa} (LOC), tanto em reputação quanto em patrimônio. O seu centro histórico é classificado pela \textbf{UNESCO} (ORG) como \textbf{Patrimônio Mundial} (ORG) . Monumentos e museus tais como os \textbf{Museus Vaticanos} (LOC) e o \textbf{Coliseu} (LOC) estão entre os destinos turísticos mais visitados do mundo, sendo que ambos os locais recebem milhões de turistas por ano. \textbf{Roma} (LOC) também sediou os \textbf{Jogos Olímpicos de Verão de 1960} (MISC) .

O festival em si realizou-se no \textbf{Studio 15 di Cinecittà} (MISC) , um complexo de teatros e estúdios situados na periferia oriental de \textbf{Roma} (LOC) (cerca de 9 quilômetros de distância) responsável pela maior parte da produção cinematográfica italiana: aí vários filmes são rodados e espetáculos televisivos são gravados. Depois da \textbf{Segunda Guerra Mundial} (MISC) a produção retomou lentamente seu ritmo, mas foi nos anos 1950 que \textbf{Cinecittà} (LOC) estabeleceu-se com um dos estúdios cinematográficos mais importantes do mundo, com as películas estadunidenses \textbf{Quo Vadis de Mervyn LeRoy} (PER) (1951) e \textbf{Ben-Hur} (MISC) de \textbf{William Wyler} (PER) (1959). Este boom teve origem na competitividade económica dos estúdios romanos, que receberam o título informal de " \textbf{Hollywood} (LOC) no \textbf{Tibre} (LOC) ".



\textbf{Formato} (PER)Originalmente, a edição era para ser realizada no \textbf{Teatro Ariston} (LOC) , em \textbf{Sanremo} (LOC) . No entanto, devido à \textbf{Guerra do Golfo} (MISC) e ao aumento da tensão que se fazia sentir na \textbf{Jugoslávia} (ORG) , a organização optou por transferir o festival para \textbf{Roma} (LOC) . Como consequência desta mudança a \textbf{RAI} (ORG) demonstrou não estar totalmente preparada para organizar este festival. Relativamente à apresentação do festival, além de ser totalmente feita em italiano e de terem sido escolhidos poucos dias antes dos ensaios, foram frequentes os erros nos nomes das canções, dos artistas e dos directores de orquestras. E, durante votação, \textbf{Gigliola Cinquetti} (PER) e \textbf{Toto Cotugno} (PER) demonstraram não estar familiarizados com o sistema, sendo por diversas vezes necessária a intervenção do supervisor da \textbf{EBU} (MISC) . Durante a actuação sueca ocorreu um problema de som dentro do \textbf{Studio 15 di Cinecittà} (MISC) que impossibilitou os presentes na sala de escutarem a canção. No entanto a emissão televisiva decorreu normalmente sem qualquer tipo de problema sonoro.

Nesta edição, \textbf{Malta} (LOC) regressou ao certame após 16 anos de interregno e só pode participar pois os \textbf{Países Baixos} (LOC) não participaram devido ao dia dedicado à memória dos que morreram durante a \textbf{Segunda Guerra Mundial} (MISC) .

Esta foi também a última edição que contou com a participação da \textbf{República Socialista Federativa da Jugoslávia} (LOC) . Devido à desintegração que se deu nos meses e anos seguintes o país desapareceu, tendo todos os seus sucessores, execpto o \textbf{Kosovo} (LOC) , participado no certame anos depois.

Os cartões postais desta edição consistiram na homenagem a uma canção italiana por parte de cada artista.



\textbf{Votação} (MISC)Cada país tinha um júri composto por 11 elementos, que atribuiu 12, 10, 8, 7, 6, 5, 4, 3, 2, 1 pontos às dez canções mais votadas.

Durante a votação, a câmera fez vários close-ups dos artistas. Em particular, \textbf{Carola Häggkvist} (PER) , \textbf{Amina Annabi} (PER) e o \textbf{Duo Datz} (PER) apareceram.

Esta votação foi marcada por vários momentos, de entre os quais as cerca de 20 correções que o supervisor executivo da \textbf{EBU} (MISC) teve de fazer, ao empate entre duas canções em 1. lugar e ao facto de \textbf{Amina Annabi} (PER) ter abraçado o \textbf{Duo Datz} (PER) quando \textbf{Israel} (LOC) outorgou 12 pontos a \textbf{França} (LOC) .

Neste festival ocorreu um inusitado empate no final. Mas, ao contrário do que aconteceu em 1969, onde 4 candidatas foram declaradas campeãs, em 1970 as regras foram alteradas para que somente pudesse existir um vencedor. E assim, pela primeira e única vez até os dias atuais, houve a necessidade de se usar o critério desempate, que seria o número de notas 12 recebidas. No caso desta edição também houve um empate no número de notas 12 para cada país, sendo 4 para ambos. Nesta situação foram usados as notas 10. A \textbf{Suécia} (LOC) tinha recebido 5 enquanto a \textbf{França} (LOC) recebeu apenas 2. Desta forma \textbf{Carola} (PER), representando a \textbf{Suécia} (LOC), venceu o \textbf{Festival Eurovisão da Canção 1991} (MISC). Se a regra em vigor atualmente fosse usada, a \textbf{França} (LOC) seria a vencedora desta edição por ter sido votada por 18 países contra 17 da \textbf{Suécia} (LOC).

A tabela seguinte apresenta, de forma mais resumida o resultado de acordo com a regra de desempate aplicadas ao longo dos anos.



\textbf{Festival} (MISC)A abertura da competição começou com um clipe de vídeo: "\textbf{Celebration} (LOC)", cantado por \textbf{Sarah Carlson} (PER) . A cantora apareceu lá, passeado pelas ruínas da \textbf{Roma antiga} (LOC) com suas dançarinas. O clipe também incluiu trechos de \textbf{Ben-Hur} (MISC) . Os apresentadores então subiram ao palco e saudaram os espectadores e o público. Um pequeno vídeo mostrou a vitória de \textbf{Toto Cutugno} (PER) em \textbf{Zagreb} (LOC) no ano anterior. O cantor cantou parte da sua canção, " \textbf{Insieme} (ORG): 1992 ". Então ele convidou \textbf{Gigliola Cinquetti} (PER) para fazer o mesmo com " \textbf{Non ho l'età} (MISC) " e acompanhou-a ao piano. Foi a primeira vez que dois ex-vencedores juntos apresentaram o concurso. Depois de outro vídeo apresentando \textbf{Sanremo} (LOC) , \textbf{Cinquetti} (LOC) e \textbf{Cutugno} (LOC) concluíram as suas introduções.

A orquestra, dirigida por \textbf{Bruno Canfora} (PER) , ficava ao fundo à direita do palco. A orquestra atraiu críticas de delegações estrangeiras. Primeiro, durante os ensaios, pois os músicos chegavam muito tarde, alegando uma interrupção do transporte público, causada pelas chuvas torrenciais que caíam sobre \textbf{Roma} (LOC) . Então, durante a transmissão, os músicos executaram muitas notas falsas. O exemplo mais notável foi o solo perdido de um dos saxofonistas, durante a canção grega.

O cenário de \textbf{Lucciano Riccieri} (PER) estava em segundo plano atrás da orquestra. cidade e no lado das enormes estátuas de um estilo antigo. Na verdade, foi a recuperação de um set de filmagem.

Os apresentadores foram \textbf{Gigliola Cinquetti} (PER) e \textbf{Toto Cutugno} (LOC) , que falaram aos espectadores quase exclusivamente em italiano , falando em inglês e francês muito esporadicamente.

Os cartões postais consistiam em breves notas por parte dos apresentadores. Depois, os artistas então apareceram num fundo de céu azul. Todos prestaram uma homenagem a uma música italiana à sua escolha. Simultaneamente, os seus nomes foram exibidos, os do seu país, bem como a sua bandeira nacional. \textbf{Cinquetti} (LOC) e \textbf{Cutugno} (LOC) então apresentaram a música e o maestro, cometendo muitos erros na pronúncia de nomes e títulos.

O intervalo foi ocupado por pequeno vídeo com o \textbf{Cinecittà} (LOC), no qual apareceram algumas estrelas de cinema, incluindo \textbf{Silvana Mangano} (PER) , \textbf{Humphrey Bogart} (PER) , \textbf{Ava Gardner} (PER) e \textbf{Gerard Depardieu} (PER) . Uma atuação do artista \textbf{Arturo Brachetti} (PER) então se seguiu. Ele apareceu sucessivamente como uma diva, na elegante \textbf{Belle Époque} (PER), como um menino, como um ator de kabuki, e como um cavalheiro de smoking preto e depois branco. No final da sua atuação, \textbf{Brachetti} (PER) cumprimentou os espectadores em inglês, francês, castelhano e italiano.

A ordem de votação no \textbf{Festival Eurovisão da Canção 1991} (MISC), foi a seguinte: [ a ]

Os países que receberam 12 pontos foram os seguintes:

Em baixo encontra-se a lista de maestros que conduziram a orquestra, na respectiva actuação de cada país concorrente.



Artistas repetentes

Alguns artistas repetiram a sua experiência \textbf{Eurovisiva} (PER). Em 1991, os repetentes foram:



Festival Eurovisão da Canção 1992



\textbf{Países Participantes Países} (LOC) que participaram no passado, mas não em 1992

O \textbf{Festival Eurovisão da Canção 1992} (MISC) (em inglês : \textbf{Eurovision Song Contest 1992} (ORG) , em francês : \textbf{Concours Eurovision} (MISC) de la chanson 1992 e em sueco : \textbf{Eurovisionens Melodi Festival 1992} (MISC) ) foi o 37 \textbf{Festival Eurovisão da Canção} (MISC) e realizou-se a 9 de maio de 1992 , na cidade de \textbf{Malmö} (LOC) , na \textbf{Suécia} (LOC) . Os apresentadores do evento foram \textbf{Harald Treutiger} (PER) e \textbf{Lydia Cappolicchio} (MISC) . \textbf{Linda Martin} (PER) foi a vencedora deste festival com a canção " \textbf{Why Me?} (MISC) " ( \textbf{Porquê Eu ?} (PER)). A canção foi escrita por \textbf{Johnny Logan} (PER) , vencedor do \textbf{Festival Eurovisão da Canção 1980} (MISC) e do \textbf{Festival Eurovisão da Canção} (MISC) 1987 , enquanto cantor.



\textbf{Local} (MISC)O \textbf{Festival Eurovisão da Canção 1992} (MISC) ocorreu em \textbf{Malmö} (LOC) , na \textbf{Suécia} (LOC) . \textbf{Malmö} (LOC) é uma cidade da \textbf{Suécia} (LOC) , capital da comuna homônima e capital da província de \textbf{Escânia} (LOC) . Em 2017 contava com uma população de 312 012 habitantes, sendo deste modo a terceira cidade mais populosa do país. O estreito de \textbf{Öre} (LOC) (\textbf{Öresund} (LOC)) separa esta cidade da cidade de \textbf{Copenhaga} (LOC) , na \textbf{Dinamarca} (LOC) . Em termos económicos, \textbf{Malmö} (LOC) é um porto exportador de lacticínios , produtos agrícolas , químicos e florestais e importador de combustíveis e maquinaria . A cidade possui indústrias de construção naval , cerveja , açúcar , maquinaria, têxteis , fosfatos , cimento , etc. Foi uma das primeiras e mais industrializadas cidades da \textbf{Suécia} (LOC). \textbf{Malmö} (LOC) foi importante porto da \textbf{Liga Hanseática} (ORG) . Os principais monumentos a visitar nesta cidade são a fortaleza de \textbf{Malmö} (LOC), o \textbf{edifício da Câmara Municipal} (LOC) e a igreja gótica de \textbf{São Pedro} (PER), construída no século XIV .

O festival em si realizou-se no \textbf{Malmö Isstadion} (LOC) , é um recinto desportivo, com uma capacidade para 4800 pessoas e foi construída em 1970 .



\textbf{Formato} (MISC)Os \textbf{Países Baixos} (LOC) regressaram ao certame e, para permitir que \textbf{Malta} (LOC) também participasse, aumentaram o número de vagas para 23, o maior número de países concorrentes até então.

A antiga \textbf{República Socialista Federativa da Jugoslávia} (LOC) deixou de existir como país, em 1992 , dando origem a quatro novas nações. No entanto, a nova \textbf{República Federativa da Jugoslávia} (LOC) foi responsável pela participação do país nesta edição do Festival Eurovisão da Canção. Este território passou a ser oficialmente designado, em 2003, como \textbf{Sérvia e Montenegro} (LOC) .

Tal como em 1990 , esta edição contou também com uma mascote: o \textbf{Eurobird} (MISC). Este serviu para animar a apresentação dos temas a concurso.

\textbf{Geraldine Olivier} (MISC) sagrou-se vencedora do \textbf{Concours Eurovision 1992} (MISC) (final nacional \textbf{suíça} (LOC)) com o tema "\textbf{Soleil, Soleil} (MISC)". No entanto, como este tema não respeitava o regulamento acabou por ser substituído por " \textbf{Mister Music Man} (MISC) " interpretado por \textbf{Daisy Auvray} (PER) que tinha alcançado o segundo lugar neste certame.



\textbf{Visual} (MISC)A abertura da competição começou com uma animação em vídeo. A câmera saiu de \textbf{Roma} (LOC) e do \textbf{Coliseu} (LOC) , circulou a \textbf{Torre de Pisa} (LOC) , sobrevoou os \textbf{Alpes} (LOC) , foi para \textbf{Paris} (LOC) e a \textbf{Torre Eiffel} (LOC) e depois contornou a costa europeia para \textbf{Malmö} (LOC) . Um vídeo com vistas turísticas da \textbf{Suécia} (LOC) e \textbf{Malmö} (LOC) seguiu-se. A câmera então revelou o palco com os dançarinos a dançarem " \textbf{Fångad} (PER) av in stormvind ", a música vencedora do ano anterior. Os apresentadores cumprimentaram os espetadores e aludiram às mudanças políticas contemporâneas e ao surgimento de novos países europeus. Os apresentadores lembraram ainda que aquele dia, 9 de maio , é o dia da \textbf{Europa} (LOC). Esta sequência terminou com dois vídeos curtos: o primeiro mostrando uma vitória sueca em uma competição de hóquei no gelo e o segundo com a vitória de \textbf{Carola} (PER) em \textbf{Roma} (LOC) no concurso anterior. \textbf{Carola} (PER) então subiu ao palco e cantou "\textbf{All the Reasons to Live} (MISC)". No final de sua atuação, ela foi acompanhada por vinte e três dançarinos, cada um carregando a bandeira de um dos países participantes.

O cenário, inspirado na história de \textbf{Malmö} (LOC) , foi preenchido pela proa de um barco viking com uma decoração iluminada de uma ponte como pano de fundo.

Os apresentadores foram \textbf{Lydia Cappolicchio} (MISC) e \textbf{Harald Treutiger} (PER) , que falaram aos espectadores em sueca , inglês e francês .

Os cartões postais desdobraram-se sobre as páginas de um álbum de lembranças. Eurobird aparecia na tela e virava as páginas com uma varinha mágica. Cada página incluía a bandeira nacional do país participante, o título da música e os nomes de seus autores e compositores. A bandeira então se expandiu, lançando um vídeo turístico sobre o país. Os cartões postais terminaram com um mapa final na página e a saída do \textbf{Eurobird} (ORG). Os apresentadores finalmente cumprimentaram o maestro.

O intervalo foi ocupado por um balé contemporâneo, intitulado "\textbf{A Century Of Dance} (MISC)" e realizado por \textbf{David Johnson Dancers} (PER). Este ballet revisitou a história da dança do século XX , em ordem cronológica de sua aparência: o \textbf{hambo} (ORG) , o \textbf{charleston} (LOC) , o rock 'n' roll , o twist , o disco , o slow e o hip-hop



\textbf{Votação} (MISC)Cada país tinha um júri composto por 11 elementos, que atribuiu 12, 10, 8, 7, 6, 5, 4, 3, 2, 1 pontos às dez canções mais votadas.

O supervisor executivo da \textbf{EBU} (MISC) foi \textbf{Frank Naef} (PER) . Os apresentadores lembraram que ele estava oficiando há quinze anos e, em seguida, anunciou que ele estava se aposentando naquele ano. \textbf{Frank Naef} (PER) agradeceu e confirmou que a \textbf{Eurovisão} (LOC) permaneceria, mas teria de se adaptar para acolher novos países que desejassem participar. Ele então felicitou calorosamente a televisão pública sueca pela sua eficiência. Concluiu agradecendo aos dois colaboradores mais próximos, que o ajudaram desde 1977: \textbf{Marie-Claire Vionnet} (PER) e \textbf{Brian Fraser} (PER) . Ele recebeu um buquê de flores de \textbf{Carola} (PER).

Durante a votação, a câmera fez vários close-ups dos artistas. Em particular, \textbf{Linda Martin} (PER) e o \textbf{Michael Ball} (PER) apareceram.



\textbf{Festival} (MISC)A ordem de votação no \textbf{Festival Eurovisão da Canção 1992} (MISC), foi a seguinte:

Os países que receberam 12 pontos foram os seguintes:

Em baixo encontra-se a lista de maestros que conduziram a orquestra, na respectiva actuação de cada país concorrente.



Artistas repetentes

Alguns artistas repetiram a sua experiência \textbf{Eurovisiva} (PER). Em 1992, os repetentes foram:



Festival Eurovisão da Canção 1993



Países que classificaram-se para a final \textbf{Países} (MISC) que não se classificaram para a final \textbf{Países} (MISC) que participaram no passado, mas não em 1993

O \textbf{Festival Eurovisão da Canção 1993} (MISC) (em inglês : \textbf{Eurovision Song Contest} (ORG) 1993 , em francês : \textbf{Concours Eurovision} (MISC) de la chanson 1993 e em irlandês : \textbf{Comórtas Amhránaíochta} (LOC) na \textbf{hEoraifíse} (MISC) 1993 ) foi o 38 Festival Eurovisão da Canção e realizou-se em 15 de maio na cidade de \textbf{Millstreet} (LOC) , na \textbf{Irlanda} (LOC) . A apresentadora foi \textbf{Fionnuala Sweeney} (PER) e a vencedora do evento foi \textbf{Niamh Kavanagh} (PER) com a canção, " \textbf{In Your Eyes} (MISC) " ( Nos teus olhos ).

Em 1993 a \textbf{União Europeia de Radiodifusão} (ORG) teve uma explosão do número de países que desejavam participar, devido não só ao desaparecimento da \textbf{URSS} (ORG) , do fim da \textbf{Checoslováquia} (LOC) e da desintegração da \textbf{Jugoslávia} (ORG) , este o único país comunista que tradicionalmente participou, tendo ganho o \textbf{Festival Eurovisão da Canção 1989} (MISC) , realizando-se em \textbf{Zagreb} (LOC) o \textbf{Festival Eurovisão da Canção 1990} (MISC) . Pela primeira vez teve de haver uma pré-qualificação apenas para os países que nunca tinham participado. Nessa pré-eliminatória participaram a \textbf{Croácia} (LOC) , a \textbf{Bósnia e Herzegovina} (LOC) , \textbf{Hungria} (LOC) , \textbf{Eslovénia} (LOC) , \textbf{Eslováquia} (LOC) , \textbf{Roménia} (LOC) e \textbf{Estónia} (LOC) . A pré-eliminatória realizou-se em \textbf{Liubliana} (LOC) , capital da \textbf{Eslovénia} (LOC) , no dia 3 de abril de 1993 . tendo sido apurados para a final os três primeiros classificados: \textbf{Bósnia e Herzegovina} (LOC) , \textbf{Croácia} (LOC) e \textbf{Eslovénia} (LOC) .



\textbf{Local} (MISC)O \textbf{Festival Eurovisão da Canção 1992} (MISC) ocorreu em \textbf{Millstreet} (LOC) , na \textbf{Irlanda} (LOC) . \textbf{Millstreet} (LOC) é uma cidade irlandesa localizada no condado de \textbf{Cork} (LOC) , província de \textbf{Munster} (LOC) . A cidade fica a 57 km de \textbf{Cork} (LOC) e está ligada à capital do condado através da linha ferroviária \textbf{Mallow-Killarney-Tralee} (ORG). Esta foi a menor cidade de sempre a receber o certame.

O festival em si realizou-se no \textbf{Green Glens Arena} (MISC) , uma arena equestre na pequena cidade. A localização escolhida para acolher este \textbf{Festival Eurovisão da Canção} (MISC) foi bastante peculiar. Além de \textbf{Millstreet} (LOC) ser uma pequena e remota cidade, com cerca de cerca de 1500 habitantes, o local escolhido para a realização deste evento foi também invulgar uma vez que a \textbf{Green Glens Arena} (ORG) era um centro equestre.



\textbf{Formato} (PER)O dono da \textbf{Green Glens Arena} (ORG) , \textbf{Noel C. Duggan} (PER), escreveu à \textbf{RTÉ} (LOC) na mesma noite da vitória irlandesa na edição de 1992, propondo o uso gratuito do local para sediar o concurso. O local, um grande centro equestre coberto e bem equipado, foi considerado mais do que adequado como a localização pela emissora RTÉ . Com enorme apoio das autoridades locais e nacionais, além de várias empresas na região, a infraestrutura da cidade foi bastante aprimorada para acomodar um evento dessa magnitude. Foi também a maior transmissão externa já tentada pela emissora estatal \textbf{RTÉ} (MISC) e foi considerada um triunfo técnico para todos os envolvidos. No entanto, foram várias as queixas feitas ao local do certame, nomeadamente o facto de ser uma cidade remota no interior. Devido à inexistência de capacidade hoteleira, os representantes dos países ficaram alojados em \textbf{Killarney} (LOC) , a 30km da cidade.

Pela primeira vez, uma mulher foi a diretora da final da competição, \textbf{Anita Notaro} (PER) .

A \textbf{EBU} (MISC) introduziu uma nova regra naquele ano para regular os muitos países que desejavam participar da competição, a repescagem. A partir de agora, os seis países que terminassem nos últimos lugares da classificação final perdiam o direito de competir no ano seguinte. Esta regra permaneceu em vigor até 2004 com a introdução das semifinais.

Para fazer face ao aumento de países interessados em participar no \textbf{Festival Eurovisão da Canção} (MISC), a \textbf{EBU} (MISC) realizou, a 3 de abril de 1993 , uma pré-selecção para o Festival Eurovisão da Canção 1993 - \textbf{Kvalifikacija za Millstreet} (PER). Esta teve lugar em \textbf{Ljubljana} (LOC) , na \textbf{Eslovénia} (LOC) , onde sete novos países da \textbf{Europa} (LOC) de Leste lutaram por um dos três lugares disponíveis para a final.

Os representantes dos sete países a concurso interpretaram, além do tema que apresentavam a concurso por um lugar no \textbf{Festival Eurovisão da Canção} (MISC), um outro tema do seu reportório. A interpretação deste último tema em nada influenciou a votação do júri.

Após uma votação bastante renhida e marcada pela votação entre países vizinhos, a \textbf{Eslovénia} (LOC), a \textbf{Bósnia-Herzegovina} (ORG) e a \textbf{Croácia} (LOC) alcançaram os três primeiros ligares que lhe garantiram o acesso à final do Festival Eurovisão da Canção 1993.



\textbf{Visual} (MISC)A abertura da competição começou com um vídeo. Na primeira parte, foi encenado um episódio retirado da mitologia celta: a história de amor entre \textbf{Eochaid Airem} (PER) e \textbf{Pewter} (LOC). Na segunda, exibições turísticas do interior da \textbf{Irlanda} (LOC) e da cidade de \textbf{Millstreet} (LOC) foram mostradas. A câmera então mostrou \textbf{Davy Spillane} (PER) , que tocou música tradicional irlandesa durante o vídeo. A abertura termina com outro vídeo, filmado durante a semana de ensaios e mostrando a recepção de delegações estrangeiras e as festividades organizadas.

O palco foi criado por \textbf{Alan Farquharson} (PER) , que também foi chefe de produção dois anos depois em \textbf{Dublin} (LOC) . O cenário era constituída por uma grande plataforma trapezoidal que jogado na extremidade inferior de uma estrutura namoro suspenso a partir do teto, e sob a superfície translúcida de ambas as estruturas de tubos de néon teve com o cenário mudou de cor em cada atuação.

A apresentadora foi \textbf{Fionnuala Sweeney} (PER) , que falou aos espectadores em irlandês , inglês e francês .

Os cartões postais mostravam os participantes à descoberta da riqueza cultural e turística da \textbf{Irlanda} (LOC) .

O intervalo foi ocupado por \textbf{Linda Martin} (PER) e \textbf{Johnny Logan} (PER) , acompanhados pelos coros da \textbf{Escola de Música de Cork} (ORG) e de \textbf{Millstreet} (LOC) .



\textbf{Votação} (MISC)Cada país tinha um júri composto por 11 elementos, que atribuiu 12, 10, 8, 7, 6, 5, 4, 3, 2, 1 pontos às dez canções mais votadas.

O supervisor executivo da \textbf{EBU} (MISC) foi, pela primeira vez \textbf{Christian Clausen} (PER) .

Durante a votação, a câmera fez vários close-ups dos artistas. Em particular, \textbf{Niamh Kavanagh} (PER) , \textbf{Anabela} (PER) e o \textbf{Annie Cotton} (PER) apareceram.

\textbf{Fionnuala Sweeney} (PER) iniciou a votação cumprimentando \textbf{Frank Naef} (PER) , que estava presente na sala.

Devido a problemas técnicos a divulgação das pontuações do \textbf{júri maltês} (MISC) foi adiada para o final da votação.

A apresentadora, quando o júri holandês atribuiu 10 pontos à \textbf{Irlanda} (LOC) , a apresentadora anunciou "\textbf{Irlanda} (LOC), 12 pontos", corrigido depois o erro.

Um momento particularmente emocionante foi quando a \textbf{Bósnia Herzegovina} (LOC) deu a sua votação. Devido à guerra civil que vivia o estado, existia o receio de não poder receber o seu voto devido a dificuldades de comunicação. Porém, quando a voz do porta-voz da \textbf{Bósnia} (LOC) foi ouvida com interferência, o público aplaudiu profusamente.



\textbf{Festival} (PER)A ordem de votação no \textbf{Festival Eurovisão da Canção 1993} (MISC), foi a seguinte:

Os países que receberam 12 pontos foram os seguintes:

Em baixo encontra-se a lista de maestros que conduziram a orquestra, na respectiva actuação de cada país concorrente.



Artistas repetentes

Alguns artistas repetiram a sua experiência \textbf{Eurovisiva} (PER). Em 1993, os repetentes foram:



Festival Eurovisão da Canção 1994



\textbf{Países Participantes Países} (LOC) rebaixados não podem participar Países que participaram no passado, mas não em 1994

O \textbf{Festival Eurovisão da Canção 1994} (MISC) (em inglês : \textbf{Eurovision Song Contest} (ORG) 1994 ; em francês : \textbf{Concours Eurovision} (MISC) de la chanson 1994 ; em irlandês : \textbf{Comórtas Amhránaíochta} (LOC) na \textbf{hEoraifíse} (MISC) 1994 ) foi o 39 edição anual do evento e teve lugar a 30 de abril de 1994 em \textbf{Dublin} (LOC) . Os apresentadores do evento foram Cynthia Ní \textbf{Mhurchú} (PER) e \textbf{Gerry Ryan} (PER) . \textbf{Paul Harrington} (PER) e \textbf{Charlie McGettigan} (PER) foram os vencedores do festival com a canção \textbf{Rock 'N' Roll Kids} (MISC) .

A \textbf{Irlanda} (LOC) foi representada por \textbf{Paul Harrington \& Charlie Mcgettigan} (PER) com " \textbf{Rock 'n' Roll Kids} (MISC) ", que alcançaram o 1 lugar com 226 pontos, ultrapassando pela primeira vez a meta de 200 pontos. Ao sagrar-se vencedora pela 6 vez, a \textbf{Irlanda} (LOC) alcançou também a sua terceira vitória consecutiva, sendo o único a atingir esta marca.

A \textbf{Polónia} (LOC) esteve envolvida num grande escândalo logo no ano da sua estreia. \textbf{Edyta Górniak} (PER) optou por cantar em inglês no ensaio geral que até 1997 era visualizado pelos júris nacionais, quebrando assim a regra que obrigava os países participantes a cantarem numa das suas línguas oficiais. Seis países propuseram que este tema fosse desclassificado e treze propuseram um impedimento definitivo à participação deste país no festival. No entanto, nada foi feito e a \textbf{Polónia} (LOC) alcançou o segundo lugar na sua primeira participação.



\textbf{Local} (MISC)O \textbf{Festival Eurovisão da Canção} (MISC) 1994 ocorreu em \textbf{Dublin} (LOC) , na \textbf{Irlanda} (LOC) . \textbf{Dublin} (LOC) é a capital e maior cidade da \textbf{Irlanda} (LOC) . O nome em inglês deriva da palavra irlandesa " \textbf{Dubhlinn} (LOC) " (ocasionalmente também grafada \textbf{Duibhlinn} (PER) ou \textbf{Dubh Linn} (PER) ), que significa "\textbf{Lago Negro} (LOC)". Localiza-se na província de \textbf{Leinster} (LOC) próxima ao ponto mediano da costa leste da \textbf{Irlanda} (LOC), sendo cortada pelo \textbf{Rio Liffey} (LOC) e o centro da região de \textbf{Dublin} (LOC). Desde 1898 possui nível administrativo de condado ( county-boroughs ). Seus limites são os condados de \textbf{Fingal} (LOC) a norte, \textbf{Dublin} (LOC) meridional a sudoeste e \textbf{Dun Laoghaire-Rathdown} (LOC) a sudeste. Tem uma população de 527.612 habitantes na cidade, e sua área metropolitana tem 1.804.156 habitantes. Fundada como um assentamento viquingue , foi o centro do \textbf{Reino de Dublin} (LOC) e se tornou a principal cidade da \textbf{Ilha} (LOC) após a invasão dos \textbf{Normandos} (LOC) . A cidade cresceu de maneira rápida durante o século XVII ; se tornou na época a segunda maior cidade do \textbf{Império Britânico} (ORG) e a quinta maior da \textbf{Europa} (LOC) . \textbf{Dublin} (LOC) entrou em um período de estagnação após o \textbf{Ato de União de 1800} (MISC) , mas continuou o centro económico da \textbf{Ilha} (LOC). Após a \textbf{Partição da Irlanda} (MISC) em 1922, virou a capital do \textbf{Estado Livre Irlandês} (MISC) , e mais tarde, da \textbf{República da Irlanda} (LOC) . \textbf{Dublin} (LOC) é reconhecida como uma cidade global , com um ranking "\textbf{Alpha-} (MISC)", colocando a cidade entre as 30 mais globalizadas do mundo. Atualmente é o principal centro histórico, cultural, económico, industrial e educacional da \textbf{Irlanda} (LOC).

O festival em si realizou-se no \textbf{Point Theatre} (LOC) (agora designado de \textbf{3Arena} (LOC)), na capital do país, \textbf{Dublin} (LOC) . O \textbf{Point Theatre} (LOC) , construída em 1998 , foi um espaço destinado a atrações públicas e festivais localizado perto do \textbf{Rio Liffey} (LOC) , em \textbf{Dublin} (LOC) , \textbf{Portugal} (LOC) . Com uma capacidade de 8 500 espectadores, o \textbf{Point Theatre} (LOC) era conhecido pela sua flexibilidade de organização de eventos, que poderiam ir desde eventos musicais , a convenções , eventos desportivos e espetáculos circenses .



\textbf{Formato Para} (PER) fazer face ao aumento de países interessados em participar no \textbf{Festival Eurovisão da Canção} (MISC) , a \textbf{EBU} (MISC) adoptou um novo método de selecção dos participantes. Assim, os 18 melhores países do ano anterior estavam directamente apurados e a estes juntaram-se os países estreantes. Os países afastados do concurso passaram a ser designados "passivos" e era-lhes assegurada a participação no ano seguinte desde que continuassem a transmitir o festival.

Foi ao fim de 38 anos de votação do júri europeu que pela primeira vez os telespectadores puderam ver a imagem do porta-voz de cada júri nacional a atribuir as suas pontuações. Esta novidade provocou algumas gargalhadas do público pois a extravagância da imagem de alguns porta-vozes era notória e propicia a uma boa gargalhada.

Apesar de as votações serem provenientes do júri nacional, em \textbf{Portugal} (LOC) foi feita um experiência com base no televoto. Assim, foram dados os números de telefone para onde os telespectadores deveriam votar. Essa votação deu o primeiro lugar a \textbf{Portugal} (LOC) , o segundo à \textbf{Alemanha} (LOC) , o terceiro à \textbf{Irlanda} (LOC) , o quarto ao \textbf{Reino Unido} (LOC) e o quinto à \textbf{Áustria} (LOC) .



\textbf{Visual} (MISC)A abertura da competição começou com um vídeo mostrando estrelas a flutuar na água, fogos de artificio e caricaturas a dançar, beber café e a andar de bicicleta. Após esse pequeno filme, foi mostrada a arena em si, onde dançarinos vestidos de branco e a usar caricaturas de figuras conhecidas da \textbf{Irlanda} (LOC), deram as boas vindas ao público, carregando bandeiras de \textbf{Países europeus} (LOC). Os apresentadores chegaram à arena numa plataforma que desceu até ao palco. Os cartões-postais desse ano tiveram um tema literário, mostrando os participantes a ler, pescar e a fazer outras atividades pela \textbf{Irlanda} (LOC).

O palco, da autoria de \textbf{Paula Farrell} (PER) , incluía no seu design uma cidade com arranha-céus com 500 lâmpadas arquitecturais, cinco projectores gigantes e duzentos ecrãs, além do pano de fundo de um céu noturno em constante mudança, baseava-se no conceito de como seria uma \textbf{Dublin} (LOC) futurista, com uma constante remanescente no \textbf{rio Liffey} (LOC). O chão foi pintado com uma tinta reflexiva azul escura para dar um efeito de água.

Os apresentadores foram Cynthia Ní \textbf{Mhurchú} (PER) e \textbf{Gerry Ryan} (PER) , que falaram aos espectadores em irlandês , inglês e francês .

Os cartões postais mostravam os participantes à descoberta da riqueza cultural e turística da \textbf{Irlanda} (LOC) .

O intervalo foi ocupado pela primeira interpretação dos \textbf{Riverdance} (MISC) , uma mistura de dança folclórica irlandesa e sapateado. Originalmente, a \textbf{RTÉ} (MISC) tinha encomendado ao compositor \textbf{Bill Whelan} (PER) , para ele escrever uma peça de música celta tradicional. A coreografia do número foi confiada aos campeões mundiais de dança \textbf{Michael Flatley} (PER) e \textbf{Jean Butler} (PER) . Na noite do concurso, \textbf{Riverdance} (MISC) recebeu uma enorme ovação pública no salão. Posteriormente, a faixa permaneceu classificada durante dezoito semanas nas paradas irlandesas e tornou-se um grande sucesso comercial, não apenas da edição de 1994 do concurso, mas também de todos os números de intervalo já apresentados. Mais tarde foi estendido para show completo que o conhecia como um grande sucesso em todo o mundo e foi visto por milhões de pessoas.



\textbf{Votação} (MISC)Cada país tinha um júri composto por 11 elementos, que atribuiu 12, 10, 8, 7, 6, 5, 4, 3, 2, 1 pontos às dez canções mais votadas.

O supervisor executivo da \textbf{EBU} (MISC) foi \textbf{Christian Clausen} (PER) .

Durante a votação, a câmera fez vários close-ups dos artistas. Em particular, \textbf{Paul Harrington} (PER) e \textbf{Charlie McGettigan} (PER) , \textbf{Edyta Górniak} (PER) , \textbf{Friderika Bayer} (PER) e \textbf{Sara Tavares} (PER) apareceram.

Devido aos avanços tecnológicos, este foi o primeiro festival em que se podia ver a imagem do porta-voz do júri nacional. Nos inícios, pensou-se na possível vitória da \textbf{Hungria} (LOC) (pontuações máximas dos 3 1s júris nacionais), mas as votações seguintes foram mais fracas e deram mais uma vitória à \textbf{Irlanda} (LOC), o primeiro país a conseguir vitória três vezes consecutivas.



\textbf{Festival} (PER)A ordem de votação no \textbf{Festival Eurovisão da Canção 1994} (MISC), foi a seguinte:

Os países que receberam 12 pontos foram os seguintes:

Em baixo encontra-se a lista de maestros que conduziram a orquestra, na respectiva actuação de cada país concorrente.



Artistas repetentes

Alguns artistas repetiram a sua experiência \textbf{Eurovisiva} (PER). Em 1994, os repetentes foram:



Festival Eurovisão da Canção 1995



\textbf{Países Participantes Países} (LOC) rebaixados não podem participar Países que participaram no passado, mas não em 1995

O \textbf{Festival Eurovisão da Canção} (MISC) 1995 (em inglês : \textbf{Eurovision Song Contest} (ORG) 1995 ; em francês : \textbf{Concours Eurovision} (MISC) de la chanson 1995 ; em irlandês : \textbf{Comórtas Amhránaíochta} (LOC) na \textbf{hEoraifíse} (MISC) 1995 ) foi o 40 edição anual do evento e teve lugar a 13 de maio de 1995 em \textbf{Dublin} (LOC) . A apresentadora do evento foi \textbf{Mary Kennedy} (PER) . O duo \textbf{Secret Garden} (PER) foi o vencedor do festival com a canção \textbf{Nocturne} (PER) , que conta com o instrumental mais longo de sempre no \textbf{Festival Eurovisão da Canção} (MISC) uma vez que este tema continha apenas 25 palavras ao longo dos 3 minutos da actuação. Ironicamente, a violinista do grupo \textbf{Secret Garden} (PER), \textbf{Fionnuala Sherry} (PER) , é de nacionalidade irlandesa.

A \textbf{Suécia} (LOC) , com a balada pop " Se på mig ", era a grande favorita à vitória, juntamente com a \textbf{Croácia} (LOC) , \textbf{Dinamarca} (LOC) , \textbf{Espanha} (LOC) , \textbf{Israel} (LOC) e a vencedora \textbf{Noruega} (LOC) .



\textbf{Local} (PER)Conseguindo tornar-se o único país a organizar a \textbf{Eurovisão} (LOC) três vezes consecutivas, o \textbf{Festival Eurovisão da Canção} (MISC) 1995 ocorreu em \textbf{Dublin} (LOC) , na \textbf{Irlanda} (LOC) . \textbf{Dublin} (LOC) é a capital e maior cidade da \textbf{Irlanda} (LOC) . O nome em inglês deriva da palavra irlandesa " \textbf{Dubhlinn} (LOC) " (ocasionalmente também grafada \textbf{Duibhlinn} (PER) ou \textbf{Dubh Linn} (PER) ), que significa "\textbf{Lago Negro} (LOC)". Localiza-se na província de \textbf{Leinster} (LOC) próxima ao ponto mediano da costa leste da \textbf{Irlanda} (LOC), sendo cortada pelo \textbf{Rio Liffey} (LOC) e o centro da região de \textbf{Dublin} (LOC). Desde 1898 possui nível administrativo de condado ( county-boroughs ). Seus limites são os condados de \textbf{Fingal} (LOC) a norte, \textbf{Dublin} (LOC) meridional a sudoeste e \textbf{Dun Laoghaire-Rathdown} (LOC) a sudeste. Tem uma população de 527.612 habitantes na cidade, e sua área metropolitana tem 1.804.156 habitantes. Fundada como um assentamento viquingue , foi o centro do \textbf{Reino de Dublin} (LOC) e se tornou a principal cidade da \textbf{Ilha} (LOC) após a invasão dos \textbf{Normandos} (LOC) . A cidade cresceu de maneira rápida durante o século XVII ; se tornou na época a segunda maior cidade do \textbf{Império Britânico} (ORG) e a quinta maior da \textbf{Europa} (LOC) . \textbf{Dublin} (LOC) entrou em um período de estagnação após o \textbf{Ato de União de 1800} (MISC) , mas continuou o centro económico da \textbf{Ilha} (LOC). Após a \textbf{Partição da Irlanda} (MISC) em 1922, virou a capital do \textbf{Estado Livre Irlandês} (MISC) , e mais tarde, da \textbf{República da Irlanda} (LOC) . \textbf{Dublin} (LOC) é reconhecida como uma cidade global , com um ranking "\textbf{Alpha-} (MISC)", colocando a cidade entre as 30 mais globalizadas do mundo. Atualmente é o principal centro histórico, cultural, económico, industrial e educacional da \textbf{Irlanda} (LOC).

O festival em si realizou-se no \textbf{Point Theatre} (LOC) (agora designado de \textbf{3Arena} (LOC)), na capital do país, \textbf{Dublin} (LOC) . O \textbf{Point Theatre} (LOC) , construída em 1998 , foi um espaço destinado a atrações públicas e festivais localizado perto do \textbf{Rio Liffey} (LOC) , em \textbf{Dublin} (LOC) , \textbf{Portugal} (LOC) . Com uma capacidade de 8 500 espectadores, o \textbf{Point Theatre} (LOC) era conhecido pela sua flexibilidade de organização de eventos, que poderiam ir desde eventos musicais , a convenções , eventos desportivos e espetáculos circenses .



\textbf{Formato} (PER)Esta foi a última edição até 2013 que só incluiu uma apresentadora.

Devido ao aumento repentino de países participantes (e interessados em participar), a \textbf{EBU} (MISC) limitou o número de países participantes a 23 por festival de forma a certificar-se que o espectáculo não excederia as 3 horas de duração, excluindo da edição a \textbf{Eslováquia} (LOC) , a \textbf{Estónia} (LOC) , a \textbf{Finlândia} (LOC) , a \textbf{Lituânia} (LOC), os \textbf{Países Baixos} (LOC) , a \textbf{Roménia} (LOC) e a \textbf{Suíça} (LOC) . Inicialmente, a \textbf{Áustria} (LOC) e a \textbf{Espanha} (LOC) também seriam excluídas, mas, devido à desistência do \textbf{Luxemburgo} (LOC) (que nunca mais voltou a participar) e da \textbf{Itália} (LOC) (que só regressaria em 1997 ), puderam participar.

Foi muita a especulação acerca da participação irlandesa pois dizia-se que a \textbf{RTÉ} (MISC) (televisão irlandesa) iria escolher um tema com poucas hipóteses de ganhar de forma a evitar ter de produzir pela quarta vez consecutiva o \textbf{Festival Eurovisão da Canção} (MISC). A escolha incidiu sobre o tema " \textbf{Dreamin} (PER)' " que foi interpretado por \textbf{Eddie Friel} (PER) alcançando o 14 lugar.

Após ganhar o \textbf{Festival Eurovisão da Canção 1994} (MISC) , a \textbf{RTÉ} (MISC) ficou preocupada se conseguiria organizar a edição de 1995, pela terceira vez consecutiva. Entretanto, a \textbf{BBC} (ORG) ofereceu-se para organizar a edição, em conjunto com a televisão irlandesa, em \textbf{Belfast} (LOC) , capital da \textbf{Irlanda do Norte} (LOC) . No entanto, colhendo os frutos do \textbf{Tigre Celta} (MISC) , o governo irlandês disponibilizou à \textbf{RTÉ} (MISC) fundos para organizar o certame europeu, que iria para 40 edição. Porém, a televisão irlandesa pediu à \textbf{União Europeia de Radiodifusão} (ORG) que, se ganhasse pela quarta vez consecutiva, não iria organizar o evento de 1996.



\textbf{Visual} (MISC)A abertura da competição começou com um vídeo turístico sobre a \textbf{Irlanda} (LOC) , misturando imagens do passado e do presente, tradições e modernidade do país. A câmara então mostrou a sala, mergulhada em meia-luz e uma tela ampla, na qual foram projetadas imagens dos vencedores anteriores do concurso. Em seguida, seguiu-se um segundo vídeo, em homenagem aos quarenta anos da \textbf{Eurovisão} (LOC). Foi uma linha do tempo das várias edições, que destacou as seis vitórias da \textbf{Irlanda} (LOC). A câmera voltou para a tela gigante que se levantou. \textbf{Mary Kennedy} (PER) fez sua entrada, descendo uma escadaria iluminada, enquanto as velas que escondiam a cena foram puxadas para os bastidores. Uma vez chegada ao palco, as escadas subiram e desapareceram, revelando o pódio numa chuva de faíscas.

O palco, da autoria de \textbf{Alan Farquharson} (PER) , que já tinha desenhado o palco do \textbf{Festival Eurovisão da Canção 1993} (MISC) , ofereceu a esta edição uma decoração bastante sombria, consistindo de uma grande cena triangular. No fundo, vários jogos animados de luzes em movimento. Também havia a forma de uma flecha gigantesca, que descia do teto e alcançava o público. Esta flecha era flanqueada por dois pódios secundários, de forma retangular e destinados a músicos e coristas. O conjunto era de cor escura e delimitado por bandas de luz. Variações foram trazidas por peças de luz, projeções de imagem, lasers verdes, bem como pela presença de vários móveis metálicos.

A apresentadora foi \textbf{Mary Kennedy} (PER) , que falou aos espectadores em irlandês , inglês e francês .

Os cartões-portais consistiam em exibições turísticas da \textbf{Irlanda} (LOC) e as suas riquezas culturais.

O intervalo consistiu em vários artistas como a banda \textbf{Clannad} (ORG) , \textbf{Brian Kennedy} (PER) (que 11 anos depois representaria a \textbf{Irlanda} (LOC)) e o compositor \textbf{Michael O'Suilleabhan} (PER) .



\textbf{Votação} (MISC)Cada país tinha um júri composto por 11 elementos, que atribuiu 12, 10, 8, 7, 6, 5, 4, 3, 2, 1 pontos às dez canções mais votadas.

O supervisor executivo da \textbf{EBU} (MISC) foi \textbf{Christian Clausen} (PER) .

Durante a votação, a câmera fez vários close-ups dos artistas. Em particular, \textbf{Secret Garden} (PER) e \textbf{Anabel Conde} (PER) apareceram.



\textbf{Festival} (MISC)A ordem de votação no \textbf{Festival Eurovisão da Canção} (MISC) 1995, foi a seguinte:

Os países que receberam 12 pontos foram os seguintes:

Em baixo encontra-se a lista de maestros que conduziram a orquestra, na respectiva actuação de cada país concorrente.



Artistas repetentes

Alguns artistas repetiram a sua experiência \textbf{Eurovisiva} (PER). Em 1995, os repetentes foram:



Festival Eurovisão da Canção 1996



Países que classificaram-se para a final \textbf{Países} (MISC) que não se classificaram para a final \textbf{Países} (MISC) que participaram no passado, mas não em 1996

O \textbf{Festival Eurovisão da Canção 1996} (MISC) (em inglês : \textbf{Eurovision Song Contest} (ORG) 1996 ; em francês : \textbf{Concours Eurovision} (MISC) de la chanson 1996 ; em norueguês : \textbf{Internasjonale} (LOC)/\textbf{Europeiske Melodi Grand Prix} (MISC) 1996 ), ou \textbf{Eurosong} (LOC) 96 foi a 41 edição anual do evento e teve lugar a 18 de maio de 1996 em \textbf{Oslo} (LOC) . Os apresentadores foram \textbf{Morten Harket} (PER) , vocalista da banda norueguesa \textbf{A-ha} (ORG) e a jornalista \textbf{Ingvild Bryn} (PER) . A cantora \textbf{Eimear Quinn} (PER) , que representou a \textbf{Irlanda} (LOC) , foi a vencedora deste ano com a canção \textbf{The voice} (MISC) .

\textbf{Portugal} (LOC) conseguiu com \textbf{Lúcia Moniz} (PER) e a sua canção O meu coração não tem cor a sua melhor classificação até à vitória em 2017 de \textbf{Salvador Sobral} (LOC) .



\textbf{Local} (MISC)O \textbf{Festival Eurovisão da Canção} (MISC) 1996 ocorreu em \textbf{Oslo} (LOC) , na \textbf{Noruega} (LOC) . \textbf{Oslolocaliza} (LOC) -se no sudeste do país e detém estatuto de comuna e condado simultaneamente. Fundada por 1048 pelo \textbf{Rei Harald III} (PER) "Hardråde" da \textbf{Noruega} (LOC), a cidade foi imensamente destruída por um incêndio em 1624. O rei dano-norueguês \textbf{Chirstian IV} (PER) reconstruiu a cidade, com o nome de \textbf{Cristiânia} (LOC) , entre 1624 e 1924 . Em 1952 a cidade foi a sede dos \textbf{Jogos Olímpicos de Inverno de 1952} (MISC) . É a cidade onde é entregue o \textbf{Prémio Nobel da Paz} (MISC) . Em 2006 \textbf{Oslo} (LOC) foi eleita pela \textbf{BBC} (ORG) como a cidade mais cara do mundo. A capital foi eleita a cidade com a 24 melhor qualidade de vida do mundo, ao passo que a \textbf{Noruega} (LOC) liderava o ranking dos países na mesma categoria. Oslo é o centro cultural, científico, econômico e governamental da \textbf{Noruega} (LOC). Tem sua atenção voltada para negociações, bancos, indústrias e navegações. É também um importante centro para indústrias marítimas e tratados marítimos na \textbf{Europa} (LOC). A cidade é o "lar" de muitas empresas voltadas para o setor marítimo, algumas delas sendo umas das maiores do mundo. Oslo é considerada uma cidade global. Por muitos anos, \textbf{Oslo} (LOC) vem sendo listada como uma das mais caras cidades do mundo, bem como outras capitais globais como \textbf{Copenhague} (LOC), \textbf{Paris} (LOC) e \textbf{Tóquio} (LOC). Em 2009, no entanto, \textbf{Oslo} (LOC) teve seu status alterado para a cidade mais cara do mundo.

O festival em si realizou-se no \textbf{Oslo Spektrum} (ORG) , construído em 1990 , é um espaço multiusos, mais conhecido por receber todos os anos o concerto do \textbf{Nobel da Paz} (MISC) .



\textbf{Formato Para} (PER) fazer face ao aumento de países interessados em participar no \textbf{Festival Eurovisão da Canção} (MISC) , a \textbf{União Europeia de Radiodifusão} (ORG) idealizou uma pré-selecção a 20 e 21 de março de 1996 . Os júris nacionais ouviram as gravações áudio dos temas a concurso e pontuaram de 1 a 12 pontos as suas dez canções favoritas. Dos 29 países a concurso apenas 22 alcançaram a final onde se juntaram à \textbf{Noruega} (LOC) que obteve qualificação directa por ter vencido o \textbf{Festival Eurovisão da Canção} (MISC) 1995 . Além disso, inscrições da \textbf{Bulgária} (LOC) , \textbf{Moldávia} (LOC) e \textbf{Ucrânia} (LOC) foram relatadas, mas não se concretizaram, com o nome da \textbf{Bulgária} (LOC) ainda aparecendo nas folhas de votação do júri de qualificação. Todos os três países fariam sua estreia nos anos 2000 .

Este foi o único ano em que a \textbf{Alemanha} (LOC) , um dos maiores financiadores do \textbf{Festival Eurovisão da Canção} (MISC), não esteve presente na final deste concurso desde 1956 . Este, aliás, foi um dos motivos pelo qual esta prática foi abandonada logo no ano seguinte.

O festival contou com a presença da rainha \textbf{Sônia da Noruega} (PER) .

Pela primeira vez, a supervisora-executiva do \textbf{Festival Eurovisão da Canção} (MISC) foi uma mulher, \textbf{Christine Marchal Ortiz} (PER) .

Uma outra grande inovação foi, pela primeira vez, um site dedicado a esta edição, que foi colocado online pela \textbf{NRK} (ORG) .

A \textbf{União Europeia de Radiodifusão} (ORG) voltou a mudar o regulamento para encurtar a lista de países participantes, fazendo com que todos os países interessados em participar, com exceção da \textbf{Noruega} (LOC) -- vencedora do ano anterior -- passassem por uma fase de qualificação, realizada em março de 1996 , para que os júris nacionais ouvissem as músicas e as classificassem. As 22 canções mais votadas passariam à final. Esta fase de qualificação nunca foi transmitida, quer na televisão , quer na rádio . Este sistema mostrou-se desde logo insustentável, em grande parte devido à eliminação da \textbf{Alemanha} (LOC) , que protestou contra esta prática, fazendo com que esta fosse abandonada logo no ano seguinte.

Todas as novas práticas introduzidas (novo nome do certame, fase de qualificação, mensagem de boa sorte e blue room), com exceção da transmissão internet, foram abandonadas a partir do ano seguinte.



\textbf{Visual} (MISC)O festival abriu com um vídeo reminiscente do passado viking e conquistador da \textbf{Noruega} (LOC) , acabando com o apresentador, \textbf{Morten Harket} (PER) , a cantar "\textbf{Heaven's Not For Saints} (MISC)". No final de sua apresentação, ele cumprimentou a orquestra e seu maestro, \textbf{Frode Thingnæs} (MISC). \textbf{Ingvild Bryn} (PER), por sua vez, fez sua aparição e juntos fizeram as apresentações habituais. Eles cumprimentaram todos os países participantes no seu idioma nacional. Concluíram mencionando o procedimento da fase de qualificação e os países não qualificados.

O palco consistia num vasto pódio de três níveis, com contornos irregulares. Duas escadas de metal em ambos os lados permitiam o acesso. A decoração consistiu de uma banda de malha suportada por pilares luminosos e acima de tudo uma ampla gama de andaimes e metálicos móvel elementos, inspirados no petróleo , uma das imagens de marca da \textbf{Noruega} (LOC) .

Os apresentadores foram \textbf{Ingvild Bryn} (PER) e \textbf{Morten Harket} (ORG) , que falou aos espectadores em norueguês , inglês e francês .

Os cartões-postais desta edição eram divididos em três partes: os participantes nos seus países, seguidos por algumas imagens de um traço característico da \textbf{Noruega} (LOC) e, finalmente, uma pequena mensagem de boa sorte para cada intérprete por um líder político de cada país participante, que podiam ser primeiros-ministros, presidentes ou outros. No caso de \textbf{Portugal} (LOC) , o então \textbf{Primeiro-Ministro António Guterres} (PER) , desejou a \textbf{Lúcia Moniz} (PER) e a toda a delegação portuguesa: Desejo felicidades à equipa portuguesa, mas desejo sobretudo, que ganhe o melhor e que a música seja cada vez mais um fator de paz e de aproximação entre os povos.

O intervalo começou com um vídeo de um norueguês em traje tradicional, cantando no topo de um pico nevado. Seguiram-se várias peças, clássicas ou contemporâneas, sempre interpretadas pelos artistas em paisagens típicas da \textbf{Noruega} (LOC) . A câmera então voltou ao palco para o balé contemporâneo \textbf{Beacon Burning} (PER), realizado pelo \textbf{Oslo Danse Ensemble} (ORG). Os dançarinos foram acompanhados por mágicos, malabaristas e comedores de fogo. No final do balé, o norueguês do vídeo apareceu no palco, para concluir o intervalo pelas mesmas notas em que havia começado.



\textbf{Votação} (MISC)Cada país tinha um júri composto por 11 elementos, que atribuiu 12, 10, 8, 7, 6, 5, 4, 3, 2, 1 pontos às dez canções mais votadas.

A supervisora executiva da \textbf{EBU} (MISC) foi \textbf{Christine Marchal-Ortiz} (PER) .

Durante a votação, a câmera fez vários close-ups dos artistas. Em particular, \textbf{Eimear Quinn} (MISC) , \textbf{Elisabeth Andreassen} (PER) , \textbf{One More Time} (MISC) e \textbf{Lúcia Moniz} (PER) apareceram.

Para a votação, foi utilizada uma "tela azul", que foi fornecido pela \textbf{Silicon Graphics} (ORG) , sendo único na história do concurso, onde eram projectadas animações \textbf{3D} (MISC) com recurso a tecnologias de realidade virtual. Além disso, a televisão norueguesa utilizou na preparação e na apresentação das canções diferentes efeitos especiais. Pela primeira vez, um porta-voz dirgiu-se ao palco, mais propriamente à blue room para divulgar os votos do júri nacional. Tratou-se da porta-voz norueguesa, \textbf{Ragnhild Sælthun Fjørtoft} (MISC) .



\textbf{Festival} (PER)Os países que receberam 12 pontos foram os seguintes:

A ordem de votação no \textbf{Festival Eurovisão da Canção 1996} (MISC), foi a seguinte:

Os países que receberam 12 pontos foram os seguintes:

Em baixo encontra-se a lista de maestros que conduziram a orquestra, na respectiva actuação de cada país concorrente.



Mensagens de boa sorte

Em 1996 , todos os participantes receberam uma pequena mensagem de boa sorte por parte de um líder político de cada país participante, que podiam ser primeiros-ministros, presidentes ou outros. Esta foi aliás, a única edição em que tal aconteceu.



Artistas repetentes

Alguns artistas repetiram a sua experiência \textbf{Eurovisiva} (PER). Em 1996, os repetentes foram:



\textbf{Festival Eurovisão da Canção 1997} (MISC)



\textbf{Países Participantes Países} (LOC) rebaixados não podem participar Países que participaram no passado, mas não em 1997

O \textbf{Festival Eurovisão da Canção 1997} (MISC) (em inglês : \textbf{Eurovision Song Contest} (ORG) 1997 , em francês : \textbf{Concours Eurovision} (MISC) de la chanson 1997 e em irlandês : \textbf{Comórtas Amhránaíochta} (LOC) na \textbf{hEoraifíse} (MISC) 1997 ) foi o 42 \textbf{Festival Eurovisão da Canção} (MISC) e realizou-se a 3 de maio de 1997 em \textbf{Dublin} (LOC) , na \textbf{Irlanda} (LOC) . Os apresentadores foram \textbf{Ronan Keating} (PER) e \textbf{Carrie Crowley} (PER) . \textbf{Katrina and the Waves} (MISC) foram os vencedores do \textbf{Festival} (LOC) desse ano com a canção \textbf{Love Shine a Light} (MISC) ( O amor ilumina uma luz ) que recebeu a pontuação máxima de 10 países.



\textbf{Local} (MISC)O \textbf{Festival Eurovisão da Canção} (MISC) 1997 ocorreu em \textbf{Dublin} (LOC) , na \textbf{Irlanda} (LOC) . \textbf{Dublin} (LOC) é a capital e maior cidade da \textbf{Irlanda} (LOC) . O nome em inglês deriva da palavra irlandesa " \textbf{Dubhlinn} (LOC) " (ocasionalmente também grafada \textbf{Duibhlinn} (PER) ou \textbf{Dubh Linn} (PER) ), que significa "\textbf{Lago Negro} (LOC)". Localiza-se na província de \textbf{Leinster} (LOC) próxima ao ponto mediano da costa leste da \textbf{Irlanda} (LOC), sendo cortada pelo \textbf{Rio Liffey} (LOC) e o centro da região de \textbf{Dublin} (LOC). Desde 1898 possui nível administrativo de condado ( county-boroughs ). Seus limites são os condados de \textbf{Fingal} (LOC) a norte, \textbf{Dublin} (LOC) meridional a sudoeste e \textbf{Dun Laoghaire-Rathdown} (LOC) a sudeste. Tem uma população de 527.612 habitantes na cidade, e sua área metropolitana tem 1.804.156 habitantes. Fundada como um assentamento viquingue , foi o centro do \textbf{Reino de Dublin} (LOC) e se tornou a principal cidade da \textbf{Ilha} (LOC) após a invasão dos \textbf{Normandos} (LOC) . A cidade cresceu de maneira rápida durante o século XVII ; se tornou na época a segunda maior cidade do \textbf{Império Britânico} (ORG) e a quinta maior da \textbf{Europa} (LOC) . \textbf{Dublin} (LOC) entrou em um período de estagnação após o \textbf{Ato de União de 1800} (MISC) , mas continuou o centro económico da \textbf{Ilha} (LOC). Após a \textbf{Partição da Irlanda} (MISC) em 1922, virou a capital do \textbf{Estado Livre Irlandês} (MISC) , e mais tarde, da \textbf{República da Irlanda} (LOC) . \textbf{Dublin} (LOC) é reconhecida como uma cidade global , com um ranking "\textbf{Alpha-} (MISC)", colocando a cidade entre as 30 mais globalizadas do mundo. Atualmente é o principal centro histórico, cultural, económico, industrial e educacional da \textbf{Irlanda} (LOC).

O festival em si realizou-se no \textbf{Point Theatre} (LOC) (agora designado de \textbf{3Arena} (LOC)), na capital do país, \textbf{Dublin} (LOC) . O \textbf{Point Theatre} (LOC) , construída em 1998 , foi um espaço destinado a atrações públicas e festivais localizado perto do \textbf{Rio Liffey} (LOC) , em \textbf{Dublin} (LOC) , \textbf{Portugal} (LOC) . Com uma capacidade de 8 500 espectadores, o \textbf{Point Theatre} (LOC) era conhecido pela sua flexibilidade de organização de eventos, que poderiam ir desde eventos musicais , a convenções , eventos desportivos e espetáculos circenses .



\textbf{Formato Depois} (PER) da controversa em torno da pré-selecção de 1996 que eliminou alguns países famosos nestes festivais, a \textbf{UER} (ORG) (\textbf{União Europeia de Radiofusão} (ORG)) introduziu um novo sistema para 1997: os países com o mais baixo score dos últimos 4 seriam excluídos de participarem no Festival Eurovisão da Canção. Aqueles países com as mais baixas pontuações nos últimos 5 anos não poderiam participar em futuros festivais. Contudo o país excluído num ano poderia participar no seguinte.

Inicialmente, \textbf{Israel} (LOC) iria particpar, mas teve que se abster, pois a data da final coincidia com o \textbf{Yom HaZikaron} (MISC) . A retirada permitiu que a \textbf{Bósnia e Herzegovina} (LOC) participasse.

Nesta edição, cinco países ( \textbf{Alemanha} (LOC) , \textbf{Áustria} (LOC) , \textbf{Suécia} (LOC) , \textbf{Suíça} (LOC) e \textbf{Reino Unido} (LOC) ) utilizaram pela primeira vez o sistema de votação por telefone como forma de contrariar os baixos índices de audiência nesses países e levar o evento a um público mais amplo e jovem. Os resultados do televoto foram, em alguns casos, diferentes daqueles entregues pelo júri profissional. A \textbf{Islândia} (LOC) , por exemplo, recebeu 16 dos seus 18 pontos graças a esse sistema.



\textbf{Visual} (MISC)A abertura da competição começou com um pequeno vídeo. Apareceram \textbf{Eimear Quinn} (MISC) , \textbf{Bjorn Ulvaeus} (PER) , \textbf{Céline Dion} (PER) e \textbf{Morten Harket} (PER) , todos desejando boa sorte aos participantes. A câmera então mostrou a sala e o palco.

O palco, projetado por \textbf{Paula Farrell} (PER) como os cenários de 1988 e 1994, foi inspirado no novo mundo das telecomunicações, com muitos monitores espalhados no fundo e no chão.

Os apresentadores foram \textbf{Carrie Crowley} (PER) e \textbf{Ronan Keating} (PER) , que falaram aos espectadores em irlandês , inglês e francês .

Os cartões postais mostravam a riqueza económica, cultural e turística da \textbf{Irlanda} (LOC) . Eles terminaram com uma apresentação visual dos artistas participantes. Dezoito ex-participantes, na sua maioria vencedores, apareceram pontualmente entre as apresentações. Eles se lembraram de sua participação no concurso e desejaram boa sorte aos concorrentes.

O intervalo foi ocupado pelo grupo \textbf{Boyzone} (ORG) , acompanhados por uma dúzia de dançarinos. A canção, "\textbf{Let The Message Run Free} (MISC)", foi especialmente escrita por \textbf{Ronan Keating} (PER) para a ocasião.



Votação

20 dos 25 países tinha um júri composto por 11 elementos, que atribuiu 12, 10, 8, 7, 6, 5, 4, 3, 2, 1 pontos às dez canções mais votadas. Os outros 5 ( \textbf{Alemanha} (LOC) , \textbf{Áustria} (LOC) , \textbf{Suécia} (LOC) , \textbf{Suíça} (LOC) e \textbf{Reino Unido} (LOC) ) usaram o televoto, onde, às 10 canções mais votadas, eram atriuídos 12, 10, 8, 7, 6, 5, 4, 3, 2, 1 ponto, com um júri de salvaguarda, em caso de erros. Um júri nacional era usado em casos de força maior, em que não se pudesse utilizar o televoto.

A supervisora executiva da \textbf{EBU} (MISC) foi \textbf{Marie-Claire Vionnet} (PER) .

Durante a votação, a câmera fez vários close-ups dos artistas. Em particular, \textbf{Katrina and the Waves} (MISC) , \textbf{Alla Pugacheva} (PER) e \textbf{Şebnem Paker} (PER) apareceram.

Devido ao televoto, pela primeira vez, houve um resumo das músicas no final das 25 atuações.

Pela primeira vez, duas pessoas anunciaram os votos. Trataram-se de \textbf{Frédéric Ferrer} (PER) e \textbf{Marie Myriam} (PER) , pela \textbf{França} (LOC).



\textbf{Prémio Barbara Dex} (MISC)Pela primeira vez, o fansite \textbf{House of Eurovision} (MISC) apresentou o \textbf{Prêmio Barbara Dex} (MISC), um prémio de humor dado ao artista mais mal vestido a cada ano no concurso. É nomeado após a artista belga, \textbf{Barbara Dex} (PER) , que ficou em último na edição de 1993, em que ela usava o seu próprio vestido projetado. A \textbf{House of Eurovision} (MISC) continuaria a fornecer o \textbf{Barbara Dex Award} (ORG) até 2016, quando outro fansite da \textbf{Eurovision} (ORG), songfestival.be, assumiu as rédeas do prémio e o apresenta todos os anos a partir de 2017.

\textbf{Debbie Scerri} (PER) , de \textbf{Malta} (LOC) foi a vencedora do prémio \textbf{Barbara Dex} (PER) de 1997.



\textbf{Festival} (PER)A ordem de votação no \textbf{Festival Eurovisão da Canção 1997} (MISC), foi a seguinte:

Os países que receberam 12 pontos foram os seguintes:

Para fazer face ao elevado número de países interessados em participar no certame, a \textbf{EBU} (MISC) resolveu aplicar uma nova regra: os 5 países que tivessem a menor média dos últimos 5 anos, não poderiam participar.

Em baixo encontra-se a lista de maestros que conduziram a orquestra, na respectiva actuação de cada país concorrente.



Artistas repetentes

Alguns artistas repetiram a sua experiência \textbf{Eurovisiva} (PER). Em 1997, os repetentes foram:



Festival Eurovisão da Canção 1998



\textbf{Países Participantes Países} (LOC) rebaixados não podem participar Países que participaram no passado, mas não em 1998

O \textbf{Festival Eurovisão da Canção 1998} (MISC) (em inglês : \textbf{Eurovision Song Contest} (ORG) 1998 e em francês : \textbf{Concours Eurovision} (MISC) de la chanson 1998 ) foi o 43 \textbf{Festival Eurovisão da Canção} (MISC) e realizou-se a 9 de maio de 1998 em \textbf{Birmingham} (LOC) , \textbf{Inglaterra} (LOC) , no \textbf{Reino Unido} (LOC) . Os apresentadores foram \textbf{Terry Wogan} (PER) e \textbf{Ulrika Jonsson} (PER) . \textbf{Dana International} (PER) , a representar \textbf{Israel} (LOC), venceu a edição deste ano com a canção " \textbf{Diva} (PER) ", escrito por \textbf{Svika Pick} (PER) e \textbf{Ginai Yoav} (PER). A cantora atraiu muita atenção, tanto em \textbf{Israel} (LOC) como na \textbf{Europa} (LOC), por causa da sua mudança de sexo em 1993, sendo a primeira artista transexual a entrar na competição.



\textbf{Local} (MISC)O \textbf{Festival Eurovisão da Canção} (MISC) 1998 ocorreu em \textbf{Birmingham} (LOC) , \textbf{Inglaterra} (LOC) , no \textbf{Reino Unido} (LOC) . \textbf{Birmingham} (LOC) é uma cidade e distrito metropolitano do condado de \textbf{Midlands Ocidentais} (LOC) , na \textbf{Inglaterra} (LOC) . É a segunda maior e mais importante cidade do \textbf{Reino Unido} (LOC) , atrás apenas da capital \textbf{Londres} (LOC) . Com uma população estimada de 1.028.701 (estimativa de 2009), integra a conurbação de \textbf{West Midlands} (LOC) , cuja população é de 2.284.093 habitantes (censo de 2001), sendo responsável por cerca 40\% da população do condado . É um centro industrial e de transportes. A cidade medieval era já um centro comercial e manufactureiro importante. Desenvolveu-se ainda mais a partir da \textbf{Revolução Industrial} (MISC) . Foi seriamente danificada durante os bombardeamentos alemães na \textbf{Segunda Guerra Mundial} (MISC) . A cidade é também é lar \textbf{Aston Villa FC} (ORG) , que dispurará a temporada a \textbf{Premier League} (MISC) 2011-12 e campeão da \textbf{UEFA} (ORG) \textbf{Champions League} (MISC) 1981-1982 , e o \textbf{Birmingham City} (LOC) , que foi rebaixado da \textbf{Premier League} (MISC) 2010-11 e disputará a \textbf{Football League Championship} (ORG) 2011-12 (segunda divisão). A cidade foi berço de algumas das bandas de rock mais influentes do heavy metal , como \textbf{Black Sabbath} (ORG) e \textbf{Judas Priest} (ORG) , e também de bandas de rock como \textbf{Duran Duran} (ORG) e \textbf{The Moody Blues} (ORG) .

O festival em si realizou-se no \textbf{National Indoor Arena} (LOC) uma arena multi-uso situada que, na época de sua inauguração, em 1991, era a maior de todo o \textbf{Reino Unido} (LOC) e, atualmente, comporta 15 800 assentos.



\textbf{Formato Este} (PER) ano foi notável por várias razões: término da orquestra, televoto em massa e suspense na votação final.

A \textbf{BBC} (ORG) esperou 16 anos para voltar a organizar o grande certame europeu. Após a vitória de \textbf{Katrina and The Waves} (MISC), foi imediatamente decidido que o evento não teria a mesma organização que as edições anteriores produzidas no \textbf{Reino Unido} (LOC). Várias cidades competiam para receber a final. Finalmente, a escolha recaiu sobre \textbf{Birmingham} (LOC) e na \textbf{National Indoor Arena} (LOC) , um estádio coberto de 4500 lugares que foi inaugurado em 1991.

\textbf{A Grécia} (LOC) , a \textbf{França} (LOC) , a \textbf{Suíça} (LOC) , a \textbf{Malta} (LOC) , o \textbf{Israel} (LOC) e a \textbf{Bélgica} (LOC) não utilizaram orquestra, cantaram com backing tracks. \textbf{França} (LOC), de facto, usou uma secção de violino da orquestra, mas como eles não trouxeram um maestro próprio, nenhum condutor foi mostrado antes da sua entrada. Por outro lado, a \textbf{Alemanha} (LOC) e a \textbf{Eslovénia} (LOC) apresentaram maestros, apesar de utilizarem faixas de apoio total e sem orquestra.

\textbf{Macedónia} (LOC) , apresentada como \textbf{Antiga República Jugoslava da Macedónia} (LOC), participou pela primeira vez. \textbf{Bélgica} (LOC) , \textbf{Finlândia} (LOC) , \textbf{Israel} (LOC) , \textbf{Roménia} (LOC) e \textbf{Eslováquia} (LOC) participaram, após a sua ruptura no concurso do ano anterior; \textbf{Áustria} (LOC) , \textbf{Bósnia e Herzegovina} (LOC) , \textbf{Dinamarca} (LOC) , \textbf{Rússia} (LOC) e \textbf{Islândia} (LOC) não participaram por causa da sua baixa pontuação média dos últimos cinco anos. A emissora italiana, \textbf{RAI} (ORG) , decidiu retirar-se do concurso, um movimento que iria ver a \textbf{Itália} (LOC) ausente da competição durante 13 anos, antes do seu retorno em 2011.

Depois dos resultados finais, ficou claro que \textbf{Israel} (LOC) , \textbf{Malta} (LOC) e \textbf{Reino Unido} (LOC) estavam a lutar pelo primeiro lugar. \textbf{Israel} (LOC) e \textbf{Malta} (LOC) foram aparentemente vinculados com 166 pontos, após a penúltima votação. Tudo se decidiu com a votação \textbf{da Macedónia} (LOC), que \textbf{Israel} (LOC) foi recompensado com 8 pontos, \textbf{Reino Unido} (LOC) 10 e 12 pontos para a \textbf{Croácia} (LOC). No 50  aniversário do estabelecimento do \textbf{estado de Israel} (LOC) , \textbf{Dana International} (LOC) levou à nação a sua 3 vitória no concurso. Além disso, \textbf{Edsilia Rombley} (PER) , que ficou em quarto lugar com 150 pontos, garantiu o melhor resultado para os \textbf{Países Baixos} (LOC), desde sua vitória em 1975.

Outros participantes foram notáveis, como a \textbf{Alemanha} (LOC) com \textbf{Guildo Horn} (PER) , cuja comédia chocante culminou. Controverso, escolhido para representar a \textbf{Alemanha} (LOC), foi criticado pela sua falta de seriedade. O seu 7 lugar foi decepcionante para alguns alemães e foi o início de quatro anos consecutivos a acabar no top-ten.

Uma marca d'água mostrando o nome do país apareceu na tela durante as apresentações pela primeira vez. Esta inovação ainda está em vigor hoje.

As regras desta edição foram publicadas em setembro de 1997 .



\textbf{Visual} (MISC)A abertura da competição começou com um vídeo sobre \textbf{Birmingham} (LOC) . Pontos de vista do passado e do presente da cidade foram justapostos para dar um vislumbre de sua história. A câmera terminou com um tiro na \textbf{Arena Indoor} (LOC). A orquestra apareceu, assim como as trombetas dos \textbf{Guardas da Vida} (MISC) que soaram no começo da retransmissão. Um breve vídeo resumindo a primeira competição organizada pela televisão pública britânica em 1960 em \textbf{Londres} (LOC) também aprareceu. Lá apareceu \textbf{Katie Boyle} (PER) , a mítica apresentadora britânica, que se encontrava na audiência.

A orquestra, que foi utilizada pela última vez, foi dirigida por \textbf{Martin Koch} (PER) .

O palco dividia-se em três partes: à esquerda, a orquestra foi instalada num buraco encimado por três arcos. Os dois primeiros arcos viram a parte superior deles com vista para o pódio central. O terceiro arco, decorado com padrões em relevo, continuou a abranger o chamado pódio central. No centro estava o espaço para intérpretes. Foi decorada com uma forma orgânica, que lembrava a cauda de uma baleia e estava incrustada com lâmpadas elétricas. No fundo, foram instalados quatro elementos metálicos, cada um composto por duas hastes e uma placa conectando-os. Finalmente, à direita, estava o espaço reservado para votação. Consistia de uma passarela de metal coberta por uma tela branca. Esta ponte foi equipada com uma mesa de controle e uma parede de telas.

Os apresentadores foram \textbf{Terry Wogan} (PER) e \textbf{Ulrika Jonsson} (PER) , que falaram aos espectadores em inglês e francês .

Os cartões postais começaram com um retrato em preto e branco dos artistas. Em seguida, seguiu um vídeo de duas partes. Um apresentando um aspecto da cultura britânica reconstituído no passado e depois mostrado no presente. Cada vídeo terminou com uma aparição surpresa na paisagem, a bandeira do país participante.

O intervalo foi intitulado de \textbf{Jupiter} (PER), \textbf{The Bringer of Joviality} (MISC) . Foi um potpourri cantado e dançado, destacando o multiculturalismo do \textbf{Reino Unido} (LOC) . Peças inspiradas nas culturas inglesa, escocesa, galesa, irlandesa, indiana e zulu foram tocadas. Lá apareceu em particular a violinista \textbf{Vanessa-Mae} (LOC) e a soprano \textbf{Lesley Garrett} (PER) .



Votação

21 dos 25 países usaram o televoto, onde, às 10 canções mais votadas, eram atriuídos 12, 10, 8, 7, 6, 5, 4, 3, 2, 1 ponto, com um júri de salvaguarda, em caso de erros. Um júri nacional era usado em casos de força maior, em que não se pudesse utilizar o televoto. Nos restantes países, foi usado o júri.

A supervisora executiva da \textbf{EBU} (MISC) foi \textbf{Christine Marchal-Ortiz} (PER) .

Durante a votação, a câmera fez vários close-ups dos artistas. Em particular, \textbf{Dana International} (LOC) , \textbf{Imaani} (LOC) e \textbf{Chiara} (LOC) apareceram.

Pela primeira vez em várias décadas, apareceu um mapa da \textbf{Europa} (LOC) ao longo da votação. Esta tradição manter-se-ia até aos dias de hoje.



Prémio \textbf{Barbara Dex} (PER)

Pela segunda vez, o fansite \textbf{House of Eurovision} (MISC) apresentou o \textbf{Prêmio Barbara Dex} (MISC), um prémio de humor dado ao artista mais mal vestido a cada ano no concurso. É nomeado após a artista belga, \textbf{Barbara Dex} (PER) , que ficou em último na edição de 1993, em que ela usava o seu próprio vestido projetado. A \textbf{House of Eurovision} (MISC) continuaria a fornecer o \textbf{Barbara Dex Award} (ORG) até 2016, quando outro fansite da \textbf{Eurovision} (ORG), songfestival.be, assumiu as rédeas do prémio e o apresenta todos os anos a partir de 2017.

\textbf{Guildo Horn} (PER) da \textbf{Alemanha} (LOC) ganhou o \textbf{Prémio Barbara Dex} (MISC) 1998 .



\textbf{Festival} (PER)A ordem de votação no \textbf{Festival Eurovisão da Canção 1998} (MISC), foi a seguinte:

Os países que receberam 12 pontos foram os seguintes:

Em baixo encontra-se a lista de maestros que conduziram a orquestra, na respectiva actuação de cada país concorrente.



Artistas repetentes

Alguns artistas repetiram a sua experiência \textbf{Eurovisiva} (PER). Em 1998, os repetentes foram:



\textbf{Festival Eurovisão da Canção 1999} (MISC)



\textbf{Países Participantes Países} (LOC) rebaixados não podem participar Países que participaram no passado, mas não em 1999

O \textbf{Festival Eurovisão da Canção de 1999} (MISC) (em inglês : \textbf{Eurovision Song Contest} (ORG) 1999 , em francês : \textbf{Concours Eurovision} (MISC) de la chanson 1999 e em hebraico :  1999 ) foi o 44 \textbf{Festival Eurovisão da Canção} (MISC) e realizou-se em 29 de maio de 1999 em \textbf{Jerusalém} (LOC) . O local escolhido para o concurso foi o \textbf{Centro Internacional de Convenções} (ORG) . Os anfitriões do espectáculo foram \textbf{Yigal Ravid} (PER) , repórter de televisão, \textbf{Dafna Dekel} (PER) , cantora que representou \textbf{Israel} (LOC) em 1992, e \textbf{Sigal Shahamon} (PER) , modelo e actriz. A vencedora do concurso foi \textbf{Charlotte Nilsson} (PER) , que representou a \textbf{Suécia} (LOC) com " \textbf{Take Me to Your Heaven} (MISC) ", que obteve 163 pontos. Esta foi a quarta vitória da \textbf{Suécia} (LOC) no concurso e a segunda na década de 1990 (após a vitória de \textbf{Carola} (PER) em 1991).



\textbf{Local Houve} (MISC) certa controvérsia na hora de designar qual seria a sede da edição de 1999, devido à complicada situação geopolítica de \textbf{Israel} (LOC) e também á oposição de grupos de judeus ultra ortodoxos, se diz que a edição de 1999 deveria-se nalguma cidade do \textbf{Reino Unido} (LOC) ou \textbf{Malta} (LOC) (respectivamente, o segundo e o terceiro lugar). Pelo que finalmente a cidade escolhida foi \textbf{Jerusalém} (LOC), e o \textbf{Centro Internacional de Convenções} (ORG) como a sede do festival.

Então, o \textbf{Festival Eurovisão da Canção} (MISC) 1999 ocorreu em \textbf{Jerusalém} (LOC) , em \textbf{Israel} (LOC) . \textbf{Jerusalém} (LOC) (em hebraico :  , \textbf{Yerushaláyim} (PER) ; em árabe :  , \textbf{al-Quds} (LOC) ; em grego : \textbf{} (LOC) , \textbf{Ierossólyma} (LOC) ), localizada em um planalto nas montanhas da \textbf{Judeia} (LOC) entre o \textbf{Mediterrâneo} (LOC) e o \textbf{mar Morto} (LOC) , é uma das cidades mais antigas do mundo . É considerada sagrada pelas três principais religiões abraâmicas -- judaísmo , cristianismo e islamismo . Israelenses e palestinos reivindicam a cidade como sua capital , mas \textbf{Israel} (LOC) mantém suas principais instituições governamentais em \textbf{Jerusalém} (LOC), enquanto o \textbf{Estado da Palestina} (LOC) , em última instância, apenas a prevê como a sua futura sede política; nenhuma das reivindicações, no entanto, é amplamente reconhecida pela \textbf{comunidade internacional} (ORG). contém um número de significativos lugares antigos cristãos, e é considerada a terceira cidade santa no \textbf{Islão} (MISC) . Apesar de possuir uma área de apenas 0,9 km 2 , a cidade antiga hospeda os principais pontos religiosos, entre eles a \textbf{Esplanada das Mesquitas} (LOC) , o \textbf{Muro das lamentações} (LOC) , o \textbf{Santo Sepulcro} (ORG) , a \textbf{Cúpula da Rocha} (LOC) e a \textbf{Mesquita de Al-Aqsa} (LOC) . A cidade antigamente murada, um patrimônio mundial, tem sido tradicionalmente dividida em quatro quarteirões, ainda que os nomes usados hoje (os bairros armênio , cristão , judeu e o muçulmano ) foram introduzidos por volta do século XIX. a \textbf{Cidade Velha} (LOC) foi indicada para inclusão na lista do patrimônio mundial em perigo pela \textbf{Jordânia} (LOC) em 1982.

O festival em si realizou-se no \textbf{Centro Internacional de Convenções} (ORG) , uma sala de concerto e centro de convenções , que é o maior do \textbf{Médio Oriente} (LOC) .



\textbf{Formato Regras} (PER) já antigas da \textbf{Eurovisão} (LOC) foram abolidas nesta edição, tais como: a regra que cada país tinha que cantar numa das suas línguas nacionais foi abolida pela primeira vez desde 1977. A maioria dos países participantes, quatorze dos vinte e três participantes, escolheu cantar, total ou parcialmente, em inglês e apenas oito inteiramente em suas respectivas línguas nacionais; \textbf{Lituânia} (LOC) , \textbf{Espanha} (LOC) , \textbf{Croácia} (LOC) , \textbf{Polónia} (LOC) , \textbf{França} (LOC) , \textbf{Chipre} (LOC) , \textbf{Portugal} (LOC) e \textbf{Turquia} (LOC) . Além disso, a música ao vivo tornou-se opcional, pela primeira vez na história da competição. A \textbf{IBA} (MISC) (estação arfintriã) aproveitou-se disso e decidiu abandonar a orquestra da \textbf{Eurovisão} (LOC), como forma de economizar dinheiro para o espetáculo. Isso significava que, pela primeira vez, todas as canções a concurso tiveram o uso de uma faixa de apoio durante suas atuações. Isso causou controvérsia para os tradicionalistas da \textbf{Eurovisão} (LOC), como o vencedor \textbf{Johnny Logan} (PER) criticando o movimento, que descreve agora o evento como um "karaoke".

Foi anunciado em 1999 que, a partir da edição de 2000 , os quatro maiores contribuintes financeiros da \textbf{União Europeia de Radiodifusão} (ORG) (\textbf{UER} (ORG)) - \textbf{Alemanha} (LOC) , \textbf{Espanha} (LOC) , \textbf{França} (LOC) e \textbf{Reino Unido} (LOC) - teriam a sua participação automática no concurso, independentemente da sua pontuações médias que tivessem ao longo dos últimos cinco anos.

A \textbf{Letónia} (LOC) tinha intenções de participar na \textbf{Eurovisão} (LOC), pela primeira vez, mas retirou-se numa fase tardia. Isso deu a \textbf{Hungria} (LOC) a oportunidade de participar; no entanto, a \textbf{Magyar Televízió} (LOC) decidiu não participar. Isso permitiu a \textbf{Portugal} (LOC) ser o 23 país a concurso nesse ano.

\textbf{Áustria} (LOC) , \textbf{Bósnia e Herzegovina} (LOC) , \textbf{Dinamarca} (LOC) e \textbf{Islândia} (LOC) voltaram ao concurso, depois de ter sido relegado de competir em 1998. A \textbf{Lituânia} (LOC) também retornou ao \textbf{Festival} (LOC), pela primeira vez em cinco anos. A delegação lituana tinha tido problemas de orçamento para enfrentar, e por isso a \textbf{UER} (ORG) permitiu que os lituanos chegassem a \textbf{Israel} (LOC) um dia mais tarde do que todos os outros. A primeira delegação, por outro lado a deslocar-se á \textbf{Terra Santa} (LOC) foi a da \textbf{Estónia} (LOC) .

Depois de ter sido relegado na edição de 1998, o \textbf{Perviy Kanal} (LOC) (televisão russa ) tinha decidido não transmitir a edição desse ano, a fim de permitir um retorno forte em \textbf{Israel} (LOC). No entanto, como apenas os países que tinham transmitidos a edição do ano anterior foram autorizados a entrar concurso do próximo ano, a \textbf{Rússia} (LOC) foi forçada a perder mais um ano. A eles se juntaram a \textbf{Finlândia} (LOC) , a \textbf{Grécia} (LOC) , a \textbf{Macedónia} (LOC) , a \textbf{Roménia} (LOC) , a \textbf{Eslováquia} (LOC) e a \textbf{Suíça} (LOC) ; os países com as menores pontuações médias ao longo dos cinco anos anteriores.

Os favoritos á vitória eram \textbf{Selma da Islândia} (PER) com " \textbf{All Out of Luck} (MISC) ", e \textbf{Marlain do Chipre} (PER) com " Tha 'Ne Erotas ", depois de uma pesquisa de internet por fãs. Mas, enquanto a \textbf{Islândia} (LOC) terminou em segundo lugar (a melhor participação do país no concurso), o \textbf{Chipre} (LOC) não conseguiu inspirar televotos, terminando em penúltimo, com apenas dois pontos, ambos do \textbf{Reino Unido} (LOC).

Na preparação para o concurso, muitos especularam que o evento não seria realizado em \textbf{Israel} (LOC), mas sim em \textbf{Malta} (LOC) ou no \textbf{Reino Unido} (LOC) (os países que completaram o top 3 do concurso em 1998). Treze participantes dos 23 escolheram para cantar a língua inglesa, total ou parcial. Além disso, a música ao vivo tornou-se opcional, pela primeira vez na história do concurso. A \textbf{IBA} (MISC) aproveitou disso e decidiu abandonar a orquestra do \textbf{Concurso} (MISC), como forma de economizar dinheiro para o espectáculo.

As regras desta edição foram publicadas em setembro de 1998 .

Os cartões postais desse ano foram filmes pré-gravados que seriam mostrados entre as canções. Subordinados ao tema da \textbf{Torá} (MISC) , dos quais muito poucas peças conhecidas são retratados na \textbf{Bíblia} (MISC) , sobretudo no \textbf{Antigo Testamento} (MISC) com base em cenas pintadas, um fragmento da \textbf{Torá} (MISC) desse ano, e ás figuras da animação que ganhavam vida seguiu-se um enredo bem-humorado. As imagens misturavam-se com filmes reais em que foram mostrados os pontos turísticos de \textbf{Israel} (LOC). As cenas normalmente não tinha qualquer ligação com o país participante que ocorreria após o cartão postal. As histórias a seguir foram abordadas (indicando a versão do \textbf{Antigo Testamento} (MISC) ):

Quatro das pinturas eram imagens da \textbf{Bíblia} (MISC) ilustrada por \textbf{Gustave Doré} (PER) . Entre as outras pinturas eram obras de vários artistas, incluindo \textbf{Michelangelo} (PER) , \textbf{Albrecht Dürer} (PER) e \textbf{Pieter Bruegel} (PER) , e imagens medievais anónimas.



Votação

19 dos 23 países usaram o televoto, onde, às 10 canções mais votadas, eram atriuídos 12, 10, 8, 7, 6, 5, 4, 3, 2, 1 ponto, com um júri de salvaguarda, em caso de erros. Um júri nacional era usado em casos de força maior, em que não se pudesse utilizar o televoto. Nos restantes países, foi usado o júri.

A supervisora executiva da \textbf{EBU} (MISC) foi \textbf{Christine Marchal-Ortiz} (PER) .

Durante a votação, a câmera fez vários close-ups dos artistas. Em particular, \textbf{Charlotte Nilsson} (PER) , \textbf{Selma} (LOC) , os \textbf{Sürpriz} (LOC) e \textbf{Rui Bandeira} (PER) apareceram.

Durante o processo de votação, os apresentadores separaram-se. \textbf{Yigal Ravid} (PER) tomou seu lugar na ponte e conduziu a votação, enquanto que \textbf{Dafna Dekel} (PER) e \textbf{Sigal Shahamon} (PER) se instalaram na \textbf{green room} (ORG) com os artistas e delegações.

Pela primeira vez, cada delegação foi individualmente seguida por uma câmara. A tela dividida permitia mostrá-los todos de uma vez. Os espectadores poderiam, portanto, ver todos os artistas. A produção também fez um plano sistemático de cada um recebendo "doze pontos."

Quando sua vitória tornou-se certa, a representante sueca, \textbf{Charlotte Nilsson} (PER) , se levantou e foi para beijar o representante islandês, \textbf{Selma} (LOC), que terminou em segundo.



\textbf{Prémio Barbara Dex} (MISC)Pela terceira vez, o fansite \textbf{House of Eurovision} (MISC) apresentou o \textbf{Prêmio Barbara Dex} (MISC), um prémio de humor dado ao artista mais mal vestido a cada ano no concurso. É nomeado após a artista belga, \textbf{Barbara Dex} (PER) , que ficou em último na edição de 1993, em que ela usava o seu próprio vestido projetado. A \textbf{House of Eurovision} (MISC) continuaria a fornecer o \textbf{Barbara Dex Award} (ORG) até 2016, quando outro fansite da \textbf{Eurovision} (ORG), songfestival.be, assumiu as rédeas do prémio e o apresenta todos os anos a partir de 2017.

\textbf{Lydia} (LOC) , de \textbf{Espanha} (LOC) ganhou o \textbf{Prémio Barbara Dex} (MISC) 1999 .



\textbf{Participações individuais} (MISC)Durante cerca de 10 meses, todos os países foram escolhendo os seus representantes, assim como as músicas que os mesmos interpretaram em \textbf{Jerusalém} (LOC). Para realizar tal selecção, cada país utilizou o seu próprio processo de selecção. Alguns optaram pela selecção interna, que consiste em a televisão organizadora daquele país fazer a escolha; no entanto, por vezes apenas o artista é seleccionado internamente, e a música não. Outros países (a maioria), utilizou um programa de televisão para seleccionar a sua entrada. Quartos de final, semifinais, second-chances e finais, foram realizados durante nos meses antecedentes do \textbf{Festival} (LOC) (até março) na maioria dos países europeus, cada um com o seu processo próprio. \textbf{Predefinição:Participações Individuais} (MISC) ESC1999



\textbf{Festival} (PER)A abertura da competição começou com uma animação digital, levando espectadores de \textbf{Birmingham} (LOC) (cidade anfitriã da edição de 1998) a \textbf{Jerusalém} (LOC) . A câmera viajou por todos os países europeus, representados por monumentos ou objetos emblemáticos. Um vídeo mostrando um helicóptero sobrevoando \textbf{Jerusalém} (LOC) à noite segue. A câmera pousou e \textbf{Dafna Dekel} (PER) e \textbf{Sigal Shahamon} (PER) saíram. Elas viajaram pelas ruas da cidade com um grupo de dançarinos e terminaram o seu percurso na esplanada em frente ao \textbf{Centro Internacional de Convenções} (ORG).

Depois de 42 anos onde a orquestra esteve sempre presente no palco de cada edição, a mesma foi abolida, sendo esta a primeira edição onde a orquestra não esteve presente, onde rendeu críticas do vencedor \textbf{Johnny Logan} (PER) , que disse que a \textbf{Eurovisão} (LOC) se tinha tornado num "\textbf{Karaoke} (MISC)".

A decoração do palco foi inspirada na chegada do novo milénio. Os seus constituintes simbolizavam a eternidade ao longo dos tempos . O palco, quadrante, estava segurado por quatro escadas forradas com tiras de luz. Dez linhas de pontos de luz percorriam o chão e deixvam um pódio de luz rectangular atravessando o palco para fazer avançar o público. Basicamente, a dois passos, foi instalado o principal elemento da decoração: um globo solar rotativo, cercado por quatro quadrantes de metal representando o zodíaco. O globo, deixou nove feixes de luz de cinco ramos cada. Estas vigas foram eles também móveis. Por trás desse elemento foi instalado um ecrã gigante. Em todos os lugares, foram colocados ou suspensos esferas e os globos, lembrando os planetas que giram ao redor do sol. As paredes e o teto do palco foram cravejados com milhares de pequenas bolhas que parecem estrelas distantes e constelações. Finalmente, para o lado direito do palco havia uma ponte que foi utilizado para o processo de votação.

Os apresentadores foram \textbf{Dafna Dekel} (PER) , \textbf{Yigal Ravid} (PER) e \textbf{Sigal Shachmon} (PER) , que falaram aos espectadores em hebraico , inglês e francês . Pela primeira vez, três pessoas apresentaram o certame europeu.

Os cartões postais consistiam em imagens bíblicas confundidas com imagens do dia a dia.

O intervalo foi fornecido por \textbf{Dana International} (PER) , que realizou um cover da música de \textbf{Stevie Wonder} (PER) "\textbf{Free} (MISC)", o que causou alguma controvérsia em \textbf{Israel} (LOC), devido à letra da canção. \textbf{Dana International} (PER) também apareceu no final do espectáculo, dando o troféu a \textbf{Nilsson} (PER). Depois de fingir que o troféu era pesado demais para o levantar, caiu em palco, derrubando os compositores. O show terminou com os três apresentadores convidando todos ao palco para cantar uma versão da versão em inglês de " \textbf{Hallelujah} (LOC) ", a canção vencedora de 1979 por \textbf{Israel} (LOC) .

A ordem de votação no \textbf{Festival Eurovisão da Canção 1999} (MISC), foi a seguinte:

Os países que receberam 5 pontos foram os seguintes:



Artistas repetentes

Alguns artistas repetiram a sua experiência \textbf{Eurovisiva} (PER). Em 1998, os repetentes foram:



Festival Eurovisão da Canção 2000



\textbf{Países Participantes} (LOC) Países rebaixados não podem participar Países que participaram no passado, mas não em 2000

O \textbf{Festival Eurovisão da Canção 2000} (MISC) (em inglês : \textbf{Eurovision Song Contest} (ORG) 2000 , em francês : \textbf{Concours Eurovision} (MISC) de la chanson 2000 e em sueco : \textbf{Eurovisionens Melodi Festival 2000} (MISC) ) foi o 45 \textbf{Festival Eurovisão da Canção} (MISC) e realizou-se em 13 de maio de 2000 em \textbf{Estocolmo} (LOC) , capital da \textbf{Suécia} (LOC) . Os apresentadores foram \textbf{Anders Lundin} (PER) e \textbf{Kattis Ahlström} (PER) . Os \textbf{Olsen Brothers} (ORG) foram os vencedores deste \textbf{Festival Eurovisão da Canção} (MISC) interpretando a canção em inglês " \textbf{Fly on the Wings of Love} (MISC) " ( \textbf{Voar} (MISC) nas asas do amor ) representando a \textbf{Dinamarca} (LOC) . No dia da sua vitória, \textbf{Jørgen Olsen} (PER) tinha 50 anos e 61 dias de idade, fazendo dele o artista mais velho a vencer o concurso. No entanto, ele só manteve o recorde durante um ano, quando \textbf{Dave Benton} (PER) triunfou em 2001, com a idade de 50 anos e 101 dias. \textbf{Portugal} (LOC) não pôde devido a não ter tido um bom resultado no ano anterior (mesmo assim, foi realizado o \textbf{Festival RTP da Canção} (MISC) desse ano ).



\textbf{Local} (MISC)O \textbf{Festival Eurovisão da Canção 2000} (MISC) ocorreu em \textbf{Estocolmo} (LOC) , na \textbf{Suécia} (LOC) . \textbf{Estocolmo} (LOC) é a capital e maior cidade da \textbf{Suécia} (LOC) . É a sede do governo sueco, representado na figura do \textbf{Riksdagen} (LOC) , o parlamento nacional do país, além de ser a residência oficial dos membros da monarquia sueca. \textbf{Estocolmo} (LOC) é o maior e mais importante centro urbano , cultural , político , financeiro , comercial e administrativo da \textbf{Suécia} (LOC) desde o século XIII . Sua localização estratégica sobre 14 ilhas no centro-sul da costa leste da \textbf{Suécia} (LOC), ao longo do \textbf{Lago Malar} (LOC) , tem sido historicamente importante. Uma vez que a capital sueca está situada sobre ilhas conhecidas por sua beleza, a cidade é destino de turistas de todo o mundo, tendo sido apelidada nos últimos anos de " \textbf{Veneza do Norte} (LOC)" . \textbf{Estocolmo} (LOC) é conhecida pelos seus edifícios e monumentos extremamente bem preservados, por seus arborizados parques , por sua riquíssima vida cultural e gastronômica, e pela gigantesca qualidade de vida que oferece a seus moradores. Há décadas, \textbf{Estocolmo} (LOC) figura como uma das cidades mais visitadas dos países nórdicos , com mais de um milhão de turistas internacionais anualmente. Nos últimos anos, tem sido citada entre as cidades mais "habitáveis" do mundo, sendo uma das mais limpas , organizadas e seguras do mundo.

O festival em si realizou-se no \textbf{Ericsson Globen} (MISC) , na época, o maior local escolhido para sediar a competição, com uma capacidade de 16 000 espectadores.



\textbf{Formato} (MISC)A SVT anunciou a 7 de julho de 1999 que o concurso seria organizado na \textbf{Ericsson Globen} (MISC) , em \textbf{Estocolmo} (LOC). Outros possíveis candidatos teriam sido \textbf{Scandinavium de Gotemburgo} (LOC) e \textbf{Malmömässan} (LOC) em \textbf{Malmö} (LOC)., que já tinham sediado o \textbf{Festival Eurovisão da Canção} (MISC) em 1985 e 1992, respectivamente. O \textbf{Globe} (MISC) foi escolhido porque \textbf{Estocolmo} (LOC) não hospedava o concurso desde 1975 e que seria um pouco mais barato do que as outras opções.

As regras desta edição foram publicadas a 23 de setembro de 1999 .

O favorito este ano era a \textbf{Estónia} (LOC) , que também era o país favorito dos fãs. No entanto, a \textbf{Dinamarca} (LOC) tomou o controlo das votações, vencendo à \textbf{Rússia} (LOC) que conseguiu um honroso segundo lugar e a \textbf{Letónia} (LOC) um terceiro.

A \textbf{Eslováquia} (LOC) , a \textbf{Grécia} (LOC) e a \textbf{Hungria} (LOC) decidiram não competir por razões financeiras. Os cinco países de resultados com a média mais baixa dos últimos cinco concursos que haviam participado ( \textbf{Bósnia e Herzegovina} (LOC) , \textbf{Lituânia} (LOC) , \textbf{Polónia} (LOC) , \textbf{Portugal} (LOC) e \textbf{Eslovénia} (LOC) ) foram excluídos, o que significaria que cinco países poderiam voltar. Estes países foram: a \textbf{Finlândia} (LOC) , a \textbf{Macedónia} (LOC) , a \textbf{Roménia} (LOC) , a \textbf{Rússia} (LOC) e a \textbf{Suíça} (LOC) . A \textbf{Letónia} (LOC) também se juntou à competição como o único país estreante.

Pela primeira vez, uma compilação de \textbf{CD} (MISC) oficial foi lançada. \textbf{Continha} (PER) todas as músicas das nações participantes e estava disponível em toda a \textbf{Europa} (LOC), tal como foi tentado no ano anterior.

\textbf{Israel} (LOC) , que abriu a competição, entrou com um grupo que agitavam as bandeiras israelenses e sírias defendendo a paz entre as duas nações. Esta foi considerada uma das piores músicas de sempre da \textbf{Eurovisão} (LOC).

Nos \textbf{Países Baixos} (LOC) , a \textbf{NOS} (ORG) decidiu tomar conta a meio da transmissão do concurso, por causa do desastre de foguetes em \textbf{Enschede} (LOC). Mais tarde, a \textbf{NOS} (ORG) declarou que era tanto por razões práticas, bem como porque achou inapropriado transmitir um programa de entretenimento, na noite de um evento tão catastrófico". Como resultado, o televoto teve que ser suspenso e os votos holandeses foram dados por um júri.

A primeira edição do terceiro milénio do festival é considerada um ponto sem retorno, já que foi apresentado como um festival modernizado da \textbf{Eurovisão} (LOC) com um design adaptado a uma grande arena e mudança de temas, estilos de design na produção televisiva, com o objetivo de atrair para um público mais jovem.

O concurso também foi transmitido no \textbf{Canadá} (LOC) , na \textbf{Austrália} (LOC) , no \textbf{Japão} (LOC) , nos \textbf{Estados Unidos} (LOC) e através da \textbf{Internet} (MISC), pela primeira vez.

O logótipo da edição de 2000 consistia num par de lábios entreabertos. Os seus projetistas a descreveram como uma boca sensual e estilizada e como a representação confusa dos conceitos de música, fala e diálogo. Este logótipo encontrou grande sucesso, ao ponto de entrar na lista três anos depois, para se tornar o logótipo oficial do concurso.



Votação

19 dos 24 países usaram o televoto, onde, às 10 canções mais votadas, eram atriuídos 12, 10, 8, 7, 6, 5, 4, 3, 2, 1 ponto, com um júri de salvaguarda, em caso de erros. Um júri nacional era usado em casos de força maior, em que não se pudesse utilizar o televoto. Nos restantes países, foi usado o júri.

A supervisora executiva da \textbf{EBU} (MISC) foi \textbf{Christine Marchal-Ortiz} (PER) .

Durante a votação, a câmera fez vários close-ups dos artistas, em particular dos \textbf{Olsen Brothers} (ORG) , \textbf{Alsou} (ORG) , \textbf{Brainstorm} (LOC) e \textbf{Ines} (LOC) .



\textbf{Participações individuais} (MISC)Durante cerca de 10 meses, todos os países foram escolhendo os seus representantes, assim como as músicas que os mesmos interpretaram em \textbf{Jerusalém} (LOC). Para realizar tal selecção, cada país utilizou o seu próprio processo de selecção. Alguns optaram pela selecção interna, que consiste em a televisão organizadora daquele país fazer a escolha; no entanto, por vezes apenas o artista é seleccionado internamente, e a música não. Outros países (a maioria), utilizou um programa de televisão para seleccionar a sua entrada. Quartos de final, semifinais, second-chances e finais, foram realizados durante nos meses antecedentes do \textbf{Festival} (LOC) (até março) na maioria dos países europeus, cada um com o seu processo próprio. \textbf{Predefinição:Participações Individuais} (MISC) ESC2000



\textbf{Festival} (PER)A abertura da competição começou com um vídeo sobre a \textbf{Suécia} (LOC) contemporânea. A boca do logótipo, sobreposta às imagens, indicava os nomes dos vinte e quatro países participantes. O vídeo terminou com uma vista aérea do \textbf{Globen} (LOC). A câmera então mostrou o interior da sala na penumbra, depois fazendo um close-up no palco. \textbf{Caroline Lundgren} (PER) , concertina da \textbf{Orquestra Sinfónica Jovenil de Estocolmo} (ORG), vestindo um traje tradicional sueco, apareceu cumprimentou a \textbf{Europa} (LOC) .

O palco consistia num pódio branco circular, afiado com retângulos brilhantes. Dois ecrãs gigantes ladeando no pódio: à esquerda, rectangular, mostrava os cartões postais, as apresentações dos artistas e o quadro de votação; um caminho reto, circular, mostrou o logotipo com as cores da bandeira nacional dos países participantes e porta-vozes nacionais. No mesmo pódio, ficavam cinco pilares brancos, móveis e rotativos. A sua frente estava equipado com uma tela vertical para imagens que projetava as apresentações. Este foi realmente a primeira vez que os artistas foram acompanhados durante as suas performances através de vídeos que ilustram o tema de cada canção. Estes cinco pilares podem ser combinadas para formar um terceiro ecrã gigante. Por fim, no topo do pódio estava pendurado uma tela circular branca e para trás, ficavam três fileiras de retângulos brilhantes.

Os apresentadores foram \textbf{Kattis Ahlström} (PER) e \textbf{Anders Lundin} (PER) , que falaram aos espectadores, maioritariamente em inglês e mais raramente em francês e sueco .

Os cartões postais consistiam em temas suecos que incorporavam cada nação em algum aspecto. Todos os cartões postais foram filmados em \textbf{Estocolmo} (LOC) , na \textbf{Suécia} (LOC) , no entanto, a única exceção foi o cartão postal para a \textbf{Suécia} (LOC) , filmado antes da \textbf{Expo 2000} (MISC) em \textbf{Hanôver} (LOC) , na \textbf{Alemanha} (LOC) .

O intervalo começou com um solo de violino, interpretado por \textbf{Caroline Lundgren} (PER) . Depois veio um vídeo intitulado \textbf{Once Upon a Time Europe Covered With Ice} (MISC) , dirigido por \textbf{Johan Söderberg} (PER) e produzido por \textbf{John Nordling} (PER) . Havia crianças e músicos, filmados em todos os países participantes, tocando a partitura de \textbf{Söderberg} (LOC) . Em paralelo, cenas da vida cotidiana foram montadas. No final do vídeo, \textbf{Caroline Lundgren} (PER) reapareceu no palco com o seu violino. Ela foi acompanhada pelos músicos de rua de \textbf{Estocolmo} (LOC), a banda \textbf{Bounce} (ORG). Juntos, eles realizaram um balé na extensão do vídeo, acompanhado pelas percussões dos \textbf{Drumcorps Strängnäs} (PER).

A ordem de votação no \textbf{Festival Eurovisão da Canção 2000} (MISC), foi a seguinte:

Os países que receberam 12 pontos foram os seguintes:

* \textbf{Croácia} (LOC) perdeu um terço da sua pontuação média devido ao facto de usarem uma faixa audio de voz na atuação, fazendo com que o país perdesse 33\% da pontuação na média.



Artistas repetentes

Alguns artistas repetiram a sua experiência \textbf{Eurovisiva} (PER). Em 2000, os repetentes foram:



Festival Eurovisão da Canção 2001



\textbf{Países Participantes Países} (LOC) rebaixados não podem participar Países que participaram no passado, mas não em 2001

O \textbf{Festival Eurovisão da Canção 2001} (MISC) (em inglês : \textbf{Eurovision Song Contest} (ORG) 2001 , em francês : \textbf{Concours Eurovision} (MISC) de la chanson 2001 e em dinamarquês : \textbf{Eurovision Melodi Grand-prix} (MISC) 2001 ) foi o 46 \textbf{Festival Eurovisão da Canção} (MISC) e teve lugar a 12 de maio de 2001 , em \textbf{Copenhaga} (LOC) na \textbf{Dinamarca} (LOC) . Os apresentadores do evento foram \textbf{Natasja Crone Back} (PER) e \textbf{Søren Pilmark} (PER) . Os vencedores deste festival foram \textbf{Tanel Padar} (PER) , \textbf{Dave Benton} (PER) e 2XL que representaram a \textbf{Estónia} (LOC) com a canção " \textbf{Everybody} (ORG) ".



\textbf{Local} (MISC)O \textbf{Festival Eurovisão da Canção} (MISC) 2001 ocorreu em \textbf{Copenhaga} (LOC) , na \textbf{Dinamarca} (LOC) . \textbf{Copenhaga} (LOC) é a capital e a maior cidade da \textbf{Dinamarca} (LOC) e dispõe de privilégios de condado . Situa-se na costa do mar Báltico em frente da \textbf{Suécia} (LOC) . A sua principal atracção turística é a escultura \textbf{Sereiazinha} (LOC) ou \textbf{Pequena Sereia} (PER) situada no porto da cidade.

O festival em si realizou-se no \textbf{Estádio Parken} (LOC) , o maior recinto da história a receber o certame, com capacidade para 35 000 espetadores, tendo sido preciso instalar uma cobertura no estádio para receber a edição. . É um estádio localizado no bairro de Østerbro em \textbf{Copenhaga} (LOC) , na \textbf{Dinamarca} (LOC) . Inaugurado em 9 de Setembro de 1992 - com o amistoso entre \textbf{Dinamarca} (LOC) e \textbf{Alemanha} (LOC) , tem capacidade para 38 065 espetadores e é casa do clube \textbf{F.C. Copenhagen} (PER) e da \textbf{Seleção Dinamarquesa de Futebol} (ORG) . É um estádio 4 Estrelas segundo a \textbf{UEFA} (ORG) desde 1993 , e recebeu as finais da \textbf{Copa da UEFA} (MISC) de 1999 - 2000 e da \textbf{Taça dos Clubes Vencedores de Taças de 1993} (MISC) - 1994 .



\textbf{Formato} (MISC)A emissora nacional dinamarquesa enfrentou alguns problemas ao organizar o concurso, como a falta de fundos e a busca de um local adequado. Finalmente, o local escolhido foi o \textbf{Estádio Parken} (LOC) , depois que a empresa que administra o estádio concordou em adicionar um teto retrátil ao edifício. Esta solução tornou-se o maior local de sempre a sediar um \textbf{Festival Eurovisão da Canção} (MISC), mas a escala não foi inteiramente um sucesso: muitas das 35 000 pessoas na plateia não puderam ver o palco, e para muitas participações o recinto parecia ser muito grande e a acústica era muito difícil.

As regras desta edição foram publicadas a 5 de outubro de 2000 .

\textbf{Mudanças} (MISC) ocorreram no processo de qualificação para a edição do ano seguinte: juntamente com os países "\textbf{Big 4} (MISC)", os 15 melhores colocados qualificariam-se para a competição do próximo ano. As outras vagas para 2002 seriam preenchidas por países que foram excluídos do concurso de 2001 por causa de sua baixa média de pontos para os anos 1996-2000.

No logótipo oficial, quatro círculos formavam um coração, e os quatro círculos estavam presentes no palco como quatro "círculos" de luzes.

Os favoritos neste concurso eram \textbf{França} (LOC) , \textbf{Grécia} (LOC) e \textbf{Dinamarca} (LOC) . No entanto, \textbf{Tanel Padar} (PER) e \textbf{Dave Benton} (PER) tornaram-se os vencedores inesperados do \textbf{Festival} (LOC), representando a \textbf{Estónia} (LOC) , sendo a primeira vez que um país da antiga \textbf{União Soviética} (ORG) ganhou a \textbf{Eurovisão} (LOC). Aliás, \textbf{Dave Benton} (PER) , nascido e criado em \textbf{Aruba} (LOC) , foi o primeiro negro e, aos 50 anos e 101 dias, o concorrente mais velho a vencer o concurso.

A canção sueca, " \textbf{Listen To Your Heartbeat} (MISC) " tornou-se parte de uma polémica depois ter sido acusada de ser um plágio da canção belga do \textbf{Festival Eurovisão da Canção 1996} (MISC) " \textbf{Liefde is een kaartspel} (MISC) ". No início isso foi negado pelos compositores suecos, \textbf{Thomas G} (PER): son e \textbf{Sethsson Henrik} (PER), mas depois de os compositores belgas os ameaçarem com um processo em tribunal, viram-se obrigados a lhes darem dinheiro. Tendo assim sido, a primeira canção do Festival Eurovisão da Canção declarada como plágio .



\textbf{Votação} (MISC)

17 dos 23 países usaram o televoto, onde, às 10 canções mais votadas, eram atriuídos 12, 10, 8, 7, 6, 5, 4, 3, 2, 1 ponto, com um júri de salvaguarda, em caso de erros. Um júri nacional era usado em casos de força maior, em que não se pudesse utilizar o televoto. Em 3 países, foi usado o júri e nos outros 3 países foi usado o método 50/50, 50\% televoto e 50\% júri.

A supervisora executiva da \textbf{EBU} (MISC) foi \textbf{Christine Marchal-Ortiz} (PER) .

Durante a votação, a câmera fez vários close-ups dos artistas. Em particular, \textbf{Tanel Padar} (PER) , \textbf{Dave Benton} (PER) \& 2XL , os \textbf{Rollo \& King} (ORG) , os \textbf{Antique} (LOC) , \textbf{Natasha St-Pier} (PER) e os \textbf{MTM} (ORG) apareceram.



\textbf{Participantes} (MISC)Durante cerca de 10 meses, todos os países foram escolhendo os seus representantes, assim como as músicas que os mesmos interpretaram em \textbf{Jerusalém} (LOC). Para realizar tal selecção, cada país utilizou o seu próprio processo de selecção. Alguns optaram pela selecção interna, que consiste em a televisão organizadora daquele país fazer a escolha; no entanto, por vezes apenas o artista é seleccionado internamente, e a música não. Outros países (a maioria), utilizou um programa de televisão para seleccionar a sua entrada. Quartos de final, semifinais, second-chances e finais, foram realizados durante nos meses antecedentes do \textbf{Festival} (LOC) (até março) na maioria dos países europeus, cada um com o seu processo próprio.



\textbf{Festival} (PER)A abertura da competição começou com a música vencedora do ano anterior. Os \textbf{Olsen Brothers} (ORG) entraram em cena, cantando " \textbf{Fly on the Wings of Love} (MISC) ". Eles imediatamente seguiram com seu novo single, \textbf{Walk Right Back} (MISC). Foram acompanhados por quatro cantores tocando gaita de foles.

O palco consistia num pódio retangular. Acima, foram suspensas uma reprodução luminosa e articulada do logótipo, bem como muitos spots. A decoração incluída na sua parte inferior, uma escada de cor branca. No centro da escada havia uma porta móvel que foi levantada no início da transmissão para os apresentadores. A parte superior consistia em pontos de suporte de andaimes, quadrados de luz e uma parede de telas. Finalmente, em ambos os lados do pódio, havia duas telas gigantes.

Os apresentadores foram \textbf{Natasja Crone Back} (PER) e \textbf{Søren Pilmark} (PER) , que falaram aos espectadores, maioritariamente em inglês e mais raramente em francês e dinamarquês .

Os cartões postais consistiam em vídeos curtos, cada um apresentando um aspecto da vida cotidiana na \textbf{Dinamarca} (LOC). Em cada um, havia um elemento que lembrava o país participante. No cartão postal final, todas as pessoas que apareceram nos vídeos voltavam para saudar os espectadores.

O intervalo foi preenchido pelo grupo dinamarquês \textbf{Aqua} (PER) , que começou com uma edição de videoclipes do grupo. Os membros, em seguida, apareceram no palco e cantaram um medley dos seus grandes sucessos: " \textbf{Around the World} (MISC) ", " \textbf{Barbie Girl} (MISC) ", " \textbf{Turn Back Time} (MISC) ", " \textbf{My Oh My} (MISC) ", " \textbf{Candyman} (MISC) ", " \textbf{Doctor Jones} (MISC) " e " \textbf{Cartoon Heroes} (MISC) ". Para esta última canção, eles se juntaram no palco pelo grupo \textbf{Safri Duo} (PER) , que cantaramum trecho de " Played-A-Live ".

A ordem de votação no \textbf{Festival Eurovisão da Canção 2001} (MISC), foi a seguinte:

Os países que receberam 12 pontos foram os seguintes:

* \textbf{Croácia} (LOC) perdeu um terço da sua pontuação média devido ao facto de usarem uma faixa audio de voz na atuação, fazendo com que o país perdesse 33\% da pontuação na média.



Festival Eurovisão da Canção 2002



\textbf{Países Participantes Países} (LOC) rebaixados não podem participar Países que participaram no passado, mas não em 2002

O \textbf{Festival Eurovisão da Canção 2002} (MISC) (em inglês : \textbf{Eurovision Song Contest} (ORG) 2002 , em francês : \textbf{Concours Eurovision} (MISC) de la chanson 2002 e em estónio : \textbf{Eurovisiooni} (PER) lauluvõistlus 2002 ) foi o 47 \textbf{Festival Eurovisão da Canção} (MISC) e teve lugar a 25 de maio de 2002 no \textbf{Saku Suurhall de Tallinn} (LOC) na \textbf{Estónia} (LOC) , tendo sido realizado pela televisão nacional da \textbf{Estónia} (LOC) ( \textbf{ETV} (MISC) ). Os apresentadores foram \textbf{Annelli Peebo} (PER) e \textbf{Marko Matvere} (PER) . O país vencedor foi a \textbf{Letónia} (LOC) , tendo vencido a cantora \textbf{Marija Naumova} (PER) ou \textbf{Marie N} (MISC) com a canção "\textbf{I Wanna} (MISC)". Participaram neste festival 24 países.



\textbf{Local} (MISC)O \textbf{Festival Eurovisão da Canção 2002} (MISC) ocorreu em \textbf{Tallin} (MISC) , na \textbf{Estónia} (LOC) . \textbf{Tallin} (LOC) é a cidade capital da \textbf{Estónia} (LOC) , localizada no \textbf{golfo da Finlândia} (LOC) , na costa norte do país junto ao \textbf{mar Báltico} (LOC) , a 80 quilómetros a sul de \textbf{Helsínquia} (LOC) . Tem cerca de 400 000 habitantes, aproximadamente um terço da população total do país. \textbf{Tallinn} (LOC) é a mais antiga capital da \textbf{Europa Setentrional} (LOC) . Era denominada \textbf{Reval do século XIII} (MISC) a 1917 e novamente entre 1941 e 1944 (invasão da \textbf{Alemanha Nazista} (ORG) ).

O festival em si realizou-se no \textbf{Saku Suurhall} (LOC), uma arena multiusos inaugurada em novembro de 2001 , a maior sala polivalente do país.



\textbf{Formato} (PER)Este foi o primeiro festival realizado num país da \textbf{Europa} (LOC) Oriental .

As regras desta edição foram publicadas a outubro de 2001 .

Em 2002, a \textbf{União Europeia de Radiodifusão} (ORG) retomou o sistema simples usado em 1994 e 1995 para determinar quais países participariam no concurso. Os países com a pontuação mais baixa (sem contar com a \textbf{França} (LOC) , a \textbf{Alemanha} (LOC) , a \textbf{Espanha} (LOC) e o \textbf{Reino Unido} (LOC) , que formam o \textbf{Big Four} (ORG)) do ano anterior foram excluídos, omitindo o número de países necessários para reduzir o número de participantes para 22. Portanto, os primeiros 15 classificados da edição de 2001 e os sete países relegados em 2001 foram classificados, mas a \textbf{União Europeia de Radiodifusão} (ORG) (\textbf{EBU} (MISC)) aumentou o seu número de participantes para 24, dando a oportunidade a \textbf{Israel} (LOC) e \textbf{Portugal} (LOC) de participarem. Em princípio, \textbf{Portugal} (LOC) tinha um lugar, mas recusou-se a participar devido a problemas internos na \textbf{RTP} (ORG) , fazendo com que a \textbf{Letónia} (LOC) participasse, que acabaria vencendo o concurso.

Esta edição teve um público de acordo com a \textbf{EBU} (MISC) e televisão estónia de 166 milhões de telespectadores, atingindo as taxas mais altas em \textbf{Malta} (LOC) (com mais de 90\% de share), \textbf{Suécia} (LOC) e \textbf{Espanha} (LOC) .

Pela primeira vez, o festival teve um tema "\textbf{A Modern Fairytale} (MISC)" ( Um conto de fadas moderno ), que visava unificar a produção do evento em torno de um tema comum. Os "postais" ou vídeos introdutórios dos participantes, por exemplo, transferiram um conto de fadas ou uma lenda popular para a moderna \textbf{Estónia} (LOC). O logótipo que acompanhava o tema era composto por 12 pedras coloridas em forma de "e" e o palco refletia as mesmas formas, com uma passerelle, que penetrava no público.

Houve preocupações no início do processo se a emissora estónia \textbf{ETV} (MISC) seria capaz de financiar o concurso, tendo os \textbf{Países Baixos} (LOC) se oferecido para organizar o certame; no entanto, as preocupações foram suspensas quando uma combinação de atividades de angariação de fundos e o governo estónio permitiram que eles organizassem o evento.

Apesar de ser uma das favoritas à vitória, a \textbf{Dinamarca} (LOC) terminou em último com apenas 7 pontos e teve que ficar de fora do concurso do ano seguinte.

A canção "\textbf{We} (MISC) all", do grupo lituano \textbf{B'Avarija} (PER) , foi desqualificada e foi substituída por "\textbf{Happy You} (MISC)", cantada por \textbf{Aivaras} (PER) .

\textbf{Controvérsias} (PER) surgiram durante a competição, devido a comentários feitos pelos comentadores sueco e belga, os quais disseram à plateia para que não votassem na canção israelita. A música recebeu zero pontos do público sueco, mas dois dos belgas, terminando em 12 no geral. Também o trio esloveno surgiu no Festival Eurovisão da Canção empregando trajes de hospedeira de bordo de um avião, com fatos/ternos vermelhos. A sua vitória na final eslovena foi fonte de controvérsia: o júri fez uma boa avaliação e houve problemas com o televoto que tinha escolhido a cantora \textbf{Karen Stavec} (PER) e foi anulado e a canção dos \textbf{Sestre} (PER) foi declarada vencedora. Surgiu ouitro movimento que defendia que se não fosse esta a canção vencedora, a \textbf{Eslovénia} (LOC) não deveria ser admitida na \textbf{União Europeia} (ORG) por causa das sua atitude discriminatória perante a homossexualidade .Foram acompanhados por três coristas, vestidos de comissário de bordo , com fatos/ternos brancos. Esses coristas eram \textbf{Mate Brodar} (PER), \textbf{Jadranka Juras} (PER) e \textbf{Anže Langus} (PER).

Os cartões postais consistiam em vídeos e desenhos animados, transpondo um famoso conto ou lenda na \textbf{Estónia} (LOC) contemporânea.



Votação

11 dos 24 países usaram o televoto, onde, às 10 canções mais votadas, eram atriuídos 12, 10, 8, 7, 6, 5, 4, 3, 2, 1 ponto, com um júri de salvaguarda, em caso de erros. Um júri nacional era usado em casos de força maior, em que não se pudesse utilizar o televoto. Em 5 países, foi usado o júri e nos outros 8 países foi usado o método 50/50, 50\% televoto e 50\% júri.

A supervisora executiva da \textbf{EBU} (MISC) foi \textbf{Christine Marchal-Ortiz} (PER) .

Durante a votação, a câmera fez vários close-ups dos artistas. Em particular, \textbf{Marie N} (MISC) , \textbf{Ira Losco} (PER) , \textbf{Jessica Garlick} (PER) e \textbf{Sahlene} (PER) apareceram.



\textbf{Participantes} (MISC)Durante cerca de 10 meses, todos os países foram escolhendo os seus representantes, assim como as músicas que os mesmos interpretaram em \textbf{Jerusalém} (LOC). Para realizar tal selecção, cada país utilizou o seu próprio processo de selecção. Alguns optaram pela selecção interna, que consiste em a televisão organizadora daquele país fazer a escolha; no entanto, por vezes apenas o artista é seleccionado internamente, e a música não. Outros países (a maioria), utilizou um programa de televisão para seleccionar a sua entrada. Quartos de final, semifinais, second-chances e finais, foram realizados durante nos meses antecedentes do \textbf{Festival} (LOC) (até março) na maioria dos países europeus, cada um com o seu processo próprio.



\textbf{Festival} (PER)A abertura da competição começou com um vídeo, lançado com estas palavras: "Era uma vez em uma terra distante..." A \textbf{Estónia} (LOC) contemporânea foi apresentada nalgumas atrações turísticas, a última mostrando o \textbf{Saku Suurhall} (ORG). A câmera revelou então o recinto, onde \textbf{Tanel Padar} (PER) e \textbf{Dave Benton} (PER) cantaram "\textbf{Everybody} (ORG)", a música vencedora do ano anterior.

O palco consistia num pódio de vários estágios, que se estendia para o público por um avanço curvo. O conjunto consistia em sete telas em movimento. Finalmente, em ambos os lados do palco, duas telas gigantes foram colocadas.

Os apresentadores foram \textbf{Annely Peebo} (PER) e \textbf{Marko Matvere} (PER) , que falaram aos espectadores, maioritariamente em inglês e francês e mais raramente e estónio .

Os cartões postais consistiam em vídeos e desenhos animados, transpondo um famoso conto ou lenda na \textbf{Estónia} (LOC) contemporânea.

Durante o intervalo comercial, \textbf{Annely Peebo} (PER) e \textbf{Marko Matvere} (PER) fizeram um dueto intitulado "\textbf{A Little Story in the Music} (MISC)" em que eles declararam o seu amor à música. Este número foi intercalado com vídeos curtos encenando-os em situações românticas. A atuação no final do desfile foi um balé contemporâneo, cantado e dançado, evocando a história da \textbf{Estónia} (LOC) e intitulado "\textbf{Rebirth} (PER)". Coreografado por \textbf{Teet Kask} (PER) , o balé foi realizado pela companhia de dança \textbf{Runo} (ORG) e pelo coral infantil estónio de televisão pública.

A ordem de votação no \textbf{Festival Eurovisão da Canção 2002} (MISC), foi a seguinte:

Os países que receberam 12 pontos foram os seguintes:



Artistas repetentes

Alguns artistas repetiram a sua experiência \textbf{Eurovisiva} (PER). Em 2002, os repetentes foram:



\textbf{Prémios Marcel Bezençon} (PER)Pela primeira vez, os \textbf{Prémios Marcel Bezençon} (ORG) foram entregues às melhores músicas concorrentes na final. Fundada por \textbf{Christer Björkman} (PER) (participante do \textbf{Festival Eurovisão da Canção de 1992} (MISC) e o atual chefe da delegação sueca) e \textbf{Richard Herrey} (PER) (um integrante da banda \textbf{Herreys} (LOC) , que foi a vencedora do \textbf{Festival Eurovisão da Canção de 1984} (MISC) ), em honra de \textbf{Marcel Bezençon} (PER) , um empresário e jornalista suíço responsável por elaborar o \textbf{Festival Eurovisão da Canção} (MISC). Os prémios são divididos em 3 categorias; \textbf{Prémio de Imprensa} (MISC); \textbf{Prémio Artístico} (MISC); e \textbf{Prémio dos Fãs} (MISC).



Festival Eurovisão da Canção 2003



\textbf{Países Participantes} (LOC) Países rebaixados não podem participar Países que participaram no passado, mas não em 2003

O \textbf{Festival Eurovisão da Canção 2003} (MISC) (em inglês : \textbf{Eurovision Song Contest 2003} (ORG) , em francês : \textbf{Concours Eurovision} (MISC) de la chanson 2003 e em letão : Eirovīzijas dziesmu konkurss 2003 ) foi o 48 \textbf{Festival Eurovisão da Canção} (MISC) e foi realizado pela televisão da \textbf{Letónia} (LOC) , \textbf{LTV} (ORG), no \textbf{Skonto Hall de Riga} (LOC) , na \textbf{Letónia} (LOC) a 24 de maio de 2003 . Os apresentadores foram \textbf{Renars Kaupers} (ORG) e \textbf{Marija Naumova} (ORG) . Participaram 26 países. A vencedora foi \textbf{Sertab Erener} (PER) que representou a \textbf{Turquia} (LOC) com a canção " \textbf{Everyway That I Can} (MISC) ". Neste festival há a destacar a canção da \textbf{Bélgica} (LOC) " \textbf{Sanomi} (LOC) ", cantada numa língua inventada e o facto de o \textbf{Reino Unido} (LOC) pela primeira vez ter obtido 0 pontos. Muitos associaram os zero pontos não à qualidade da canção, mas à política: o primeiro-ministro \textbf{Tony Blair} (PER) (\textbf{Reino Unido} (LOC)) apoiou \textbf{George W. Bush} (PER) na invasão do \textbf{Iraque} (LOC) . Pela primeira vez \textbf{Portugal} (LOC) enviou uma canção que tinha umas partes finais em inglês . A edição de 2003 foi a última a realizar-se numa única noite. A \textbf{União Europeia de Radiodifusão} (ORG) (\textbf{EBU} (MISC)) revelou que adicionaria uma semifinal à competição, a fim de acomodar o crescente número de países interessados que desejavam participar do concurso. O concurso também marcou a primeira vez na história da competição, onde todos os participantes estavam participando pela primeira vez.



Local

Em 22 de agosto de 2002 , o organismo de televisão letão \textbf{Latvijas Televīzija} (PER) (\textbf{LTV} (ORG)) anunciou que tinha escolhido o \textbf{Skonto Hall} (LOC) , em \textbf{Riga} (LOC) , como sede do concurso de 2003.

Sendo assim, o \textbf{Festival Eurovisão da Canção 2003} (MISC) ocorreu em \textbf{Riga} (LOC) , na \textbf{Letónia} (LOC) . \textbf{Riga} (LOC) é a capital e a maior cidade da \textbf{Letónia} (LOC) . Está localizada na nordeste da \textbf{Europa} (LOC) , sendo banhada pelo \textbf{mar Báltico} (LOC) e situado-se no coração do golfo de \textbf{Riga} (LOC) , na foz do rio \textbf{Duína Ocidental} (LOC) . É a mais importante cidade da \textbf{Letónia} (LOC), principal centro político, cultural, populacional e económico do país. Segundo dados de 2013, possuí uma população de 643 615 habitantes, ou aproximadamente 1/3 da população letã no mesmo período. É a mais populosa das capitais bálticas , e a segunda maior área metropolitana da região, depois da região metropolitana de \textbf{Vilnius} (LOC) . Para fins administrativos, \textbf{Riga} (LOC) é uma cidade independente e está localizada no distrito de \textbf{Riga} (LOC) . O centro histórico de \textbf{Riga} (LOC) foi declarado \textbf{Patrimônio da Humanidade} (PER) pela \textbf{UNESCO} (ORG) , e a cidade é particularmente notável por sua arquitetura \textbf{Art Nouveau} (PER) ( \textbf{Jugendstil} (MISC) ), comparável em importância a \textbf{Viena} (LOC) , \textbf{São Petersburgo} (LOC) e \textbf{Barcelona} (LOC) .

O festival em si realizou-se no \textbf{Skonto Hall} (LOC) uma arena multiusos usado principalmente para diferentes exposições. O \textbf{Skonto Hall} (LOC) tem capacidade para 6 500 pessoas.



\textbf{Formato} (MISC)A \textbf{EBU} (MISC) divulgou as regras para a edição de 2003 em novembro de 2002 , que detalhava que 26 países participariam, tornando-se o maior número de participantes a participar do concurso até então. As regras também modificaram os critérios de elegibilidade para inscrições, alterando a data limite de divulgação de canções de 1 de janeiro de 2003 para 1 de outubro de 2002 . Houve também uma mudança na regra de desempate, que declararia vencedor o país que tivesse sido votado por maior número de países. O sorteio da ordem de atuação foi feita a 29 de novembro de 2002 , em \textbf{Riga} (LOC) , apresentado por \textbf{Marie N} (PER) e Renārs Kaupers .

No dia do concurso, a \textbf{Rússia} (LOC) era a grande favorita para vencer o concurso, juntamente com a \textbf{Espanha} (LOC) .

Uma coletânea oficial, com todas as 26 participações do concurso, foi lançada pela primeira vez com o selo \textbf{EMI} (ORG) / \textbf{CMC} (ORG) .

Os preparativos completos para o concurso de 2003 começaram a 18 de maio de 2003 , no \textbf{Skonto Hall} (LOC). Houve ensaios, entrevistas coletivas e participantes também participaram em conversas na internet. Dois ensaios foram realizados no dia 23 de maio , diante de cerca de 12 000 pessoas. Os organizadores do concurso realizaram uma conferência de imprensa. Uma das questões levantadas foi a falta de convites para a pós-festa. O ensaio final foi realizado a 24 de maio , dia do concurso. Também foi realizada uma simulação do processo de votação, na qual os apresentadores se conectaram com todos os 26 países via satélite pela primeira vez.



Votação

23 dos 26 países usaram o televoto, onde, às 10 canções mais votadas, eram atriuídos 12, 10, 8, 7, 6, 5, 4, 3, 2, 1 ponto, com um júri de salvaguarda, em caso de erros. Um júri nacional era usado em casos de força maior, em que não se pudesse utilizar o televoto. Os outros usaram o júri

Durante a votação, a câmera fez vários close-ups dos artistas. Em particular, \textbf{Sertab Erener} (PER) , os \textbf{Urban Trad} (ORG) e as \textbf{t.A.T.u} (PER). apareceram.



\textbf{Participações individuais} (MISC)Durante cerca de 10 meses, todos os países elegeram os seus representantes, assim como as músicas que os mesmos viriam a interpretar em \textbf{Riga} (LOC). Para realizar tal selecção, cada país utilizou o seu próprio processo de selecção. Alguns optaram pela selecção interna, que consiste em a televisão organizadora daquele país, é quem faz a escolha; no entanto, por vezes apenas o artista é seleccionado internamente, e a música não. Outros países (a maioria), utiliza um programa de televisão para seleccionar a sua entrada. Quartos de final, semifinais, second-chances e finais, foram realizadas durante 10 meses na maioria dos países europeus, cada um com o seu processo próprio (também a internet foi utilizada na fase de escolhas). No conjunto de artigos em baixo, é possível ler mais ao pormenor o tipo de processo que cada país utilizou, assim como os resultados e reacções.



\textbf{Festival} (PER)A abertura da competição começou com um vídeo. A câmera viajou pelo espaço até entrar na \textbf{Via Láctea} (LOC) e no sistema solar . Então apareceu um planeta de flores, simbolizando a \textbf{Letónia} (LOC) . O concurso contou com convidados especiais que comunicaram com os anfitriões via satélite: \textbf{Lys Assia} (PER) , vencedora de 1956 saudou os anfitriões e espectadores desde \textbf{Nicósia} (LOC) , \textbf{Elton John} (PER) falou com os apresentadores ao vivo do \textbf{Life Ball} (MISC) , em \textbf{Viena} (LOC) e um astronauta e um cosmonauta - \textbf{Ed Lu} (MISC) e \textbf{Yuri Malenchenko} (PER) - deram as suas saudações desde a \textbf{Estação Espacial Internacional} (MISC) .

O palco, projetado por \textbf{Ozoliņš Aigars} (PER) foi inspirado no conceito \textbf{Planeta Letónia} (LOC), que foi chamado de "mundo da música".

Os apresentadores foram \textbf{Marie N} (PER) e Renārs Kaupers , que falaram aos espectadores em inglês e francês .

Os cartões postais, dirigidos por \textbf{Ugis Brikmanis} (PER) e contaram com a participação dos artistas que competiram no concurso interagindo com as várias paisagens da \textbf{Letónia} (LOC): florestas, rios, lagos e cidades.

Durante o intervalo comercial, dois jornalistas, \textbf{Ilze Jaunalksne} (ORG) e \textbf{Dīvs Reiznīeks} (ORG) , falaram da green room. A primeira trocou algumas palavras com a representante islandesa, \textbf{Birgitta Haukdal} (PER) e o segundo, com o representante cipriota, \textbf{Stelios Constantas} (PER) . A atuação no final do desfile foi uma curta-metragem dirigido por \textbf{Anna Viduleja} (PER) que contou com uma sequência de atuações do grupo pós-folclore letão \textbf{Iļģi} (LOC), banda de \textbf{Renārs Kaupers} (LOC) , \textbf{Brainstorm} (PER) , \textbf{Marie N} (PER) e o pianista \textbf{Raimonds Pauls} (MISC) .

Uma nova funcionalidade chamada de "\textbf{Green Room} (MISC)" (sala verde, em português ), a sala onde os artistas são posicionados para ver as votações, era visível junto do público para descobrir o fundo do palco antes da votação.

A ordem de votação no \textbf{Festival Eurovisão da Canção 2003} (MISC), foi a seguinte:

Os países que receberam 12 pontos foram os seguintes:



\textbf{Prémios Marcel Bezençon} (PER)Pela segunda vez, os \textbf{Prémios Marcel Bezençon} (ORG) foram entregues às melhores músicas concorrentes na final. Fundada por \textbf{Christer Björkman} (PER) (participante do \textbf{Festival Eurovisão da Canção de 1992} (MISC) e o atual chefe da delegação sueca) e \textbf{Richard Herrey} (PER) (um integrante da banda \textbf{Herreys} (LOC) , que foi a vencedora do \textbf{Festival Eurovisão da Canção de 1984} (MISC) ), em honra de \textbf{Marcel Bezençon} (PER) , um empresário e jornalista suíço responsável por elaborar o \textbf{Festival Eurovisão da Canção} (MISC). Os prémios são divididos em 3 categorias; \textbf{Prémio de Imprensa} (MISC); \textbf{Prémio Artístico} (MISC); e \textbf{Prémio dos Fãs} (MISC).



Festival Eurovisão da Canção 2004



Países que classificaram-se para a final \textbf{Países} (MISC) que não se classificaram para a final \textbf{Países} (MISC) que participaram no passado, mas não em 2004

O \textbf{Festival Eurovisão da Canção 2004} (MISC) (em inglês : \textbf{Eurovision Song Contest 2004} (ORG) , em francês : \textbf{Concours Eurovision} (MISC) de la chanson 2004 e em turco : \textbf{Eurovision Şarkı Yarışması 2004} (MISC) ) foi o 49 \textbf{Festival Eurovisão da Canção} (MISC) e realizou-se nos dias 12 e 15 de maio de 2004 , em \textbf{Istambul} (LOC) , na \textbf{Turquia} (LOC) . Os apresentadores desta edição foram \textbf{Korhan Abay} (PER) e \textbf{Meltem Cumbul} (PER) que falaram aos espectadores em inglês , francês e turco .

A vencedora foi \textbf{Ruslana} (PER) que representou a \textbf{Ucrânia} (LOC) , com a canção " \textbf{Wild dances} (MISC) " ( Danças selvagens ). Esta foi a primeira vitória da \textbf{Ucrânia} (LOC) na competição, após apenas 1 ano de participação. Devido à expansão do concurso, este ano foi a primeira vez em que um não vencedor obteve mais de 200 pontos. Antes desse concurso, apenas as vencedoras em 1994 e 1997 haviam passado dessa marca. Neste concurso, as 3 músicas mais votadas obtiveram mais de 200 pontos.

Devido ao número elevado de países que desejavam participar no evento, foi, pela primeira vez, instituída uma semifinal. Esta foi também a primeira vez que foi instituído o famosos símbolo eurovisivo do coração no centro da palavra "\textbf{Eurovisão} (LOC)", com a bandeira em forma de coração no centro que deve ser alterada para competições futuras. O slogan do concurso de \textbf{Istambul} (LOC) era \textbf{Under The Same Sky} (MISC) , que comunicava a importância de uma \textbf{Europa} (LOC) unida e a integração turca .

Participaram 36 países. 14 países acederam diretamente à final (os 4 grandes financiadores do evento: \textbf{Reino Unido} (LOC) , \textbf{França} (LOC) , \textbf{Espanha} (LOC) e \textbf{Alemanha} (LOC) e o top 10 do ano anterior). Os restantes participaram numa semifinal, constituída por 22 países, sendo apurados apenas os 10 primeiros que se juntaram aos 14 já apurados. Participaram na final 24 países.

Foi a primeira vez que a \textbf{Turquia} (LOC) sediou o concurso, 29 anos após a estreia do país, e também foi a primeira vez desde a edição de 1998 em \textbf{Birmingham} (LOC) que não foi sediado na capital do país anfitrião. Foi o primeiro \textbf{Festival Eurovisão da Canção} (MISC) realizado num país e cidade transcontinental, num país de maioria muçulmana e num país de língua turcomana .

Um \textbf{CD} (MISC) oficial foi lançado e, pela primeira vez, todo o concurso foi lançado em DVD.

Um novo vídeo dos \textbf{ABBA} (ORG) foi mostrado na semifinal, descrevendo brevemente como os \textbf{ABBA} (ORG) começaram e qual foi a resposta da primeira gravadora abordada. Apresentava pequenos bonecos da banda a tocar trechos das suas canções (as vozes sendo as da banda) e apresentava \textbf{Rik Mayall} (PER) como o gerente da gravadora. O vídeo contou com uma participação especial dos membros originais da banda ( \textbf{Benny Andersson} (PER) , \textbf{Agnetha Fältskog} (PER) , \textbf{Anni-Frid Lyngstad} (PER) e \textbf{Björn Ulvaeus} (PER) ), bem como \textbf{Cher} (PER) . Além disso, todas as falas foram tiradas de canções do \textbf{ABBA} (ORG) . Esta parte foi cortada do DVD do Festival Eurovisão da Canção e lançado separadamente. As referências ao vídeo que antecederam a exibição também foram cortadas.

A organização do festival foi uma oportunidade para a \textbf{Turquia} (LOC) , face ao seu processo de adesão à \textbf{União Europeia} (ORG) , mostrar uma imagem de modernidade e capacidade organizativa através das duas galas transmitidas pela \textbf{Europa} (LOC) e outros países do mundo. A vitória de 2003 e a organização da edição de 2004 tiveram, segundo algumas análises, uma resposta na economia, especialmente no turismo, o que foi potencializado na organização do festival, onde foram veiculadas imagens promocionais da \textbf{Turquia} (LOC).

Houve grandes problemas com o som, principalmente na semifinal, pois o áudio estava muito alto ou muito baixo, problema que melhorou um pouco na final.

Os favoritos para vencer o festival eram principalmente \textbf{Ucrânia} (LOC) e \textbf{Chipre} (LOC) e, em menor medida, \textbf{Espanha} (LOC) , \textbf{Sérvia e Montenegro} (LOC) , \textbf{Grécia} (LOC) e a anfitriã \textbf{Turquia} (LOC) .



\textbf{Local} (PER)

Tendo ganho a edição anterior , a produção hesitou sobre a cidade onde seria realizada a competição. Depois de contemplar por um tempo \textbf{Ancara} (LOC) , a capital do país, \textbf{Istambul} (LOC) foi escolhida. Isso permitiu enfatizar o simbolismo da cidade-ponte, conectando os dois continentes ( \textbf{Europa} (LOC) e \textbf{Ásia} (LOC) ).

Assim, o \textbf{Festival Eurovisão da Canção 2004} (MISC) ocorreu em \textbf{Istambul} (LOC) , na \textbf{Turquia} (LOC) . Istambul , a antiga \textbf{Bizâncio} (LOC) e \textbf{Constantinopla} (LOC) , é a maior cidade da \textbf{Turquia} (LOC) e a quarta maior do mundo , rivalizando com \textbf{Londres} (LOC) como a mais populosa da \textbf{Europa} (LOC) , com 15 029 231 habitantes na sua área metropolitana (2017). A grande maioria da população é muçulmana , mas também há um grande número de laicos e uma ínfima minoria de cristãos e judeus . Foi a capital do \textbf{Império Romano do Oriente} (MISC) e do \textbf{Império Otomano} (LOC) até 1923, cujo governante máximo, o sultão , foi durante séculos reconhecido como califa , o chefe supremo de todos os muçulmanos, o que fazia da cidade uma das mais importantes de todo o \textbf{Islão} (MISC). Atualmente, embora a capital do país seja \textbf{Ancara} (LOC) , \textbf{Istambul} (LOC) continua a ser o principal polo industrial, comercial, cultural e universitário (aí estão sediadas mais de uma dezena de universidades ) do país. É a sede do \textbf{Patriarcado Ecumênico de Constantinopla} (LOC) , sede da \textbf{Igreja Ortodoxa} (MISC) . A cidade ocupa ambas as margens do \textbf{estreito do Bósforo} (LOC) e do norte do \textbf{mar de Mármara} (LOC) , os quais separam a \textbf{Ásia} (LOC) da \textbf{Europa} (LOC) no sentido norte-sul, uma situação que faz de \textbf{Istambul} (LOC) a única cidade que ocupa dois continentes. A parte central da parte europeia é por sua vez dividida pelo estuário do \textbf{Corno de Ouro} (LOC) . É usual dizer-se que a cidade tem dois ou três centros , conforme se considere ou não que na parte asiática também existe um centro . No lado europeu há duas zonas com mais destaque em termos de movimento de pessoas e património cultural: o mais antigo, onde se situava o núcleo da antiga \textbf{Bizâncio} (LOC) e \textbf{Constantinopla} (LOC), correspondente ao atual distrito de \textbf{Fatih} (LOC) , fica a sul do \textbf{Corno de Ouro} (LOC), enquanto que \textbf{Beyoğlu} (LOC) , a antiga \textbf{Pera} (LOC) e onde se situava o bairro europeu medieval de \textbf{Gálata} (LOC) , fica a norte. O centro da parte asiática tem contornos menos precisos, e ocupa parte dos distritos de \textbf{Üsküdar} (LOC) e \textbf{Kadıköy} (LOC) . Algumas zonas históricas da parte europeia de Istambul foram declaradas \textbf{Património Mundial} (MISC) pela \textbf{UNESCO} (ORG) em 1985. Em 2010, a cidade foi a \textbf{Capital Europeia da Cultura} (ORG) . Devido à sua dimensão e importância, \textbf{Istambul} (LOC) é considerada uma megacidade e uma cidade global .

A produção, então, escolheu o \textbf{Mydonose Showland} (LOC), um centro de entretenimento na forma de uma tenda piramidal gigante perto do \textbf{Aeroporto Internacional Atatürk} (LOC) , como local, depois mudou para a \textbf{Abdi İpekçi Arena} (PER) , uma arena multiusos com capacidade para 12 500 espectadores. Esta última, inaugurada em 1986 e sede de competições de basquetebol , foi escolhida por ser maior e mais prática para um evento dessa importância. Foi demolida em 2018 .



\textbf{Formato} (MISC)A Eurovisão deste ano foi o primeiro a ser um evento de dois dias, com uma semi-final realizada numa quarta-feira e a grande final realizada no sábado seguinte. Sob este novo formato, os "\textbf{Big 4} (MISC)" ( \textbf{França} (LOC) , \textbf{Alemanha} (LOC) , \textbf{Espanha} (LOC) e o \textbf{Reino Unido} (LOC) ) e o top 10 do concurso de 2003 tiveram classificação automática para a final.

Esta foi também a primeira vez que foi instituído o famosos símbolo eurovisivo do coração no centro da palavra "\textbf{Eurovisão} (LOC)", com a bandeira em forma de coração no centro que deve ser alterada para competições futuras. O slogan do concurso de \textbf{Istambul} (LOC) era \textbf{Under The Same Sky} (MISC) , que comunicava a importância de uma \textbf{Europa} (LOC) unida e a integração turca .

Pela primeira vez, um DVD foi lançado após a competição, contendo a semifinal completa e a final. Além disso, o sítio www.eurovision.tv foi oficialmente assumido pela \textbf{União Europeia de Radiodifusão} (ORG) , passando a ser a sua plataforma de comunicação.

Para lhe conferir uma identidade mais permanente e duradoura, a \textbf{EBU} (MISC) decidiu disponibilizar ao concurso um logótipo genérico. Este usava a palavra "Eurovisão" de forma estilizada. A letra "v" central tinha a forma de um coração, enquadrando a bandeira nacional do país anfitrião.

O slogan adotado foi \textbf{Under The Same Sky} (MISC) , servindo de inspiração para o design do palco. Isso consistia em dois níveis embutidos um no outro. O nível inferior, de formato oval, tinha uma tela embutida no solo. O nível superior, que se estendia ao público, possuía três pódios circulares pretos com tela embutida na face superior. A decoração consistia numa forma orgânica opaca pendendo do palco e dotada de uma cúpula luminosa na parte superior. O cenário era feito de telas de \textbf{LED} (ORG) .

Os cartões postais começaram com uma projeção do logótipo da competição tendo como pano de fundo o horizonte da cidade de \textbf{Istambul} (LOC). Seguiu-se um vídeo mostrando as riquezas culturais e turísticas da \textbf{Turquia} (LOC). Finalmente, \textbf{Azra Akin} (PER) , \textbf{Miss Mundo 2002} (MISC) , apareceu na tela, soprando pétalas de rosa das quais se projetava a bandeira nacional do país participante.



\textbf{Votação Todos} (MISC) os países participantes da competição, incluindo aqueles que não se classificaram para a final, puderam votar. Depois da apresentação de todas as músicas, cada país abriu o televoto para permitir que os seus espectadores votassem na sua música favorita, excepto no seu próprio país. Cada país atribuiu pontos com base no número de votos dados para cada música: a música que recebeu mais votos do televoto recebeu 12 pontos, a segunda 10 pontos, a terceira 8 pontos, e assim sucessivamente até 1 ponto. Pela primeira vez também, a contagem regressiva do tempo restante de votação foi exibida na tela.

Este foi o primeiro ano em que todos os 36 países participantes votaram com base no televoto. Nos casos em que o televoto não pode ser utilizado, um júri de reserva foi usado.No entanto, \textbf{França} (LOC) , \textbf{Polónia} (LOC) e \textbf{Rússia} (LOC) não transmitiram a semifinal (porque não participaram dela) e, portanto, não deram votos como os outros trinta e três países. Na \textbf{Bélgica} (LOC) , o \textbf{RTBF} (ORG) de língua francesa não transmitiu a semifinal, mas a VRT de língua holandesa sim.

Em caso de empate, o número de países que votaram nas canções empatadas seria contado, e a canção com mais países dando pontos a ela seria a vencedora. No caso de um novo empate, o método anteriormente usado de contagem regressiva do número de 12 pontos, 10 pontos, etc., seria usado para encontrar um eventual vencedor.

O supervisor executivo da \textbf{EBU} (MISC) foi \textbf{Svante Stockselius} (PER) .



\textbf{Participantes} (MISC)Durante cerca de 10 meses, todos os países foram escolhendo os seus representantes, assim como as músicas que os mesmos interpretaram em \textbf{Istambul} (LOC) . Para realizar tal selecção, cada país utilizou o seu próprio processo de selecção. Alguns optaram pela selecção interna, que consiste na televisão organizadora daquele país fazer a escolha; no entanto, por vezes apenas o artista é seleccionado internamente, e a música não. Outros países (a maioria), utilizou um programa de televisão para seleccionar a sua entrada. Quartos de final, semifinais, second-chances e finais, foram realizados durante nos meses antecedentes do \textbf{Festival} (LOC) (até março) na maioria dos países europeus, cada um com o seu processo próprio.

\textbf{Andorra} (LOC) , \textbf{Albânia} (LOC) , \textbf{Bielorrússia} (LOC) e \textbf{Sérvia e Montenegro} (LOC) fizeram a sua estreia, com o \textbf{Mónaco} (LOC) a regressar após uma ausência de 25 anos. \textbf{Luxemburgo} (PER) planeou regressar após uma ausência de 11 anos, mas saiu mais tarde após o surgimento de questões financeiras entre a \textbf{RTL} (MISC) e a \textbf{EBU} (MISC) .



\textbf{Festival} (PER)A abertura da competição começou com uma breve animação do logo. A câmera mostrou o público na sala. Uma breve demonstração de dervixes rodopiantes se seguiu. Os apresentadores entraram em cena para as apresentações habituais, agradecendo à \textbf{Letónia} (LOC) pela organização da competição no ano anterior. A abertura ficou concluída com a apresentadora \textbf{Meltem Cumbul} (PER) a cantar " \textbf{Nel} (ORG) blu dipinto di blu (\textbf{Volare} (PER)) " para o público.

Durante o primeiro intervalo comercial, \textbf{Meltem Cumbul} (ORG) trocou algumas palavras com os artistas na \textbf{green room} (ORG), nomeadamente \textbf{Piero Esteriore} (PER) e \textbf{David D'Or} (PER) .

Durante o segundo intervalo, \textbf{Korhan Abay} (PER) falou com \textbf{Željko Joksimović} (PER) , o grupo \textbf{Re-union} (ORG) e \textbf{Deen} (ORG) .

Após a apresentação dos 22 semifinalistas, as filas foram abertas nos países votantes. Este ano a transmissão da semifinal foi opcional para os classificados directamente para a final, o que fez com que \textbf{França} (LOC), \textbf{Polónia} (LOC) e \textbf{Rússia} (LOC) não a transmitissem e, portanto, não votassem; dando assim um total de 33 países votantes.

O intervalo foi ocupado por uma curta atuação do grupo de ballet turco \textbf{Fire of Anatolia} (MISC) . Em seguida, seguiu a exibição do \textbf{ABBA} (ORG) : \textbf{The Last Video} (MISC), para comemorar o trigésimo aniversário da vitória do grupo no concurso, em 1974 . O vídeo apresentou uma breve aparição dos membros originais da banda ( \textbf{Benny Andersson} (PER) , \textbf{Agnetha Fältskog} (PER) , \textbf{Anni-Frid Lyngstad} (PER) e \textbf{Björn Ulvaeus} (PER) ), bem como \textbf{Cher} (PER) . Além disso, todas as réplicas foram tiradas das músicas do \textbf{ABBA} (ORG) .

O intervalo termina com a apresentação dos quatorze países já qualificados para a final.

Os resultados foram revelados ao vivo por \textbf{Korhan Abay} (PER) e \textbf{Meltem Cumbul} (LOC) , que abriram aleatoriamente dez envelopes vermelhos, contendo os nomes dos dez países qualificados. Excepcionalmente, os resultados obtidos pelos países não qualificados foram exibidos na tela.

Os países que receberam 12 pontos foram os seguintes:

A final começou com um pequeno vídeo mostrando as vistas aéreas de \textbf{Istambul} (LOC) , dia e noite e concluindo com um mapa da \textbf{Abdi İpekçi Arena} (MISC) . A vencedora do ano passado, \textbf{Sertab Erener} (PER) , entrou em palco e cantou " \textbf{Everyway That I Can} (MISC) ", a música vencedora do ano anterior, e "\textbf{Leave} (LOC)".

Os apresentadores então subiram ao palco e agradeceram a \textbf{Erener} (PER). Fizeram as apresentações habituais e, como durante a semifinal, a apresentadora Meltem Cumbul cantou " \textbf{Nel} (ORG) blu dipinto di blu (\textbf{Volare} (PER)) " para o público

Durante o primeiro intervalo comercial, os apresentadores entraram contacto via satélite com \textbf{Hamburgo} (LOC) e \textbf{Thomas Anders} (LOC) . Mas apesar das tentativas de produção, não conseguiram entrar em contacto com \textbf{Istambul} (LOC) e \textbf{Las Palmas} (LOC) .

Durante o segundo intervalo, no meio do processo de votação, \textbf{Sertab Erener} (PER) apareceu na \textbf{green room} (LOC), trocando algumas palavras com os membros do grupo \textbf{Athena} (ORG) . Os apresentadores, por sua vez, puderam finalmente juntar-se a \textbf{Las Palmas} (LOC) .

O intervalo foi preenchido pelo grupo de balé \textbf{Fire of Anatolia} (MISC) . \textbf{Sertab Erener} (PER) então reapareceu na tela, trocando algumas palavras com \textbf{Ruslana} (PER) , \textbf{Max} (PER) , \textbf{Stefan Raab} (PER) e o grupo \textbf{Re-union} (ORG) .

Para a votação do primeiro festival que realizou uma semifinal, ficou decidido que a ordem de votação seria por ordem alfabética pelo \textbf{Código ISO} (ORG) , começando pela \textbf{Andorra} (LOC) e terminando na \textbf{Ucrânia} (LOC) .

O festival foi uma das edições mais renhidas, reflectindo-se nas pontuações altas dos primeiros lugares e nos poucos pontos que a parte inferior obteve. Também foi demonstrado no número de países que votaram nas canções nas primeiras posições, dando quase um monopólio a este respeito: \textbf{Sérvia e Montenegro} (LOC) juntamente com a \textbf{Grécia} (LOC) foram votados por todos (algo que não se repetiria até 2009).

A ordem de votação no \textbf{Festival Eurovisão da Canção 2004} (MISC), foi a seguinte:

Os países que receberam 12 pontos foram os seguintes:



\textbf{Factos} (MISC) e \textbf{Controvérsias} (MISC)Uma hora depois de a semifinal ter sido transmitida, a \textbf{União Europeia de Radiodifusão} (ORG) descobriu que houve problemas com a contagem dos votos do \textbf{Mónaco} (LOC) e \textbf{Croácia} (LOC) . Digame, uma afiliada da \textbf{Deutsche Telekom} (ORG) , que era responsável pelo processamento de todos os votos, relatou que encontrou problemas com o software de cálculo e houve um problema com a votação do text message na \textbf{Croácia} (LOC). Consequentemente, alguns votos não foram contabilizados nos resultados anunciados ao final da transmissão da semifinal. Quando os resultados foram corrigidos para incluir esses votos adicionais, eles não afetaram os países que se qualificaram para a final.

\textbf{Vale} (ORG) destacar também que \textbf{Sérvia e Montenegro} (LOC) terminou em 1 lugar na \textbf{Semifinal} (LOC) com 263 pontos e depois em 2 na \textbf{Final} (LOC) com 263 pontos.

Este ano também foi notável, pois foi o primeiro ano em que a \textbf{Turquia} (LOC) votou no \textbf{Chipre} (LOC) e o segundo ano consecutivo em que o \textbf{Chipre} (LOC) votou na \textbf{Turquia} (LOC). No entanto, num movimento que irritou alguns cipriotas, quando o país apresentou os seus votos nenhum mapa da ilha foi mostrado (todos os outros foram precedidos com o seu país a ser destacado no mapa). Isso ocorreu devido ao reconhecimento pela \textbf{Turquia do Chipre do Norte} (LOC) como uma república independente (não reconhecida por nenhum outro \textbf{Estado} (LOC)). É provável que a \textbf{Turquia} (LOC) tenha desistido de mostrar o mapa porque teria apenas destacado a parte sul da ilha, irritando assim a comunidade internacional.

Este também foi o primeiro ano em que os votos foram relidos pelos anfitriões em apenas um idioma. Antes de 2004, todos os pontos eram repetidos em francês e inglês , mas devido à votação de 36 países, e mais nos próximos anos, em 2004 para economizar tempo, os anfitriões apenas releram cada pontuação num idioma, o oposto do que o porta-voz do país original falou.

Nos anos seguintes, os primeiros sete pontos já seriamcolocados no quadro e apenas os últimos pontos seriam lidos para economizar tempo.

Houve problemas técnicos quando, num curto intervalo no meio das músicas, (usado para o intervalo publicitário), os anfitriões tentaram contactar várias partes da \textbf{Europa} (LOC). Tentaram entrar em contato com \textbf{Alemanha} (LOC) , \textbf{Espanha} (LOC) e \textbf{Turquia} (LOC) , mas no final só conseguiram uma resposta da \textbf{Alemanha} (LOC) . Também durante a introdução do cartão postal da \textbf{Romênia} (LOC) , a informação para a participação romena apareceu na tela, mas foi rapidamente retirada.

Pouco antes da apresentação da \textbf{Eslovênia} (LOC) , a emissora turca acidentalmente fez um intervalo comercial, o que significa que a música eslovena não foi ouvida pelos espectadores turcos e, consequentemente, a \textbf{Turquia} (LOC) não deu votos para a música.



Festival Eurovisão da Canção 2005



Países que classificaram-se para a final \textbf{Países} (MISC) que não se classificaram para a final \textbf{Países} (MISC) que participaram no passado, mas não em 2005

O \textbf{Festival Eurovisão da Canção 2005} (MISC) (em inglês : \textbf{Eurovision Song Contest} (ORG) 2005 ; em francês : \textbf{Concours Eurovision} (MISC) de la chanson 2005 ; em ucraniano : \textbf{} (PER)  \textbf{} (PER) 2005 ) foi o 50 \textbf{Festival Eurovisão da Canção} (MISC) que teve lugar a 19 de maio (semifinal) e a 21 de maio de 2005 (final). Foi organizado pela televisão estatal da \textbf{Ucrânia} (LOC) \textbf{NTU} (ORG) e teve lugar no \textbf{Palácio dos Desportos de Kiev} (LOC) . Os apresentadores do evento foram \textbf{DJ Pasha} (PER) e \textbf{Masha Efrosinina} (PER) .

A vencedora foi \textbf{Helena Paparizou} (PER) , com a canção " \textbf{My} (MISC) Number \textbf{One} (MISC) ", em representação da \textbf{Grécia} (LOC) , naquela que foi a primeira e, até hoje, a única vitória do país no certame europeu.

A \textbf{Bulgária} (LOC) e a \textbf{Moldávia} (LOC) estrearam-se no concurso, com a \textbf{Hungria} (LOC) a regressar após uma ausência de sete anos. Inicialmente, o \textbf{Líbano} (LOC) também iria fazer a sua estreia, com a canção " \textbf{Quand} (LOC) tout s'enfuit ", interpretada por \textbf{Aline Lahoud} (PER) . No entanto, o facto de haverem dúvidas sobre a transmissão de todo o concurso por parte da emissora pública libanesa - devido ao não-reconhecimento da existência do \textbf{Estado de Israel} (LOC) -, fizeram com que a emissora optasse por desistir em março de 2005 .

Participaram 39 países. 14 países acederam diretamente à final (os 4 grandes financiadores do evento: \textbf{Reino Unido} (LOC) , \textbf{França} (LOC) , \textbf{Espanha} (LOC) e \textbf{Alemanha} (LOC) e o top 10 do ano anterior). Os restantes participaram numa semifinal, constituída por 25 países, sendo apurados apenas os 10 primeiros que se juntaram aos 14 já apurados. Participaram na final 24 países. Curiosamente, os \textbf{Big 4} (MISC) acabaram todos no chamado "\textbf{Bottom 4} (MISC)", ou seja, nos quatro últimos lugares da tabela.



\textbf{Organização} (ORG)Ganhando o direito de organizar a edição de 2005 após ter ganho o \textbf{Festival Eurovisão da Canção 2004} (MISC) , a televisão pública ucraniana \textbf{NTU} (ORG) escolheu a capital \textbf{Kiev} (LOC) a 6 de setembro de 2004 . No entanto, o local escolhido, o \textbf{Palácio dos Desportos de Kiev} (LOC) , teve de passar por um processo de renovação para cumprir os requisitos designados pela \textbf{União Europeia de Radiodifusão} (ORG) , que começou em \textbf{Dezembro de 2004} (MISC) , e terminou em maio de 2005 .

A organização da competição encontrou um grande hiato: a \textbf{Revolução Laranja} (MISC) . Em novembro de 2004 , foram realizadas eleições presidenciais na \textbf{Ucrânia} (LOC) , após as quais o candidato \textbf{Viktor Yanukovych} (PER) foi declarado vencedor. A oposição liberal, liderada por \textbf{Viktor Yushchenko} (PER) , exigiu imediatamente a anulação dos resultados e novas eleições, suspeitando de inúmeras fraudes. \textbf{Kiev} (LOC) foi palco de manifestações incessantes até que a oposição ganhou o caso. As novas eleições confirmaram a vitória dos liberais e \textbf{Viktor Yushchenko} (PER) tornou-se o novo presidente ucraniano. Mas durante estes episódios, a organização da competição permaneceu paralisada. O novo regime, uma vez instalado, fez tudo o que estava ao seu alcance para recuperar o tempo perdido e dar uma imagem positiva da \textbf{Ucrânia} (LOC) na cena internacional.

Por outro lado, a \textbf{União Europeia de Radiodifusão} (ORG) introduziu uma nova regra relativa ao televoto, em que um número específico de votos telefónicos teria de ser registado a nível nacional para que o televoto do país em causa fosse validado. Abaixo deste limite, o televoto seria invalidado e os resultados determinados por um júri de substituição. A regra foi aplicada na semifinal para a \textbf{Albânia} (LOC) , \textbf{Andorra} (LOC) e \textbf{Mónaco} (LOC) , e na final para a \textbf{Andorra} (LOC) , \textbf{Moldávia} (LOC) e \textbf{Mónaco} (LOC) .

\textbf{Kiev} (LOC) é a capital e a cidade mais populosa da \textbf{Ucrânia} (LOC) . Fica no centro-norte do país, ao longo do \textbf{rio Dnieper} (LOC) . É um importante centro industrial, científico, educacional e cultural da \textbf{Europa Oriental} (LOC) . É o lar de muitas indústrias de alta tecnologia , instituições de ensino superior e marcos históricos. A cidade possui um extenso sistema de transporte público e infraestrutura, incluindo o \textbf{metro de Kiev} (LOC) .

Diz-se que o nome da cidade deriva do nome de \textbf{Kyi} (MISC), um de seus quatro fundadores lendários. Durante sua história, \textbf{Kiev} (LOC), uma das cidades mais antigas do \textbf{Leste Europeu} (LOC), passou por vários estágios de destaque e obscuridade. A cidade provavelmente existia como um centro comercial já no século V. Um assentamento eslavo na grande rota comercial entre a \textbf{Escandinávia} (ORG) e \textbf{Constantinopla} (LOC) , \textbf{Kiev} (LOC) era um afluente dos cazares , até sua captura pelos varangianos ( vikings ) em meados do século IX. Sob o domínio varangiano, a cidade tornou-se a capital da \textbf{Rússia de Kiev} (LOC) , o primeiro estado eslavo oriental . Completamente destruída durante as invasões mongóis em 1240 , a cidade perdeu a maior parte de sua influência nos séculos seguintes. Era uma capital de província de importância marginal na periferia dos territórios controlados por seus poderosos vizinhos, primeiro a \textbf{Lituânia} (LOC) , depois a \textbf{Polónia} (LOC) e finalmente a \textbf{Rússia} (LOC) .

A cidade prosperou novamente durante a \textbf{Revolução Industrial do Império Russo} (MISC) no final do século XIX. Em 1918, depois que a \textbf{República Popular da Ucrânia} (LOC) declarou independência da \textbf{Rússia Soviética} (LOC) , \textbf{Kiev} (LOC) tornou-se sua capital. A partir de 1921, \textbf{Kiev} (LOC) foi uma cidade da \textbf{Ucrânia Soviética} (LOC) , proclamada pelo \textbf{Exército Vermelho} (ORG) e, a partir de 1934, \textbf{Kiev} (LOC) tornou-se sua capital. A cidade foi quase completamente arruinada durante a \textbf{Segunda Guerra Mundial} (MISC) , mas rapidamente se recuperou nos anos do pós-guerra, permanecendo a terceira maior cidade da \textbf{União Soviética} (ORG) .

Após o colapso da \textbf{União Soviética} (ORG) e a independência da \textbf{Ucrânia} (LOC) em 1991, \textbf{Kiev} (LOC) permaneceu como capital e experimentou um fluxo constante de migrantes ucranianos de outras regiões do país. Durante a transformação do país em uma economia de mercado e democracia eleitoral , \textbf{Kiev} (LOC) continuou a ser a maior e mais rica cidade da \textbf{Ucrânia} (LOC). Sua produção industrial dependente de armamentos caiu após o colapso soviético, afetando negativamente a ciência e a tecnologia, mas novos setores da economia, como serviços e finanças, facilitaram o crescimento de salários e investimentos de \textbf{Kiev} (LOC), além de fornecer financiamento contínuo para o desenvolvimento de moradias e infraestrutura urbana. \textbf{Kiev} (LOC) emergiu como a região mais pró- ocidental da \textbf{Ucrânia} (LOC); partidos que defendem uma integração mais estreita com a \textbf{União Europeia} (ORG) dominam durante as eleições.

Os quartos de hotel eram escassos porque os organizadores do concurso pediram ao governo ucraniano que bloqueasse as reservas que não controlavam através de alocações oficiais de delegações ou pacotes turísticos. Isto levou ao cancelamento das reservas de hotel de muitas pessoas.

Os organizadores esperavam que, ao acolher a \textbf{Eurovisão} (LOC), impulsionasse a imagem da \textbf{Ucrânia} (LOC) no estrangeiro e aumentasse o turismo, enquanto o novo governo do país esperava que também desse um impulso modesto ao objetivo a longo prazo de adquirir a adesão à \textbf{União Europeia} (ORG) .



\textbf{Formato} (PER)O logótipo oficial do concurso permaneceu o mesmo da edição anterior , com a bandeira do país no coração a ser alterada, inserindo-se a da \textbf{Ucrânia} (LOC) . Depois de "\textbf{Under The Same Sky" de Istambul} (MISC), o slogan da edição de 2005 foi "\textbf{Awakening} (MISC)", que simbolizava o despertar do país e da cidade prontos para se apresentarem à \textbf{Europa} (LOC) .

Esta foi a primeira edição a ser transmitida em formato widescreen 16:9 .

Os anfitriões do Festival Eurovisão da Canção foram a apresentadora de televisão \textbf{Maria "Masha" Efrosinina e o DJ Pavlo "Pasha" Shylko} (MISC) . A vencedora do ano anterior \textbf{Ruslana} (LOC) também estava prevista apresentar, mas acabou por desistir por vários motivos, incluindo o seu fraco domínio do inglês . No entanto, atuou durante a final e entrevistou os artistas na green room. Os boxeadores ucranianos \textbf{Vitali} (LOC) e \textbf{Wladimir Klitschko} (PER) abriram o televoto, enquanto um troféu especial foi entregue ao vencedor pelo presidente da \textbf{Ucrânia} (LOC) , \textbf{Viktor Yushchenko} (PER) .

Um \textbf{CD} (MISC) e DVD oficiais foram lançados, juntamente com distintivos em forma de coração com as bandeiras de todos os trinta e nove países participantes. A \textbf{União Europeia de Radiodifusão} (ORG) também encomendou o livro "\textbf{The Eurovision Song Contest - The Official History} (MISC)" do autor britânico/americano \textbf{John Kennedy O'Connor} (PER) para comemorar o quinquagésimo aniversário do concurso. O livro foi apresentado na no intervalo da final entre as músicas da \textbf{Sérvia e Montenegro} (LOC) e \textbf{Dinamarca} (LOC) . O livro foi publicado em inglês , alemão , francês , holandês , sueco , dinamarquês e finlandês .

Durante a semifinal, houve algumas falhas sonoras, principalmente durante a música norueguesa, logo após a introdução e também durante a música irlandesa. Estes não foram corrigidos para o lançamento do DVD.

Os cartões postais (pequenos clipes exibidos entre as canções) desta edição ilustraram a cultura e a herança da \textbf{Ucrânia} (LOC) , juntamente com um lado mais moderno e industrial do país.



Países participantes

Participaram neste festival 39 países: 14 dos quais tinham acesso direto à final do dia 21 de maio , enquanto que os restantes 24 países tiveram de participar numa semifinal que apurou 10 destes.

O \textbf{Líbano} (LOC) manifestou o interesse em participar tendo mesmo escolhido uma canção e uma intérprete. Porém, ocorreu um problema político: O \textbf{Líbano} (LOC) aceitara as regras da \textbf{EBU} (MISC) mas não tencionava emitir a participação de \textbf{Israel} (LOC) dadas as divergências políticas entre estes dois países (problema palestiniano). Uma das exigências da \textbf{UER} (ORG) é que todos os países são obrigados a transmitir todas as canções, e assim o \textbf{Líbano} (LOC) desistiu para não ser desqualificado.

Os únicos países europeus que não participaram foram \textbf{Itália} (LOC) , \textbf{Eslováquia} (LOC) , \textbf{República Checa} (LOC) (apesar de ter submetido a sua candidatura provisória), \textbf{Luxemburgo} (LOC) e \textbf{São Marino} (LOC) .



\textbf{Festival} (PER)A semifinal teve lugar no dia 19 de maio de 2005 , às 21h \textbf{CET} (ORG) . Votaram todos os países participantes da edição, incluindo os que não concorreram na semifinal.

Os países que receberam 12 pontos foram os seguintes:

A final teve lugar no dia 21 de maio de 2005 , às 21h \textbf{CET} (ORG) . Votaram todos os países participantes da edição, incluindo os que não se qualificaram para a final.

Os países que receberam 12 pontos foram os seguintes:



Festival Eurovisão da Canção 2006



Países que classificaram-se para a final \textbf{Países} (MISC) que não se classificaram para a final \textbf{Países} (MISC) que participaram no passado, mas não em 2006

O \textbf{Festival Eurovisão da Canção de 2006} (MISC) foi o 51. \textbf{Festival Eurovisão da Canção} (MISC) e teve lugar nos dias 18 de Maio (\textbf{Semifinal} (LOC)) e 20 de \textbf{Maio de 2006} (MISC) (\textbf{Final} (LOC)). Foi organizado pela televisão estatal da \textbf{Grécia} (LOC) , \textbf{ERT} (LOC). A competição teve como espaço o \textbf{Olympic Indoor Hall de Atenas} (LOC) , na \textbf{Grécia} (LOC). Os apresentadores do evento foram \textbf{Sakis Rouvas} (PER) e \textbf{Maria Menounos} (PER) .

A \textbf{Finlândia} (LOC) destrona os seus concorrentes, com uma actuação ousada, (com um grupo de heavy metal e com visual de horror ) vencendo assim o 51. Festival Eurovisão da Canção, com 292 pontos, organizando assim, o \textbf{Festival Eurovisão da Canção 2007} (MISC) .



\textbf{Formato} (MISC)O Formato é o mesmo utilizado desde 2004, com uma semi-final de 23 países , de onde são qualificados 10 para a final, que é constituída por 20 países.

O logotipo oficial da competição permaneceu o mesmo de 2004 e 2005 com a bandeira ucraniana no coração sendo alterado. O sub-logótipo foi criado pela empresa de design \textbf{Karamela} (LOC) para a televisão grega. Foi aparentemente baseado no disco de \textbf{Phaistos} (LOC) que é um símbolo popular da \textbf{Grécia} (LOC) antiga. De acordo com o \textbf{TRE} (ORG), foi "inspirado pelo vento e o mar, a luz dourada do sol e o brilho da areia". Na sequência de \textbf{Istambul} (LOC) "sob o mesmo céu" e \textbf{Kiev} (LOC) "Despertar", o slogan para o show de 2006 foi "\textbf{Feel The Rhythm} (MISC)". Este tema foi também a base para os cartões-postais para o show de 2006, que enfatizava a importância histórica da \textbf{Grécia} (LOC), além de ser um importante destino turístico moderno.



Países participantes

Participaram neste festival 37 países: 14 países acederam directamente à final de 20 de \textbf{Maio} (PER), os restantes 23 países tiveram de participar numa semifinal que apurou 10 países, que se juntaram posteriormente aos 14 já apurados, perfazendo assim um total de 24 países participantes na final do festival.

A \textbf{Arménia} (LOC) entra pela primeira vez no \textbf{Festival Eurovisão da Canção} (MISC) .

A \textbf{Áustria} (LOC) e \textbf{Hungria} (LOC) não participaram este ano no festival. Este último devido a razões financeiras, anunciadas pela televisão húngara.

Apesar do seu interesse em entrar no festival, a \textbf{Geórgia} (LOC) não constou na lista dos participantes da \textbf{Eurovisão} (LOC) deste ano.

Mais uma vez os únicos países europeus que não participaram são \textbf{Itália} (LOC) , \textbf{Eslováquia} (LOC) , \textbf{República Checa} (LOC) , \textbf{Luxemburgo} (LOC) , \textbf{São Marino} (LOC) e \textbf{Vaticano} (LOC) .

O \textbf{Liechtenstein} (LOC) não pode participar por não pertencer à \textbf{UER} (ORG) ou \textbf{EBU} (MISC) (\textbf{União Europeia de Radiotelevisão} (ORG)) e não ter televisão pública.

A \textbf{Sérvia e Montenegro} (LOC) decidiu retirar-se do festival devido a irregularidades nas votações, causadas pelos problemas internos dos dois estados que completam o país. Assim sendo, a \textbf{Croácia} (LOC) passou directamente da semifinal à final completando os 14 países pré-qualificados para a final. Contudo, a \textbf{Sérvia e Montenegro} (LOC) votou para a escolha do país vencedor da eurovisão.

\textbf{Portugal} (LOC), mais uma vez falhou a qualificação para a final, tendo a canção sido eliminada.



\textbf{Festival} (PER)Devido à grande quantidade de países votantes na semi-final, a tabela de resultados foi dividida em duas, de forma a não ultrapassar os limites do artigo, e de forma a não estragar o layout e forma do artigo.

Os 10 participantes mais votados na semifinal participaram igualmente na final.



Festival Eurovisão da Canção 2007



Países que classificaram-se para a final \textbf{Países} (MISC) que não se classificaram para a final \textbf{Países} (MISC) que participaram no passado, mas não em 2007

O \textbf{Festival Eurovisão da Canção de 2007} (MISC) (em inglês : \textbf{Eurovision Song Contest 2007} (ORG) , em francês : \textbf{Concours Eurovision} (MISC) de la chanson 2007 e em finlandês : \textbf{Eurovision Laulukilpailu} (PER) 2007 ) foi a 52. edição do \textbf{Festival Eurovisão da Canção} (MISC) e teve como data de realização os dias \textbf{10 de Maio} (MISC) (\textbf{Semifinal} (LOC)) e 12 de \textbf{Maio} (LOC) (\textbf{Final} (LOC)) de 2007 , na \textbf{Finlândia} (LOC) , país vencedor da \textbf{Eurovisão} (LOC) em 2006 com o grupo \textbf{Lordi} (PER) . O espaço que recebeu o evento é o \textbf{Hartwall Areena} (LOC) , na \textbf{Helsínquia} (LOC) , estando toda a organização a cargo da estação de televisão pública finlandesa \textbf{Yle} (PER) e da \textbf{EBU} (MISC), \textbf{European Broadcasting Union} (MISC). \textbf{Jaana Pelkonen} (PER) e \textbf{Mikko Leppilampi} (ORG) foram os apresentadores desta edição da \textbf{Eurovisão} (LOC).

Em \textbf{Março de 2006} (MISC), anunciou-se que a qualificação para a final de 2007 mudaria relativamente aos anos precedentes, de modo que somente os 6 melhores países da final de 2006, juntamente com os 4 contribuintes financeiros da \textbf{EBU} (MISC) (\textbf{the Big Four} (MISC)) seriam qualificados automaticamente, e os restantes lugares seriam ocupados pelos 14 melhores países da semifinal de 2007. Entretanto, numa conferência de imprensa imediatamente antes da competição de 2006, a \textbf{EBU} (MISC) indicou que estas mudanças ao sistema da qualificação não ocorreriam, mantendo-se assim o sistema dos anos anteriores, ou seja, 14 países (os 10 melhores países da final de 2006 e os \textbf{Quatro Grandes} (LOC)) estão automaticamente qualificados para a final de 12 de \textbf{Maio} (LOC), os restantes países têm de participar numa semifinal que selecciona 10 países, que se juntam posteriormente aos 14 já apurados, perfazendo assim um total de 24 países participantes na final do festival.



\textbf{Formato} (MISC)A 12 de \textbf{Março de 2007} (MISC) , os sorteios para a ordem de apresentação para o procedimento semi-final, final e a votação teve lugar. Um novo recurso permitiu cinco wild-card países a partir dos países semi-final e três do final para escolher a sua posição inicial. Os chefes de delegação subiu ao palco e escolheu o número que iria tomar. Na semi-final, a \textbf{Áustria} (LOC), \textbf{Andorra} (LOC), \textbf{Turquia} (LOC), \textbf{Eslovénia} (LOC) e \textbf{Letónia} (LOC) foram capazes de escolher suas posições. No final, \textbf{Armênia} (LOC), \textbf{Ucrânia} (LOC) e \textbf{Alemanha} (LOC) foram capazes de exercer esse privilégio. Todos os países optaram por dar pontos na segunda metade de ambas as noites. Logo após o sorteio, as entradas foram aprovadas pelo \textbf{EBU} (MISC), acabando com a possibilidade de desclassificação para a canção de \textbf{Israel} (LOC). O \textbf{Reino Unido} (LOC) escolheu a sua entrada depois do prazo de limite, porque eles foram concedidos providência especiais da \textbf{EBU} (MISC) .

O concurso viu algumas pequenas alterações ao voto de tempo. O vídeo síntese de compilação de todas as entradas, incluindo números de telefone foi mostrado duas vezes. O processo de votação foi a mesma de 2006, excepto havia 15 minutos para votar, um aumento de 5 na edição de 2006. No final, os resultados de cada país foram mais uma vez mostrado 1-7 pontos automaticamente na tabela e apenas 8, 10 e 12 pontos foram lidos pelos porta-vozes. Pela primeira vez, o vencedor foi premiado com uma promoção turnê pela \textbf{Europa} (LOC) , visitando \textbf{Dinamarca} (LOC) , \textbf{Espanha} (LOC) , \textbf{Suécia} (LOC) \textbf{Holanda} (LOC) , \textbf{Grécia} (LOC) e \textbf{Alemanha} (LOC). A turnê foi realizada entre 16 de maio e 21 de maio . O evento foi patrocinado pela \textbf{European Communications Group TeliaSonera} (ORG) , e - como acontece com várias edições anteriores - \textbf{Nobel Biocare} (MISC). \textbf{Apocalyptica} (LOC) actuaram no intervalo.

O logotipo oficial da competição permaneceu o mesmo em 2006, a bandeira no centro do coração foi mudado para a bandeira finlandesa. A \textbf{União Europeia de Radiodifusão} (ORG) e \textbf{Yle} (PER) anunciou que o tema para a edição de 2007 seria "\textbf{Fantasy True} (MISC)", que abraçou a \textbf{Finlândia} (LOC) e "\textbf{Finnishness} (MISC)" em termos de polaridades associadas com o país. The design agency \textbf{Dog Design} (PER) was responsible for \textbf{the design of the} (MISC) visual theme of the contest \textbf{which incorporated vibrant kaleidoscopic} (PER) patterns formed from various symbols \textbf{including exclamation marks and the letter} (MISC) F. o palco era em forma de um kantele , um instrumento tradicional da \textbf{Finlândia} (LOC). Em 20 Fevereiro 2007 um site oficial reformulado para o concurso foi lançado marcando a primeira exposição pública do tema deste ano. Um \textbf{CD} (MISC) oficial e DVD foram lançados (mas não HD DVD ou Blu-ray, apesar de o evento ser transmitido em alta definição pela primeira vez). Um livro de fãs oficial também foi lançada. Os temas dos cartões postais (vídeos curtos entre os atos) eram histórias curtas acontecendo em diferentes lugares da \textbf{Finlândia} (LOC). De salientar ter sido o primeiro festival \textbf{Eurovisivo} (MISC) com acabamentos em HD, sendo a sua maioria gravado em \textbf{Alta Definicao} (LOC).



\textbf{Sistema de Votação} (MISC)

O sistema de votação foi efectuado por televoto, onde cada país votou nas suas dez canções preferidas, numa pontuação que vai de 1 a 7, e 8, 10 e 12 pontos. A votação por televoto, foi efectuada em ambas as semifinais e na final, porém nas semifinais, havia um sistemade júris, que escolheram o décimo país que passaria directamente à final.



\textbf{Festival} (PER)Qualquer membro da \textbf{União Europeia de Rádio e Televisão} (ORG) pode participar no \textbf{Festival} (LOC). A 1 de Dezembro de 2006 foi anunciado que o limite máximo de países participantes no total (\textbf{Semi-Final + Final} (LOC)) seria 43, dado que houve um crescente interesse por parte de alguns países em participar no festival de 2007. Após a confirmação da desistência do \textbf{Mónaco} (LOC), o \textbf{Festival Eurovisão da Canção 2007} (MISC) contou com a presença de 42 países.

Pela primeira e última vez 28 países participaram na semifinal.

Devido à grande quantidade de países votantes na semi-final, a tabela de resultados foi dividida em duas, de forma a não ultrapassar os limites do artigo, e de forma a não estragar o layout e forma do artigo.



Festival Eurovisão da Canção 2008



Países que classificaram-se para a final \textbf{Países} (MISC) que não se classificaram para a final \textbf{Países} (MISC) que participaram no passado, mas não em 2008

O \textbf{Festival Eurovisão da Canção de 2008} (MISC) (em inglês : \textbf{Eurovision Song Contest} (ORG) 2008 , em francês : \textbf{Concours Eurovision} (MISC) de la chanson 2008 e em sérvio : \textbf{  2008} (PER) ) foi a 53 edição do \textbf{Festival Eurovisão da Canção} (MISC) , que foi sediada na capital da \textbf{Sérvia} (LOC) , a \textbf{Cidade de Belgrado} (LOC) . A primeira semifinal foi realizada em 20 de \textbf{Maio de 2008} (MISC) , a segunda semifinal em 22 de \textbf{Maio} (LOC) e, por fim, a final ocorreu a 24 de Maio do mesmo ano. Todos os três espetáculos ocorreram na \textbf{Beogradska Arena} (LOC) , uma das maiores arenas da \textbf{Europa} (LOC), com uma capacidade total de 23.000 lugares. O país vencedor foi a \textbf{Rússia} (LOC) , com o cantor \textbf{Dima Bilan} (PER) a interpretar a canção " \textbf{Believe} (PER) ", com 272 pontos (ganhando por uma margem de 42 pontos). Na primeira semifinal, o vencedor foi \textbf{a Grécia} (LOC) , representada por \textbf{Kalomira} (LOC) , com a canção " \textbf{Secret Combination} (MISC) ", recebendo 156 pontos (vencendo com uma margem de 17 pontos). Já na segunda semifinal, a vencedora foi a \textbf{Ucrânia} (LOC) , representada pela cantora \textbf{Ani Lorak} (PER) , com o tema " \textbf{Shady Lady} (MISC) ". no final, a cantora recebeu 152 pontos, vencendo a segunda semifinal por 32 pontos. Na edição do Festival Eurovisão da Canção 2007 , a \textbf{Sérvia} (LOC) ganhou o direito de realizar o \textbf{Festival Eurovisão da Canção 2008} (MISC), com a música \textbf{Molitva} (MISC) , interpretada pela cantora \textbf{Marija Šerifović} (PER) , tornando-se assim a primeira vez que o país recebe e organiza o evento. Neste ano participaram quarenta e três países, um número recorde nunca antes atingido no festival. Dado ao elevado número de participantes, o responsável pelo \textbf{Festival Eurovisão da Canção} (MISC), \textbf{Svante Stockselius} (PER) propôs (ainda primeiramente para o ano de 2009) a realização de duas semifinais e de uma final após ter recebido várias queixas por partes de fãs, jornalistas e participantes relacionados com o "voto vizinho" (votos dos blocos dos países do \textbf{Leste europeu} (LOC), dos \textbf{Balcãs} (LOC) e da \textbf{Escandinávia} (LOC) ), o "voto emigrante" e a política. Mais tarde, depois da apresentação de uma hipotética realização de tal formato, o director geral responsável pela \textbf{Eurovisão} (LOC), confirmou a realização de duas semifinais e uma final. A realização do concurso no país foi também marcada pelo conflito entra a \textbf{Sérvia} (LOC) e a independência do \textbf{Kosovo} (LOC) . Dados os conflitos e as manifestações ocorridas nas ruas da capital, a \textbf{EBU} (MISC) declarou numa conferência por telefone a possibilidade de anular o \textbf{Festival} (LOC) e transferi-lo para outra cidade tendo mesmo recebido propostas por parte da \textbf{Ucrânia} (LOC) (segunda classificada em 2007) e da \textbf{Finlândia} (LOC) (país anfitrião do ano passado), e ainda da \textbf{Grécia} (LOC) (organizadora do festival em 2006). Porém, o governo sérvio e a sua televisão estatal \textbf{RTS} (ORG) prometeram uma forte medida de segurança para os participantes e para os espectadores que chegassem à cidade. As delegações da \textbf{Croácia} (LOC) , \textbf{Albânia} (LOC) e \textbf{Israel} (LOC) tiveram medidas especiais por parte da segurança. O tema deste ano foi a "\textbf{Confluência de Sons} (MISC)" ( \textbf{The Confluence of Sounds} (MISC) ) que está relacionado com o facto da cidade sérvia estar localizada ao pé de uma confluência de dois grandes rios da \textbf{Europa} (LOC): o \textbf{Rio Danúbio} (LOC) e o \textbf{Rio Sava} (LOC) . O festival foi apresentado por \textbf{Jovana Janković} (PER) e \textbf{Željko Joksimović} (PER) (antigo representante do país em 2004). Também o sistema de votação foi ligeiramente alterado, adicionando um júri na primeira e na segunda semifinal, para selecionar um décimo país a passar à grande finalíssima.



\textbf{Local} (MISC)A \textbf{Sérvia} (LOC) ganhou o direito de organizar o \textbf{Festival Eurovisão da Canção 2008} (MISC) após a vitória de \textbf{Marija Šerifović} (PER) no Festival Eurovisão da Canção 2007 em \textbf{Helsinque} (LOC) , na \textbf{Finlândia} (LOC) , com a música " \textbf{Molitva} (MISC) ". A regra implicou que o festival de 2008 deveria ocorrer no país. A \textbf{Belgrade Arena} (LOC) foi escolhida como local para o festival, já que era o único local disponível no país para acolher a disputa. No dia 14 de Setembro de 2007 , \textbf{Jussi Pajunen} (PER) , \textbf{Prefeito de Helsinque} (LOC) , entregou as "chaves da \textbf{Eurovisão} (LOC)" a prefeito de \textbf{Belgrado} (LOC). Este simbolismo significou o início de mais uma tradição na \textbf{Eurovisão} (LOC) para o futuro, e o chaveiro contém uma tranca de cada cidade que alguma vez organizou o festival.

Devido a problemas e conflitos em \textbf{Belgrado} (LOC), por conta da declaração unilateral da independência do \textbf{Kosovo} (LOC) a 22 de fevereiro de 2008 , a \textbf{EBU} (MISC) realizou uma conferência via telefone para decidir se o concurso deveria ou não ser transferido para outro país. A \textbf{Ucrânia} (LOC) foi desde logo considerada como uma opção, e a mais fiável, por ter ficado em segundo lugar no festival do ano anterior . A \textbf{YLE} (MISC) foi outra opção ponderada, visto a estação televisiva ter organizado o concurso em 2007, em \textbf{Helsinque} (LOC) , \textbf{Finlândia} (LOC) . A estação televisiva grega \textbf{Ellinikí Radiofonía Tileórasi} (PER) (\textbf{ERT} (ORG)) também se ofereceu (à \textbf{EBU} (MISC)) para realizar o concurso em \textbf{Atenas} (LOC) , \textbf{Grécia} (LOC) (à semelhança do que já havia acontecido no passado, em 2006 ). Mais tarde, acabou por ser decidido que o festival continuaria em \textbf{Belgrado} (LOC), com a \textbf{EBU} (MISC) a dar todo o suporte necessário para a realização do mesmo.

Além da \textbf{Belgrade Arena} (LOC), o festival pode ser assistido na principal praça da cidade. A praça do \textbf{Velho Palácio} (LOC) , praça a qual, é a da câmara municipal da cidade.

Devido aos problemas relacionados com a segurança, mais tarde, para além da normal segurança para o evento, a \textbf{RTS} (ORG), televisão anfitriã, teve a garantia de segurança total, por parte do governo sérvio, para todos os visitantes e participantes do concurso. As delegações da \textbf{Albânia} (LOC) , \textbf{Croácia} (LOC) e \textbf{Israel} (LOC) tiveram segurança especial, por serem considerados "países de risco".

\textbf{Belgrado} (LOC) conta com 112 rotas de autocarro, 12 linhas de eléctricos e 8 linhas de autocarros elétricos (género monorail ). Os bilhetes podem ser adquiridos diretamente com os motoristas dos próprios transportes (com um custo de 48 RSD, cerca de 0,5 EUR) ou em quiosques (cerca de 28 RSD, equivalente a 0,25 EUR). O controlo de passagens é feito esporadicamente pelos funcionários das respectivas companhias operadoras. \textbf{Belgrado} (LOC) concentra um nó ferroviário que permite conexões com as principais capitais da \textbf{Europa} (LOC). O \textbf{Aeroporto Internacional Nikola Tesla} (LOC) voltou a receber um número progressivamente maior de voos depois da retirada, em 2000 , das sanções econômicas impostas à antiga \textbf{Jugoslávia} (ORG) pela \textbf{ONU} (ORG) . Estradas de grande porte oferecem acesso a \textbf{Novi Sad} (LOC) (norte), \textbf{Niš} (LOC) (sul), e \textbf{Zagreb} (LOC) (oeste, na \textbf{Croácia} (LOC) ). \textbf{Belgrado} (LOC) possui quatro pontes sobre os rios \textbf{Sava} (LOC) e \textbf{Danúbio} (LOC) , sendo as mais importantes a \textbf{Ponte Branko} (LOC) e a \textbf{Ponte Gazela} (LOC) , que conectam o centro da cidade à \textbf{Nova Belgrado} (LOC). O \textbf{Porto de Belgrado} (LOC) encontra-se na costa do \textbf{rio Danúbio} (LOC). A via fluvial é usada predominantemente para o transporte de mercadorias.

Para o festival em si, a organização, à semelhança de anos anteriores, alugou vários autocarros para o transporte de todas as quarenta e três delegações pela \textbf{Cidade de Belgrado} (LOC) . Durante os quinze dias em que ocorreram ensaios e os grandes espetáculos, a cidade de \textbf{Belgrado} (LOC) viu uma grande enchente de turistas vindos de toda a \textbf{Europa} (LOC), não só para ver os três espetáculos mais vistos na \textbf{Europa} (LOC), mas também para frequentarem as festas organizadas directamente relacionadas com o tema \textbf{Eurovisão} (LOC) , durante a \textbf{Semana Eurovisiva} (MISC) .



\textbf{Formato} (PER)Durante uma conferência de imprensa em \textbf{Helsínquia} (LOC) , em \textbf{Maio de 2007} (MISC) , \textbf{Svante Stockselius} (PER) , executivo supervisor do festival pela \textbf{EBU} (MISC), anúnciou que o formato da competição deveria ser estendido para duas semifinais em 2008 ou 2009. A \textbf{28 de Setembro de 2007} (MISC) , foi confirmado que a \textbf{EBU} (MISC) havia aprovado o plano para a realização de duas semifinais já em 2008.

Baseado na pesquisa conduzida pela parceira da \textbf{EBU} (MISC) que trata do sistema do televoto da Eurovisão, a \textbf{Digame} (LOC) , os semifinalistas foram colocados em dois potes a quando da divisão por outros seis potes. Esta primeira divisão foi feita para manter os países que têm uma história de votação significativa, longe dos países em que mais votaram ao longo da história. Cada televisão presente no festival, tinha obrigatoriamente que emitir a semifinal em que entraram/participaram, com a transmissão da outra semifinal em formato opcional, isto é, apenas transmitido caso o canal televisivo assim o entendesse. O sorteio para dividir os países para as semifinais, ocorreu no \textbf{Belgrade City Hall} (ORG) (em português conhecido por \textbf{Câmara Municipal} (PER)), na segunda-feira do dia 28 de \textbf{Janeiro} (LOC) de 2008 , às 13:00 da tarde \textbf{CET} (ORG) e a cerimónia foi conduzida por dançarinos da \textbf{National Dance Ensemble KOLO} (ORG). Primeiro, dois envelopes com '\textbf{Semi-Final} (ORG) 1' e '\textbf{Semi-Final 2} (MISC)' foram sorteados. Depois, três países de cada pote foram escolhidos aleatoriamente para participarem na primeira semifinal e outros três na segunda. O país que sobrou no \textbf{Pote 5} (LOC), participou no primeiro envelope que foi anteriormente sorteado. Enquanto que o país deixado no \textbf{Pote 6} (LOC), concorreu na segunda semifinal.

Os países automaticamente qualificados para a final (os finalistas), escolheram se transmitiriam ambas as semifinais, ou apenas uma delas, mas os telespectadores destes cinco países (os \textbf{Big 4} (MISC) e o país anfitrião), apenas poderiam votar numa das semifinais. Do mesmo sorteio realizado para dividir os países semifinalistas, também foi decidido quais dos cinco países finalistas transmitiriam e teriam direitos de votos num dos dois eventos antecessores da grande final. As semifinais e a final, para além da transmissão televisiva, também foram transmitidas a nível mundial, através de um webcast live , efectuado pela Eurovision.tv . Os nove países com mais votos no televoto qualificaram-se directamente para a grande final, a realizar-se no sábado daquela semana (\textbf{Eurovisiva} (MISC)), e um décimo país foi escolhido/determinado por um júri independente. Ficou então estipulado que na final apenas participariam vinte e cinco países, não repetindo assim grandes finais ocorridas no passado com mais de trinta participantes no mesmo espetáculo.

A \textbf{RTS} (ORG) , televisão anfitrião, lançou uma competição com o objectivo da criação e invenção de um tema, de um logo e do design para um palco. O tema do \textbf{Festival Eurovisão da Canção 2008} (MISC), andou à volta da "\textbf{confluence of sound} (MISC)" ( \textbf{Confluência do Som} (MISC) ). Este factor, estava carregado de simbolismo, visto que a \textbf{Cidade de Belgrado} (LOC), encontra-se na confluência de dois famosos rios \textbf{Europeus} (LOC), o \textbf{rio Sava} (LOC) e o \textbf{rio Danúbio} (LOC) . O logótipo escolhido, uma clave de sol formou a base gráfica do design criado por \textbf{Boris Miljković} (PER) .

Os cartões postais na primeira e segunda semifinal, assim como na final, basearam-se à volta da criação da bandeira da nação que actuaria de seguida, através de actuações artísticas como dança, levantamento de pesos, malabarismos, hotelaria, amor, ginástica, desfiles de moda, entre outros. Cada cartão postal continha uma pequena história que relembrava \textbf{Belgrado} (LOC) e as suas pessoas. Durante a transmissão de cada cartão postal, uma pequena carta, escrita pelo músico do próximo país a atuar, foi demonstrada no ecrã. Todas e cada carta, foram escrita na língua materna do país do artista que representaria o seu país a seguir, à exceção do cartão da \textbf{Sérvia} (LOC), que consistiu em dar as boas-vindas a \textbf{Belgrado} (LOC) e à \textbf{Sérvia} (LOC) (no ecrã apareceu "\textbf{Welcome to Belgrade} (MISC)" e "\textbf{Welcome to Serbia} (MISC)" em várias línguas), e o cartão da \textbf{Bélgica} (LOC) que foi escrito na mesma língua feita depropositadamente para a sua entrada. Todos os cartões postais acabaram, sendo estampados com o logótipo do Festival Eurovisão da Canção 2008.r De acordo com a \textbf{RTS} (ORG), o palco representava identidades nativas, história e temas modernos, símbolos e mensagens reconhecidas universalmente. O palco baseado no próprio tema do festival, "\textbf{Confluence of Sound} (MISC)", conteve um largo número de televisões e de LCD display screens . O palco tinha ainda alguma das mais recentes novidades em possibilidades electrónicas, incluindo a capacidade de algumas partes do palco se moverem Toda a infra-estrutura, foi desenhado pelo arquitecto de \textbf{Chicago} (LOC) , \textbf{David Cushing} (PER) .

A primeira semifinal foi criada em torno do tema da própria cidade. O concurso foi aberto com o panorama da cidade de \textbf{Belgrado} (LOC) formando no painel mais atrás do palco, duas ondas que "escorreram" pelo palco abaixo, até se tocarem no centro de toda a estrutura - na confluência, o principal tema do festival.

A segunda semifinal, já teve como tema mais próprio a água, que foi enaltecido pelo visual do palco durante o acto do intervalo, onde a água formou as cores principais do palco.

Por fim, a final baseou-se essencialmente no tema da confluência. A construção do palco, demorou vários dias e foi realizada por várias equipas provenientes de toda a \textbf{Europa} (LOC). A utilização de fogos de artifício foi fortemente utilizada nas entradas da \textbf{Bósnia e Herzegovina} (LOC) , da \textbf{Finlândia} (LOC) , da \textbf{Turquia} (LOC) , da \textbf{República Checa} (LOC) , da \textbf{Bulgária} (LOC) e da \textbf{Suíça} (LOC) . O palco recebeu "nota positiva" por parte dos mídia e dos fãs em geral, descrevendo o palco da quinquagésima terceira edição do festival, como "um dos melhores e mais bonitos palcos na história de toda a competição".



\textbf{Sistema de Votação} (MISC)

O sistema de votação foi efectuado por televoto, onde cada país votou nas suas dez canções preferidas, numa pontuação que vai de 1 a 7, e 8, 10 e 12 pontos. A votação por televoto, foi efectuada em ambas as semifinais e na final, porém nas semifinais, havia um sistema de \textbf{júris} (LOC), que escolheram o décimo país que passaria diretamente à final.



\textbf{Participações individuais} (MISC)Durante cerca de 10 meses, todos os 43 países elegeram os seus representantes, assim como as músicas que os mesmos viriam a interpretar em \textbf{Moscovo} (LOC). Para realizar tal seleção, cada país utilizou o seu próprio processo de seleção. Alguns optaram pela seleção interna, que consiste em a televisão organizadora daquele país, é quem faz a escolha; no entanto, por vezes apenas o artista é selecionado internamente, e a música não. Outros países (a maioria), utiliza um programa de televisão para selecionar a sua entrada. Quartos de final, semifinais, second-chances e finais, foram realizadas durante 10 meses na maioria dos países europeus, cada um com o seu processo próprio (também a internet foi utilizada na fase de escolhas). No conjunto de artigos em baixo, é possível ler mais ao pormenor o tipo de processo que cada país utilizou, assim como os resultados e reações.



\textbf{Participantes} (MISC)A 21 de Dezembro de 2007 , a \textbf{EBU} (MISC) confirmou que estariam presentes 43 países, no \textbf{Festival Eurovisão da Canção 2008} (MISC), em \textbf{Belgrado} (LOC), \textbf{Sérvia} (LOC). O \textbf{São Marino} (LOC) , como o mais recente membro ativo da \textbf{EBU} (MISC) , e o \textbf{Azerbaijão} (LOC) , realizaram a sua estreia na \textbf{Eurovisão} (LOC) no concurso de 2008. Por sua vez, a \textbf{Áustria} (LOC) retirou-se do festival; a sua televisão nacional responsável pela representação do país na \textbf{Eurovisão} (LOC), \textbf{ORF} (ORG) , constatou "we've already seen in 2007 that \textbf{it's not the quality of the song} (MISC), \textbf{but the country} (MISC) of origin that determines \textbf{the decision} (MISC)" (Nós já vimos em 2007 que não é a qualidade da música que conta, mas sim o país de origem que determina a decisão). A \textbf{Itália} (LOC), que já não competia desde 1997, e a qual teria sido uma finalista direta, como parte dos \textbf{Big 4} (MISC) (que caso \textbf{Itália} (LOC) participasse passaria a chamar-se \textbf{Big 5} (MISC)), voltou a não marcar presença na \textbf{Eurovisão} (LOC) pelo décimo primeiro ano consecutivo, numa que é das mais longas ausências feitas por um país na \textbf{Eurovisão} (LOC) ainda existente. Por parte da \textbf{Eslováquia} (LOC), a explicação para a sua não participação, baseou-se em problemas financeiros, relacionados com o início da grande crise financeira que se abateu sobre o \textbf{Mundo} (MISC) durante o ano de 2008. As formas de seleção para os representantes de cada país, estiveram a cargo de cada uma das quarenta e três estações televisivas que se inscreveram no festival. A maior parte dos países, utilizou um programa de seleção em aberto, isto é, com a possibilidade de todo e qualquer cidadão, nacionalizado naquele determinado país, pudesse submeter uma música própria. Outra forma de selecionar, é internamente (meio optado por algumas televisões, também de forma a poupar algum dinheiro). Em 2008 , o primeiro país a escolher o seu representante, foi \textbf{Andorra} (LOC) , ainda em dezembro de 2007 , o país escolheu o seu representante. Por sua vez, o último país a fazê-lo foi a \textbf{Suécia} (LOC) , a realizar a sua final da sua tradicional e famosa forma de apresentação, a meio do mês de \textbf{Março} (PER) de 2008 . Depois de todas as seleções realizadas, surgiram na lista de entradas, dois artistas que participaram pela segunda vez no evento. A representante da \textbf{Suécia} (LOC), \textbf{Charlotte Perrelli} (PER) ganhou o \textbf{Festival Eurovisão da Canção} (MISC) 1999 representando o seu país na altura, e novamente em 2008; por sua vez, a \textbf{Rússia} (LOC) volta a enviar um mesmo representante, \textbf{Dima Bilan} (PER) , que alcançou o segundo lugar no \textbf{Festival Eurovisão da Canção 2006} (MISC) ,e que viria a ganhar a edição de 2008. Por outro lado, um dos autores do tema que representou \textbf{Israel} (LOC), é a famosa artista \textbf{Dana International} (PER) , vencedora do \textbf{Festival Eurovisão da Canção 1998} (MISC) . A cantora \textbf{Gisela} (LOC) , representante de \textbf{Andorra} (LOC), formou parte do coro da representante espanhola \textbf{Rosa López} (PER) no Festival Eurovisão da Canção 2002 . Durante a edição de 2008, foi apresentada, a milésima centésima música apresentada na \textbf{Eurovisão} (LOC), que contava já com cinquenta e três edições em 2008. O \textbf{Festival Eurovisão da Canção 2008} (MISC), marcou ainda a história do concurso ainda mais profundamente, através da implementação de mais uma semifinal, perfazendo o total de duas. Sendo assim, o festival deixaria de ter apenas duas noites (como o que já vinha a acontecer desde 2004 ), e passou a ter três noites (o que perfez um total de catorze dias seguidos, conhecidos como \textbf{Época Eurovisiva} (PER) , onde entra a já famosa \textbf{Semana Eurovisiva} (MISC) ).



\textbf{Organização} (ORG)Uma vez concluídas as seleções internas de cada país em relação ao seu representante, a 17 de março de 2008 realizou-se um sorteio onde se determinou a ordem de apresentação das canções em cada semifinal, e a ordem de votação, para a grande final.

A 24 de Janeiro de 2008 , todos os 38 países nas semifinais foram separados da seguinte maneira, pelo seu histórico de votação e pela sua localização geográfica:

Depois do sorteio, ambas as semifinais, foram organizadas da seguinte forma:

A ordem de entrada dos países, quer nas semifinais, quer na grande final, foi escolhida, a 17 de \textbf{Março 2008} (MISC) , numa reunião com os altos responsáveis do \textbf{Festival} (LOC) e da \textbf{EBU} (MISC).



\textbf{Festival} (PER)A primeira semifinal do Festival Eurovisão da Canção 2008 realizou-se na terça-feira, dia 20 de \textbf{Maio de 2008} (MISC) . Dezanove países participaram neste evento, cada um deles com direito a voto uma vez finalizada a apresentação de todos os temas, e a juntar-se a eles, estavam a \textbf{Alemanha} (LOC) e a \textbf{Espanha} (LOC) , que participaram apenas como votantes.

O evento iniciou-se com uma apresentação intitulada por \textbf{Video} (PER) killed a radio star ("O vidio matou uma estrela da rádio", interpretado em inglês , inspirado no título de uma famosa canção do grupo britânico \textbf{The Buggles} (PER) ), depois desta apresentação,apareceram os apresentadores e deu-se lugar à parte mais aguardada do evento, a apresentação das canções a concurso. Uma vez concluída a apresentação das dezanove canções, o tenista sérvio \textbf{Novak Đoković} (PER) , lançou uma bola com a sua raquete, para o meu dos espectadores presentas na \textbf{Belgrade Arena} (LOC), a fim de dar início à votação por parte dos países aptos a votar nesta semifinal, que duraria quinze minutos. Enquanto a votação decorria, foram recordadas as dezanove apresentações mais que uma vez, ao fim do tempo de votação, realizou-se um intervalo, para que os votos pudessem ser analisados. Durante este intervalo, os artistas foram os \textbf{Metropole Orkest} (PER) , \textbf{Slobodan Trkulja y el grupo Balkanopolis} (PER) .

Os resultados da primeira semifinal, apenas dados a conhecer após o final do \textbf{Festival Eurovisão da Canção} (MISC), afirmaram que a cantora grega \textbf{Kalomira} (PER) , tinha ficado em primeiro lugar com a sua canção " \textbf{Secret combination} (MISC) ", uma música pop e \textbf{R\&B} (LOC) , a qual alcançou 156 pontos (incluindo 4 países que lhe deram o máximo possível de votos). " \textbf{The fire in your} (MISC) eyes " de \textbf{Boaz Mauda} (PER), que representou \textbf{Israel} (LOC) , obteve o segundo lugar com 139 pontos, porém, recebeu os famosos doze pontos por parte de cinco países. A balada " \textbf{Believe} (LOC) " do russo \textbf{Dima Bilan} (PER) acabou em terceiro lugar com apenas menos quatro pontos que o segundo lugar, ficou assim com 135 pontos. No festival deste ano, também houve uma novidade no sistema de votação, a implementação de um júri autónomo, com poder para escolher o décimo classificado desta semifinal, escolheu a artista \textbf{Isis Gee} (PER) , representante da \textbf{Polónia} (LOC) , que coincidentemente obteve o décimo lugar na votação dos espectadores.

Países com uma grande história na \textbf{Eurovisão} (LOC), como a \textbf{Bélgica} (LOC) , a \textbf{Irlanda} (LOC) e a \textbf{Holanda} (LOC) voltaram a ficar de fora da final, e a nenhum deles chegar/ficar sequer no top ten da escolha do público. No entanto, os dois países estreantes, tiveram resultados bastante diferentes: o \textbf{Azerbaijão} (LOC) classificou-se para a grande final com o tema " \textbf{Day after day} (MISC) ", onde \textbf{Elnur Hüseynov} (PER) e \textbf{Samir Javadzadeh} (PER) representavam musicalmente a luta entre um anjo e um demónio ; já o tema em italiano , " \textbf{Complice} (PER) ", que representou \textbf{San Marino} (LOC) , ficou em último lugar da tabela.

A segunda semifinal realizou-se na quinta-feira, dia 22 de \textbf{Maio de 2008} (MISC) e, igualmente à primeira semifinal, contou com dezanove participantes. Três classificados diretamente para a final, a \textbf{França} (LOC) , o \textbf{Reino Unido} (LOC) e a \textbf{Sérvia} (LOC) , exerceram o seu direito de voto nesta semifinal, como havia ficado estipulado durante o sorteio para tal.

O espetáculo iniciou-se com uma apresentação denominada de \textbf{Srbija za početnike} (PER) ("\textbf{Sérvia} (LOC) para principiantes" executada em sérvio ) que contou com o bailarino e ator \textbf{Aleksandar Josipović} (PER) como mestre de cerimónias. Depois da abertura, que esteve bastante ligada à água como o principal tema do evento, foi a vez de as dezanove canções a concurso se fazerem ouvir. O período de quinze minutos para as votações telefónicas foi aberto simbolicamente com a presença de \textbf{Lys Assia} (PER) , vencedora do Festival Eurovisão da Canção 1956 , o primeiro alguma vez realizado (ocorreu, fez em 2008, há 53 anos). Uma apresentação do \textbf{Teatro Nacional de Belgrado} (LOC) deu passo finalmente à entrega dos resultados e dos dez classificados.

Nesta semifinal, \textbf{Ani Lorak} (PER) ganhou com o seu tema de dance-pop " \textbf{Shady lady} (MISC) ", com um total de 152 pontos. A \textbf{Ucrânia} (LOC) , que teve uma apresentação com quatro bailarinos que acompanharam a cantora junto a uma jaula de espelhos, luzes e vidros, obteve a pontuação máxima de seis nações. \textbf{Ani} (PER) foi seguida pela portuguesa \textbf{Vânia Fernandes} (PER) com a canção " \textbf{Senhora do Mar} (MISC) (\textbf{Negras Águas} (LOC)) ", um fado que relatava o drama das viúvas de marinheiros, e que se converteria na melhor posição de \textbf{Portugal} (LOC) na história do concurso até à data (anteriormente alcançou o sexto lugar no \textbf{Festival Eurovisão da Canção 1996} (MISC) ). O terceiro lugar foi para a \textbf{Croácia} (LOC) e o quarto para a \textbf{Dinamarca} (LOC) .

\textbf{Charlotte Perrelli} (PER) , representante da \textbf{Suécia} (LOC) no \textbf{Festival Eurovisão da Canção 2008} (MISC) com o tema pop schlager " \textbf{Hero} (PER) ", foi a primeira "salvada" pelo jurado, encarregue de escolher o décimo país a passar à final. Apesar de ter ganho o \textbf{Festival Eurovisão da Canção 1999} (MISC) e de ser uma das favoritas da imprensa (juntamente com \textbf{Portugal} (LOC)), obteve apenas 54 pontos e acabou em 12. lugar. Posto isto, e ao ser a selecionada dos jurados, conseguiu assim a sua passagem à final juntamente com as nove melhor músicas classificadas pelo público através da votação que ocorreu anteriormente. Outro dos favoritos, era o suíço \textbf{Paolo Meneguzzi} (PER) (vencedor do \textbf{Festival Internacional da Canção de Viña del Mar} (ORG) em 1996 ), não respondeu às expetativas como um dos favoritos e a sua canção " Era stupendo ", alcançou o 13. lugar, ficando assim fora da grande final dois dias depois da realização da segunda semifinal.

A final ocorreu a 24 de maio de 2008 , e na sua abertura atuou a vencedora do Festival Eurovisão da Canção 2007 , que interpretou a sua música vencedora, em versão remix. O evento foi vencido pela \textbf{Rússia} (LOC) . A vitória russa foi mais tarde questionada, em maio do mesmo ano, por oficiais da \textbf{Ucrânia} (LOC) , cuja representante ficou em segundo lugar.

Os seguintes países competiram nas duas semifinais, que foram realizadas/transmitidas na terça-feira, dia 20 de maio de 2008 , e na quinta-fira, dia 24 de maio de 2008 . A estes vinte países provenientes das semifinais, juntaram-se os finalistas automáticos, a \textbf{Alemanha} (LOC) e \textbf{Espanha} (LOC) , que tiveram poder de voto na primeira semifinal, e o \textbf{Reino Unido} (LOC) , a \textbf{França} (LOC) e a \textbf{Sérvia} (LOC) , que votaram na segunda semifinal. A \textbf{Espanha} (LOC) e a \textbf{França} (LOC), transmitiram apenas a semifinal, em que tinham o direito de voto, enquanto que a \textbf{Alemanha} (LOC), a \textbf{Sérvia} (LOC) e o \textbf{Reino Unido} (LOC), transmitiram ambas as semifinais (mas este último, em delay).

A \textbf{Polónia} (LOC) , o \textbf{Reino Unido} (LOC) e a \textbf{Alemanha} (LOC) , todos receberam catorze pontos no toal. As regras do concurso, são claras quanto aos empates. Quando existem dois ou mais países com a mesma pontuação, fica à frente o que teve mais douze pontos, se dois ou mais tiverem o mesmo número de douze pontos, passam a contar os dez pontos, e assim sucessivamente. Com as regras aplicadas, e com a \textbf{Alemanha} (LOC) a perfazer o maior número de douze pontos recebidos, a \textbf{Alemanha} (LOC) acabou em 23. lugar. A \textbf{Polónia} (LOC) obteve o segundo maior número de douze pontos, acabando em 24. lugar, e por fim o \textbf{Reino Unido} (LOC), que foi o país dos três que recebeu menos douze pontos, acabando assim em 25. lugar, o último da tabela. A tabela presente na página oficial do concurso, em Eurovision.tv , mostra as linhas destes países a vazio, no entanto, em anos anteriores como 2007 , 2004 e 2002 , as mesmas linhas estavam vazias, sendo assim erro da própria organização do festival no lançamento dos dados.

O sorteio para a ordem de votação dos quarenta e três países a concurso deu-se a 17 de \textbf{Março de 2008} (MISC) .

Em baixo, estão representados todos os países a que foram atribuídos 12 pontos, assim como o número de vezes:



\textbf{Cobertura} (MISC) televisiva

\textbf{O Festival Eurovisão da Canção} (MISC) , é o programa não-desportivo mais visto em todo o \textbf{Mundo} (MISC) (para lá dos \textbf{Jogos Olímpicos} (MISC) , e finais de \textbf{Europeus} (LOC) e \textbf{Mundiais de Futebol} (MISC) ). Todos os anos, mais de 100 milhões de pessoas por toda a \textbf{Europa} (LOC) vêem o espetáculo, traduzido em direto, e comentado por um comentador escolhido pela sua televisão transmissora, membro da \textbf{EBU} (MISC) . Porém, o festival para além de ser transmitido pela TV também é transmitido através de estações de rádio, e mais recentemente em direto (e posteriormente em vídeo on-demand), no site oficial do festival, através da eurovision.tv . Todos os anos, cerca de 50 televisões, 12 rádios e um canal oficial de televisão on-line ( eurovision.tv ), transmitem o festival para mais de 100 milhões de pessoas, no entanto o número vem a aumentar todos os anos, principalmente por países que já não chegavam à final., chegarem lá, tal aconteceu em 2008 , quando \textbf{Portugal} (LOC) passou pela primeira vez à final, desde a implantação de uma semifinal. Com todo o sucesso que o festival começou a voltar a ter, é esperado que nos próximos anos a final do concurso seja vista por mais de um bilião de pessoas em todo o mundo.

Todas as televisões que enviam um representante para uma edição do \textbf{Festival Eurovisão da Canção} (MISC), são obrigados a transmitir pelo menos a semifinal em que entram e a final. No entanto, os países podem escolher também transmitir a semifinal onde não participam, o que é adoptado pela maioria, mesmo que por vezes a semifinal onde não participam, seja transmitida noutro horário. Como exemplo mais "aficcionado", temos a \textbf{BBC} (ORG) do \textbf{Reino Unido} (LOC) , que transmitiu os três espetáculos em direto. O \textbf{Reino Unido} (LOC) ,, \textbf{San Marino} (LOC) , \textbf{Grécia} (LOC) , \textbf{Croácia} (LOC) , \textbf{Irlanda} (LOC) , \textbf{Alemanha} (LOC) , \textbf{Holanda} (LOC) , \textbf{Noruega} (LOC) , \textbf{Malta} (LOC) , \textbf{Sérvia} (LOC) , \textbf{Finlândia} (LOC) , \textbf{Bósnia e Herzegovina} (LOC) , \textbf{Dinamarca} (LOC) , \textbf{Portugal} (LOC) , \textbf{Chipre} (LOC) , \textbf{Israel} (LOC) , \textbf{Estónia} (LOC) , \textbf{Turquia} (LOC) , \textbf{Letónia} (LOC) , \textbf{Eslovénia} (LOC) , \textbf{Rússia} (LOC) , \textbf{Ucrânia} (LOC) , \textbf{Lituânia} (LOC) , \textbf{República Checa} (LOC) , \textbf{Andorra} (LOC) , \textbf{Albânia} (LOC) , \textbf{Bulgária} (LOC) , \textbf{Islândia} (LOC) , \textbf{Suécia} (LOC) , \textbf{Roménia} (LOC) e a \textbf{Austrália} (LOC) confirmaram que iriam transmitir ambas as semifinais de 2008 (algumas televisões fizeram-no em delay). Em baixo está uma lista com as televisões membros da \textbf{EBU} (MISC), que têm os direitos e o dever de transmitir o \textbf{Festival Eurovisão da Canção 2008} (MISC), assim como os espetáculos que transmitiram e em que tempo (directo ou em delay):

A \textbf{RTS} (ORG) transmitiu o evento em 1080i high-definition (HD) e em 5.1 \textbf{Surround Sound} (LOC). O novo sistema de televisão de alta-definição, foi instalado na totalidade na \textbf{Beograska Arena} (LOC) , em \textbf{Abril de 2008} (LOC) . Este é o segundo ano que o evento é emitiu em alta-definição, em directo. A \textbf{BBC} (ORG) HD transmitiu o concurso em \textbf{High Definition} (MISC) no \textbf{Reino Unido} (LOC) e na \textbf{República da Irlanda} (LOC) . A televisão sueca \textbf{SVT} (ORG) , emitiu ambas as semifinais e a final no seu canal de alta definição, SVT HD . A televisão da \textbf{Lituânia} (LOC) LRT transmitiu os três espetáculos em 1080i high-definition (HD) no seu canal \textbf{LTV} (ORG). O mesmo ocorreu no canal suíço de alta definição HD suisse ; neste canal, os telespectadores também estavam aptos a escolher o idioma dos comentários realizados ao longo dos eventos do \textbf{Festival Eurovisão da Canção 2008} (MISC). Contudo, todos os outros países transmitiram o espetáculo apenas através das suas definições standard, e o evento ficou/está disponível para compra, apenas num DVD de definição standard, e não foi realizado em \textbf{Blu-ray Disc} (PER) . Em baixo, encontra-se a lista das televisões que transmitiram em alta-definição o festival:

O site oficial do \textbf{Festival Eurovisão da Canção} (MISC) , eurovision.tv , também transmitiu o espetáculo em directo para todo o \textbf{Mundo} (MISC) (mas sem qualquer tradução, ou número de telefone para votação), através de \textbf{peer to peer} (MISC) meio \textbf{Octoshape} (ORG) .

Algumas televisões também utilizaram os sus canais de rádio para difundir o festival. Em baixo fica a lista de canais de rádio que emitiram o festival em 2008:

Como já é hábito, desde o início das transmissões do festival, cada televisão membro, escolhe um ou mais comentadores, que em directo, traduzem o festival para a língua materna do país a que pertencem. Na ligação acima, é possível observar uma tabela com todos os comentadores que cada país levará a \textbf{Moscovo} (LOC), para reportar o festival, que é apenas apresentado em inglês e francês, no local onde está a decorrer (por vezes surgem apresentações na língua do país anfitrião, ou em inglês).

Tal como em todos os anos, desde que o festival se rege por um sistema de televoto, cada televisão anfitriã nacional, terá que escolher um porta-voz, que divulgará os votos do televoto do seu país. As pontuações de 1 a 7 apareceram automaticamente no quadro de votações, e posteriormente, o porta voz anunciará os países que receberam os 8, 10 e 12 pontos do seu país. Os resultados do \textbf{Festival Eurovisão da Canção} (MISC), serão posteriormente colocados no site oficial, pelo \textbf{Executivo Superior} (LOC) da \textbf{EBU} (MISC). A lista das pessoas que transmitiram os pontos atribuídos pelo seu país, aos outros países participantes no \textbf{Festival Eurovisão da Canção 2008} (MISC), de cada país, pode ser observada na ligação no início desta secção (porta-vozes do televoto).



\textbf{CD} (MISC) e DVD oficial

Tal como acontece desde 2000 , foi editado um \textbf{CD} (MISC), com todas as quarenta e três canções que participaram no concurso. Primeiramente, as músicas foram colocadas à venda, individualmente no site oficial da \textbf{Eurovisão} (LOC), e só posteriormente (aseguir ao festival), é que foi lançado o \textbf{CD} (MISC) duplo oficial do \textbf{Festival Eurovisão da Canção 2008} (MISC). Com este \textbf{CD} (MISC), foi também lançado um DVD, com os três espetáculos gravados, assim como outras partes bónus (ensaios, etc.). Para além do \textbf{CD} (MISC) e do DVD, foi também lançada toda uma gama de produtos exclusiva e alusiva ao festival de 2008. Foram fabricados objectos que vão desde canecas, a tapetes para ratos, camisolas e t-shirts, porta-chaves, moedas comemorativas, entre muitos outros.



\textbf{Controvérsias} (PER)Uma vez finalizado o \textbf{Festival Eurovisão da Canção 2007} (MISC) , a \textbf{Sérvia} (LOC) começou a planear os preparativos para a celebração do evento do ano seguinte. Desde o começo, a \textbf{Beogradska Arena de Belgrado} (LOC) surgiu como um dos melhores sítios para ser a sede do festival, ao ser uma das maiores arenas de toda \textbf{Europa} (LOC), com uma capacidade para mais de 20.000 espectadores. Em \textbf{Maio de 2007} (MISC) , a empresa de cosméticos \textbf{Wella} (PER) anunciou que sería o patrocinador oficial do evento de 2008 juntamente com o banco \textbf{Raiffeisen Bank} (LOC) . No entanto, logo de início as criticas começaram.

Depois da edição de 2007 , diversas críticas foram realizadas ao festival pois muitos consideraram haver uma excessiva influência de factores não musicais nos resultados, como por exemplo o produzido pelos "blocos de votos" durante as votações (onde os países votavam e votam nos seus países, independentemente da qualidade musical da música do país em questão). O director da \textbf{União Europeia de Rádiodifusão} (ORG) , \textbf{Svante Stockselius} (PER) , anunciou que algumas mudanças no sistema de votação e no próprio formato do festival poderiam ser implementadas em 2009 para resolver este problema, e que na edição de 2008, manter-se-ia o esquema/formato, utilizado no concurso desde 2004 . Uma das possibilidades era implementar um sistema com duas semifinais, devido ao crescente número de competidores. Posto isto, pelos finais de \textbf{Julho} (LOC) de 2007 , colocou-se a possibilidade de se realizarem estas alterações logo para a edição de 2008, a realizar-se em \textbf{Belgrado} (LOC) , passando a haver assim duas semifinais, e só o grupo dos \textbf{Big 4} (MISC) ( \textbf{Alemanha} (LOC) , \textbf{Espanha} (LOC) , \textbf{França} (LOC) e \textbf{Reino Unido} (LOC) ), e o país anfitrião, é que teriam a oportunidade e classificariam-se directamente para a grande final. Todas estas alterações, acabaram por ser aceites dois meses mais tarde. A \textbf{28 de Setembro de 2007} (MISC) , num encontro da \textbf{EBU} (MISC) , realizado em \textbf{Verona} (LOC) , \textbf{Itália} (LOC) , os altos comissários da \textbf{organização televisiva europeia} (ORG), confirmaram a realização de duas semifinais no \textbf{Festival Eurovisão da Canção 2008} (MISC).

A 14 de Setembro de 2007 deu-se início a uma nova tradição no mundo \textbf{Eurovisivo} (MISC), quando o presidente da câmara de \textbf{Helsínquia} (LOC) entregou as "chaves do concurso" ao seu homólogo belgradês. Com isto, os preparativos para o evento començaram a adquirir maior importância e levou a um desenrolar de novas situações e notícias: a \textbf{Beogradska Arena} (PER) foi relativamente cedo confirmmada como sede e local principal para a realização do festival e \textbf{Radio-televizija Srbije} (MISC) (a televisão organizadora do festival) realizou um concurso para elegor o tema que seria executado no e durante o evento. No concurso gannhou um desenho controvertido em que constão várias imagens abstractas, representando vários elementos. Posto isto, foi eleito um tema diferente, intitulado por \textbf{Confluência dos Sons} (MISC) , referindusse à posição geográfica da cidade anfitriã na confluência dos rios \textbf{Danúbio} (LOC) e \textbf{Sava} (LOC) ; no desenho criado por \textbf{Boris Miljković} (PER) ver-se-ia representado logo no logotipo do evento, composto por linhas/rios de água a azul e roxo sobre um fundo branco (de forma similar à \textbf{Bandeira da Sérvia} (LOC) ) que chocam formando uma clave de sol . O tema da "confluência" também serviu de inspiração para o cenário principal, a qual possuía uma extensão que entrava pela "público adentro", marcada por dois rios e um conjunto de ecrãs de cristal líquido no fundo e nos lados do palco.

Para a animação do evento, a cadena televisiva organizadora designou a apresentadora \textbf{Jovana Janković} (PER) e o músico \textbf{Željko Joksimović} (PER) , o qual obteve o segundo lugar durante o \textbf{Festival Eurovisão da Canção 2004} (MISC) representando a \textbf{Sérvia e Montenegro} (LOC) e a mesma posição como como compositor durante o \textbf{Festival Eurovisão da Canção 2006} (MISC) do tema da \textbf{Bósnia e Herzegovina} (LOC) . Depois da eleição dos apresentadores, surgiu uma controvérsia devido a um possível conflito de interesses por parte de \textbf{Joksimović} (LOC), que era o compositor do tema " \textbf{Oro} (MISC) ", representante do país anfitrião.

A 17 de Fevereiro de 2008 , o território do \textbf{Kosovo} (LOC) declarou unilateralmente a sua independência em relação â \textbf{Sérvia} (LOC) , desencadeando uma série de protestos dentro do território sérvio, especialmente contra aqueles países que apoiaram a independência kosovar. Antes desta situação, a \textbf{EBU} (MISC) estudou a possibilidade de transferir a sede do concurso para outro país; dentro das melhores opções encontravam-se os três últimos organizadores dos três últimos festival: \textbf{Helsínquia} (LOC) , \textbf{Atenas} (LOC) e \textbf{Kiev} (LOC) . Finalmente, o comité organizador liderado pela \textbf{RTS} (ORG) recebeu a garantia governamental, de que todos e quaisquer visitantes e participantes do concurso, seriam protegido pelas forças da autoridade, a fim de providenciar uma boa estadia em \textbf{Belgrado} (LOC). Por isso, e devido a este factor, a \textbf{EBU} (MISC) decidiu manter o seu apoio à realização do festival na \textbf{Sérvia} (LOC).

Às possibilidades de manifestações de carácter político, somaram-se os receios de actos descriminatórios homofóbicos contra os fãs e seguidores do evento, devido a que o \textbf{Festival Eurovisão da Canção} (MISC) é considerado como um dos principais ícones gays para as comunidades \textbf{LGBT} (ORG) da \textbf{Europa} (LOC),, devido ao facto de o festival ser uma das poucas ocasiões em que um país tem a oportunidade para mostrar se os direitos dos cidadãos são respeitados ou não, e em parte por várias personalidades transexuais e homossexuais terem representado alguns países e obtido uma boa qualificação (como o caso de \textbf{Dana International} (LOC) , que venceu o \textbf{Festival Eurovisão da Canção 1998} (MISC) ). Segundo alguns mídia sérvios, estimavasse que na altura do concurso, mais de vinte mil pessoas \textbf{LGBT} (MISC) assistissem ao festival por toda a \textbf{Europa} (LOC). Diversas organizações \textbf{LGBT} (ORG) manifestaram o seu receio logo, pelos e de que grupos extremistas anunciaram que repeleriam violentamente qualquer amostra pública de homossexualidade nas ruas de \textbf{Belgrado} (LOC), de igual forma como haviam atacado noutras ocasiões anteriores marchas do orgulho gay no capital \textbf{Sérvia} (LOC). Antes desta situação, \textbf{Stockselius} (PER) afirmou que o \textbf{Presidente} (MISC) \textbf{Boris Tadić} (PER) assegurou-lhe que não haveria qualquer tipo de discriminação para os seguidores do evento, independentemente da sua orientação sexual.



Artistas repetentes

Ao longo da história, vários artistas repetiram a sua experiência eurovisiva uma ou mais vezes. Em 2008, os repetentes foram:



Festival Eurovisão da Canção 2009



Países que classificaram-se para a final \textbf{Países} (MISC) que não se classificaram para a final \textbf{Países} (MISC) que participaram no passado, mas não em 2009

O \textbf{Festival Eurovisão da Canção 2009} (MISC) (em inglês : \textbf{Eurovision Song Contest} (ORG) 2009 , em francês : \textbf{Concours Eurovision} (MISC) de la chanson 2009 e em russo : \textbf{  -2009} (MISC) ) foi a 54. edição do Festival Eurovisão da Canção , que é recebida pela \textbf{cidade de} (LOC) \textbf{Moscovo} (LOC) , na \textbf{Rússia} (LOC) . Em 2008 , com a canção " \textbf{Believe} (PER) ", de \textbf{Dima Bilan} (PER) , a \textbf{Rússia} (LOC) ganhou o \textbf{Festival} (LOC), e com isso os direitos de realizar o evento que decorreu de 12 de \textbf{Maio} (LOC), com a 1. semifinal, até 16 de Maio de 2009, com a final. Pelo meio, em 14 de \textbf{Maio} (LOC), foi realizada a 2. semifinal, no \textbf{Olimpiisky Indoor Arena} (MISC) .

O grande vencedor foi \textbf{Alexander Rybak} (PER) , com a música " \textbf{Fairytale} (PER) ", dando assim a oportunidade de a \textbf{Noruega} (LOC) organizar o \textbf{Festival Eurovisão da Canção 2010} (MISC) . \textbf{Alexander Rybak} (PER) bateu o recorde de pontos obtidos na \textbf{Festival} (LOC), com 387, e ao mesmo tempo, venceu com a maior margem de diferença para o segundo classificado. Antes do grande vencedor da final, o vencedor da 1. semifinal, foi a representante da \textbf{Islândia} (LOC) , \textbf{Yohanna} (LOC) , com a canção " \textbf{Is It True?} (MISC) ", e o vencedor da 2. semifinal foi o representante da \textbf{Noruega} (LOC) , \textbf{Alexander Rybak} (PER) , com a música " \textbf{Fairytale} (PER) ".

Mudanças no sistema de votação reintroduziram o sistema de \textbf{júris} (LOC) , que possui 50\% da votação final (os outros 50\%, são efectuados por todo e qualquer cidadão que esteja nos países que concorrem ao \textbf{Festival} (LOC)). Confirmaram a sua entrada e participaram no concurso 42 países. Com a saída da \textbf{Geórgia} (LOC), que viu a sua música desqualificada e não quis seleccionar outra, o \textbf{Festival} (LOC) acabou por ficar reduzido. A \textbf{Eslováquia} (LOC) anunciou que regressaria ao \textbf{Festival} (LOC), enquanto que \textbf{São Marino} (LOC) desistiu, devido a dificuldades financeiras.

Outros países que pareciam querer regressar no edição de 2009 eram a \textbf{Itália} (LOC) e \textbf{Marrocos} (LOC). A \textbf{Geórgia} (LOC) e a \textbf{Letónia} (LOC), afirmaram primeiramente um boicote ao concurso, devido ao conflito no \textbf{Cáucaso} (LOC) mas, mais tarde, a \textbf{European Broadcasting Union} (ORG) (\textbf{EBU} (MISC)) veio confirmar que ambos os países participariam na edição de 2009.

Pela primeira vez na \textbf{História de todos os Festival} (MISC), existiram apresentadores diferentes para as semifinais e para a \textbf{Final} (LOC). Para apresentar as semifinais, foram escolhidos o apresentador de televisão, \textbf{Andrey Malakhov} (PER) e a modelo \textbf{Natalia Vodyanova} (PER) , que vive actualmente em \textbf{Londres} (LOC) . A \textbf{Grande Final} (LOC) de 16 de Maio , foi apresentada pela personalidade mediática (conhecida por já ter feito anúncios televisivos, etc.) \textbf{Ivan Urgant} (PER) , juntamente com \textbf{Alsou} (PER) , que representou a \textbf{Rússia} (LOC) no \textbf{Eurofestival de 2000} (LOC) , em \textbf{Estocolmo} (LOC) , acabando em 2 lugar com a canção " \textbf{Solo} (PER) ", segundo lugar este que foi o melhor resultado do país até aquela altura. Os quatro apresentadores foram apresentados no dia 7 de Maio de 2009 .

O \textbf{Festival Eurovisão da Canção de 2009} (MISC), ficou ainda marcado pela tentativa da entrada do \textbf{Kosovo} (LOC) , porém, devido ao não reconhecimento de alguns países participantes no \textbf{Festival} (LOC), a entrada do \textbf{Kosovo} (LOC) tornou-se mais difícil, tendo mesmo a \textbf{Sérvia} (LOC) afirmado que, caso o \textbf{Kosovo} (LOC) entrasse no concurso, se retiraria do certame.

Outros países que demonstraram interesse em participar no \textbf{Festival} (LOC) foram o \textbf{Cazaquistão} (LOC) , o \textbf{Líbano} (LOC) e a \textbf{Escócia} (LOC) . Pela primeira vez esteve também em cima da mesa a possibilidade de extinguir os \textbf{Big 4} (MISC) ( \textbf{Reino Unido} (LOC) , \textbf{França} (LOC) , \textbf{Espanha} (LOC) e \textbf{Alemanha} (LOC) ), obrigando-os a ter que passar por uma das semifinais, tal como todos os outros países concorrentes, à excepção do país anfitrião. Porém, tal não aconteceu, e os quatro países continuam a ter a passagem directa para a final.

A 30 de \textbf{Janeiro} (LOC) de 2009 realizou-se o sorteio para determinar em qual das semifinais cada país teria que concorrer, assim como o sorteio para determinar qual a \textbf{Semifinal} (LOC) a transmitir por cada um dos cinco países que passam directamente na final (para além da transmissão, estes também poderão votar nessa mesma \textbf{Semifinal} (LOC)). Posteriormente, a 16 de \textbf{Março} (PER) , realizou-se o sorteio para decidir em que posição cada país actuaria nas semifinais, e na \textbf{Final} (LOC), assim como a ordem para as votações.

Pela primeira vez desde a sua introdução em 2007, não houve o já famoso "\textbf{Eurovision winner's tour} (MISC)", que consistia em o vencedor do \textbf{Festival} (PER) andar a viajar durante quase um ano, pela \textbf{Europa} (LOC), a divulgar a sua música vencedora. O facto de este tour ter acabado deve-se às altas despesas que o mesmo acarreta para a \textbf{EBU} (MISC) (patrocinadora de todo o tour ), que garante não ter os lucros necessários para continuar com a iniciativa.

Com a conclusão de todas as infraestruturas necessárias para a realização do evento, a 2 de Maio de 2009, os primeiros eventos na cidade de \textbf{Moscovo} (LOC) iniciaram-se também no dia 2 de Maio de 2009, a partir das 8h30 (hora de \textbf{Moscovo} (LOC)). A abertura da época eurovisiva deu-se oficialmente no dia 2 de Maio de 2009, como já se esperava. Os primeiros ensaios ocorreram na \textbf{arena} (PER), no mesmo dia, com uma grande cobertura mediática, e expectativa por parte de fãs de todo o \textbf{Mundo} (MISC), porém estes ensaios eram os das aberturas e intervalos dos três espetáculos, as delegações começaram a fazer os seus ensaios no dia 3 de \textbf{Maio de 2009} (MISC) .

A famosa \textbf{Eurovision Song Contest Opening Ceremony} (ORG) (\textbf{Cerimónia de Abertura do Festival Eurovisão da Canção} (MISC)) deu-se a \textbf{10 de Maio} (MISC) , um dia antes do início dos ensaios finais, nos quais já houve audiência no local. Para o espectáculo que antecede a semana mais importante de toda a \textbf{Eurovisão} (LOC) , estiveram presentes vários artistas de edições anteriores, assim como a presença dos responsáveis pelo \textbf{Festival} (ORG) em \textbf{Moscovo} (LOC) , da \textbf{EBU} (MISC) , etc..



\textbf{Formato} (MISC)A discussão sobre o formato do \textbf{Festival Eurovisão da Canção 2009} (MISC), ocorreu em \textbf{Atenas} (LOC) , \textbf{Grécia} (LOC) , em \textbf{Junho} (LOC) de 2008. Foi feita então uma proposta, que faria com que os países do \textbf{Big 4} (MISC) ( \textbf{França} (LOC) , \textbf{Alemanha} (LOC) , \textbf{Espanha} (LOC) e o \textbf{Reino Unido} (LOC) ) perdessem o seu lugar automático na final do concurso. Contudo, acabou por ser confirmado que os países do \textbf{Big 4} (MISC) , continuariam a estar na \textbf{Final do Festival 2009} (MISC).

Em 30 de \textbf{Janeiro} (LOC) de 2009 teve lugar o sorteio dos países que iriam participar na 1. ou \textbf{2. Semifinal} (MISC). De acordo com o sistema do \textbf{Festival de 2008} (MISC) , todos os países são separados em seis potes individuais baseados nas votações efectuadas por cada um dos \textbf{Festivais} (MISC) anteriores. O sorteio foi criado para assegurar que os países que têm uma maior probabilidade de dar pontos a outros na competição, não participem na mesma \textbf{Semifinal} (LOC) (onde só podem votar os países que participam na mesma). Depois do sorteio, foram então conhecidos os países que iriam participar em cada \textbf{Semifinal} (LOC). Mais tarde foi também feito outro sorteio, para determinar em que \textbf{Semifinal} (LOC) e em que países iriam votar o grupo dos \textbf{Big 4} (MISC) ( \textbf{Reino Unido} (LOC) , \textbf{Espanha} (LOC) , \textbf{Alemanha} (LOC) e \textbf{França} (LOC) ) e o país anfitrião ( \textbf{Rússia} (LOC) ). Os \textbf{Big 4} (MISC) e o país anfitrião já estão automaticamente qualificados para a \textbf{Final} (LOC) do concurso. O sorteio para a ordem de actuação nas semifinais e na \textbf{Final} (LOC), assim como o sorteio para a ordem na qual cada país, através de um porta-voz apresenta as suas pontuações atribuídas, ocorreram em \textbf{Março de 2009} (MISC).

Em 2009 , pela terceira vez consecutiva desde 2007 , a pequena produção realizada pela \textbf{EBU} (MISC) , em conjunto com a emissora anfitriã do concurso, intitulada de \textbf{Eurovision Countdown} (MISC) , foi para o ar, em três episódios de meia hora cada. Em cada episódio, foram apresentados todos os concorrentes, assim como reveladas algumas das surpresas guardadas até ao último momento. Os apresentadores do programa foram \textbf{Yana Churikova} (PER), que também apresentou os dois sorteios \textbf{Eurovisivos} (PER) de 2009, e que também apresentou a \textbf{Final nacional russa} (MISC) para escolher o seu representante, juntamente com \textbf{Jovan Radomir} (PER), que apresentou o \textbf{Festival} (LOC) russo de 2008.

Pela primeira vez, também o \textbf{You Tube} (MISC) , um dos sites mais visitados do mundo, aderiu ao espírito eurovisivo, transformando o seu logotipo, com símbolos da \textbf{Eurovisão 2009} (MISC). O canal da Eurovisão no \textbf{You Tube} (MISC), teve mais de 20 milhões de visitas apenas em cinco dias, tornando-o assim o segundo canal mais visitado do mundo.

Também pela primeira vez na história do festival, o evento foi ligado em directo com a \textbf{Estação Espacial Internacional} (MISC), onde os astronautas agradeceram o maravilhoso espetáculo, e deram a ordem para a \textbf{Europa} (LOC) começar a votar.

Durante o sorteio para a divisão dos países pelas duas semifinais e \textbf{Final} (MISC), a televisão anfitriã, o \textbf{Channel One Russia} (LOC) apresentou o sub-logo e o tema do \textbf{Festival 2009} (MISC). O sub-logo é baseado numa espécie de " \textbf{Fantasy Bird} (MISC) ", que pode ser utilizado com muitas cores. Tal como nos anos anteriores, o sub-logo será apresentado durante e ao lado da apresentação do logótipo genérico. O sub-logo foi desenhado pela \textbf{The Red Square Company} (ORG) . Esta foi a primeira vez, que não houve um slogan para o concurso, desde 2001 .

O palco foi desenhado pelo famoso designer nova-iorquino \textbf{John Casey} (PER) , e é baseado à volta do tema contemporâneo da vanguarda russa . \textbf{Casey} (PER), que desenhou anteriormente o palco para o \textbf{Festival Eurovisão da Canção 1997} (MISC) em \textbf{Dublin} (LOC), na \textbf{Irlanda} (LOC), esteve também envolvido nos trabalhos para as edições de 1994 e 1995 . Ele explicou que "mesmo antes dele ter trabalhado com os \textbf{Russos} (LOC) nos \textbf{TEFI Awards} (ORG) em \textbf{Moscovo} (LOC) em 1998 ele foi inspirado pela e desenhou para a arte através do período vanguardista russo, especialmente os construtivistas… Ele tentou desenhar o palco com um design teatral para o concurso, que incorpora a arte vanguardista russa numa mistura com a contemporânea, quase inteiramente feito com diferentes tipos de ecrãs \textbf{LED} (ORG)." \textbf{Casey} (PER) explicou que juntos, os vários painéis de \textbf{LED} (ORG) formam o produto final. Para além disso, grandes secções do palco têm a habilidade e a capacidade de se moverem, incluindo o portão circular central feito de ecrãs \textbf{LED} (ORG) curvados, que pode ser movido a fim de dar um efeito e um toque especial/sentimento para cada uma das 42 músicas a concurso. A construção do palco iniciou-se a 31 de \textbf{Março} (PER) , estando pronta a \textbf{2 de Maio} (LOC) , um dia antes dos primeiros ensaios ocorrerem na \textbf{arena} (PER).

Há cerca de cem anos, o químico russo \textbf{Dmitri Mendeleiev} (PER) criou a tabela periódica , que é utilizada por todo o \textbf{Mundo} (MISC). Como tal, os organizadores do \textbf{Festival Eurovisão da Canção 2009} (MISC), não quiseram deixar escapar uma das maiores invenções do seu país e criaram uma tabela periódica, com todos os 42 concorrentes deste ano, baseando-se no seu número de participações e de vitórias ao longo da história, por mais de cinquenta anos. Também é através desta tabela periódica que os designers da famosa \textbf{Green Room} (ORG) se inspiraram para decorar e iluminar a mesma, através das cores a que cada país corresponde (tal como os elementos químicos têm uma cor na tabela periódica).

Pela primeira vez no \textbf{Festival} (LOC) a introdução do país a actuar não foi feita através de um postal onde se mostram imagens de cada país a concurso, do país anfitrião (ou da cidade), nem através de mensagens ou actos artísticos (como em \textbf{Belgrado} (LOC) , no último ano). Os cartões postais contaram com imagens feitas a computador, que representaram vários símbolos nacionais de cada país, porém através de objectos (predominantemente musicais). Exemplo disso, é o \textbf{Arco do Triunfo} (LOC) , em \textbf{Paris} (LOC) , que aparecerá recriado por um sistema de colunas de som. Para além desta nova forma de apresentar cada país, a \textbf{Rússia} (LOC) também inseriu nestes cartões postais, com uma duração de aproximadamente 30 segundos, a modelo e \textbf{Miss Mundo 2008} (MISC) (e anteriormente \textbf{Miss Rússia} (PER) 2007 ), \textbf{Ksenia Sukhinova} (PER) , que aparecerá nos ecrãs de todo o \textbf{Mundo} (MISC), com as cores da bandeira de cada país, e 42 penteados e estilos diferentes.

Os famosos postcards , foram então constituidos por:

Também do lado de fora do pavilhão foram colocados vários "objectos" alusivos ao evento. Várias flores (e relva) foram colocadas no exterior da arena, formando a palavra \textbf{Eurovision} (MISC) , e dando assim um ar mais verde e ecológico no exterior. À volta de todo o edifício da arena, foram colocados vários painéis gigantes, que circundam todo o estádio, e formam um gigantesco anel, com as imagens da \textbf{Miss Mundo} (ORG), com as cores representativas de cada país (imagens essas que também aparecem nos cartões postais).

Pela televisão russa passou também um anúncio televisivo, com uma duração de cerca de 30 segundos, alusivo ao \textbf{Festival} (LOC), indicando as datas das semifinais e da \textbf{Final} (LOC), e os principais patrocinadores do evento. A principal imagem foi a modelo que "vestiu" as 42 nações, que aparece apenas "vestida" com a \textbf{Rússia} (LOC). Também por toda a cidade foram colocados painéis publicitários, cada um com a \textbf{Miss Mundo} (ORG) caracterizada de cada nação, com a bandeira da mesma por detrás.



\textbf{Local} (MISC)O \textbf{Festival Eurovisão da Canção} (MISC) 2009 será realizado na \textbf{Rússia} (LOC) , depois da sua vitória no \textbf{Festival Eurovisão da Canção 2008} (MISC) em \textbf{Belgrado} (LOC) , \textbf{Sérvia} (LOC) , com a canção de \textbf{Dima Bilan} (PER) , " \textbf{Believe} (LOC) ". \textbf{Vladimir Putin} (PER) , \textbf{Primeiro Ministro da Rússia} (MISC) , constatou que o \textbf{Festival} (LOC) seria realizado na capital russa, \textbf{Moscovo} (LOC) , e foi proposto pela televisão responsável pelo concurso, o \textbf{Channel One Russia} (LOC) , que o \textbf{Festival} (LOC) se realizasse no \textbf{Olimpiisky Indoor Arena} (MISC) , o antigo pavilhão \textbf{Olímpico de Moscovo} (MISC)). Esta proposta foi avaliada pela \textbf{European Broadcasting Union} (ORG) (\textbf{EBU} (MISC)) , e confirmada a \textbf{13 de Setembro de 2008} (MISC) . O responsável pelo edifício, \textbf{Vladimir Churilin} (PER), veio a público refutar rumores de que era necessária uma reconstrução urgente da infra-estrutura, dizendo que: "\textbf{It will not be required for the Eurovision Song Contest} (MISC). \textbf{We} (MISC) now can take up to 25 thousand spectators" ("Não será exigido para o \textbf{Festival Eurovisão da Canção} (MISC). Neste momento nós podemos acomodar 25 mil espectadores"). No entanto, o estádio recebeu obras de melhoramento, num valor aproximado de 10 milhões de euros, [ carece de fontes ? ] tendo uma capacidade para cerca de 25 000 espectadores, e entre 3 mil a 5 mil lugares, exclusivos para jornalistas de todo o \textbf{Mundo} (LOC), perfazendo assim um total de 28 a 30 mil lugares para o \textbf{Festival de 2009} (MISC). Esta é, ao mesmo tempo, a segunda maior arena do \textbf{Festival} (LOC) na história, tendo sido apenas ultrapassada em 2001 , pelo \textbf{Parken Stadium} (LOC) , onde cerca de de 40 mil espectadores assistiram ao concurso ao vivo. No fim, a remodelação do pavilhão, levou à substituição total de todos os assentos do mesmo, à criação de 24 novas salas de maquilhagem, uma nova sala vip, uma \textbf{Green Room} (ORG) completamente remodelada e uma nova sala de conferências, com trezentos computadores, instalados em cento e sessenta secretárias.

A construção do palco e de todas as infra-estruturas relacionadas com o \textbf{Festival} (LOC), iniciou-se cinquenta dias antes da realização do mesmo. A 31 de \textbf{Março} (PER) de 2009 , os primeiros camiões chegaram ao \textbf{Olimpiisky Indoor Arena} (MISC) , onde a construção começou de imediato. O final das obras aconteceu a 2 de Maio de 2009, perto do início dos ensaios por parte dos quarenta e dois países no próprio palco. Enquanto a construção decorreu, vários engenheiros e peritos informáticos criaram os mais de quarenta efeitos de luz, som, \textbf{pirotecnia} (LOC), etc., para que tudo estivesse pronto até ao \textbf{Festival} (LOC). Em 2 de Maio a última peça do palco para o \textbf{Festival Eurovisão da Canção} (MISC) 2009 foi colocada, e no dia seguinte começaram os ensaios.

Para além do \textbf{Olimpiisky Indoor Arena} (MISC), a organização do \textbf{Festival} (LOC) utilizou a \textbf{Praça Vermelha} (LOC) como local ao ar livre para comemoração do \textbf{Festival} (MISC) (como já vem sendo hábito), assim como um grande edifício ao lado da mesma, onde ficou instalado o \textbf{Euroclub} (ORG) e se realizaram concertos e festas. Esse grande edifício, situado ao lado da mais famosa \textbf{praça russa} (LOC), conhecida por \textbf{Manege de Moscovo} (LOC) , foi inaugurado como \textbf{Euroclub} (LOC) a 6 de \textbf{Maio de 2009} (MISC). Este é o maior espaço onde o evento já foi realizado. O espaço esteve aberto das 18h00 até às 03h00 do dia seguinte, e por vezes o horário era prolongado até às 06h00. Pelo espaço passaram vários artistas que já participaram na \textbf{Eurovisão} (LOC), e realizaram-se várias conferências de imprensa, festas, etc..

\textbf{O Ministério do Interior Russo} (MISC), representado pelo deputado chefe do ministério do departamento de segurança, \textbf{Major General Nikolai Trifonov} (MISC) anunciou a 28 de Abril de 2009 , que mais de 8500 polícias e seguranças, circularão pela \textbf{Cidade de Moscovo} (LOC) , durante o tempo do festival. A segurança será reforçada principalmente junto ao \textbf{Olimpiisky Indoor Arena} (MISC) , e ao pé dos principais hotéis (onze) onde ficarão instaladas as claques de apoio, assim como as próprias comitivas dos países (que podem ultrapassar as 20 pessoas apenas num país). No comunicado também foi expresso que o \textbf{Ministério} (LOC) não espera nenhum ataque terrorista, mas que prefere não correr o risco de um ataque vir a ocorrer durante o festival, assim como algum protesto que ponha em causa o evento. A segurança também foi reforçada, devido às datas do festival coincidirem com dois dos principais feriados da \textbf{Rússia} (LOC) , o \textbf{Dia da Primavera} (MISC) e do \textbf{Trabalhador} (ORG) a 1 de \textbf{Maio} (MISC) , e o \textbf{Dia da Vitória} (MISC), a \textbf{9 de Maio} (LOC) . Ambos os feriados atraem gigantescos números de pessoas à cidade, assim como um acréscimo nas manifestações e protestos. Para além do grande reforço policial, na entrada da arena/pavilhão onde se realiza o festival, existirão vários detectores de metais, e pela primeira vez serão feitos testes ao nível de álcool no sangue. Qualquer tipo de bebida alcoólica é proibida no local do festival, mesmo em que reduzidas quantidades. Outra das muitas medidas de segurança tomadas pelas autoridades russas, é o corte do trânsito ao pé da arena, nas ruas \textbf{Samarskaya} (LOC), \textbf{Durova} (LOC) e \textbf{Shchepkina} (LOC), desde o dia 11 de Maio até ao dia 16 do mesmo mês. Estas ruas sã conhecidas pelos seus grandes engarrafamentos, e a fim de evitar atrasos e possíveis problemas em nível de segurança, as autoridades resolveram cortar o trânsito por seis dias (o que não é muito normal em \textbf{Moscovo} (LOC).

As equipas especiais russas evitaram um ataque terrorista, durante a primeira semifinal do Festival Eurovisão da Canção 2009. A notícia foi avançada a 14 de Maio de 2009 , pela site da mídia russa \textbf{Life.ru} (MISC) . Empregados da \textbf{Agência Federal de Segurança Russa} (ORG) descobriram no interior do estádio olímpico, onde o festival estava a decorrer, um esconderijo onde se encontrava um pacote com seis balas de 9 mm. A descoberta ocorreu horas antes de a semifinal ser posta no "ar", para milhões de telespectadores. Os agentes federais disseram ainda que consideravam que as balas eram destinadas a ser utilizadas num atentado contra a vida de um dos intérpretes de um dos 42 países a concurso (18 naquele dia). O artista apontado foi o representante da \textbf{Bulgária} (LOC) , porém ainda não foi dito por fontes oficiais a quem era dirigido o ataque.

A \textbf{cidade de Moscovo} (LOC) dispõem de um gigantesco sistema de metropolitano . O \textbf{Metropolitano de Moscovo} (LOC) chega a quase todas as partes da cidade. Nos três dias do festival, o horário do metropolitano será alargado até às 3h30m da manhã, o que acarreta um custo de 10 milhões de rublos para o sistema de transportes moscovita. O \textbf{Metropolitano de Moscovo} (LOC) foi inaugurado em 1935 , sendo o maior do mundo por densidade de passageiros, transportando por volta de 3340 milhões de pessoas por ano e cerca de 9,2 milhões de pessoas por dia. A sua rede é composta por 173 estações, distribuídas por 12 linhas através de 282,4 km (sendo assim o sexto mais extenso do mundo). Apesar do metropolitano ser bastante utilizado, os fãs eurovisivos foram aconselhados a não utilizar o mesmo, por este ser grande e confuso, onde uma pessoa facilmente se perde. Para facilitar uma pouco a "vida" de todos os visitantes, a \textbf{Câmara de Moscovo} (ORG), colocou em todas as ruas da capital, placas com os nomes das ruas no alfabeto latino , assim como o sistema de metropolitano fez no nome das estações.

Vários autocarros foram alugados e postos ao serviço do festival. Depois de terem "sofrido" uma pequena remodelação no seu exterior, colocando-se assim imagens elucidativas ao festival, os autocarros estão ao serviço de todos, com rotas desde o \textbf{Olimpiisky Indoor Arena} (MISC) até ao \textbf{Euro Club} (MISC) , passando por todos os hotéis oficiais e reconhecidos pela organização do concurso.

Para transportar os artistas e as suas delegações a organização do Festival Eurovisão da Canção 2009, alugou várias camionetas, carros, etc, tudo com os símbolo do festival a percorrer todo o seu comprimento. Sempre que uma delegação executa uma viagem terá sempre polícia a acompanhá-la, de modo a evitar problemas de maior.



Votação

Em resposta a algumas televisões responsáveis pela divulgação do concurso nos países europeus, que continuavam a queixar-se do voto político, vizinho e diáspora, a \textbf{EBU} (MISC) avaliou o sistema de votação usado no concurso, com a possibilidade de fazer uma alteração no sistema de votação, já para o \textbf{Festival Eurovisão da Canção 2009} (MISC). Os organizadores do festival enviaram então, um questionário sobre o sistema de votação às televisões difusoras, e um grupo de referência incorporou as respostas dos questionários numa sugestão para a alteração do sistema de votação no próximo ano. A \textbf{Telewizja Polska} (MISC) (\textbf{TVP} (ORG)) , da \textbf{Polónia} (LOC), propôs que se adicionasse o júri internacional, similar ao que foi usado no \textbf{Festival Eurovisão da Dança 2008} (MISC) no \textbf{Festival Eurovisão da Canção} (MISC) , para reduzir o impacto do chamado block voting e avaliar a qualidade do artista e da sua música. Desde então, foi confirmado que o voto seria feito através do sistema 50 por 50, em que metade dos votos seria atribuído pelo televoto e outra metade pelo júri nacional. O método de selecção dos países da semifinal para a final, será o mesmo, contudo, apenas nove países passarão à próxima fase através de televoto, dado que o décimo finalista será escolhido através de um sistema de jurados. Pormenores sobre o método ficaram combinados após reunião do "\textbf{Reference Group} (MISC)" em Dezembro de 2008. Os júris nacionais, foram originalmente implantados no \textbf{Festival Eurovisão da Canção} (MISC), em 1997 . Porém, apenas seis anos depois, em 2003 , a votação passou a ser realizada através de 100\% televoto.

\textbf{Edgar Böhm} (PER), director de entretenimento do canal público \textbf{Österreichischer Rundfunk} (MISC) (\textbf{ORF} (ORG)), da \textbf{Áustria} (LOC) , constatou que o formato de 2008, com duas semifinais "continuaria a incorporar uma mistura de países que iriam certamente ser politicamente favorecidos no sistema de votação", e "que, a menos que um guia de instruções simples e claro sobre como estão organizadas as semifinais feitas pela \textbf{EBU} (MISC), a \textbf{Áustria} (LOC) não participará no \textbf{Festival de Moscovo} (MISC) 2009 ". Apesar da inclusão de um júri na final, a \textbf{Áustria} (LOC) não participará no festival, mas emitirá o mesmo no seu país, tal como em 2008.

As novas regras de votação para o \textbf{Festival Eurovisão da Canção} (MISC), são as seguintes:

Para assegurar que as regras são cumpridas, a \textbf{EBU} (MISC) enviará para cada país, um auditor externo, que monitorizará tudo o que ocorrer em relação ao festival.

No fim, depois da votação do público e dos júris, os 12 países com mais votos, passam a ter uma pontuação de 1 a 12 respeitando os votos anteriores. No fim soma-se ambos os valores que vão de 1 a 12 nas pontuações, e volta-se a pôr por ordem de pontuações determinando assim o grande vencedor. Para se ter uma ideia melhor de como será o sistema de votação, tomem-se os exemplos em baixo:

\textbf{Predefinição:Tabelas de Exemplo} (MISC) do \textbf{Novo Sistema de Votação do Festival Eurovisão da Canção} (LOC)

Em caso de dois ou mais países terem igual número de pontos no resultado definitivo, tanto na final como na semifinal, utiliza-se como regra de desempate o número de votos entregues com o máximo de pontuação. Se dois ou mais países tiverem igual o número de votos com 12 pontos, então analisa-se o número de votos com 10 pontos, e assim sucessivamente até acabar com o empate.

A edição de 2009, terá pela primeira vez desde 2002, um grupo de jurados, no entanto, desta vez o grupo de jurados será o maior de sempre. Um total de 210 jurados irão julgar as entradas do Festival Eurovisão da Canção 2009 (5 jurados por país). Os júris devem consistir na maior variedade possível de profissões relacionadas com o mundo do espetáculo (cantores, produtores, etc), e deverá ser diversificado em género, sexo, idade, estatuto social, etc. Todos os membros dos \textbf{júris} (LOC) têm que ser cidadãos do país em que são os júris, e não podem estar ligados de qualquer forma aos representantes do seu país. Caso contrário, o seu voto pode ser anulado.

Tal como em todos os anos, desde que o festival se rege por um sistema de televoto, cada televisão anfitriã nacional, terá que escolher um porta-voz, que divulgará os votos do júri e do televoto do seu país. As pontuações de 1 a 7 apareceram automaticamente no quadro de votações, e posteriormente, o porta voz anunciará os países que receberam os 8, 10 e 12 pontos do seu país. Os resultados do \textbf{Festival Eurovisão da Canção} (MISC), serão posteriormente colocados no site oficial, pelo \textbf{Executivo Superior} (LOC) da \textbf{EBU} (MISC). Os resultados serão dados a conhecer apenas com as pontuações do público, do júri, e no fim dos dois juntos. A lista das pessoas que transmitiram os pontos atribuídos pelo seu país e pelo júri, aos outros países participantes no \textbf{Festival Eurovisão da Canção} (MISC) 2009, de cada país, pode ser observada na ligação no início desta secção (porta-vozes dos júris e televoto).



Países classificados para a final

A possibilidade dos quatro países ( \textbf{França} (LOC) , \textbf{Reino Unido} (LOC) , \textbf{Espanha} (LOC) e \textbf{Alemanha} (LOC) ) passarem directamente para a final foi questionada tendo mesmo sido notificado a obrigatoriedade destes países passar por uma das semifinais, porém a 14 de Setembro de 2008 , os altos comissários da \textbf{EBU} (MISC), vieram a público dar a conhecer a sua decisão na sua última reunião: que os \textbf{Big 4} (MISC) (\textbf{Reino Unido} (LOC), \textbf{Espanha} (LOC), \textbf{França} (LOC) e \textbf{Alemanha} (LOC)), não perderiam o seu privilégio de não passarem por uma das semifinais, mas sim diretamente para a final.



\textbf{Participações individuais} (MISC)Durante cerca de dez meses, todos os países elegeram os seus representantes, assim como as músicas que os mesmos viriam a interpretar em \textbf{Moscovo} (LOC). Para realizar tal selecção, cada país utilizou o seu próprio processo de selecção. Alguns optaram pela selecção interna, que consiste em a televisão organizadora daquele país, é quem faz a escolha; no entanto, por vezes apenas o artista é seleccionado internamente, e a música não. Outros países (a maioria), utiliza um programa de televisão para seleccionar a sua entrada. Quartos de final, semifinais, second-chances e finais, foram realizadas durante dez meses na maioria dos países europeus, cada um com o seu processo próprio (também a internet foi utilizada na fase de escolhas). No conjunto de artigos em baixo, é possível ler mais ao pomenor o tipo de processo que cada país utilizou, assim como os resultados e reacções.



\textbf{Participantes} (MISC)Seguindo a publicação da lista de todos os países participantes pela \textbf{EBU} (MISC), quarenta e três países confirmaram a sua participação no Festival Eurovisão da Canção , incluindo a \textbf{Eslováquia} (LOC) , que regressa ao concurso, após onze anos consecutivos de ausência. Mais tarde o número de concorrentes baixou para os quarenta e dois. A \textbf{Geórgia} (LOC) anunciou originalmente que não participaria no festival, como consequência da \textbf{Guerra na Ossétia do Sul} (MISC) , em protesto contra os polícias estrangeiros da \textbf{Rússia} (LOC), mas desde o fim da guerra, a \textbf{Geórgia} (LOC) resolveu não "faltar" ao festival, decisão que também foi inspirada com a sua vistória no \textbf{Festival Eurovisão da Canção Júnior 2008} (MISC) , e onde a \textbf{Rússia} (LOC) também atribuíu 12 pontos respectivamente à \textbf{Geórgia} (LOC) (a pontuação mais alta que um país pode dar a outro). \textbf{Svante Stockselius} (PER) , o fiscal da \textbf{EBU} (MISC) para a \textbf{Eurovisão} (LOC), disse que um número recorde de países estaria a concorrer em \textbf{Moscovo} (LOC), em 2009, dizendo que 44 ou mais países entrariam no concurso, no entanto isto acabou por não ocorrer, e o número de entradas no \textbf{Festival Eurovisão da Canção 2009} (MISC), é igual ao número de entradas do festival de 2008.

Surgiram rumores acerca da participação e do retorno de \textbf{São Marino} (LOC) e do \textbf{Mónaco} (LOC) . A \textbf{Télé Monte Carlo} (LOC) (\textbf{TMC} (ORG)), a televisão oficial do \textbf{Mónaco} (LOC), confirmou mesmo que estariam a decorrer conversações entre eles e a \textbf{EBU} (MISC), para o regresso do \textbf{Mónaco} (LOC) ao festival em 2009. Ao mesmo tempo, surgiram rumores de que \textbf{São Marino} (LOC) , ou melhor, a sua televisão oficial, a \textbf{Radiotelevisione della Repubblica di San Marino} (PER) (\textbf{SMRTV} (MISC)), iria retirar-se do festival em 2009, face ao mau resultado no último ano. No fim, depois de ter confirmado a sua intenção de participar em \textbf{Moscovo} (LOC), a \textbf{SMRTV} (MISC) foi forçada a desistir do evento, devido a dificuldades financeiras que puseram em causa uma segunda participação.

Durante este período também se falou da possibilidade do ingresso de novos países. Depois de vários anos sem uma cadeia de radiodifusão própria, a 1FLTV foi criada no \textbf{Liechtenstein} (LOC) , a qual está qualificada para entrar na \textbf{EBU} (MISC) , requisito fundamental para se poder participar na \textbf{Eurovisão} (LOC) , algo que já tinha causado intenções no pequeno principado mais do que uma vez durante vários anos.

A televisão oficial da \textbf{Letónia} (LOC), a \textbf{LTV} (ORG), a 17 de Dezembro de 2008, afirmou a sua retirada do Festival Eurovisão da Canção, três dias depois da participação final estar deliniada. Isso aconteceu devido a cortes orçamentais de mais de 2 milhões de \textbf{Lats} (PER) no orçamento da \textbf{LTV} (ORG), prejudicando assim, a sua capacidade de pagar as propinas para a participação no concurso. A \textbf{LTV} (ORG) confirmou que havia informado a \textbf{EBU} (MISC) da sua intenção de desistir, baseada apenas em dificuldades financeiras. A \textbf{LTV} (ORG) esteve em discussão com a \textbf{EBU} (MISC) numa reunião, para encontrar então uma solução para manter o país no Festival Eurovisão da Canção. A 20 de Dezembro de 2008, a \textbf{LTV} (ORG) anunciou por fim, a sua retirada oficial do festival, e que tanto a \textbf{EBU} (MISC) como o \textbf{Channel One} (LOC), concordaram em não forçar uma penalização mais tarde pela sua desistência como televisão emissora do festival de 2009. A \textbf{LTV} (ORG), aproveitou também, para expressar a sua intenção de participar no \textbf{Festival Eurovisão da Canção 2010} (MISC) . Contudo, a 12 de \textbf{Janeiro} (LOC) de 2009 foi confirmado que a \textbf{Letónia} (LOC) participaria no festival de 2009. A 9 de \textbf{Março} (PER) de 2009, a \textbf{EBU} (MISC) resolveu desqualificar a música georgiana, dando a hipótese à \textbf{Geórgia} (LOC) de escolher outra, no entanto oo país não aceitou isso, e resolveu boicotar a edição de 2009, a 11 de \textbf{Março} (PER) de 2009, onde faria a sua terceira aparição no festival.

O primeiro artista seleccionado para participar no festival foi \textbf{Sakis Rouvas} (PER) , eleito pela cadeia televisiva ERT como representante da \textbf{Grécia} (LOC) , no dia 15 de \textbf{Julho} (LOC) de 2008 . O cantor, que obteve o terceiro lugar na \textbf{Eurovisão 2004} (MISC) apresentou a música " \textbf{This is our night} (MISC) " para a \textbf{Eurovisão} (LOC), a ser eleita por populares entre mais outras duas músicas. Dentro dos artistas seleccionados, destaca-se a participação de grandes figuras de diversos meios sociais e incluindo alguns artistas já de renome europeu, especialmente entre os países da \textbf{Europa Ocidental} (LOC), que nos últimos anos tinham tido artistas de curta trajectória. É neste grupo que se encontram casos como o de \textbf{Sakis Rouvas} (PER), a belga de origem turca \textbf{Hadise} (LOC) , a francesa \textbf{Patricia Kaas} (PER) e o grupo holandês \textbf{De Toppers} (MISC) , para além de estrelas originárias de reality shows como a espanhola \textbf{Soraya Arnelas} (PER) e a ucraniana \textbf{Anastasiya Prykhodko} (ORG) , representando esta última o país anfitrião . Novos talentos surgiram durante os processos de selecção, como o cantor norueguês \textbf{Alexander Rybak} (PER) que conseguiu seis vezes mais votos que o adversário que ficou em segundo lugar, durante o famoso \textbf{Melodi Grand Prix} (MISC) . \textbf{Alexander Rybak} (PER) é dado, desde que foi eleito para representar a \textbf{Notuega} (LOC) , como o grande vencedor de 2009 , estando sempre em primeiro na grande maioria das votações organizadas por vários sites; ou ainda o caso de \textbf{Jade Ewen} (PER) , representante do \textbf{Reino Unido} (LOC) . Outras figuras com grande reconhecimento a nível internacional na música, participaram (de forma indirecta), como \textbf{Sir Andrew Lloyd Webber} (PER) , que organizou todo o processo de selecção no \textbf{Reino Unido} (LOC) e compôs a canção para o país, intitulada de " \textbf{My time} (MISC) "; \textbf{Ronan Keating} (PER) que escreveu a letra de " \textbf{Believe} (LOC) again " (interpretada por \textbf{Niels Brinck} (PER) , proveniente da \textbf{Dinamarca} (LOC) ) e o sueco-iraniano \textbf{Arash} (PER) , que cantará junto à estreante \textbf{Aysel Teymurzadeh} (PER) o tema representante do \textbf{Azerbaijão} (LOC) . A selecção do duo isrealita composto por \textbf{Noa} (PER) e \textbf{Mira Awad} (LOC) no meio do conflito da \textbf{Faixa de Gaza} (LOC) gerou grande impacto mediático devido à origem árabe de \textbf{Awad} (LOC), que é a primeira representante israelita na \textbf{Eurovisão} (LOC) proveniente da dita comunidade.

Na tabela em baixo, é possível observar (por ordem alfabética), as participações de cada país (a sua música, artista, evento de selecção do representante e data do mesmo). Para mais informações sobre cada país no Festival Eurovisão da Canção 2009, basta seguir a ligação efectuada pelo nome dos países na tabela, ou na secção de participações individuais mais acima neste artigo.



\textbf{Organização} (ORG)A 30 de \textbf{Janeiro} (LOC) de 2009 , o sorteio para decidir que países participariam na primeira ou na segunda semifinal, teve lugar no \textbf{Hotel Marriott de Moscovo} (PER) , em \textbf{Moscovo} (LOC) , \textbf{Rússia} (LOC) . O sorteio esteve a cargo da apresentadora russa \textbf{Yana Churikova} (PER). Baseando-se no mesmo sistema de eventos do ano passado , todos os países (trinta e oito ao todo) foram separados em seis potes individuais baseados no seu historial de votação de concursos anteriores e na sua localização geográfica. O sorteio foi implementado (ou criado), para garantir que os países que estão dispostos a dar a sua pontuação uns aos outros na competição, não participem na mesma semifinal. Destes seis potes, saíram então as listas dos países que participariam na primeira e na segunda semifinal respectivamente, ficando assim todos os países divididos em dois grupos, da maneira mais justa possível. Ficou também determinado que a \textbf{Alemanha} (LOC), \textbf{Espanha} (LOC) e o \textbf{Reino Unido} (LOC) transmitiriam e votariam na primeira semifinal, enquanto que a \textbf{França} (LOC) e a \textbf{Rússia} (LOC), transmitiriam e votariam na segunda semifinal. Para tal distribuição de transmissões e de votos, foi criado um sétimo pote, apenas com estes cinco países, e apenas com o intuito de definir que semifinal é que teriam que transmitir, e poderiam votar.

Uma das regras, que também ficou estabelecida, desde a criação de duas semifinais, é que cada país é obrigado a transmitir a final em que ficou colocado, e pode se quiser, transmitir, e apenas isso, a outra semifinal. O sorteio para a ordem de actuação nas semifinais e na final, e ainda a ordem de votação de cada país, apenas acontecerá em \textbf{Março de 2009} (MISC). Em baixo estão todos os seis potes, com os respectivos países que integraram cada um, para o sorteio da distribuição das semifinais:

Após o sorteio das semifinais, a 30 de \textbf{Janeiro} (LOC) de 2009 , foi possível efectuar a divisão de trinta e oito países em duas semifinais. Em baixo, encontra-se uma tabela, onde estão todos os países que foram a sorteio, e encontram-se ao mesmo tempo, pela ordem de que saíram dos potes onde foram divididos posteriormente. Por baixo ainda, encontram-se os cinco países com passaporte directo para a final, que terão que transmitir, e poderão votar numa das semifinais. A sua ordem, está colocada, pela forma como "sairam" dos potes, no sorteio das semifinais. A 28 de Abril a \textbf{EBU} (MISC) aceitou/aprovou o pedido que a \textbf{Espanha} (LOC) tinha feito, que consistia na troca da transmissão e votação da semifinal. Assim, e ao contrário do que inicialmente estava definido, a \textbf{Espanha} (LOC) transmitiu a segunda semifinal, e votou na mesma, tendo assim mais tempo para promover o próprio festival em si (o que a \textbf{EBU} (MISC) achou por bem) Tal troca deveu-se ao facto de a \textbf{Espanha} (LOC) se encontrar a transmitir vários espetáculos relacionados com a \textbf{Eurovisão} (LOC), e esses mesmos espetáculos têm a grande final agendada para o dia da 1. semifinal, para além disso, desta forma a \textbf{Espanha} (LOC) ganha mais algum tempo para promover a sua música. No entanto, também foi afirmado de que a \textbf{Espanha} (LOC) teria pedido para trocar o dia de transmissão, devido ao debate anual que se realiza no \textbf{Parlamento Espanhol} (MISC) todos os anos, debate o qual, a \textbf{TVE} (ORG) é a responsável por emitir para toda a \textbf{Espanha} (LOC). Passados dois dias da aprovação por parte da \textbf{EBU} (MISC), vários países criticaram tal acção, principalmente \textbf{Andorra} (LOC), que diz ser bastante prejudicada sem a \textbf{Espanha} (LOC) a votar na semifinal em qual o país entra, porém a representante de \textbf{Andorra} (LOC) diz-se confiante, e afirma mesmo que o país avançará para a grande final

A ordem de entrada de cada país apenas foi conhecida depois da divisão dos países pelas duas semifinais, o sorteio para a ordem de entrada dos países, foi efectuado a 16 de \textbf{Março} (PER) de 2009 . O segundo sorteio eurovisivo esteve a cargo da apresentadora russa \textbf{Yana Churikova} (PER), de um famoso comediante russo e do vencedor da última edição do festival, \textbf{Dima Bilan} (PER) . Os resultados podem ser vistos nas tabelas das semifinais e das finais. Para determinar a ordem das actuações na final, foram levados a cabo mais dois sorteios, logo após cada uma das semifinais, com os dez países que avançam para a final.



Ensaios

Com a conclusão da construção do palco para o \textbf{Festival Eurovisão da Canção 2009} (MISC), a 2 de Maio de 2009 , no dia a seguir, dia 3 de \textbf{Maio} (PER) , a época \textbf{Eurovisiva} (PER) foi inaugurada na cidade de \textbf{Moscovo} (LOC) . Cada país teve a oportunidade de ensaiar três vezes, e os que passaram à final ensaiaram mais uma vez, desta vez para a final. Os ensaios iniciaram-se logo no mesmo dia, sendo os primeiros países a ensaiar, \textbf{Montenegro} (LOC) , \textbf{República Checa} (LOC) , \textbf{Bélgica} (LOC) , \textbf{Bielorrússia} (LOC) , \textbf{Suécia} (LOC) , \textbf{Arménia} (LOC) , \textbf{Andorra} (LOC) , \textbf{Suíça} (LOC) e \textbf{Turquia} (LOC) (aqui dados pela ordem de actuação/ensaio). Ao todo, foram nove os países que realizaram o seu primeiro ensaio e a sua primeira conferência de imprensa. Com o final do primeiro dia, todos os fãs por toda a \textbf{Europa} (LOC) e por todo o \textbf{Mundo} (MISC), puderam ver e rever como seria a actuação de cada um dos nove países, incluindo o cenário que cada concorrente terá em palco no próprio festival. No segundo dia de ensaios foram outros nove os países que pisaram o palco moscovita. \textbf{Israel} (LOC) , \textbf{Bulgária} (LOC) , \textbf{Islândia} (LOC) , \textbf{Macedónia} (LOC) , \textbf{Roménia} (LOC) , \textbf{Finlândia} (LOC) , \textbf{Portugal} (LOC) , \textbf{Malta} (LOC) e \textbf{Bósnia e Herzegovina} (LOC) , foram os países que ensaiaram no dia 4 de \textbf{Maio} (MISC) , e também no mesmo dia deram as suas primeiras entrevistas. No terceiro dia, começaram os ensaios dos países que concorreriam na segunda semifinal, \textbf{Croácia} (LOC) , \textbf{Irlanda} (LOC) , \textbf{Letónia} (LOC), \textbf{Sérvia} (LOC) , \textbf{Polónia} (LOC) , \textbf{Noruega} (LOC) , \textbf{Chipre} (LOC) , \textbf{Eslováquia} (LOC) e \textbf{Dinamarca} (LOC) foram os países que realizaram o primeiro ensaio para cada um. O terceiro dia também foi o que mais atenção suscitou um pouco por todo o lado, devido ao ensaio da \textbf{Noruega} (LOC), país este que desde a escolha do seu representante é tido vencedor da \textbf{Eurovisão 2009} (MISC). No dia seguinte, 6 de \textbf{Maio de 2009} (MISC) , foi a vez da \textbf{Eslovénia} (LOC) , \textbf{Hungria} (LOC) , \textbf{Azerbaijão} (LOC) , \textbf{Grécia} (LOC) , \textbf{Lituânia} (LOC), \textbf{Moldávia} (LOC) , \textbf{Albânia} (LOC) , \textbf{Ucrânia} (LOC) , \textbf{Estónia} (LOC) e \textbf{Holanda} (LOC) realizarem o seu ensaio. Com o fim do quarto dia de ensaios, concluiram-se todos os primeiros ensaios de todos os países que actuarão na primeira semifinal. Durante os ensaios, vários sites (oiktimes, estoday, eurovision.tv, etc), realizaram vários arrtigos que acompanharam os ensaios hora a hora, e realizando sempre ao fim do dia um apanhado geral do que aconteceu na arena onde se realizará o evento. Nos 5., 6., 7. e 8. dias de ensaios, todos os países voltaram a ensaiar pela segunda vez no palco, já com as suas indicações para o melhoramento de erros encontrados no primeiro ensaio. No 5. dia foi a vez dos seguintes países ensaiarem pela segunda vez: \textbf{Montenegro} (LOC) , \textbf{República Checa} (LOC) , \textbf{Bélgica} (LOC) , \textbf{Bielorrússia} (LOC) , \textbf{Suécia} (LOC) , \textbf{Eslováquia} (LOC) , \textbf{Andorra} (LOC) , \textbf{Suíça} (LOC) , \textbf{Turquia} (LOC) , \textbf{Israel} (LOC) , \textbf{Bulgária} (LOC) , \textbf{Islândia} (LOC) e \textbf{Macedónia} (LOC) . No 6 dia de ensaios, dia 8 de \textbf{Maio} (LOC) , foi a vez de os últimos sete países da 1. semifinal e os seis primeiros da segunda ensaiarem. Ao todo, os treze países são: \textbf{Roménia} (LOC) , \textbf{Finlândia} (LOC) , \textbf{Portugal} (LOC) , \textbf{Malta} (LOC) , \textbf{Bósnia e Herzegovina} (LOC) , \textbf{Croácia} (LOC) , \textbf{Irlanda} (LOC) , \textbf{Letónia} (LOC), \textbf{Sérvia} (LOC) , \textbf{Polónia} (LOC) , \textbf{Noruega} (LOC) , \textbf{Chipre} (LOC) e a \textbf{Arménia} (LOC) .

Para o sétimo dia, ficaram guardados alguns dos restantes treze países a concurso na semifinal, e pela primeira vez, os cinco países que estão directamente na final, ensaiaram. \textbf{Dinamarca} (LOC) , \textbf{Eslovénia} (LOC) , \textbf{Hungria} (LOC) , \textbf{Azerbaijão} (LOC) , \textbf{Grécia} (LOC) , \textbf{França} (LOC) , \textbf{Rússia} (LOC) , \textbf{Alemanha} (LOC) , \textbf{Reino Unido} (LOC) e \textbf{Espanha} (LOC) , são os dez países que ensaiaram no dia \textbf{9 de Maio} (LOC) . O dia 9 de Maio, sétimo dia de ensaios, fica ainda marcado, pela visita inesperada do \textbf{Primeiro Ministro Russo} (MISC) , \textbf{Vladimir Putin} (PER) . \textbf{Putin} (PER) esteve na arena durante pouco tempo, no entanto assistiu ao ensaio do \textbf{Azerbaijão} (LOC) , no qual os representantes do país ficaram supreendidos. Depois de ver a entrada \textbf{Azerbaijã} (LOC), \textbf{Putin} (PER) deu as boas-vindas a \textbf{Sir Andrew Lloyd Webber} (MISC) , compositor da música representante do \textbf{Reino Unido} (LOC) , e assegurou-lhe que iria votar na sua música. \textbf{Putin} (PER) dirigiu-se ao \textbf{Olympiyski Arena} (PER), juntamente com \textbf{Konstantin Ernst} (PER) , \textbf{Director Geral do Channel One Russia} (MISC) (o canal anfitrião). Ambos foram recebidos pelo \textbf{Executivo Supervisor do Festival} (LOC) Eurovisão da Canção Svante Stockselius , que mais tarde em entrevista ao eurovision.tv (site oficial do \textbf{ESC} (MISC)), disse que "He (\textbf{Vladimir Putin} (PER)) was very much interested in the technical details, and \textbf{wanted to know if} (MISC) there were any challenges or problems," ( Ele (\textbf{Vladimir Putin} (PER)) estava muito interessado nos detalhes técnicos, e queria saber se havia alguma hipótese de algo correr mal ). O 8 dia marca o fim dos ensaios primários. No dia 11 de Maio , nono dia de ensaios, começam os ensaios gerais (incluindo o guarda-roupa completo utilizado no próprio festival). Este é também o último dia em que os ensaios serão à porta fechada, pois a partir do 9 dia, já existirá audiência na \textbf{arena} (PER), com bilhetes pagos. Para o 8 dia, os últimos países que restam para ensaiar são: \textbf{Lituânia} (LOC), \textbf{Moldávia} (LOC) , \textbf{Albânia} (LOC) , \textbf{Ucrânia} (LOC) , \textbf{Estónia} (LOC) e \textbf{Holanda} (LOC) (da segunda semifinal), e os restantes cinco países da final: \textbf{França} (LOC) , \textbf{Rússia} (LOC) , \textbf{Alemanha} (LOC) , \textbf{Reino Unido} (LOC) e \textbf{Espanha} (LOC) No mesmo dia deu-se a \textbf{Cerimónia de Abertura do Festival} (MISC) Eurovisão da Canção , no \textbf{Euro Club} (MISC) situado na \textbf{Manege de Moscovo} (LOC) . Foi uma das maiores cerimónias de abertura de sempre. No dia 11 de Maio de 2009 , oficialmente o 9. dia, começou a semana eurovisiva. Neste dia foram levados a cabo dois ensaios gerais (incluindo guarda-roupa). Os ensaios foram divididos em dois, um ao 12:00 \textbf{CET} (ORG) e outro às 20:00 do mesmo fuso horário. Na primeira parte, ensaiaram todos os dezoito países a concurso da primeira semifinal \textbf{Montenegro} (LOC) , \textbf{República Checa} (LOC) , \textbf{Bélgica} (LOC) , \textbf{Bielorrússia} (LOC) , \textbf{Suécia} (LOC) , \textbf{Arménia} (LOC) , \textbf{Andorra} (LOC) , \textbf{Suíça} (LOC) , \textbf{Turquia} (LOC) , \textbf{Israel} (LOC) , \textbf{Bulgária} (LOC) , \textbf{Islândia} (LOC) , \textbf{Macedónia} (LOC) , \textbf{Roménia} (LOC) , \textbf{Finlândia} (LOC) , \textbf{Portugal} (LOC) , \textbf{Malta} (LOC) e \textbf{Bósnia e Herzegovina} (LOC) . Na segunda parte do dia, todos os países voltaram a ensaiar, havendo dois intervalos nos ensaios, um para mostrar os artistas a chegar a \textbf{Moscovo} (LOC) e outro para os representantes conhecerem a \textbf{Green Room} (ORG) . Também no mesmo dia, o \textbf{Channel One da Rússia} (LOC) (\textbf{C1R} (ORG)), também deu uma conferência de imprensa revelando os últimos pormenores para os três espetáculos, incluindo alguns postcards que são utilizados antes da entrada de cada páis. No 10. dia, dia 12 de \textbf{Maio de 2009} (MISC) , dia da \textbf{1. Semifinal do Festival Eurovisão da Canção 2009} (MISC) , todos os países a concorrer nesta semifinal voltaram a ensaiar pela última vez antes do espetáculo, às 17:00 \textbf{CET} (ORG). Mais tarde, às 21h (\textbf{CET} (ORG)) deu-se a 1. semifinal. Posteriormente à final, os vencedores da semifinal deram a sua conferência de imprensa, e o sorteio para as posições dos países na final, foi levado a cabo. Ao 11. dia, foi a vez dos países concorrentes na 2 semifinal realizarem os seus dois ensaios gerais do dia. \textbf{Croácia} (LOC) , \textbf{Irlanda} (LOC) , \textbf{Letónia} (LOC), \textbf{Sérvia} (LOC) , \textbf{Polónia} (LOC) , \textbf{Noruega} (LOC) , \textbf{Chipre} (LOC) , \textbf{Eslováquia} (LOC) , \textbf{Dinamarca} (LOC) , \textbf{Eslovénia} (LOC) , \textbf{Hungria} (LOC) , \textbf{Azerbaijão} (LOC) , \textbf{Grécia} (LOC) , \textbf{Lituânia} (LOC), \textbf{Moldávia} (LOC) , \textbf{Albânia} (LOC) , \textbf{Ucrânia} (LOC) , \textbf{Estónia} (LOC) e \textbf{Holanda} (LOC) , foram os dezanove países a ensaiar, da primeira vez às 12:00 \textbf{CET} (ORG) e outro às 20:00 \textbf{CET} (ORG). No dia 14 de \textbf{Maio de 2009} (MISC) , deu-se a 2 semifinal do Festival Eurovisão da Canção 2009 (21:00 \textbf{CET} (ORG)), onde houve mais um ensaio antes da semifinal, às 17:00 \textbf{CET} (ORG). No dia a seguir à 2 semifinal, todos os vinte e cinco países concorrentes na final, ensaiaram duas vezes (às 14:00 e 21:00 \textbf{CET} (ORG)). No dia 16 de Maio de 2009 , o 14 dia de \textbf{Eurovisão} (LOC), deu-se a \textbf{grande Final do Festival} (MISC) Eurovisão da Canção 2009 , onde todos os países tiveram a oportunidade de pisar o palco mais uma vez antes do grande momento às 14:00 \textbf{CET} (ORG) (7 horas antes da final). De lembrar que durante todos os nove ensaios gerais, também os actos das aberturas, dos intervalos e os próprios apresentadores foram ensaiados.



\textbf{Festival} (PER)Com quarenta e dois participantes no total, trinta e sete dos mesmo tiveram de passar por uma das duas semifinais, dezoito na primeira e dezanove na segunda. Os cinco países restantes (os \textbf{Big4} (MISC) e a \textbf{Rússia} (LOC)), tiveram acesso directo à final. Com a conclusão da construção do palco para o festival, no dia 2 de Maio de 2009 , depois de um mês de construção dia e noite, os ensaios começaram, na capital russa no dia 3 de maio de 2009. Todas as cooperativas de cada país viajaram para \textbf{Moscovo} (LOC) , assim como vários repórteres e jornalistas de todo o \textbf{Mundo} (LOC), entre os dias 1 , 2 e 3 de \textbf{Maio} (PER). As boas-vindas serão oficialmente dadas, no \textbf{Centro de Imprensa} (LOC) e no \textbf{Centro de Acreditação} (LOC), que será inaugurado no dia 2 de Maio, às 8:30, hora local. A \textbf{Rússia} (LOC) (o \textbf{C1R} (ORG) ), país anfitrião, em conferência de imprensa contou que o país não queria ganhar esta edição, nem ansiava ganhar outra tão cedo. O canal explicou que durante quinze anos a \textbf{Rússia} (LOC) concorreu para que pudesse vencer e organizar o evento em \textbf{Moscovo} (LOC) , como acontece neste ano, e explicou ainda que o país não está no festival para competir com outros países no número de vitórias. Anteriormente também o \textbf{Reino Unido} (LOC) havia dito que um segundo lugar seria o mesmo que ganhar o concurso, visto que os custos da organização do mesmo são muito elevados para o país suportar numa altura de crise como se vive. Os resultados de ambas as semifinais só serão dados a conhecer ao público, juntamente com os resultados da final, através do site oficial da eurovisão.

A final do \textbf{Festival Eurovisão da Canção} (MISC) 2009 realizou-se a 16 de Maio de 2009. Os finalistas serão:

Nota: As pontuações estão divididas em \textbf{Total} (LOC) (T) , \textbf{Televoto} (PER) (\textbf{Tv} (LOC).) e \textbf{Júri} (MISC) (\textbf{J} (LOC)) .

A ordem de divulgação dos votos, na final do \textbf{Festival Eurovisão da Canção 2009} (MISC), foi decidida através de um sorteio que se realizou no dia 16 de \textbf{Março} (PER) em \textbf{Moscovo} (LOC). Na noite da final, a \textbf{16 de Maio} (LOC), depois dos 15 minutos de votação, e de mais 15 minutos de intervalo, a partir de \textbf{Moscovo} (LOC), da arena onde vai ser realizado o festival, vai ser realizada uma ligação em directo para cada país participante na edição de 2009 (incluindo os que não passaram à final). Geralmente as ligações são feitas para a capital do país em questão, que tem o seu porta-voz num estúdio, com um fundo que consiste numa fotografia ou numa imagem real do panorama da cidade de onde estão ligados. Assim, cada porta-voz tem cerca de dois minutos para divulgar as votações nacionais, o que demora mais de uma hora.

Como vem sendo hábito, e por ser a pontuação máxima, os 12 pontos tornaram-se um símbolo do Festival Eurovisão da Canção , sendo a votação que determina o favorito de um país, e que pode virar o resultado a qualquer momento, assim como desempatar países que estejam empatados. Em baixo, é possível ver uma tabela com o número de 12 pontos que cada país recebeu, e os respectivos países que deram a sua pontuação mais alta a determinado país.



Cobertura rádio-televisiva

\textbf{O Festival Eurovisão da Canção} (MISC) , é o programa não-desportivo que consegue juntar o maior número de telespectadores ao mesmo tempo em todo o \textbf{Mundo} (MISC) (para lá dos \textbf{Jogos Olímpicos} (MISC) , e finais de \textbf{Europeus} (LOC) e \textbf{Mundiais de Futebol} (MISC) ). Todos os anos, mais de 100 milhões de pessoas por toda a \textbf{Europa} (LOC) vêem o espetáculo, traduzido em directo por um comentador escolhido pela sua televisão transmissora, membro da \textbf{EBU} (MISC) . Porém, o festival para além de ser transmitido televisivamente, também o é transmitido através de estações de rádio, e mais recentemente em directo (e posteriormente em vídeo on-demand), no site oficial do festival, através da eurovision.tv . Todos os anos, cerca de 50 televisões, 12 rádios e um canal oficial de televisão on-line ( eurovision.tv ), transmitem o festival para mais de 100 milhões de pessoas, no entanto o número vem a aumentar todos os anos, principalmente por países que já não chegavam à final, chegarem lá, tal aconteceu em 2008 , quando \textbf{Portugal} (LOC) passou pela primeira vez à final, desde a implantação de uma semifinal. A \textbf{RTP} (ORG) , televisão portuguesa que emite o festival, teve na final do concurso, mais de 3 milhões de telespectadores (num país com pouco mais de 10 milhões de pessoas). O mesmo aconteceu na \textbf{Sérvia} (LOC) . Para 2009 , é esperada uma audiência de 200 milhões de pessoas espalhadas por todo o \textbf{Mundo} (MISC) , mas uma maioria na \textbf{Europa} (LOC). Nas secções em baixo, é possível ver que televisões emitem festival, e que eventos emite assim como se são em directo ou não. Também as televisões que transmitiram para todo o \textbf{Mundo} (LOC), encontram-se na lista. Assim que estiver disponível, também será editada uma lista com as emissoras de rádio.

Todas as televisões que enviam um representante para uma edição do \textbf{Festival Eurovisão da Canção} (MISC), são obrigados a transmitir pelo menos a semifinal em que entram e a final (à excepção dos \textbf{Big4} (MISC), que se quiserem podem apenas transmitir a final). No entanto, os países podem escolher também transmitir a semifinal onde não participam, o que é adoptado pela maioria, mesmo que por vezes a semifinal onde não participa, seja transmitida noutro horário. Como exemplo mais "aficcionado", temos a \textbf{BBC} (ORG) do \textbf{Reino Unido} (LOC) , que tranmitirá os três espetáculos em directo, e no dia a seguir aos mesmos em diferido (às 00:45 \textbf{UGT} (MISC)). Em baixo está uma lista com as televisões membros da \textbf{EBU} (MISC), que têm os direitos e o dever de transmitir o \textbf{Festival Eurovisão da Canção 2009} (MISC):

Através da internet, será possível ver o concurso em qualquer parte do \textbf{Mundo} (LOC). Para além disso, várias televisões de vários países não europeus, transmitiram o festival. Também várias televisões Europeias transmitiram o festival nos seus canais internacionais. Também é afirmado que nos \textbf{Estados Unidos da América} (LOC) e na \textbf{Índia} (LOC) , o festival vai ser transmitido, porém nada disso está confirmado oficialmente. A lista de televisões transmissoras do festival a nível internacional pode ser observada em baixo.

O site oficial do \textbf{Festival Eurovisão da Canção} (MISC) também irá transmitir o festival via \textbf{peer to peer} (MISC) meio \textbf{Octoshape} (MISC) . No ano passado a audiência através da internet chegou aos 74 mil telespectadores, para 2009, mais de 100 mil são esperados.

\textbf{Em Portugal} (LOC), 3 canais distintos da \textbf{RTP} (ORG) transmitirão festival, e apesar de a \textbf{Madeira} (LOC) ser território nacional, está incluída nas transmissões internacionais, por ser uma transmissão além fronteiras \textbf{Europeias} (LOC) (de lembrar que a \textbf{Madeira} (LOC) pode ser considerada parte integrante do \textbf{Continente Africano} (LOC)).

Pela terceira vez, o festival também será transmitido em canais de alta definição. Alguns países, através do seu canal de alta definição, também possibilitaram que outros países assistam ao festival em HD. O canal \textbf{RTS} (ORG) da \textbf{Sérvia} (LOC), em 2009 apostou na difusão do festival, transmitindo-o de quatro maneiras diferentes. Esta grande transmissão deve-se ao facto de no último ano as audiências terem aumentado significativamente (de lembrar que a \textbf{Sérvia} (LOC) foi o país anfitrião em 2008).

Pela primeira vez (oficial), algumas televisões irão transmitir o festival através da internet. Para além disso, o site oficial do Festival Eurovisão da Canção , eurovision.tv , também irá transmitir o espetáculo em directo para todo o \textbf{Mundo} (MISC) (mas sem qualquer tradução, ou número de telefone para votação).

Pela primeira vez também, existem canais televisivos que transmitem o festival, em canais próprios para telemóveis.

No Festival Eurovisão da Canção 2008 , 12 emissoras de rádio emitiram o festival, este ano espera-se que o mesmo aconteça.

Até ao momento, apenas as rádios listadas abaixo confirmaram a emissão do festival (rádios de países que concorrem ao festival):

Como já é hábito, desde o início das transmissões do festival, cada televisão membro, escolhe um ou mais comentadores, que em directo, traduzem o festival para a língua materna do país a que pertencem. Na ligação acima, é possível observar uma tabela com todos os comentadores que cada país levará a \textbf{Moscovo} (LOC), para reportar o festival, que é apenas apresentado em inglês e francês, no local onde está a decorrer (por vezes surgem apresentações na língua do país anfitrião, ou em inglês).



\textbf{CD} (MISC) e DVD oficial

O \textbf{CD} (MISC) oficial do \textbf{Festival Eurovisão da Canção 2009} (MISC), ficou disponível a 1 de \textbf{Maio de 2009} (MISC) . O \textbf{CD} (MISC) é um \textbf{CD} (MISC) duplo, onde estão gravadas digitalmente todas as 42 músicas a concurso, custando um total de 19,95 euros. Porém o download de todas as músicas do festival (em versão normal e karaoke), ficou disponível a \textbf{12 de Abril de 2009} (MISC) , no site oficial, a 0,99 euros cada. O DVD apenas estará disponível a 5 de \textbf{Junho} (PER) de 2009, e será um DVD triplo, com os três espetáculos e extras, e será vendido a um preço de 24,95 euros. Para além do \textbf{CD} (MISC) e DVD, outros produtos foram lançados para o evento. Pela primeira vez, foi criado um livro com todo o programa do festival. Intitulado por \textbf{The official programme of the Eurovision Song Contest} (MISC) 2009 inclui os discursos de boas-vindas, perfis de todos os países a concurso, fala sobre o vencedor do festival em 2008 e o tour efectuado pelo mesmo, tem artigos sobre \textbf{Moscovo} (LOC) , sobre o palco, sobre os apresentadores, todas as letras das músicas, etc, com um custo de 19,95 euros Para além destes três produtos, muitos outros foram editados. Desde canecas, porta-chaves, camisolas/\textbf{T-shirts} (MISC), tapetes para ratos, a pins, posters, bonés, mochilas, imans, etc, tudo a preços que rodam os 20 euros. Para além disso, também foram criados vários pacotes, sendo o mais importante deles, o \textbf{Red Square Packtage} (ORG) , o pacote inclui o \textbf{CD} (MISC), o DVD, o programa oficial e as moedas do Festival Eurovisão da Canção 2009 (com um custo de 44,95 euros).



\textbf{Factos} (MISC) e controvérsias

A \textbf{Eslováquia} (LOC) , voltou ao concurso no \textbf{Festival Eurovisão da Canção 2009} (MISC), após onze anos de ausências (desde 1998 ). Tal facto deveu-se à conversa/reunião que houve entre a televisão eslovaca \textbf{STV} (MISC) e entre a \textbf{EBU} (MISC) . Mais tarde o director da \textbf{EBU} (MISC), \textbf{Bjørn Erichsen} (PER) , confirmou a participação/regresso deste país no festival .

Já no ano passado, a \textbf{Eslováquia} (LOC) demonstrou um interesse em participar, tal como havia acontecido nos anos anteriores desde 1998. Porém, nunca havia sido possível regressar ao festival devido a dificuldades finaceiras. No entanto, em 2008 , o país já deveria ter participado, no entanto não conseguiu, devido aos elevados custos para a transmissão dos \textbf{Jogos Olímpicos} (MISC) , que a \textbf{STV} (MISC) teve de pagar.

Devido às acções da \textbf{Rússia} (LOC) durante a guerra na \textbf{Geórgia} (LOC) , em \textbf{Agosto de 2008} (MISC) , havia muitos países - especialmente na \textbf{Europa Oriental} (LOC) - que reflectiram, de modo a que a concorrência boicotasse o festival (pelo mesmo se se realizasse na \textbf{Rússia} (LOC)). Estes países incluíram: a \textbf{Geórgia} (LOC) , que já havia confirmado que não iria participar, mas, em \textbf{Dezembro} (LOC), após a vitória na \textbf{Festival Eurovisão da Canção Júnior 2008} (MISC) , decidiu participar e os \textbf{Estados Bálticos da Estónia} (LOC) , \textbf{Letónia} (LOC) e \textbf{Lituânia} (LOC). Desde então, o \textbf{Ministro da Cultura da Estónia} (MISC) , \textbf{Laine Janes} (PER) , veio a sugerir que, juntamente com a \textbf{Letónia} (LOC) e a \textbf{Lituânia} (LOC), deveria-se realizar um boicote em sinal de protesto. A \textbf{Lituânia} (LOC) manteve o boicote como uma opção, o conflito no \textbf{Cáucaso} (LOC) poderia persistir. No entanto, na \textbf{Letónia} (LOC), tanto a televisão estatal \textbf{Latvijas Televizija} (MISC), assim como o eminente compositor e ex-ministro da cultura da \textbf{Letónia Raimonds Pauls} (LOC) , estavam contra uma retirada, e por fim, a \textbf{VTL} (ORG) finalmente confirmou a sua participação no concurso, em \textbf{Moscovo} (LOC) .

Além disso, a televisão oficial da \textbf{Estónia} (LOC), a \textbf{Eesti Rahvusringhääling} (MISC) , confirmou participação no concurso, uma vez que num estudo, 66\% dos estónios afirmaram querer o seu país representado na \textbf{Eurovisão} (LOC) . A \textbf{31 de Outubro} (MISC) , foi a vez de a \textbf{Lietuvos Nacionalinis Radijas ir Televizija} (LOC) , a televisão da \textbf{Lituânia} (LOC) membra da \textbf{EBU} (MISC) , confirmar a sua presença em \textbf{Moscovo} (LOC).

Logo no início, a \textbf{Geórgia} (LOC) manteve-se hesitante em enviar representação para \textbf{Moscovo} (LOC) , devido à recente guerra entre a \textbf{Rússia} (LOC) (país anfitrião) e o país . No entanto, a \textbf{Geórgia} (LOC) resolveu participar devido à sua recente vitória na \textbf{Eurovisão Júnior} (MISC) . Durante dois meses, o canal georgiano, \textbf{GPB} (LOC) , organizou a sua final, para escolher o representante. A final deu-se, de onde saíram vencedoras, as \textbf{Stefane \& 3G} (MISC) , com a música " \textbf{We don't wanna put in} (MISC) ". A partir de aqui toda a polémica anteriormente ultrapassada voltou, deste vez por causa da música. Ao se ler o título, apenas o interpretamos como um normal, em que a tradução significa apenas \textbf{Nós Não Queremos Pôr} (MISC) , no entanto, o título da música foi considerado um jogo de palavras, desmentido várias vezes pelo grupo representante. Se se analisar o título pelo jogo de palavras, a tradução fica \textbf{Nós} (MISC) Não \textbf{Queremos Putin} (PER) . As últimas palavras put in , aos olhos dos críticos, referem-se ao \textbf{Primeiro-Ministro da Rússia Vladimir Putin} (ORG) , que juntamente com a primeira parte do título torna-se mais séria. Também durante a música, várias partes foram entendidas como ofensas ao \textbf{Primeiro-Ministro} (PER), que tinha lançado guerra na \textbf{Geórgia} (LOC). A música foi assim considerada uma ofensa pessoal a \textbf{Putin} (PER), e devido a toda a polémica que originou, a \textbf{EBU} (MISC) desqualificou-a, referindo que não correspondia às regras do concurso. O grupo que canta a música desmentiu todas as acusações, afirmando que a música não tem nada de político e é apenas uma música de festival. A \textbf{Geórgia} (LOC) teve a oportunidade de escolher outra música ou de alterar a sua música para a representar em \textbf{Moscovo} (LOC), mas o país não aceitou a proposta e resolveu boicotar o festival em \textbf{Moscovo} (LOC). No entanto, e depois mesmo de ter sido desqualificada e de se ter retirado, a \textbf{Geórgia} (LOC) tem sido convidada por vários países europeus, para interpretar a sua música, tal como outros concorrentes fazem para divulgar a música que cantaram no festival. O grupo georgiano escolhido pela população do país actuou em quase todas as capitais da \textbf{Europa} (LOC), e em várias cidades do \textbf{Reino Unido} (LOC). E enquanto este tour acontecia, o grupo voltava a afirmar que a música não tinha sido criada para resultar no que aconteceu, que era apenas uma música de estilo disco e nada mais. No entanto, em \textbf{10 de Maio} (MISC) , dois dias antes da primeira semifinal do festival, a comitiva georgiana acabou por confirmar que a sua música era de carácter político e sátiro, contra o \textbf{Primeiro Ministro Russo} (MISC), \textbf{Vladimir Putin} (PER). A \textbf{Geórgia} (LOC) afirmou ainda que nunca teve intenções de actuar em \textbf{Moscovo} (LOC) , e agradeceu por todo o mediatismo dado pelos meios de comunicação, a fim de divulgarem a sua música, e pôr a descoberto o regime que ainda rege na \textbf{Rússia} (LOC) .

Pela primeira vez no festival, houve acusações de plágios como nunca se vira. Quase todos os países viram os seus representantes acusados de plágio, que incluíam plágios de músicas de cantores famosos, e até mesmo plágios de músicas criadas por \textbf{Ministros} (MISC). No entanto, a maioria veio-se a revelar falsa, apenas com o intuito de denegrir a imagem do país. Apenas a \textbf{Hungria} (LOC), teve a escolha de representante mais atribulada, apenas tendo conseguido efectuá-la à terceira vez (na primeira vez, a música já havia sido apresentada antes do prazo que poderia ser, e na segunda, a cantora escolhida resolveu não representar o país, dois dias depois de ter sido escolhida).

No dia 28 de Abril de 2009 , a \textbf{EBU} (MISC) aceitou o pedido de \textbf{Espanha} (LOC) , para passar os seus direitos de transmissão e de votação para a 2. semifinal, ao invés da primeira como já estava destinado desde \textbf{Março} (PER) do mesmo ano. Tal pedido de troca de emissões e votações deveu-se ao facto de o canal espanhol estar a emitir um programa dedicado à \textbf{Eurovisão} (LOC), que acabará com uma grande gala no dia da 1. semifinal, ou seja, dia 12 de \textbf{Maio de 2009} (MISC). No entanto, o entrave maior à sua emissão foi o facto de no mesmo dia estar marcado o mais importante debate político espanhol, que ocorre apenas uma vez por ano, o \textbf{Debate} (ORG) sobre el Estado de la Nación . No entanto esta troca suscitou algumas críticas e mal estar a alguns participantes, principalmente a \textbf{Portugal} (LOC) e a \textbf{Andorra} (LOC) , dois dos países que geralmente mais beneficiam dos votos espanhóis. \textbf{Andorra} (LOC), mesmo assim, foi o país que levou o ocorrido mais a sério, no entanto diz não precisar dos votos de \textbf{Espanha} (LOC) para passar à final. No dia 5 de \textbf{Maio} (MISC) , a representante espanhola fez uma conferência de imprensa, dirigida aos fãs da \textbf{Eurovisão portugueses} (MISC) e andorrinos, explicando o que se passou, e que não haveria qualquer problema com aquela decisão. Já no dia 6 de \textbf{Maio de 2009} (MISC), a \textbf{TVE} (ORG) afirmou que também transmitirá a 1. semifinal em directo, porém apenas no seu web site RTVE.es . Para piorar mais ainda a situação da televisão espanhola, no dia da segunda semifinal, o espetáculo foi para o ar uma hora mais tarde do que o previsto, pelo que a televisão não transmitiu a semifinal em directo. Em vez de transmitir a semifinal, a \textbf{TVE} (ORG) transmitiu um torneio de ténis. As razões para tal (justificadas pelo canal televisivo), foram dificuldades técnicas em transmitir o evento em directo. As linhas telefónicas para a votação foram desactivadas, e a votação de \textbf{Espanha} (LOC) ficou a cargo do seu grupo de cinco jurados apenas. Com este "incidente", a \textbf{Espanha} (LOC) desrrespeitou uma das regras do festival, que diz que todos os países devem transmitir pelo menos a semifinal onde participam ou votam (como é o caso da \textbf{Espanha} (LOC)). Porém, também em \textbf{França} (LOC) houve problemas com as transmissões, pelo que o problema poderia ter sido por parte do canal russo \textbf{C1R} (ORG) . No entanto, no dia a seguir à semifinal, a \textbf{TVE} (ORG) confirmou que preferiu transmitir o torneio de ténis em vez do espetáculo eurovisivo, e que esperava que a \textbf{EBU} (MISC) não a "castigasse", o que não aontecerá, pois a \textbf{EBU} (MISC) exigiu uma justificação oficial por parte do canal espanhol, e decidirá o que fazer ao mesmo no seu próximo encontro. No entanto, tal sanção não prejudicou em nada a representante espanhola de 2009 . No entanto, e no final acabou por ser confirmado que a \textbf{EBU} (MISC) não iria sancionar \textbf{Espanha} (LOC). Por isso, \textbf{Espanha} (LOC) já não é considerada como um país que possa desistir do Festival Eurovisão da Canção 2010 .

Durante os ensaios, o desgosto de alguns artistas com os fundos da sua canção e o trabalho de câmaras é normal, no entanto em 2009 , esses "arrufos" entre artistas e o canal russo \textbf{C1R} (ORG) , foram elevados a outro nível. Alguns países apenas pediram para os fundos serem mudados, e foram "atendidos" pela organização (como no caso de \textbf{Malta} (LOC) , por exemplo). No entanto, países como a \textbf{Irlanda} (LOC) e a \textbf{Lituânia} (LOC) não tiveram tanta "sorte". A \textbf{Irlanda} (LOC) mostrou o seu desagrado em público, sendo o primeiro país de toda a história eurovisiva a trazer uma história deste género para a imprensa. Os representantes irlandeses queixaram-se do trabalho efectuado pelas câmaras no primeiro ensaio, de forma a mudar o mesmo já para o segundo, no entanto os operadores de câmara fizeram exactamente o mesmo que no primeiro ensaio. Por sua vez, a \textbf{Lituânia} (LOC) também não gostou do trabalho das câmaras (o que também não foi mudado de um ensaio para o outro), e também pediu para o fundo que aparece durante a sua apresentação ser mudado, pedido este que não foi atendido. No entanto, o representante lituano foi mais longe, e afirma boicotar o festival caso os seus pedidos não sejam ouvidos.

Pela primeira vez na história, o \textbf{Festival Eurovisão da Canção} (MISC) também chegou aos \textbf{Estados Unidos da América} (LOC) . No dia 8 de \textbf{Maio de 2009} (MISC) , o programa mais visto da \textbf{América} (LOC), foi dedicado ao festival, passando vários momentos históricos do concurso e com a actuação especial do representante alemão , \textbf{Oscar Loya} (PER) (de lembrar que \textbf{Loya} (PER) cresceu na \textbf{Califórnia} (LOC) e por isso e tido como cidadão norte-americano) Mais tarde, a 11 de Maio de 2009 , a cadeia de televisão \textbf{NRK} (ORG) confirmou (depois de ter obtido as autorizações prévias para tal), que iria passar o vídeo do representante norueguês no programa de \textbf{Oprah Winfrey} (PER) .

O activista russo pelos direitos gays, \textbf{Nikolai Alekseev} (PER) utilizou a presença do \textbf{Festival} (PER) na \textbf{Rússia} (LOC) como uma plataforma para promover a posição do país em relação aos direitos das pessoas \textbf{LGBT} (ORG) , contrariando o prefeito de \textbf{Moscovo} (LOC), \textbf{Yuri Lujkov} (PER) , que afirmou que a homossexualidade é " satânica ". \textbf{Alekseev} (LOC) anunciou que a edição 2009 da \textbf{Marcha do Orgulho Gay de Moscovo} (MISC) , a parada gay anual da cidade, iria coincidir com a final do concurso, no dia 16 de Maio, o dia antes do \textbf{Dia Internacional} (ORG) contra a \textbf{Homofobia} (LOC) . A parada também foi chamada de "\textbf{Parada Gay Eslávica} (MISC) ", para promover os direitos e a cultura através de toda a \textbf{Região Eslávica da Europa} (LOC) . O desfile foi não obteve autorização por \textbf{Moscovo} (LOC) (pelos responsáveis da cidade), com base de que seria uma maneira de "destroir a moralidade da sociade" e foram emitidos avisos e declarações, afirmando que os manifestantes que se mivimentassem no sentido de realizar o desfile seriam tratados de forma "rigorosa", e essas "medidas rigorosas" seriam enfrentadas por qualquer um que se junta-se à manifestação. No evento, 20 protestantes foram presos pela polícia de \textbf{Moscovo} (LOC) no dia da final, incluindo o próprio \textbf{Nikolai Alexeyev} (PER) e o activista pelos direitos humanos, \textbf{Peter Tatchell} (PER) , que exclamou "isto mostra que as pessoas na \textbf{Rússia} (LOC) não são livres" enquanto era levado pela políca. A representante da \textbf{Suécia} (LOC), \textbf{Malena Ernman} (PER) consagrou a causa dizendo que ela não é homossexual, mas ficaria orgulhosa de se chamar isso a ela própria para ajudar os seus fãs. Ela constatou que ficou triste por causa do governo de \textbf{Moscovo} (LOC) não permitir um "tributo ao amor". A porta-voz \textbf{Sueca} (MISC), \textbf{Sarah Dawn Finer} (PER) ao ler os votos, utilizou um pin ao \textbf{peit} (LOC), com um cão desenhado com as cores do arco-íris, em forma de solariedade. O vencedor norueguês do evento, \textbf{Alexander Rybak} (PER) também se referiu à controvérsia quando disse que o \textbf{Festival Eurovisão da Canção} (MISC) é por si só a maior parada gay.

Uma série de controvérsias entre a \textbf{Arménia} (LOC) e o \textbf{Azerbaijão} (LOC) tiveram lugar durante as semifinais e final. O 'postcard' introdotório da \textbf{Arménia} (LOC) (o pequeno vídeo clip que foi para o ar antes da actuação armeniana) mostrou, para além de outros monumentos, a \textbf{We Are Our Mountains} (MISC) , uma estátua localizada em \textbf{Stepanakert} (LOC) , a cidade capital não-reconhecida de \textbf{Nagorno-Karabakh} (LOC) , que em de jure constitui parte do \textbf{Azerbaijão} (LOC). A estátua foi construída durante os tempos soviéticos para comemorar a ascensão dos territórios armenianos. O \textbf{Azerbaijão} (LOC) fez uma queixa à \textbf{EBU} (MISC) relativamente ao vidio, dizendo que o mesmo era inaceitável, visto/baseado no facto de \textbf{Nagorno-Karabakh} (LOC) ser parte oficial do \textbf{Azerbaijão} (LOC), e que foi substancialmente editado pela televisão para a final. Em retaliação, a apresentadora dos votos armenianos, , \textbf{Sirusho} (MISC) , tinha atrás de si, um quadro com a foto do monumento, repetido várias vezes enquanto ela revelava os votos, e por fundo, encontrava-se um ecrã com a principal praça da capital, para mostrar a disputa pelo monumento. Controvérsia à parte, os votos da \textbf{Arménia} (LOC) foram revelados, e nenhum voto foi atribuído ao \textbf{Azerbaijão} (LOC) por parte do seu país vizinho.

O \textbf{Azerbaijão} (LOC) mostrou os monumentos \textbf{Maqbaratoshoara} (LOC) e \textbf{Segonbad} (LOC) , os símbolos das cidades de \textbf{Tabriz} (LOC) e \textbf{Urumieh} (LOC) no \textbf{Azerbaijão} (LOC) região do \textbf{Irão} (LOC) . A imprensa armeniana, posteriormente criticou o facto de que enquanto a demostração de um monumento por parte da \textbf{Arménia} (LOC), situado no \textbf{Azerbaijão} (LOC), na região de \textbf{Karabakh} (LOC) depois do protesto do governo do \textbf{Azerbaijão} (LOC), este permitiu a introdução de monumentos \textbf{Iranianos} (LOC) no seu cartão postal.

No final de \textbf{Maio} (MISC), a \textbf{EBU} (MISC) pronúnciou-se pela primeira vez sobre o assunto, constatando que seriam anunciadas sanções para o \textbf{Azerbaijão} (LOC), em \textbf{Junho} (LOC). Tal como a referência constata, a televisão \textbf{Azerbaijã} (MISC) falhou na parte de mostrar a representação armeniana, e em vez disso, colocou no ar, a actuação de outros representantes. Para além disso, durante os quinze minutos de votação, não apareceu nenhum número na televisão, para o público do \textbf{Azerbaijão} (LOC) poder votar pela \textbf{Arménia} (LOC). O \textbf{Azerbaijão} (LOC) negou todas as acusações, e promoveu um vidio da sua televisão, em suporte da sua posição.

Todos os anos, várias estatísticas são feitas, de forma a comparar as várias edições realizadas ao longo da história . Principalmente estas estatísticas incluem o número de línguas utilizadas no geral, os mais compridos e pequenos títulos das canções e o número de artistas.

Tal como em todas as edições passadas, desde a retirada da regra da língua (que consistia em cada país ter que cantar pelo menos parte da música na sua língua oficial), em 1999 , o inglês é novamente a língua mais usada no palco da \textbf{Eurovisão} (LOC) outra vez. 32 de 42 países cantaram as suas músicas parcial ou inteiramente em inglês (24 das entradas, serão interpretadas apenas em inglês). Com um grande "espaço", a segunda língua mais utilizada na edição deste ano, é o russo , que será ouvido em três músicas (contando com a parte da canção proveniente da \textbf{Lituânia} (LOC) , que será interpretada em russo). Em terceiro lugar, a língua mais utilizada é o francês , utilizado pela \textbf{França} (LOC) e pela \textbf{Suécia} (LOC) . Nove países utilizaram obviamente, mais do que uma língua em palco. A língua portuguesa será apenas utilizada na música representante de \textbf{Portugal} (LOC) , como é hábito desde a entrada do país no festival há quase 50 anos . Assim como o português, há outras línguas que serão apenas utilizadas por um país concorrente.

18 países utilizaram uma (ou mais) das suas línguas nacionais oficiais em palco. Por outro lado, existem algumas músicas que não são interpretadas nem em inglês, nem noutra língua que seja a oficial do país que representam, como é o caso do romeno utilizado pela \textbf{República Checa} (LOC) , o francês na música da \textbf{Suécia} (LOC), o russo interpretado na canção da \textbf{Letónia} (LOC) e o ucraniano utilizado pela \textbf{Rússia} (LOC) (facto que causou algum transtorno no início). Também se pode mencionar o uso do latim na entrada irlandesa , no entanto, a expressão \textbf{Et} (MISC) cetera (mais conhecida como etc ), tornou-se muito comum em muitas línguas.

O maior título de canção do \textbf{Festival Eurovisão da Canção} (MISC) deste ano, é o da canção ucraniana : "\textbf{Be my Valentine!} (MISC) (\textbf{Anti-crisis girl} (MISC))" com 34 caracteres o que o aproxima do \textbf{recorded} (MISC) instituido em 1964 , pela música alemã , " \textbf{Man gewöhnt sich so schnell an das Schöne} (MISC) " com 41 caracteres. O título da música isrealita seria ainda maior do que o da música proveniente da \textbf{Ucrânia} (LOC), caso o título em hebreu da música tivesse sido acrescentado em parêntises ao actual título, contudo, e oficialmente a música é apenas intitulada como " \textbf{There must be another way} (MISC) ". Os títulos das músicas mais pequenas, correspondem às músicas da \textbf{Lituânia} (LOC) e da \textbf{Rússia} (LOC) , com as músicas " \textbf{Love} (PER) " e " \textbf{Mamo} (LOC) " respectivamente. No entanto, ambas não são os títulos com menos caracteres numa música a concurso no festival, pois esse recorde cabe a algumas músicas que apenas são intituladas por dois caracteres, como por exemplo, "\textbf{El} (LOC)" em 1982 e "\textbf{Go} (MISC)" em 1988 .

Sem contar com os artistas que são considerados sub-cantores/dançarinos, os seguintes actos irão competir em \textbf{Moscovo} (LOC) :

O número de artistas femininos volta mais uma vez a ser superior ao de artistas maculinos, facto que já ocorre há várias edições do festival, no entanto nota-se um acréscimo de duos a concurso.

Ao longo da história, vários artistas repetiram a sua experiência eurovisiva uma ou mais vezes. Em 2009, os repetentes foram:

Ao longo dos anos, mais particularmente a partir de 2000 , a \textbf{Eurovisão} (LOC) ultrapassou a barreira do ser um festival musical, passando assim a utilizar grandes arenas para a realização do concurso, assim como a construção de grandes palcos, etc. Em baixo podem ser vistos alguns dos números do Festival Eurovisão da Canção 2009, que conseguiu bater alguns recordes:



Festival Eurovisão da Canção 2010



Países que classificaram-se para a final \textbf{Países} (MISC) que não se classificaram para a final \textbf{Países} (MISC) que participaram no passado, mas não em 2010

O \textbf{Festival Eurovisão da Canção de 2010} (MISC) (em inglês : \textbf{Eurovision Song Contest} (ORG) 2010 , em francês : \textbf{Concours Eurovision} (MISC) de la chanson 2010 e em norueguês : \textbf{Internasjonale} (LOC)/\textbf{Europeiske Melodi Grand Prix} (MISC) 2010 ) foi o 55 \textbf{Festival Eurovisão da Canção} (MISC) na história. Com a vitória de \textbf{Alexander Rybak} (PER) pela \textbf{NRK} (ORG) na \textbf{Eurovisão 2009} (MISC) com a canção " \textbf{Fairytale} (PER) ", a \textbf{Noruega} (LOC) venceu o direito de organizar o festival, que se realizou na cidade de \textbf{Baerum} (LOC) ,um subúrbio de \textbf{Oslo} (LOC) na \textbf{Telenor Arena} (LOC) (com uma capacidade superior a vinte mil espectadores). As datas oficiais foram apresentadas no dia 27 de maio de 2009 . Nos dias 25 de maio e 27 de maio realizaram-se as duas semifinais, respectivamente, e no dia 29 de maio teve lugar a grande final. Assim o \textbf{Festival} (PER) realizou-se não só na data mais tardia desde que há semifinais, como também na data mais tardia desde o festival de 1999 (a final nesse ano foi em 29 de Maio, ou seja, o mesmo dia). No entanto, as datas previstas para o \textbf{Festival} (LOC) nem sempre foram estas, visto que tiveram de ser alteradas. As datas preliminares eram 18 de maio (1 semifinal), 20 de maio (2 semifinal) e dia 22 de maio (grande final). No entanto, a data da final gerou alguma "confusão", pois no dia 22 de \textbf{Maio} (LOC), também teria lugar a final da \textbf{Champions League} (MISC) , por isso estas datas iniciais do festival foram alteradas, e antes disso, houve várias especulações sobre uma provável realização do festival em 2 e 4 de \textbf{Junho} (PER) (semifinais) e a final no dia 6 de junho. Devido ao fator da data provisória chocar com a da \textbf{Final da Champions League} (MISC), várias televisões (como a \textbf{RTP} (ORG) de \textbf{Portugal} (LOC), por exemplo), teriam que transmitir o festival no seu canal secundário ou então posteriormente.

O formato do festival foi o mesmo de 2008 e 2009, com duas semifinais e, posteriormente, uma final. O sistema de votação, sofreu apenas uma alteração em relação aos festivais de 2008 e 2009 : as semifinais passam também a ter um sistema de votação 50\% televoto e 50\% júri, como acontecia na final, deixando de haver 9 classificados por televoto e um por júri para a final. Já o ato de votação acarreta a maior mudança na \textbf{Eurovisão} (LOC) desde a implementação do sistema de votação por televoto. Ao contrário dos últimos anos, onde eram utilizados 15 minutos para votar no fim do espetáculo, em 2010 a votação passou a acontecer durante todo o espetáculo (nos três eventos), mas mesmo assim no final foi feita a habitual contagem decrescente dos quinze minutos finais de votação. Para a edição deste ano, havia a expectativa de um novo recorde de países. Era esperada a participação de 40 a 45 nações europeias (o máximo estabelecido pela \textbf{EBU} (MISC) . A \textbf{União Europeia de Radiodifusão} (ORG) anunciou, logo após o festival de 2009, que trabalharia fortemente para que antigos países como a \textbf{Itália} (LOC) , o \textbf{Mónaco} (LOC) e a \textbf{Áustria} (LOC) retornassem ao concurso já na edição de 2010, porém tal não aconteceu. Aquele organismo anunciou também que o \textbf{Cazaquistão} (LOC), a \textbf{Palestina} (LOC) e o \textbf{Qatar} (LOC), apesar de quererem entrar no concurso, o mais provável de acontecer seria não terem sua participação autorizada, já que se localizam fora da zona geográfica pretendida para o festival e além disso não têm uma televisão local afiliada à \textbf{EBU} (MISC).

Pela primeira vez, em vários países, fãs criaram petições online para levar um determinado artista ao festival, como é o caso de \textbf{Portugal} (LOC), com o seu ano pioneiro nesta matéria, onde foram criadas dezenas de petições e publicidades. A primeira reunião oficial entre a \textbf{EBU} (MISC) e a televisão norueguesa, \textbf{NRK} (ORG), responsável pelo festival, ocorreu a 10 e 11 de junho de 2009 . Nessa reunião, \textbf{Jon Ola Sand} (PER) foi nomeado \textbf{Produtor Executivo do Festival Eurovisão da Canção 2010} (MISC), assim como a restante equipa apresentada. Para esta edição do festival, \textbf{União Europeia de Radiodifusão} (ORG) anunciou que convidou 54 países para participar no evento, no entanto com um máximo de 45 países. Em 4 de dezembro de 2009 ocorreu a entrega da chave da cidade anfitriã eurovisiva em \textbf{Oslo} (LOC), onde também foi apresentado o tema do festival: " \textbf{Share The Moment} (MISC) ", assim como os traços gerais para a realização do concurso. No dia 31 de dezembro de 2009 , foi revelada a lista final de participantes pela \textbf{União Europeia de Radiodifusão} (ORG), onde constava a \textbf{Lituânia} (LOC) que se havia retirado do festival. No fim, ficaram confirmados 39 países confirmados para 2010, três a menos do que a edição anterior e quatro a menos do que as edições de 2008 e a posterior em 2011.



\textbf{Formato} (PER)O \textbf{Ministro Norueguês da Cultura} (MISC), \textbf{Trond Giske} (PER) , e o director-executivo da \textbf{NRK} (ORG) , \textbf{Hans-Tore Bjerkaas} (PER) aprovaram um orçamento de 17 milhões de euros (150 milhões de \textbf{Kroners Norueguesas} (MISC)) para a realização do evento. O valor é bem menor do que todos os gastos totais na edição de 2009 , em \textbf{Moscovo} (LOC) (que teve o custo total de 40 milhões de euros), mas no entanto, é bem maior do que o dedicado ao \textbf{Festival Eurovisão da Canção 2007} (MISC) , em \textbf{Helsínquia} (LOC) . No dia 12 de junho de 2009 , depois de uma reunião nos dois dias anteriores entre a \textbf{EBU} (MISC) e a \textbf{NRK} (ORG), foi apresentada a equipa gestora do \textbf{Festival Eurovisão da Canção 2010} (MISC). O \textbf{Produtor Executivo} (MISC) da edição 2010 da \textbf{Eurovisão} (LOC) foi \textbf{Jon Ola Sand} (PER) ; \textbf{Hasse Lindmo} (PER) foi o \textbf{Produtor Televisivo} (LOC); enquanto que \textbf{Stian Malme} (PER) foi o \textbf{Manager} (MISC) do local onde ocorreu o festival (este último, já é conhecido do público eurovisivo, especialmente o norueguês, por ser o \textbf{Chefe da Delegação Norueguesa} (MISC) em outros festivais); finalmente, \textbf{Evy Hauffen} (PER) foi a \textbf{Line Producer} (PER) (ela foi a manager durante o \textbf{Festival Eurovisão da Canção 1996} (MISC), que também aconteceu no país). O festival teve pelo terceiro ano consecutivo duas semifinais, de onde apenas vinte países concorrentes (10 de cada semifinal) passaram à final, na qual concorreram vinte e cinco países (os \textbf{Big 4} (MISC) e o país organizador tiveram, como já é hábito, apuramento directo para a final). No 23 de setembro de 2009 , a \textbf{NRK} (ORG) foi obrigada a rever os gastos para o Festival Eurovisão 2010, elevando assim o custo do festival para 24 milhões de euros (211 milhões de \textbf{Coroas Norueguesas} (LOC)). No entanto a televisão norueguesa já pensava seriamente em aumentar os custos para 30 milhões de euros (10 milhões a menos do que a edição anterior). O orçamento final para o \textbf{Festival Eurovisão da Canção 2010} (MISC), foi aprovado pelo \textbf{Parlamento Norueguês} (ORG) em 13 de outubro de 2009. O custo final do \textbf{Festival Eurovisão da Canção 2010} (MISC), ficou aproximadamente em 25.3 milhões de euros. Para cobrir os custos extras, o governo e a \textbf{NRK} (ORG) concordaram que os cidadãos noruegueses pagassem 11,8 euros (99 \textbf{Coroas Noruegueses} (LOC)). A receita geral do festival está estimada em 9 milhões de euros (75 milhões de \textbf{Coroas Noruegueses} (LOC)). Assim, a edição de 2010 do \textbf{Festival Eurovisão da Canção} (MISC) custou cerca de 15 milhões de euros a menos do que a edição do ano anterior (depois de ter sido oficialmente anunciado que a edição de 2009 havia custado 43 milhões de euros, e não apenas 40 como anunciado de início). Em 15 de Dezembro de 2009, depois da publicação da lista final de participantes, teve lugar um sorteio entre as escolas de \textbf{Oslo} (LOC), para as quais foi determinado um país, do qual terão de trabalhar em vários aspectos. Este projeto, pensado pela \textbf{NRK} (ORG), visou tornar o Festival Eurovisão da Canção mais educativo e abrangente, voltando a atrair a atenção dos mais jovens.

Durante os dias 10 e 11 de unho de 2009, a equipa de produção para o festival, realizou a sua primeira reunião com a \textbf{União Europeia de Radiodifusão} (ORG) (\textbf{EBU} (MISC)) para uma troca de informações sobre o que seria feito até à data, para trocar ideias e estabelecer as bases de uma cooperação frutuosa e intensa. Depois de uma sessão, com a participação de pessoas de uma variedade de áreas, pequenos grupos foram formados para falar sobre cada parte da organização do evento, como aspectos técnicos do festival, o conteúdo dos três espectáculos transmitidos ao vivo, infraestruturas, marketing , \textbf{branding} (LOC) , nova mídia, facilidades de imprensa e de comunicação. A 3 de outubro de 2009 , o site oficial da \textbf{Eurovisão} (LOC) só divulgou poucas informações que foram confirmadas pela televisão norueguesa e pela \textbf{EBU} (MISC) . A esta altura já tinha sido escolhida a temática, mas só foi divulgada apenas a informação de que os cartões postais de cada país quebrariam a tradição de serem cartões que apresentam os pontos do país, não tendo mais explicações. Por sua vez, a \textbf{EBU} (MISC) afirmou estar muito satisfeita com todo o trabalho desenvolvido pela \textbf{Noruega} (LOC), e afirmou ainda que o planeamento da \textbf{Eurovisão} (LOC) levou apenas 5 meses, enquanto os russos apresentaram seu planeamento em 8 meses. É considerado que isto se deve ao fato de a televisão norueguesa já ter organizado o festival duas vezes. No entanto, mais nenhuma informação foi divulgado. Wm 26 de outubro foram reveladas as empresas responsáveis pelo design de todo o festival, desde o logotipo até às telas que iriam estar na sala verde. Foram três empresas que desenvolveram todo o projeto do evento - a \textbf{Snøhetta} (LOC) , a \textbf{Gosu Design} (LOC) e a \textbf{Handicraft} (PER) . A 4 de dezembro de 2009 , foi realizada a cerimônia de entrega da \textbf{Host City Insignia} (ORG) para as cidades de \textbf{Osloe de Bærum} (LOC) , sedes da \textbf{Eurovisão de 2010} (MISC). Os prefeitos de câmara de ambas as cidades, \textbf{Fabian Stang} (PER) e \textbf{Odd Reinsfelt} (PER) , receberam a chave da \textbf{Eurovisão de Yury Luzhkov} (LOC) , vice-presidente da câmara de \textbf{Moscovo} (LOC), sede da última edição. No mesmo dia, foi apresentada toda a temática para o \textbf{Festival Eurovisão da Canção 2010} (MISC). O tema do festival foi escolhido como \textbf{Share The Moment} (PER) , em português \textbf{Compartilhe o Momento} (PER) . O logotipo do festival está em volta de bolas brancas e roxas, como as cores que apareceram na apresentação da \textbf{Noruega} (LOC) em 2009. A concepção de arte do tema do \textbf{Festival Eurovisão da Canção 2010} (MISC) foi desenhada pelas premiadas empresas de design norueguês \textbf{Gosu} (PER), \textbf{Handverk} (PER) e \textbf{Snøhetta} (PER), com base em pequenos aspectos dados pela \textbf{NRK} (ORG). " Queremos partilhar o \textbf{Festival Eurovisão da Canção} (MISC), ao invés de apenas fazer a organização ", constatou \textbf{Hasse Lindmo} (PER), produtor de TV do concurso para a \textbf{NRK} (ORG), em entrevista ao \textbf{Eurovision} (ORG) " Uma coisa especial que faz com que o \textbf{Festival Eurovisão da Canção} (MISC) seja um evento tão único é o fato de que é a única data do ano, na \textbf{Europa} (LOC), que cerca de 125 milhões de pessoas fazem exatamente a mesma coisa, exactamente na mesma hora. Quer se esteja na \textbf{Islândia} (LOC) ou na \textbf{Arménia} (LOC), em \textbf{Portugal} (LOC) ou na \textbf{Finlândia} (LOC). Ou então aqui mesmo, em \textbf{Oslo} (LOC) ", explicou \textbf{Lindmo} (PER). " Como é que são 125 milhões de pessoas? Essa foi uma das questões principais. " Cada uma das esferas coloridas, em forma de balão, representam a reunião de pessoas e da diversidade das emoções que cercam o \textbf{Festival Eurovisão da Canção} (MISC). A escolha do tema \textbf{Share The Moment} (MISC) foi estritamente "encaixado" nas duas \textbf{Semifinais} (MISC) e na \textbf{Final do Festival} (MISC) Eurovisão da Canção 2010, bem como em todos os eventos relacionados nas semanas e meses anteriores, como os sorteios, as reuniões dos \textbf{Chefes das Delegações em Março} (MISC), a recepção de boas-vindas, em meados de \textbf{Maio} (PER) e depois nas festas que ocorrem durante o período eurovisivo e após o mesmo. Porém, nas horas seguintes ao anúncio/apresentação, as opiniões dos fãs recaíram na sua grande maioria para o desgosto pelo logotipo, sendo as principais críticas, a sua simplicidade e o facto de parecer ter sido feito por um fã, e não pelas três empresas de design responsáveis por ele.

No dia 8 de fevereiro de 2010 , com a colocação dos ingressos á venda, foi também divulgada a planta para a \textbf{Telenor Arena} (LOC) , com a disposição do palco e das arquibancadas. Através desta imagem, foi possível observar que o palco seria redondo, e teria uma pequena saliência que "seria interna" da primeira multidão na arena. Esta saliência era feita através de uma pequeno corredor semicircular e acabava noutro palco mais pequeno, em forma circular. No fundo esteve novamente montado uma gigantesco painel de "leds". A sala de espera dos artistas ficou novamente por detrás do palco. A 19 de março de 2010 foi revelada a primeira imagem do palco, sendo possível observar um palco simples, com a parte central de forma redonda e com dois níveis, o fundo teria um painel \textbf{LED} (ORG), porém com vários recortes de vários tamanhos com o feitio de circunferências. O palco teve também um pequeno anexo em frente do mesmo, ao qual se acedia através de uma pequena pista, e também ele tinha forma redonda.

Um grande número de fãs começou uma campanha no site social \textbf{Facebook} (MISC) , para o retorno da orquestra ao \textbf{Festival Eurovisão da Canção em Oslo} (MISC), pela primeira vez, desde o Festival Eurovisão da Canção 1998 , com perto de quatro mil assinaturas já conseguidas. A orquestra utilizada desde o primeiro concurso em 1956 , deixou de estar presente na \textbf{Eurovisão} (LOC), depois do concurso de 1998 (42 anos depois), devido aos rápidos avanços tecnológicos na área da música, o que tornou muito mais útil o uso do "playback". O debate sobre a possibilidade de um regresso da orquestra ao concurso tem sido altamente discutido entre os fãs, uma vez que, se a \textbf{EBU} (MISC) apoiar a ideia, será "renascida" a famosa orquestra \textbf{Eurovisiva} (PER), que durante 42 anos marcou presença no concurso. \textbf{Jan Fredrik Heyerdahl} (PER), da \textbf{Norwegian Radio Orchestra} (ORG), constatou que eles estariam interessados em participar no festival do próximo ano, perante um acordo entre a \textbf{EBU} (MISC) e a \textbf{NRK} (ORG). No entanto, o principal desafio apresentado à orquestra, é o número de países/entradas por concurso. A última vez que uma orquestra marcou presença no concurso, concorreram apenas vinte e cinco países, em apenas um evento televisivo, no entanto, o número de participantes aumentou de vinte e cinco para cerca de quarenta, desde esse ano, existindo também agora três eventos televisivos, o que implica a disponibilidade da orquestra por mais de um mês, inteiramente exclusivo à \textbf{Eurovisão} (LOC) (fora os ensaios que a orquestra teria que organizar entre si).

O diretor-geral do \textbf{Festival Eurovisão da Canção 2010} (MISC), anunciou que os cartões postais do festival, que antecedem cada um dos países/entradas participantes, não serão sobre lendas nem fábulas sobre a \textbf{Noruega} (LOC), nem sobre tradições viquingues , ou qualquer outro tipo de história sobre a \textbf{Noruega} (LOC). O próprio diretor-geral da \textbf{Eurovisão} (LOC) já havia mostrado algum receio de os cartões postais serem feitos sempre do mesmo jeito. Ao contrário, foram usadas paisagens de toda a \textbf{Noruega} (LOC), de forma a coincidir com o tema do festival \textbf{Compartilhe} (LOC) o momento . Como \textbf{Jon Olan Dand} (PER) descreveu (diretor do festival deste ano), os postais serão "quimicamente livres para a beleza da \textbf{Noruega} (LOC)". A 19 de março de 2010 , foi revelado o verdadeiro e total conteúdo dos cartões postais. Pela primeira vez desde 1996, os artistas voltaram a aparecer nos cartões. Os postais foram então divididos em duas partes, a primeira destinada aos artistas, e a segunda, a parte em que os artistas seriam guiados por famílias norueguesas, numa descoberta das tradições do país.

A maior novidade na \textbf{Eurovisão} (LOC) para 2010, foi a realização de um gigantesco \textbf{Flash Mob Dance} (MISC) . Ao todo, dez cidades europeias viram a iniciativa, onde qualquer cidadão comum pôde participar, e foi feito então um vídeo com as dez cidades, a culminar na dança das 18 mil pessoas na \textbf{Telenor Arena} (LOC). A música a dançar foi "\textbf{Glow} (PER)", dos \textbf{Madcon} (MISC) , composta despropositadamente para a ocasião. Esta foi uma das iniciativas sob o tema do festival, e que serviu para compartilhar o momento musical que foi o flash mob.

Antes do anúncio oficial dos apresentadores por parte da \textbf{NRK} (ORG), \textbf{Jon Almaas} (PER) e \textbf{Fredrik Skavlan} (PER) estavam bem colocados para serem eleitos. O famoso duo TV 2 , \textbf{Thomas Numme} (PER) e \textbf{Harald Rønneberg} (PER), ficaram em primeiro durante uma campanha de votação levada a cabo pela revista \textbf{Dagbladet} (PER) , na sua página da \textbf{internet} (MISC), onde os leitores puderam votar na celebridade que mais queriam que apresentasse o festival. Também na rede televisiva da \textbf{TV2} (MISC), a personalidade televisiva \textbf{Dorthe} (PER) Skappel demonstrou interesse em apresentar o grandioso espetáculo, e ficou em segundo na votação da \textbf{Dagbladet} (PER). A escolha e o anúncio oficial de quem seriam os apresentadores foi realizado durante um encontro entre os dirigentes do festival, a 22 de março de 2010 , no entanto a revelação dos apresentadores foi feita antes.

Assim sendo, a \textbf{NRK} (ORG) anunciou os apresentadores do festival a 10 de março de 2010 : \textbf{Erik Solbakken} (PER) , \textbf{Haddy N'jie} (PER) e \textbf{Nadia Hasnaoui} (PER) . Os três apresentadores abriram os três eventos em conjunto; \textbf{Solbakken} (PER) e \textbf{N'jie} (PER) introduziram os artistas e fizeram reportagens a partir da sala de espera dos artistas durante a votação. \textbf{Hasnaoui} (PER) apresentou os apurados para a final (nas semifinais), a secção de votação e afins, relacionados com o quadro de votações na final. Esta será a segunda vez na história que mais de dois apresentadores estarão responsáveis por conduzir o espetáculo, tendo a última vez ocorrido no \textbf{Festival Eurovisão da Canção 1999} (MISC) , sediado em \textbf{Jerusalém} (LOC) , \textbf{Israel} (LOC) (apesar de em 2009 terem estado quatro apresentadores, apenas dois apresentaram cada evento). Outro fator interessante dos apresentadores, é o fato de \textbf{Nadia Hasnaoui} (PER) e \textbf{Haddy} (PER) Njie terem nascido fora da \textbf{Noruega} (LOC) e serem filhas de norueguesas e \textbf{Haddy} (MISC) Njie também ter sido a primeira negra a apresentar o festival.

Tal como ocorre em todas as edições do \textbf{Festival Eurovisão da Canção} (MISC), nos últimos anos, existiu um patrocinador principal do concurso. A edição de 2010 teve como patrocinador principal o \textbf{Telenor Group} (ORG) , o principal grupo de telecomunicações norueguês, presente em oito países da \textbf{Europa} (LOC) e em outros espalhados pelo mundo. O \textbf{Telenor Group} (ORG), é também o dono da \textbf{Telenor Arena} (LOC) , local onde aconteceu o festival. Durante a apresentação do grupo como patrocinador, também foi revelada a criação de um telemóvel/celular dedicado à \textbf{Eurovisão} (LOC), com conteúdo relacionados e ligado ao festival.



\textbf{Local} (MISC)O \textbf{Festival Eurovisão da Canção 2010} (MISC) foi realizado em \textbf{Oslo} (LOC) , na \textbf{Noruega} (LOC) , depois da sua vitória no \textbf{Festival Eurovisão da Canção 2009} (MISC) em \textbf{Moscovo} (LOC) , \textbf{Rússia} (LOC) , com a canção de \textbf{Alexander Rybak} (PER) , " \textbf{Fairytale} (PER) ". Em 27 de maio de 2009 , a \textbf{NRK} (ORG) , televisão oficial e organizadora do \textbf{Festival Eurovisão da Canção 2010} (MISC), constatou que o \textbf{Festival} (LOC) seria realizado na região metropolitana de Oslo (esta foi a quarta vez que o país acolheu o evento, caso se conte com a realização do \textbf{Festival Eurovisão da Canção Júnior 2004} (MISC) também no país). E posteriormente a \textbf{NRK} (ORG) fez a proposta da realização do \textbf{Festival na Telenor Arena} (MISC) , uma arena nova no subúrbio de \textbf{Bærum} (LOC) , ou no já conhecido \textbf{Oslo Spektrum} (MISC) . Ambas as propostas foram avaliadas pela \textbf{European Broadcasting Union} (ORG) (\textbf{EBU} (MISC)), e o local sede do festival foi confirmado. Antes da revelação do local para o evento, o maior jornal \textbf{Norueguês} (MISC), \textbf{Aftenposten} (PER) constatou que a procura por uma arena que satisfizesse as exigências do evento começou poucos dias depois da vitória em \textbf{Moscovo} (LOC). A \textbf{NRK} (ORG) e o \textbf{Ministério da Cultura} (ORG) tiveram conversas com a \textbf{EBU} (MISC) , e decidiram que não haveria a necessidade da construção de uma nova arena para o \textbf{ESC} (MISC) (até porque não haveria tempo para isso), e em vez disso, seria utilizada uma infraestrutura já existente no país. Junto com a declaração, foram listadas quatro arenas, duas em \textbf{Oslo} (LOC),uma em \textbf{Bærum} (LOC) e outra em \textbf{Hamar} (LOC) . Eram elas: \textbf{Telenor Arena} (PER) (inaugurada em \textbf{Fevereiro de 2009} (LOC), com uma capacidade para mais de vinte mil espectadores em concertos, e a mais indicada desde o início para sediar o festival, está em \textbf{Bærum} (LOC) , a quinze minutos de carro do centro de \textbf{Oslo} (LOC) ); o \textbf{Oslo Spektrum} (ORG) (inaugurado em 1990, sediou o \textbf{Festival Eurovisão da Canção 1996} (MISC) , era considerada pequena demais para grandes eventos, no entanto tem uma capacidade para 10 mil pessoas, porém, e apesar de se ter tornado pequena para o tamanho atual do \textbf{Eurovisão} (LOC), vários manifestaram o seu apreço por esta arena sediar o festival por já o ter feito uma vez, e além disso, é a arena com a melhor acústica, de entre todas as opções); o \textbf{Vallhall} (LOC) (também uma arena situada em \textbf{Oslo} (LOC), com capacidade para 12 500 pessoas em concertos) e a \textbf{Vikingskipet Olympic Arena} (PER) (foi a única opção oficial fora da cidade de \textbf{Oslo} (LOC), situa-se em \textbf{Hamar} (LOC), e tem uma capacidade para 15 mil pessoas). Vários cidadãos noruegueses manifestaram o seu descontentamento, ao verem a região metropolitana de Oslo ser escolhida para local do festival, dizendo que a \textbf{Noruega} (LOC) é muito mais do que a capital. No entanto, e para além das quatro propostas apresentadas, nenhum local do país tem a infraestrutura necessária para sediar o evento. Apesar disso, e como foi dito desde o início, o local escolhido foi a \textbf{Telenor Arena} (PER), que tem como único problema a acústica, problema que se esperia ser resolvido durante o ano de preparações que havia até ao evento. Em 3 de julho de 2009 , foi decidido que o festival seria realizado na \textbf{Telenor Arena} (LOC) no município vizinho de \textbf{Bærum} (LOC), a cerca de 15 minutos de \textbf{Oslo} (LOC) (do centro da cidade). \textbf{Hasse Lindmo} (PER) , produtor do próximo \textbf{Festival da Eurovisão} (MISC) na \textbf{NRK} (ORG), afirmou: " Acreditamos que esta é uma boa escolha. A \textbf{Fonebu Arena} (LOC) foi construída para acolher grandes eventos - e este é definitivamente um evento grandioso. Esta escolha assegura grandes condições para a produção ". O \textbf{Oslo Spektrum} (ORG) acabou foi eliminado como sede e principal do festival, seu tamanho.s

No dia 27 de \textbf{Janeiro} (LOC) de 2010, foi apresentado o plano oficial para as \textbf{Cidades de Bærum e Oslo} (LOC) durante a \textbf{Eurovisão} (LOC). Após vários meses de conversações com diversas instalações hoteleiras, a \textbf{EBU} (MISC) apresentou a sua lista oficial de hotéis para a \textbf{Eurovisão} (LOC) e respectivos preços. Estes hotéis seriam destinados às delegações, jornalistas e fãs que se dirigissem à cidade para ver o espetáculo. No mesmo dia foi anunciado que os ingressos para o \textbf{Festival Eurovisão da Canção 2010} (MISC), seriam colocados à venda no dia 8 de Fevereiro de 2010. Foram vendidos ingressos para as semifinais, e os dois últimos ensaios gerais para a final, todos estes a um preço de 600, 450 e 300 \textbf{NOK} (ORG), o que equivale respectivamente a 73.25, 54.93 e 36.63 euros, e para a final, aonde os preços serão de 1600, 1200 e 800 \textbf{NOK} (ORG), o que equivale respectivamente a 195.38, 146.54 e 97.67 euros. Ao todo, estiveram disponíveis 90 mil ingressos para estes espetáculos. Nas semifinais e final, houve uma espectáculo durante 45 minutos, antes do início da cerimónia eurovisiva. No dia em que os ingressos foram colocados à venda, cerca das 9h00 do dia 8 de Fevereiro de 2010, bastaram apenas vinte minutos para que os ingressos para a final do concurso se esgotassem completamente. Já os bilhetes para os restantes espectáculos também tiveram grande procura, não atingindo o máximo durante a manhã do dia 8, mas tendo esgotado pouco tempo depois. A disposição da sala de espetáculo também foi revelada, tendo esta sido dividida em 21 partes na arena, mais 12 nas bancadas em redor e por fim, mais 7 camarotes \textbf{VIP} (ORG). Também a forma do palco foi revelado no desenho, sendo que este terá forma circular, e entraria pela audiência.

A \textbf{Cidade de Oslo} (LOC) , é servida por um sistema de metropolitano que possui 84,2 km de extensão, divididos por 6 linhas, que têm 105 estações no total (no entanto apenas dezesseis delas são subterrâneas ou dentro de outro tipo de instalações). O \textbf{metropolitano de Oslo} (MISC) é operado pela \textbf{Oslo T-banedrift} (MISC) com contrato da \textbf{Autoridade de Transportes Ruter} (LOC). As linhas de metropolitano consistem em seis linhas, que passam todas pelo centro da cidade. Diariamente, 268 mil pessoas utilizam o sistema. Para que haja o serviço em todos os quinze bairros de \textbf{Oslo} (LOC), existem duas linhas que vão para \textbf{Bærum} (LOC) . A última expansão do sistema data de 2006 , e até 2010 toda a frota de comboios/carruagens será substituída. O sistema foi inicialmente inaugurado, em 31 de \textbf{Maio de 1898} (MISC) , como sendo um sistema de trens suburbanos, e apenas em 22 de maio de 1966 , é que passou a ser um sistema de metropolitano. Além do metropolitano, \textbf{Oslo possíu} (PER) um sistema de ônibus/autocarros elétricos. Todos eles, incluindo o metropolitano, são geridos pela \textbf{AS Oslo Sporveier} (MISC). Oslo possui três aeroportos, sendo o principal deles Oslo lufthavn, \textbf{Gardermoen} (LOC) . Os outros dois são o \textbf{Aeroporto de Sandefjord} (LOC), \textbf{Torp} (LOC) e o \textbf{Aeroporto de Moss} (LOC), \textbf{Rygge} (PER) . Além desses serviços, \textbf{Oslo} (LOC) também tem linhas do \textbf{TGV} (MISC) , assim como todo o resto do país. A \textbf{Estação Central de Oslo} (LOC) situa-se no "coração" da cidade (ao lado do \textbf{Oslo Specktrum} (MISC) ). Oslo também é servida por várias estradas, tanto nacionais como internacionais. Um programa de aluguel/aluguer de bicicletas está operacional desde \textbf{Abril de 2002} (MISC). Com um cartão de crédito, os usuários do sistema têm acesso a bicicletas em mais de noventa estações da cidade.

Nos vários dias de espetáculos e ensaios, entre 15 e 30 de \textbf{Maio} (LOC), circulariam pela \textbf{Cidade de Oslo} (LOC) vários ônibus/autocarros que estiveram ao dispor de todas as delegações, jornalistas, fãs e outros. Estes fizeram o percurso entre os muitos hotéis oficiais da \textbf{Eurovisão 2010} (MISC) e a \textbf{Telenor Arena} (PER) (e vice versa).



Votação

Em 2010 o sistema de votação foi modificado, porém só nas semifinais. Em 2008 e 2009 , nas semifinais, apenas existia um wildcard do júri, isto é, o público escolhia nove países para a final, e o júri apenas um. A partir deste ano, a votação das semifinais adoptou o mesmo sistema da final, sendo a votação 50\% televoto e 50\% júri. Outra regra que foi alterada, e que marca a maior mudança na \textbf{Eurovisão} (LOC) desde a substituição do antigo sistema de votação, é o facto de a votação deixar de ser efectuada apenas durante 15 minutos no fim do espetáculo. Tal como ocorre desde 2003 no Festival Eurovisão da Canção Júnior , a votação iniciou-se assim que a primeira música soou no palco de \textbf{Oslo} (LOC), até ao fim da actuação da última entrada, mais os habituais quinze minutos no final do espectáculo, cronometrados. Este novo modelo foi utilizado nos três espectáculos, e visa aumentar o número de votos, de tempo de votação e de interesse pelo festival. Sendo assim, as regras que valem a partir da edição de 2010, são as seguintes :

Para assegurar que as regras são cumpridas, a \textbf{EBU} (MISC) enviará para cada país, um auditor externo, que irá monitorar tudo o que ocorrer em relação ao festival.

No fim, depois da votação do público e dos júris, os 10 países com mais votos, passam a ter uma pontuação de 1 a 12 respeitando os votos anteriores. No fim soma-se ambos os valores que vão de 1 a 12 nas pontuações, e volta-se a pôr por ordem de pontuações determinando assim o grande vencedor. Para se ter uma ideia melhor de como será o sistema de votação, observe os seguintes exemplos

\textbf{Predefinição:Tabelas de Exemplo} (MISC) do \textbf{Novo Sistema de Votação do Festival Eurovisão da Canção} (LOC)

Em caso de dois ou mais países terem igual número de pontos no resultado definitivo, tanto na final como na semifinal, é usada a regra de desempate aonde aquela canção que mais países votaram. Se isso persistir o critério de desempate será o número de notas 12 e indo pras notas 10 por diante.

De lembrar que todas estas medidas foram aplicadas em resposta a algumas televisões responsáveis pela divulgação do concurso nos países europeus, que continuavam a queixar-se do voto político, vizinho e diáspora. Com estas meidas, a \textbf{União Europeia de Radiodifusão} (ORG) espera tornar o concurso mais justo e imparcial, de modo a evitar conflitos devido às classificações.

Tal como em todos os anos, desde que o festival se rege por um sistema de televoto, cada televisão anfitriã nacional, teve que escolher um porta-voz, que divulgou os votos do júri e do televoto do seu país. As pontuações de 1 a 7 apareceram automaticamente no quadro de votações, e posteriormente, o porta voz anunciou os países que receberam os 8, 10 e 12 pontos do seu país. Os resultados do \textbf{Festival Eurovisão da Canção} (MISC), foram posteriormente colocados no site oficial, pelo \textbf{Executivo Superior} (LOC) da \textbf{EBU} (MISC). Os resultados foram dados a conhecer apenas com as pontuações do público e do júri juntas, e à posteriori , cada país deverá revelar as suas pontuações individuais do júri e do televoto, sendo que o primeiro a fazê-lo foi a \textbf{Roménia} (LOC). A lista das pessoas que transmitiram os pontos atribuídos pelo seu país e pelo júri, aos outros países participantes no \textbf{Festival Eurovisão da Canção} (MISC) 2009, de cada país, pode ser observada na ligação no início desta secção (porta-vozes dos júris e televoto).



Países classificados para a final

Tal como em anos anteriores, para a final, cinco países foram automaticamente confirmados na final, sendo eles:



\textbf{Participações individuais} (MISC)Durante cerca de 10 meses, todos os países foram escolhendo os seus representantes, assim como as músicas que os mesmos interpretaram em \textbf{Oslo} (LOC). Para realizar tal selecção, cada país utilizou o seu próprio processo de selecção. Alguns optaram pela selecção interna, que consiste em a televisão organizadora daquele país fazer a escolha; no entanto, por vezes apenas o artista é seleccionado internamente, e a música não. Outros países (a maioria), utilizou um programa de televisão para seleccionar a sua entrada. Quartos de final, semifinais, second-chances e finais, foram realizados durante nos meses antecedentes do \textbf{Festival} (LOC) (até \textbf{Março} (LOC)) na maioria dos países europeus, cada um com o seu processo próprio (também a internet foi utilizada na fase de escolhas, como foi o caso de \textbf{Portugal} (LOC) ). \textbf{Predefinição:Participações Individuais} (MISC) ESC2010



\textbf{Participantes} (MISC)O primeiro país a confirmar a sua participação no \textbf{Festival Eurovisão da Canção 2010} (MISC), foi a \textbf{Holanda} (LOC) , ainda antes da realização do \textbf{Festival Eurovisão da Canção 2009} (MISC) (com a confirmação da sua participação para 2010, o país também afirmou que participará em 2011 e 2012). Tal confirmação, para três anos seguidos do festival, deve-se ao facto da alteração da televisão transmissora do país. Outros países que confirmaram a sua participação ainda antes da realização do evento do ano passado, foram a \textbf{Eslováquia} (LOC) e a \textbf{Grécia} (LOC) . Já durante os meses de \textbf{Maio} (MISC), \textbf{Junho} (LOC), \textbf{Julho} (LOC), \textbf{Agosto} (LOC), \textbf{Setembro} (LOC), \textbf{Outubro} (MISC) e Novembro, a maioria dos países confirmou se participaria ou não na edição de 2010, a realizar-se em \textbf{Oslo} (LOC) . O dia 15 de dezembro de 2009 , foi definido como o último dia para confirmações e desistências sem nenhuma penalização financeira. As primeiras selecções nacionais para os representantes dos países, começaram entre \textbf{Setembro e Outubro} (MISC) de 2009, com centenas de castings realizados por várias televisões, assim como o recebimento de várias maquetes para representar os vários países. Vários artistas manifestaram o seu interesse por representar o seu país, como aconteceu na \textbf{Grécia} (LOC), em \textbf{Andorra} (LOC) , ou na \textbf{Moldávia} (LOC) , e sucedendo-se um pouco por vários países na \textbf{Europa} (LOC) (incluindo \textbf{Espanha} (LOC) , \textbf{Reino Unido} (LOC) , \textbf{Turquia} (LOC) , \textbf{Arménia} (LOC) , etc). Pela primeira vez, em vários países, fãs criaram petições on-line para levar um determinado artista ao festival, como por exemplo \textbf{Portugal} (LOC) (pioneiro nesta matéria). No entanto, 2010 foi um ano em que a maioria dos países utilizará processos de selecção livre sendo mais difícil, assim, a qualquer artista individual conseguir representar o país, tendo que passar pela votação e escolha a ser realizada pelo público. Pela primeira vez, em alguns países (como na \textbf{Bielorrússia} (LOC)), o representante será escolhido pelo \textbf{Presidente} (MISC) do próprio país (ou pelo menos este terá uma última palavra quanto à escolha). O primeiro país a revelar o seu represente foi a \textbf{Bulgária} (LOC) , a 18 de \textbf{Outubro} (MISC). \textbf{Miroslav Kostadinov} (PER) tornou-se assim o primeiro artista conhecido oficialmente do \textbf{Festival Eurovisão da Canção 2010} (MISC). Países antigos, como a \textbf{Suécia} (LOC), \textbf{Noruega} (LOC), \textbf{Portugal} (LOC) e \textbf{Holanda} (LOC), continuaram com os seus tradicionais e conhecidos processos de selecção, que acompanham os países desde a sua estreia na \textbf{Eurovisão} (LOC). Outros países utilizaram novos e/ou recentes modelos de selecção, como é o caso da \textbf{Espanha} (LOC), do \textbf{Reino Unido} (LOC), da \textbf{Alemanha} (LOC), da \textbf{Finlândia} (LOC), entre outros. Em 2010, e apesar do aumento do número de selecções abertas, o número de países a utilizar escolhas internas foi de seis, entre eles a \textbf{França} (LOC), com mais de cinquenta participações, quase todas elas seleccionadas internamente. O país que causou maior impacto com o seu modelo de selecção para 2010, foi \textbf{Portugal} (LOC), ao introduzir, 45 anos depois da primeira edição, duas semifinais no \textbf{Festival RTP da Canção} (MISC), e ao passá-lo para a segunda maior arena de espetáculos \textbf{em Portugal} (LOC), o \textbf{Campo Pequeno} (ORG). Apesar de algumas selecções antes de Dezembro, a maioria dos países escolheu o seu representante em \textbf{Dezembro de 2009} (MISC), \textbf{Janeiro} (LOC), \textbf{Fevereiro} (LOC) e \textbf{Março de 2010} (MISC), sendo que o maior número de finais se realizou em \textbf{Março de 2010} (MISC), com \textbf{Israel} (LOC) a ser o último país a escolher o seu representante, no dia 14 do mesmo mês, e a \textbf{Bósnia e Herzegovina} (LOC), que no mesmo dia apresenta publicamente a sua canção representante. No dia 31 de dezembro de 2009 , a \textbf{União Europeia de Radiodifusão} (ORG) revelou a lista final de participantes no Festival Eurovisão da Canção 2010, que através das confirmações apresentadas durante mais de um ano, contou com 39 países integrantes, apesar dos 54 convidados pela \textbf{EBU} (MISC). Com esta revelação, encerrou-se a época de confirmações de participações e não participações para o \textbf{Festival Eurovisão da Canção 2010} (MISC), a 55 edição do mesmo. O último país a pronunciar-se para a primeira edição da segunda década do século \textbf{XXI} (MISC), foi a \textbf{Lituânia} (LOC).



\textbf{Organização} (ORG)A 7 de fevereiro de 2010 , realizou-se o sorteio para decidir que países participarão na primeira ou na segunda semifinal, em \textbf{Oslo} (LOC) , \textbf{Noruega} (LOC) , já a ordem de votação para a final do festival, foi apenas sorteada a 22 de março de 2010 . Baseando-se no mesmo sistema de eventos do Festival Eurovisão da Canção 2008 , todos os países foram separados em cinco potes individuais baseados no seu historial de votação de concursos anteriores e na sua localização geográfica. O sorteio foi implementado (ou criado), para garantir que os países que dispostos a dar a sua pontuação uns aos outros na competição, não participem na mesma semifinal. Destes cinco potes, saíram então as listas dos países que participariam na primeira e na segunda semifinal respectivamente, ficando assim todos os países divididos em dois grupos, da maneira mais justa possível. Foi ainda criado um sexto pote, apenas para a \textbf{Alemanha} (LOC) , \textbf{Espanha} (LOC) , \textbf{França} (LOC) , \textbf{Noruega} (LOC) e \textbf{Reino Unido} (LOC) , apenas com o intuito de definir que semifinal é que teriam que transmitir, na qual poderiam votar. Uma das regras, que também ficou estabelecida, desde a criação de duas semifinais, é que cada país é obrigado a transmitir a semifinal em que ficou colocado, e pode, se quiser, transmitir a outra semifinal.

Em baixo estão todos os seis potes, com os respectivos países que integraram cada um, para o sorteio da distribuição das semifinais:

Após o sorteio das semifinais, a 7 de fevereiro de 2010 , foi possível efectuar a divisão dos países em duas semifinais. O sorteio ocorreu no centro de \textbf{Oslo} (LOC), e foi presidido por pelo director de imprensa da televisão anfitriã \textbf{Peter Svaar} (PER) . As bolas foram retiradas de cada pote pelo chefe da polícia de \textbf{Oslo} (LOC). O sorteio foi supervisionado pelo responsável máximo do \textbf{Festival Eurovisão da Canção} (MISC), \textbf{Svante Stockselius} (PER) e pela directora jurídica da \textbf{EBU} (MISC) .

Apesar do festival ter o seu primeiro espetáculo apenas no dia 25 de maio de 2010 , os trabalhos na \textbf{Telenor Arena} (LOC) começaram a 20 de abril de 2010, com o início da montagem do palco para a \textbf{Eurovisão 2010} (MISC). Toda a infraestrutura esteve concluída no dia 10 de maio de 2010, dia em que se iniciaram os trabalhos de aprimoração do som e criação dos gráficos para as actuações dos 39 países a concurso, assim como para as aberturas e intervalos dos três espectáculos. No dia seguinte à final, a 30 de maio de 2010 iniciou-se a desmontagem de todas as estruturas no local. No dia 3 de junho de 2010, a \textbf{NRK} (ORG) entregou a \textbf{Telenor Arena} (PER) ao \textbf{Telenor Group} (ORG) , empresa responsável pela \textbf{arena} (PER).



\textbf{Eurovision Week} (MISC)A \textbf{Eurovision Week} (MISC) (ou \textbf{Semana Eurovisão} (MISC)), é a famosa semana em que decorre o evento Eurovisão em questão. Para 2010, a \textbf{NRK} (ORG), televisão responsável pela organização do \textbf{Festival Euroisão da Canção} (MISC) do mesmo ano, apresentou a \textbf{Câmara Municipal de Oslo} (LOC) como sede local para a festa de boas-vindas da \textbf{Eurovisão 2010} (MISC), que teve lugar no dia 23 de \textbf{Maio de 2010} (MISC), domingo, um dia antes do primeiro ensaio geral para a 1 semifinal. As famosas after-parties (depois da festa) tiveram lugar noutra parte da cidade, onde as delegações de cada país puderam festejar depois dos espetáculos, assim como os fãs. Durante o resto da semana ocorreram outras festas, assim como os ensaios para o festival.

A época eurovisiva foi aberta oficiarlmente no dia 12 de \textbf{Maio de 2010} (MISC), com os primeiros ensaios a decorrer no mesmo dia e até ao dia 22 de \textbf{Maio} (LOC). No dia 24 de \textbf{Maio} (LOC), foi realizado o primeira e o segundo ensaio para a 1 semifinal do concurso, no dia seguinte foi realizado um terceiro ensaio para a mesma semmifinal, e nesse dia ocorreu o primeiro espectáculo de 2010. Durante o dia 26 foi a vez dos dois primeiros ensaios da 2 semifinal, e no dia 27 teve lugar um terceiro ensaio, com o evento a ter ocorrido nesse mesmo dia. Para a final, os dois primeiros ensaios gerais foram realizados no dia 28 de \textbf{Maio} (LOC), e o último ensaio geral de 2010 foi a \textbf{29 de Maio} (MISC), dia da final do \textbf{Festival Eurovisão da Canção 2010} (MISC).

A primeira celebração oficial do \textbf{Festival Eurovisão da Canção 2010} (MISC), foi a já famosa \textbf{Eurovision Opening Party} (MISC) , sediada, em 2010, na \textbf{Câmara Municipal de Oslo} (LOC), e decorreu no dia 23 de \textbf{Maio de 2010} (MISC). A cerimónia contou com vários artistas, ex-representantes na \textbf{Eurovisão} (LOC), e teve como participante principal, a \textbf{Orquestra Metropolitana de Oslo} (LOC).



\textbf{Festival Nota} (PER): As pontuações estão divididas em \textbf{Total} (LOC) (T) , \textbf{Televoto} (PER) (\textbf{Tv} (LOC).) e \textbf{Júri} (MISC) (\textbf{J} (LOC)) .

A ordem de divulgação dos votos, na final do \textbf{Festival Eurovisão da Canção 2010} (MISC), foi decidida através de um sorteio que se realizou no dia 23 de \textbf{Março em Oslo} (MISC). Na noite da final, a \textbf{29 de Maio} (MISC), depois dos últimos 15 minutos de votação, e de mais 15 minutos de intervalo, a partir de \textbf{Oslo} (LOC), da arena onde vai ser realizado o festival, vai ser realizada uma ligação em directo para cada país participante na edição de 2010 (incluindo os que não passaram à final). Geralmente as ligações são feitas para a capital do país em questão, que tem o seu porta-voz num estúdio, com um fundo que consiste numa fotografia ou numa imagem real do panorama da cidade de onde estão ligados. Assim, cada porta-voz tem cerca de dois minutos para divulgar as votações nacionais, o que demora mais de uma hora.



\textbf{Transmissão do Festival} (MISC)O \textbf{Festival Eurovisão da Canção} (MISC) , é o programa não-desportivo que consegue juntar o maior número de telespectadores ao mesmo tempo em todo o \textbf{Mundo} (LOC) (somente tem audiência menor do que as \textbf{Cerimónias dos Jogos Olímpicos} (MISC) , e as finais de \textbf{Campeonatos Europeus de Futebol} (MISC) e \textbf{Mundiais} (MISC) ). Todos os anos, mais de 125 milhões de pessoas por toda a \textbf{Europa} (LOC) e de todo mundo vêem o espetáculo, traduzido em directo por um comentador escolhido pela sua televisão transmissora, membro da \textbf{EBU} (MISC) . Porém, o festival para além de ser transmitido televisivamente, também o é transmitido através de estações de rádio, e mais recentemente em directo (e posteriormente em vídeo on-demand), no site oficial do festival, através da eurovision.tv . Todos os anos, cerca de 50 televisões, 12 rádios e um canal oficial de televisão on-line ( eurovision.tv ), transmitem o festival para mais de 100 milhões de pessoas, no entanto o número vem a aumentar todos os anos.

Todas as televisões que enviam um representante para uma edição do \textbf{Festival Eurovisão da Canção} (MISC), são obrigados a transmitir pelo menos a semifinal em que entram e a final (à excepção dos \textbf{Big4} (MISC), que se quiserem podem apenas transmitir a final). No entanto, os países podem escolher também transmitir a semifinal onde não participam, o que é adoptado pela maioria, mesmo que por vezes a semifinal onde não participa, seja transmitida noutro horário. Em baixo encontra-se a lista das cadeias televisivas e respectivos canais, que emitirão o festival:

Foi possível assistir ao concurso em qualquer parte do mundo, através da internet. Para além disso, várias televisões de vários países não europeus, transmitiram o festival. Também várias televisões Europeias transmitiram o festival nos seus canais internacionais. Em baixo é possível observar uma lista com os canais que transmitiram o festival, para além fronteiras europeias:

Pela terceira vez, o festival também foi transmitido em canais de alta definição. Alguns países, através do seu canal de alta definição, possibilitaram que outros países assistam ao festival em HD.

Pela segunda vez, algumas televisões transmitiram o festival através da internet. Para além disso, o site oficial do Festival Eurovisão da Canção , eurovision.tv , também transmitiu os três espectáculos em directo para todo o \textbf{Mundo} (MISC) (mas sem qualquer tradução, ou número de telefone para votação).

Tal como no ano anterior, algumas cadeias de telemóveis e de televisões irão transmitir o \textbf{Festival Eurovisão da Canção} (MISC), em directo para telemóveis, através dos seus canais de \textbf{Mobile TV} (MISC) . As operadoras móveis que o farão são as seguintes:

Apesar de já não ser a maneira mais utilizada para emitir nem seguir o festival, algumas cadeias de rádio (pertencentes às cadeias televisivas), transmitem o festival em directo. Tal como em todos os anos anteiores, 2010 não foi excepção, e também teve cobertura via rádio.

Como já é hábito, desde o início das transmissões do festival, cada televisão membro, escolhe um ou mais comentadores, que em directo, traduzem o festival para a língua materna do país a que pertencem. Na ligação acima, é possível observar uma tabela com todos os comentadores que cada país levou a \textbf{Oslo} (LOC), para reportar o festival, que é apenas apresentado em inglês e francês (porém, marioritariamente em inglês), no local onde está a decorrer (por vezes surgem apresentações na língua do país anfitrião, ou em inglês).



Produtos oficiais

O principal patrocinador do \textbf{Festival Eurovisão da Canção 2010} (MISC), foi o \textbf{Telenor Group} (ORG) , uma das maiores companhias de redes móveis da \textbf{Europa} (LOC). Como tal, o grupo chegou ao acordo com a \textbf{EBU} (MISC) , para produzir e vender pelos seus mercados (oito na \textbf{Europa} (LOC)), o \textbf{Telemóvel Oficial do Festival Eurovisão da Canção} (LOC) . Os seus mercados \textbf{Europeus} (LOC) são a \textbf{Noruega} (LOC), \textbf{Dinamarca} (LOC), \textbf{Suécia} (LOC), \textbf{Finlândia} (LOC), \textbf{Hungria} (LOC), \textbf{Montenegro} (LOC), \textbf{Sérvia} (LOC), \textbf{Ucrânia} (LOC) e \textbf{Rússia} (LOC). O \textbf{CD} (MISC) oficial do festival, com todas as 39 músicas a concurso foi colocado à venda a 17 de \textbf{Maio} (LOC), enquanto que o DVD será lançado apenas a 18 de \textbf{Junho} (LOC). Outros produtos relacionados com o eventos (livros, roupa, souvenirs, etc), ficaram à venda a 23 de \textbf{Março} (PER).



\textbf{Factos} (MISC) e \textbf{Controvérsias} (MISC)Inicialmente as datas anúnciadas para o festival eram 18 de maio (1 semifinal), 20 de maio (2 semifinal) e 22 de maio (grande final) de 2010 . O anúncio foi feito durante a último conferência de imprensa do Festival Eurovisão da Canção 2009, logo a seguir à final do mesmo. No entanto, a data da final gerou polémica desde a sua apresentação, pois no dia 22 de \textbf{Maio} (LOC), seria também a final da \textbf{Champions League} (MISC) , uma das competições de futebol mais importantes do mundo. A final da \textbf{Liga dos Campeões} (MISC) geralmente é efectuada numa quarta-feira, no entanto em 2010, a organização resolveu mudar a data para um sábado. Para que não houvesse problemas com as televisões europeias quanto às prioridades de transmissão, a \textbf{EBU} (MISC) viu-se obrigada a alterar as data do Festival Eurovisão da Canção 2010. Caso o dia da final se mantivesse, várias televisões europeis (como a \textbf{RTP} (ORG) de \textbf{Portugal} (LOC), por exemplo), teriam de desistir do festival, para transmitir o evento desportivo devido aos contratos previamente acordados, o que não acontece com a \textbf{Eurovisão} (LOC), em que todos os anos tem de ser realizada uma inscrição no concurso. Uma das hipóteses para a não desistência de muitos países seria a de estes emitirem o festival não no seu canal principal, mas sim num secundário (como por exemplo, em \textbf{Portugal} (LOC), o concurrso em vez de ser emitido pela \textbf{RTP 1} (MISC) , seria emitido pela \textbf{RTP 2} (MISC) , ambas pertencentes à \textbf{RTP} (ORG)).

Ao longo de oito meses, mais de cinquenta países europeus confirmaram ou desconfirmaram a sua participação no Festival Eurovisão da Canção 2010. As seguintes categorias incluem as referências às estreias, regressos e retiradas esperadas, anúnciadas, confirmadas e não confirmadas para a \textbf{LV} (ORG) edição da \textbf{Eurovisão da Canção} (MISC).

Ao longo da história, mais de 1500 artistas pisaram o palco \textbf{Eurovisão} (LOC), vários destes artistas repetiram a sua experiência eurovisiva uma ou mais vezes. Em 2010, os repetentes foram:



Outros \textbf{Factos de Registo} (MISC)O site aqa.63336.com Arquivado em 22 de julho de 2010, no \textbf{Wayback Machine} (LOC) . fez uma análise a quem mais e menos falhou notas durante a final do \textbf{Festival} (MISC). Caso o vencedor fosse quem menos falhou, teria sido a \textbf{Bélgica} (LOC). Abaixo, uma lista, ordenada de quem mais falhou para quem menos falhou:



Festival Eurovisão da Canção 2011



Países que classificaram-se para a final \textbf{Países} (MISC) que não se classificaram para a final \textbf{Países} (MISC) que participaram no passado, mas não em 2011

O \textbf{Festival Eurovisão da Canção de 2011} (MISC) (em inglês : \textbf{Eurovision Song Contest} (ORG) 2011 ; em francês : \textbf{Concours Eurovision} (MISC) de la chanson 2011 ; em alemão : \textbf{Liederwettbewerb der Eurovision} (MISC) 2011 ) foi o \textbf{56. Festival} (MISC) Eurovisão da Canção na história. O festival será realizado na \textbf{Esprit Arena} (LOC) em \textbf{Düsseldorf} (LOC) , com as datas oficiais, avançadas pela \textbf{NDR} (ORG) e pela \textbf{EBU} (MISC) de 10 e 12 de maio para as duas semifinais e 14 de maio de 2011 para a grande final. Originalmente, o concurso seria realizado a 17, 19 e 21 de maio, mas foram modificadas devido à final da \textbf{Taça da Alemanha} (MISC) ser realizada no dia 21 de maio . Em 2011, pela segunda vez na história da \textbf{Eurovisão} (LOC), um cantor vencedor da edição anterior voltou a concorrer pelo seu país. \textbf{Lena Meyer-Landrut} (PER), que havia vencido a edição de Oslo, representou, novamente, a \textbf{Alemanha} (LOC) em 2011. Tal tinha acontecido apenas com \textbf{Lys Assia} (PER) a primeira vencedora da \textbf{Eurovisão} (LOC), em 1956, que representou o seu país (a \textbf{Suíça} (LOC) ) em 1957 e 1958. O vencedor do \textbf{Festival} (LOC) foi o \textbf{Azerbaijão} (LOC) , que pela primeira vez venceu o \textbf{Festival} (LOC), à sua quarta participação. No seu regresso após 14 anos de ausência, a \textbf{Itália} (LOC) conseguiu um brilhante segundo lugar, a apenas 32 pontos do vencedor



\textbf{Importância} (PER) histórica

Esta foi a terceira vez que a \textbf{Alemanha} (LOC) sediou o \textbf{Festival Eurovisão da Canção} (MISC) , uma vez que sediou em 1957 (por indicação, porque, na primeira edição, a regra de que o vencedor da edição anterior sedia o próximo festival não existia) e em 1983, quando ganhou o \textbf{Festival Eurovisão da Canção de 1982} (MISC) . Esta foi também a primeira vez desde 1998 que um país da \textbf{Europa Ocidental} (LOC) sediou o \textbf{Festival} (LOC), por que aconteceu, duas vezes no \textbf{Médio Oriente} (LOC) , quatro vezes na \textbf{Escandinávia} (LOC) e seis vezes na \textbf{Europa} (LOC) Oriental desde 1999. Esta foi, também, a primeira vez desde a introdução da regra dos \textbf{Big 5} (MISC) , em 1999, que um membro do "grupo" sediou o evento, e foi o primeiro país do grupo a sediar um festival desde 1998. Foi ainda a primeira vez, na história, que a \textbf{Alemanha} (LOC) sediou o \textbf{Festival da Eurovisão} (MISC) como um país unificado.



\textbf{Local} (MISC)A \textbf{Alemanha} (LOC) possuía várias cidades com condições reais de sediar o evento, numa situação inédita. A cidade-sede da próxima edição do festival não foi anunciada durante a conferência de imprensa dos vencedores da edição anterior. A última edição na qual não foi realizada numa capital foi a anterior em 2010 em \textbf{Baerum} (LOC) , um subúrbio de \textbf{Oslo} (LOC) ,a capital da \textbf{Noruega} (LOC) . Após a vitória do ano anterior, foram consideradas várias hipóteses, desde \textbf{Berlim} (LOC) , passando por \textbf{Hanôver} (LOC) , \textbf{Gelsenkirchen} (LOC) , \textbf{Colônia} (LOC) , \textbf{Düsseldorf} (LOC) , \textbf{Frankfurt Am Main} (LOC) , \textbf{Munique} (LOC) e \textbf{Hamburgo} (LOC) , tendo cada uma o apoio também de vários políticos locais. Esperava-se o anúncio da cidade-sede no final de junho de 2010. Mas, a 31 de \textbf{Maio} (PER) se teve o conhecimento de que a escolha final seria \textbf{Hamburgo} (LOC) ou \textbf{Hanôver} (LOC), devido ao fato de que as arenas das duas cidades estariam vagas para as datas do \textbf{Festival} (PER). Esta é a primeira vez em anos recentes que um país tem tantas cidades com condições reais para sediar o festival. Em 21 de Agosto de 2010 foram anunciadas as quatro cidades candidatas para sediar o \textbf{Festival de 2011} (MISC). São elas :

Foi reportado na mesma notícia que as cidades de \textbf{Gelsenkirchen} (LOC) , \textbf{Colônia} (LOC) , \textbf{Frankfurt} (LOC) e \textbf{Munique} (LOC) também mostraram interesse em sediar o \textbf{Festival} (LOC), mas não inscreveram uma candidatura.

Os planos de cada cidade candidata eram os seguintes :

A 2 de outubro de 2010 , foi anunciada a eliminação de \textbf{Hamburgo} (LOC),por não ter condições financeiras de arcar com o evento.

Finalmente, no dia 12 de outubro de 2010, \textbf{Düsseldorf} (LOC) foi anunciada como a escolhida devido a ESPRIT \textbf{Arena} (PER) ter as melhores condições de suportar o evento. O local foi inaugurado em 2004 para substituir o antigo \textbf{Rheinstadion} (LOC) e tem uma capacidade máxima para 60 mil pessoas. Durante os dias do \textbf{Festival} (LOC) e para os ensaios,a capacidade será reduzida a 35 mil pessoas.



Votação

Em 2010 o sistema de votação foi modificado, porém só nas semifinais. A partir desse ano, a votação das semifinais adotou o mesmo sistema da final, sendo a votação 50\% televoto e 50\% júri.< Outra regra que foi alterada, e que marca a maior mudança na \textbf{Eurovisão} (LOC) desde a substituição do sistema antigo de votação, é o fato do televoto ter somente 15 minutos disponível. Segue-se o mesmo modelo adotado no Festival Eurovisão da Canção Júnior , a votação teve início no momento em que a primeira música soou no palco de \textbf{Oslo} (LOC), até ao fim da última música, mais os habituais quinze minutos no final do espectáculo, cronometrados. Este novo modelo foi utilizado nos três espectáculos, visando um aumento expressivo do número de votos, e do interesse pelo festival. Sendo assim, as regras que valem a partir da edição de 2010, são as seguintes :

Para assegurar que as regras são cumpridas, a \textbf{EBU} (MISC) enviará para cada país, um auditor externo, que irá monitorar tudo o que ocorrer em relação ao festival.

Ao, depois da votação do público e dos júris, os 10 países com mais votos, passam a ter uma pontuação de 1 a 12 respeitando os votos anteriores. No fim soma-se ambos os valores que vão de 1 a 12 nas pontuações, e volta-se a pôr por ordem de pontuações determinando assim o grande vencedor. Para se ter uma ideia melhor de como será o sistema de votação, observe os seguintes exemplos

\textbf{Predefinição:Tabelas de Exemplo} (MISC) do \textbf{Novo Sistema de Votação do Festival Eurovisão da Canção} (LOC)

Em caso de dois ou mais países terem igual número de pontos no resultado definitivo, tanto na final como na semifinal, é usada a regra de desempate aonde aquela canção que mais países votaram. Se isso persistir o critério de desempate será o número de notas 12 e indo pras notas 10 por diante.

A final edição de 2011 terá, pela terceira vez consecutiva, um grupo de jurados. Um total de mais de 200 jurados avaliaram as canções do Festival Eurovisão da Canção 2011 (5 jurados por país). Os júris devem consistir na maior variedade possível de profissões relacionadas com o mundo do showbizz (cantores, produtores, etc), e deve ser diversificado em gênero, sexo, idade, status social, etc. Todos os membros dos \textbf{júris} (LOC) têm que ser cidadãos do país que está votando, e não podem estar ligados de qualquer forma aos representantes do seu país. Caso contrário, o seu voto pode ser anulado.

Tal como em todos os anos, desde que o festival se rege por um sistema de televoto, cada televisão anfitriã nacional, escolheu um porta-voz, que divulgou os votos do júri e do televoto do seu país. As pontuações de 1 a 7 apareceram automaticamente no quadro de votações, e posteriormente, o porta voz anunciou os países que recebem os 8, 10 e 12 pontos do seu país. Os resultados do \textbf{Festival Eurovisão da Canção} (MISC), foram posteriormente colocados no site oficial, pelo \textbf{Executivo Superior} (LOC) da \textbf{EBU} (MISC). Os resultados foram dados a conhecer apenas com as pontuações do público e do júri juntas, e à posteriori , cada país deverá revelar as suas pontuações individuais do júri e do televoto. A lista das pessoas que transmitiram os pontos atribuídos pelo seu país e pelo júri, aos outros países participantes no \textbf{Festival Eurovisão da Canção} (MISC) 2011, de cada país, está observada na ligação no início desta seção (porta-vozes dos júris e televoto).



\textbf{Participações individuais} (MISC)Durante cerca de 10 meses, todos os países escolheram os seus representantes, assim como as músicas que os mesmos interpretarão no festival. Para realizar tal selecção, cada país utilizou o seu próprio processo de selecção. Por vezes é usado o método da selecção interna (a televisão organizadora daquele país faz a escolha). Neste método acontece também que se escolha apenas o artista, e não a música (como aconteceu com a \textbf{Alemanha} (LOC), que escolheu \textbf{Lena} (PER) mas a música foi escolhida pelos espectadores. Ma a maior parte dos países utiliza um programa de televisão para seleccionar a sua entrada. Quartos-de-final, semifinais, second-chances e finais, foram realizadas durante os meses antecedentes do \textbf{Festival} (LOC) (a última em \textbf{Março de 2011} (MISC)) na maioria dos países europeus, cada um com o seu processo próprio (também a internet é utilizada na fase de escolhas, como tem sido o caso de \textbf{Portugal} (LOC) nos últimos dois anos).

O primeiro país a confirmar a sua presença nesta edição do \textbf{Festival Eurovisão} (MISC), foi a \textbf{Holanda} (LOC) , ao mesmo tempo que o país ficou confirmado para 2010 e 2012. Esta confirmação, para três anos, aconteceu graças à mudança da emissora responsável pela selecção e transmissão do \textbf{Concurso} (MISC). A \textbf{Bielorrússia} (LOC) , segundo país a confirmar, iria mudar de televisão em 2011, mas a emissora que iria ser responsável pelo novo processo não foi aceite pela \textbf{EBU} (MISC). O primeiro país a apresentar os planos para 2011 foi o \textbf{Chipre} (LOC) , onde um evento televisivo de talentos escolheu o representante do país para 2011. O primeiro país do grupo dos \textbf{Big 5} (MISC), com passagem directa à final, a confirmar sua participação, foi a \textbf{Espanha} (LOC) . Tal confirmação foi feita pelo director do canal televisivo espanhol, responsável pela representação do país, a \textbf{TVE} (ORG) , que assegurou na mesma altura que seriam feitas alterações no processo de selecção do representante em 2011, pois aconteceram polémicas na selecção nos três anos anteriores. Foi criado um padrão de qualidade das canções para o canal. Sendo assim, além da confirmação do país para o 56 \textbf{Festival Eurovisão da Canção} (MISC), ficou também confirmada a utilização de uma selecção aberta para a escolha do representante, como acontece desde 2007 no país.



\textbf{Festival Nota} (PER): As pontuações estão divididas em \textbf{Total} (LOC) (T) , \textbf{Televoto} (PER) (\textbf{Tv} (LOC).) e \textbf{Júri} (MISC) (\textbf{J} (LOC)) .



\textbf{Organização} (ORG)A 17 de \textbf{Janeiro} (LOC) de 2011 , realizou-se o sorteio para decidir que países participarão na primeira ou na segunda semifinal, em \textbf{Düsseldorf} (LOC) , \textbf{Alemanha} (LOC) , já a ordem de votação para a final do festival será realizado posteriormente. Todos os países foram separados em seis potes individuais baseados no seu historial de votação de concursos anteriores e na sua localização geográfica. O sorteio foi implementado (ou criado), para garantir que os países que dispostos a dar a sua pontuação uns aos outros na competição não participem na mesma semifinal. Destes seis potes, sairam então as listas dos países que participarão na primeira e na segunda semifinal, respectivamente, ficando assim todos os países divididos em dois grupos, da maneira mais justa possível. Foi ainda criado um sétimo pote, apenas para a \textbf{Alemanha} (LOC) , \textbf{Espanha} (LOC) , \textbf{França} (LOC) , \textbf{Itália} (LOC) e \textbf{Reino Unido} (LOC) , apenas com o intuito de definir que semifinal é que teriam que transmitir, na qual poderiam votar. Uma das regras, que também ficou estabelecida, desde a criação de duas semifinais, é que cada país é obrigado a transmitir a semifinal em que ficou colocado, e pode, se quiser, transmitir a outra semifinal.

Em baixo estão todos os sete potes, com os respectivos países que integraram cada um, para o sorteio da distribuição das semifinais :

Após o sorteio das semifinais, a 17 de \textbf{Janeiro} (LOC) de 2011 , foi possível efectuar a divisão dos países em duas semifinais.



Transmissão

Em baixo encontra-se a lista das cadeias televisivas e respectivos canais, que emitirão o festival:



\textbf{Factos} (MISC) e controvérsias

Ao longo da história, mais de 1500 artistas pisaram o palco \textbf{Eurovisão} (LOC).

Vários destes artistas repetiram a sua experiência eurovisiva uma ou mais vezes.



Festival Eurovisão da Canção 2012



Países que classificaram-se para a final \textbf{Países} (MISC) que não se classificaram para a final \textbf{Países} (MISC) que participaram no passado, mas não em 2012

O \textbf{Festival Eurovision de 2012} (MISC) (em inglês : \textbf{Eurovision Song Contest 2012} (ORG) ; em francês : \textbf{Concours Eurovision} (MISC) de la chanson 2012 ; em azeri : 2012 \textbf{Avroviziya Mahnı Müsabiqsi} (ORG) ) foi a 57 edição anual do evento. O festival foi realizado em \textbf{Baku} (LOC) , \textbf{Azerbaijão} (LOC) , depois da vitória desse país em \textbf{Düsseldorf} (LOC) , sendo representados pelo duo \textbf{Ell \& Nikki} (MISC) , com a música \textbf{Running Scared} (MISC) . O evento foi realizado no \textbf{Baku Crystal Hall} (LOC) , uma nova arena que foi construída na cidade. As duas semifinais foram realizadas em 22 e 24 de maio, enquanto a final foi em 26 de maio.

O país vencedor do festival foi a \textbf{Suécia} (LOC), com a canção "\textbf{Euphoria} (PER)" interpretada por \textbf{Loreen} (PER).



\textbf{Local Diferente da Alemanha} (MISC) , onde havia sido realizada a edição do ano anterior , o \textbf{Azerbaijão} (LOC) é um país pequeno. A única alternativa plausível seria realizar a próxima edição na sua capital, \textbf{Baku} (LOC) , mas o país ainda não tinha um local para sediar o concurso. Previa-se a construção de uma nova arena especialmente para o evento no centro da cidade. No entanto, em 19 de maio de 2011, os organizadores anunciaram que provavelmente seria usado o \textbf{Estádio Tofiq Bahramov} (LOC) , que tem capacidade para 37 mil pessoas sentadas, ou então o \textbf{Palácio Heydar Aliyev} (LOC) . A confirmação do local que iria sediar o \textbf{Concurso} (MISC) foi anunciada em 25 de janeiro de 2012, exatamente no dia do sorteio da definição das semifinais e de suas respectivas datas.



\textbf{Organização} (ORG)O sorteio para definir em qual semifinal cada país iria se apresentar foi realizado em \textbf{25 de Janeiro} (LOC) de 2012. Os países participantes (exceto \textbf{Alemanha} (LOC) , \textbf{Azerbaijão} (LOC) , \textbf{Espanha} (LOC) , \textbf{França} (LOC) , \textbf{Itália} (LOC) e \textbf{Reino Unido} (LOC) ) foram divididos em seis grupos, baseados na votação dos concursos anteriores. Em cada grupo, a primeira metade se apresenta na primeira semifinal, enquanto a segunda metade na segunda parte. Esse sorteio funciona como uma espécie de rascunho da ordem de apresentação nas semifinais, já que ocorre um segundo sorteio para definir a ordem de participação de cada um dos países participantes e seus respectivos ensaios.



\textbf{Festival Fonte} (MISC) (dos resultados): Página Oficial do Festival Eurovisão da Canção (\textbf{Azerbaijão} (LOC), \textbf{Itália} (LOC) e \textbf{Espanha} (LOC) votam nesta semifinal)

Fonte (dos resultados): Página Oficial do Festival Eurovisão da Canção (\textbf{França} (LOC), \textbf{Alemanha} (LOC) e \textbf{Reino Unido} (LOC) votam nessa semifinal. Nota: a \textbf{Alemanha} (LOC) fez um pedido especial para votar nessa semifinal, que foi deferido.)

Fonte (dos resultados): Página Oficial do Festival Eurovisão da Canção



Factos e \textbf{Controvérsias} (MISC)Ao longo da história, mais de 1500 artistas pisaram o palco \textbf{Eurovisão} (LOC), vários destes artistas repetiram a sua experiência eurovisiva uma ou mais vezes. Em 2012, os repetentes foram:



Festival Eurovisão da Canção 2013



Países que classificaram-se para a final \textbf{Países} (MISC) que não se classificaram para a final \textbf{Países} (MISC) que participaram no passado, mas não em 2013

O \textbf{Festival Eurovisão da Canção de 2013} (MISC) foi a 58 edição anual do evento. Foi realizado pela segunda vez em \textbf{Malmö} (LOC) , na \textbf{Suécia} (LOC) nos dias 14, 16 e 18 de maio de 2013. A \textbf{Suécia} (LOC) venceu a edição de 2012 , com a canção " \textbf{Euphoria} (PER) ", de \textbf{Loreen} (LOC) e foi, por isso, a organizadora da edição de 2013. Pela primeira vez desde 1995, teve apenas uma apresentadora,a comendiante \textbf{Petra Mede} (PER) .

As emissoras nacionais de televisão 39 países da \textbf{União Europeia de Radiodifusão} (ORG) participaram. A \textbf{Bósnia e Herzegovina} (LOC), \textbf{Eslováquia} (LOC), \textbf{Portugal} (LOC) e \textbf{Turquia} (LOC) anunciaram a sua retirada nesta edição.

O slogan foi "\textbf{We Are One} (MISC)" (Nós somos um), acompanhado por uma borboleta como logótipo



\textbf{Local} (MISC)A \textbf{Malmö Arena} (LOC) , em \textbf{Malmö} (LOC), foi anunciada pela host-broadcaster (emissora anfitriã) \textbf{Sveriges Television} (MISC) (\textbf{SVT} (ORG)) como o local do \textbf{Eurovision Contest} (ORG) 2013 em 8 de \textbf{Julho} (LOC) de 2012. Esta foi a quinta vez depois de 1975, 1985, 1992 e 2000, que o festival é realizado na \textbf{Suécia} (LOC) e a segunda vez na história que foi realizado em \textbf{Malmö} (LOC).

Após a final do festival em 2012, na cole(c)tiva de imprensa após a vitória de \textbf{Loreen} (LOC) , a dire(c)tora da televisão sueca \textbf{SVT Eva Hamilton} (ORG) disse aos repórteres que as cidades suecas de \textbf{Solna} (LOC) , \textbf{Gotemburgo} (LOC) e \textbf{Malmö} (LOC) seriam consideradas como possíveis locais para a edição de 2013. A arena mais recente, o \textbf{Ericsson Globe} (MISC) , utilizada no \textbf{Festival Eurovisão da Canção 2000} (MISC) , devido ao facto de estar alugada para o \textbf{Campeonato Mundial de Hóquei sobre o Gelo de 2013} (MISC) durante as datas preliminares do festival, e assim foi descartada.

As cidades candidatas, provisoriamente reservaram as as suas arenas e vários quartos de hotel, como parte de suas propostas para sediar a edição de 2013.

A SVT avaliou todas as possibilidades sobre a cidade-sede do festival em 2013. Uma alternativa levantada pelo jornal \textbf{Expressen} (ORG) foi a de fazer a competição em três locais diferentes - as semifinais em \textbf{Malmö} (LOC) e \textbf{Gotemburgo} (LOC) e a final em \textbf{Estocolmo} (LOC). A alternativa foi abortada por questões de logística envolvidas ao festival.

Em 20 de junho de 2012, a \textbf{SVT} (ORG) divulgou a dessistência de \textbf{Gotemburgo} (LOC), devido ao tradicional \textbf{Göteborg Horse Show} (ORG), evento tradicional de hipismo realizado no final de abril. Devido as demandas da \textbf{UER} (ORG), a \textbf{Arena} (PER) deveria estar desocupada no mesmo período. Houve também um consenso sobre as vagas disponíveis de hotéis no período e o número de eventos na cidade. O anúncio final do local, cidade e datas finais foi feito em 8 de julho de 2012.

Finalmente, em 8 de julho de 2012, a cidade de \textbf{Malmö} (LOC) foi anunciada como a cidade-sede do \textbf{Eurovision Song Contest} (ORG) 2013, que será na \textbf{Malmö Arena} (LOC) . As razões alegadas pela \textbf{SVT} (ORG) para a escolha da cidade foram o fa(c)to da infraestrutura estar em excelente condições e o desejo da \textbf{UER} (ORG) de voltar a realizar o festival em uma arena de tamanho menor do que as anteriores,para diminuir a escala do evento e melhorar a sua qualidade.

Tal como foi feito nos anos anteriores, desde a adoção do logotipo genérico do festival, a organização usará uma identidade visual própria para o desenvolvimento desta edição, a qual foi apresentada em 17 de janeiro de 2013 . O emblema principal desta é uma borboleta composta de diferentes cores e texturas, que segundo os organizadores pretendem demonstrar a diversidade cultural dos 39 países participantes,baseados no efeito borboleta . Junto ao logo,está o slogan \textbf{We} (MISC) are one (Somos todos um em português). A agência \textbf{Happy F\&B} (PER) foi a responsável pela criação do slogan, a qual teve a responsabilidade de criar uma identidade inovadora , na qual posteriormente será integrada ao site oficial,a decoração urbana, campanhas publicitárias e os gráficos da transmissão.

O design do palco, de autoria de \textbf{Viktor Brattström} (PER) e \textbf{Frida Arvidsson} (PER), é inspirado por suas palavras em diferentes aspectos das borboletas, arquitetura moderna e a alta costura. O cenário, circular e encimada por uma estrutura em forma de arco terá uma extensão de passarela que vai ligar com palco circular menor. O palco e a passarela vão ocupare a maior parte da superfície da \textbf{arena} (PER) e estão rodeados por espectadores em pé, "a interação entre espectadores e o público", segundo o produtor executivo do festival \textbf{Martin Österdahl} (PER). Outra diferença das edições anteriores é que o cenário não tem um fundo de leds, mas eles vão ficar no chão para criar desenhos em \textbf{3D} (MISC), no fundo serão utilizado 28 projetores de última geração para criar diferentes ambientes. Além disso, entre outros aspectos técnicos destaca a presença de 210 altofalantes, 40 quilômetros de cabos elétricos, 1.243 lâmpadas para iluminação e 24 quilômetros de cabos de fibra ótica.

Durante o festival, cada apresentação será precedida por uma breve introdução em vídeo (conhecido como "\textbf{postcard} (MISC)" ou "cartão postal") que mostrarão os artistas em seu dia a dia,em seus respectivos países.



\textbf{Participações Para saber} (MISC) os resultados da semifinal 2 está aqui- http://www.eurovision.tv/page/history/by-year/contest?event=1793

Para saber os resultados da grande final veja aqui- \textbf{http://www.eurovision.tv/page/history/by-year/contest?event=1773} (MISC)



Outros países

Ao longo da história, mais de 1500 artistas pisaram o palco \textbf{Eurovisão} (LOC), vários destes artistas repetiram a sua experiência eurovisiva uma ou mais vezes. Em 2013, os repetentes foram:



Festival Eurovisão da Canção 2014



Final : \textbf{Emmelie} (MISC) de \textbf{Forest} (LOC) (vencedora do ano passado 2013) e \textbf{Gaia Cauchi} (PER) (vencedora do \textbf{JESC 2013} (MISC))

\textbf{Países} (LOC) que classificaram-se para a final \textbf{Países} (MISC) que não se classificaram para a final \textbf{Países} (MISC) que participaram no passado, mas não em 2014

O \textbf{Festival Eurovisão da Canção de 2014} (MISC) (em inglês : \textbf{Eurovision Song Contest} (ORG) 2014 ; em francês : \textbf{Concours Eurovision} (MISC) de la chanson 2014 ; em dinamarquês : \textbf{Europæisk Melodi Grand Prix} (PER) 2014 ) foi a 59 edição anual do \textbf{Festival Eurovisão da Canção} (MISC) . Consistiu em duas semifinais nos nos dias 6 e 8 de maio de 2014 , e na final no 10 de maio. Esta foi a terceira vez que o festival aconteceu na \textbf{Dinamarca} (LOC) , em \textbf{Copenhaga} (LOC) , sendo que a última vez foi em 2001 . A \textbf{Croácia} (LOC) , a \textbf{Bulgária} (LOC) o \textbf{Chipre} (LOC) e a \textbf{Sérvia} (LOC) decidiram abandonar a competição mas poderão participar numa próxima edição. \textbf{Polónia} (LOC) e \textbf{Portugal} (LOC) regressaram com dois e um ano de ausência, respectivamente. A edição contou com trinta e sete participantes, o número mais baixo desde o \textbf{Festival Eurovisão da Canção 2007} (MISC) .



\textbf{Local} (MISC)O \textbf{Festival Eurovisão da Canção} (MISC) foi realizado no estaleiro \textbf{B\&W em Copenhaga} (MISC) . O antigo estaleiro passou por diversas melhorias necessárias para ser capaz de acolher o evento. Localizado na área de \textbf{Refshaleøen de Amager} (LOC), em \textbf{Copenhaga} (LOC), este complexo industrial de idade tornou-se um "playground para a indústria de cultura e experiências".

Com os seus espaços interiores enormes, o escritório de turismo da cidade, o \textbf{Wonderful Copenhagen} (LOC) afirmou que o local pode ser especificamente concebido para atender às necessidades da \textbf{DR} (ORG) e da \textbf{UER} (ORG) , responsáveis pela organização do festival.

Estes foram os locais candidatos para sediar o \textbf{Espetáculo :} (MISC)Em 17 de junho de 2013 foi divulgada a desistência de \textbf{Aalborg} (LOC) , porque segundo as autoridades municipais, não tinha capacidade para sediar o evento.

Em 20 de junho de 2013, foi divulgado que a cidade de \textbf{Frederícia} (LOC) desistiu do evento devido a sua péssima infraestrutura.

Em 28 de junho de 2013, os diretores do \textbf{Parken Stadium} (PER) anunciaram que o estádio estava se retirando da candidatura. Os jogos de futebol do clube local, \textbf{Copenhaga} (LOC), e da seleção dinamarquesa não podiam ser desmarcados ou transferidos para outro recinto, pelo calendário da \textbf{UEFA} (ORG) e da \textbf{FIFA} (ORG) .

As três cidades candidatas oficiais e os possíveis locais foram :

Em \textbf{2 de Setembro de 2013} (MISC) foi divulgado que a cidade de \textbf{Copenhaga} (LOC) iria sediar o evento.



\textbf{Participantes Valentina Monetta} (PER) , que já havia representado \textbf{San Marino} (LOC) em 2012 e 2013 , retornou para a competição pela terceira vez seguida. Isso faz de \textbf{Monetta} (PER) a quarta cantora principal a competir em três edições consecutivas do festival (e a única entre eles a nunca vencer em uma das ocasiões), seguindo \textbf{Lys Assia} (PER) (que representou a \textbf{Suíça} (LOC) em 1956 , 1957 e 1958 ), \textbf{Corry Broken} (PER) (que representou os \textbf{Países Baixos} (LOC) em 1956, 1957 e 1958) e \textbf{Udo Jürgens} (PER) (que representou a \textbf{Áustria} (LOC) em 1964 , 1965 e 1966 ). \textbf{Monetta} (PER) voltaria à competição mais uma vez, em 2017 . \textbf{Paula Seling} (PER) e \textbf{Ovi} (ORG) já haviam representado a \textbf{Roménia} (LOC) na edição de 2010 . As \textbf{Irmãs Tolmachevy} (PER) (\textbf{Tolmachevy Sisters} (PER)) , já haviam representado e ganhado pela \textbf{Rússia} (LOC) no \textbf{Eurovisão Júnior} (MISC) em 2006 .

Além disso, \textbf{Tamara Todevska} (PER) , que estava fazendo vocais de apoio para a \textbf{Macedônia} (LOC), já havia representado o país em 2008 . Ela representaria a \textbf{Macedônia do Norte} (LOC) anos depois, na edição de 2019 . \textbf{Martina Majerle} (PER) , que estava fazendo vocais de apoio para \textbf{Montenegro} (LOC) , já havia representado a \textbf{Eslovênia} (LOC) em 2009 e feito vocais de apoio para a \textbf{Croácia} (LOC) em 2003 , \textbf{Montenegro} (LOC) em 2008 e \textbf{Eslovênia} (LOC) em 2007 , 2011 e 2012.



Sistema de votação

Os telespectadores e os júris profissionais de cada país irão determinar, numa proporção de 50\%/50\% o resultado das duas semi-finais do concurso. Em cada semi-final, dez países qualificar-se-ão para a final.

Os espectadores de cada país participante em cada semi-final poderão votar via telefone, \textbf{SMS} (LOC) e aplicação oficial do evento (disponível em várias plataformas móveis). A votação será aberta após todos os países terem apresentado a sua performance, sendo que só se pode votar até vinte vezes pelo mesmo dispositivo e não se pode votar dentro do próprio país. A votação será encerrada 15 minutos após a abertura. O tele-voto determinará 50\% do resultado, sendo que os júris profissionais determinarão os restantes 50\%. O júri está impedido pelos regulamentos de votar no seu próprio país; este júri nacional é composto por cinco membros (incluindo um presidente) e é o mesmo júri que vai votar na \textbf{Final} (LOC).

O resultado final é a junção de tele-voto mais júri, e, em caso de empate entre músicas prevalecerá o tele-voto.

Os espectadores da \textbf{Itália} (LOC), \textbf{Dinamarca} (LOC), \textbf{Reino Unido} (LOC), \textbf{Espanha} (LOC), \textbf{Alemanha} (LOC) e \textbf{França} (LOC) também votam nas semifinais, sendo 3 dos finalistas em cada semi-final (o sorteio determinou que a \textbf{Espanha} (LOC), \textbf{França} (LOC) e \textbf{Dinamarca} (LOC) votarão na primeira, ao passo que a \textbf{Alemanha} (LOC), \textbf{Itália} (LOC) e \textbf{Reino Unido} (LOC) na segunda). Os dez países qualificados serão anunciados no final de cada semi-final, por ordem aleatória, não reflectindo o resultado do quadro de pontuação.

Para todos os países finalistas o resultado final da votação será classificado pelo tele-voto e júri. Os telespectadores poderão votar através das mesmas formas como nas semi-finais (aplicativo oficial, telefone e \textbf{SMS} (LOC)). A votação será aberta após a última canção ser apresentada e terminará quinze minutos depois. Estes votos determinarão metade do resultado e a outra metade restante pelos júris nacionais. O júri será o mesmo que votou numa das semi-finais; este irá assistir ao vivo e classificar todas as músicas com base no segundo ensaio final (\textbf{Dress Rehearsal} (MISC)).

A votação final nacional é obtido fundindo as duas classificações, e vai atribuir 12 pontos ao país com a melhor combinação do júri e do tele-voto, precedido de 10 pontos para o país com a segunda melhor classificação combinada, e assim sucessivamente até ao país que ficar em 10  no ranking combinado que receberá 1 ponto. Tradicionalmente apenas 1-12 pontos serão dadas, os países classificados fora do top-10 não receberá pontos. Além disso os telespectadores e \textbf{júris} (PER) não poderão votar no país que representam. Tal como nas semi-finais, se houver um empate entre dois ou mais músicas no ranking combinado entre tele-voto e júri, a canção que obter uma melhor classificação do tele-voto irá prevalecer sobre o outro. Com base na classificação total combinada de todas as músicas cada país irá distribuir 1 a 8, 10 e 12 pontos, e estes serão apresentados em directo durante a final por porta-vozes de todos os países participantes . Como de costume os pontos de um a sete aparecerão no ecrã e o porta-voz revelará oito, dez e doze pontos.

O ranking da divisão combinada dos votos dos telespectadores e do júri serão publicado pela autoridade executiva da \textbf{EBU do Festival} (MISC) Eurovisão da Canção no site oficial do evento dentro de quatro semanas após a final. Para proteger a justiça da votação, a \textbf{EBU} (MISC) não mostrará o ranking de divisão de tele-voto e júri por país. A publicação desses números, apenas irá destacar explicitamente o número mínimo de tele-votos necessários para tornar um resultado estatisticamente válido, para que não leve a uma influência desproporcional indesejada nos telespectadores nesses países nos próximos anos de concurso.

A votação do júri é monitorizada por um notório independente em cada país, alem deste observador da votação a \textbf{EBU} (MISC) tem o direito de enviar um observador independente adicional para a sessão do júri. Os júris nacionais são compostos por uma variedade de membros em termos de idade, sexo e de ocupação. Todos os membros do júri devem ser cidadãos do país que estão a representar, e nenhum membro do júri deve estar ligado a qualquer um dos participantes,tanto musical quanto artisticamente, de tal maneira que não podem votar independentemente. As emissoras participantes devem enviar uma carta de cumprimento das instruções de voto, juntamente com declarações assinadas por cada membro do júri afirmando que irão votar de forma independente.

Os nomes dos membros dos \textbf{júris} (LOC) devem ser revelados pelas emissoras participantes durante a \textbf{Final} (LOC), sendo que cada membro de cada júri nacional deve classificar todas as músicas do evento. O ranking combinado dos membros do júri de cada país determina o resultado do júri do referido país. Ao classificar cada música cada membro do júri terá que ter em conta a capacidade vocal do(s) artista(s), o desempenho em palco, a composição e originalidade da música, e a impressão geral da performance.

Sobre o sistema de televotos: A empresa com sede na \textbf{Alemanha} (LOC) GmbH tem sido parceira da votação a mais de anos,notavelmente desde 2004. A empresa reúne todos os votos dos telespectadores e votos do júri de todos os países e vai sendo acompanhado de perto por dois observadores independentes. Os seus sistemas são construídos para lidar com todos os votos recebidos de acordo com as regras do \textbf{Festival} (LOC), e pode detectar e excluir qualquer irregularidade sistemática.



Países classificados para a final

Tal como em anos anteriores, para a final, seis países foram automaticamente confirmados na final, sendo eles:



\textbf{Apresentadores} (PER)Depois de \textbf{Petra Mede} (PER) ter conduzido sozinha e com maestria a apresentação do Festival Eurovisão da Canção 2013 , a emissora dinamarquesa anunciou no dia 4 de Fevereiro de 2014 os três apresentadores da edição de 2014: \textbf{Nikolaj Koppel} (PER), \textbf{Lise Rønne} (PER) e \textbf{Pilou Asbæk} (PER) foram os escolhidos.Será a primeira vez que o trio irá apresentar o certame,será composto por dois homens e uma mulher. O anúncio foi feito no \textbf{Odd Fellow Mansion} (MISC), no centro da capital dinamarquesa, pela produtora executiva do certame, \textbf{Pernille Gaardbo} (PER). "Estou muito feliz em anunciar a nossa decisão. O \textbf{ESC} (MISC) é o maior concurso de música do mundo, o que coloca exigências excepcionais sobre os apresentadores, sendo que estes precisam manter a calma e colocar o profissionalismo no mais alto nível... Estamos confiantes que \textbf{Nikolaj} (PER), \textbf{Pilou} (PER) e \textbf{Lise} (PER) irão preencher totalmente estes requisitos e que os europeus vão sentir a energia positiva que eles transmitem." afirmou a produtora após o anúncio dos nomes.

\textbf{Nikolaj Koppel} (PER) e \textbf{Lise Rønne} (PER), que já são apresentadores da \textbf{DR} (ORG) , trazem uma enorme experiência na apresentação de programas em horário nobre, mas ambos concordam que o \textbf{ESC} (MISC) "joga num campeonato próprio". "Não posso comparar o trabalho de ser anfitrião do \textbf{ESC} (MISC) com mais nenhum trabalho" afirmou \textbf{Nikolaj Koppel} (PER), "É um trabalho de outra dimensão, pois o \textbf{ESC} (MISC) joga num campeonato único. Mas garanto que vou fazer o meu melhor para que todos possam sentir bem vindos na disputa!". \textbf{Lise Rønne} (PER), por outro lado, falou sobre a felicidade e o medo de ter sido escolhida como apresentadora: "Este é o meu trabalho dos sonhos! Eu estou tão animada, e um pouco assustada, com o trabalho que será feito, que mal posso esperar para começar a trabalhar. Estou ansiosa para receber todos na festa!". Por outro lado, \textbf{Pilou Asbæk} (PER) irá desempenhar as funções de apresentador pela primeira vez. Apesar de ser bastante conhecido do público dinamarquês pela sua carreira de ator, \textbf{Pilou} (PER) afirmou que nem pensou duas vezes aquando do convite: "Quando me fizeram a proposta, não tive qualquer sombra de dúvidas sobre aceitar. Estou nervoso, é um fato, mas vou fazer de tudo para me tornar um grande apresentador do \textbf{ESC} (MISC). Estou certo que vai ser uma experiência fantástica e uma edição que ficará na história do concurso."

\textbf{Koppel} (LOC) é um pianista treinado da \textbf{Academia Real Dinamarquesa de Música} (ORG) e já se apresentou em todo o mundo. De 2000, \textbf{Koppel} (PER) trabalhou como jornalista, e 2005-2010, ele foi o \textbf{Diretor Musical} (MISC) para os famosos \textbf{Jardins Tivoli de Copenhaga} (LOC).

Desde 2011, ele tem sido o \textbf{Editor de Comissionamento} (MISC) para a estação de rádio do \textbf{DR P2} (MISC) clássica e editor da \textbf{Rádio P1} (ORG). \textbf{Koppel} (PER) foi o anfitrião do dinamarquês \textbf{TV Aftenshowet} (MISC) ( \textbf{The Evening Mostrar} (ORG) ) e tem uma vasta experiência como uma série de programas de TV.

\textbf{Asbæk} (PER) é um ator dinamarquês, formado pela \textbf{Escola Nacional Dinamarquês de Artes Cênicas} (ORG) (2008). \textbf{Asbæk} (PER) já ganhou inúmeros prêmios em sua curta carreira. Ele ganhou o \textbf{Prêmio Shooting Stars no Festival de Cinema de Berlim} (MISC) (2011), e ambos \textbf{Prêmio dos Críticos de Cinema Dinamarquês} (MISC), \textbf{Bodil} (LOC), e \textbf{Prêmio do Danish Film Academy} (MISC), \textbf{Robert} (PER) - tanto para \textbf{Melhor Artista Masculino} (MISC) de chumbo no filme "R" (2010) .

\textbf{Asbæk} (PER) recebeu elogios da crítica por seu trabalho em filmes como Para \textbf{Verdner} (PER) ( \textbf{Worlds Apart} (PER) , de 2010), \textbf{Kapringen} (LOC) ( um sequestro , 2012), \textbf{Spies \& Glistrup} (ORG) ( \textbf{Sex} (LOC), \textbf{Drugs \& Tributação} (ORG) , 2013) e não menos importante o seu papel de liderança como spin doctor \textbf{Kasper Juul} (PER) na própria série drama político do \textbf{DR Borgen} (MISC) (2010, 2011, 2013). \textbf{Borgen} (PER) cativou os telespectadores de todo o mundo, e \textbf{Asbæk} (PER) ganhou televisão de concessão de \textbf{TV Festival de Copenhaga} (MISC) por seu papel de liderança.

\textbf{Asbæk} (PER) tem sua estréia como um apresentador de TV no \textbf{Festival da Canção 2014 Eurovision} (PER).

\textbf{Lise Rønne} (PER) é bacharel em estudos de mídia e um diploma em jornalismo. \textbf{Rønne} (LOC) começou sua carreira na \textbf{TV MTV} (MISC) em \textbf{Londres} (LOC) como assistente de \textbf{Produtor} (MISC).

Desde 2002, ela tem sido a \textbf{DR} (ORG). \textbf{Rønne} (LOC) tem experiência de vários papéis de hospedagem. Ela já recebeu a seleção nacional dinamarquesa para o \textbf{Festival Eurovisão da Canção} (MISC) por duas vezes (2011 e 2013), e foi por quatro anos série de \textbf{X-Factor} (PER) , na \textbf{Dinamarca} (LOC) (2008, 2009, 2011, 2012). De 2011-2013, ela ocupou o cargo de editor no \textbf{DR Media} (MISC).



\textbf{Participações confirmadas Artista} (MISC) : 4 de Outubro de 2013 Música : 18 de \textbf{Março de 2014} (MISC)\textbf{Artista : 25 de Novembro de 2013 Música} (MISC) : 4 de \textbf{Março de 2014} (MISC)\textbf{Artista : 11 de Janeiro de 2014 Música} (MISC) : 5 de \textbf{Março de 2014} (MISC)\textbf{Artista} (MISC) : 28 de \textbf{Agosto de 2013 Música} (MISC) : 22 de \textbf{Fevereiro de 2014} (LOC)

\textbf{Artista} (MISC) : 19 de \textbf{Novembro de 2013 Música} (MISC) : 9 de \textbf{Março de 2014} (MISC)\textbf{Artista} (MISC) : 19 de \textbf{Junho de 2013 Música} (PER) : 14 de \textbf{Março} (PER) de 2014



Artistas repetentes

\textbf{San Marino} (LOC) - \textbf{Valentina Monetta} (PER) pela 3vez consecutiva vai representar o país em 2012 em \textbf{Baku} (LOC) no \textbf{Azerbaijão} (LOC) não conseguiu passar a final , em 2013 tentou mais uma vez em \textbf{Malmõ} (PER) na \textbf{suécia} (LOC) e não conseguiu passar a final.

\textbf{Roménia} (LOC) - Paula \textbf{Seling \&} (ORG) Ovi representarão mais uma vez o país , em 2010 participaram em \textbf{Oslo na Noruega} (MISC) e conseguiram chegar ao 3Lugar.

\textbf{Rússia} (LOC) - \textbf{As Gêmeas Tolmachevy} (MISC) representaram o país no \textbf{JESC 2006} (MISC) e ganharam.



Sorteio das semifinais

Dia vinte de janeiro decorreu o sorteio de atuação dos participantes , os \textbf{Big 5} (MISC) e o \textbf{Anfitrião} (PER) podem votar nas semifinais. desasseis países vão participar na 1 semifinal e na 2 semifinal conta com um total de 15 países , num total de trinta e um países que vão participar nas duas semifinais.

A distribuição dos países em \textbf{Pote} (LOC) são as seguintes:.



\textbf{Resultados Dos} (MISC) 31 países que vão participar no festival deste ano, 16 países participarão na primeira semi-final a realizar-se no dia 6 de \textbf{Maio} (MISC) (\textbf{Terça Feira} (LOC)), e 15 países na segunda semi-final a realizar-se dia 8 de \textbf{Maio} (LOC) (\textbf{Quinta Feira} (LOC)). Em cada semi-final irão passar os 10 países que tiverem melhor pontuação. Na primeira semi-final votarão os países participantes mais 3 finalistas: \textbf{Espanha} (LOC) , \textbf{França} (LOC) e \textbf{Dinamarca} (LOC) . Na segunda semi-final votarão também os países participantes, tal como na primeira semi-final, assim como a \textbf{Alemanha} (LOC) , \textbf{Itália} (LOC) e \textbf{Reino Unido} (LOC) . Em cada espectáculo (semi-finais e grande final) nenhum país poderá obviamente votar em si próprio. A grande final decorrerá dia 10 de Maio (\textbf{Sábado} (LOC)) contando com 26 países concorrentes, e apenas um país sairá vencedor, sendo este o que tiver melhor pontuação, ganhando o privilégio de realizar a próxima edição (anfitrião) e por conseguinte ganha o passaporte para a final do próximo ano.

O \textbf{Local} (ORG) será na arena situado na capital da \textbf{Dinamarca} (LOC) em \textbf{Copenhaga} (LOC) .

Está semi-final conta com uma convidada especial \textbf{Jessica Mauboy} (PER), representante da \textbf{Austrália} (LOC).

A \textbf{Final} (LOC) conta com duas convidadas ; \textbf{Emmelie de Forest} (PER) ; vencedora do \textbf{ESC} (MISC) 2013 e \textbf{Gaia Cauchi} (PER) ; vencedora do \textbf{JESC 2013} (MISC) .



Outras \textbf{Votações} (LOC)As \textbf{Marcel Bezençon Awards} (PER) foram entregues pela primeira vez durante o \textbf{Festival Eurovisão da Canção 2002} (MISC) , em \textbf{Tallinn} (LOC), \textbf{Estónia} (LOC) , honrando as músicas mais concorrentes na final. Fundada por \textbf{Christer Björkman} (PER) ( \textbf{Suécia} (LOC) representante está no \textbf{Eurovision Song Contest} (ORG) 1992 eo actual \textbf{Chefe de Delegação para a} (MISC) \textbf{Suécia} (LOC)) e \textbf{Richard Herrey} (PER) (membro da \textbf{Herreys eo Festival Eurovisão da Canção 1984} (MISC) vencedor de \textbf{Suécia} (LOC)), os prêmios são nomeados após o criador da competição anual, \textbf{Marcel Bezençon} (PER) . Os prêmios são divididos em três categorias:. \textbf{Prêmio de Melhor Musica} (MISC) , \textbf{Prêmio Artístico} (MISC) e \textbf{Prêmio Compositor} (MISC):

A \textbf{OGAE Internacional ESC Poll} (ORG) revelou a grande opinião de fãs da \textbf{Eurovisão} (LOC) em toda a \textbf{Europa} (LOC) e do \textbf{Mundo sobre o Festival Eurovisão da Canção 2014} (MISC) . Todos os anos, todos os OGAE-clubes organizar uma votação interna dos anos atuais canções \textbf{Eurovision} (MISC). Os resultados serão publicados no site, bem como em nossos sites parceiros. O objetivo deste voto é a voz da opinião fãs ao longo dos anos atuais canções.



\textbf{Cobertura} (MISC) televisiva das semifinais e final

Em baixo encontra-se a lista das cadeias televisivas que emitirão o festival ; ainda não se sabe se estes países tramitaram os três espetáculos ou só dois o só um.

Foi possível assistir ao concurso em qualquer parte do mundo, através da internet. Para além disso, várias televisões de vários países não europeus, transmitiram o festival. Também várias televisões Europeias transmitiram o festival nos seus canais internacionais. Em baixo é possível observar uma lista com os canais que transmitiram o festival, para além fronteiras europeias:



\textbf{Factos} (MISC) e \textbf{Controvérsias} (MISC)As datas previstas para o \textbf{Festival da Canção 2014} (PER) era dia 13 , 15 e 17 de \textbf{Maio} (LOC) foi dado pela apresentadora do Festival Eurovisão da Canção 2013 Petra Mede no fim do espetáculo depois foi confirmado que as novas datas seriam mais cedo ; 6 , 8 e \textbf{10 de Maio} (MISC).

Este ano ouveram 4 saídas e 2 regressos, por outro lado houve um país que desconfirmou e depois confirmou; houve dois paíes que tentaram participar mas, foi recusada:

\textbf{Sobre a Crimeia} (MISC) O maior certame de música europeu será mais um prova de fogo à \textbf{Crimeia} (LOC) e ao conflito \textbf{Rússia} (LOC)/\textbf{Ucrânia} (LOC). A questão que impera, no momento, é porque país votará a província. Qualquer decisão terá, indubitavelmente, um significado político.

A \textbf{UER} (ORG) admite a dificuldade da decisão, embora remeta o assunto para as questões técnicas. \textbf{Jon Ola Sand} (PER) disse que a realidade está para lá do nosso controlo. O jornalista sueco \textbf{Stig Fredriksson} (PER), correspondente da agência \textbf{TT} (ORG) em \textbf{Moscovo} (LOC), explicou que o dilema é muito interessante. Nenhum país reconheceu a anexação. Aceitar que a \textbf{Crimeia} (LOC) vote pela \textbf{Rússia} (LOC) é aceitar a sua incorporação.

O supervisor executivo do \textbf{ESC} (MISC) prefere não misturar política com o evento e diz que é um desafio técnico. É algo que vai ser decidido quando soubermos a infraestrutura do território e quem é o operador que vai trabalhar no local, se é russo ou ucraniano. A \textbf{UER} (ORG) está a trabalhar neste assunto com a televisão pública ucraniana. Uma solução seria não permitir que a \textbf{Crimeia} (LOC) votasse. No entanto, esta solução não agrada à \textbf{UER} (ORG) e a \textbf{Jon Ola Sand} (PER).

\textbf{Loreen} (LOC), vencedora do certame em 2012, defendeu em entrevista a um tablóide sueco que a \textbf{Crimeia} (LOC) é território ucraniano. Junta-se assim a outra vencedora do \textbf{ESC} (MISC), \textbf{Ruslana} (LOC), que também defende que a anexação da região é ilegal. As duas cantoras são conhecidas pelas suas lutas pelos direitos humanos.

\textbf{Conchita Wurst} (LOC) deu à \textbf{Áustria} (LOC) a primeira vitória desde 1966 e ao agradecer, sublinhou que é uma vitória da liberdade e da tolerância. Alguns países contestaram a participação da "mulher barbuda" no \textbf{Festival da Eurovisão} (MISC).

\textbf{Conchita Wurst} (LOC) não passa despercebida. É um travesti com barba, um rosto delicado e maquilhado coberto por uma barba preta e espessa, cabelos longos, vestido de cauda e brincos compridos. \textbf{Conchita} (LOC) é o alter ego do austríaco \textbf{Thomas Neuwirth} (PER), de 25 anos, e venceu o festival com a música "\textbf{Rise Like a Phoenix} (MISC)".



Festival Eurovisão da Canção 2015



Países que classificaram-se para a final \textbf{Países} (MISC) que não se classificaram para a final \textbf{Países} (MISC) que participaram no passado, mas não em 2015

O \textbf{Festival Eurovisão da Canção de 2015} (MISC) (em inglês: \textbf{Eurovision Song Contest} (ORG) 2015 ; em francês: \textbf{Concours Eurovision} (MISC) de la chanson 2015 ; em alemão: \textbf{Liederwettbewerb der Eurovision 2015} (MISC)) foi a 60 edição anual do evento. Depois da marcação de datas provisórias ficou decidido que as semifinais serão realizadas nos dias 19 de 21 de \textbf{Maio} (LOC), e a final no dia 23 de \textbf{Maio} (PER). Esta foi a segunda vez que a \textbf{Áustria} (LOC) sediou o festival, sendo a última em 1967 . Depois de uma reunião entre a emissora da \textbf{Áustria} (LOC) ( \textbf{ORF} (ORG) ) e a \textbf{EBU} (MISC) , decidiu-se que o produtor executivo deste ano seria \textbf{Edgar Böhm} (PER) . A capital do país, \textbf{Viena} (LOC) , foi o local do evento.

Enquanto \textbf{Chipre} (LOC) , \textbf{República Checa} (LOC) e \textbf{Sérvia} (LOC) regressaram a \textbf{Ucrânia} (LOC) decidiu retirar-se devido à crise geopolítica que o país enfrenta. Em 2015 a \textbf{Austrália} (LOC) participou na competição pela primeira vez, tendo sido convidada a participar pela \textbf{EBU} (MISC), pelo interesse que o país tem pelo concurso e porque o transmite há 30 anos. Devido a isto ficou classificada automaticamente para a final.

Após uma das votações mais tensas da história, baseada em 50\% júri e 50\% televoto,a \textbf{Suécia} (LOC) ,ganhou o festival pela sexta vez com \textbf{Måns Zelmerlöw} (LOC) e a música " \textbf{Heroes} (MISC) ". O país foi o primeiro país a ganhar o \textbf{Festival} (PER) duas vezes no século XXI e no atual formato, além disso essa foi a segunda vitória do país em três anos.A \textbf{Itália} (LOC) ganhou no televoto,a \textbf{Rússia} (LOC) ficou em segundo e a \textbf{Suécia} (LOC) em terceiro. A \textbf{Suécia} (LOC) ganhou no júri,a \textbf{Letônia} (LOC) ficou em segundo lugar e a \textbf{Rússia} (LOC) em terceiro. Essa foi a primeira vez no atual formato em que o ganhador não ganhou no televoto. Também pela primeira vez na história,os três primeiros colocados tiveram 280 pontos ou mais.A canção russa " \textbf{A Million Voices} (MISC) " foi a primeira não-vencedora a passar de 300 pontos. \textbf{Áustria} (LOC) e \textbf{Alemanha} (LOC) se tornaram as primeiras canções desde desde 2003 a zeraram no placar. Além disso,a \textbf{Áustria} (LOC) foi o primeiro país anfitrião na história do \textbf{Festival} (LOC) a zerar o placar. Dentre desses fatos inéditos,cabem ainda as melhores participações de \textbf{Montenegro} (LOC) e da \textbf{República Tcheca} (LOC),na história do concurso.



\textbf{Local} (PER)

Para mais detalhes sobre o país anfitrião, consulte \textbf{Áustria} (LOC) .

As três cidades candidatas oficiais e os possíveis locais foram :

No dia 6 de Agosto de 2014, foi revelado que \textbf{Viena} (LOC) foi escolhida para sede da \textbf{Eurovisão 2015} (MISC).



\textbf{Formato} (PER)O formato do concurso foi mantido e espera-se um orçamento de 45 milhões, dos quais 40\% será investido pela \textbf{EBU} (MISC) e 60\% pela \textbf{ORF} (ORG), no entanto, \textbf{Jon Ola Sand} (PER) disse que vai ser algo que vai devidamente homenagear o concurso. Nos finais de 2011 a \textbf{UER} (ORG) anunciou que está em curso uma digitalização e arquivação de todas as edições, desde a primeira em 1956 até à mais recente. Em \textbf{Junho de 2013} (MISC) foi anunciado que o arquivo estava prestes a terminar para celebrar o 60 aniversário. Este estará acessível ao público através do site oficial do concurso.



\textbf{Equipa} (PER)

A emissora pública austríaca \textbf{ORF} (ORG) e a \textbf{União Europeia de Radiodifusão} (ORG) revelaram os nomes da equipa que trabalhará no evento durante os próximos meses. \textbf{Alexander Wrabetz} (PER), o diretor geral da \textbf{ORF} (ORG), disse que tiveram sorte em ter um grande número de personalidades com grande experiência na produção de grandes eventos de TV, que predominantemente também trabalharam o tema do \textbf{Festival} (MISC). Ainda citou que estava convencido que a \textbf{ORF} (ORG), com uma equipa de excelência e em conjunto com os principais parceiros da \textbf{UER} (ORG), bem como outros parceiros públicos e privados da \textbf{Áustria} (LOC), farão com que o \textbf{Festival Eurovisão da Canção 2015} (MISC) seja um evento de alta classe realizado no coração da \textbf{Europa} (LOC).

\textbf{Jon Ola Sand} (PER), mostrou-se satisfeito e constatou que com esta equipa, a \textbf{ORF} (ORG) estabeleceu uma base profissional para a organização do Festival Eurovisão da Canção de 2015.



Regras

Para mais detalhes sobre as regras, consulte: \textbf{Regras do Festival Eurovisão da Canção} (LOC) .

Para mais detalhes sobre a votação, consulte: \textbf{Votação do Festival Eurovisão da Canção} (MISC) .



\textbf{Apresentadores} (PER)Pela primeira vez na história, três mulheres irão apresentar o festival, todas muito conhecidas nos países tectónicos : \textbf{Mirjam Weichselbraun} (ORG), \textbf{Alice Tumbler} (PER) e \textbf{Arabella Kiesbauer} (PER) irão conduzir as galas, enquanto que \textbf{Conchita Wurst} (PER) ,como vencedora do ano anterior irá ancorar a \textbf{green room} (LOC).

\textbf{Mirjam Weichselbraun} (ORG) é uma das apresentadoras mais famosas nos países e língua alemã, tendo apresentado eventos como a \textbf{MTV European Music Awards} (MISC) e o \textbf{Rock Am Ring} (MISC). Atualmente apresenta a versão austríaca do \textbf{Dancing With The Stars} (MISC).

Filha de pai austríaco e mãe francesa vinda da \textbf{Martinica} (LOC), \textbf{Alice Tumler} (PER) fala seis línguas, tem passagem por várias emissoras por toda \textbf{Europa} (LOC). É conhecida por apresentar programas musicais ao vivo, bem como por fazer documentários culturais para a emissora.

\textbf{Arabella Kiesbauer} (PER) começou a trabalhar na \textbf{ORF} (ORG) no final da década de 80. Filha de pai ganês e mãe alemã, nasceu em \textbf{Viena} (LOC) onde foi criada pela sua avó. Em 1994 saiu da \textbf{ORF} (ORG) para apresentar o seu próprio talk-show no canal alemão \textbf{ProSieben} (LOC), que durou 10 anos, voltando para a \textbf{ORF} (ORG) em 2002. Foi nomeada embaixadora da \textbf{Luta contra o Racismo} (MISC) pelo \textbf{Ministério das Relações Internacionais e Integração} (LOC), recebendo a \textbf{Medalha de Ouro} (LOC) da \textbf{República da Áustria} (LOC) pelo seu compromisso de tolerância e contra a desigualdade.



Países classificados para a final

Tal como em anos anteriores, para a final, seis países foram automaticamente confirmados na final, sendo eles:



\textbf{Slogan} (PER)

O \textbf{Slogan} (PER) deste ano foi, "\textbf{Building Bridges} (PER)" (A construir \textbf{Pontes} (LOC)).



\textbf{Logótipo} (PER)Devido aos 60 anos do festival celebrados em 2015, a \textbf{União Europeia de Radiodifusão} (ORG) decidiu atualizar o logótipo genérico para assinalar a data. Segundo \textbf{Jon Ola Sand} (PER), supervisor executivo do \textbf{ESC} (MISC): "Foi realizada uma extensa pesquisa entre as emissoras participantes, milhares de espectadores regulares, fãs e outras partes interessadas. Um dos resultados claros, era que o \textbf{logótipo} (LOC) é cada vez mais reconhecido, mas achamos que merecia uma reformulação depois de ter sido usado em 10 edições".



\textbf{Artistas Repetentes Arménia} (PER) - \textbf{Inga Arshakyan} (PER) na edição de 2009, em \textbf{Moscovo} (LOC) , \textbf{Rússia} (LOC), representou a \textbf{Arménia} (LOC) com a sua irmã \textbf{Anush} (PER) , ficando em 10 lugar, este ano regressa como membro da banda, \textbf{Genealogy} (PER).

\textbf{Azerbaijão} (LOC) - \textbf{Elnur Hüseynov} (PER) representou o país na edição de 2008 (a primeira participação do \textbf{Azerbaijão} (LOC)), em \textbf{Belgrado} (LOC) , em 2008, na \textbf{Sérvia} (LOC) , em conjunto com \textbf{Samir Javadzadeh} (PER) , este ano regressa como artista solo.

\textbf{Malta} (LOC) - \textbf{Amber} (LOC) será a representante deste ano de \textbf{Malta} (LOC). Na edição de 2012 foi corista de \textbf{Kurt Calleja} (PER).

\textbf{San Marino} (LOC) - \textbf{Michelle Perniola} (LOC) e \textbf{Anita Simoncini} (PER) serão os representantes de \textbf{São Marino} (LOC). Foram os representantes do país no \textbf{JESC2013} (ORG) e \textbf{JESC2014} (LOC) respectivamente, apesar de \textbf{Anita Simoncini} (PER) ser membro do grupo \textbf{The Peppermints} (ORG).

\textbf{Bielorrússia} (LOC) - \textbf{Uzari} (LOC) e \textbf{Maiumuna} (LOC) serão os representantes da \textbf{Bielorrússia} (LOC). \textbf{Uzari} (LOC) foi corista de \textbf{Anastacia Vinnikova} (PER) na edição de 2011.



Sorteio das semifinais

\textbf{Sorteio das semifinais} (MISC) ocorreu em 26 de \textbf{Janeiro} (LOC) de 2015.



\textbf{Resultados Dos} (MISC) 33 países que participaram no festival nas semifinais, deste ano, 16 países participaram na primeira semifinal realizada no dia 19 de \textbf{Maio} (LOC) (\textbf{Terça Feira} (LOC)), e 17 países na segunda semifinal a realizada no dia 21 de \textbf{Maio} (LOC) (\textbf{Quinta Feira} (LOC)). Em cada semifinal passaram os 10 países que tiveram melhor pontuação. Na primeira semifinal votaram os países participantes mais 3 finalistas: \textbf{Espanha} (LOC) , \textbf{França} (LOC) e \textbf{Áustria} (LOC) . Na segunda semifinal votaram também os países participantes, tal como na primeira semifinal, assim como \textbf{Alemanha} (LOC) , \textbf{Itália} (LOC) e \textbf{Reino Unido} (LOC) . Em cada show (semifinais e grande final) nenhum país podia obviamente votar em si próprio. A \textbf{Austrália} (LOC) como país convidado voltou nas duas semifinais. A grande final decorreu dia 23 de \textbf{Maio} (MISC) (\textbf{Sábado} (LOC)) contando com 27 países concorrentes, e apenas um país saiu vencedor, sendo aquele que acumulou mais pontos, ganhando o privilégio de realizar a próxima edição e estar automaticamente classificado para a final do ano que vem.

A ordem de atuação foi revelada no dia 23 de \textbf{Março} (PER).

A ordem de atuação foi revelada no dia 23 de \textbf{Março} (PER).

A ordem de atuação foi revelada no dia 23 de \textbf{Março} (PER).

A ordem de atuação será revelada na madrugada do dia 22 de \textbf{Maio de 2015} (MISC).



Outras \textbf{Votações} (LOC)A OGAE Internacional ESC Poll esteve a revelar a grande opinião de fãs da \textbf{Eurovisão} (LOC) em toda a \textbf{Europa} (LOC) e do \textbf{Mundo sobre o Festival Eurovisão da Canção 2015} (MISC). Todos os anos, todos os OGAE-clubes organiza uma votação interna dos anos atuais canções \textbf{Eurovision} (MISC). Os resultados serão publicados no site, bem como nossos seus sites parceiros. O objetivo deste voto é a voz da opinião fãs ao longo dos anos atuais.



Festival Eurovisão da Canção 2016



Países que classificaram-se para a final \textbf{Países} (MISC) que não se classificaram para a final \textbf{Países} (MISC) que participaram no passado, mas não em 2016

O \textbf{Festival Eurovisão da Canção de 2016} (MISC) (em inglês : \textbf{Eurovision Song Contest} (ORG) 2016 e em sueco : \textbf{Eurovision Song Contest} (ORG) 2016 ) foi a 61. edição anual do evento. Realizou-se, pela 6. vez, na \textbf{Suécia} (LOC) , depois de \textbf{Mans Zelmerlow} (PER) ter ganho a edição anterior com a canção \textbf{Heroes} (MISC) . O evento ocorreu pela 3. vez em \textbf{Estocolmo} (LOC) e pela 2. vez no \textbf{Ericsson Globe} (MISC) , que também que já havia sido sede do festival em 2000 . A última vez até à data em que o festival se tinha realizado em solo sueco deu-se na cidade de \textbf{Malmö} (LOC) , em 2013 .

Este festival foi notável, pois foi a primeira vez desde 1975 em que o procedimento da votação mudou consideravelmente: cada país deu 1-8, 10 e 12 por duas vezes, uma vez pelo júri específico, outra vez pelo tele voto. Na final, os resultados do júri foram anunciados da forma tradicional, enquanto que o resultado do tele voto foi anunciado de uma só vez. Cada resultado teve um peso de 50\% no resultado final. O slogan foi \textbf{Come} (PER) together .

A cantora \textbf{Jamala} (PER) , representando a \textbf{Ucrânia} (LOC) , venceu o \textbf{Festival Eurovisão da Canção} (MISC) com o tema " 1944 ", inspirada na história real da sua família, que foi vítima da deportação a força dos \textbf{Tártaros da Crimeia} (MISC) por \textbf{Stalin} (PER) , durante a \textbf{Segunda Guerra Mundial} (MISC) , como vingança pelo apoio da etnia aos nazistas . Nascida no \textbf{Quirguistão} (LOC) , \textbf{Jamala} (LOC) é etnicamente tártara da \textbf{Crimeia} (LOC).



\textbf{Formato} (PER)As datas preliminares foram anunciadas na reunião da delegação, realizada a 16 de março de 2015, em \textbf{Viena} (LOC). As meias-finais foram, a 10 e 12 de maio de 2016, e a final no dia 14 de maio de 2016.



\textbf{Apresentadores} (PER)Na noite da vitória de 2015, \textbf{Måns Zelmerlöw} (LOC) , em \textbf{Viena} (LOC) , foi questionado se queria ser o apresentador da edição de 2016, e ele aceitou. \textbf{Zelmerlow} (PER) apresentou o \textbf{Melodifestivalen} (PER) (seleção sueca para o \textbf{Festival} (LOC)) em 2010 e o popular espetáculo sueco \textbf{Allsång på Skansen} (PER) .

Em 25 de maio de 2015, \textbf{Christer Björkman} (PER) contou ao jornal \textbf{Expressen} (ORG) que \textbf{Sanna Nielsen} (PER) , \textbf{Gina Dirawi} (PER) e \textbf{Petra Mede} (PER) eram potenciais apresentadoras do festival. Já a 1. de junho de 2015, o mesmo jornal comentou que a \textbf{SVT} (ORG) estava considerando o ator \textbf{Dolph Lundgren} (PER) , juntamente com \textbf{Måns} (LOC), como o duo de apresentadores.

Em 14 de dezembro de 2015 foi anunciado o retorno de \textbf{Petra Mede} (PER) (apresentadora solo em 2013 ), juntamente com \textbf{Måns Zelmerlöw} (LOC) , à apresentação do evento.



Regras

A \textbf{UER} (ORG) anunciou que o sistema de votação a passará pela maior mudança desde 1975 em que se estabeleceu o formato clássico do festival. Desde (2011-2015), cada país entregava de 1 a 8, 10 e 12 pontos a suas 10 canções favoritas, a partir da combinação de seu júri (50\%) e do seu televoto (50\%). A partir de 2016, cada jurado profissional de cada país e cada televoto de cada país otorgará por separado 1 a 8, 10 e 12 pontos as suas 10 canções favoritas.

O formato de anúncio das semifinais é mantido, aonde os 10 agora finalistas são anunciados de forma aleatória. Na final, haverá duas fases de anúncio dos pontos. Na primeira se mantém a dinâmica tradicional aonde os resultados dos júris são anunciados (12 pontos para o mais votado, 10 para o segundo e assim sucessivamente). Na segunda parte, os apresentadores revelarão os resultados do televoto, mas a partir de agora revelados globalmente (que representa o total dos 12, 10 e de 8 a 1 ponto de cada país) em ordem decrescente (do posto 26 ao 1). O vencedor da \textbf{Eurovisão} (LOC) não será conhecido até à última e derradeira votação. Tanto o público como os jurados continuarão com o peso de 50\%.

Com o novo sistema o ganhador não foi conhecido até último minuto,o que deu mais emoção. "Este é um grande passo para fazer um melhor 'show' televisivo e para garantir mais emoção", assegurou \textbf{Jon Ola Sand} (PER) , supervisor executivo do certame.



\textbf{Sedes} (PER)A SVT anunciou a 24 de Maio que a sua primeira escolha para a \textbf{Eurovisão} (LOC) tinha sido a \textbf{Tele2 Arena} (ORG) em \textbf{Estocolmo} (LOC). Entretanto, outras cidades na \textbf{Suécia} (LOC) poderiam apresentar uma candidatura. Depois da vitória da \textbf{Suécia} (LOC) em \textbf{Viena} (LOC), as cidades teriam 3 semanas para apresentar a sua candidatura e que o anúncio seria feito a meio do \textbf{Verão} (MISC). O \textbf{Ericsson Globe} (MISC) em \textbf{Estocolmo} (LOC) foi anunciado a 8 de Julho de 2015. Esta será a segunda vez que a \textbf{arena} (PER) sediará o evento. A primeira foi em 2000 .



\textbf{Participantes} (MISC)Os países que competiram em \textbf{Maio} (LOC) na \textbf{Suécia} (LOC), foram revelados na \textbf{primeira semana de Janeiro} (MISC), os seguintes países revelaram publicamente a sua participação na eurovisão: \textbf{Albânia} (LOC) , \textbf{Alemanha} (LOC) , \textbf{Áustria} (LOC) , \textbf{Azerbaijão} (LOC) , \textbf{Bélgica} (LOC) , \textbf{Bielorrússia} (LOC) , \textbf{Bósnia e Herzegovina} (LOC) , \textbf{Bulgária} (LOC) , \textbf{Chipre} (LOC) , \textbf{Dinamarca} (LOC) , \textbf{Eslovénia} (LOC) , \textbf{Espanha} (LOC) , \textbf{Estónia} (LOC) , \textbf{Finlândia} (LOC) , \textbf{França} (LOC) , \textbf{Geórgia} (LOC) , \textbf{Grécia} (LOC) , \textbf{Holanda} (LOC) , \textbf{Hungria} (LOC) , \textbf{Irlanda} (LOC) , \textbf{Israel} (LOC) , \textbf{Itália} (LOC) , \textbf{Letónia} (LOC) , \textbf{Lituânia} (LOC) , \textbf{Macedónia} (LOC) , \textbf{Malta} (LOC) , \textbf{Noruega} (LOC) , \textbf{Polónia} (LOC) , \textbf{Reino Unido} (LOC) , \textbf{República Checa} (LOC) , \textbf{Roménia} (LOC) , \textbf{Suécia} (LOC) (anfitrião), \textbf{Suíça} (LOC) e além disso, a \textbf{Turquia} (LOC) e a \textbf{Ucrânia} (LOC) confirmaram o seu regresso, a \textbf{Turquia} (LOC) deixou de participar a partir do ano de 2013 e a \textbf{Ucrânia} (LOC) só não participou em 2015. Nos regressos, a \textbf{Croácia} (LOC) é o nome mais falado em regressar, também a \textbf{Bulgária} (LOC) e a \textbf{Eslováquia} (LOC) poderão estar de volta. As emissoras dos países de \textbf{Andorra} (LOC) , \textbf{Luxemburgo} (LOC) , \textbf{Portugal} (LOC) e \textbf{Mónaco} (LOC) já revelaram que não irão participar em 2016.

Originalmente, o grupo italiano \textbf{Stadio} (PER) representaria o país por ter vencido a edição de 2016 do \textbf{Festival de San Remo} (PER) , porém renunciaram a participação. A emissora italiana \textbf{RAI} (ORG) , decidiu por mérito dar a chance à cantora que ficou na segunda posição no festival, a cantora \textbf{Francesca Michielin} (PER) .



\textbf{Resultados Dos} (MISC) 36 países que participarão no festival nas semifinais, deste ano, 18 países participarão na primeira semifinal realizada no dia \textbf{10 de Maio} (MISC) (\textbf{Terça Feira} (LOC)), e 18 países na segunda semifinal a realizada no dia 12 de \textbf{Maio} (LOC) (\textbf{Quinta Feira} (LOC)). Em cada semifinal passarão 10 países que tiveram melhor pontuação. Na primeira semifinal votarão os países participantes mais 3 finalistas: \textbf{Espanha} (LOC) , \textbf{França} (LOC) e \textbf{Suécia} (LOC) . Na segunda semifinal votarão também os países participantes, tal como na primeira semifinal, assim como \textbf{Alemanha} (LOC) , \textbf{Itália} (LOC) e \textbf{Reino Unido} (LOC) . Em cada show (semifinais e grande final) nenhum país poderá obviamente votar em si próprio. A grande final decorrerá dia 14 de Maio (\textbf{Sábado} (LOC)) contando com 26 países concorrentes, e apenas um país sairá vencedor, sendo aquele que acumulará mais pontos, ganhando o privilégio de realizar a próxima edição e estar automaticamente classificado para a final do ano que vem.

A ordem de atuação foi revelada no dia 23 de \textbf{Março} (PER).

A ordem de atuação foi revelada no dia 23 de \textbf{Março} (PER).

A \textbf{Roménia} (LOC) originalmente seria o décimo segundo país a competir nesta semifinal, mas, devido à dívida que a emissora pública tinha com a \textbf{EBU} (MISC) , foi forçada a desistir.



Outras \textbf{Votações} (LOC)A OGAE Internacional ESC Poll esteve a revelar a grande opinião de fãs da \textbf{Eurovisão} (LOC) em toda a \textbf{Europa} (LOC) e do \textbf{Mundo sobre o Festival Eurovisão da Canção 2015} (MISC) . Todos os anos, todos os OGAE-clubes organiza uma votação interna dos anos atuais canções \textbf{Eurovision} (MISC). Os resultados serão publicados no site, bem como nossos seus sites parceiros. O objetivo deste voto é a voz da opinião fãs ao longo dos anos atuais.

*\textbf{Resultados} (MISC) refletidos da votação de 42 membros \textbf{OGAE} (MISC).



\textbf{Comentadores e Porta-Vozes} (MISC)

SVT anunciou a 22 de abril de 2016 que pela primeira vez foram transmitidos em \textbf{Língua Gestual Internacional} (MISC) para deficientes auditivos. Todos os três espetáculos foram traduzidos por \textbf{Julia Kankkonen} (PER). As atuações serão interpretadas por dez intérpretes de linguagem gestual e os diálogos dos anfitriões por três intérpretes de linguagem gestual:

As transmissões internacionais em linguagem gestual serão transmitidos online ao lado dos três espetáculos ao vivo, com os seguintes países também a transmitirem na televisão:

Os porta-vozes anunciarão a pontuação de 12 pontos do júri nacional do respectivo país na seguinte ordem:



Festival Eurovisão da Canção 2017



Países que classificaram-se para a final \textbf{Países} (MISC) que não se classificaram para a final \textbf{Países} (MISC) que participaram no passado, mas não em 2017

O \textbf{Festival Eurovisão da Canção de 2017} (MISC) (em inglês : \textbf{Eurovision Song Contest} (ORG) 2017 ; em francês : \textbf{Concours Eurovision} (MISC) de la chanson 2017 ; em ucraniano : \textbf{} (PER)  \textbf{} (PER) 2017 ) foi a 62. edição anual do evento. O festival foi realizado pela segunda vez em \textbf{Kiev} (LOC) , \textbf{Ucrânia} (LOC) , depois de \textbf{Jamala} (LOC) ter ganhado a edição anterior com a canção 1944 . As semifinais realizaram-se nos dias 9 e 11 de maio, enquanto que a final se deu no dia 13 de maio. A última vez que o \textbf{Festival} (LOC) adulto foi realizado em solo ucraniano foi em 2005 , também em \textbf{Kiev} (LOC) .

\textbf{Portugal} (LOC) e \textbf{Roménia} (LOC) regressaram este ano depois de estarem ausentes na edição de 2016 . A \textbf{Bósnia e Herzegovina} (LOC) retirou-se por problemas financeiros, enquanto que a \textbf{Rússia} (LOC) anunciou a sua retirada a 13 de abril de 2017. Esta retirada foi anunciada depois que a sua representante, \textbf{Yulia Samoylova} (PER) , foi banida da \textbf{Ucrânia} (LOC) , após alegações de que ela tenha entrado ilegalmente na \textbf{Crimeia} (LOC) em 2015.

O cantor \textbf{Salvador Sobral} (PER) , representando \textbf{Portugal} (LOC) , venceu esta edição com o tema " Amar pelos \textbf{Dois} (LOC) ", com 758 pontos, a maior pontuação da história e a primeira vitória de \textbf{Portugal} (LOC) no Festival Eurovisão da Canção que, ao longo de 53 anos, nunca tinha conquistado nenhum top 5. Esta é também a primeira vitória para um país a cantar na sua língua oficial desde 2007 , ano em que a \textbf{Sérvia} (LOC) ganhou o certame com " \textbf{Molitva} (MISC) ". O top 3 - \textbf{Portugal} (LOC), \textbf{Bulgária} (LOC) e \textbf{Moldávia} (LOC) -, conquistaram todos as suas melhores posições da história.



Regras

Para mais detalhes sobre as regras, consulte: \textbf{Regras do Festival Eurovisão da Canção} (LOC) .

Para mais detalhes sobre a votação, consulte: \textbf{Votação do Festival Eurovisão da Canção} (MISC) .



Preparação do Festival

Nove cidades mostraram interesse em acolher o \textbf{Festival: Dnipropetrovsk} (MISC) , \textbf{Kherson} (LOC) , \textbf{Carcóvia} (LOC) , \textbf{Kiev} (LOC) , \textbf{Lviv} (LOC) , \textbf{Odessa} (LOC) , \textbf{Tcherkássi} (LOC) , \textbf{Vinnystia} (LOC) e \textbf{Irpin} (LOC) . Porém apenas seis seguiram as condições propostas pela \textbf{NTU} (ORG).

A 22 de julho, três cidades foram anunciadas como finalistas para sediar o \textbf{Festival} (LOC), depois da emissora ucraniana \textbf{NTU} (ORG) ter eliminado, por diversos fatores, as cidades de \textbf{Kherson} (LOC) , \textbf{Carcóvia} (LOC) e \textbf{Lviv} (LOC) .

A data prevista para o anúncio da cidade-sede e do local foi 1 de agosto de 2016. No entanto, a \textbf{NTU} (ORG) confirmou em 26 de julho de 2016 que a cidade seria anunciada no dia seguinte, o que acabou não acontecendo. Com isso, o anúncio foi realocado novamente para 1 de agosto. Neste dia, a \textbf{NTU} (ORG) referiu que não tinha ainda decidido e que estava considerando as três propostas, entendendo que faltavam algumas condições em cada uma. Em 10 de agosto, \textbf{Oleksandr Kharebin} (PER), \textbf{Diretor-Geral Adjunto} (PER) da \textbf{NTU} (ORG) , anunciou que a cidade sede seria anunciada até dia 24 de agosto.

Após estes sucessivos adiamentos sobre a decisão da cidade sede, o prefeito de Dnipro lançou um apelo público, em 19 de agosto, pela sua conta oficial do \textbf{Facebook} (MISC) , ao time responsável pelo processo, para que o fizessem o mais rapidamente possível e referiu que, se não se soubesse a cidade sede dentro de uma semana, \textbf{Dnipro} (PER) estaria desistindo oficialmente do evento.

Minutos depois, \textbf{Oleksandr Kharebin} (PER) respondeu, pedindo que o político mantivesse a calma, citando que a cidade sede da \textbf{Eurovisão} (LOC) em 2017 seria anunciada em 23 de agosto.

Durante os dois dias seguintes, não se houve conhecimento de nenhuma nova informação relativa ao processo. Até que, em 25 de agosto, foi-se anunciada a realização de uma entrevista coletiva para o anúncio da cidade-sede. Alguns minutos antes desta coletiva começar, contudo, ela foi cancelada, com o argumento de que ainda faltavam alguns detalhes finais para o anúncio.

O então \textbf{Diretor Geral} (PER) da \textbf{NTU} (ORG), \textbf{Zurab Alasania} (PER), mais tarde anunciou que \textbf{Dnipropetrovsk} (LOC) estava oficialmente fora da corrida, sendo que as finalistas eram \textbf{Odessa} (LOC) e \textbf{Kiev} (LOC) .

A 9 de setembro, foi anunciado que \textbf{Kiev} (LOC) e o \textbf{Centro Internacional de Exibições} (ORG) (\textbf{IEC} (ORG)) iriam sediar o \textbf{Festival Eurovisão da Canção 2017} (MISC).

Em 20 de julho, as seis cidades, agora candidatas, se apresentaram oficialmente ao vivo em um programa de televisão chamado \textbf{Batalha das Cidades} (MISC) (em inglês : \textbf{City Battle} (ORG) ), no qual se incluía a discussão entre os membros do \textbf{Comitê Organizador} (ORG), representantes da \textbf{União Europeia de Radiodifusão} (ORG) , músicos profissionais da \textbf{Ucrânia} (LOC) e fãs de toda a parte. O programa foi apresentado por \textbf{Timur Miroshnychenko} (PER) e transmitido no horário nobre do canal 1 da \textbf{NTU} (ORG) , na \textbf{Radio Ukraine} (MISC) e pelo \textbf{YouTube} (MISC) , através do canal oficial do festival. Cada proposta foi apresentada ao longo do programa pelos representantes de cada cidade:

\textbf{Carcóvia} (LOC) : \textbf{Igor Terekhov} (PER) | Vice-Presidente da Câmara

Cherson : \textbf{Volodymyr Mykolaienko} (ORG) | Presidente da Câmara

Kiev : \textbf{Oleksii Reznikov | Vice-chefe} (LOC) do departamento administrativo da cidade

\textbf{Lviv} (LOC) : \textbf{Andrii Moskalenko | Vice-Presidente da Câmara} (PER)Odessa : \textbf{Pavel Vugelman |} (LOC) Vice-Presidente da Câmara

 Cidade anfitriã  Finalistas  Cidades que se retiraram

Precisarava de algumas obras, incluindo uma cobertura.

Em certas cidades, faltavam infraestruturas necessárias para acolher o \textbf{Festival} (LOC) : os locais propostos eram a céu aberto, o que é incapacitante para acolhê-lo. Tcherkássi , \textbf{Vinnystsia} (LOC) e \textbf{Irpin} (LOC) foram as cidades que nunca apresentaram uma candidatura oficial.

A 8 de junho, foi realizada a primeira reunião da organização do \textbf{Festival} (LOC), de acordo com a chefe da delegação ucraniana e posteriormente produtora executiva do evento, \textbf{Viktoria Romanova} (PER). Durante esse tempo, a \textbf{UER} (ORG) e a \textbf{NTU} (ORG) estipulou as regras e as condições que as cidades candidatas teriam de respeitar para sediar o evento. Já o primeiro-ministro \textbf{Volodymyr Groysman} (PER) , anunciou, a 18 de maio de 2016, que provavelmente um processo de escolha da sede seria aberto, pois o evento tornara-se uma prioridade para o país. As características impostas pela \textbf{UER} (ORG) foram:

Em 24 de junho, a \textbf{UER} (ORG) revelou o calendário de seleção da cidade anfitriã. A seleção da cidade foi entre os dias 24 de junho e 1 de agosto e consistiu em 4 etapas :

As datas provisórias iniciais para a 62. edição foram anunciadas a 14 de março de 2016 durante a tradicional reunião dos \textbf{Chefes de Delegação} (MISC), na sede do festival do ano anterior, com as semifinais originalmente marcadas para os dias 16 e 18 de maio e a final para 20 de maio de 2017. Essas datas provisórias foram escolhidas pela \textbf{UER} (ORG) com o objetivo de evitar potenciais choques de datas com outros eventos, tal como a final da \textbf{Liga dos Campeões da UEFA} (MISC) . No entanto, logo em junho, a \textbf{NTU} (ORG) propôs novas datas: antecipar o festival em uma semana, para os dias 9, 11 e 13 de maio de 2017. Estas mesmas datas estiveram sujeitas a alterações, mas foram confirmadas no mesmo dia em que foi anunciado onde se iria realizar o festival. Outra alternativa foram os dias 23, 25 e 27 de maio de 2017.

O orçamento previsto para o \textbf{Festival de 2017} (MISC) foi de  663 milhões (aproximadamente 23,2 milhões de  ou 24 milhões de dólares). Desta soma,  457 milhões (aproximadamente 16 milhões  ou \textbf{US\$} (MISC) 17 milhões) provieram do governo ucraniano. A cidade de \textbf{Kiev} (LOC) alocou  208 milhões (aproximadamente 7,2 milhões  ou \textbf{US\$} (MISC) 7,6 milhões) à organização do \textbf{Festival} (LOC) e previu repartir as suas despesas entre 2016 e 2017, utilizando  51 milhões (aproximadamente 1.8 milhão  ou \textbf{US\$} (MISC) 1.9 milhão) no período pré-evento, e os restantes  154 milhões (5,4 milhões  ou \textbf{US\$} (MISC) 5.7 milhões) para serem usados durante os meses anteriores ao festival.

A 5 de janeiro de 2017, a emissora \textbf{NTU} (ORG) anunciou os membros da equipa principal que estaria encarregada do \textbf{Festival de 2017} (MISC). A emissora nomeou \textbf{Aleksandr Kharebin} (PER) e \textbf{Viktorya Romanova} (PER) como produtores executivos. Porém, uma parte da equipe renunciou aos cargos em 13 de fevereiro de 2017, forçando a \textbf{NTU} (ORG) a nomear substitutos.

A emissora ucraniana encontrou muitas dificuldades quanto à preparação do \textbf{Festival} (LOC). A 27 de novembro de 2016 , \textbf{Aleksandr Kharebin} (PER), diretor-geral da emissora, anunciou que os preparativos tinham tido meses de atraso. Graças à, de acordo com ele, "burocracia, uma legislação inadequada e uma atitude insolente de certos oficiais" . Mais concretamente, a alocação do orçamento pelo governo ucraniano foi atrasado e cada etapa dos preparativos deveria, como diz a lei ucraniana, passar por uma chamada de ofertas públicas -- a escolha do tema e do slogan do \textbf{Festival} (LOC), por exemplo, foi feita por ofertas públicas.

Estes problemas acumulados poderiam tirar os direitos de organização do festival da \textbf{NTU} (ORG). Porém, a 28 de novembro de 2016, \textbf{Aleksandr Kharebin} (PER) anunciou que a \textbf{Ucrânia} (LOC) iria, mesmo assim, organizar o \textbf{Festival} (LOC), e, apesar dos rumores publicados na imprensa ucraniana, progressos significativos teriam sido feitos a propósito da organização.

A preparação do \textbf{Festival} (LOC) foi igualmente marcada por diversas demissões. Inicialmente, \textbf{Zurab Alasania} (PER), diretor-geral da emissora, apresentou a sua demissão a 1 de novembro de 2016, em sinal de protesto contra o baixo orçamento alocado pelo governo ucraniano à emissora e a pressão para não revelar os problemas encontrados pela organização. Seguinte a esta demissão, uma pessoa -- não nomeada na carta de demissão --, designou um novo diretor para o \textbf{Festival} (LOC), que assumiu o controlo total da sua preparação, anulando as ações da equipa principal encarregada da organização. Este bloqueio levou à demissão de 21 membros da equipa a 13 de fevereiro de 2017 , incluindo os produtores executivos \textbf{Viktorya Romanova} (PER) e \textbf{Aleksandr Kharebin} (PER), a diretora comercial \textbf{Iryna Asman} (PER), o gerente geral do evento \textbf{Denys Bloshchenskyi} (PER) e o diretor da segurança \textbf{Aleksii Karaban} (PER). A \textbf{UER} (ORG) , depois disso, insistiu com a emissora sobre a importância de envolver rapidamente e eficazmente os elementos já aprovados para o \textbf{Festival} (LOC) e respeitar o planeamento estabelecido, apesar das mudanças em cima da hora.

O sorteio para determinar a atribuição dos países participantes às respectivas semifinais teve lugar no \textbf{Column Hall} (LOC) a 31 de janeiro de 2017 , apresentado por \textbf{Timur Miroshnychenko} (PER) e \textbf{Nika Konstantinova} (PER). Os 37 semifinalistas foram alocados em seis potes, com base em padrões históricos de votação, conforme calculado pelo parceiro oficial do televoto, \textbf{Digame} (PER) .

O logótipo e o slogan foram revelados a 30 de janeiro de 2017 . O slogan para o \textbf{Festival de 2017} (MISC) é \textbf{Celebrate Diversity} (PER) (em português , \textbf{Celebra} (LOC) a \textbf{Diversidade} (LOC)). O slogan é a mensagem central para o evento ano e é complementada por um logótipo criativo baseado no tradicional colar de grânulos ucraniano conhecido como \textbf{Namysto} (ORG) . Mais do que uma peça de joelheira, \textbf{Namysto} (LOC) é um amuleto protetor e um símbolo de beleza e saúde. É composto de diversas pérolas diferentes, cada uma com o seu próprio design, e celebra a diversidade e individualidade.



Promoção

Foi anunciado que a 30 de abril que as equipas criativas da \textbf{Eurovisão} (LOC) e do \textbf{Twitter} (MISC) estariam a trabalhar juntas para criar 3 emojis que acompanharam os hashtags promocionais durante a edição. O emoji de coração aparecereu juntamente com \#ESC2017 e \#\textbf{Eurovision} (MISC) , enquanto que o emoji do troféu do vencedor foi usado para \#12Points e \#douzepoints . O último emoji é o logótipo do certame, que apareceu junto ao \textbf{hashtag} (MISC) \#\textbf{CelebrateDiversity} (MISC) , o slogan desta edição.

A 5 de maio, o \textbf{Banco Central da Ucrânia} (LOC) anunciou o lançamento de uma moeda especial, em prata, para marcar o facto da \textbf{Ucrânia} (LOC) receber e organizar a 62. edição do \textbf{Festival da Eurovisão} (MISC). A moeda de 5 grívnia mostra o céu de \textbf{Kiev} (LOC) de um lado e o logótipo da \textbf{Eurovisão 2017} (MISC) do outro. O logo contém a bandeira ucraniana. Foram feitas 40 mil moedas, com metade delas a serem vendidas no website do \textbf{Banco Central} (LOC), onde puderam ser compradas até 15 de maio.



\textbf{Participantes Dezoito} (PER) países participaram na 1. semifinal. A \textbf{Itália} (LOC) , a \textbf{Espanha} (LOC) e a \textbf{Reino Unido} (LOC) votaram nesta semifinal.

\textbf{Dezanove} (PER) países iriam participam na 2. semifinal, até que a cantora da \textbf{Rússia} (LOC), \textbf{Julia Samoylova} (PER), foi expulsa do concurso. A \textbf{França} (LOC) , a \textbf{Alemanha} (LOC) e a \textbf{Ucrânia} (LOC) votaram nesta semifinal.

Os países a seguir qualificam-se diretamente para a final do \textbf{Festival:} (MISC)\textbf{Resultados Dos} (MISC) 36 países que participaram no festival nas semifinais deste ano, 18 países participaram na primeira semifinal, realizada no dia 9 de maio (na terça-feira), enquanto que os outros 18 participaram na segunda semifinal, realizada no dia 11 de maio (na quinta-feira). Em cada semifinal, classificaram-se para a grande final os 10 países que tiveram melhor pontuação. Na primeira, votaram os países participantes e mais 3 finalistas: \textbf{Itália} (LOC) , \textbf{Espanha} (LOC) e \textbf{Reino Unido} (LOC) . Na segunda, também votaram os países participantes, bem outros 3 finalistas: \textbf{França} (LOC) , \textbf{Alemanha} (LOC) e \textbf{Ucrânia} (LOC) . Em cada show (semifinais e grande final), nenhum país pôde votar em si próprio. A grande final ocorreu no dia 13 de maio (sábado), contando com 26 países concorrentes (10 de cada semifinal, os países do \textbf{Big 5} (MISC) (\textbf{Alemanha} (LOC), \textbf{Espanha} (LOC), \textbf{França} (LOC), \textbf{Itália} (LOC) e \textbf{Reino Unido} (LOC)) e o país anfitrião, \textbf{Ucrânia} (LOC)).

A ordem de atuação foi revelada no dia 31 de \textbf{Março} (PER).



Outros países

Para um país para ser elegível para participar no \textbf{Festival Eurovisão da Canção} (MISC) , ele precisa ser um membro ativo da \textbf{União Europeia de Radiodifusão} (ORG) (\textbf{UER} (ORG)), ou, para membros associados, ter sua participação aprovada pela instituição. A \textbf{UER} (ORG) confirmou, inicialmente, que quarenta e três países fariam parte do festival. Contudo, com a desistência da \textbf{Rússia} (LOC) , o número foi reduzido para quarenta e dois.



Álbum oficial

\textbf{Eurovision Song Contest} (ORG): \textbf{Kyiv 2017} (MISC) é o álbum de compilação oficial do festival, lançado pela \textbf{União Europeia de Radiodifusão} (ORG) e lançado pela \textbf{Universal Music Group} (ORG) a \textbf{21 de Abril} (ORG) e fisicamente a 28 de Abril de 2017. O álbum conta com todas as 42 participações, incluindo os semifinalistas que não se classificaram para a final, e a canção russa, que se retirou do festival a 13 de abril de 2017.



Festival Eurovisão da Canção 2018



Semi-final 2: "\textbf{ESClopedia} (MISC)" (parte 2); Danças a representar vencedores anteriores do concurso pelas apresentadoras ( \textbf{Catarina Furtado} (PER) , \textbf{Daniela Ruah} (PER) , \textbf{Sílvia Alberto} (PER) , \textbf{Filomena Cautela} (LOC) ); "\textbf{Planet Portugal} (MISC)" (parte 2). Final: "Mano a mano" e "\textbf{Amar} (LOC) pelos dois" - interpretação de \textbf{Salvador Sobral} (PER) (com \textbf{Caetano Veloso} (PER) e \textbf{Júlio Resende} (PER) ); \textbf{Mayra Andrade} (PER) , \textbf{Plutónio} (PER) , \textbf{Dino D'Santiago} (PER) e \textbf{Sara Tavares} (PER) interpretam música contemporânea portuguesa no "\textbf{Interval Act} (MISC)" (com \textbf{Branko} (PER) ).

Países que classificaram-se para a final \textbf{Países} (MISC) que não se classificaram para a final \textbf{Países} (MISC) que participaram no passado, mas não em 2018

O \textbf{Festival Eurovisão da Canção de 2018} (MISC) (em inglês : \textbf{Eurovision Song Contest} (ORG) 2018 ; em francês : \textbf{Concours Eurovision de la Chanson 2018} (MISC) ) foi a 63. edição do certame musical anual, que se realizou pela primeira vez \textbf{em Portugal} (LOC) após a vitória de \textbf{Salvador Sobral} (PER) na edição anterior , em \textbf{Kiev} (LOC), \textbf{Ucrânia} (LOC), com a canção \textbf{Amar} (ORG) pelos \textbf{Dois} (ORG) .

Os espetáculos do evento decorreram na \textbf{Altice Arena} (LOC) , em \textbf{Lisboa} (LOC) . As semifinais ocorreram nos dias 8 e 10 de maio e a final aconteceu a 12 de maio .

O slogan escolhido para a edição - \textbf{All Aboard!} (MISC) (em português : \textbf{Todos a Bordo!} (MISC) ) - foi revelado a 7 de novembro de 2017 . Por esta altura, a \textbf{RTP} (ORG) havia já confirmado a presença de 43 países no evento.



Localização

O local escolhido para realizar os três espetáculos (semifinais e final) foi a \textbf{Altice Arena} (PER) (anteriormente designada de \textbf{MEO Arena} (LOC) e também conhecida como \textbf{Pavilhão Atlântico} (LOC) ), situada na capital do país, \textbf{Lisboa} (LOC) .

O pavilhão, construído para a \textbf{Exposição Mundial de 1998} (MISC) (abreviadamente conhecida como \textbf{Expo '98} (MISC) ), é um espaço destinado a atrações públicas, festivais e outros espetáculos. Localizado na freguesia do \textbf{Parque das Nações} (LOC) , tem capacidade para 20 000 espectadores, sendo o maior pavilhão de espetáculos de \textbf{Portugal} (LOC) e o terceiro maior da \textbf{Europa} (LOC).

Este tem ligações de metro ao \textbf{Aeroporto Humberto Delgado} (LOC) e comboio nas proximidades, mais precisamente pela \textbf{Gare do Oriente} (LOC) , que o conecta ao resto do país e à \textbf{União Europeia} (ORG) . É também acessível por outros meios de transporte, como o autocarro ou o táxi.

A \textbf{Altice Arena} (PER) oferece ainda serviços como a bilheteira \textbf{Blueticket} (PER) , bares e multibanco e está adaptada a pessoas com mobilidade condicionada, visto que tem lugares específicos para estas assistirem aos eventos. O parque de estacionamento mais próximo é o parque TEJO , a cerca de 5 minutos do local, sendo que a zona dispõe de vários outros parques.

Após a \textbf{RTP} (ORG) ter aceitado o desafio de organizar a edição de 2018, o que se confirmou na conferência de imprensa a seguir à final da edição de 2017, não houve qualquer fase oficial de licitação para a seleção da cidade anfitriã, pelo que, a 15 de maio desse ano, a estação pública de rádio e televisão de \textbf{Portugal} (LOC) realizou uma reunião de urgência onde foi decidido que \textbf{Lisboa} (LOC) seria a cidade anfitriã do evento.

Porém, logo após esse anúncio, várias vozes críticas se levantaram a contestarem a decisão da estação. Os críticos queriam que a emissora levasse a cabo um concurso oficial para decidir qual a cidade-sede do \textbf{Festival} (LOC), utilizando o argumento de que o processo acontecia desde 2011 e que se havia repetido todos os anos nos países anfitriões do certame.

Mesmo assim, a \textbf{RTP} (ORG) e a \textbf{UER} (ORG) (organismo responsável pelo concurso) resolveram avaliar 4 cidades \textbf{em Portugal} (LOC) com vista a selecionar a melhor para acomodar um evento desta dimensão: \textbf{Braga} (LOC) , \textbf{Gondomar} (LOC) , \textbf{Lisboa} (LOC) e \textbf{Santa Maria da Feira} (LOC) . A análise destas cidades seguiu critérios de avaliação específicos e ajustados aos padrões da entidade promotora para este evento e que não passaram só pela existência de uma sala de espetáculos, mas também por:

Durante o processo, os presidentes dos municípios de \textbf{Guimarães} (LOC) , \textbf{Faro} (LOC) e \textbf{Santa Maria da Feira} (LOC) revelaram o seu interesse em concorrer para organizar a \textbf{Eurovisão} (LOC) . O município do \textbf{Porto} (LOC) , considerado um possível candidato a receber o certame, declarou na voz do seu presidente, \textbf{Rui Moreira} (PER) , que não se mostraria interessado numa potencial corrida à organização do \textbf{Festival} (LOC), apontando como motivos os custos avultados que o evento poderia acarretar para a cidade e que a \textbf{Câmara Municipal} (LOC) não estaria disposta a pagar.

Por sua vez, o \textbf{Conselho Metropolitano do Porto} (ORG) , através do seu presidente, defendeu a realização do evento na área metropolitana portuense por considerar uma igualdade de circunstâncias entre \textbf{Lisboa} (LOC) e a cidade nortenha e preferir a descentralização do evento. \textbf{Emídio Sousa} (PER), presidente do organismo, declarou que "um equipamento de excelência como o \textbf{Europarque} (LOC) , localizado em \textbf{Santa Maria da Feira} (LOC) , poderia vir a receber o \textbf{Festival da Eurovisão} (MISC)", uma vez que, no raio de 30 quilómetros, contaria com "mais de 8000 quartos", referindo ainda a proximidade ao aeroporto e às grandes linhas ferroviárias.

A candidatura desta cidade ganhou ainda mais força a 31 de maio quando \textbf{Emídio Sousa} (PER) revelou que a candidatura do \textbf{Europarque} (PER) estaria a ser trabalhada e que "já existiam algumas conversas com a \textbf{RTP} (ORG) ". No mesmo dia, \textbf{Espinho} (PER) junta-se a \textbf{Santa Maria da Feira} (LOC) e \textbf{Gondomar} (LOC) entre os interessados da \textbf{Área Metropolitana do Porto} (LOC) a acolher o evento e o presidente da \textbf{Câmara Municipal} (LOC) da cidade, \textbf{Pinto Moreira} (PER), garante que "a nave polivalente de \textbf{Espinho} (LOC) " teria "uma área útil superior à \textbf{Altice Arena} (PER) ".

\textbf{Dias} (PER) mais tarde, a 2 de junho , \textbf{Frank-Dieter Freiling} (ORG), chefe do \textbf{Grupo de Referência da UER} (ORG) , afirmou que queria ver um concurso público para selecionar a cidade anfitriã do evento, apesar de saber que \textbf{Lisboa} (LOC) seria muito provavelmente a escolhida.

O presidente da câmara de Braga , \textbf{Ricardo Rio} (PER) , admitiu a 5 de junho que estariam em curso conversações com a \textbf{RTP} (ORG) para que o \textbf{Parque de Exposições de Braga} (LOC) (posteriormente renomeado \textbf{Fórum Braga} (ORG) ) fosse o local a sediar a edição do concurso. Segundo o autarca, a suposta candidatura não se trataria de " nenhuma pretensão concreta", uma vez que seria preciso saber primeiro as condições exigidas e, só depois, o município tomaria uma decisão sobre essa matéria, considerando, contudo, que a renovação prevista para aquele espaço colocaria \textbf{Braga} (PER) à frente de outras cidades interessadas em receber a \textbf{Eurovisão} (LOC) , nomeadamente a cidade vizinha de \textbf{Guimarães} (LOC) , que já havia manifestado interesse em acolher o evento no seu \textbf{Pavilhão Multiusos} (MISC).

A primeira reunião entre os responsáveis da \textbf{RTP} (ORG) e o \textbf{Grupo de Referência da UER} (ORG) teve lugar a 13 de junho , em \textbf{Genebra} (LOC) , onde os responsáveis portugueses entregaram uma pasta com as possíveis locais que poderiam acolher a competição.

Finalmente, foi a 25 de julho de 2017 que a \textbf{RTP} (ORG) e a \textbf{UER} (ORG) , em conferência de imprensa, revelaram que \textbf{Lisboa} (LOC) seria a cidade a sediar o \textbf{Festival Eurovisão da Canção 2018} (MISC) . \textbf{Jon Ola Sand} (PER) , supervisor executivo do concurso desde 2011 , indicou que "a cidade apresentou uma proposta exemplar" e que estaria ansioso para trabalhar em conjunto com a emissora portuguesa no sentido de "fazer a \textbf{Eurovisão} (LOC) mais emocionante de sempre".

No dia em que a cidade anfitriã e o local dos espetáculos foram revelados, houve também a confirmação dos locais onde iriam ocorrer outros eventos relacionados com o concurso, como a \textbf{Blue Carpet} (ORG) (em português: \textbf{Passadeira Azul} (LOC)), o local da \textbf{Eurovision Village} (MISC) (em português: \textbf{Aldeia Eurovisiva} (PER)) ou o \textbf{EuroClub} (MISC) (em português: \textbf{Clube Eurovisivo} (ORG)).

O presidente da \textbf{RTP} (ORG) , \textbf{Gonçalo Reis} (PER), confirmou que a \textbf{Câmara Municipal de Lisboa} (LOC) cederia a \textbf{Praça do Comércio} (LOC) (mais conhecida como \textbf{Terreiro do Paço} (LOC)) durante dez dias, mais concretamente de 4 de maio a 13 de maio de 2018 , para aí ser instalada a \textbf{Eurovision Village} (ORG) , com stands de promoção nacional e internacional. A \textbf{Eurovision Village} (ORG) esteve aberta ao público diariamente, de forma gratuita, das 15:00 às 23:00.

A cerimónia de abertura e a chegada à "\textbf{Passadeira Azul} (LOC)" (os organizadores optaram por substituir a tradicional cor vermelha pela azul, devido à proximidade do \textbf{rio Tejo} (LOC)) ocorreu junto ao moderno \textbf{Museu de Arte} (LOC), \textbf{Arquitetura e Tecnologia} (PER) (\textbf{MAAT} (LOC)), situado na freguesia de \textbf{Belém} (LOC), e foi apresentada por \textbf{Inês Lopes Gonçalves} (PER), \textbf{Cláudia Semedo} (PER) , \textbf{Pedro Granger} (PER) e \textbf{Pedro Penim} (PER).

O \textbf{EuroClub} (ORG), espaço social de festa mas também de descanso para as delegações, artistas, intérpretes e fãs devidamente acreditados, esteve inicialmente para ser recebido num espaço localizado no \textbf{Cais do Sodré} (LOC) , a boate "\textbf{Lust} (ORG) in Rio". Porém, a 9 de março de 2018 , a \textbf{RTP} (ORG) informou que o local iria ser alterado para a \textbf{Praça do Comércio} (LOC) , ficando no espaço do "\textbf{MoMe Lisboa} (LOC)".



\textbf{Formato} (MISC)A 16 de outubro de 2017 , \textbf{Gonçalo Reis} (PER), presidente do \textbf{Conselho de Administração da RTP} (ORG) , garantiu que \textbf{Portugal} (LOC) iria organizar "o \textbf{Festival da Eurovisão} (MISC) mais barato de sempre" mas que teria, mesmo assim, uma proposta criativa das mais interessantes alguma vez feitas, com capacidade para "projetar \textbf{Portugal} (LOC) ". "Posso hoje afirmar que vamos fazer o \textbf{Festival da Eurovisão} (MISC) mais barato de sempre e que estou crente que as nossas equipas criativas vão desenvolver uma das '\textbf{Eurovisões} (LOC)' mais interessantes de sempre", afirmou.

" As nossas equipas criativas estão a desenvolver o conceito [para] transmitir a imagem de um \textbf{Portugal} (LOC) contemporâneo, de um \textbf{Portugal} (LOC) aberto ao mundo, de uma cultura de inclusão, de uma cultura de tolerância, de uma cultura positiva ", disse também o gestor.

\textbf{Gonçalo Reis} (PER) defendeu que este iria "ser um projeto mobilizador para todas as indústrias criativas, para \textbf{Lisboa} (LOC) e para \textbf{Portugal} (LOC) " e que iria projetar o país , avisando que este se tratava de um acontecimento " que ultrapassa a \textbf{RTP} (ORG) ", apesar de concluir que a estação pública teria " todas as condições para produzir um grande evento de televisão, para promover as nossas indústrias criativas, a indústria do entretenimento " e que " o [\textbf{Festival da]} (MISC) \textbf{Eurovisão} (LOC) é algo que envolve a cidade , o turismo, todo o país .

A 21 de outubro de 2017 , \textbf{Daniel Deusdado} (PER), então diretor de programas da \textbf{RTP} (ORG) , revelou que a administração não tencionava cortar valores na grelha de programação devido à organização do evento internacional. " A administração já disse que tentará fazer o mais barato \textbf{Festival da Eurovisão} (PER) dos últimos anos (...) também já foi dito às equipas que planeiam o orçamento para 2018 que o valor importado às grelhas de programas não será afetado pela \textbf{Eurovisão ...} (MISC) mas são previsões. A determinação da equipa que está a gerir é jogar com este equilíbrio. ", afirmou.

\textbf{Gonçalo Madaíl} (PER), subdiretor da \textbf{RTP1} (ORG) , garantiu na altura que existiria um programa de voluntários para o \textbf{Festival Eurovisão 2018} (MISC) . " Um projeto como este requer a participação de muitos voluntários a vários níveis com alguns requisitos, como o dominío da língua inglesa (...) Nesta fase não está previsto o modelo, mas será comunicado em breve ", declarou.

Depois da polémica em torno da escolha da cidade anfitriã do certame em 2017, do banimento da cantora russa \textbf{Julia Samoylova} (PER) e da retirada da \textbf{Rússia} (LOC) do concurso, a \textbf{UER} (ORG) lançou, a 31 de julho de 2017 , um novo conjunto de regras do \textbf{Festival da Eurovisão} (MISC) que estariam em vigor na edição de 2018 , sediada \textbf{em Portugal} (LOC) .

A partir desse momento, a emissora anfitriã do \textbf{Festival} (PER) teria a obrigação de "garantir que todos os concorrentes possam atuar pessoalmente no palco do concurso", sendo que todos os organismos de radiodifusão estão proibidos de enviar membros da delegação que "tenham antecedentes susceptíveis de induzir as autoridades nacionais do país anfitrião a negar o acesso ao território".

As alterações às regras, enviadas a todas as emissoras nacionais e publicadas pela emissora alemã \textbf{ARD} (ORG) , são uma resposta específicas às ações das delegações da \textbf{Ucrânia} (LOC) e da \textbf{Rússia} (LOC) na temporada passada sendo que, segundo as novas regras, os dois países seriam declarados culpados. Além disso, a \textbf{UER} (ORG) acrescentou vários parágrafos sobre a natureza não política da competição, algo que foi violado pelos dois países na edição passada, bem como a imparcialidade dos membros do júri, sendo uma reação às críticas das votações da \textbf{Arménia} (LOC) , \textbf{Azerbaijão} (LOC) e \textbf{Ucrânia} (LOC) em \textbf{Kiev} (LOC) .

A 16 de maio de 2017 , o presidente da \textbf{RTP} (ORG) , \textbf{Gonçalo Reis} (PER), afirmou que a estação pública iria "encontrar condições para fazer a \textbf{Eurovisão} (LOC) sem excessos" e que seria a "Eurovisão mais económica dos últimos anos". A 26 de julho de 2017 , o \textbf{Presidente da Câmara Municipal de Lisboa} (MISC) , \textbf{Fernando Medina} (PER) , anunciou que a taxa turística, em conjunto com a \textbf{UER} (ORG) , iria financiar o evento, cujo orçamento não foi divulgado.

A 21 de outubro de 2017, \textbf{Daniel Deusdado} (PER), diretor de programas da \textbf{RTP} (ORG) , questionado sobre os custos da organização, revelou que " a \textbf{Eurovisão} (LOC) enquanto organismo traz 5 milhões de receitas e \textbf{Lisboa} (LOC) contribui com o local do evento e outras funcionalidades ", destacando também a contribuição do \textbf{Turismo de Portugal} (ORG) e os fundos oriundos das receitas publicitárias. Contudo, " o orçamento só será finalizado no final do ano ". [ carece de fontes ? ]

A 12 de novembro de 2017, foi revelado que o orçamento para o \textbf{Festival Eurovisão 2018} (MISC) seria um total de 23 milhões de euros, repartidos pelos diversos organismos. A \textbf{RTP} (ORG) entraria diretamente com 12 milhões de euros e o restante seria financiado pela \textbf{Câmara Municipal de Lisboa} (LOC) , pelo \textbf{Turismo de Portugal} (ORG) , pela \textbf{EBU} (MISC) e por patrocinadores.

Já em 15 de março de 2018, a \textbf{Câmara Municipal de Lisboa} (LOC) aprovou uma proposta que aponta para uma despesa por parte do município de 5 milhões de euros, dos quais 2 milhões serão alocados só ao aluguer do \textbf{Altice Arena} (PER) .

O slogan para o festival, \textbf{All Aboard!} (MISC) (em português : \textbf{Todos a Bordo!} (MISC)) , foi revelado a 7 de novembro de 2017.

\textbf{Portugal} (LOC), como país, sempre conectou a \textbf{Europa} (LOC) com o resto do mundo através do oceano e, há 500 anos, \textbf{Lisboa} (LOC) foi o centro de muitas das rotas marítimas mais importantes do mundo. Hoje, \textbf{Lisboa} (LOC) usou a conectividade do oceano como inspiração para o slogan \textbf{All Aboard!} (MISC) , convidando a comunidade internacional para se juntar na edição do festival da \textbf{Eurovisão} (LOC) deste ano.

A identidade visual é caracterizada por motivos oceânicos que aludem à localização de \textbf{Lisboa} (LOC) junto à costa e ao \textbf{oceano Atlântico} (LOC) . Em paralelo com o logótipo principal, que retrata uma concha estilizada, foram projetados doze logótipos suplementares para simbolizar os diferentes aspetos do ecossistema marinho .

A 8 de Janeiro de 2018, a \textbf{RTP} (ORG) e a \textbf{UER} (ORG) anunciaram que o concurso seria apresentado, pela primeira vez, por quatro apresentadoras do sexo feminino, sendo o conjunto composto pelas apresentadoras da \textbf{RTP Catarina Furtado} (ORG) , \textbf{Filomena Cautela} (MISC) e \textbf{Sílvia Alberto} (LOC) , juntamente com a atriz \textbf{Daniela Ruah} (PER) . Seria a primeira vez desde 2013 que o concurso não incluiria um apresentador do sexo masculino e o segundo ano consecutivo que os apresentadores seriam do mesmo sexo.

O sorteio para determinar em que semifinal cada país participaria decorreu a 29 de janeiro de 2018, nos \textbf{Paços do Concelho da Câmara Municipal de Lisboa} (LOC) . Os 37 semifinalistas foram distribuídos por seis potes, com base em padrões de votação históricos calculados pelo parceiro oficial do televoto do concurso, a empresa Digame Arquivado em 25 de março de 2018, no \textbf{Wayback Machine} (LOC) .. A distribuição dos diferentes potes ajuda a reduzir a chance de um chamado "voto de bloco" e aumenta o suspense nas meias-finais. A cerimónia foi apresentada por \textbf{Sílvia Alberto} (PER) e \textbf{Filomena Cautela} (LOC) e incluiu a passagem do testemunho da \textbf{Eurovisão de Vitali Klitschko} (MISC) , presidente da câmara de \textbf{Kiev} (LOC) (sede do festival anterior), para \textbf{Fernando Medina} (PER) , presidente da câmara de \textbf{Lisboa} (LOC) .

A equipa técnica do evento contou com 23 nomes, sendo que mais de 70\% eram portugueses. Houve ainda responsáveis vindos da \textbf{Ucrânia} (LOC) , \textbf{Bélgica} (LOC) , \textbf{Suécia} (LOC) e \textbf{Estados Unidos da América} (LOC) , entre outros países. Mais de 700 pessoas estiveram diretamente envolvidas no evento, incluindo cerca de 400 voluntários.



\textbf{Participantes} (MISC)A \textbf{UER} (ORG) anunciou a 7 de novembro de 2017 que 42 países participariam na edição do \textbf{Festival} (MISC). A \textbf{Rússia} (LOC) confirmou o regresso depois de uma saída na edição precedente , onde a participação da \textbf{Macedónia} (LOC) (\textbf{ARJM} (ORG)) foi bloqueada provisoriamente pela \textbf{UER} (ORG) devido a dívidas não pagas pela estação nacional do país. Entretanto, dez dias depois, a \textbf{UER} (ORG) anunciou que a \textbf{ARJ da Macedónia} (LOC) (atualmente \textbf{Macedónia do Norte} (LOC)) estaria autorizada a participar no \textbf{Festival} (LOC), aumentando assim o número de países participantes para 43, o maior desde as edições de 2008 e 2011 .

Os países do \textbf{Big Five} (MISC) ( \textbf{Alemanha} (LOC) , \textbf{França} (LOC) , \textbf{Itália} (LOC) , \textbf{Espanha} (LOC) , \textbf{Reino Unido} (LOC) ) e o país anfitrião \textbf{Portugal} (LOC) , desde que participam, estão diretamente qualificados para a final. Além disso, dez países vêm das duas semi-finais, de modo a ter 26 países na final. Na final, todos os países participantes no concurso terão o direito de votar.



\textbf{Festival Dos 37} (MISC) países que participaram nas semifinais desta edição, 19 países participaram na primeira semifinal, realizada no dia 8 de maio (terça-feira), enquanto que os outros 18 participaram na segunda semifinal, realizada no dia 10 de maio (quinta-feira).

Em cada semifinal, classificaram-se para a final os 10 países que tiveram melhor pontuação. Na primeira, votaram os países participantes e mais 3 finalistas: \textbf{Portugal} (LOC) , \textbf{Espanha} (LOC) e \textbf{Reino Unido} (LOC) . Na segunda, também votaram os países participantes, bem como os outros 3 finalistas: \textbf{França} (LOC) , \textbf{Alemanha} (LOC) e \textbf{Itália} (LOC) . Em cada espetáculo (semifinais e final), nenhum país pôde votar em si próprio.

A final ocorreu no dia 12 de maio (sábado), contando com 26 países concorrentes (10 de cada semifinal + países do \textbf{Big 5} (MISC) : \textbf{Alemanha} (LOC), \textbf{Espanha} (LOC), \textbf{França} (LOC), \textbf{Itália} (LOC) e \textbf{Reino Unido} (LOC) + país anfitrião: \textbf{Portugal} (LOC)).

A ordem de atuação foi revelada no dia 3 de abril de 2018.

Os países em negrito deram o máximo de 24 pontos (12 pontos do júri e do televoto) ao participante especificado. Os países que receberam 12 pontos foram os seguintes:

Os países em negrito deram o máximo de 24 pontos (12 pontos do júri e do televoto) ao participante especificado. Os países que receberam 12 pontos foram os seguintes:

Os países em negrito deram o máximo de 24 pontos (12 pontos do júri e do televoto) ao participante especificado. Os países que receberam 12 pontos foram os seguintes:

Os países em negrito deram o máximo de 24 pontos (12 pontos do júri e do televoto) ao participante especificado. Os países que receberam 12 pontos foram os seguintes:

Os países em negrito deram o máximo de 24 pontos (12 pontos do júri e do televoto) ao participante especificado. Os países que receberam 12 pontos foram os seguintes:

Os países em negrito deram o máximo de 24 pontos (12 pontos do júri e do televoto) ao participante especificado. Os países que receberam 12 pontos foram os seguintes:



Outros países

Os membros associados são as agências audiovisuais de caráter nacional que operam fora da \textbf{Área Europeia de Radiodifusão} (ORG) .

No comunicado pode ler-se: " a participação no Festival Eurovisão da Canção é limitada a membros da \textbf{UER} (ORG) ou associados convidados especialmente. Associados da \textbf{UER} (ORG) podem ser elegíveis para participar na \textbf{Eurovisão} (LOC) , isto é decidido pelo \textbf{Grupo de Referência} (ORG), o corpo governativo do \textbf{Festival da Eurovisão} (MISC) , numa base de dia para dia. O anúncio dos países participantes em 2018 será feito em tempo oportuno ".

No início do ano a \textbf{Khabar Agency} (PER) , televisão pública cazaque, afirmou não saber ainda se se estrearia na \textbf{Eurovisão 2018} (MISC) mas tem vontade que tal aconteça. A \textbf{Khabar Agency} (PER) , tornou-se membro associado da \textbf{UER} (ORG) a 1 de janeiro de 2016 .

No entanto, tal não acontecerá: a \textbf{União Europeia de Radiodifusão} (ORG) decidiu pôr um travão a esta participação. A RTK fez a sua pré-inscrição para o Festival Eurovisão da Canção 2018 e terá recebido feedback positivo da \textbf{RTP} (ORG) e do \textbf{Grupo de Referência} (ORG). Mas, mais tarde, a \textbf{UER} (ORG) recuou e decidiu excluir o país por este não ser membro da \textbf{Organização das Nações Unidas} (ORG) . A 27 de setembro de 2017 , \textbf{Mentor Shala} (PER), diretor-geral da \textbf{RTK} (ORG) , explicou, numa carta enviada ao site \textbf{ESCToday} (MISC), que " a RTK respondeu a todos os requisitos da \textbf{UER} (ORG) para participar na \textbf{Eurovisão} (LOC) . No ano passado, o \textbf{Grupo de Referência} (ORG) deu-nos uma resposta positiva sobre a nossa participação com a condição de que o país organizador reconhecesse a nossa independência. Visto que a \textbf{Ucrânia} (LOC) não nos reconhece como independentes, tal não aconteceu. Após a vitória de \textbf{Portugal} (LOC) , tínhamos a certeza que iríamos participar na \textbf{Eurovisão} (LOC) visto que \textbf{Portugal} (LOC) reconhece a nossa independência . O \textbf{Grupo de Referência} (ORG) voltou a dar-nos uma resposta positiva. Inscrevemo-nos oficialmente e obtivemos o apoio de muitas entidades. Estávamos confiantes que a resposta seria positiva. Mas ficamos muito surpreendidos quando recebemos uma carta da \textbf{UER} (ORG) a dizer que o \textbf{Kosovo} (LOC) não pode participar na \textbf{Eurovisão} (LOC) , por enquanto, pois não é membro da \textbf{ONU} (ORG) . Isto é uma razão absurda visto que o \textbf{Kosovo} (LOC) foi admitido em todas as federações e outras organizações mundiais como a \textbf{FIFA} (ORG) , \textbf{UEFA} (ORG) , \textbf{FIBA} (ORG) , \textbf{WB} (MISC) , \textbf{Jogos Olímpicos} (MISC) , \textbf{FMI} (ORG) , etc. \textbf{Sabemos} (MISC) que não podemos ser aceites como membros de pleno direito na \textbf{UER} (ORG) porque não pertencemos à \textbf{ONU} (ORG) mas nós só queríamos participar na \textbf{Eurovisão !} (MISC) Não é de todo justo e não há razão para bloquear a participação de um país num festival destes. Lamentamos que uma organização não política como a \textbf{UER} (ORG) esteja a fazer política no caso do \textbf{Kosovo} (LOC) . Não desistiremos e tenho a certeza que um dia verão o \textbf{Kosovo} (LOC) na \textbf{Eurovisão} (LOC) ".

A australiana SBS , por se encontrar fora da zona de influência da \textbf{UER} (ORG) , tem apenas o estatuto de membro associado mas participa na \textbf{Eurovisão} (LOC) pois recebe um convite especial todos os anos. O \textbf{Kosovo} (LOC) esperava ter o mesmo tratamento. A \textbf{União Europeia de Radiodifusão} (ORG) já reagiu a este assunto dizendo que "a RTK não pode participar na \textbf{Eurovisão} (LOC) pois não é membro nem associada da \textbf{UER} (ORG) . Os estatutos da \textbf{UER} (ORG) dizem que os membros têm de fazer parte da \textbf{União Internacional das Telecomunicações} (ORG) ou ser membros do \textbf{Conselho da Europa} (ORG) . O \textbf{Kosovo} (LOC) não faz parte de qualquer uma destas. A \textbf{UER} (ORG) ajudou a estabelecer a \textbf{RTK} (ORG) em 1999 e continua a trabalhar com a emissora para proteger o serviço público no país . O \textbf{Grupo de Referência} (ORG) tem, no entanto, um mandato para criar exceções às regras ".



Álbum oficial

\textbf{Eurovision Song Contest} (ORG): Lisbon 2018 é o álbum de compilação oficial da edição do \textbf{Festival} (MISC), lançado pela \textbf{União Europeia de Radiodifusão} (ORG) e distribuído pela \textbf{Universal Music Group} (ORG) digitalmente a 20 de Abril de 2018. O álbum conta com todas as 43 participações, incluindo os semifinalistas que não se classificaram para a final.



Festival Eurovisão da Canção 2019



Final: "\textbf{Heroes} (MISC)" interpretado por \textbf{Conchita Wurst} (PER) "\textbf{Fuego} (PER)" interpretado por \textbf{Måns Zelmerlöw} (PER) "\textbf{Dancing Lasha Tumbai} (MISC)" interpretado por \textbf{Eleni Foureira "Toy} (PER)" interpretado por \textbf{Verka Serduchka "} (PER)Hallelujah" interpretado por \textbf{Gali Atari} (PER) "\textbf{Bo'ee} (ORG) - \textbf{Come to Me} (MISC)" interpretado por \textbf{Idan Raichel Dana International Gal Gadot} (PER) " \textbf{Like a Prayer} (MISC) " interpretado por \textbf{Madonna Lior Suchard Países} (PER) que classificaram-se para a final \textbf{Países} (MISC) que não se classificaram para a final \textbf{Países} (MISC) que participaram no passado, mas não em 2019

O \textbf{Festival Eurovisão da Canção de 2019} (MISC) (em inglês : \textbf{Eurovision Song Contest} (ORG) 2019 ; em francês : \textbf{Concours Eurovision} (MISC) de la chanson 2019 ; em hebraico :  2019 ) foi a 64. edição anual do evento. O festival foi realizado pela terceira vez em \textbf{Israel} (LOC) , sendo que esta foi a primeira vez que o evento foi realizado em \textbf{Tel Aviv} (LOC) , depois de \textbf{Netta Barzilai} (MISC) ter ganhado a edição anterior com a canção \textbf{Toy} (MISC) . As semifinais realizaram-se nos dias 14 e 16 de maio, enquanto que a final se deu no dia 18 de maio. \textbf{Israel} (LOC) já havia sediado duas edições anteriores em 1979 e em 1999 em \textbf{Jerusalém} (LOC) .

Quarenta e um países participam no \textbf{Festival} (LOC), com a \textbf{Bulgária} (LOC) e a \textbf{Ucrânia} (LOC) ausentes, ambos pela primeira vez desde 2015 .



Localização

A última vez que o \textbf{Festival} (LOC) foi realizado em solo israelita foi em 1999 , em \textbf{Jerusalém} (LOC) . A primeira vitória foi em 1979 .

O local escolhido para realizar os três espetáculos (semifinais e final) foi a \textbf{Expo Tel Aviv} (MISC) , situada em \textbf{Tel Aviv} (LOC) .

A 11 de julho de 2018, a emissora israelita \textbf{KAN} (ORG) - iria organizadar, a \textbf{Eurovisão} (LOC), após o desaparecimento da \textbf{IBA} (MISC) em 2017 - publicou os critérios técnicos para a escolha da cidade sede:

A 13 de maio de 2018, um dia depois do final da edição de 2018 do \textbf{Festival} (LOC), vencido pela representante de \textbf{Israel} (LOC), o primeiro-ministro israelita \textbf{Benyamin Netanyahu} (PER) anunciou no \textbf{Twitter} (MISC) que a competição seria realizada em \textbf{Jerusalém} (LOC) . Esta informação não foi, contudo, confirmada pela \textbf{UER} (ORG) .

O município de \textbf{Jerusalém} (LOC) mostrou rapidamente grande interesse em sediar a competição e propôs, a partir de 13 de maio de 2018, que o \textbf{Festival de 2019} (MISC) fosse realizado no \textbf{Pais Arena} (LOC) ou no \textbf{Teddy Stadium} (LOC). O presidente da câmara de \textbf{Tel Aviv} (LOC) declarou a 13 de maio de 2018 que a sua cidade não apresentaria uma proposta para sediar a competição. Mais tarde, a administração de \textbf{Tel Aviv} (LOC) explicou que esse anúncio foi feito em respeito ao governo israelita e por forma a deixá-lo decidir, e que \textbf{Tel Aviv} (LOC) estaria pronta para sediar o evento se \textbf{Jerusalém} (LOC) não fosse escolhida.

A 16 de maio de 2018, o presidente da câmara da cidade de \textbf{Petah Tikva} (LOC) anunciou o interesse da cidade em sediar o \textbf{Festival de 2019} (MISC). No entanto, a cidade não possui uma sala grande o suficiente para o evento.

A 10 de junho de 2018, a cidade de \textbf{Haifa} (LOC) anunciou o seu interesse em sediar a competição. No mesmo dia, \textbf{Eilat} (PER) é anunciada como uma possível cidade anfitriã para a competição. No entanto, como \textbf{Petah Tikva} (MISC), não dispõe de infraestruturas suficientes para sediar a competição, estas teriam que ser construídas de raiz, caso a cidade fosse selecionada.

A 18 de junho de 2018, o primeiro-ministro \textbf{Benjamin Netanyahu} (PER) declarou que \textbf{Israel} (LOC) se havia comprometido a permanecer em conformidade com as regras da \textbf{UER} (ORG) relativas à constituição de emissoras membros, por forma a assegurar a realização do \textbf{Festival} (PER) em solo israelita. As regras de estabelecimento da \textbf{KAN} (ORG) incluíam agora a condição de que a programação de notícias fosse posteriormente delegada a uma segunda entidade pública de radiodifusão, o que violará no futuro as regras da \textbf{EBU} (MISC), as quais exigem que os radiodifusores membros tenham os seus próprios departamentos internos de notícias.

A 19 de junho de 2018, uma reunião entre o \textbf{Grupo de Referência da UER} (ORG) e a emissora israelita \textbf{KAN} (ORG) foi realizada em \textbf{Genebra} (LOC) , na sede da \textbf{UER} (LOC), para iniciar os preparativos para a competição. Após esta reunião, \textbf{Israel} (LOC) foi oficialmente confirmado como o país anfitrião para a edição de 2019 da competição - rumores circulavam de que a \textbf{Áustria} (LOC), que ficou em terceiro lugar em 2018 e sediou a edição de 2015, seria designada como país anfitrião. No final desta reunião, a \textbf{KAN} (ORG) e a \textbf{UER} (ORG) anunciaram que a cidade-sede e as datas da competição seriam anunciadas até setembro de 2018.

No dia seguinte, \textbf{Israel} (LOC) foi oficialmente confirmado como o país anfitrião. Assim, na semana seguinte, o processo de escolha da cidade-sede foi aberto .

A 28 de julho de 2018, o ministro israelita \textbf{Michael Oren} (PER) , o qual está intimamente ligado ao primeiro-ministro \textbf{Netanyahu} (LOC), declarou que \textbf{Jerusalém} (LOC) não tinha condições financeiras e organizacionais para sediar o evento e que ele provavelmente seria realizado pela primeira vez em \textbf{Tel Aviv} (LOC). Logo após, relatórios do governo foram divulgados e comprovaram que o governo local não poderia depositar a caução de 12 milhões de euros. Assim, a \textbf{KAN} (ORG) foi forçada a pedir um empréstimo do valor para poder arcar com as despesas relativas a esta caução.

Depois de diversos impasses entre a \textbf{KAN} (ORG) e o governo central, era praticamente certo de que o festival não seria realizado em solo israelita. Faltando apenas algumas horas para que as condições por parte da emissora tivessem de ser garantidas, foi alcançado um acordo entre as duas partes. O \textbf{Ministério da Agricultura} (ORG) autorizou o empréstimo de 12 milhões de euros, que serão reservados para eventuais despesas organizacionais como a hospedagem e a locomoção das delegações pela cidade. Se o acordo não fosse alcançado, o presidente da câmara de \textbf{Tel Aviv} (LOC) já havia avisado que a cidade já tinha a quantia necessária separada e iria ceder o \textbf{Centro de Convenções da Cidade} (LOC) para a realização do festival .

Na semana de 27 de agosto de 2018, o supervisor executivo do evento \textbf{Jon Ola Sand} (PER) coordenou uma visita de uma delegação oficial da entidade por diversas cidades em \textbf{Israel} (LOC) para avaliar as propostas de \textbf{Jerusalém} (LOC) e \textbf{Tel Aviv} (LOC), além de assistir uma apresentação da cidade portuária de \textbf{Eilat} (PER) . A 30 de agosto de 2018, \textbf{Sand} (PER) declarou numa entrevista a \textbf{KAN} (ORG), que \textbf{Eilat} (PER) havia retirado a sua candidatura, deixando-a entre \textbf{Jerusalém} (LOC) e \textbf{Tel Aviv} (LOC) .

A 13 de setembro de 2018, a \textbf{EBU} (MISC) anunciou o \textbf{Centro de Convenções de Tel Aviv} (LOC) como a sede do certame .

O \textbf{Eurovision Village} (MISC) é a área oficial dos fãs e patrocinadores do \textbf{Festival Eurovisão da Canção} (MISC) durante a semana dos eventos, onde é possível assistir a apresentações de artistas locais, bem como a transmissão de espectaculoso ao vivo do local principal. Ele ficará localizado no \textbf{parque Charles Clore} (LOC) em \textbf{Tel Aviv} (LOC) .

O \textbf{EuroClub} (ORG) é o palco para as after-parties oficiais e apresentações privadas dos participantes do \textbf{Festival} (LOC). Ao contrário do \textbf{Eurovision Village} (MISC), o acesso ao \textbf{EuroClub} (ORG) é restrito aos fãs, delegações e imprensa credenciados. Ele está localizado no \textbf{Hangar 11} (LOC) no \textbf{Porto de Tel Aviv} (LOC) .

O evento "\textbf{Orange Carpet} (PER)" e a \textbf{Cerimonia de Abertura} (MISC), onde todos os participantes e as suas delegações são apresentados à imprensa credenciada e fãs, aconteceu na \textbf{Praça Habima} (LOC) , no centro de \textbf{Tel Aviv} (LOC), no dia 12 de maio de 2019 às 19:00 \textbf{IDT} (MISC) (17:00 UTC+1 hora de \textbf{Lisboa} (LOC)), seguido por um pequeno evento no \textbf{Auditório Charles Bronfman} (MISC) .



\textbf{Formato} (PER)O slogan deste ano é \textbf{"Dare to Dream} (MISC)" (em português europeu : Atreva-se a sonhar e em hebraico  ).

"Este slogan serve como uma inspiração e representa e simboliza tudo o que o \textbf{Festival Eurovisão da Canção} (MISC) é. É sobre inclusão. É sobre diversidade. É sobre a unidade", disse \textbf{Jon Ola Sand} (PER) durante o anúncio em \textbf{Tel Aviv} (LOC). "Estar nesse palco, ousar sonhar, pensar que você pode estar aqui, ser corajoso o suficiente, estar confiante estar lá se apresentando para uma audiência mundial é algo sonhado por muitos. Foi o que \textbf{Netta} (LOC) fez no ano passado, quando se apresentou em \textbf{Lisboa} (LOC). Ela foi a esse palco com um sonho - o sonho para trazer o \textbf{Festival Eurovisão da Canção} (MISC) de volta a \textbf{Israel} (LOC). E ela conseguiu. E neste ano, em maio, em \textbf{Tel Aviv} (LOC), todos nos reuniremos novamente para celebrar os bons valores do \textbf{Festival Eurovisão da Canção} (MISC) e faremos isso aqui, com a ajuda da emissora \textbf{KAN} (ORG) e da equipe israelita\textbf{"} (MISC), concluiu. O slogan apela à diversidade, inclusão e união.

\textbf{Inspirado} (LOC) pelo slogan do evento, \textbf{Dare to Dream} (MISC) , a \textbf{UER} (ORG) e a emissora anfitriã apresentaram o projeto visual do evento de 2019; três triângulos que, quando unidos, brilham juntos para criar uma estrela dourada. Ao revelar o projeto escolhido para este ano, o canal anfitrião divulgou o porque da escolha do logótipo :

A identidade gráfica serviu de inspiração para o palco e a \textbf{green room} (ORG) desenhados pelo designer de \textbf{Florian Wieder} (LOC) . O palco terá 1.000 metros quadrados de \textbf{LED} (ORG), juntamente com 300 triângulos móveis. Ao contrário de em anos anteriores, o \textbf{LED} (ORG) será o principal protagonista do palco e, assim, serão usados aproximadamente 1 quilómetro quadrado de \textbf{LEDs} (ORG). Devido ao tamanho do \textbf{Centro} (LOC) haverá ainda mais 540 metros quadrados, na \textbf{green Room} (MISC) , que ficará noutro pavilhão, devido ao pequeno espaço disponível para o palco principal.

A 25 de janeiro de 2019, a \textbf{UER} (ORG) e a emissora \textbf{KAN} (ORG) anunciam os nomes dos quatro apresentadores desta edição. São duas mulheres e dois homens:a top model internacional \textbf{Bar Refaeli} (PER) , que apresenta as versões locais dos programas \textbf{Million Dollar Shooting Star} (MISC) e \textbf{The X Factor} (MISC) , \textbf{Erez} (PER) Tal ,um dos apresentadores mais conhecidos do país e que foi o comentarista do \textbf{Festival de 2018} (MISC); \textbf{Assi Azar} (PER) , também apresentador de televisão,mais conhecido por apresentar a versão original do \textbf{Rising Star} (ORG) e \textbf{Lucy Ayoub} (PER),que apresenta diversos programas relacionados a cultura na emissora.

Com cerca de uma centena de conferências de imprensa agendadas para as duas semanas eurovisivas no \textbf{Centro de Imprensa do Festival} (LOC) Eurovisão 2019, Sivan Avrahami e \textbf{Nadav Embon} (PER) foram os escolhidos para a condução das mesmas.

O sorteio das semi-finais aconteceu a 28 de janeiro de 2019 às 15:00 (GMT) no \textbf{Museu de Arte de Tel Aviv} (LOC) . Os trinta e seis semi-finalistas foram distribuídos por seis potes, com base nos padrões de votação históricos calculados pelo parceiro oficial de voto do \textbf{Festival} (LOC), \textbf{Digame} (LOC).

Para distribuir os países nas diferentes semi-finais, foi feito um sorteio. No mesmo sorteio, os países que se qualificam automaticamente para a final - \textbf{Israel} (LOC), país anfitrião, bem como os cinco principais contribuintes financeiros da \textbf{UER} (ORG): \textbf{Alemanha} (LOC) , \textbf{Espanha} (LOC) , \textbf{França} (LOC) , \textbf{Itália} (LOC) e o \textbf{Reino Unid} (LOC) o - são também atribuídos à sorte, por sua vez, à semi-final durante a qual terão direito a voto .

Este sorteio foi, então, baseado nos potes seguintes:

A 30 de março de 2019, a \textbf{EBU} (MISC) anunciou que a apresentação do resultado do televoto durante a grande final mudaria pela primeira vez desde que o atual sistema de votação foi introduzido em 2016. A apresentação dos resultados do júri permanecerá a mesma com um porta-voz em direto em cada país participante, revelando a música de topo do seu júri nacional que ganhou 12 pontos. Numa mudança em relação aos anos anteriores, o resultado do televoto será revelado na ordem da posição dos jurados, do mais baixo ao mais alto.

A 1 de março de 2019, foi anunciado que o mentalista \textbf{Lior Suchard} (PER) se apresentará como convidado especial durante a final. Foi confirmado a 8 de abril de 2019 que \textbf{Madonna} (PER) apresentará duas músicas durante a atuação do intervalo da final. A 15 de abril de 2019, a \textbf{UER} (ORG) divulgou mais informações sobre os atos de abertura e intervalo. A primeira semifinal será aberta por \textbf{Netta Barzilai} (PER) , que fará uma nova versão de sua música vencedora, \textbf{" Toy} (MISC) ", além de um single ainda a ser lançado. Na segunda semi-final, os \textbf{Shalva Band} (PER) irão trazer a música \textbf{"Million Dreams} (MISC)". A \textbf{Grande Final} (LOC) inclui performances de seis ex-participantes do \textbf{Festival Eurovisão da Canção} (MISC): \textbf{Conchita Wurst} (PER) , que interpretará " \textbf{Heroes} (MISC) ", \textbf{Måns Zelmerlöw} (MISC) com \textbf{"} (MISC)\textbf{Fuego} (MISC)", \textbf{Eleni Foureira} (PER) com \textbf{"Dancing Lasha Tumbai} (PER)", \textbf{Verka Serduchka com "Toy} (PER)", \textbf{Gali Atari} (PER) com a sua música vencedora " \textbf{Hallelujah " e Dana International} (MISC) com uma canção ainda não anunciada. \textbf{Idan Raichel} (PER) canta a música \textbf{"Bo'ee - Come To} (MISC) Me" junto com os 26 finalistas, enquanto a atriz e modelo israelita \textbf{Gal Gadot} (PER) também fará uma aparição na \textbf{Grande Final} (LOC).



Países participantes

O número total de países será quarenta e um (menos 2 do que na edição anterior, devido à saída da \textbf{Bulgária} (LOC) , por motivos financeiros, e da \textbf{Ucrânia} (LOC) , por motivos diversos ). Das quarenta e uma músicas, vinte e seis são em inglês, dez em idiomas nativos e cinco em inglês com variações noutros idiomas.

A primeira semi-final aconteceu a 14 de maio de 2019 às 22:00h na hora de \textbf{Israel} (LOC) (20:00h UTC+1 ) . Dezassete países participaram na primeira semi-final. \textbf{Espanha} (LOC) , \textbf{França} (LOC) e \textbf{Israel} (LOC) votam nesta semi-final. A \textbf{Ucrânia} (LOC) estava colocada na segunda parte da semi-final, no entanto anunciou a sua retirada devido a controvérsias com a emissora.

A segunda semi-final aconteceu a 16 de maio de 2019 às 22:00h na hora de \textbf{Israel} (LOC) (20:00h UTC+1 ) . Dezoito países participaram na segunda semi-final. \textbf{Alemanha} (LOC) , \textbf{Itália} (LOC) e \textbf{Reino Unido} (LOC) votam nesta semi-final. A \textbf{Suiça} (LOC) realizou um pedido para atuar na segunda semi-final .

A final será a sábado 18 de maio às 22:00h na hora de \textbf{Israel} (LOC) (20:00h UTC+1 ) . Vinte e seis países participam na final, com todos os 41 países participantes elegíveis a votarem.



Outros países

Para um país para ser elegível para a participar no \textbf{Festival Eurovisão da Canção} (MISC) , precisa de ser um membro ativo da \textbf{União Europeia de Radiodifusão} (ORG) (\textbf{UER} (ORG)) ou, para membros associados, ter a sua participação aprovada pela instituição.



\textbf{Transmissão do Festival} (MISC)A maioria dos países enviará comentadores para \textbf{Tel Aviv} (LOC) ou comentar no seu próprio país, a fim de acrescentar informações aos participantes e, se necessário, fornecer informações sobre as votações.

Os porta-vozes anunciam a pontuação de 12 pontos do júri nacional do seu respetivo país.



Incidentes e controvérsias

Depois de diversos problemas financeiros e de gestão, em 2014, o governo israelita decidiu apresentar um projeto de lei que determinava que a \textbf{IBA} (MISC) fosse dissolvida e uma nova empresa pública fosse criada para substituir as suas funções, a qual chamar-se-ia \textbf{Empresa Pública Israelita de Radiodifusão} (ORG) . A 1 de outubro de 2016 esta lei foi aprovada pelo parlamento local. No entanto, a 19 de julho de 2016, o \textbf{Primeiro Ministro Benjamin Netanyahu} (MISC) anunciou que independentemente da aprovação ou não da nova lei pelo \textbf{Knesset} (ORG) , o lançamento e a implementação da nova emissora, por questões técnicas e comerciais, teria que ser adiada para 2018, o que levou a severas críticas por parte de diversos setores da sociedade israelita. O modelo de funcionamento escolhido pelo governo israelita também sofreu severas críticas por parte da \textbf{União Europeia de Radiodifusão} (ORG) (\textbf{EBU} (MISC)), que não aprovou a forma de financiamento da nova emissora (o orçamento disponível para a emissora estaria vinculado à legislação orçamental do país) e também e a ausência planeada de serviços noticiosos. As regras da \textbf{EBU} (MISC) determinam que uma emissora pública que se torne membro da associação deverá ter uma percentagem diária da sua programação relacionada com a transmissão de programas jornalísticos. Mesmo com essas críticas, o governo local sofreu diversas pressões por parte de membros do \textbf{Knesset} (ORG) e foi forçado a antecipar o encerramento da \textbf{IBA} (MISC) para abril de 2017. Porém, o festival daquele ano já tinha suas datas marcadas e, com diversos compromissos comerciais relativos ao evento, o encerramento do canal foi marcado para a uma da manhã local e as atividades do canal foram encerradas imediatamente após o encerramento da transmissão do certamente (esta foi a última transmissão e produção da emissora). Alguns dias antes, a 9 de maio, a equipa da emissora foi informada do fim do canal, faltando apenas duas horas antes da transmissão do último programa. A programação parcial no \textbf{Canal 1} (LOC) foi retomada no dia seguinte, sem programação inédita, enquanto o \textbf{Canal} (ORG) 33 foi encerrado com apenas um slide explicando o encerramento em árabe. Todas as estações de rádio da \textbf{IBA} (MISC) continuaram a funcionar normalmente até ao dia 15 de maio de 2017.

Para a garantia da participação de \textbf{Israel} (LOC), uma pequena equipa composta por vinte pessoas permaneceu a trabalhar para assegurar a transmissão da segunda semifinal do certame, transmitida a 11 de maio de 2017, com a grande final também transmitida dois dias depois. O programa foi emitido no \textbf{Canal 1} (LOC), sem comentários e com a apresentação de legendas durante o programa para fins da votação. A estação de rádio 88FM da \textbf{IBA} (MISC) também transmitiu o programa, com comentários em direto de \textbf{Kobi Menora} (PER), \textbf{Dori Ben Ze'ev} (PER) e \textbf{Alon Amir} (PER). Durante o anúncio dos votos do júri de \textbf{Israel} (LOC), o porta-voz do júri israelita, o jornalista \textbf{Ofer Nachshon} (PER), foi o responsável pela despedida oficial do canal e alguns minutos depois do final do espectáculo o sinal foi desligado.

A 14 de Maio de 2018, \textbf{Yaakov Litzman} (PER) , líder do ultra-ortodoxo partido \textbf{United Torah Judaism} (MISC) e ministro da \textbf{Saúde} (LOC) israelita enviou uma carta aos ministros do turismo, comunicação e cultura e desporto no qual ele pediu que o evento não viole as regras religiosas. Ele escreve :

De acordo com a lei religiosa judaica, o \textbf{Shabat} (LOC) é observado a partir de sexta à noite antes do pôr do sol até sábado à noite. Se a transmissão da final, que começa às 22h (hora local), não for afetada pelo \textbf{Shabbat} (LOC), os ensaios de sexta à noite e todo dia de sábado o serão. O presidente do comité da Eurovisão da UER, \textbf{Frank-Dieter Freiling} (ORG), indicou estar plenamente consciente das tensões e ter planos para as remediar durante as suas comunicações com o canal israelita.

Em setembro de 2018, 140 artistas de renome internacional assinam uma tribuna apelando ao boicote da edição de 2019 do \textbf{Festival} (LOC). Nas assinaturas deste apelo, lançado em apoio aos artistas palestinianos, foram notados os nomes de \textbf{Roger Waters} (PER) , \textbf{Ken Loach} (PER) , \textbf{Mike Leigh} (PER) , \textbf{Aki Kaurismäki} (PER) , \textbf{Yann Martel} (PER) , \textbf{Alia Shawkat} (PER) , \textbf{Jacques Tati} (PER) , \textbf{Alain Guichardet} (PER) e \textbf{Elli Medeiros} (PER) . De acordo com as assinaturas, \textbf{"} (PER)Enquanto os palestinianos não puderem aproveitar a liberdade, a justiça e a igualdade de direitos, não poderá haver negócios com o \textbf{Estado} (LOC) que os recusa os seus direitos fundamentais."

Durante a final da seleção nacional, foi anunciado que a emissora havia se reservado ao direito de mudar a decisão tomada pelo júri e pelo público ucraniano. Após a vitória de \textbf{Maruv} (PER), foi relatado que a emissora havia enviado um contrato à sua agência, exigindo que \textbf{Maruv} (LOC) cancelasse todas as aparições e atuações na \textbf{Rússia} (LOC) para se tornar a representante ucraniana. À artista foram também dadas 48 horas para assinar o contrato, caso contrário seria substituída. No dia seguinte, \textbf{Maruv} (PER) revelou que o contrato da emissora também a proibiu de improvisar no palco e de se comunicar com qualquer jornalista sem a permissão da emissora, exigindo que ela cumprisse integralmente com os pedidos da emissora. Mais tarde, a emissora publicou um comunicado explicando cada linha do contrato. Se a artista não seguisse nenhuma dessas cláusulas, seria multada em  2 milhões. \textbf{Maruv} (PER) também afirmou que a emissora não lhe daria nenhuma compensação financeira pela competição e não pagaria pela viagem a \textbf{Tel Aviv} (LOC).

A 25 de fevereiro de 2019, tanto \textbf{Maruv} (LOC) como a emissora confirmaram que ela não representaria a \textbf{Ucrânia} (LOC) em \textbf{Israel} (LOC) devido a disputas dentro do contrato, e que outro representante seria escolhido. A finalista da \textbf{National Jazz} (ORG), a \textbf{Freedom Jazz} (ORG), anunciou a 26 de fevereiro de 2019 que havia rejeitado a proposta da emissora de representar a \textbf{Ucrânia} (LOC). O terceiro colocado \textbf{Kazka} (PER) confirmou que havia também rejeitado a oferta no dia seguinte.

O incidente reuniu a cobertura mediática de importantes emissoras internacionais, como o \textbf{The New York Times} (ORG) , o \textbf{Washington Post} (ORG) , o \textbf{Billboard} (ORG) , o \textbf{The Independent} (MISC) , o \textbf{SBS News} (MISC) , o \textbf{Irish Independent} (MISC) , o \textbf{Le Figaro} (ORG) , o \textbf{Cosmopolitan} (MISC) , e o \textbf{ABC} (ORG) .

A 27 de fevereiro de 2019, foi anunciado que a \textbf{Ucrânia} (LOC) não participaria no Festival Eurovisão da Canção 2019.

Os preços dos bilhetes para o evento deste ano provocou criticas, tanto em \textbf{Israel} (LOC) como no exterior, com \textbf{The Times of Israel} (ORG) dizendo que são "provavelmente os mais caros de sempre da \textbf{Eurovisão"} (MISC). As razões para a atribuição dos preços altos deve-se ao alto custo de vida em \textbf{Israel} (LOC) e ao facto do governo de \textbf{Israel} (LOC) não estar a financiar a produção da \textbf{Eurovisão} (LOC) deste ano. Apesar do local poder acolher 10.000 pessoas, apenas 7.300 lugares estão disponíveis devido ao tamanho do palco e ao equipamento técnico, bem como as medidas de segurança. Dos 7.300 lugares, 3.000 foram reservados pela \textbf{UER} (ORG) , ficando apenas 4.300 disponíveis para os fãs .

A 3 de março de 2019, a venda dos bilhetes tinha sido congelada devido a irregularidades notificadas pelo \textbf{Comité de Supervisão} (ORG) da emissora israelita \textbf{KAN} (ORG). A comunicação social israelita em língua hebraica relatou que os bilhetes foram revendidos ilegalmente por 2,5 vezes o preço original. O ministro da Segurança Pública, \textbf{Gilad Erdan} (PER), instruiu uma investigação .

\textbf{Kan} (PER) sofreu um ataque cibernético por um grupo de hackers que afetou a vida de acessibilidade da emissora na primeira semifinal. Os hackers puderam mostrar brevemente declarações anti-israelenses sobre os fluxos como \textbf{"Israel} (PER) não é seguro, você verá" e \textbf{"Risco} (PER) de ataque com mísseis, por favor, se abrigue". O incidente está sendo investigado pela emissora e pela \textbf{European Broadcasting Union} (ORG) (\textbf{EBU} (MISC)). \textbf{Kan} (PER) divulgou uma declaração sobre o incidente dizendo: "O problema foi resolvido rapidamente, e parece que durante as primeiras semifinais um site foi hackeado aqui por alguns minutos, e acreditamos que as mensagens não foram vistas por muitas pessoas."

Várias emissoras em toda a \textbf{Europa} (LOC) relataram vários problemas durante a transmissão ao vivo da primeira semifinal. As questões variavam de uma perda de comentários de \textbf{Tel Aviv} (LOC) para a \textbf{Holanda} (LOC) e a \textbf{Macedônia do Norte} (LOC). A emissora polonesa, \textbf{TVP} (MISC), teve que substituir seu comentarista regular \textbf{Artur Orzech} (PER), que estava baseado em \textbf{Tel Aviv} (LOC) com outra pessoa que estava baseada em \textbf{Varsóvia} (LOC), porque os telespectadores não puderam ouvir \textbf{Orzech} (LOC). A \textbf{Alemanha} (LOC) e o \textbf{Reino Unido} (LOC) perderam uma parte do show. No \textbf{Reino Unido} (LOC), o programa recomeçou quando a recapitulação das eliminatórias da primeira semifinal começou a ser jogada, substituída por uma mensagem "Lamentamos a quebra deste programa e estamos tentando corrigir a falta\textbf{"} (MISC) . enquanto em \textbf{França} (LOC) a emissora \textbf{France Televisions} (ORG) teve problemas de áudio durante as apresentações portuguesas e belgas .

Em 18 de maio de 2019, a \textbf{União Européia de Radiodifusão} (ORG) (\textbf{EBU} (MISC)) confirmou que o júri bielorrusso foi demitido da \textbf{Grande Final} (LOC) e seus votos não contam para o resultado final. Esta decisão vem depois que os membros do júri revelaram durante uma entrevista quais países votaram durante a primeira semifinal. Isso vai contra as regras do concurso, que afirmam que os resultados da semifinal não serão revelados antes da \textbf{Grande Final} (LOC). A \textbf{UER} (ORG) também declarou que os votos da \textbf{Bielorrússia} (LOC) serão alocados com base em um resultado agregado aprovado pelos auditores.

A holding da \textbf{Eurovisão} (LOC) em \textbf{Israel} (LOC) havia enfrentado protestos devido à ocupação em curso do \textbf{Estado da Palestina} (LOC) por \textbf{Israel} (LOC). Duas manifestações \textbf{pró-Palestina} (MISC) notáveis ocorreram durante a final: durante as apresentações de \textbf{Madonna} (PER) como "\textbf{Like a Prayer} (MISC)" e "\textbf{Future} (MISC)", o cantor dirigiu um monólogo para dançarinos usando máscaras de gás entre as duas canções, aludindo à "tempestade" dentro de nós ", e afirmando que" eles acham que não estamos cientes de seus crimes. Nós sabemos, mas não estamos prontos para agir "- interpretado como referência ao conflito. Durante "\textbf{Future} (MISC)", duas dançarinas usando bandeiras israelenses e palestinas nas costas de suas fantasias foram vistas quando o vocalista convidado \textbf{Quavo} (PER) cantou a letra "Nem todo mundo está vindo para o futuro, nem todo mundo está aprendendo com o passado".

Durante o desvelamento dos resultados, membros da entrada islandesa \textbf{Hatari} (PER) também foram vistos apresentando uma bandeira contendo a bandeira da \textbf{Palestina} (LOC). Houve anteriormente preocupações de que o auto-descrito grupo anticapitalista usasse seu desempenho para protestar contra a ocupação israelense da \textbf{Palestina} (LOC), e a banda já havia recebido avisos da \textbf{EBU} (MISC) sobre as declarações que haviam feito antes da competição. Após o concurso, a \textbf{EBU} (MISC) afirmou que "as conseqüências desta ação serão discutidas pelo \textbf{Grupo de Referência} (ORG) (o conselho executivo do \textbf{Concurso} (MISC)) após o \textbf{Concurso} (MISC)".

A \textbf{Campanha Palestina} (MISC) pelo \textbf{Boicote Acadêmico e Cultural de Israel} (LOC) pediu um completo boicote à \textbf{Eurovisão} (LOC) em \textbf{Israel} (LOC). Em resposta às ações de \textbf{Hatari} (PER), a organização declarou que "artistas que insistem em cruzar a linha de piquete palestino, jogando em \textbf{Tel Aviv} (LOC) desafiando nossos apelos, não podem compensar o mal que causam à nossa luta pelos direitos humanos" equilibrando "seu ato cúmplice com algum projeto com os palestinos. A sociedade civil palestina rejeita esmagadoramente esta folhagem de folhas de figo. " \textbf{Hatari} (PER) posteriormente anunciou uma colaboração com o artista palestino \textbf{Bashar Murad} (PER) para seu próximo single.



Outros prémios

Para além do troféu principal do vencedor, o \textbf{Prémio Marcel Bezençon} (MISC) e o \textbf{Prémio Barbara Dex} (MISC) foram contestados durante o \textbf{Festival Eurovisão da Canção 2019} (MISC). Além disso, a votação da \textbf{OGAE} (MISC) teve lugar antes do concurso.

O \textbf{Prêmio Marcel Bezençon} (MISC) foi entregue pela primeira vez durante o \textbf{Festival Eurovisão da Canção 2002} (MISC) em \textbf{Tallinn} (LOC), \textbf{Estônia} (LOC), homenageando as melhores músicas concorrentes na final. Fundada por \textbf{Christer Björkman} (PER) (representante da \textbf{Suécia} (LOC) no \textbf{Eurovision Song Contest} (ORG) 1992 e atual chefe de delegação da \textbf{Suécia} (LOC)) e \textbf{Richard Herrey} (PER) (membro do \textbf{Herreys} (LOC) e do \textbf{Eurovision Song Contest} (ORG) 1984 vencedor da \textbf{Suécia} (LOC)), os prêmios receberam o nome do criador do a competição anual \textbf{Marcel Bezençon} (PER). Os prêmios são divididos em três categorias: \textbf{Prêmio de Imprensa} (MISC), \textbf{Prêmio Artístico} (MISC) e \textbf{Prêmio Compositor} (MISC). Os vencedores são revelados pouco antes da final da \textbf{Eurovisão} (LOC).

A \textbf{Organização Geral dos Amadores da Eurovisão} (ORG) (\textbf{OGAE} (MISC)) é um dos maiores clubes de fãs da \textbf{Eurovisão} (LOC). Ela tem secções em toda a \textbf{Europa} (LOC) e foi criada em 1984 na \textbf{Finlândia} (LOC) . Todos os países que participam ou participaram na competição anual podem ter a sua própria secção \textbf{OGAE} (MISC). Os outros paises do mundo são, quanto a eles, agrupados na \textbf{"OGAE Resto do Mundo} (MISC)", criado em 2004. Cada ano, os clubes têm a oportunidade de votar, seguindo o modelo da \textbf{Eurovisão} (LOC), para os seus favoritos.

A tabela seguinte mostra os dez primeiros do voto \textbf{OGAE} (MISC) 2019, os resultados finais foram anunciados a 30 de abril de 2019.



Álbum oficial

\textbf{Eurovision Song Contest} (ORG): \textbf{Tel Aviv} (LOC) 2019 é o álbum de compilação oficial do festival, lançado pela \textbf{União Europeia de Radiodifusão} (ORG) e distribuído pela \textbf{Universal Music Group} (ORG) digitalmente a 12 de Abril de 2019. O álbum conta com todas as 41 participações, incluindo os semifinalistas que não se classificaram para a final.



Festival Eurovisão da Canção 2020



Países que planejaram participar em 2020 e se classificariam automaticamente para a final \textbf{Países} (MISC) que planejavam participar em 2020 e teriam competido nas semifinais \textbf{Países} (LOC) que participaram no passado, mas não planeavam participar em 2020

O \textbf{Festival Eurovisão da Canção 2020} (MISC) (em inglês : \textbf{Eurovision Song Contest} (ORG) 2020 ; em francês : \textbf{Concours Eurovision} (MISC) de la chanson 2020 ; em neerlandês : \textbf{Eurovisiesongfestival 2020} (MISC) ) seria a 65. edição anual do evento , que iria acontecer na arena \textbf{Ahoy} (PER) , em \textbf{Roterdão} (LOC) , nos \textbf{Países Baixos} (LOC) , entre 12 e 16 de maio de 2020. Esta seria a quinta vez na história do certame que o evento teria sido realizado naquele país. A última vez que a versão adulta do festival foi realizada nos \textbf{Países Baixos} (LOC) foi em 1980 . Além disso, teria sido o primeiro evento a ser realizado sob a bandeira da rede \textbf{Eurovision} (MISC) no país desde a edição de 2012 da versão infantil do festival . As semifinais seriam realizadas nos dias 12 e 14 de maio e a final no dia 16 de maio.

A 18 de março de 2020 , foi anunciado que a \textbf{UER} (ORG) tinha cancelado o evento devido à pandemia de \textbf{COVID-19} (LOC) .



Localização

O \textbf{Festival} (PER) estava previsto para acontecer nos \textbf{Países Baixos} (LOC), que venceram a edição anterior . Tal como o protocolo da \textbf{UER} (ORG) orienta, os preparativos para esta edição começaram alguns minutos após a vitória de \textbf{Duncan Laurence} (PER) em \textbf{Tel Aviv} (LOC) , quando \textbf{Jon Ola Sand} (PER) , o supervisor executivo do \textbf{Festival} (LOC), entregou uma pilha de papéis e uma "pen drive" aos executivos neerlandeses, todos relacionados com a produção e a execução da próxima edição. Este festival será uma coprodução entre as três emissoras públicas do país e cada uma deverá assumir uma área da produção: a \textbf{AVROTROS} (ORG) e os seus dois canais irmãos a \textbf{Nederlandse Publieke Omroep} (PER) (\textbf{NPO} (MISC)) e a \textbf{Nederlandse Omroep Stichting} (ORG) (\textbf{NOS} (ORG)), cada uma assumindo uma área de produção.

Seguindo os critérios usados nas edições anteriores, as cidades interessadas deverão cumprir as seguintes exigências:

Mesmo antes do início oficial da fase de seleção da cidade anfitriã, várias cidades declararam o seu interesse. De acordo com o primeiro-ministro dos \textbf{Países Baixos} (LOC) , \textbf{Mark Rutte} (PER) poucas horas após a vitória de \textbf{Duncan} (PER), várias autoridades locais já tinham anunciado as suas intenções de sediar o evento.

Já a 19 de maio de 2019, um dia após a vitória de \textbf{Duncan Laurence} (PER), várias cidades declararam que estavam interessadas em sediar a competição. A capital, \textbf{Amesterdão} (LOC), já tinha indicado para a receber em abril de 2019 se os \textbf{Países Baixos} (LOC) vencessem a edição de 2019. A cidade portuária de \textbf{Roterdão} (LOC) também se ofereceu para sediar a competição. Tanto \textbf{Amesterdão} (LOC) quanto \textbf{Roterdão} (LOC) deram um passo em frente em relação as outras candidatas, por terem uma infraestrutura de telecomunicações, transporte e acessibilidade melhores do que as outras e também por terem experiência gigantesca de eventos de grande dimensão. Mesmo já com a existência de duas pré-favoritas, as cidades de \textbf{Arnhem} (LOC) , \textbf{Haia} (LOC) , \textbf{Leuvarda} (LOC) , \textbf{Maastricht} (LOC) , \textbf{Utrecht} (LOC) e \textbf{Zwolle} (LOC) também manifestaram interesse no dia 19 de maio. Por outro lado, a província de \textbf{Brabante do Norte} (LOC) expressou no mesmo dia estar interessada sem indicar um local específico. A 20 de maio de 2019, a cidade de \textbf{Den Bosch} (ORG) , localizada nesta província, também apresentou uma proposta.

Em 20 de maio de 2019, o diretor comercial do \textbf{Aeroporto de Enschede-Twente} (LOC) anunciou que estava interessado em sediar a competição num dos hangares do aeroporto. Ele também afirmou que, se o local preenchesse todos os critérios da emissora, uma candidatura oficial seria confirmada posteriormente.

A 24 de maio de 2019, o diretor do \textbf{IJsselhallen} (MISC), uma arena que tinha sido proposta para sediar a competição em \textbf{Zwolle} (LOC), anunciou que não apresentaria uma candidatura oficial, uma vez que o espaço não preenchia os critérios exigidos.

A 25 de maio de 2019, a cidade de \textbf{Breda} (LOC) , no norte de \textbf{Brabante} (LOC), propõe a sua candidatura em nome da província.

A 29 de maio de 2019, a \textbf{NPO} (MISC) revelou à imprensa holandesa que o processo oficial de candidaturas para a escolha da cidade-sede do Festival Eurovisão 2020 estava aberto e que as interessadas teriam por volta de um mês para elaborar as suas propostas até final de julho. O anúncio da sede foi realizado em 30 de agosto de 2019, com a escolha de \textbf{Roterdão} (LOC).

Inicialmente, \textbf{Zwolle} (LOC) também considerou a possibilidade de apresentar uma candidatura. Mas após a divulgação dos critérios exigidos para o processo de candidatura, o \textbf{IJsselhallen} (MISC), que seria o local proposto, não estava dentro das exigências da \textbf{EBU} (MISC) e a cidade acabou retirando a sua proposta. \textbf{Haia} (LOC) também excluiu a inscrição do \textbf{World Forum} (ORG), que sediou a competição em 1976 e 1980 , por causa da baixa capacidade de público do centro de convenções. Empresários da cidade de \textbf{Enschede} (LOC) também manifestaram interesse em sediar o evento dentro de um hangar desocupado do \textbf{Aeroporto da cidade} (LOC) , mas o conselho executivo da cidade disse que não apoiava a ideia.



\textbf{Formato} (PER)O slogan desta edição, \textbf{" Open Up} (MISC) ", foi revelado no dia 24 de outubro de 2019. A logo oficial e o \textbf{branding} (ORG) foram revelados no dia 28 de novembro de 2019. Desenhado por CLEVER\textbf{FRANKE} (ORG), é "uma representação abstrata das cores das bandeiras dos 41 países participantes da edição de 2020 ordenados conforme sua primeira participação no festival\textbf{"} (MISC).

A \textbf{EBU} (MISC) revelou o design do palco para o Festival Eurovisão da Canção 2020 em dezembro de 2019. O design é inspirado no slogan \textbf{" Open} (MISC) Up " e na típica paisagem plana neerlandesa. O palco do festival foi projetado pelo cenógrafo alemão \textbf{Florian Wieder} (PER), que também projetou os palcos das edições do \textbf{Eurovisão de 2011} (MISC), 2012, 2015 e de 2017 a 2019. Ao contrário da edição anterior, o \textbf{Green Room} (ORG) foi colocado dentro do próprio local do evento.

O festival terá três apresentadores: \textbf{Chantal Janzen} (PER) , atriz e apresentadora de televisão; \textbf{Jan Smit} (PER) , cantor e comentarista do festival; e \textbf{Edsilia Rombley} (PER) , cantora que representou os \textbf{Países Baixos} (LOC) nas edições do \textbf{Eurovisão de 1998} (MISC) e de 2007 . A vlogger de beleza \textbf{Nikkie de Jager} (PER) (\textbf{NikkieTutorials} (MISC)) será apresentadora do conteúdo online do festival, incluindo uma série no \textbf{Youtube} (MISC) "por trás das câmeras\textbf{"} (MISC), gravada com os participantes do festival. Ela também fará a cobertura do \textbf{tapete vermelho} (ORG) durante a cerimônia de abertura, além de fazer uma aparição em todas as três apresentações, nos dias 12, 14 e 16 de maio de 2020. \textbf{Roos Moggré} (LOC) e \textbf{Andrew Makkinga} (PER) apresentarão a conferência de imprensa do festival.

O sorteio para dividir os países participantes de cada semi-final ocorreu no dia 28 de janeiro de 2020, no \textbf{Rotterdam's City Hall} (ORG). Os trinta e cinco semi-finalistas foram divididos em cinco potes, com base nos padrões de votação históricos calculados pelo parceiro oficial de \textbf{televoto do Festival} (LOC), \textbf{Digame} (LOC). O sorteio com diferentes potes ajudou a reduzir as chances do chamado \textbf{" bloc voting} (MISC) " e aumenta o suspense da semi-final. O sorteio também determinou qual das semi-finais cada um dos seis países de qualificação automática - o \textbf{Big Five} (ORG) mais os \textbf{Países Baixos} (LOC) - transmitirá e participará da votação. A cerimônia foi apresentada por \textbf{Chantal Janzen} (PER), \textbf{Jan Smit} (PER) e \textbf{Edsilia Rombley} (PER), e contou com a passagem da insígnia de anfitrião de \textbf{Zippi Brand Frank} (PER), vice-prefeito de \textbf{Tel Aviv} (LOC) (cidade anfitrião da última edição do festival), para \textbf{Ahmed Aboutaleb} (PER), prefeito de \textbf{Rotterdam} (LOC) .

O chefe da delegação espanhola revelou no dia 22 de outubro de 2019, que a \textbf{EBU} (MISC) consultará as delegações sobre potenciais mudanças no sistema de votação. O chefe da delegação grega revelou no dia 30 de outubro de 2019 que a maioria das delegações (80\%) votou a favor de manter o sistema de votação atual.

O conceito de postcards para 2020 continua com o tema do evento, \textbf{" Open Up "} (MISC). Cada artista vai visitar uma parte diferente dos \textbf{Países Baixos} (LOC) , se conectando com a população local e participando de atividades, tradições ou hobbies típicos neerlandeses.

A segunda semi-final será aberta por uma performance do breakdancer \textbf{Redo} (PER). O ato de intervalo da grande final será uma apresentação de alguns dos mais importantes vencedores anteriores do festival.Estão confirmadas as participações de \textbf{Gigliola Cinquetti} (PER) , com \textbf{"Non ho l'età} (PER)", \textbf{Lenny Kuhr} (PER) , com "De troubadour", \textbf{Getty Kaspers} (PER) (representando o grupo \textbf{Teach-In} (ORG) ), com "\textbf{Ding-a-dong} (PER)", \textbf{Sandra Kim} (PER) , com "\textbf{J'aime la vie} (MISC)", \textbf{Paul Harrington} (PER) e \textbf{Charlie McGettigan} (PER) , com "\textbf{Rock 'n' Roll Kids} (MISC)", \textbf{Alexander Rybak} (PER) , com "\textbf{Fairytale} (PER)" e \textbf{Duncan Laurence} (PER) , com "\textbf{Arcade} (LOC)".

Uma performance conjunta de uma orquestra sinfônica composta por jovens músicos com o DJ \textbf{Peter Gabriel} (PER) abrirá a final e produzirá a música que irá acompanhar a entrada dos 26 finalistas na \textbf{arena} (PER).



\textbf{Países Participantes} (LOC)A \textbf{EBU} (MISC) anunciou no dia 13 de novembro de 2019 que quarenta e um países vão participar do festival, com \textbf{Bulgária} (LOC) e \textbf{Ucrânia} (LOC) retornando após suas ausências na edição de 2019 do festival, e com \textbf{Hungria} (LOC) e \textbf{Montenegro} (LOC) se retirando.

As 3 integrantes do grupo sueco-americano \textbf{The Mammas} (MISC) , foram as backvocals de \textbf{John Lundvik} (PER) quando ele representou a \textbf{Suécia} (LOC) no ano anterior em \textbf{Tel Aviv} (LOC) Isso também aconteceu com \textbf{Vasil} (LOC) que também foi um dos backvocals de sua compatriota \textbf{Tamara Todevska} (PER) e com a vencedora do \textbf{Festival Eurovisão da Canção Júnior 2015} (MISC) ,a maltesa \textbf{Destiny Chukunyere} (MISC) ,que fez parte da equipe de palco que que acompanhou a sua compatriota \textbf{Michaela Pace} (PER) também no ano anterior. Um outro participante que também fez backvocal em edições anteriores é o austríaco \textbf{Vincent Bueno} (PER) , que acompanhou o seu compatriota \textbf{Nathan Trent} (PER) , em \textbf{Kiev} (LOC) 2017 . Uma situação interessante é a da local \textbf{Stefania Liberakakis} (PER) ,filha de pais gregos,que representou o seu país na versão júnior do festival em 2016 ,como uma das 3 integrantes do grupo \textbf{Kisses} (LOC) .Porém,quatro anos,mais tarde,ela irá representar o seu país de origem,\textbf{a Grécia} (LOC). Também retornará ao festival,a moldova \textbf{Natalia Gordienko} (PER) que agora estará como artista solo.\textbf{Gordienko} (PER) foi a representante de seu país em \textbf{Atenas} (LOC) 2006 ao lado de dois outros artistas: \textbf{Arsenium} (MISC) e Connect-R . A sérvia \textbf{Sanja Vučić} (PER) foi artista solo em \textbf{Estocolmo 2016} (MISC) e retorna ao festival como integrante do grupo \textbf{Huricane} (LOC),ao lado dela estará a montenegrina \textbf{Ksenija Knežević} (PER),que foi backvocal de seu pai, \textbf{Knez} (PER) ,em no ano anterior em \textbf{Viena} (LOC) . A única artista que retornará como solo é a ítalo-eriteia \textbf{Senhit} (LOC) ,que também representou \textbf{San Marino} (LOC) em \textbf{Düsseldorf} (LOC) 2011 .

A primeira semi-final estava marcada para a noite do dia dia 12 de maio de 2020,começando as 21:00 (\textbf{CEST} (ORG)). Dezessete países participariam da primeira semi-final. Esses países, mais \textbf{Alemanha} (LOC) , \textbf{Itália} (LOC) e \textbf{Países Baixos} (LOC) , votariam nessa semi-final.

A primeira semi-final está marcada para a noite do dia dia 12 de maio de 2020,começando as 21:00 (\textbf{CEST} (ORG)).. Dezoito países participarão da segunda semi-final. Esses países, mais \textbf{França} (LOC) , \textbf{Espanha} (LOC) e \textbf{Reino Unido} (LOC) , votarão nesta semi-final.



Outros países

A elegibilidade para uma potencial participação no Festival Eurovisão da Canção requer uma emissora nacional com participação ativa na \textbf{UER} (ORG), que poderá transmitir o \textbf{Festival} (LOC) através da rede \textbf{Eurovisão} (LOC) . A \textbf{UER} (ORG) irá convidar a participação no festival para todos os cinquenta e seis membros ativos. Em contraste com os anos anteriores, o membro associado \textbf{Austrália} (LOC) não precisará de um convite para o concurso de 2020, pois recebeu permissão para participar até 2023.



Ataque terrorista frustrado

\textbf{Em Fevereiro de 2020} (MISC), as autoridades luxemburguesas detiveram \textbf{Alexander H.} (PER), um cidadão sueco, por planear um ataque terrorista no Festival Eurovisão da Canção agendado para \textbf{Maio} (LOC) de 2020 em \textbf{Roterdão} (LOC), \textbf{Países Baixos} (LOC). A conspiração, descoberta pelos serviços de segurança do \textbf{Luxemburgo} (LOC) , envolveu \textbf{Alexander H.} (PER), então com 17 anos, e um cúmplice neerlandês não identificado. Um documento intitulado \textbf{"Tempo de diversão} (MISC) para a \textbf{Eurovisão 2020} (MISC) - Para um futuro melhor e com menos aceitação excessiva", encontrado no computador portátil de \textbf{H.} (PER), delineava planos para envenenar os participantes com cianeto ou ricina , distribuir gás cloro através de sistemas de ventilação ou foguetes personalizados e bloquear saídas de emergência para maximizar o número de vítimas. Os materiais para a produção de cloro e os protótipos de foguetões foram apreendidos na casa de H. em \textbf{Strassen} (LOC), no \textbf{Luxemburgo} (LOC). \textbf{H.} (PER), membro do grupo neonazi \textbf{The Base} (ORG) , foi motivado por uma ideologia extremista que visava símbolos multiculturais, detestando grupos como mulheres, pessoas de cor e homossexuais. As autoridades holandesas e os organizadores do festival desconheciam a conspiração, que foi frustrada antes da execução. A \textbf{Eurovisão de 2020} (MISC) foi finalmente cancelada devido à pandemia de \textbf{COVID-19} (LOC). Os procuradores solicitaram uma pena de 12 anos, com decisão judicial prevista para novembro de 2025.



Festival Eurovisão da Canção 2021



\textbf{Final: Desfile} (MISC) de bandeiras com apresentação dos 26 países finalistas com música de \textbf{DJ Pieter Gabriel Final: "} (PER)\textbf{Rock the Roof} (MISC)" interpretado por \textbf{Lenny Kuhr} (PER) , \textbf{Teach-In} (MISC) com \textbf{Getty Kaspers} (PER), \textbf{Sandra Kim} (PER) , \textbf{Helena Paparizou} (PER) , \textbf{Lordi} (PER) e \textbf{Måns Zelmerlöw "Music Binds Us} (PER)" interpretado por \textbf{Afrojack} (PER) , \textbf{Glennis Grace} (PER) e \textbf{Wulf " Arcade} (PER) " interpretado por \textbf{Duncan Laurence Países} (PER) que classificaram-se para a final \textbf{Países} (MISC) que não se classificaram para a final \textbf{Países} (MISC) que participaram no passado, mas não em 2021

O \textbf{Festival Eurovisão da Canção 2021} (MISC) (em inglês : \textbf{Eurovision Song Contest} (ORG) 2021 ; em francês : \textbf{Concours Eurovision} (MISC) de la chanson 2021 ; em neerlandês : \textbf{Eurovisiesongfestival 2021} (MISC) ) foi a 65. edição do Festival Eurovisão da Canção . O evento decorreu nos dias 18 , 20 e 22 de maio sendo realizado no \textbf{complexo Ahoy de Roterdão} (LOC) nos \textbf{Países Baixos} (LOC) . Após o cancelamento do Festival Eurovisão da Canção 2020 devido à pandemia de \textbf{COVID-19} (LOC) , a \textbf{União Europeia de Radiodifusão} (ORG) e o grupo de referência decidiram como seria a edição de 2021, nesse ínterim a câmara municipal de \textbf{Roterdão} (LOC) decidiu manter a cidade como sede da competição e, desta forma, os \textbf{Países Baixos} (LOC) tornariam-se anfitriões do festival por 5. ocasião na história, sendo a última vez em 1980 .

Os apresentadores foram os mesmos designados para a \textbf{Eurovisão: Europe Shine} (MISC) \textbf{A Light} (MISC) : \textbf{Edsilia Rombley} (PER) , \textbf{Nikkie de Jager} (LOC) , \textbf{Jan Smit} (PER) e \textbf{Chantal Janzen} (PER). O lema da edição foi \textbf{Open Up} (MISC) (em português : Abra-se), que faz referência à abertura dos \textbf{Países Baixos} (LOC) face ao mundo. O logo desta edição foi o mesmo que era suposto ser utilizado em 2020 e representa de maneira abstrata todas as cores das bandeiras dos países participantes.

Nesta edição houve 39 nações participantes. \textbf{Hungria} (LOC) e \textbf{Montenegro} (LOC) decidiram não participar nesta ocasião, assim como também não eram suposto fazê-lo em 2020. A 5 de março de 2021, \textbf{Arménia} (LOC) decidiu retirar-se da competição pelo conflito do \textbf{Alto Carabaque} (LOC) e a 26 de março, a \textbf{Bielorrússia} (LOC) foi desqualificada por desrespeitar a regra de não incluir mensagens políticas nas letras das canções.



Localização

Após o cancelamento do Festival Eurovisão da Canção 2020 , a \textbf{União Europeia de Radiodifusão} (ORG) esteve em negociações com as emissoras neerlandesas \textbf{AVROTROS} (ORG) , \textbf{NPO} (MISC) e \textbf{NOS} (MISC) , bem como com a câmara municipal de \textbf{Roterdão} (LOC) , que organizaria o festival de 2020, sobre a possibilidade de organizar a edição de 2021 na mesma cidade.

Segundo os média dos \textbf{Países Baixos} (LOC), a câmara municipal de \textbf{Roterdão} (LOC) iria precisar de 6.700.000 adicionais para a organização do evento em 2021. O governo da cidade aprovou o orçamento da proposta a 23 de abril de 2020 com 33 votos a favor e 8 contra. Desta forma, mantiveram o \textbf{Ahoy} (PER) \textbf{Rotterdam} (LOC) como sede, complexo com espaço para 16.000 pessoas, que já recebeu o \textbf{Festival Eurovisão da Canção Júnior 2007} (MISC) , o \textbf{Nationaal Songfestival} (LOC) em 2000, 2001 e 2003; e os \textbf{MTV Europe Music Awards} (ORG) em 1997 y 2016. No decorrer da emissão do programa especial \textbf{Eurovisão: Europe Shine} (MISC) \textbf{A Light} (MISC) , os apresentadores confirmaram que a sede iria manter-se a mesma.



\textbf{Organização} (ORG)O edição de 2021 será uma co-produção entre três organizações de rádio e televisão dos \textbf{Países Baixos} (LOC): AVROTROS , \textbf{Nederlandse Omroep Stichting} (PER) (\textbf{NOS} (ORG)) e \textbf{Nederlandse Publieke Omroep} (PER) (\textbf{NPO} (MISC)), sendo que cada uma delas terá um papel diferente. \textbf{Sietse Bakker} (PER) e \textbf{Astrid Dutrénit} (PER) serão produtores executivos enquanto \textbf{Emilie Sickinghe} (PER) e \textbf{Jessica Stam} (PER) serão produtores adjuntos.

Em janeiro de 2020, a \textbf{União Europeia de Radiodifusão} (ORG) anunciou que \textbf{Martin Österdahl} (PER) tornaria-se o supervisor executivo do Festival Eurovisão da Canção após a edição de 2020 , sucedendo \textbf{Jon Ola Sand} (PER) . Antes da nomeação, \textbf{Österdahl} (MISC) havia sido produtor executivo das edições de 2013 e 2016 e havia sido membro do grupo de referência do \textbf{Festival Eurovisão da Canção} (MISC) entre 2012 e 2018.



Países participantes

A \textbf{UER} (ORG) anunciou inicialmente no dia 26 de outubro de 2020 que quarenta e um países iriam participar do concurso, apresentando a mesma formação de países que participaria da edição de 2020, que foi cancelada. \textbf{Bulgária} (LOC) e \textbf{Ucrânia} (LOC) estão retornando ao festival após suas ausências na edição de 2019, enquanto \textbf{Hungria} (LOC) e \textbf{Montenegro} (LOC) confirmaram que não retornarão ao festival, tendo ambas como última participação a edição de 2019.

Em março de 2021, \textbf{Arménia} (LOC) e \textbf{Bielorrússia} (LOC) confirmaram a não participação no concurso; A \textbf{Arménia} (LOC) se retirou devido a crises sociais e políticas após a guerra no \textbf{Alto Carabaque} (LOC) em 2020 , enquanto a \textbf{Bielorrússia} (LOC) foi desqualificada após apresentar uma inscrição em violação às regras, reduzindo assim o número de países participantes para 39.

Após o cancelamento do festival em 2020, as emissoras participantes de 24 países anunciaram que selecionariam internamente os mesmos artistas previamente já selecionados para competir em 2021. Além disso, os cantores que já tinham sido selecionados para as representações da \textbf{Estónia} (LOC) e da \textbf{Lituânia} (LOC) no ano passado também ganharam as suas seleções nacionais em 2021.

Excluindo 2020, o festival contou com três representantes que também atuaram anteriormente como vocalistas do mesmo país, e cinco artistas que participaram de outros eventos do \textbf{Eurovisão} (LOC) ou como coristas para o mesmo ou para outro país.

A primeira semi-final ocorreu no dia 18 de maio de 2021 às 21:00 ( \textbf{CEST} (ORG) ). Dezasseis países participaram na primeira semi-final. Esses países, mais \textbf{Alemanha} (LOC) , \textbf{Itália} (LOC) e \textbf{Países Baixos} (LOC) , votaram nesta semi-final. A \textbf{Bielorrússia} (LOC) foi originalmente alocada para participar na primeira metade da semifinal, mas foi desqualificada do concurso após enviar uma inscrição em violação às regras.

A segunda semi-final ocorreu no dia 20 de maio de 2021 às 21:00 ( \textbf{CEST} (ORG) ). Dezassete países participaram na segunda semi-final. Esses países, mais \textbf{França} (LOC) , \textbf{Espanha} (LOC) e \textbf{Reino Unido} (LOC) , votaram nesta semi-final. A \textbf{Arménia} (LOC) , que inicialmente estava confirmada para participar do festival entrando na segunda parte da segunda semi-final, desistiu de participar devido a crise social e politica resultante da \textbf{Guerra de Nagorno-Karabakh} (MISC) em 2020.

A final ocorreu no dia 22 de maio de 2021, com início às 21:00 (\textbf{CEST} (ORG)). Vinte e seis países participaram da final, composta pelo país anfitrião (\textbf{Países Baixos} (LOC)), os \textbf{Big Five} (ORG) e os dez melhores colocados de cada uma das duas semi-finais. Todos os trinta e nove países participantes votaram na final.



Outros países

A elegibilidade para uma potencial participação no Festival Eurovisão da Canção requer uma emissora nacional com participação ativa na \textbf{UER} (ORG), que poderá transmitir o \textbf{Festival} (LOC) através da rede \textbf{Eurovisão} (LOC) . A \textbf{UER} (ORG) irá convidar a participação no festival para todos os cinquenta e seis membros ativos. Em contraste com os anos anteriores, o membro associado \textbf{Austrália} (LOC) não precisará de um convite para o concurso de 2020, pois recebeu permissão para participar até 2023.



Festival Eurovisão da Canção 2022



Países que classificaram-se para a final \textbf{Países} (MISC) que não se classificaram para a final \textbf{Países} (MISC) que participaram no passado, mas não em 2022

O \textbf{Festival Eurovisão da Canção 2022} (MISC) (em inglês : \textbf{Eurovision Song Contest} (ORG) 2022 ; em francês : \textbf{Concours Eurovision} (MISC) de la chanson 2022 ; em italiano : \textbf{Concorso Eurovisione della Canzone} (PER) ) foi a 66 edição do Festival Eurovisão da Canção . O festival realizou-se em \textbf{Itália} (LOC) , após a vitória do país na edição de 2021 , com a música \textbf{Zitti} (PER) e \textbf{Buoni} (PER), dos \textbf{Måneskin} (LOC) . Foi a terceira vez que \textbf{Itália} (LOC) realizou o festival, depois de \textbf{Nápoles} (LOC) , em 1965 , e \textbf{Roma} (LOC) , em 1991 . O evento decorreu nos dias 10 , 12 e 14 de maio , sendo realizado na arena \textbf{Pala Alpitour} (LOC) , em \textbf{Turim} (LOC) .



Localização

- Localização das cidades pré-selecionadas - Outras cidades em licitação

\textbf{O Festival} (MISC) seria em \textbf{Itália} (LOC), país vencedor da edição anterior . Tal como o protocolo da \textbf{UER} (ORG) orienta, os preparativos para esta edição começaram alguns minutos após a vitória dos \textbf{Måneskin} (PER) em \textbf{Roterdão} (LOC) , quando \textbf{Martin Österdahl} (PER), o supervisor executivo do \textbf{Festival} (LOC), entregou uma pilha de papéis e uma "pen drive" aos executivos italianos, todos relacionados com a produção e a execução da próxima edição.

Mesmo antes do início oficial da fase de seleção da cidade anfitriã, várias cidades declararam o seu interesse em receber o festival. No dia seguinte à vitória italiana, os responsáveis políticos de \textbf{Bolonha} (LOC) , \textbf{Milão} (LOC) , \textbf{Nápoles} (LOC) , \textbf{Pésaro} (LOC) e \textbf{Turim} (LOC) manifestaram publicamente o interesse em sediar o evento.

A 8 de outubro de 2021, foi reportado que \textbf{Turim} (LOC) teria sido a cidade selecionada para acolher o \textbf{Festival} (LOC) em 2022.

As seguintes cidades expressaram o interesse em acolher:

\textbf{†} (LOC) Cidade selecionada ‡ Pré-selecionada Apresentou o livro de lances



Produção

O \textbf{Festival Eurovisão da Canção 2022} (MISC) foi produzido pela emissora italiana \textbf{RAI} (ORG) sendo \textbf{Claudio Fasulo} (PER) e \textbf{Simona Martorelli} (PER) os produtores executivos, \textbf{Cristian Biondani} (PER) e \textbf{Duccio Forzano} (PER) os diretores das três emissões em direto, \textbf{Claudio Santucci} (PER) foi a cabeça do espétaculo, e \textbf{Emanuele Cristofoli} (PER) a diretora artística as atuações da abertura e dos intervalos.

O governo italiano contribuiu com 1,5 milhões de euros como parte do orçamento necessário para acolher o evento, enquanto que o município de \textbf{Turim} (LOC) e o governo regional de \textbf{Piedmont} (LOC) contribuiu com 10 milhões de euros no total. O orçamento inicial do evento foi de 16,3 milhões de euros.

O tema artistico e o slogan para o \textbf{Festival} (LOC), "\textbf{The Sound of Beauty} (MISC)" (em português : "O som da beleza"), foi revelado a 21 de janeiro de 2022. Criado pelo estudio \textbf{Flopicco} (PER) situado em \textbf{Roma} (LOC), a arte foi construida à volta de estruturas simétricas e padrões de cimáticas para transmitir as propriedades do som, mas também refletir a forma dos jardins italianos, jà a tipográfica foi inspirada dos posters italianos do inicio do século XX. As cores foram retiradas das da bandeira da \textbf{Itália} (LOC) .

O apresentador de televisão \textbf{Alessandro Cattelan} (PER) e os cantores \textbf{Laura Pausini} (PER) e \textbf{Mika} (PER) foram os apresentadores do \textbf{Festival de 2022} (MISC). Foram também nomeados como anfitriões pela agência de noticias \textbf{Adnkronos} (LOC) e pela revista de televisão \textbf{TV Sorrisi} (MISC) e \textbf{Canzoni} (PER) sendo oficialmente confirmados durante a noite do \textbf{Sanremo Music Festival} (MISC) 2022 a 2 de fevereiro, depois de terem aparecido nessa emissão como convidados especiais.

A "\textbf{Turquoise Carpet} (MISC)" (em português : "\textbf{Passadeira} (LOC) turquesa") e a cerimonia de abertura foram apresentados por \textbf{Gabriele Corsi} (PER), \textbf{Carolina Di Domenico} (PER), \textbf{Mario Acampa} (MISC) e \textbf{Laura Carusino} (PER). Os três apresentados moderaram ainda a conferência de impressa do \textbf{Festival} (LOC).

O palco do \textbf{Festival de 2022} (PER) foi revelado a 18 de fevereiro de 2022. Criado pela designer italiana \textbf{Francesca Montinaro} (PER) e nomeado "\textbf{The Sun Within} (MISC)" (em português : "O sol dentro"), a estrutura do palco foi inspirada pelos movimentos da luz de um sol cinético, com a capacidade pretendida de mostrar o movimento teatral. O design também apresentava uma cascata de água em funcionamento e uma green room modelada num jardim italiano.

A 30 de abril de 2022, a \textbf{UER} (ORG) informou sobre as atuações de abertura e intervalos.

A primeira semifinal foi aberta com um espetáculo de ingenuidade e criatividade italiana, acompanhado pelo hino oficial do concurso, "\textbf{The Sound of Beauty} (MISC)", interpretado por \textbf{Sherol Dos Santos} (PER), enquanto o intervalo contou com um medley de "\textbf{Horizon in Your Eyes} (MISC)", "\textbf{Satisfaction} (MISC)" e "\textbf{Golden Nights} (MISC)" interpretadas por \textbf{Dardust} (PER), \textbf{Benny Benassi} (PER) e \textbf{Sophie and the Giants} (MISC) com a maestrina \textbf{Sylvia Catasta} (ORG), uma breve homenagem a \textbf{Raffaella Carrà} (PER) realizada pelos apresentadores do \textbf{Festival} (LOC), e \textbf{Diodato} (PER) interpretando "Fai rumore".

A segunda semifinal foi aberta por "\textbf{The Italian Way} (MISC)", um ato construído em volta da improvisação italiana realizada pelo co-apresentador \textbf{Alessandro Cattelan} (PER), enquanto o intervalo contou com um medley de "\textbf{Fragile} (ORG)" e "\textbf{People Have the Power} (MISC)" realizado pelos co-apresentadores \textbf{Laura Pausini} (PER) e \textbf{Mika} (ORG) , e \textbf{Il Volo} (PER) interpretando uma nova versão de "Grande amore".

A final foi aberta com o tradicional desfile de bandeiras, com apresentação de todos os vinte e cinco finalistas, acompanhados pelo \textbf{Rockin} (PER)' 1000 interpretando "\textbf{Give Peace} (LOC) a \textbf{Chance} (MISC)" e a co-apresentadora \textbf{Laura Pausini} (PER) interpretando um medley de \textbf{" Benvenuto} (LOC) ", " \textbf{Io} (LOC) canto ", \textbf{" La solitudine} (MISC) ", " \textbf{Le cose che vivi} (MISC) " e "\textbf{Scatola} (MISC)". Os atos deointervalo incluíram \textbf{Måneskin} (PER) com apresentação do seu novo single "\textbf{Supermodel} (PER)" e "\textbf{If I} (MISC) Can Dream", \textbf{Gigliola Cinquetti} (PER) a cantar a sua música vencedora "\textbf{Non ho l'età} (MISC)", e a co-apresentadora \textbf{Mika} (ORG) apresentando um medley de " \textbf{Love Today} (MISC) ", " \textbf{Grace Kelly} (PER) ", o seu novo single "\textbf{Yo Yo} (MISC)" e "\textbf{Happy Ending} (MISC)". A astronauta italiana \textbf{Samantha Cristoforetti} (PER) também apareceunuma mensagem pré-gravada da \textbf{Estação Espacial Internacional} (MISC) .



\textbf{Participantes Nesta} (MISC) edição, fizeram parte 40 países, com o regresso da \textbf{Arménia} (LOC) e do \textbf{Montenegro} (LOC) e a suspensão da \textbf{Rússia} (LOC) .



Canções participantes

A primeira semifinal aconteceu a 10 de maio de 2022 . Dezassete países participaram na primeira semifinal. Esses países mais \textbf{França} (LOC) e \textbf{Itália} (LOC) votaram nesta semifinal. Sombreados, encontram-se os países que se qualificaram para a final.

A segunda semifinal aconteceu em 12 de maio de 2022 às 21:00 (\textbf{CEST} (ORG)). Dezoito países participaram na segunda semifinal que, junto com a \textbf{Alemanha} (LOC) , \textbf{Espanha} (LOC) e \textbf{Reino Unido} (LOC) , votaram nesta semifinal. Sombreados, encontram-se os países que se qualificaram para a final.

A final aconteceu a 14 de maio de 2022 . Vinte e cinco países participaram na final, compondo os "\textbf{Big Five} (MISC)" (entre os quais o país anfitrião \textbf{Itália} (LOC)) e dez das entradas mais bem classificadas de cada uma das duas semifinais. Todos os quarenta países participantes votaram na final. \textbf{Sombreado} (LOC), encontra-se o país vencedor.



Outros países

A elegibilidade para uma potencial participação no Festival Eurovisão da Canção requer uma emissora nacional com participação ativa na \textbf{UER} (ORG) , que poderá transmitir o \textbf{Festival} (LOC) através da rede \textbf{Eurovisão} (LOC) . A \textbf{UER} (ORG) irá convidar a participação no festival para todos os cinquenta e seis membros ativos. Em contraste com os anos anteriores, o membro associado \textbf{Austrália} (LOC) não precisará de um convite para o concurso de 2022, pois recebeu permissão para participar até 2023.



\textbf{Transmissão do Festival Todas} (MISC) as emissoras participantes podem optar por comentar no site ou remotamente para dar ao público local uma visão geral do programa e informações de votação. Embora tenham que transmitir pelo menos as semifinais onde votam e as finais, a maioria das emissoras exibe os três programas com planos de programação diferentes. Algumas emissoras não participantes podem querer continuar querer transmitir o \textbf{Festival} (LOC). Estas são as emissoras que confirmaram transmitir e/ou os seus comentadores:

Notas: (a) Inicialmente, a \textbf{Ucrânia} (LOC) escolheu \textbf{Alina} (LOC) Pash com a canção " \textbf{Shadows of Forgotten Ancestors} (MISC) " para representar o país. No entanto, após controvérsia sobre uma suposta violação das regras que impedem artistas ucranianos de se apresentaram na \textbf{Rússia} (LOC) ou na \textbf{Crimeia} (LOC) de competir na seleção nacional ucraniana, a vencedora da seleção, \textbf{Alina Pash} (PER), desistiu da competição.



Festival Eurovisão da Canção 2023



Países que classificaram-se para a final \textbf{Países} (MISC) que não se classificaram para a final \textbf{Países} (MISC) que participaram no passado, mas não em 2023

O \textbf{Festival Eurovisão da Canção 2023} (MISC) (em inglês : \textbf{Eurovision Song Contest} (ORG) 2023 ; em francês : \textbf{Concours Eurovision} (MISC) de la chanson 2023 ) foi a 67. edição do \textbf{Festival Eurovisão da Canção} (MISC) . Embora a \textbf{Ucrânia} (LOC) tenha vencido a edição de 2022 - com a música  \textbf{Stefania} (LOC) , dos \textbf{Kalush Orchestra} (MISC) - e, como tal, por tradição, o país teria acolhido e organizado o festival do ano seguinte, a \textbf{União Europeia de Radiodifusão} (ORG) determinou que a \textbf{Ucrânia} (LOC) não poderia acolher devido a questões de segurança relacionadas com a invasão russa da \textbf{Ucrânia} (LOC) em 2022 .

O \textbf{Festival} (LOC) foi, portanto, realizado no \textbf{Reino Unido} (LOC) , na sequência do segundo lugar do país no festival de 2022. Esta edição será organizada pela \textbf{UER} (ORG) e pela \textbf{BBC} (ORG) , em colaboração com a rádio e televisão ucraniana UA:\textbf{PBC} (ORG) , e decorreu no \textbf{Liverpool Arena} (LOC) , em \textbf{Liverpool} (LOC) , nos dias 9 , 11 e 13 de maio de 2023. Esta foi a nona vez que o \textbf{Reino Unido} (LOC) acolheu o \textbf{Festival} (LOC), tendo a última realização acontecido em 1998 . Trinta e sete países participaram no \textbf{Festival} (LOC), com a \textbf{Bulgária} (LOC) , \textbf{Macedónia do Norte} (LOC) e \textbf{Montenegro} (LOC) ausentes devido a razões económicas, sobretudo causadas pela crise energética mundial de 2021-2022.



Localização

O \textbf{Festival de 2023} (MISC) foi realizado em \textbf{Liverpool} (LOC) no \textbf{Reino Unido} (LOC) . Foi a nona vez que o país acolheu o \textbf{Festival} (LOC), tendo-o feito anteriormente em 1960 , 1963 , 1968 , 1972 , 1974 , 1977 , 1982 e a última em 1998 . O \textbf{Liverpool Arena} (LOC) foi o local escolhido para acolher o evento. Localizado no complexo ACC \textbf{Liverpool} (LOC), esta infraestrutura tem capacidade para 11.000 pessoas sentadas e é frequentemente palco de diversos concertos e eventos desportivos. No local chegou a ser organizado o \textbf{MTV Europe Music Awards} (ORG) de 2008 , o evento \textbf{Personalidade Desportiva do Ano da BBC} (MISC) de 2008 e 2017 e mais recentemente o \textbf{Campeonato Mundial de Ginástica Artística} (MISC) de 2022

Os reis \textbf{Carlos III} (PER) e \textbf{Camila} (PER) , no dia 28 de Abril de 2023 inauguraram o palco do festival. No palco foi instalado um dispositivo com um botão para que o soberano e a mulher fizessem as honras e acendessem as luzes, dando o 'pontapé de saída' para a 67. edição da \textbf{Eurovisão} (LOC).

A edição de 2022 foi ganha pela \textbf{Ucrânia} (LOC) , que, de acordo com a tradição eurovisiva, seria convidada a sediar o próximo \textbf{Festival} (LOC) pela tradição da \textbf{Eurovisão} (LOC), com o \textbf{Presidente da Ucrânia} (MISC) \textbf{Volodymyr Zelensky} (PER) a manifestar a vontade de organizar o certame numa \textbf{Mariupol} (MISC) reconstruída. No entanto, à luz da \textbf{invasão da} (MISC) \textbf{Ucrânia} (LOC) pela \textbf{Rússia} (LOC) em 2022 , especulou-se que o país não seria capaz de organizar o evento, e devido a tal situação, vários países expressaram o interesse em acolher o \textbf{Festival} (PER) caso a \textbf{Ucrânia} (LOC) não o pudesse fazer: \textbf{Bélgica} (LOC) , \textbf{Espanha} (LOC) (que rapidamente desistiu) \textbf{Itália} (LOC) , \textbf{Países Baixos} (LOC) , \textbf{Polónia} (LOC) , \textbf{Reino Unido} (LOC) e \textbf{Suécia} (LOC) .

A 16 de maio de 2022, o diretor da emissora ucraniana \textbf{Suspilne} (PER) declarou que desejava acolher o \textbf{Festival numa Ucrânia} (MISC) em paz e que esperava que o pais seria capaz de garantir a segurança de todos os participantes e das suas delegações durante o evento. A 20 de maio de 2022, \textbf{Chernotytskyi} (PER) declarou mais tarde que a emissora tinha começado as discussões com a \textbf{UER} (ORG) para organizar o evento.

Vários políticos ucranianos defendiam a realização do \textbf{Festival na Ucrânia} (MISC). O presidente ucraniano \textbf{Volodymyr Zelensky} (PER) , declarou que esperava que um dia houvesse a possibilidade de organizar o \textbf{Festival em Mariupol} (MISC) . A 26 de maio de 2022, o primeiro vice-chefe da cidade de \textbf{Kiev} (LOC) , \textbf{Kyivde Mykola Povoroznyk} (PER), declarou que \textbf{Kiev} (LOC) estava pronta para acolher o \textbf{Festival} (LOC) se necessário. A 3 de junho de 2022, o ministro da \textbf{Cultura da Ucrânia} (LOC), \textbf{Oleksandr Tkachenko} (PER), declarou ter a intenção de discutir algumas mudanças com a \textbf{UER} (ORG) para poder acolher o \textbf{Festival no país} (MISC). A 10 de junho de 2022, a representante do governo ucraniano, \textbf{Verkhovna Rada} (MISC), declarou que um comité tinha sido criado para organizar o \textbf{Festival} (LOC).

A 17 de junho de 2022 foi anunciado através de um comunicado oficial da \textbf{UER} (ORG), que a \textbf{Ucrânia} (LOC) não irá acolher o \textbf{Festival de 2023} (MISC) depois do \textbf{Grupo de Referência} (ORG) do certame ter avaliado as garantias de segurança apresentadas pela emissora ucraniana para receber a edição do próximo ano. Ainda no comunicado, a \textbf{UER} (ORG) confirmou que "para garantir a continuidade do evento, a \textbf{UER} (ORG) irá iniciar as discussões com a \textbf{BBC} (ORG), enquanto vice-campeã do concurso deste ano, para potencialmente sediar o \textbf{Festival Eurovisão 2023} (MISC) no \textbf{Reino Unido} (LOC)" , garantindo também que "é nossa total intenção que a vitória da \textbf{Ucrânia} (LOC) se reflita no concurso do próximo ano", deixando a porta aberta para uma possível qualificação automática para a final .

Com o \textbf{Reino Unido} (LOC) a concordar em acolher o \textbf{Festival de 2023} (MISC), foi a quinta vez que teve a ocasião de o fazer depois de um pais vencedor não o conseguir acolher seguindo dos \textbf{Países Baixos} (LOC) em 1960 , a \textbf{França} (LOC) em 1963 , o \textbf{Mónaco} (LOC) em 1972 e o \textbf{Luxemburgo} (LOC) em 1974 . Nesta ocasião, foi a nona vez que o \textbf{Festival} (LOC) se realizado em solo britânico batendo assim o próprio recorde. No entanto, após o anúncio do início das conversações com a \textbf{BBC} (ORG) , o \textbf{Ministro da Cultura da Ucrânia} (MISC) \textbf{Oleksandr Tkachenko} (PER) e os anteriores vencedores ucranianos da \textbf{Eurovisão} (LOC) pediram à \textbf{UER} (ORG) a revisão da decisão, com o próprio \textbf{Primeiro-ministro} (LOC) do \textbf{Reino Unido Boris Johnson} (LOC) a defender a realização do certame na \textbf{Ucrânia} (LOC). Para além disso, o jornal britânico \textbf{The Guardian} (ORG) avançou com a hipótese de, perante uma possível recusa do \textbf{Reino Unido} (LOC), a \textbf{Eurovisão} (LOC) se realizar em \textbf{Bruxelas} (LOC) .

Logo depois do anúncio de que a \textbf{Ucrânia} (LOC) não organizaria o evento, varias cidades do \textbf{Reino Unido} (LOC) expressaram o interesse em acolher: \textbf{Aberdeen} (LOC) , \textbf{Belfast} (LOC) , \textbf{Birmingham} (LOC) , \textbf{Brighton} (LOC) , \textbf{Bristol} (LOC) , \textbf{Cardiff} (LOC) , \textbf{Edimburgo} (LOC) , \textbf{Glasgow} (LOC) , \textbf{Leeds} (LOC) , \textbf{Liverpool} (LOC) , \textbf{Londres} (LOC) , \textbf{Manchester} (LOC) , \textbf{Newcastle} (LOC) , \textbf{Sheffield} (LOC) , \textbf{Sunderland} (LOC) , e \textbf{Wolverhampton} (PER) .

A 25 de julho de 2022, a \textbf{UER} (ORG) confirmou a realização do \textbf{Festival de 2023} (MISC) no \textbf{Reino Unido} (LOC) com a seleção da cidade-sede no mesmo dia. Na primeira fase do processo de seleção, as cidades podem apresentar oficialmente a sua candidatura, tendo as cidades de \textbf{Aberdeen} (LOC) , \textbf{Leeds} (LOC) , \textbf{Liverpool} (LOC) , \textbf{Manchester} (LOC) apresentado as suas candidaturas no mesmo dia. A 26 de julho seguiram-se as candidaturas de \textbf{Bristol} (LOC) , \textbf{Glasgow} (LOC) e \textbf{Sheffield} (LOC) com \textbf{Birmingham} (LOC) a oficializar no dia seguinte a 27 de julho . A 3 de agosto , a capital galesa , \textbf{Cardiff} (LOC) , retirou o interesse em acolher o \textbf{Festival} (PER). A 12 de agosto de 2022, a \textbf{BBC} (ORG) anunciou que entre as 20 cidades interessadas, apenas 7 seriam pré-selecionadas: \textbf{Birmingham} (LOC), \textbf{Glasgow} (LOC), \textbf{Leeds} (LOC), \textbf{Liverpool} (LOC), \textbf{Manchester} (LOC), \textbf{Newcastle} (LOC), e \textbf{Sheffield} (LOC). Essas cidades foram para a segunda fase de seleção, onde desenvolveram detalhadamente suas propostas para passarem pela avaliação da \textbf{BBC} (ORG) com base numa lista de critérios. A 27 de setembro de 2022, a \textbf{BBC} (ORG) anunciou uma atualização da lista de cidades pré-selecionadas, reduzindo-a a apenas duas cidades, \textbf{Glasgow} (LOC) e \textbf{Liverpool} (LOC), e a 7 de outubro, \textbf{Liverpool} (LOC) foi a cidade selecionada.

\textbf{†} (LOC) Cidade selecionada ‡ \textbf{Cidades} (LOC) finalistas ‡ Cidades pré-selecionadas Cidades não selecionadas

\textbf{Cidades} (LOC) não selecionadas \textbf{Cidades finalistas} (LOC) Cidades que demonstraram interesse, mas não se candidataram

Além do local principal, a cidade anfitriã também organizou eventos paralelos em relação ao \textbf{Festival} (LOC).

A " \textbf{Eurovision Village} (MISC)" (\textbf{Aldeia da Eurovisão} (LOC)) foi o espaço oficial para os fãs e patrocinadores durante as semanas do evento. Nesta área foi possível assistir a performances dos participantes do \textbf{Festival} (LOC) e de artistas locais e também ver as três emissões emitidas em direto a partir do local principal. A aldeia foi localizada na \textbf{Pier Head} (LOC), um conjunto de edifícios na margem do \textbf{rio Mersey} (LOC) e esteve aberto dos dias 5 a 13 de maio de 2023. O \textbf{EuroClub} (MISC) foi organizado no \textbf{Camp and Furnance} (MISC), onde ocorreu realizada a festa privada para delegados, imprensa e fãs credenciados do Festival Eurovisão da Canção.



Produção

O \textbf{Festival Eurovisão da Canção 2023} (MISC) foi realizado pela emissora britânica \textbf{BBC} (ORG) em colaboração com a emissora ucraniana \textbf{UA} (ORG):\textbf{PBC} (ORG) . As três emissões foram produzidas pela \textbf{BBC Studios} (ORG) e pela subsidiária comercial da \textbf{BBC} (ORG). \textbf{Rachel Ashdown} (PER) liderou a comissão encarregue da produção no que diz respeito à aparência e tom do \textbf{Festival} (LOC), incluindo o sorteio que decidiu em quais semifinais cada país atuaria. \textbf{Martin Green CBE} (PER) foi nomeado diretor administrativo e liderou uma equipa responsável por supervisionar todos os aspetos do \textbf{Festival} (PER). Já \textbf{Andrew Cartmell} (PER) foi o escolhido para produtor executivo do \textbf{Festival} (LOC), assumindo a responsabilidade de ambos os espetáculos. Em termos de liderança criativa, \textbf{Lee Smithurst} (PER) foi a escolhida para ser o "cérebro do espetáculo", com a responsabilidade de criar conteúdo editorial. Para o cargo de "cérebro do \textbf{Festival} (LOC)" está \textbf{Twan van de} (PER) \textbf{Nieuwenhuijzen} (LOC) ; ocupando o mesmo cargo nos \textbf{Festivais} (MISC) de 2021 e 2022 , o criador de conteúdo teve a responsabilidade de trabalhar com as delegações nacionais para propor ideias criativas em palco e para os artistas. Finalmente, \textbf{James O'Brien} (PER) juntou-se à equipa da \textbf{BBC Studios} (ORG) como \textbf{Executivo Encarregado da Produção} (MISC) e supervisionou uma equipa responsável pelo resultado técnico de ambos os espetáculos e transmissão do \textbf{Festival} (LOC).

A 7 de outubro, ao mesmo tempo que o anúncio da cidade anfitriã, a \textbf{UER} (ORG) revelou o logo padrão para o \textbf{Festival de 2023} (MISC). O coração da \textbf{Eurovisão} (LOC), contém, por tradição, a bandeira do país vencedor da edição anterior, neste caso a bandeira da \textbf{Ucrânia} (LOC) . Debaixo de "\textbf{Song Contest} (LOC)" (\textbf{Festival da Canção} (PER)) está excecionalmente designado "\textbf{United Kingdom} (ORG)" (\textbf{Reino Unido} (LOC)) seguindo-se da cidade anfitriã, "\textbf{Liverpool} (LOC) 2023".

A identidade gráfica e o slogan "\textbf{United by Music} (ORG)" (Unidos pela música), foram revelados a 31 de janeiro de 2023. Criado pela consultora de marcas londrina \textbf{Superunion} (MISC) e a empresa de produção ucraniana \textbf{Starlight Media} (MISC) , a arte foi construída à volta de uma linha de corações em duas dimensões parecendo um electrocardiograma , representando o ritmo e o som, enquanto que as cores foram inspiradas das bandeiras ucraniana e britânica . Já a topografia foi inspirada dos painéis toponímicos da ruas de \textbf{Liverpool} (LOC) no século 20 e a herança musical da cidade.

A 22 de fevereiro de 2023, os apresentadores foram anunciados. A cantora britânica \textbf{Alesha Dixon} (PER) , a atriz britânica \textbf{Hannah Waddingham} (PER) e a cantora ucraniana \textbf{Julia Sanina} (PER) apresentaram as três emissões, com o apresentador britânico da televisão irlandesa, \textbf{Graham Norton} (PER) , a juntar-se ao grupo na final. \textbf{Timur Miroshnychenko} (PER) (que co-apresentou o \textbf{Festival de 2017} (MISC)) e \textbf{Sam Quek} (PER) apresentaram a cerimónia de abertura do evento.

O design do palco para a edição de 2023 foi revelado a 2 de fevereiro de 2023. Criado pelo designer nova-iorquino \textbf{Julio Himede} (PER), o design do palco baseia-se "dos princípios de união, celebração e comunidade", tendo inspiração num grande abraço e em "aspetos culturais e similaridades entre a \textbf{Ucrânia} (LOC), o \textbf{Reino Unido} (LOC) e especificamente a cidade de \textbf{Liverpool} (LOC)". O palco teva largura de 450m e cerca de 220 m de ecrãs \textbf{LED} (ORG) independentes, entre eles mais de 700 ecrãs \textbf{LED} (ORG) no solo e 1500 metros de luzes \textbf{LED} (ORG).

Durante o período de votação, \textbf{Sam Ryder} (PER), artista que representou o \textbf{Reino Unido} (LOC) na edição de 2022, subiu a palco para apresentar o single "\textbf{Mountain} (PER)" e contou com o convidado \textbf{Roger Taylor} (PER), dos \textbf{Queen} (ORG). Para homenagear \textbf{Liverpool} (LOC), vários antigos participantes do festival da \textbf{Eurovisão Mahmood} (MISC), \textbf{Netta} (LOC), \textbf{Cornelia Jacobs} (PER), \textbf{Dadi Freyr} (PER), \textbf{Sonia} (PER) e \textbf{Duncan Laurence} (PER)) juntaram-se em palco para um medley com canções de artistas da cidade.



\textbf{Formato} (MISC)A 22 de novembro de 2022, a \textbf{UER} (ORG) anunciou mudanças importantes no sistema de voto para o \textbf{Festival de 2023} (MISC). Os resultados das semifinais seriam determinado somente pelo televoto, já os resultados da final seriam determinados pelo conjunto de júris nacionais e pelo televoto, como em anos anteriores. No caso eventual de um júri estar indisponível para entregar o resultado do televoto para as semifinais, o resultado de um júri de reserva seria utilizado. Os espetadores de países não participantes estavam também aptos a votar em todos os espetáculos, com os seus votos a serem agregados e apresentados como um conjunto de pontos individual.

Pelo terceiro ano consecutivo os coros pré-gravados foram permitidos, assim, cada delegação pôde usar cantores de coro no palco ou fora dele, ou uma combinação de interpretação em direto com coros pré-gravados.

A \textbf{BBC} (ORG) contratou um empresa de produção independente para produzir o sorteio das semifinais, que inclui a passagem da chave da cidade da cidade anfitriã precedente \textbf{Turim} (LOC) à cidade de \textbf{Liverpool} (LOC).

O sorteio para determinar em que semifinal os país participantes atuariam teve lugar a 31 de janeiro de 2023 na \textbf{Saint George's Hall} (LOC) . Os 31 países semifinalistas foram divididos em 5 potes, baseados em padrões de votação históricos calculados pelo parceiro de televoto oficial do \textbf{Festival} (LOC), \textbf{Digame} (LOC). A razão da utilização do sistema de potes foi a necessidade de reduzir as chances de "bloc-voting" (votação em bloco) para aumentar o suspense nas semifinais. O sorteio determinou ainda em quais semifinais votariam os 6 países participantes qualificados automaticamente, sendo eles o país vencedor da edição anterior, a \textbf{Ucrânia} (LOC) e os países do "\textbf{Big Five} (MISC)", \textbf{Alemanha} (LOC), \textbf{França} (LOC), \textbf{Itália} (LOC), \textbf{Espanha} (LOC) e o \textbf{Reino Unido} (LOC). A cerimonia foi apresentada pelos apresentadores britânicos \textbf{AJ Odudu} (PER) e \textbf{Rylan Clark} (PER), e incluiu a passagem da insígnia da cidade do presidente da câmara de \textbf{Turim} (LOC), \textbf{Stefano Lo Russo} (LOC), para a presidente da câmara de \textbf{Liverpool} (LOC), \textbf{Joanne Anderson} (PER).

A produtora londrina \textbf{ModestTV} (MISC) foi comissionada para produzir a transmissão da cerimonia. O evento foi transmitido no canal oficial do \textbf{YouTube do Festival} (MISC), na \textbf{BBC Two} (MISC) no \textbf{Reino Unido} (LOC), e na \textbf{Suspilne Kultura} (MISC) na \textbf{Ucrânia} (LOC) com comentários de \textbf{Timur Miroshnychenko} (PER).

Como para a organização do sorteio das semifinais, a \textbf{BBC} (ORG) contratou uma empresa de produção independente para criar o conceito para e produzir os " postcards" (cartas postais) que seriam difundidos antes de cada atuação. A gravação dos mesmos ocorreu entre janeiro de abril de 2023.



Países participantes

A 20 de outubro de 2022, a \textbf{UER} (ORG) anunciou que 37 países marcariam presença no \textbf{Festival de 2023} (MISC), o número de participantes mais baixo desde a edição de 2014 , pois a \textbf{Bulgária} (LOC) , \textbf{Macedónia do Norte} (LOC) e \textbf{Montenegro} (LOC) optaram por não participar devido a razões financeiras.

A 10 de fevereiro de 2023, foi anunciado que a \textbf{República Checa} (LOC) participaria pela primeira vez como o seu nome curto em inglês, " \textbf{Czechia} (MISC)" (\textbf{Chéquia} (LOC)).

\textbf{Participaram do Festival} (MISC) quatro artistas que já atuaram como vocalista principal do mesmo país: a sueca \textbf{Loreen} (MISC) venceu a edição de 2012 , na qual também participou \textbf{Pasha Parfeny} (MISC), representante da \textbf{Moldávia} (LOC) que em 2013 foi coro para \textbf{Aliona Moon} (PER) . \textbf{Marco Mengoni} (PER) representou a \textbf{Itália} (LOC) em 2013 , e a lituana \textbf{Monika Linkytė} (MISC) representou seu país em 2015 junto com \textbf{Vaidas Baumila} (PER) . Além disso, o belga \textbf{Gustaph} (ORG) foi coro para \textbf{Sennek} (LOC) em 2018 e para os \textbf{Hooverphonic} (ORG) em 2021 , e \textbf{Iru Khechanovi} (MISC) , a representante da \textbf{Geórgia} (LOC), venceu o \textbf{Festival Eurovisão da Canção Júnior} (MISC) 2011 como membro do grupo \textbf{Candy} (PER).

As semifinais realizaram-se nos dias 9 e 11 de maio de 2023 pelas 20h00 ( UTC+1 ). A ordem de atuação dos países foi revelada a 31 de janeiro.

A final realizou-se no dia 13 de maio de 2023 pelas 20:00 ( UTC+1 ). Vinte e seis países participaram na final, composta pela \textbf{Ucrânia} (LOC), país vencedor da edição anterior, pelos \textbf{Big Five} (MISC) (que inclui o \textbf{Reino Unido} (LOC)), e as dez melhores atuações de cada semifinal. Todos os países votaram na final.



Festival Eurovisão da Canção 2024



Países que classificaram-se para a final \textbf{País} (LOC)/es que se classificou/\textbf{ram} (MISC) nas semifinais, mas foi/foram desclassificado/s antes da final \textbf{Países} (MISC) que não se classificaram para a final \textbf{Países} (MISC) que participaram no passado, mas não em 2024

O \textbf{Festival Eurovisão da Canção 2024} (MISC) [ a ] foi a 68. edição do Festival Eurovisão da Canção . Realizou-se em \textbf{Malmo} (LOC) , \textbf{Suécia} (LOC) e o vencedor foi \textbf{Nemo} (PER) com a música " \textbf{The Code} (MISC) ". Foi organizado pela \textbf{União Europeia de Radiodifusão} (ORG) (\textbf{UER} (ORG)/\textbf{EBU} (MISC)) em conjunto com a emissora anfitriã, a \textbf{Televisão Pública da Suécia} (MISC) (SVT). As duas semifinais tiveram lugar nos dias 7 e 9 de maio , enquanto a final decorreu a 11 de maio .

A 14 de Novembro de 2023, a \textbf{União Europeia de Radiodifusão} (ORG) anunciou que "\textbf{United By Music} (MISC)" ("unidos através da música") tinha sido escolhido como o slogan da edição de 2024, o mesmo que no ano anterior, slogan esse que se manterá nas próximas edições do festival.



Localização

O festival realiza-se na \textbf{Suécia} (LOC) , após a vitória do país na edição de 2023 com a música " \textbf{Tattoo} (MISC) ", de \textbf{Loreen} (LOC) . É a sétima vez que o país acolhe o festival, depois de \textbf{Estocolmo} (LOC) em 1975 , 2000 e 2016 , \textbf{Gotemburgo} (LOC) em 1985 e \textbf{Malmö} (LOC) em 1992 e 2013 .

Imediatamente após a vitória da \textbf{Suécia} (LOC) no Festival Eurovisão da Canção 2023 , as primeiras cidades a manifestar interesse em organizar a edição de 2024 foram \textbf{Estocolmo} (LOC) , \textbf{Gotemburgo} (LOC) e \textbf{Malmö} (LOC) , as três maiores cidades do país, que já haviam sediado o certame em ocasiões anteriores (\textbf{Estocolmo} (LOC) em 1975 , 2000 e 2016 , \textbf{Gotemburgo} (LOC) em 1985 e \textbf{Malmö} (LOC) em 1992 e 2013 ). Para além destas, várias outras cidades também manifestaram a intenção de concorrer nos dias que se seguiram à vitória de 2023, nomeadamente \textbf{Eskilstuna} (LOC) , \textbf{Jönköping} (LOC) , \textbf{Örnsköldsvik} (LOC) , \textbf{Partille} (LOC) e \textbf{Sandviken} (LOC) .

A SVT estabeleceu como prazo 12 de junho de 2023 para as cidades interessadas se inscreverem formalmente. \textbf{Estocolmo} (LOC) e \textbf{Gotemburgo} (LOC) anunciaram oficialmente as suas propostas a 7 e 10 de junho , respectivamente, seguidos por \textbf{Malmö} (LOC) e \textbf{Örnsköldsvik} (LOC) a 13 de junho . A 8 de junho , o município de \textbf{Sandviken} (LOC) decidiu desistir, e em \textbf{Jönköping} (LOC) , \textbf{Helene Berg} (PER), diretora administrativa da agência de turismo de \textbf{Småland} (LOC) , pediu às autoridades locais que apresentassem uma candidatura, o que náo aconteceu, por falta de meios financeiros e de uma arena que cumprisse os requisitos de capacidade da \textbf{União Europeia de Radiodifusão} (ORG) .

Pouco antes do encerramento do período de inscrições, a \textbf{SVT} (ORG) revelou ter recebido propostas de várias cidades, embora não tenha especificado quais. A 7 de julho , \textbf{Malmö} (LOC) foi escolhida como a cidade anfitriã.

Chave: \textbf{†} (LOC) Cidade selecionada ‡ \textbf{Cidades} (LOC) finalistas ‡ Cidades pré-selecionadas Cidades não selecionadas



Países participantes

A 5 de dezembro de 2023 , foram revelados que estariam presentes 37 países, incluindo o \textbf{Luxemburgo} (LOC) , que participa pela primeira vez desde 1993 . No entanto, a participação da \textbf{Roménia} (LOC) só foi decidida a 25 de janeiro de 2024 , com a confirmação da sua desistência.

As semifinais realizaram-se nos dias 7 e 9 de maio de 2024 pelas 20h00 ( UTC+1 ). A ordem de atuação dos países foi revelada a 30 de janeiro .

Os países participantes, em conjunto com a \textbf{Alemanha} (LOC) , o \textbf{Reino Unido} (LOC) , a \textbf{Suécia} (LOC) e o \textbf{Resto do Mundo} (MISC) (votação agregada dos países não-participantes), votam nesta semifinal. Além das canções concorrentes, o \textbf{Reino Unido} (LOC) , a \textbf{Alemanha} (LOC) e a \textbf{Suécia} (LOC) apresentaram as suas músicas durante a semifinal, aparecendo no palco após as participações da \textbf{Irlanda} (LOC) , \textbf{Islândia} (LOC) e \textbf{Moldávia} (LOC) , respetivamente.

A primeira semifinal foi aberta pelos ex-participantes \textbf{Eleni Foureira} (PER) , \textbf{Eric Saade} (PER) e \textbf{Chanel} (PER) , que interpretaram as respetivas músicas concorrentes - " \textbf{Fuego} (PER) " ( \textbf{Chipre 2018} (MISC) ), " \textbf{Popular} (MISC) " ( \textbf{Suécia} (LOC) 2011 ) e " \textbf{SloMo} (MISC) " ( \textbf{Espanha} (LOC) 2022 ). Como atos de intervalo, a semifinal contará com o tricampeão da \textbf{Irlanda} (LOC) ( 1980 , 1987 e 1992 ) \textbf{Johnny Logan} (PER) apresentando a música vencedora sueca de 2012 " \textbf{Euphoria} (LOC) ", e o participante sueco de 2018 \textbf{Benjamin Ingrosso} (PER) apresentando um medley das suas canções "\textbf{Who's Laughing Now} (MISC)", "\textbf{Kite} (MISC)" e "\textbf{Honey Boy} (MISC)". De acordo com um relatório do \textbf{Aftonbladet} (LOC) , a artista original \textbf{Loreen} (MISC) deveria inicialmente se apresentar com \textbf{Johnny Logan} (PER) , mas decidiu não fazê-lo.

Os países participantes, em conjunto com a \textbf{Espanha} (LOC) , a \textbf{França} (LOC) , a \textbf{Itália} (LOC) e o \textbf{Resto do Mundo} (MISC) (votação agregada dos países não-participantes), votam nesta semifinal. Além das canções concorrentes, a \textbf{França} (LOC) , a \textbf{Espanha} (LOC) e a \textbf{Itália} (LOC) apresentaram as suas canções durante o espetáculo, aparecendo no palco após as músicas da \textbf{Chéquia} (LOC) , \textbf{Letónia} (LOC) e \textbf{Estónia} (LOC) , respetivamente.

O intervalo da segunda semifinal contou com \textbf{Helena Paparizou} (PER) , \textbf{Sertab Erener} (PER) e \textbf{Charlotte Perrelli} (PER) que interpretaram as respetivas músicas concorrentes vencedoras - " \textbf{My Number One} (MISC) " ( Grécia 2005 ), " \textbf{Everyway That I Can} (MISC) " ( \textbf{Turquia} (LOC) 2003 ) e " \textbf{Take Me to Your Heaven} (MISC) " ( \textbf{Suécia} (MISC) 1999 ) - com o público a cantar em conjunto.

A final realizou-se no dia 11 de maio de 2024 pelas 20h00 ( UTC+1 ).

A final foi aberta por \textbf{Björn Skifs} (PER) , que interpretou " \textbf{Hooked on a Feeling} (MISC) ", seguida do desfile das bandeiras, acompanhada por um medley de canções suecas. Os intervalos foram protagonizados pela banda sueca \textbf{Alcazar} (LOC) , que interpretou "\textbf{Crying at the Discoteque} (MISC)"; um tributo à música vencedora da edição de 1974 " \textbf{Waterloo} (MISC) ", dos \textbf{ABBA} (ORG) , protagonizado por \textbf{Carola} (PER) ( \textbf{Suécia} (LOC) 1991 ), \textbf{Charlotte Perrelli} (PER) (\textbf{Suécia} (LOC) 1999 ) e \textbf{Conchita Wurst} (PER) ( \textbf{Áustria} (LOC) 2014 ), antecedido por um segmento pré-gravado com os \textbf{ABBA} (ORG), sobre a sua experiência eurovisiva; e \textbf{Loreen} (MISC) , que apresentou o seu novo single "Forever" e " \textbf{Tattoo} (MISC) ", a canção vencedora do \textbf{Festival Eurovisão da Canção 2023} (MISC) .

Apesar de se ter qualificado para a final, os \textbf{Países Baixos} (LOC) foram desqualificados, devido a um episódio de agressão física, protagonizado pelo artista holandês \textbf{Joost Klein} (PER) , a um membro da produção. No entanto, a ordem de atuação manterve-se a mesma, assim como o direito de voto dos \textbf{Países Baixos} (LOC) .



\textbf{Transmissão do Festival Todas} (MISC) as emissoras participantes podem optar por terem comentadores a informar sobre todos os acontecimentos ao seu público local. Embora devam transmitir pelo menos a semifinal em que votam e a final, a maioria das emissoras transmite todos os três espetáculos com planos de programação diferentes. Para além disso, algumas emissoras não participantes transmitem o concurso. O canal oficial da \textbf{Eurovisão} (LOC) no \textbf{YouTube} (MISC) oferece transmissões internacionais em direto sem comentários de todos os espetáculos. A seguir estão as emissoras que confirmaram total ou parcialmente seus planos de transmissão:



\textbf{Controvérsias} (PER)

Desde a eclosão da guerra \textbf{Israel-Hamas} (LOC) , a 7 de outubro de 2023 , têm sido feitos apelos crescentes para que \textbf{Israel} (LOC) seja excluído do certame, com base na crise humanitária resultante das operações militares israelitas na \textbf{Faixa de Gaza} (LOC) ; isto incluiu protestos e petições dirigidas a emissoras nacionais em vários países participantes, nomeadamente na \textbf{Finlândia} (LOC) , \textbf{Islândia} (LOC) , \textbf{Noruega} (LOC) , \textbf{Irlanda} (LOC) , entre outros, exigindo que se retirassem ou pressionassem a \textbf{União Europeia de Radiodifusão} (ORG) para excluir \textbf{Israel} (LOC) . A emissora islandesa RÚV deverá discutir a sua participação com o artista vencedor da sua final nacional. Até janeiro de 2024 , nenhuma emissora manifestou oposição à participação do país.

Em novembro de 2023 , a equipa de produção da \textbf{SVT} (ORG) declarou a sua intenção de aumentar as medidas de segurança e de manter contacto com a autoridade policial de \textbf{Malmö} (LOC) durante o concurso, citando o risco de potenciais ataques terroristas como uma repercussão da guerra.

Durante o mês de janeiro de 2024 , a \textbf{UER} (ORG) anunciou que a \textbf{Eurovisão} (LOC) "é uma competição entre organizações de serviço público de toda a \textbf{Europa} (LOC) e do \textbf{Médio Oriente} (LOC), que são membros da \textbf{União Europeia de Radiodifusão} (ORG). É uma competição para emissoras - não para governos", rejeitando a exclusão de \textbf{Israel} (LOC) do certame.

No final de fevereiro de 2024 , vários relatos dos mídia israelitas afirmaram que duas músicas foram selecionadas para consideração, intituladas "\textbf{October Rain} (PER)" e "\textbf{Dance Forever} (MISC)", e que ambas foram submetidas à \textbf{UER} (ORG) para avaliação, mas foram rejeitados por conterem mensagens políticas. A emissora israelita \textbf{IPBC} (ORG)/\textbf{Kan} (PER) confirmou estas decisões a 3 de março , ao mesmo tempo que afirmou que pediu aos escritores de ambas as canções que "fizessem os ajustes necessários" para que fossem elegíveis. A música selecionada, intitulada " \textbf{Hurricane} (MISC) ", foi aprovada pela \textbf{UER} (ORG) no dia 7 de março e revelada três dias depois.

Várias seleções nacionais foram alvo de protestos a favor do boicote de \textbf{Israel} (LOC) do certame, como na \textbf{Noruega} (LOC) , \textbf{Suécia} (LOC) , \textbf{Dinamarca} (LOC) , e na \textbf{Finlândia} (LOC) . Para além disso, a 11 de março de 2024 , a sinalização digital instalada no \textbf{Malmö Live} (LOC) em preparação para o concurso foi salpicada com tinta vermelha e a sua base foi com spray com as palavras " \textbf{Gaza Livre} (MISC) ".

A polícia de \textbf{Malmö} (LOC) teria recebido três pedidos de protestos a serem organizados a partir de 13 de março de 2024 : um contra a participação israelita na competição, planeada para incluir 10 mil pessoas numa caminhada de \textbf{Stortorget} (LOC) a \textbf{Möllevångstorget} (LOC); outro contra a guerra em \textbf{Gaza} (LOC), a ser organizado fora da \textbf{Malmö Arena} (LOC) ; e um para mostrar apoio à delegação israelita da \textbf{Eurovisão} (LOC).

Durante o primeiro ensaio geral para a final o representante holandês, \textbf{Joost Klein} (PER) , não compareceu para a sua atuação, apesar de ter estado presente durante o desfile das bandeiras. A \textbf{União Europeia de Radiodifusão} (ORG) declarou num comunicado de imprensa que estava "a investigar um incidente que lhe foi comunicado e que envolvia o artista neerlandês", bem como que "este não iria ensaiar até nova ordem". O artista também não esteve presente no espetáculo do júri, tendo sido utilizada uma gravação da sua atuação na segunda semi-final. O incidente envolveu uma funcionária que apresentou queixa contra \textbf{Klein} (PER). As emissoras holandesas \textbf{AVROTROS} (ORG) e \textbf{NPO} (MISC) tiveram uma discussão com a \textbf{UER} (ORG) e este incidente levou à desqualificação do participante neerlandês.

Várias delegações, incluindo \textbf{Portugal} (LOC) , denunciaram que a comitiva israelita estava a cometer a assediar e perseguir membros das comitivas da \textbf{Irlanda} (LOC) , \textbf{Grécia} (LOC) , \textbf{Suíça} (LOC) e \textbf{Países Baixos} (LOC) . A comitiva portuguesa, não sendo alvos dos assédios, por ser mais discreta na oposição à participação de \textbf{Israel} (LOC), afirma ter testemunhado os mesmos. Por esse motivo, \textbf{Portugal} (LOC) terá pedido à \textbf{União Europeia de Radiodifusão} (ORG), organizadora do festival, uma reunião de emergência que terá tido lugar no dia da final.

A equipa da \textbf{Irlanda} (LOC) também apresentou queixa em relação ao assédio do que estava a ser vítima por parte de \textbf{Israel} (LOC).

A comitiva neerlandesa também afirma que o representante apenas respondeu aos assédios dos quais estava a ser vítima. Um membro do evento terá alegadamente filmado \textbf{Joost Klein} (PER) sem o seu consentimento múltiplas vezes, após o mesmo ter pedido que parasse. O incidente não foi considerado grave pelas autoridades suecas.

A \textbf{EBU} (MISC) retirou das redes sociais a atuação da representante de \textbf{Portugal} (LOC), \textbf{Iolanda} (LOC) , porque a artista mostrou as unhas pintadas com o padrão do keffiyeh , símbolo de resistência da \textbf{Palestina} (LOC) e por ter dito no final da canção: "A paz irá prevalecer". O vídeo da atuação foi o último a ser colocado, sendo o único que não seguiu a ordem das apresentações.

A comitiva francesa também temeu censura e sanções após \textbf{Slimane} (LOC) ter feito um discurso acerca da paz durante os ensaios gerais.

Uma bandeira da causa não-binária também foi proibida de entrar no recinto. Mas tarde, a bandeira foi alegadamente erguida durante a atuação da \textbf{Suíça} (LOC). \textbf{Nemo} (PER) identifica-se como pessoa não binária .



Festival Eurovisão da Canção 2025



Países que classificaram-se para a final \textbf{Países} (MISC) que não se classificaram para a final \textbf{Países} (MISC) que participaram no passado, mas não em 2025

O \textbf{Festival Eurovisão da Canção 2025} (MISC) foi a 69. edição da competição organizada pela \textbf{União Europeia de Radiodifusão} (ORG) (\textbf{UER} (ORG)) e pelo organismo de radiodifusão anfitrião, a \textbf{Sociedade Suíça de Radiodifusão} (PER) e \textbf{Televisão} (LOC) ( \textbf{SRG SSR} (ORG) ). O certame realizou-se na cidade de \textbf{Basileia} (LOC) , na \textbf{Suíça} (LOC) , e o vencedor foi \textbf{Johannes Pietsch} (PER) , conhecido pelo seu nome artístico "\textbf{JJ} (PER)", com a música " \textbf{Wasted Love} (MISC) ". Esta foi a terceira vez que a \textbf{Suíça} (LOC) organizou o concurso, tendo-o feito anteriormente nas edições de 1956 e 1989 , realizados em \textbf{Lugano} (LOC) e \textbf{Lausanne} (LOC) , respetivamente.

A edição foi apresentada pela humorista \textbf{Hazel Brugger} (PER) e pela representante da \textbf{Suíça} (LOC) na edição de 1991 \textbf{Sandra Studer} (PER) . Ambas foram acompanhadas pela apresentadora italiana \textbf{Michelle Hunziker} (PER) na apresentação da final.



Localização

Após a vitória da \textbf{Suíça} (LOC) em 2024 , as autoridades locais de \textbf{Genebra} (LOC) manifestaram interesse em acolher a edição de 2025 no \textbf{Palexpo} (MISC) e apresentaram uma candidatura formal. No mesmo dia, o presidente do governo de \textbf{Basileia} (LOC) , \textbf{Conradin Cramer} (PER), também manifestou interesse em que \textbf{Basileia} (LOC) sediasse o festival. A 12 de maio , o \textbf{Olma Hall} (LOC) em \textbf{São Galo} (LOC) foi proposto como um local potencial, \textbf{A 13} (MISC) de maio , \textbf{Lugano} (LOC) descartou uma candidatura para sediar a edição de 2025, enquanto o presidente do governo cantonal de \textbf{Berna} (LOC) , \textbf{Phillipp Müller} (PER), expressou a sua relutância em sediar a competição na capital suíça de facto , citando "um aumento do antissemitismo " no concurso. \textbf{Müller} (PER) foi contrariado publicamente por todo o governo cantonal de \textbf{Berna} (LOC) em 15 de maio; o conselho parabenizou \textbf{Nemo} (PER) e anunciou seu apoio na organização do concurso em \textbf{Berna} (LOC). Enquanto isso, foi relatado que o conselho de \textbf{Zurique} (LOC) realizou uma reunião de "alta prioridade" para discutir uma candidatura. A 14 de maio , \textbf{Lausana} (LOC) descartou candidaturas para sediar a edição de 2025, alegando falta de infraestrutura. No dia 15 de maio , \textbf{Bienna} (LOC) , cidade natal de \textbf{Nemo} (PER), declarou interesse em ser associada e co-anfitriã do evento. Em 17 de maio, o governo local de \textbf{Friburgo} (LOC) declarou que estava examinando uma potencial candidatura para sediar o concurso. Em 5 de junho, o governo de \textbf{Basileia-Cidade} (PER) confirmou a sua candidatura para sediar o concurso, propondo o \textbf{St} (MISC). \textbf{Jakobshalle} (LOC) e o \textbf{St} (MISC). \textbf{Jakob-Park} (MISC) como locais possíveis. Em 6 de junho, os municípios de \textbf{Bienna} (LOC) e \textbf{Berna} (LOC) anunciaram uma potencial candidatura conjunta. No dia 12 de junho, a cidade de \textbf{São Galo} (LOC) anunciou que não apresentará candidatura por não cumprir os requisitos para sediar o concurso.

A emissora anfitriã \textbf{SRG SSR} (ORG) lançou o processo licitatório na semana de 27 de maio de 2024, emitindo uma lista de requisitos para as cidades interessadas. \textbf{Basileia} (LOC), \textbf{Berna} (LOC), \textbf{Genebra} (LOC) e \textbf{Zurique} (LOC) declararam oficialmente o seu interesse e finalizaram as suas propostas em 28 de junho. Representantes da emissora anfitriã visitarão as quatro cidades candidatas no início de julho e selecionarão duas delas até o final do mês. A 30 de agosto de 2024, foi anunciado que a edição de 2025 será em \textbf{Basileia} (LOC).

\textbf{† Cidade Escolhida * Pré-selecionada  Enviou} (LOC) uma proposta



Países participantes

A elegibilidade para participação no \textbf{Festival Eurovisão da Canção} (MISC) requer uma emissora nacional com membros ativos da \textbf{UER} (ORG) capaz de receber o concurso através da \textbf{Rede Eurovisão} (MISC) e transmiti-lo ao vivo para todo o país. A \textbf{UER} (ORG) emite um convite para participar no concurso a todos os membros ativos. Os países que fazem parte dos "\textbf{Big Five} (MISC)" e o país anfitrião, a \textbf{Suíça} (LOC) nesta ocasião, receberam automaticamente uma vaga na final do concurso, enquanto todos os outros países foram colocados em duas semifinais.

As semifinais realizaram-se nos dias 13 e 15 de maio de 2025 pelas 20h00 ( UTC+1 ). A ordem de atuação dos países foi revelada a 28 de janeiro .

Os países participantes, em conjunto com a \textbf{Espanha} (LOC) , \textbf{Itália} (LOC) , \textbf{Suíça} (LOC) e o \textbf{Resto do Mundo} (LOC) (votação agregada dos países não-participantes), votaram nesta semifinal.

Os países participantes, em conjunto com a \textbf{Alemanha} (LOC) , \textbf{França} (LOC) , \textbf{Reino Unido} (LOC) e o \textbf{Resto do Mundo} (MISC) (votação agregada dos países não-participantes), votaram nesta semifinal.

A final realizou-se a 17 de maio de 2025 pelas 20h00 ( UTC+1 ).



Outros países

Os organismos de radiodifusão de membros activos da \textbf{União Europeia de Radiodifusão} (ORG) em \textbf{Andorra} (LOC) , \textbf{Bósnia e Herzegovina} (LOC) e \textbf{Eslováquia} (LOC) confirmaram até agora a sua não participação antes do anúncio da lista de participantes pela \textbf{UER} (ORG).



\textbf{Transmissão do Festival Todos} (MISC) os organismos de radiodifusão participantes podem optar por ter comentadores no local ou à distância que forneçam informações sobre a votação ao seu público local. Embora sejam obrigados a transmitir a final e a semi-final em que o seu país vota, a maior parte dos organismos de radiodifusão cobre os três programas. Alguns organismos de radiodifusão não participantes também transmitem o concurso. O \textbf{canal do Festival Eurovisão da Canção} (LOC) no \textbf{YouTube} (MISC) sempre oferece transmissões internacionais em direto, sem comentários, de todos os espectáculos.



Produção

O \textbf{Festival Eurovisão da Canção 2025} (MISC) foi produzido pela \textbf{Sociedade Suíça de Radiodifusão} (PER) e \textbf{Televisão} (LOC) ( \textbf{SRG SSR} (ORG) ). A equipa principal foi constituída por \textbf{Reto Peritz} (PER) e \textbf{Moritz Stadler} (PER), como produtores executivos, e \textbf{Yves Schifferle} (PER), como chefe do espetáculo. Repetindo as suas funções da edição de 2024, \textbf{Christer Björkman} (PER) foi o chefe do concurso e \textbf{Tobias Åberg} (PER) o chefe da produção, com outros elementos da equipa de produção, incluindo \textbf{Nadja Burkhardt-Tracol} (PER) como chefe do evento, \textbf{Manfred Winz} (PER) como chefe das finanças, \textbf{Aurore Chatard} (PER) como chefe da segurança e \textbf{Kevin Stuber} (PER) como chefe do departamento jurídico.

A organização do concurso foi reestruturada para 2025, o que foi anunciado pela \textbf{UER} (ORG) a 1 de julho de 2024, na sequência de uma análise das controvérsias do concurso de 2024. O cargo de supervisor executivo, ocupado por \textbf{Martin Österdahl} (PER) desde a edição de 2021 , foi renomeado para supervisor do espetáculo, tendo sido criados dois novos cargos: o de diretor da \textbf{Eurovisão} (LOC) e o de chefe de marca e comercial. O diretor do festival supervisionará o trabalho do supervisor do espetáculo e do chefe da marca e do comercial.

\textbf{O Conselho Executivo de Basileia-Cidade} (MISC) contribuiu com 35 milhões de francos suíços (cerca de 37,3 milhões de euros) para o orçamento do concurso.



Festival Eurovisão da Canção 2026



O \textbf{Festival Eurovisão da Canção 2026} (MISC) será a 70. edição do Festival Eurovisão da Canção . Será organizado pela \textbf{União Europeia de Radiodifusão} (ORG) (\textbf{UER} (ORG)) e pela estação de radiodifusão anfitriã, que se prevê ser a austríaca \textbf{Österreichischer Rundfunk} (PER) (\textbf{ORF} (ORG)), na sequência da vitória de \textbf{Johannes Pietsch} (PER) com a canção  \textbf{Wasted Love} (ORG)  no concurso de 2025. O festival terá lugar em \textbf{Viena} (LOC) , no \textbf{Wiener Stadthalle} (PER) , e espera-se que decorra em maio de 2026.

Até 10 de dezembro de 2025, 34 radiodifusores de serviço público confirmaram provisoriamente a intenção de participar no concurso. A \textbf{Islândia} (LOC), a \textbf{Irlanda} (LOC), os \textbf{Países Baixos} (LOC), a \textbf{Eslovénia} (LOC) e a \textbf{Espanha} (LOC) estão a boicotar em protesto contra a inclusão de \textbf{Israel} (LOC) no contexto da guerra em \textbf{Gaza} (LOC), marcando o maior boicote desde a data de 1970, enquanto a \textbf{Bulgária} (LOC), a \textbf{Moldávia} (LOC) e a \textbf{Roménia} (LOC) regressam após ausências em edições recentes.



Localização

O concurso de 2026 terá lugar em \textbf{Viena} (LOC) , \textbf{Áustria} (LOC) , na sequência da vitória do país no concurso de 2025 com a canção  \textbf{Wasted Love} (ORG) , interpretada por \textbf{JJ} (PER) . Será a terceira vez que a \textbf{Áustria} (LOC) acolhe o concurso, depois de já o ter feito em 1967 e 2015 , ambas as vezes também em \textbf{Viena} (LOC). O local escolhido para o concurso é a \textbf{Wiener Stadthalle} (PER) , com capacidade para 16 152 espectadores, que já tinha acolhido o concurso em 2015.

Além do local principal, a \textbf{Rathausplatz} (LOC) será o local da \textbf{Eurovision Village} (MISC), que acolhe atuações dos participantes do concurso e de artistas locais, assim como transmissões dos espetáculos em direto para o público em geral. A \textbf{Câmara Municipal de Viena} (LOC) acolherá o \textbf{EuroClub} (MISC), que organiza as festas oficiais e atuações privadas dos participantes do concurso, bem como o \textbf{Tapete Turquesa} (PER) e a cerimónia de abertura a 10 de maio de 2026, onde os concorrentes e as suas delegações serão apresentados à imprensa acreditada e aos fãs. Ambos os locais repetirão as mesmas funções que desempenharam em 2015.

Após a vitória da \textbf{Áustria} (LOC) em 2025, o diretor da \textbf{Österreichischer Rundfunk} (MISC) ( \textbf{ORF} (ORG) ), \textbf{Roland Weißmann} (PER), destacou a adequação do local e a proximidade de aeroportos como critérios-chave no processo de selecção da cidade anfitriã para 2026, enquanto a directora de programação da \textbf{ORF} (ORG), \textbf{Stefanie Groiss-Horowitz} (PER), destacou a falta de grandes arenas recentemente construídas, mas incentivou os municípios com planos viáveis a apresentarem candidaturas.

Várias cidades austríacas manifestaram interesse em sediar o concurso de 2026 poucos dias após a vitória de 2025. A 18 de maio de 2025, o presidente da câmara de \textbf{Viena} (LOC) , \textbf{Michael Ludwig} (PER), confirmou a intenção da cidade em apresentar candidatura. No mesmo dia, \textbf{Graz} (LOC) afirmou estar a estudar uma possível candidatura, com a presidente da câmara, \textbf{Elke Kahr} (PER), apontando o \textbf{Stadthalle Graz} (LOC) como um local adequado. O \textbf{Schwarzl Freizeit Zentrum} (ORG), também em \textbf{Graz} (LOC), foi proposto como possível local pelo seu gestor de concertos e operador, \textbf{Klaus Leutgeb} (PER). Também a 18 de maio, \textbf{Insbruque} (LOC) e \textbf{Wels} (PER) confirmaram que apresentariam candidatura com a \textbf{Olympiahalle} (LOC) e uma nova sala de exposições, respetivamente. \textbf{Oberwart} (PER) também manifestou interesse em acolher o concurso. A 19 de maio, o presidente da câmara de \textbf{Sankt Pölten} (PER) , \textbf{Matthias Stadler} (PER), propôs o \textbf{VAZ} (MISC) de \textbf{St} (MISC). \textbf{Pölten} (PER) como possível local. A 26 de maio, \textbf{Ebreichsdorf} (PER) apresentou uma proposta para acolher o concurso num espaço temporário.

A \textbf{ORF} (ORG) lançou o processo de candidatura a 2 de junho de 2025, abrindo um período para que cidades e municípios manifestassem o seu interesse. Esses candidatos receberam os documentos detalhados do concurso e tinham até 4 de julho para submeter as suas propostas. \textbf{Ebreichsdorf} (PER) retirou-se do processo a 15 de junho, seguida por \textbf{Oberwart} (PER) a 21 de junho, \textbf{Graz} (LOC) a 27 de junho e \textbf{Wels} (PER) a 1 de julho. \textbf{Viena} (LOC) e \textbf{Insbruque} (LOC) foram as únicas cidades a apresentar candidaturas até ao prazo limite. A 20 de agosto, a \textbf{UER} (ORG) e a \textbf{ORF} (ORG) anunciaram \textbf{Viena} (LOC) como a cidade anfitriã.



Lista de participantes

A participação no Festival Eurovisão da Canção está condicionada à existência de uma emissora nacional membro ativo da \textbf{União Europeia de Radiodifusão} (ORG), com capacidade para receber o concurso pela rede \textbf{Eurovisão} (LOC) e transmiti-lo em direto em todo o território nacional. A organização dirige convites de participação a todos os membros ativos.

As semifinais realizar-se-ão nos dias 12 e 14 de maio de 2026 pelas 20h00 ( UTC+1 ), no horario de \textbf{Portugal} (LOC) continental . A ordem de atuação dos países foi revelada a 12 de janeiro de 2026.

\textbf{Quinze} (LOC) países competirão nesta semifinal. Esses países, juntamente com a \textbf{Alemanha} (LOC) e a \textbf{Itália} (LOC) , bem como países não participantes incluídos num voto online agregado denominado  \textbf{Resto do Mundo} (MISC) , votarão nesta semifinal.

\textbf{Quinze} (LOC) países competirão nesta semifinal. Esses países, juntamente com a \textbf{Áustria} (LOC) , a \textbf{França} (LOC) e o \textbf{Reino Unido} (LOC) , bem como países não participantes incluídos num voto online agregado denominado  \textbf{Resto do Mundo} (MISC) , votarão nesta semifinal.

A final realizar-se-á a 16 de maio de 2026, às 20:00 (UTC+1) no horario de \textbf{Portugal} (LOC) continental , e contará com 25 países concorrentes: o país anfitrião, a \textbf{Áustria} (LOC), os \textbf{Big Four} (MISC) e as dez canções mais bem classificadas de cada uma das duas semifinais. Todos os 35 países participantes, com júri e televoto, bem como os países não participantes incluídos num voto online agregado denominado \textbf{Resto do Mundo} (MISC), votarão na final.

A participação de \textbf{Israel} (LOC) na \textbf{Eurovisão} (LOC) tornou-se altamente polémica devido à guerra em \textbf{Gaza} (LOC) . As canções de 2024 e 2025 enfrentaram críticas por alegada carga política, e o governo israelita foi acusado de lançar campanhas para influenciar o televoto. Apesar das polémicas, \textbf{Israel} (LOC) teve bons resultados.

Vários canais europeus exigiram à \textbf{União Europeia de Radiodifusão} (ORG) que excluísse \textbf{Israel} (LOC) e ameaçaram abandonar o concurso. A \textbf{UER} (ORG) adiou decisões, cancelou uma assembleia extraordinária por causa de um cessar-fogo e, no final, aprovou mudanças no sistema de votação que permitiram manter \textbf{Israel} (LOC) na competição. Como reação, \textbf{Irlanda} (LOC), \textbf{Países Baixos} (LOC), \textbf{Eslovénia} (LOC) e \textbf{Espanha} (LOC) decidiram não participar em 2026.

\textbf{Israel} (LOC) escolheu o seu representante de 2026 através do programa \textbf{HaKokhav HaBa} (MISC) , mas a canção foi selecionada internamente.

As seguintes emissoras anunciaram a sua retirada, em consequência da autorização dada a \textbf{Israel} (LOC) para participar :



Produção e formato

\textbf{O Festival Eurovisão da Canção 2026} (MISC) será produzido pela emissora nacional austríaca ORF. A equipa principal é composta por \textbf{Michael Krön} (PER) como produtor executivo, \textbf{Stefan Zechner} (PER) como produtor do espetáculo, \textbf{Daniel Hack} (PER) como responsável de produção, \textbf{Christine Tichy} (PER) como responsável técnica, \textbf{Roman Horacek} (PER) como chefe de comunicação, \textbf{Iris Keutter} (PER) como responsável de marketing, \textbf{Oliver Lingens} (PER) como gestor de eventos, \textbf{Christina Lassnig} (PER) como assistente executiva, \textbf{Christina Heinzle-Conrad} (PER) como secretária-geral e \textbf{Martin Szerencsi} (PER) como assessor jurídico. \textbf{Zechner} (LOC), \textbf{Tichy} (LOC), \textbf{Horacek} (LOC), \textbf{Keutter} (LOC), \textbf{Lingens} (PER) e \textbf{Szerencsi} (PER) ocuparam previamente cargos semelhantes ou equivalentes na edição de 2015, realizada em \textbf{Viena} (LOC).

Em junho de 2025, a \textbf{UER} (ORG) anunciou que \textbf{Martin Österdahl} (PER) deixaria o seu cargo de supervisor executivo do concurso, sendo que o diretor do \textbf{Festival Eurovisão da Canção} (MISC), \textbf{Martin Green} (PER), assumiria temporariamente as suas funções.

Um estudo do instituto de investigação \textbf{EcoAustria} (MISC) estima que o orçamento do concurso ascenda a 36 milhões de euros, sendo que o \textbf{Conselho Municipal} (LOC) e o \textbf{Landtag de Viena} (LOC) destinam 22,6 milhões de euros, enquanto a \textbf{UER} (ORG) deverá contribuir com cerca de 5 milhões de euros.

A 18 de agosto de 2025, a \textbf{UER} (ORG) apresentou uma versão renovada do logótipo genérico, concebida pelo estúdio de \textbf{branding Pals} (MISC), sediado em \textbf{Sheffield} (LOC), para assinalar o 70. aniversário do Festival Eurovisão da Canção. Foi também introduzido um novo elemento de design, denominado \textbf{Chameleon Heart} (MISC), composto por 70 camadas do coração da \textbf{Eurovisão} (LOC) em \textbf{3D} (MISC), adaptável às necessidades dos futuros países anfitriões.



\textbf{Transmissão do Festival Todos} (MISC) os organismos de radiodifusão participantes podem optar por ter comentadores no local ou à distância que forneçam informações sobre a votação ao seu público local. Embora sejam obrigados a transmitir a final e a semi-final em que o seu país vota, a maior parte dos organismos de radiodifusão cobre os três programas. Alguns organismos de radiodifusão não participantes também transmitem o concurso. O \textbf{canal do Festival Eurovisão da Canção} (LOC) no \textbf{YouTube} (MISC) sempre oferece transmissões internacionais em direto, sem comentários, de todos os espectáculos.





\vspace{1cm}
\begin{thebibliography}{9}
\bibitem{eurovisao_todos}
    Fonte analisada: \texttt{eurovisao\_todos}. Texto processado automaticamente.
\end{thebibliography}

\end{document}
